% =========================================================================
%  DERECHO CIVIL Y DERECHO PROCESAL CIVIL Y COMERCIAL
%  Documento único e integrado de estudio (apuntes_integrados.tex)
%  Compilar con: pdflatex apuntes_integrados.tex (dos o tres pasadas)
% =========================================================================
\documentclass[12pt,oneside]{book}

% -------------------------------------------------------------------------
% METADATOS
% -------------------------------------------------------------------------
\newcommand{\universidad}{Universidad de Buenos Aires (UBA)}
\newcommand{\facultad}{Facultad de Derecho}
\newcommand{\materia}{Derecho Civil y Derecho Procesal Civil y Comercial}
\newcommand{\titulo}{Derecho Civil y Derecho Procesal Civil y Comercial}
\newcommand{\subtitulo}{Apuntes integrados de estudio}
\newcommand{\autor}{Isaac Edgar Camacho Ocampo}
\newcommand{\anio}{2026}

% -------------------------------------------------------------------------
% PAQUETES
% -------------------------------------------------------------------------
\usepackage[utf8]{inputenc}
\usepackage[T1]{fontenc}
\usepackage[spanish,es-nodecimaldot]{babel}
\usepackage[
    top=2.5cm,
    bottom=2.5cm,
    left=3cm,
    right=2.5cm,
    headheight=15pt
]{geometry}
\usepackage{amsmath}
\usepackage{amssymb}
\usepackage{mathtools}
\usepackage{graphicx}
\usepackage{float}
\usepackage{caption}
\usepackage{subcaption}
\usepackage{booktabs}
\usepackage{array}
\usepackage{longtable}
\usepackage{tabularx}
\usepackage{enumitem}
\usepackage[table]{xcolor}
\usepackage[most]{tcolorbox}
\usepackage{fancyhdr}
\usepackage{ragged2e}
\usepackage[
    colorlinks=true,
    linkcolor=blue!50!black,
    citecolor=green!40!black,
    urlcolor=blue!60!black,
    pdftitle={Derecho Civil y Derecho Procesal Civil y Comercial: apuntes integrados de estudio},
    pdfauthor={Isaac Edgar Camacho Ocampo},
    pdfsubject={Derecho Civil (Introducción y Parte General) y Derecho Procesal Civil y Comercial},
    bookmarks=true
]{hyperref}

% -------------------------------------------------------------------------
% CONFIGURACIÓN GENERAL
% -------------------------------------------------------------------------
\graphicspath{{./img/}{./graficos/}{./figuras/}}
\setcounter{tocdepth}{2}
\setcounter{secnumdepth}{2}

\setlength{\parindent}{0pt}
\setlength{\parskip}{0.5em}
\setlength{\emergencystretch}{2em}
\setlength{\LTpre}{6pt}
\setlength{\LTpost}{6pt}

\setlist[itemize]{topsep=0.3em, itemsep=0.2em}
\setlist[enumerate]{topsep=0.3em, itemsep=0.2em}

\clubpenalty=10000
\widowpenalty=10000

% -------------------------------------------------------------------------
% TÍTULOS DE PARTE Y CAPÍTULO (tamaños moderados) Y TEXTO DE PARTE
% -------------------------------------------------------------------------
\makeatletter
\renewcommand\@makechapterhead[1]{%
  \vspace*{0.5cm}%
  {\parindent\z@ \raggedright \normalfont
   \ifnum \c@secnumdepth >\m@ne
     \large\bfseries\color{black!65} \@chapapp\space\thechapter \par\nobreak\vskip 4\p@
   \fi
   \interlinepenalty\@M \Large\bfseries #1\par\nobreak
   \vskip 6\p@ \hrule height 0.6pt \vskip 22\p@}}
\renewcommand\@makeschapterhead[1]{%
  \vspace*{0.5cm}%
  {\parindent\z@ \raggedright \normalfont
   \interlinepenalty\@M \Large\bfseries #1\par\nobreak
   \vskip 6\p@ \hrule height 0.6pt \vskip 22\p@}}
\newcommand{\partintro}[1]{\gdef\@partintro@text{#1}}
\gdef\@partintro@text{}
\renewcommand\part{%
  \clearpage \thispagestyle{plain}%
  \null\vfil
  \secdef\@part\@spart}
\def\@part[#1]#2{%
  \ifnum \c@secnumdepth >-2\relax
    \refstepcounter{part}%
    \addcontentsline{toc}{part}{\thepart\hspace{1em}#1}%
  \else
    \addcontentsline{toc}{part}{#1}%
  \fi
  \markboth{}{}%
  {\centering \interlinepenalty \@M \normalfont
   \ifnum \c@secnumdepth >-2\relax
     \large\bfseries \partname\nobreakspace\thepart \par\vskip 12\p@
   \fi
   \LARGE\bfseries #2\par}%
  \@endpart}
\def\@endpart{%
  \vskip 2em
  {\centering\begin{minipage}{0.86\textwidth}\normalfont\normalsize\RaggedRight\@partintro@text\end{minipage}\par}%
  \gdef\@partintro@text{}%
  \vfil\newpage}
\makeatother

% -------------------------------------------------------------------------
% ENCABEZADOS Y PIES
% -------------------------------------------------------------------------
\pagestyle{fancy}
\fancyhf{}
\fancyhead[L]{\small\nouppercase{\leftmark}}
\fancyhead[R]{\small Derecho Civil y Procesal}
\fancyfoot[C]{\thepage}
\renewcommand{\headrulewidth}{0.4pt}
\renewcommand{\footrulewidth}{0pt}
\fancypagestyle{plain}{%
  \fancyhf{}
  \fancyfoot[C]{\thepage}
  \renewcommand{\headrulewidth}{0pt}
}

% -------------------------------------------------------------------------
% COMANDOS MATEMÁTICOS
% -------------------------------------------------------------------------
\newcommand{\abs}[1]{\left|#1\right|}
\newcommand{\norm}[1]{\left\lVert#1\right\rVert}
\newcommand{\vect}[1]{\mathbf{#1}}
\newcommand{\deriv}[2]{\frac{d#1}{d#2}}
\newcommand{\pderiv}[2]{\frac{\partial#1}{\partial#2}}

% -------------------------------------------------------------------------
% TIPOS DE COLUMNA Y ENTORNO PARA TABLAS
% -------------------------------------------------------------------------
\newcolumntype{Y}{>{\RaggedRight\arraybackslash}X}
\newcolumntype{B}[1]{>{\bfseries\RaggedRight\arraybackslash}p{#1}}
\newcolumntype{P}[1]{>{\RaggedRight\arraybackslash}p{#1}}
\newenvironment{tabla}{%
    \par\medskip\begingroup\small\renewcommand{\arraystretch}{1.3}\noindent\ignorespaces
}{%
    \par\endgroup\medskip
}

% -------------------------------------------------------------------------
% CUADRO CORNELL (columna izquierda estrecha / columna derecha amplia)
% -------------------------------------------------------------------------
\newcommand{\cornellhead}{%
\toprule
\textbf{Preguntas / palabras clave} & \textbf{Notas / respuestas} \\
\midrule
}
\newcolumntype{Q}{>{\raggedright\arraybackslash}p{0.26\textwidth}}
\newcolumntype{Z}{>{\raggedright\arraybackslash}p{0.68\textwidth}}
\newcolumntype{V}{>{\raggedright\arraybackslash}p{\dimexpr0.94\textwidth+2\tabcolsep+\arrayrulewidth\relax}}
\newcommand{\cornellinicio}{\begingroup\small\renewcommand{\arraystretch}{1.4}}
\newcommand{\cornellfin}{\endgroup}
\newcommand{\cornellpregunta}[2]{\textbf{#1} & #2 \\ \hline}
\newcommand{\cornellresumen}[1]{\rowcolor{blue!5}\multicolumn{2}{|V|}{\textbf{Resumen:} #1} \\ \hline}

\newlength{\cornellL}
\newlength{\cornellR}
\setlength{\cornellL}{0.27\textwidth}
\setlength{\cornellR}{0.63\textwidth}

% -------------------------------------------------------------------------
% CAJAS ACADÉMICAS (uso semántico)
% -------------------------------------------------------------------------
\newtcolorbox{definition}[1]{
    enhanced, breakable,
    title={Definición: #1},
    fonttitle=\bfseries,
    colback=blue!3, colframe=blue!50!black,
    boxrule=0.8pt, arc=2mm
}
\newtcolorbox{important}{
    enhanced, breakable,
    title={Importante},
    fonttitle=\bfseries,
    colback=yellow!8, colframe=orange!70!black,
    boxrule=0.8pt, arc=2mm
}
\newtcolorbox{example}[1]{
    enhanced, breakable,
    title={Ejemplo: #1},
    fonttitle=\bfseries,
    colback=green!4, colframe=green!40!black,
    boxrule=0.8pt, arc=2mm
}
\newtcolorbox{formula}[1]{
    enhanced, breakable,
    title={Fórmula / Regla: #1},
    fonttitle=\bfseries,
    colback=purple!4, colframe=purple!50!black,
    boxrule=0.8pt, arc=2mm
}
\newtcolorbox{exercise}[1]{
    enhanced, breakable,
    title={Ejercicio: #1},
    fonttitle=\bfseries,
    colback=gray!5, colframe=gray!60!black,
    boxrule=0.8pt, arc=2mm
}
\newtcolorbox{note}{
    enhanced, breakable,
    title={Nota},
    fonttitle=\bfseries,
    colback=white, colframe=gray!50,
    boxrule=0.6pt, arc=2mm
}
% Transcripción de disposiciones normativas citadas
\newtcolorbox{norma}[1]{
    enhanced, breakable, frame hidden,
    borderline west={2.5pt}{0pt}{blue!50!black},
    colback=gray!6, boxrule=0pt, arc=0pt,
    left=8pt, right=6pt,
    before upper={\textbf{\upshape #1}\par\smallskip\itshape}
}
\newtcolorbox{articulo}[1]{
    enhanced, breakable, frame hidden,
    borderline west={2pt}{0pt}{gray!60},
    colback=gray!4, boxrule=0pt, sharp corners,
    left=10pt, top=3pt, bottom=3pt,
    before upper={\textbf{\upshape #1.}\ \itshape}
}

% -------------------------------------------------------------------------
% COMPLEMENTOS PARA EL BLOQUE DE DERECHO PROCESAL
% -------------------------------------------------------------------------
\captionsetup{font=small, labelfont=bf}
\newcolumntype{L}[1]{>{\raggedright\arraybackslash}p{#1}}
\newtcolorbox{normativa}[1]{
    enhanced, breakable, frame hidden,
    borderline west={2.5pt}{0pt}{blue!50!black},
    colback=gray!6, boxrule=0pt, arc=0pt,
    left=8pt, right=6pt,
    before upper={\textbf{\upshape #1}\par\smallskip\itshape}
}
\newcommand{\cornelltitle}[1]{\par\bigskip\noindent\textbf{\sffamily Cuadro Cornell: #1}\par\nopagebreak\smallskip}
\newcommand{\cqrow}[2]{\textbf{#1} & #2 \\[0.5em]}
\newcommand{\cornellblock}[1]{%
  \noindent\begin{tabularx}{\textwidth}{@{}>{\raggedright\arraybackslash}p{0.27\textwidth}!{\vrule width 0.4pt}>{\raggedright\arraybackslash}X@{}}
  \toprule
  \textbf{Preguntas / palabras clave} & \textbf{Notas / respuestas} \\
  \midrule
  #1
  \bottomrule
  \end{tabularx}\par}
\newcommand{\cornellblocksum}[2]{%
  \noindent\begin{tabularx}{\textwidth}{@{}>{\raggedright\arraybackslash}p{0.27\textwidth}!{\vrule width 0.4pt}>{\raggedright\arraybackslash}X@{}}
  \toprule
  \textbf{Preguntas / palabras clave} & \textbf{Notas / respuestas} \\
  \midrule
  #1
  \midrule
  \multicolumn{2}{@{}p{\textwidth}@{}}{\textbf{Resumen:} #2} \\
  \bottomrule
  \end{tabularx}\par}


\begin{document}

\begin{titlepage}
\thispagestyle{empty}
\begin{center}
\vspace*{1cm}
{\large \textsc{\universidad}}\\[0.3cm]
{\large \textsc{\facultad}}

\vspace{3cm}
{\Huge \textbf{Derecho Civil}}\\[0.4cm]
{\LARGE y}\\[0.4cm]
{\Huge \textbf{Derecho Procesal Civil y Comercial}}

\vspace{1.2cm}
\rule{0.70\textwidth}{0.5pt}\\[0.5cm]
{\large \subtitulo}\\[0.5cm]
\rule{0.70\textwidth}{0.5pt}

\vfill
{\large \autor}\\[0.3cm]
{\large \anio}
\end{center}
\vfill
\begin{flushleft}
\textit{Este es un modesto aporte a los estudiantes de ``Derecho Civil --- Parte General'' y de ``Derecho Procesal Civil y Comercial''; no pretende sustituir el material oficial de las materias.}
\end{flushleft}
% TODO: verificar consistencia entre fuentes originales. Los archivos fuente indican universidad y año de forma no uniforme (años entre 2020 y 2026; el curso procesal figura como <<Derecho Procesal Civil>> y como <<Derecho Procesal Civil y Comercial>>); se adoptó la versión predominante de la fuente civil.
\end{titlepage}

\frontmatter
\tableofcontents

\chapter{Introducción general}

Este volumen reúne dos bloques de estudio: el \textbf{derecho civil argentino} (introducción y parte general), tomando como eje el Código Civil y Comercial de la Nación (CCyC), y el \textbf{derecho procesal civil y comercial}, que sigue el desarrollo del proceso conforme al Código Procesal Civil y Comercial de la Nación (CPCCN). Los capítulos de cada bloque se agrupan en partes numeradas de forma continua (Partes~I a VII para el derecho civil; Partes~VIII a XII para el derecho procesal).

\section*{Bloque de derecho civil (Partes I a VII)}

La organización de este bloque sigue los elementos de la relación jurídica: el \textbf{sujeto} (quién entabla la relación), el \textbf{objeto} (la materia sobre la que recae) y la \textbf{causa} (el origen del que emergen las relaciones jurídicas), y se completa con la \textbf{extinción} de esas relaciones, es decir, con el modo en que termina aquello que nació.

\begin{table}[H]
\centering
\small
\renewcommand{\arraystretch}{1.25}
\begin{tabularx}{\textwidth}{>{\raggedright\arraybackslash}p{0.17\textwidth} >{\raggedright\arraybackslash}p{0.33\textwidth} >{\raggedright\arraybackslash}X}
\toprule
\textbf{Parte} & \textbf{Cuestión que responde} & \textbf{Contenido} \\
\midrule
I. Fundamentos & ¿De dónde proviene el derecho aplicable y dónde están los límites de su ejercicio? & Fuentes del derecho; ley, costumbre, jurisprudencia y doctrina; abuso del derecho. \\
II. La persona humana & ¿Quién es sujeto de derechos y cómo se protege? & Capacidad y sus restricciones; derechos personalísimos; fin de la existencia; ausencia y presunción de fallecimiento. \\
III. Las personas jurídicas & ¿Qué otros sujetos pueden ser titulares de derechos? & Concepto, régimen, responsabilidad, disolución, fundaciones y asociaciones. \\
IV. Patrimonio, derechos y bienes & ¿Sobre qué recaen las relaciones y con qué responde el deudor? & Patrimonio; clasificación de los derechos; bienes y cosas; garantía patrimonial y acciones de los acreedores. \\
V. Hechos y actos jurídicos & ¿Cómo nacen y qué efectos producen las relaciones jurídicas? & Hechos jurídicos y acto voluntario; acto jurídico y sus elementos; forma e instrumentos; clasificación, modalidades y efectos. \\
VI. Vicios de los actos jurídicos & ¿Qué ocurre cuando el acto está afectado en su voluntad, en su equilibrio o en su sinceridad? & Error, dolo, violencia; lesión; simulación y fraude. \\
VII. Extinción & ¿Cómo terminan las relaciones jurídicas? & Hechos y actos extintivos; prescripción y caducidad. \\
\bottomrule
\end{tabularx}
\caption{Mapa general del libro}
\end{table}

\section*{Bloque de derecho procesal civil y comercial (Partes VIII a XII)}

Este bloque reúne y ordena el desarrollo del proceso civil y comercial argentino desde la producción de la prueba hasta la ejecución de la sentencia y los sistemas cautelares. La lectura sigue la progresión natural del proceso: primero la \textbf{prueba y la conclusión de la causa} (Parte~VIII); luego las formas en que el proceso termina y la \textbf{sentencia} (Parte~IX); después los mecanismos para atacar las decisiones, es decir, la \textbf{impugnación y los recursos} (Parte~X); a continuación la \textbf{ejecución de la sentencia} y el juicio ejecutivo (Parte~XI); y, finalmente, los \textbf{sistemas cautelares} (Parte~XII), que permiten asegurar el resultado del proceso en cualquiera de sus etapas.

\begin{tabularx}{\textwidth}{>{\raggedright\arraybackslash}p{0.17\textwidth} >{\raggedright\arraybackslash}X}
\toprule
\textbf{Parte} & \textbf{Contenido} \\
\midrule
VIII. Prueba y conclusión & Prueba testimonial, confesión, pericial, indiciaria y reconocimiento judicial; alegatos; etapa conclusional; valoración de la prueba. \\
IX. Terminación y sentencia & Modos anormales de terminación; caducidad de la instancia; resoluciones judiciales; cosa juzgada; sentencia civil y sentencia penal. \\
X. Impugnación y recursos & Teoría general; requisitos; doble instancia; aclaratoria y reposición; nulidad; apelación. \\
XI. Ejecución y juicio ejecutivo & Naturaleza y requisitos de la ejecución; procedimiento; condena ilíquida; obligaciones de hacer, no hacer y dar; juicio ejecutivo; sentencia de remate. \\
XII. Sistemas cautelares & Fundamentos y clasificación; presupuestos; sistema asegurativo genérico; tutela anticipada; embargo; secuestro e intervención; anotación de litis; familia. \\
\bottomrule
\end{tabularx}

Cada capítulo del bloque procesal se cierra con uno o más cuadros Cornell, que formulan preguntas de recuperación, sus respuestas y un resumen, para facilitar el estudio activo.

Cada parte se despliega en capítulos que avanzan desde los conceptos fundamentales hacia sus aplicaciones, y la mayoría de ellos se cierra con un cuadro Cornell de repaso que permite recuperar las ideas centrales mediante preguntas. Las cajas de color se reservan para definiciones, reglas, ejemplos y advertencias, y el texto corrido sigue siendo el soporte principal de la exposición.

\mainmatter
\partintro{Antes de estudiar a las personas, los bienes y los actos, esta parte presenta el marco desde el cual se razona jurídicamente: de dónde proviene el derecho aplicable (las fuentes), con qué jerarquía y desde cuándo rige, cómo debe interpretarse y cuáles son los límites al modo de ejercer los derechos.}
\part{Fundamentos: las fuentes del derecho y el ejercicio de los derechos}
\chapter[Concepto y clasificación de las fuentes]{Concepto y clasificación de las fuentes del derecho}

Los seis primeros capítulos de esta parte abordan las fuentes del derecho en el derecho civil argentino, con especial énfasis en el Código Civil y Comercial de la Nación (Ley 26.994, en adelante CCyC). Se estudian la clasificación de las fuentes (formales y materiales), la ley como fuente principal (su concepto, clasificación, jerarquía y vigencia), la costumbre jurídica, la jurisprudencia (incluidos los fallos plenarios) y la doctrina como fuente material.

\begin{note}
El conocimiento de las fuentes del derecho constituye una competencia fundamental para abogados, jueces, escribanos, funcionarios públicos y asesores legales: todo el razonamiento jurídico descansa en saber dónde buscar la norma aplicable, cuál tiene mayor jerarquía y cómo interpretarla cuando el texto no es claro. Por ejemplo, quien redacta un contrato debe saber si una cláusula contradice una norma imperativa (y por tanto es nula) o si simplemente desplaza una norma supletoria (lo cual está permitido). Un juez que resuelve un litigio debe ordenar las fuentes: primero la Constitución y los tratados de derechos humanos, luego las leyes, luego los decretos, y sólo entonces la costumbre o la doctrina como orientación no vinculante. Dominar el tema evita también errores graves: citar jurisprudencia como si fuera obligatoria cuando no lo es, ignorar un fallo plenario que sí lo es, o desconocer que una costumbre mercantil puede integrar válidamente un contrato comercial.
\end{note}

\section{Concepto de fuentes del derecho}

Las fuentes son los \textbf{modos de expresión del derecho}, es decir, las maneras en que el derecho se exterioriza. Su estudio comprende el concepto y la clasificación de las fuentes, las disposiciones generales del Código Civil y Comercial y el desarrollo de la principal de ellas, que es la ley.

\section{Fuentes formales y fuentes materiales}

Es habitual y clásica la distinción entre fuentes formales y fuentes materiales.

\begin{definition}{Fuentes formales}
Son aquellas que son \textbf{obligatorias} y \textbf{vinculan al intérprete}; constituyen el modo de manifestarse externamente del derecho positivo.
\end{definition}

\begin{definition}{Fuentes materiales}
Son aquellas que \textbf{no son obligatorias}. Dan el contenido; persuaden por su razonabilidad; ofrecen razonamientos y argumentos en torno a lo justo, pero no son vinculantes.
\end{definition}

Quien tiene a su cargo interpretar la ley (por ejemplo, un juez) no está obligado a seguir una fuente material. Sin embargo, la fuente material muchas veces \textbf{nutre a las fuentes formales}: la jurisprudencia, que durante mucho tiempo es una fuente material, cuando se consolida puede dar lugar a que se sancione, por ejemplo, una ley que recepta algo que la jurisprudencia ya venía sosteniendo.

\subsection{Fuentes formales en el derecho argentino}

\begin{table}[H]
\centering
\begin{tabularx}{\textwidth}{>{\raggedright\arraybackslash}p{0.27\textwidth} >{\raggedright\arraybackslash}X}
\toprule
\textbf{Fuente formal} & \textbf{Definición} \\
\midrule
La ley en sentido material &
Norma escrita, obligatoria, de carácter general, dictada por quien tiene el cuidado de una sociedad en orden al bien común. Esta norma de carácter general, escrita, promulgada y conocida es una fuente en sentido formal. \\
\addlinespace
La costumbre jurídica &
Conducta repetida de manera generalizada, en forma uniforme, por una porción importante de una comunidad, con la convicción de que se trata de un deber jurídico. No es un mero repetir de un hecho como un uso puramente social: existe una convicción jurídica. \\
\addlinespace
Los fallos plenarios &
Sentencias emanadas de una cámara de apelaciones reunida en pleno, cuando resuelve una cuestión en torno a la cual existen sentencias contradictorias de salas de esa misma cámara. Establece una doctrina que se vuelve obligatoria por disposición del Código Procesal Civil y Comercial de la Nación. \\
\bottomrule
\end{tabularx}
\end{table}

\subsection{Fuentes materiales en el derecho argentino}

\begin{table}[H]
\centering
\begin{tabularx}{\textwidth}{>{\raggedright\arraybackslash}p{0.27\textwidth} >{\raggedright\arraybackslash}X}
\toprule
\textbf{Fuente material} & \textbf{Descripción} \\
\midrule
La jurisprudencia en general & El conjunto de sentencias que dictan los jueces. \\
\addlinespace
La doctrina & La opinión de los autores sobre los distintos temas jurídicos. \\
\addlinespace
El derecho comparado & El recurso al estudio y análisis del ordenamiento jurídico de otros países, de sus leyes y de las decisiones de sus tribunales. \\
\addlinespace
Los principios y valores jurídicos & Enunciados en el artículo 2 del CCyC. \\
\bottomrule
\end{tabularx}
\end{table}

\section{El desborde de las fuentes}

En el momento actual se asiste a un cierto \textbf{desborde de las fuentes del derecho}. Las razones que se señalan son las siguientes:

\begin{itemize}
    \item La relevancia otorgada a la Constitución, a los tratados de derechos humanos y a los principios jurídicos fundamentales hace que muchas veces se invoquen tratados y que la ley pierda relevancia o fuerza.
    \item Esta situación se acompaña de una cierta crisis del positivismo jurídico y de un avance de teorías interpretativas, como el neoconstitucionalismo en sus distintas variantes. Estas teorías sostienen que el juez debe realizar la mejor interpretación posible, mirar el derecho bajo su mejor luz y tratar de descubrir la mejor solución en cada caso, lo cual quita relevancia a la ley.
    \item La existencia de comités de derechos humanos y la abundante jurisprudencia de la Corte Interamericana de Derechos Humanos, que se pronuncia sobre muchísimos temas (opiniones consultivas, sentencias de homologación de acuerdos, de supervisión del cumplimiento de sentencias, de interpretación de sentencias), genera un enorme material y plantea la duda de si debe o no aplicarse.
    \item Aparece el denominado \textit{soft law}: derecho no tan rígido ni tan formalizado, que surge de recomendaciones de un comité o de algún considerando de la Corte Interamericana en un caso determinado, respecto del cual se plantea cómo debe ser aplicado.
\end{itemize}

Todo ello tiene como consecuencia una cierta \textbf{inseguridad jurídica}. Es el contexto en el que se ejerce el derecho en el siglo XXI: un contexto de creciente \textbf{constitucionalización del derecho privado}.

\begin{important}
Para el derecho civil, el Código Civil y Comercial sigue siendo la fuente principal (fuente formal especial), aunque exista una presencia muy fuerte de los tratados de derechos humanos. Esto marca nuevos desafíos para el intérprete.
\end{important}

\section{Cuadro Cornell: concepto y clasificación de las fuentes}

\begin{longtable}{@{}>{\raggedright\arraybackslash}p{0.25\textwidth} >{\raggedright\arraybackslash}p{0.68\textwidth}@{}}
\toprule
\textbf{Preguntas / palabras clave} & \textbf{Notas / respuestas} \\
\midrule
\endfirsthead
\toprule
\textbf{Preguntas / palabras clave} & \textbf{Notas / respuestas} \\
\midrule
\endhead
\bottomrule
\endfoot

\textbf{¿Qué son las fuentes del derecho?} &
Son los modos de expresión del derecho; las maneras en que el derecho se exterioriza. Se clasifican en fuentes formales (obligatorias, vinculantes para el intérprete, como la ley, la costumbre y los fallos plenarios) y fuentes materiales (no vinculantes, pero que orientan por su razonabilidad, como la jurisprudencia general, la doctrina y el derecho comparado). \\[4pt]

\textbf{¿Cuál es la diferencia entre fuente formal y fuente material?} &
La fuente formal obliga al intérprete: si existe una ley aplicable, el juez debe aplicarla. La fuente material persuade pero no obliga: un juez puede considerar la doctrina de un autor reconocido o un fallo de otro tribunal, pero no está obligado a seguirlos. La jurisprudencia, salvo los fallos plenarios, es fuente material en Argentina. \\[4pt]

\textbf{¿Cuáles son las fuentes formales y materiales en el derecho argentino?} &
Formales: la ley en sentido material, la costumbre jurídica y los fallos plenarios. Materiales: la jurisprudencia en general, la doctrina, el derecho comparado y los principios y valores jurídicos (art.~2 CCyC). \\[4pt]

\textbf{¿Qué se entiende por desborde de las fuentes?} &
La relevancia dada a la Constitución, a los tratados de derechos humanos y a los principios jurídicos, junto con la crisis del positivismo, el avance de teorías interpretativas como el neoconstitucionalismo y la abundante producción de comités y de la Corte Interamericana (incluido el \textit{soft law}), hace que la ley pierda relevancia y genera inseguridad jurídica. \\

\midrule
\multicolumn{2}{p{0.93\textwidth}}{%
\textbf{Resumen:} Las fuentes son los modos de expresión del derecho. Las formales (ley en sentido material, costumbre jurídica y fallos plenarios) obligan al intérprete; las materiales (jurisprudencia general, doctrina, derecho comparado y principios y valores jurídicos) persuaden por su razonabilidad pero no obligan, aunque pueden nutrir a las formales. En el contexto actual de desborde de las fuentes y constitucionalización del derecho privado, el CCyC sigue siendo la fuente principal del derecho civil.} \\
\end{longtable}


% ----------------------------------------------------------
\chapter[Arts. 1 y 2 del Código Civil y Comercial]{El Código Civil y Comercial: fuentes, aplicación e interpretación}

\section{Artículo 1: fuentes y aplicación}

El artículo 1, titulado ``Fuentes y aplicación'', marca el tono del tema de las fuentes.

\begin{quote}
\textbf{Art.~1 CCyC --- Fuentes y aplicación.} Los casos que este código rige deben ser resueltos según las leyes que resulten aplicables, conforme con la Constitución Nacional y los tratados de derechos humanos en los que la República sea parte. A tal efecto se tendrá en cuenta la finalidad de la norma. Los usos, prácticas y costumbres son vinculantes cuando las leyes o los interesados se refieren a ellos o en situaciones no regladas legalmente, siempre que no sean contrarios a derecho.
\end{quote}

De la primera parte surge la \textbf{ley como fuente formal}. La segunda parte refiere a la \textbf{costumbre}: los usos, prácticas y costumbres son vinculantes (fuente en sentido formal) cuando las leyes o los interesados se refieren a ellos, o en situaciones no regladas legalmente, siempre que no sean contrarios a derecho.

Respecto de la redacción, la palabra ``casos'' parecería remitir siempre a una conflictividad, cuando a veces puede tratarse de una situación o relación jurídica (que es lo que propiamente trata el derecho civil y comercial) sin que haya un caso en sentido jurídico-técnico propio. Se trata de una disquisición sin mayor relevancia.

\subsection{Tratados que no son de derechos humanos}

El artículo 1 tampoco menciona a los tratados y convenciones que no son de derechos humanos. La materia que rige el CCyC es mucho más amplia que los tratados de derechos humanos, y existen tratados y convenciones que, aunque no sean de derechos humanos, son aplicables a materias civiles y especialmente comerciales. Tienen un rango superior a la ley, por lo que son fuente formal aunque no tengan jerarquía constitucional.

\begin{quote}
\textbf{Art.~2594 CCyC --- Normas aplicables (Derecho Internacional Privado).} Las normas jurídicas aplicables a situaciones vinculadas con varios ordenamientos jurídicos nacionales se determinan por los tratados y las convenciones internacionales vigentes de aplicación en el caso y, en defecto de normas de fuente internacional, se aplican las normas del derecho internacional privado argentino de fuente interna.
\end{quote}

\section{La jurisprudencia en el artículo 1: el debate legislativo}

La comisión redactora del anteproyecto, integrada por quien era presidente de la Corte Suprema, por la magistrada que integró la Corte y por la autora del proyecto, terminó su tarea en marzo de 2012 y elevó el anteproyecto. En esa versión, el artículo 1 incluía a la jurisprudencia como fuente: establecía que los casos debían resolverse teniendo en cuenta la jurisprudencia en consonancia con las circunstancias del caso.
% TODO: contenido incompleto o ambiguo en las notas originales (redacción exacta de la frase suprimida y composición de la comisión)

Esto dio lugar a un debate en la comisión bicameral (el Ejecutivo envía el proyecto al Senado y se conforma una comisión bicameral). En el debate en el Congreso, los legisladores decidieron quitar esa frase.

El fundamento dado por la comisión bicameral en su dictamen final es que se optó por suprimir la referencia a la jurisprudencia porque esta importa un criterio vinculado al análisis del asunto judicial en cuestión, en procura de soluciones jurídicas más adecuadas que se fundamentan en el antecedente judicial: algo más que un criterio de interpretación directa. Básicamente se quiso señalar que \textbf{no se quería dar a la jurisprudencia el carácter de fuente formal}.

\begin{important}
Si la jurisprudencia hubiera figurado en el artículo 1, podría haber sido considerada fuente formal y podría haberse pensado en un pasaje hacia un sistema donde la jurisprudencia fuera obligatoria. La jurisprudencia es una fuente material y no tiene carácter de fuente formal, a diferencia de otros sistemas con un deber mucho más estricto de seguir el precedente, que no está presente en el sistema argentino.
\end{important}

\section{Artículo 2: interpretación de la ley}

El artículo 2 brinda criterios para interpretar la ley, tema importante porque la ley siempre está sujeta a interpretación. Se reconoce como antecedente el artículo 16 del Código Civil de Vélez.

\begin{quote}
\textbf{Art.~2 CCyC --- Interpretación.} La ley debe ser interpretada teniendo en cuenta sus palabras, sus finalidades, las leyes análogas, las disposiciones que surgen de los tratados sobre derechos humanos, los principios y los valores jurídicos, de modo coherente con todo el ordenamiento.
\end{quote}

\begin{table}[H]
\centering
\begin{tabularx}{\textwidth}{>{\raggedright\arraybackslash}p{0.24\textwidth} >{\raggedright\arraybackslash}X}
\toprule
\textbf{Criterio} & \textbf{Desarrollo} \\
\midrule
Palabras de la ley &
Se atiende a qué palabras usa la norma. Por ejemplo, el art.~765 dice ``puede liberarse'': el ``puede'' indica una opción o posibilidad del deudor; no dice que el acreedor deba aceptar. \\
\addlinespace
Finalidad de la norma &
Es un concepto difícil de determinar. Por ejemplo, la finalidad de la ley de alquileres podría ser favorecer una mayor seguridad jurídica y la protección de los vulnerables. \\
\addlinespace
Leyes análogas &
Cuando un tema no puede resolverse por el texto ni por la finalidad de la ley, se recurre a otras leyes. Durante años no hubo normas sobre consentimiento informado para actos médicos en general, pero sí una sobre consentimiento informado para trasplante de órganos; por analogía se aplicaba la ley de trasplante de órganos hasta que en 2009 llegó la Ley de Derechos del Paciente (26.529). \\
\addlinespace
Tratados de derechos humanos &
Criterio importante de interpretación, en base al proceso de constitucionalización y a la idea de que los tratados tienen rango constitucional. \\
\addlinespace
Principios y valores jurídicos &
Ha dado lugar a importantes discusiones. Los sectores positivistas señalan que esto abre demasiado la cuestión interpretativa. Una visión iusnaturalista diría que son principios que emanan de la razón humana y están por encima de la ley positiva; una visión neoconstitucionalista, que emanan del consenso internacional (como el principio de dignidad humana en la Declaración Universal de Derechos Humanos). Hay una coincidencia: al final hay una referencia a principios superiores a la ley positiva. \\
\addlinespace
Coherencia del ordenamiento &
Todo debe hacerse de forma coherente con todo el ordenamiento jurídico, tratando de salvar su coherencia interna. \\
\bottomrule
\end{tabularx}
\end{table}

\section{Cuadro Cornell: artículos 1 y 2 del Código}

\begin{longtable}{@{}>{\raggedright\arraybackslash}p{0.25\textwidth} >{\raggedright\arraybackslash}p{0.68\textwidth}@{}}
\toprule
\textbf{Preguntas / palabras clave} & \textbf{Notas / respuestas} \\
\midrule
\endfirsthead
\toprule
\textbf{Preguntas / palabras clave} & \textbf{Notas / respuestas} \\
\midrule
\endhead
\bottomrule
\endfoot

\textbf{¿Qué establece el art.~1 del CCyC sobre las fuentes?} &
Titulado ``Fuentes y aplicación'', dispone que los casos se resuelven según las leyes aplicables, conforme a la Constitución Nacional y los tratados de derechos humanos. También establece que los usos, prácticas y costumbres son vinculantes cuando las leyes o los interesados se refieren a ellos, o en situaciones no regladas, siempre que no sean contrarios a derecho. \\[4pt]

\textbf{¿Por qué se eliminó la jurisprudencia del art.~1?} &
La comisión bicameral, en su dictamen final, optó por suprimir la referencia a la jurisprudencia para no darle el carácter de fuente formal, evitando la posibilidad de pasar a un sistema de precedente obligatorio. La jurisprudencia es fuente material. \\[4pt]

\textbf{¿Qué ocurre con los tratados que no son de derechos humanos?} &
El art.~1 no los menciona, pero existen tratados y convenciones aplicables a materias civiles y comerciales que tienen rango superior a la ley; son fuente formal aunque no tengan jerarquía constitucional. \\[4pt]

\textbf{¿Cómo se interpreta la ley según el art.~2 CCyC?} &
Teniendo en cuenta: 1) las palabras de la ley; 2) su finalidad; 3) las leyes análogas; 4) las disposiciones de los tratados de derechos humanos; 5) los principios y valores jurídicos; 6) todo ello de modo coherente con el ordenamiento jurídico en su conjunto. \\

\midrule
\multicolumn{2}{p{0.93\textwidth}}{%
\textbf{Resumen:} El art.~1 del CCyC reconoce como fuentes formales la ley (conforme a la Constitución y a los tratados de derechos humanos) y la costumbre (cuando la ley o los interesados se refieren a ella o en situaciones no regladas, si no es contraria a derecho). La jurisprudencia fue suprimida del texto para no otorgarle carácter de fuente formal. El art.~2 establece los criterios de interpretación: palabras, finalidad, leyes análogas, tratados de derechos humanos, principios y valores jurídicos y coherencia con todo el ordenamiento.} \\
\end{longtable}


% ----------------------------------------------------------
\chapter[La ley como fuente formal]{La ley como fuente formal}

\section{Concepto de ley: definición iusnaturalista}

\begin{definition}{Ley (Tomás de Aquino, iusnaturalismo)}
La ley es una \textbf{ordenación de la razón} en orden al \textbf{bien común}, \textbf{promulgada} por quien tiene el cuidado de la comunidad.
\end{definition}

La ley marca y prescribe conducta, regula, establece qué se puede hacer, señala aquello que debe realizarse, prohíbe conductas e indica sanciones para quien quebranta una disposición.

\begin{itemize}
    \item \textbf{Ordenación de la razón:} se espera que la ley exprese fundamentos, se derive de principios lógicos y tenga una lógica interna que mueva a la voluntad. Las personas se mueven según fines y, para ver esos fines, cuentan con la racionalidad, que permite comprender qué es aquello que se debe hacer. La ley puede disponer algo que parezca irrazonable, lo que da lugar a la discusión sobre las leyes injustas; pero en principio se supone que toda ley prescribe algo para el bien común.
    \item \textbf{Bien común:} es la realización del conjunto de condiciones que permiten que las personas, las familias y las asociaciones puedan alcanzar su perfección. Para ello se necesita una autoridad que disponga qué debe hacer cada uno, señale los límites de las conductas y lo debido al otro.
    \item \textbf{Promulgada:} la ley necesita ser conocida; no puede quedar en la mente de quien gobierna o del Congreso, sino que debe ser dada a conocer.
    \item \textbf{Proviene de la autoridad:} quien no tiene autoridad no puede lograr, por ejemplo, que le transfieran un impuesto. La Constitución establece los mecanismos formales por los que se establecen las leyes, los decretos y las distintas normas.
\end{itemize}

\begin{example}{Alimentos en la familia}
En la familia se deben alimentos: los padres hacia los hijos. Es algo que hace al bien común porque resulta bueno que haya un deber de acompañar a los niños en su crianza, lo cual tiene implicancias económicas. La ley prescribe que se paguen alimentos, y eso apunta al bien común: si los niños crecen y se desarrollan, la sociedad progresa, con más personas libres, sanas y capaces de entablar relaciones y buscar el bien entre todos.
\end{example}

\section{Ley en sentido formal y ley en sentido material}

Esta clasificación suele generar muchas confusiones.

\begin{table}[H]
\centering
\begin{tabularx}{\textwidth}{>{\raggedright\arraybackslash}p{0.17\textwidth} >{\raggedright\arraybackslash}X >{\raggedright\arraybackslash}p{0.30\textwidth}}
\toprule
\textbf{Tipo} & \textbf{Concepto} & \textbf{Ejemplo} \\
\midrule
Ley en sentido \textbf{formal} &
Actos que emanan del Poder Legislativo y que han sido elaborados por el procedimiento establecido en la Constitución en el título de la formación y sanción de leyes, que requiere también la participación del Poder Ejecutivo, que promulga o veta. &
El Código Civil y Comercial (Ley N.° 26.994): el Congreso, bajo el procedimiento establecido, emite estos actos, los numera y los publica en el Boletín Oficial. \\
\addlinespace
Ley en sentido \textbf{material} &
Norma escrita, general y obligatoria que emana de la autoridad estatal competente. También es una fuente formal; la clave es que es norma escrita en general. Puede ser un decreto del Poder Ejecutivo o una resolución. El carácter de ley se lo dan la generalidad y la obligatoriedad. &
El Decreto 1089/2012, que reglamenta la Ley 26.529 de Derechos del Paciente: no es ley en sentido formal (no fue aprobada por el Congreso) pero sí en sentido material. \\
\bottomrule
\end{tabularx}
\end{table}

\begin{example}{Derechos del paciente}
Ante un problema como paciente (por ejemplo, en el contexto de la pandemia, por no haber recibido un trato adecuado), se recurre a la Ley de Derechos del Paciente (26.529, ley en sentido formal) y a su decreto reglamentario. Ambas son leyes aplicables al caso: leyes en sentido material.
\end{example}

Generalmente la ley en sentido formal es también ley en sentido material. Muy pocas veces hay leyes en sentido formal que no lo son: a veces se aprueba una pensión para alguien o alguna disposición muy específica que no tiene el alcance general de las leyes en sentido material. Por lo tanto, \textbf{no toda ley en sentido formal es ley en sentido material}.
% TODO: contenido incompleto o ambiguo en las notas originales (la frase "toda ley en sentido material formal puede ser también ley en sentido material" es confusa)

\section{Caracteres de la ley como fuente formal}

\begin{table}[H]
\centering
\begin{tabularx}{\textwidth}{>{\raggedright\arraybackslash}p{0.19\textwidth} >{\raggedright\arraybackslash}X}
\toprule
\textbf{Carácter} & \textbf{Explicación} \\
\midrule
Obligatoriedad &
Es imperativa; se impone a la voluntad porque ha sido establecida por la autoridad que tiene la competencia dada constitucionalmente. El art.~4 CCyC establece: ``Las leyes son obligatorias para todos los que habitan el territorio de la República, sean ciudadanos o extranjeros, residentes, domiciliados o transeúntes, sin perjuicio de lo que dispongan las leyes especiales''. \\
\addlinespace
Generalidad &
Rige a un número indeterminado de personas: a quienquiera que esté comprendido en el ámbito de su aplicación. Esto la distingue de los actos administrativos, donde el Estado decide respecto de un particular concreto (por ejemplo, al adjudicar una licitación). \\
\addlinespace
Origen público &
Emana de la autoridad con competencia en la materia. En sentido formal, del Congreso; en sentido material, de cada autoridad competente (por ejemplo, resoluciones de organismos como la Inspección General de Justicia), siempre dentro del marco de la ley y la Constitución. \\
\addlinespace
Coactividad &
Ante el incumplimiento existen maneras de forzar el cumplimiento. Por ejemplo, en el derecho civil, si alguien comete un ilícito civil y causa un daño, se puede lograr coactivamente que pague, embargando sus bienes y ejecutándolos. Las sanciones son resarcitorias o represivas. \\
\bottomrule
\end{tabularx}
\end{table}

\section{Procedimiento de formación de las leyes}

\begin{table}[H]
\centering
\begin{tabularx}{\textwidth}{>{\centering\arraybackslash}p{0.06\textwidth} >{\raggedright\arraybackslash}p{0.16\textwidth} >{\raggedright\arraybackslash}X}
\toprule
\textbf{N.°} & \textbf{Etapa} & \textbf{Descripción} \\
\midrule
1° & Iniciativa &
Presentar o proponer un proyecto de ley. Están facultados los miembros de cada cámara y el Poder Ejecutivo; también puede haber iniciativa popular si se dan ciertos requisitos que la Constitución establece con una ley reglamentaria. \\
\addlinespace
2° & Discusión &
Debate al interior de cada cámara. Hay cámara de origen y cámara revisora; ciertas materias deben comenzar por una cámara y no por la otra. Dentro de cada cámara hay comisiones de debate, se elaboran dictámenes, los dictámenes conforman el ``orden del día'', este se lleva al plenario y el plenario vota. \\
\addlinespace
3° & Sanción &
Acto por el cual el Poder Legislativo aprueba el proyecto. La cámara de origen aprueba; la otra revisa y, si hace cambios, vuelve a la cámara de origen. Si la cámara revisora aprobó sus cambios con dos terceras partes, la cámara de origen debe insistir con esa mayoría para vencer a la cámara revisora. \\
\addlinespace
4° & Promulgación &
El Poder Ejecutivo tiene la atribución de promulgar u observar (vetar) parcial o totalmente. La promulgación puede ser expresa (mediante un decreto) o tácita (si transcurren diez días sin que el Ejecutivo se pronuncie, se considera promulgada). \\
\addlinespace
5° & Publicación &
Acto por el cual se da a conocer la ley a través del Boletín Oficial, instancia en la que se puede conocer el texto completo. Ha ocurrido que alguien busca la ley en Google y obtiene un texto que no es el correcto (una versión previa o el proyecto originario); la versión vigente debe buscarse en el Boletín Oficial y en los sitios oficiales. \\
\bottomrule
\end{tabularx}
\end{table}

\section{Clasificación de las leyes}

\begin{table}[H]
\centering
\begin{tabularx}{\textwidth}{>{\raggedright\arraybackslash}p{0.17\textwidth} >{\raggedright\arraybackslash}p{0.18\textwidth} >{\raggedright\arraybackslash}X}
\toprule
\textbf{Criterio} & \textbf{Tipo} & \textbf{Descripción y ejemplo} \\
\midrule
Por su contenido dispositivo & Prohibitiva &
Ejemplo: art.~57 CCyC: ``Queda prohibida toda práctica destinada a producir una alteración genética del embrión que se transmita a su descendencia'' (normas de ingeniería genética sobre embriones, para preservar la integridad de la especie humana y evitar actos eugenésicos). \\
\addlinespace
& Dispositiva &
Regula un hecho. Ejemplo: el plazo mínimo de un contrato de locación según la nueva ley de alquileres del año 2020 debe ser de tres años. \\
\midrule
Por la materia & De fondo (sustantivas) &
Regulan materias que hacen a los derechos fundamentales. Según el art.~75 inc.~12 de la Constitución Nacional, dictarlas es competencia del Congreso de la Nación (como el CCyC). Las provincias no podrían dictar sus propios códigos civiles. \\
\addlinespace
& De forma (adjetivas) &
Códigos procesales: son competencia de las provincias. El sistema federal hace que las provincias conserven el poder que no han delegado en el gobierno central. \\
\midrule
Por su imperatividad & Imperativas (indisponibles) &
Obligatorias; se imponen a las partes, que no pueden dejarlas de lado en sus convenciones. Ejemplo: el plazo mínimo del contrato de locación (3 años); no puede pactarse por 1 año aunque las partes quieran. \\
\addlinespace
& Supletorias (disponibles) &
Las partes pueden dejarlas de lado. Rigen por defecto, cuando las partes no dispusieron nada. Ejemplo: si el código permite la sublocación salvo acuerdo en contrario y el contrato nada dice, rige supletoriamente la permisión de sublocar. \\
\bottomrule
\end{tabularx}
\end{table}

\begin{quote}
\textbf{Art.~962 CCyC --- Normas indisponibles y supletorias en contratos.} Las normas legales relativas a los contratos son supletorias por regla general, en materia de contratos. La excepción se da cuando del modo de expresión, de su contenido o de su contexto resulta su carácter indisponible.
\end{quote}

\subsection{Ejemplo problemático: el artículo 765 CCyC}

\begin{quote}
\textbf{Art.~765 CCyC --- Concepto de obligación dineraria (texto parcial).} Si por el acto por el que se ha constituido la obligación se estipuló dar moneda que no sea de curso legal en la República, la obligación debe considerarse como de dar cantidades de cosas y el deudor puede liberarse dando el equivalente en moneda de curso legal.
\end{quote}

La regla según la cual ``el deudor puede liberarse'' ha dado lugar a dos interpretaciones: una parte de la doctrina argentina considera que es una norma \textbf{supletoria} (porque dice ``puede liberarse''); otra parte sostiene que es \textbf{imperativa}.

\begin{example}{Consecuencias prácticas de la calificación del art.~765}
Supóngase que se prestaron dólares y el contrato dice que quien los recibe se obliga a devolver en dólares y renuncia a la potestad de entregar el equivalente en moneda de curso legal. Si se interpreta que la norma es imperativa, esa cláusula no tendría ningún valor y el deudor podría, de todos modos, entregar pesos. Si se la interpreta como supletoria, la cláusula es válida.
\end{example}

La mayoría en las Jornadas de Derecho Civil del año 2015 opinó que esta norma era de carácter supletorio, y esta es la interpretación que ha primado hasta el momento.

\section{Orden jerárquico de las leyes en el derecho argentino}

La norma suprema es la Constitución Nacional. Los tratados internacionales de derechos humanos con jerarquía constitucional están dotados de esa jerarquía por el artículo 75, inciso 22, y se interpretan como complementarios de las normas de la primera parte de la Constitución.

\begin{table}[H]
\centering
\begin{tabularx}{\textwidth}{>{\centering\arraybackslash}p{0.09\textwidth} >{\raggedright\arraybackslash}X >{\raggedright\arraybackslash}p{0.32\textwidth}}
\toprule
\textbf{Rango} & \textbf{Norma} & \textbf{Base constitucional} \\
\midrule
1° & Constitución Nacional y tratados de derechos humanos con jerarquía constitucional & Art.~75 inc.~22 CN: norma suprema. \\
2° & Tratados internacionales y concordatos (no de derechos humanos), con jerarquía supralegal & Art.~75 inc.~22 CN: superiores a las leyes, inferiores a la Constitución. \\
3° & Leyes nacionales & Art.~75 inc.~12 y otros. \\
4° & Decretos del Poder Ejecutivo & Art.~99 CN. \\
5° & Resoluciones y reglamentos de organismos de la administración pública & Marco dado por las leyes. \\
\bottomrule
\end{tabularx}
\end{table}

Esta misma estructura se puede repetir en cada provincia con sus propias constituciones provinciales, leyes y decretos.

\begin{example}{Resolución de la Inspección General de Justicia}
Al constituirse una asociación civil resulta aplicable una resolución de la Inspección General de Justicia (IGJ) en el ámbito de Buenos Aires, que debe ser obedecida. Pero esa resolución no puede contradecir ni el Código Civil ni la Constitución Nacional. Si, en su manera de interpretar este marco, la IGJ los contradijera, estaría ante algo inválido que se podría impugnar.
\end{example}

\section{Cuadro Cornell: la ley como fuente formal}

\begin{longtable}{@{}>{\raggedright\arraybackslash}p{0.25\textwidth} >{\raggedright\arraybackslash}p{0.68\textwidth}@{}}
\toprule
\textbf{Preguntas / palabras clave} & \textbf{Notas / respuestas} \\
\midrule
\endfirsthead
\toprule
\textbf{Preguntas / palabras clave} & \textbf{Notas / respuestas} \\
\midrule
\endhead
\bottomrule
\endfoot

\textbf{¿Cuál es la definición de ley según Tomás de Aquino?} &
Ordenación de la razón en orden al bien común, promulgada por quien tiene el cuidado de la comunidad. Esta definición iusnaturalista destaca que la ley debe ser racional (derivada de principios lógicos), orientada al bien común, proveniente de quien tiene autoridad y promulgada (conocida). \\[4pt]

\textbf{¿Cuál es la diferencia entre ley formal y ley material?} &
Ley en sentido formal: emana del Poder Legislativo siguiendo el procedimiento constitucional (con participación del Poder Ejecutivo, que promulga o veta). Ley en sentido material: cualquier norma escrita, general y obligatoria emanada de autoridad estatal competente, aunque no sea del Congreso (decretos, resoluciones, etc.). Generalmente la ley formal es también material, pero no toda ley formal lo es (por ejemplo, la que aprueba una pensión para una persona determinada). \\[4pt]

\textbf{¿Cuáles son los caracteres de la ley?} &
Obligatoriedad (art.~4 CCyC), generalidad (rige a un número indeterminado de personas), origen público (autoridad competente) y coactividad (existen medios para forzar el cumplimiento; las sanciones son resarcitorias o represivas). \\[4pt]

\textbf{¿Cómo se forma una ley?} &
Iniciativa, discusión, sanción, promulgación (expresa o tácita a los diez días) y publicación en el Boletín Oficial. \\[4pt]

\textbf{¿Qué son las leyes imperativas y supletorias?} &
Las imperativas (indisponibles) se imponen a las partes y no pueden ser dejadas de lado por convención privada (ej.: el plazo mínimo de 3 años en la locación). Las supletorias (disponibles) rigen por defecto cuando las partes no dispusieron nada y pueden ser desplazadas por acuerdo (ej.: la sublocación, permitida salvo pacto en contrario). El art.~962 CCyC establece que en materia de contratos la regla es que las normas son supletorias. \\[4pt]

\textbf{¿Cuál es el orden jerárquico normativo en Argentina?} &
1°) Constitución Nacional y tratados de derechos humanos con jerarquía constitucional (art.~75 inc.~22). 2°) Tratados internacionales sin jerarquía constitucional (supralegales, pero inferiores a la Constitución). 3°) Leyes nacionales. 4°) Decretos del Poder Ejecutivo. 5°) Resoluciones y reglamentos de organismos administrativos. Una resolución que contradiga las normas superiores es inválida y puede impugnarse. \\

\midrule
\multicolumn{2}{p{0.93\textwidth}}{%
\textbf{Resumen:} La ley, definida como ordenación de la razón al bien común promulgada por quien cuida de la comunidad, puede ser formal (emanada del Congreso por el procedimiento constitucional) o material (norma escrita, general y obligatoria de autoridad competente). Sus caracteres son obligatoriedad, generalidad, origen público y coactividad. Se clasifica por su contenido (prohibitiva, dispositiva), por la materia (de fondo, de forma) y por su imperatividad (imperativas y supletorias). La jerarquía normativa va de la Constitución y los tratados de derechos humanos con jerarquía constitucional hasta las resoluciones y reglamentos administrativos.} \\
\end{longtable}


% ----------------------------------------------------------
\chapter[Vigencia de la ley y cómputo de plazos]{Vigencia temporal de la ley y cómputo de plazos}

\section{Vigencia temporal de la ley}

\begin{quote}
\textbf{Art.~5 CCyC --- Vigencia.} Las leyes rigen después del octavo día de su publicación oficial, o desde el día que ellas determinen.
\end{quote}

La ley rige desde el octavo día de su publicación en el Boletín Oficial (desde el día de la publicación se cuentan ocho días y allí comienza su vigencia), salvo que la propia ley determine algo distinto. A veces la ley dice ``esta ley rige desde el día de su publicación'' o ``a los 90 días'' o ``a los 180 días''.

\begin{example}{Entrada en vigencia del Código Civil y Comercial}
El Código Civil y Comercial iba a regir inicialmente a partir del 1° de enero de 2016, pero otra ley dispuso que rigiera a partir del 1° de agosto de 2015. El Código fue sancionado en octubre de 2014 y publicado en el Boletín Oficial el 8 de octubre de 2014; desde octubre hasta agosto no estuvo vigente, porque la ley determinó que entrara en vigencia el 1° de agosto. El legislador puede establecer cuándo rige la ley.
\end{example}

\subsection{Modos de derogación}

\begin{table}[H]
\centering
\begin{tabularx}{\textwidth}{>{\raggedright\arraybackslash}p{0.27\textwidth} >{\raggedright\arraybackslash}X}
\toprule
\textbf{Tipo} & \textbf{Descripción} \\
\midrule
Expresa &
La ley posterior dispone explícitamente la derogación de la anterior. Ejemplo: el CCyC (Ley 26.994) derogó expresamente el Código Civil de Vélez Sársfield (Ley 340). \\
\addlinespace
Tácita por incompatibilidad &
Cuando la ley posterior y la anterior resultan incompatibles entre sí, la nueva ley deroga a la anterior. \\
\addlinespace
Ley especial deroga ley general &
Cuando hay una ley especial en la materia, se aplica la especial. Ejemplo: en materia de sociedades, la ley de sociedades es la ley especial respecto del CCyC (ley general). \\
\addlinespace
Ley general posterior no deroga ley especial &
Salvo que se pueda percibir la voluntad derogatoria del legislador. \\
\addlinespace
Caducidad &
A veces la propia ley establece su duración temporal; o pierde fuerza porque hay una costumbre contraria a ella o porque cambiaron las circunstancias y ya no tiene ese valor. Esto suele pasar muy raramente. \\
\bottomrule
\end{tabularx}
\end{table}

\section{Modo de contar los intervalos del derecho (art.~6 CCyC)}

\begin{quote}
\textbf{Art.~6 CCyC --- Modo de contar los intervalos del derecho.} El modo de contar los intervalos del derecho es el siguiente: día es el intervalo que corre de medianoche a medianoche. En los plazos fijados en días, a contar desde uno determinado, queda éste excluido del cómputo, el cual debe empezar al día siguiente. Los plazos de meses o años se computan de fecha a fecha. Cuando en el mes del vencimiento no hubiere día equivalente al inicial del cómputo, se entiende que el plazo expira el último día de ese mes. Los plazos vencen a la hora veinticuatro del día del vencimiento respectivo. El cómputo civil de los plazos es de días completos y continuos, y no se excluyen los días inhábiles o no laborables. En los plazos fijados en horas, a contar desde una hora determinada, queda ésta excluida del cómputo, el cual debe empezar desde la hora siguiente.
\end{quote}

\begin{table}[H]
\centering
\begin{tabularx}{\textwidth}{>{\raggedright\arraybackslash}p{0.21\textwidth} >{\raggedright\arraybackslash}X >{\raggedright\arraybackslash}p{0.28\textwidth}}
\toprule
\textbf{Tipo de plazo} & \textbf{Regla de cómputo} & \textbf{Ejemplo} \\
\midrule
En días &
Se excluye el día del cómputo; empieza a contarse desde el día siguiente. Un día se toma de medianoche a medianoche. &
Plazo de 15 días desde el 15/08: no cuenta el 15; comienza el 16; vence el 30/08 a las 24 h. \\
\addlinespace
En meses &
Se cuentan de fecha a fecha (mes a mes), no por días exactos. &
3 meses desde el 15/08: vence el 15/11, aunque tengan distintas cantidades de días. \\
\addlinespace
En años &
Se cuentan de fecha a fecha (año a año). &
1 año desde el 15/08/2020 vence el 15/08/2021. \\
\addlinespace
Mes de vencimiento sin día equivalente &
Si en el mes del vencimiento no existe esa fecha, se toma el último día del mes. &
Plazo que vence el 31/02 (inexistente): vence el 28/02 (o 29 en año bisiesto). \\
\addlinespace
Días inhábiles &
El cómputo civil es de días completos y continuos; no se excluyen los días inhábiles o no laborables (admite prueba en contrario por ley o acuerdo de partes). &
A diferencia del cómputo procesal, que sí excluye inhábiles. \\
\bottomrule
\end{tabularx}
\end{table}

\begin{note}
Estas normas son de carácter supletorio, porque las partes pueden disponer un cómputo distinto en sus contratos o una ley especial puede establecer otras reglas.
\end{note}

\section{Cuadro Cornell: vigencia de la ley y cómputo de plazos}

\begin{longtable}{@{}>{\raggedright\arraybackslash}p{0.25\textwidth} >{\raggedright\arraybackslash}p{0.68\textwidth}@{}}
\toprule
\textbf{Preguntas / palabras clave} & \textbf{Notas / respuestas} \\
\midrule
\endfirsthead
\toprule
\textbf{Preguntas / palabras clave} & \textbf{Notas / respuestas} \\
\midrule
\endhead
\bottomrule
\endfoot

\textbf{¿Cuándo entra en vigencia una ley?} &
Desde el octavo día de su publicación en el Boletín Oficial, salvo que la propia ley determine algo distinto (desde el día de su publicación, a los 90 días, a los 180 días, o una fecha determinada). \\[4pt]

\textbf{¿Cómo deja de regir una ley?} &
Por derogación expresa (la nueva ley lo indica), tácita (incompatibilidad entre leyes), o mediante la regla de que la ley especial desplaza a la general; una ley general posterior no deroga a la especial, salvo voluntad derogatoria del legislador. También por caducidad, situación que suele ocurrir muy raramente. \\[4pt]

\textbf{¿Cómo se cuentan los plazos según el art.~6 CCyC?} &
Un día va de medianoche a medianoche. Los plazos en días excluyen el día del cómputo (se empieza a contar desde el día siguiente). Los plazos en meses y años se cuentan de fecha a fecha. Si en el mes de vencimiento no existe el día equivalente, se toma el último día de ese mes. El cómputo civil es de días completos y continuos, sin excluir inhábiles. Estas reglas son supletorias: las partes o la ley especial pueden pactar otra cosa. \\

\midrule
\multicolumn{2}{p{0.93\textwidth}}{%
\textbf{Resumen:} Las leyes rigen desde el octavo día de su publicación oficial o desde la fecha que ellas determinen, y dejan de regir por derogación expresa o tácita, por la primacía de la ley especial o por caducidad. Los plazos civiles se computan de días completos y continuos, sin excluir inhábiles, salvo acuerdo o ley especial en contrario.} \\
\end{longtable}


% ----------------------------------------------------------
\chapter[La costumbre jurídica]{La costumbre como fuente formal}

\section{Concepto y elementos}

\begin{definition}{Costumbre}
Observancia constante y uniforme de un cierto comportamiento por parte de los miembros de una comunidad social, con la convicción de que ese comportamiento responde a una necesidad jurídica.
\end{definition}

La costumbre es una conducta que se repite, generalizada y uniforme, pero no cualquier conducta repetida: existe una \textbf{convicción de que es un deber jurídico}. Este elemento es muy importante, porque distingue la costumbre de un mero uso social.

\begin{table}[H]
\centering
\begin{tabularx}{\textwidth}{>{\raggedright\arraybackslash}p{0.16\textwidth} >{\raggedright\arraybackslash}p{0.20\textwidth} >{\raggedright\arraybackslash}X}
\toprule
\textbf{Elemento} & \textbf{Denominación} & \textbf{Descripción} \\
\midrule
Objetivo & \textit{Corpus} & Repetición de un comportamiento por cierto tiempo razonablemente prolongado, por parte de una comunidad. \\
\addlinespace
Subjetivo & \textit{Animus} (\textit{opinio iuris}) & Convicción de los miembros de la comunidad de la obligatoriedad de esa conducta. Permite distinguir la costumbre de cualquier otro uso social. \\
\bottomrule
\end{tabularx}
\end{table}

\begin{example}{Costumbre y uso social}
Un ejemplo trivial: una persona tiene la costumbre de ir siempre a un mismo café, pero ello no responde a un deber jurídico; no se obligó a consumir en ese café todos los días. Para demostrar que existe una costumbre jurídica hay que demostrar la observancia constante de un comportamiento, pero con la convicción de su obligatoriedad.
% TODO: contenido incompleto o ambiguo en las notas originales (el desarrollo del ejemplo del café y del cobro del doble es confuso)
\end{example}

\section{Importancia histórica de la costumbre}

La costumbre fue durante mucho tiempo la fuente más importante en la producción normativa. Ha influido en institutos típicos del derecho francés, que son la base del ordenamiento argentino:

\begin{table}[H]
\centering
\begin{tabularx}{\textwidth}{>{\raggedright\arraybackslash}p{0.34\textwidth} >{\raggedright\arraybackslash}X}
\toprule
\textbf{Instituto} & \textbf{Descripción} \\
\midrule
El patrimonio como garantía de los acreedores & La expectativa de poder exigir y ejecutar coactivamente el patrimonio (los bienes) del deudor. \\
\addlinespace
La fuerza vinculante de los contratos (\textit{pacta sunt servanda}) & Los pactos deben cumplirse y honrarse entre las partes. \\
\addlinespace
La revocabilidad de las donaciones & Instituto que surgió de la costumbre. \\
\addlinespace
La prescripción desde la exigibilidad & El cómputo de la prescripción comienza el día de la exigibilidad de una obligación. \\
\bottomrule
\end{tabularx}
\end{table}

Todas estas cuestiones surgieron del derecho consuetudinario en Francia, fueron recogidas por el Código Francés y pasaron así a los distintos ordenamientos. Además, la costumbre tiene importancia, y sigue teniéndola, en materia comercial y en muchas materias vinculadas con los contratos.

\section{Tipos de costumbre}

\begin{table}[H]
\centering
\begin{tabularx}{\textwidth}{>{\raggedright\arraybackslash}p{0.20\textwidth} >{\raggedright\arraybackslash}p{0.14\textwidth} >{\raggedright\arraybackslash}X}
\toprule
\textbf{Tipo} & \textbf{Latín} & \textbf{Descripción} \\
\midrule
Costumbre según la ley & \textit{Secundum legem} &
Rige cuando la ley delega en la costumbre la atribución de regular una cierta situación; su autoridad está de alguna manera delegada por la propia ley. Ejemplo: el art.~1255 CCyC sobre el precio de los contratos de obra y servicios, regulado por los usos y costumbres. \\
\addlinespace
Costumbre en ausencia de ley & \textit{Praeter legem} &
Rige en ausencia de disposición legal: en cierta materia nunca hubo una ley que regulara un tema, pero existió una costumbre. Ejemplo: antes de la Ley 18.248 (sobre nombre de personas físicas) existían costumbres que regulaban el nombre. \\
\addlinespace
Costumbre contra la ley & \textit{Contra legem} &
La que contradice una disposición expresa; en principio no es válida. Ejemplo histórico: una norma en Francia que prohibía a las mujeres usar pantalones, vigente hasta 2010; la costumbre ya la había dejado obsoleta aunque no hubiera sido derogada formalmente. \\
\bottomrule
\end{tabularx}
\end{table}

\section{La costumbre en el Código Civil y Comercial}

En el Código Civil de Vélez, el artículo 17 decía que los usos y costumbres no pueden crear derechos sino cuando las leyes se refieren a ellos o en situaciones no regladas legalmente: una fórmula negativa que el nuevo Código modifica. En el Código de Comercio, antes de la unificación, también había una gran importancia de la costumbre; su artículo 2 decía que en las materias en que las convenciones particulares pueden derogar la ley, la tradición (costumbre) autoriza al juez a indagar si la esencia del acto utiliza la costumbre.
% TODO: contenido incompleto o ambiguo en las notas originales (paráfrasis del art. 2 del Código de Comercio)

El CCyC amplía la noción respecto del código anterior: antes se hablaba de ``usos y costumbres''; ahora se agrega la expresión ``prácticas''. La costumbre es vinculante (fuente formal) cuando se dan los siguientes supuestos, conforme a la última parte del art.~1:

\begin{table}[H]
\centering
\begin{tabularx}{\textwidth}{>{\raggedright\arraybackslash}p{0.27\textwidth} >{\raggedright\arraybackslash}X}
\toprule
\textbf{Supuesto} & \textbf{Explicación} \\
\midrule
La ley se refiere a ellos & Costumbre \textit{secundum legem}. Igual que en el código anterior. \\
\addlinespace
Las partes se refieren a ellos & Novedad del CCyC: las partes pueden pactar que su relación jurídica se rija en ciertos aspectos por las costumbres. Tiene especial relevancia en materia comercial. \\
\addlinespace
Situaciones no regladas legalmente & Costumbre \textit{praeter legem}. Ya estaba en el código anterior. \\
\addlinespace
Que no sean contrarios a derecho & Condición agregada por el CCyC (no estaba en el código anterior). Es una fórmula de mayor amplitud que simplemente ``no oponerse a la ley'': la costumbre no puede ser contraria al derecho en su conjunto. \\
\bottomrule
\end{tabularx}
\end{table}

\subsection{Presencia de la costumbre en materia de contratos}

\begin{quote}
\textbf{Art.~964 CCyC --- Integración del contrato.} El contenido del contrato se integra con: a) las normas indisponibles, que se aplican en sustitución de las cláusulas incompatibles con ellas; b) las normas supletorias; c) los usos y prácticas del lugar de celebración, en cuanto sean aplicables porque hayan sido declarados obligatorios por las partes o porque sean ampliamente conocidos y regularmente observados en el ámbito en que se celebra el contrato, siempre que su aplicación sea razonable.
\end{quote}

\begin{quote}
\textbf{Art.~970 CCyC --- Contratos nominados e innominados (parcial).} Los contratos innominados están regidos, en el siguiente orden, por: a) la voluntad de las partes; b) las normas generales sobre contratos y obligaciones; c) los usos y prácticas del lugar de celebración [\ldots].
\end{quote}

\begin{quote}
\textbf{Art.~1063 CCyC --- Significado de las palabras (interpretación).} Las palabras empleadas en el contrato deben entenderse en el sentido que les da el uso general, excepto que tengan un significado específico que surja de la ley, del acuerdo de las partes o de los usos y prácticas del lugar de celebración.
\end{quote}

\begin{table}[H]
\centering
\begin{tabularx}{\textwidth}{>{\raggedright\arraybackslash}p{0.16\textwidth} >{\raggedright\arraybackslash}X >{\raggedright\arraybackslash}p{0.14\textwidth}}
\toprule
\textbf{Operación} & \textbf{En qué consiste} & \textbf{Art.} \\
\midrule
Integrar & Determinar qué cláusulas existen. Si un contrato no dice nada sobre un tema, se recurre a la ley imperativa, luego a la supletoria y luego a los usos y prácticas, que pasan a ser parte del contenido del contrato. & 964 \\
\addlinespace
Interpretar & Determinar el alcance y sentido de una cláusula existente. Uno de los criterios para interpretar las palabras de un contrato son los usos y prácticas del lugar. & 1063 \\
\bottomrule
\end{tabularx}
\end{table}

\section{Cuadro Cornell: la costumbre jurídica}

\begin{longtable}{@{}>{\raggedright\arraybackslash}p{0.25\textwidth} >{\raggedright\arraybackslash}p{0.68\textwidth}@{}}
\toprule
\textbf{Preguntas / palabras clave} & \textbf{Notas / respuestas} \\
\midrule
\endfirsthead
\toprule
\textbf{Preguntas / palabras clave} & \textbf{Notas / respuestas} \\
\midrule
\endhead
\bottomrule
\endfoot

\textbf{¿Cuáles son los elementos de la costumbre jurídica?} &
\textit{Corpus} (elemento objetivo): repetición constante y uniforme de un comportamiento por parte de una comunidad, durante un tiempo razonablemente prolongado. \textit{Animus} u \textit{opinio iuris} (elemento subjetivo): convicción de los miembros de la comunidad de que esa conducta responde a un deber jurídico. Sin el segundo elemento, es un mero uso social, no costumbre jurídica. \\[4pt]

\textbf{¿Cuáles son los tipos de costumbre?} &
\textit{Secundum legem}: la ley remite a la costumbre para regular cierta situación (ej.: precio de los contratos de obra). \textit{Praeter legem}: rige en ausencia de ley. \textit{Contra legem}: contradice una disposición legal expresa; en principio inválida. El CCyC exige que la costumbre no sea contraria a derecho. \\[4pt]

\textbf{¿Cuándo es vinculante la costumbre según el CCyC?} &
Cuando las leyes o los interesados se refieren a ella, o en situaciones no regladas legalmente, siempre que no sea contraria a derecho. Novedades respecto del código anterior: se agrega el término ``prácticas'', la posibilidad de que las partes se refieran a ellas y la condición de no ser contrarios a derecho. \\[4pt]

\textbf{¿Qué diferencia hay entre integrar e interpretar un contrato?} &
Integrar (art.~964) es determinar qué cláusulas existen: ley imperativa, luego supletoria y luego usos y prácticas. Interpretar (art.~1063) es determinar el alcance y sentido de una cláusula existente, para lo cual los usos y prácticas del lugar son un criterio. \\

\midrule
\multicolumn{2}{p{0.93\textwidth}}{%
\textbf{Resumen:} La costumbre jurídica requiere \textit{corpus} y \textit{animus}. Puede ser \textit{secundum}, \textit{praeter} o \textit{contra legem}. El art.~1 del CCyC la reconoce como vinculante cuando la ley o los interesados se refieren a ella o en situaciones no regladas, siempre que no sea contraria a derecho, y tiene presencia en materia de contratos (arts.~964, 970 y 1063).} \\
\end{longtable}


% ----------------------------------------------------------
\chapter[Jurisprudencia, plenarios y doctrina]{Jurisprudencia, fallos plenarios y doctrina}

\section{La jurisprudencia como fuente material}

La jurisprudencia es el conjunto de sentencias que dictan los tribunales de justicia. Es indudablemente una \textbf{fuente material}, porque no es vinculante.

La sentencia es la decisión judicial adoptada ante un caso concreto y, por lo tanto, tiene fuerza vinculante para las partes de ese caso. Por una aspiración de justicia, si ante una situación análoga se resolvió de una manera, existe la expectativa de obtener la misma resolución ante una situación parecida: el sentido de justicia lleva a preguntarse cómo puede ser que ante el mismo caso se resuelva distinto.

En el sistema argentino, sin embargo, la jurisprudencia no obliga: el hecho de que un juez haya dicho algo no obliga al siguiente, que puede resolver distinto. Tiene autoridad (un juez que ve que varios casos se resolvieron de una manera determinada puede considerar que no debe apartarse), pero no está obligado. Si hay una línea jurisprudencial muy uniforme o francamente mayoritaria, es difícil cambiarla, salvo que existan factores fundamentales superiores vinculados con principios fundamentales del ordenamiento jurídico o con la dignidad de la persona.

Como se vio en el capítulo sobre el CCyC, el artículo 1 no menciona la jurisprudencia. En el proyecto, elaborado por una comisión creada por decreto presidencial en febrero de 2011 y que entregó el anteproyecto en marzo de 2012, el artículo 1 la incluía como criterio para interpretar la ley y resolver un caso. Los legisladores de la comisión bicameral quitaron esa frase para evitar que la jurisprudencia pasara a ser fuente formal con obligatoriedad general.

\section{Los fallos plenarios como fuente formal}

Salvo los fallos plenarios, la jurisprudencia en principio no es fuente formal.

\begin{definition}{Fallo plenario}
Sentencia dictada por una cámara de apelaciones en pleno que resuelve una cuestión en torno a la cual existen sentencias contradictorias dictadas por salas del mismo tribunal, y que se convierte en \textbf{obligatoria} para todas las salas y tribunales inferiores de ese fuero, en esa jurisdicción y competencia.
\end{definition}

Para comprenderlo es necesario entender cómo funciona un tribunal. Los tribunales están organizados en instancias; se toma como ejemplo el Código Procesal Civil y Comercial de la Nación, en el ámbito de la Ciudad de Buenos Aires:

\begin{table}[H]
\centering
\begin{tabularx}{\textwidth}{>{\raggedright\arraybackslash}p{0.19\textwidth} >{\raggedright\arraybackslash}p{0.31\textwidth} >{\raggedright\arraybackslash}X}
\toprule
\textbf{Instancia} & \textbf{Órgano} & \textbf{Función} \\
\midrule
Primera instancia & Juzgados de primera instancia (unipersonales) & Tramita la demanda, contestación, producción de prueba, alegatos y sentencia. \\
\addlinespace
Segunda instancia & Cámara de Apelaciones (organizada en salas de 3 miembros cada una: Sala A, B, C\ldots hasta la M) & Revisa por vía de apelación las decisiones de primera instancia. Cada sala resuelve de manera independiente. \\
\addlinespace
Tribunal en pleno & Todas las salas de la Cámara reunidas & Resuelve cuando hay contradicción entre salas sobre una misma cuestión jurídica. El fallo resultante es obligatorio. \\
\bottomrule
\end{tabularx}
\end{table}

\begin{example}{Contradicción entre salas que motiva un fallo plenario}
Siguiendo con el tema del art.~765 (dólares): en una causa toca la Sala E, que sostiene que el art.~765 es supletorio, de modo que puede pactarse que la liberación sea sólo en dólares. La parte queda conforme porque pactó que el deudor no podría liberarse entregando pesos. En otro juicio se apela y la causa va a la Sala C, que entiende que la norma es imperativa y que las partes no pueden pactar que se entreguen siempre dólares: el deudor siempre tiene la potestad de entregar pesos. Hay contradicción: la Sala E dice una cosa y la Sala C otra.

Ante el fallo de la Sala C que perjudica a la parte, esta interpone el \textbf{recurso de inaplicabilidad de la ley}, que solicita a la Sala C que convoque el plenario de todas las salas. Todas las salas (A, B, C, \ldots, M) se reúnen en plenario y resuelven el tema, unificando la interpretación. Esa posición se convierte en obligatoria para todas las salas y para todos los tribunales inferiores de ese fuero, en esa jurisdicción y competencia.
\end{example}

\begin{quote}
\textbf{Arts.~300 y 303 del CPCCN --- Efectos del fallo plenario.} La interpretación de la ley establecida en una sentencia plenaria será obligatoria para la misma Cámara y para los jueces de primera instancia respecto de los cuales aquélla sea tribunal de alzada, sin perjuicio de que los jueces dejen a salvo su opinión personal. Sólo podrá modificarse dicha doctrina por medio de una nueva sentencia plenaria.
\end{quote}

El fundamento de los fallos plenarios es \textbf{unificar la jurisprudencia} y evitar criterios dispares, porque ofenden el sentido de justicia. El recurso se interpone dentro de los diez días de la sentencia definitiva, mencionando el precedente contradictorio y exponiendo los fundamentos; el tribunal resuelve si es procedente. Si se dicta un fallo plenario, este establece la doctrina legal aplicable y sólo puede ser cambiado por otro plenario.

\begin{important}
Un plenario de la Cámara Civil de Buenos Aires obliga a los jueces civiles de la Ciudad de Buenos Aires, pero no puede obligar a las cámaras de Córdoba o de Mendoza, que tendrán sus propios plenarios de acuerdo con sus códigos procesales.
\end{important}

\section{Difusión de la jurisprudencia}

\begin{table}[H]
\centering
\begin{tabularx}{\textwidth}{>{\raggedright\arraybackslash}p{0.32\textwidth} >{\raggedright\arraybackslash}X}
\toprule
\textbf{Fuente de difusión} & \textbf{Descripción} \\
\midrule
Publicaciones oficiales de la CSJN & La Corte Suprema de Justicia de la Nación tiene su sitio oficial (www.csjn.gov.ar) con sus fallos. \\
\addlinespace
Centro de Información Judicial (CIJ) & Publica a nivel nacional las sentencias que los tribunales tienen la obligación de remitir, con los datos debidamente analizados en los casos donde puede haber una cuestión de intimidad. \\
\addlinespace
Sistema Argentino de Información Jurídica (SAIJ) & También recoge mucha jurisprudencia, la publica y la difunde junto con distintos artículos y revistas. \\
\addlinespace
Revistas especializadas & ``La Ley'', ``El Derecho'', ``Jurisprudencia Argentina'' y muchas otras, varias de base online, que difunden la jurisprudencia. \\
\bottomrule
\end{tabularx}
\end{table}

\section{La doctrina como fuente material}

La doctrina es la opinión de los juristas. Es una fuente material y, obviamente, no es obligatoria.

Hubo un momento histórico, en el primer milenio romano, en que ciertos juristas tenían el \textit{ius respondendi ex auctoritate principis}: la autoridad para responder cuestiones y decir el derecho con la autoridad del príncipe. Fue un momento histórico determinado y ya no sería posible.

La doctrina tiene valor material: los autores escriben y publican en las revistas, esos textos llegan a los jueces, que los citan muchas veces al interpretar. Las partes también citan en sus escritos de demanda y contestación a autores que refuercen su postura. Por ejemplo, al interpretar el art.~765, se cita a un autor que dice que la norma es imperativa y a otro que dice que es supletoria; por la autoridad de ese jurista puede convenir citarlo, y así se convence a los jueces. El juez verá si los argumentos son convincentes, porque es una fuente material y por lo tanto no es obligatoria.

\section{Otras fuentes materiales: derecho comparado y principios jurídicos}

\begin{table}[H]
\centering
\begin{tabularx}{\textwidth}{>{\raggedright\arraybackslash}p{0.22\textwidth} >{\raggedright\arraybackslash}X}
\toprule
\textbf{Fuente material} & \textbf{Descripción y alcance} \\
\midrule
Derecho comparado &
Estudio y análisis de la legislación positiva existente en otros países, así como de la interpretación y aplicación que hacen sus propios tribunales. Tiene un valor relativo: que un juez de otro país haya dicho algo no implica que el juez argentino deba seguirlo. A veces tiene más valor si la ley argentina sigue a una ley extranjera (por ejemplo, el error reconocible del CCyC viene del Código italiano, caso en que el derecho comparado tiene más valor interpretativo). No es fuente formal: el juez no está obligado porque en España se resuelva de determinada manera. \\
\addlinespace
Principios y valores jurídicos &
El art.~2 señala que para interpretar la ley hay que tener en cuenta los principios y valores jurídicos. Es un criterio de interpretación que también de alguna manera es fuente material: informa la manera de determinar lo que es justo. Hay debate: los positivistas señalan que abre demasiado la cuestión interpretativa; el iusnaturalismo dice que son principios que emanan de la razón humana y están por encima de la ley positiva; el neoconstitucionalismo dice que emanan del consenso internacional (como la dignidad humana en la Declaración Universal de Derechos Humanos). \\
\bottomrule
\end{tabularx}
\end{table}

\section{Cuadro Cornell: jurisprudencia, plenarios y doctrina}

% TODO: contenido incompleto o ambiguo en las notas originales (el desarrollo cita los arts. 300 y 303 del CPCCN; el resumen original agrega "arts. 288 y ss.")
\begin{longtable}{@{}>{\raggedright\arraybackslash}p{0.25\textwidth} >{\raggedright\arraybackslash}p{0.68\textwidth}@{}}
\toprule
\textbf{Preguntas / palabras clave} & \textbf{Notas / respuestas} \\
\midrule
\endfirsthead
\toprule
\textbf{Preguntas / palabras clave} & \textbf{Notas / respuestas} \\
\midrule
\endhead
\bottomrule
\endfoot

\textbf{¿Qué significa que la jurisprudencia es fuente material?} &
Que los fallos de los jueces orientan pero no obligan a otros jueces. Cada juez puede resolver distinto del anterior, incluso en casos análogos. La sentencia es vinculante sólo para las partes del juicio. La excepción son los fallos plenarios, cuya doctrina sí es obligatoria. \\[4pt]

\textbf{¿Qué es un fallo plenario y por qué es fuente formal?} &
Es una sentencia dictada por una cámara de apelaciones reunida en pleno, convocada mediante el recurso de inaplicabilidad de la ley, cuando existen criterios contradictorios entre salas. Unifica la interpretación y su doctrina es obligatoria para todas las salas y jueces de primera instancia de ese fuero y jurisdicción. Sólo puede ser modificada por otro plenario. Se regula en los arts.~288 y ss. y 300 a 303 del CPCCN. \\[4pt]

\textbf{¿Qué valor tiene la doctrina como fuente del derecho?} &
Es la opinión de los juristas. Es fuente material: no obliga pero puede persuadir. Los jueces pueden citarla y las partes la utilizan en sus escritos para fundar sus posiciones. Su peso depende de la autoridad del jurista y de sus argumentos. El \textit{ius respondendi} fue un momento histórico que ya no sería posible. \\[4pt]

\textbf{¿Qué valor tienen el derecho comparado y los principios jurídicos?} &
El derecho comparado tiene un valor relativo y no es fuente formal; tiene más valor interpretativo cuando la ley argentina sigue a una ley extranjera. Los principios y valores jurídicos (art.~2) informan la manera de determinar lo que es justo y son objeto de debate entre positivistas, iusnaturalistas y neoconstitucionalistas. \\

\midrule
\multicolumn{2}{p{0.93\textwidth}}{%
\textbf{Resumen:} La jurisprudencia es fuente material: no obliga más allá de las partes del caso, salvo los fallos plenarios, que unifican criterios contradictorios entre salas y son obligatorios para las salas y los tribunales inferiores del fuero. La doctrina, el derecho comparado y los principios y valores jurídicos son fuentes materiales que orientan y persuaden, pero no obligan.} \\
\end{longtable}


% ----------------------------------------------------------

\section{Síntesis de las fuentes del derecho}

Las fuentes del derecho son el punto de partida de todo razonamiento jurídico: indican de dónde proviene la norma que debe aplicarse y qué jerarquía tiene frente a otras.

En el derecho argentino, el Código Civil y Comercial (art.~1) reconoce como fuentes formales (vinculantes) la ley (en sentido formal y material), la costumbre jurídica (cuando la ley o las partes se refieren a ella, o en situaciones no reguladas, siempre que no contravenga el derecho) y los fallos plenarios (sentencias de cámaras en pleno que unifican criterios contradictorios entre salas y resultan obligatorias para los tribunales inferiores del fuero). Son fuentes materiales (orientadoras pero no obligatorias) la jurisprudencia general, la doctrina, el derecho comparado y los principios y valores jurídicos (art.~2).

La ley, como fuente principal, puede ser formal (emanada del Congreso siguiendo el procedimiento constitucional) o material (cualquier norma escrita, general y obligatoria emanada de autoridad competente). Se clasifica, además, en imperativa (cuyas disposiciones las partes no pueden derogar) y supletoria (que rige por defecto cuando las partes no acordaron otra cosa). La interpretación de la ley exige integrar sus palabras con su finalidad, las leyes análogas, los tratados de derechos humanos y los principios jurídicos, todo de manera coherente con el ordenamiento.

Las leyes rigen desde el octavo día de su publicación oficial, salvo que establezcan otra fecha, y los plazos se computan de días completos y continuos, sin excluir inhábiles, salvo acuerdo en contrario. Comprender este sistema jerárquico y las reglas de vigencia e interpretación es indispensable para cualquier profesional del derecho: permite saber qué norma aplicar, con qué fuerza, desde cuándo y cómo.
\chapter{El problema conceptual del abuso del derecho}
% ============================================================

\section{Introducción}

El abuso del derecho es uno de los institutos más importantes del Derecho Civil argentino. Su estudio comprende tres cuestiones: qué significa ejercer un derecho de forma ``abusiva'', cuáles son los criterios teóricos para identificarlo y cómo fue incorporado progresivamente al ordenamiento jurídico argentino, desde el Código Civil de Vélez Sársfield hasta el Código Civil y Comercial de 2015.

\begin{important}
El artículo 10 del Código Civil y Comercial se aplica en todas las ramas del derecho: civil, comercial, laboral, administrativo y ambiental. En contratos, sociedades, propiedad horizontal o relaciones laborales, el abuso del derecho es un límite que los jueces aplican cotidianamente para evitar injusticias.
\end{important}

Desde el punto de vista profesional, el abuso del derecho constituye:
\begin{itemize}
    \item una \textbf{herramienta de defensa y ataque}: puede invocarse para proteger a un cliente cuando alguien ejerce un derecho de manera dañina, inequitativa o contraria a la buena fe;
    \item un \textbf{límite a las pretensiones desmedidas}, aplicado por los jueces en distintos ámbitos;
    \item un \textbf{principio transversal}, que rige también las relaciones de vecindad, de consumo, comerciales y familiares.
\end{itemize}

\begin{example}{hilo conductor del capítulo}
Un vecino construyó una parrilla en el espacio común del edificio. Otro copropietario tiene derecho a reclamarle que la retire. Pero si este último hizo algo similar y no le reclama al resto, cabe preguntarse si el derecho que ejerce sigue siendo justo. La pregunta que recorre el tema es: \textbf{¿cuándo el ejercicio de un derecho se vuelve abusivo?}
\end{example}

Los temas se encadenan lógicamente: primero se analiza el problema conceptual, luego los criterios teóricos y, finalmente, cómo el derecho argentino fue respondiendo a este desafío.

\section{El dilema del abuso del derecho}

La expresión \textbf{abuso del derecho} plantea, desde el comienzo, una aparente contradicción lógica: ¿cómo es posible que algo que de por sí es justo --- es decir, que es derecho, algo exigible, una potestad conforme a una cierta regla --- se torne abusivo? Este es el gran dilema de una noción que se presenta como contradictoria.

\subsection{La noción de derecho subjetivo como punto de partida}

La posición que se adopte frente al concepto de \textbf{derecho subjetivo} determina si se admite o no la posibilidad del abuso del derecho: según la noción de derecho subjetivo que se tenga, se tendrá una visión del abuso del derecho, admitiéndolo o no admitiéndolo.

\section{Caso práctico: propiedad horizontal}

Para ilustrar el concepto se presenta un caso concreto resuelto por los tribunales argentinos.

\begin{example}{reclamo selectivo en propiedad horizontal}
Un copropietario tiene la potestad de exigir que otro condómino que construyó algo en violación del reglamento de copropiedad lo derribe. Sin embargo, si todos los copropietarios también construyeron en infracción al reglamento, y quien reclama (a) solo le exige a uno y no al resto, y además (b) él mismo tampoco cumplió con el reglamento, los tribunales consideraron que ese ejercicio del derecho era abusivo. La pretensión fue declarada inadmisible.
\end{example}

El interrogante que plantea el caso es hasta qué punto se tiene derecho a exigir a otro que cumpla con el reglamento de copropiedad cuando quien reclama tampoco lo ha cumplido. Si bien existe un derecho a reclamar en una propiedad horizontal --- por ejemplo, contra el condómino que construyó en exceso, incluso invadiendo zonas comunes ---, en este caso se consideró que quien pedía también había realizado actos que violentaban el reglamento, por lo que su reclamo constituía un abuso del derecho.

\section{Críticas a la noción de abuso del derecho}

La noción de abuso del derecho ha recibido críticas doctrinarias, en particular desde la visión que concibe al derecho subjetivo como una potestad absoluta de la voluntad.

% TODO: contenido incompleto o ambiguo en las notas originales (la cita de Planiol figura como "confiada a la mitad de la voluntad")
\textbf{Planiol} sostuvo que, si el derecho es una potestad confiada a la voluntad, la expresión ``abuso del derecho'' choca contra sí misma: no podría haber abuso, porque la persona siempre estaría ejerciendo algo que de por sí es legítimo. Además, si es la voluntad la que determina el contenido de la noción de derecho, es la propia voluntad la que fija dicho contenido.

Esta posición individualista extrema tuvo su recepción en la redacción original del Código Civil argentino, aunque luego fue modificada con la reforma de 1968 y con el Código Civil y Comercial de 2015.

\section{Terminología técnica correcta}

\begin{important}
La terminología exacta es \textbf{``ejercicio abusivo de un derecho''}, y no ``abuso del derecho''. La pregunta correcta es \emph{cuándo el ejercicio de un derecho se torna abusivo}: el derecho en sí mismo es legítimo; lo abusivo es la manera en que se lo ejerce.
\end{important}

\section{Cuadro Cornell: el problema conceptual del abuso}

\begin{longtable}{
    >{\raggedright\arraybackslash}p{0.27\textwidth}
    >{\raggedright\arraybackslash}p{0.63\textwidth}
}
\toprule
\textbf{Preguntas / palabras clave} & \textbf{Notas / respuestas} \\
\midrule
\endfirsthead
\toprule
\textbf{Preguntas / palabras clave} & \textbf{Notas / respuestas} \\
\midrule
\endhead
\bottomrule
\endfoot

\textbf{¿Qué es el abuso del derecho?} &
Es el ejercicio de un derecho subjetivo legítimo de una manera que se torna abusiva. La terminología exacta es ``ejercicio abusivo de un derecho'', porque el derecho en sí mismo es legítimo; lo abusivo es la forma en que se lo ejerce. \\[0.8em]

\textbf{¿Por qué es una noción contradictoria?} &
Porque si un derecho es algo justo y exigible, resulta paradójico que pueda tornarse abusivo. Este es el gran dilema de la noción. \\[0.8em]

\textbf{¿De qué depende que se admita o no el abuso del derecho?} &
De la noción de derecho subjetivo que se adopte: según ella se tendrá una visión del abuso del derecho, admitiéndolo o no admitiéndolo. \\[0.8em]

\textbf{¿Cómo resolvieron los tribunales el caso del reclamo selectivo en propiedad horizontal?} &
Consideraron abusivo el reclamo de quien exigía a un solo condómino la demolición de lo construido en infracción al reglamento, cuando todos los copropietarios habían construido en infracción y quien reclamaba tampoco lo había cumplido. La pretensión fue declarada inadmisible. \\[0.8em]

\textbf{¿Cuál fue la crítica de Planiol?} &
Que si el derecho es una potestad confiada a la voluntad, la expresión ``abuso del derecho'' choca contra sí misma, porque la persona siempre estaría ejerciendo algo legítimo. \\

\midrule
\multicolumn{2}{p{\dimexpr0.90\textwidth+2\tabcolsep\relax}}{%
\textbf{Resumen:} el abuso del derecho plantea la paradoja de que un derecho legítimo pueda ejercerse de modo abusivo. La posibilidad de admitirlo depende de la noción de derecho subjetivo que se adopte, y fue criticado desde la visión de la voluntad como potestad absoluta. Por eso la expresión técnicamente exacta es ``ejercicio abusivo de un derecho''. Un caso de propiedad horizontal muestra cómo los tribunales rechazaron un reclamo selectivo por considerarlo abusivo.
} \\
\end{longtable}

% ============================================================
\chapter{Criterios teóricos para determinar el abuso}
% ============================================================

La doctrina elaboró criterios para identificar cuándo hay ejercicio abusivo de un derecho. Se distingue entre criterios \textbf{subjetivos} y criterios \textbf{objetivos}, y existe además una postura que los combina.

\section{Criterios subjetivos}

Los criterios subjetivos miran a las condiciones del agente que ejerce el derecho. Son tres:

\begin{table}[H]
\centering
\begin{tabularx}{\textwidth}{
    >{\centering\arraybackslash}p{1cm}
    >{\raggedright\arraybackslash}p{4.2cm}
    >{\raggedright\arraybackslash}X
}
\toprule
\textbf{N.°} & \textbf{Criterio} & \textbf{Descripción} \\
\midrule
1.° & \textbf{Intención de dañar} & El ejercicio es abusivo cuando la finalidad del agente es causar un daño. La intención dañosa es el elemento determinante. \\
\addlinespace
2.° & \textbf{Negligencia en el ejercicio} & El ejercicio es abusivo cuando el derecho se ejerce de forma negligente, sin tomar las diligencias que le son propias. \\
\addlinespace
3.° & \textbf{Sin interés alguno} & El ejercicio es abusivo cuando se realiza sin ningún interés detrás; es decir, sin un interés legítimo que lo justifique. \\
\bottomrule
\end{tabularx}
\caption{Criterios subjetivos de determinación del abuso.}
\end{table}

% TODO: contenido incompleto o ambiguo en las notas originales (en la transcripción oral, el segundo criterio subjetivo figura enunciado como "no es abusivo cuando se ejerce de forma negligente"; el cuadro del apunte y la conclusión indican que la negligencia torna abusivo el ejercicio)

\section{Criterios objetivos}

Los criterios objetivos atienden a circunstancias, elementos o factores objetivos, más allá de las disposiciones subjetivas del agente que ejerce el derecho.

\begin{table}[H]
\centering
\begin{tabularx}{\textwidth}{
    >{\centering\arraybackslash}p{1cm}
    >{\raggedright\arraybackslash}p{4.2cm}
    >{\raggedright\arraybackslash}X
}
\toprule
\textbf{N.°} & \textbf{Criterio} & \textbf{Descripción} \\
\midrule
1.° & \textbf{Fin económico y social} & Será abusivo el ejercicio que va en contra del fin económico y social que tiene ese derecho. \\
\addlinespace
2.° & \textbf{Fin de su institución} & Es abusivo el ejercicio contrario a los fines para los cuales fue instituido el derecho. Este criterio fue receptado por la Ley 17.711. \\
\addlinespace
3.° & \textbf{Buena fe y buenas costumbres} & Criterio de Borda: es abusivo el ejercicio que excede los límites de la moral, la buena fe y las buenas costumbres. Ejemplo: intereses usurarios. \\
\bottomrule
\end{tabularx}
\caption{Criterios objetivos de determinación del abuso.}
\end{table}

% TODO: contenido incompleto o ambiguo en las notas originales (en la transcripción oral, el criterio de Borda se enuncia con la expresión "buena fe y las malas costumbres")
\begin{example}{intereses usurarios}
El caso típico y habitual se da cuando se reclaman intereses que llegan a ser usurarios. La \textbf{usura} es una tasa de interés desmedida: se acepta una tasa de interés razonable, pero si esta se convierte en desmedida, se están excediendo los límites de la moral. Se trata de un exceso respecto de la moral; no de una contrariedad con la buena fe y las buenas costumbres, pues de ser así no sería derecho.
\end{example}

\section{Criterio mixto}

Existe también una postura que combina ambos tipos de criterios: será el juez, en cada caso, quien mediante una ponderación de los elementos objetivos y subjetivos evalúe si el ejercicio del derecho ha sido abusivo.

\section{Cuadro Cornell: criterios para determinar el abuso}

\begin{longtable}{
    >{\raggedright\arraybackslash}p{0.27\textwidth}
    >{\raggedright\arraybackslash}p{0.63\textwidth}
}
\toprule
\textbf{Preguntas / palabras clave} & \textbf{Notas / respuestas} \\
\midrule
\endfirsthead
\toprule
\textbf{Preguntas / palabras clave} & \textbf{Notas / respuestas} \\
\midrule
\endhead
\bottomrule
\endfoot

\textbf{¿Qué miran los criterios subjetivos y cuáles son?} &
Miran las condiciones del agente que ejerce el derecho. Son: 1.° intención de dañar (finalidad de causar un perjuicio); 2.° negligencia en el ejercicio (no se toman las diligencias debidas); 3.° ausencia de todo interés legítimo que justifique el ejercicio. \\[0.8em]

\textbf{¿Qué miran los criterios objetivos y cuáles son?} &
Miran circunstancias o factores objetivos, más allá de las disposiciones subjetivas del agente. Son: 1.° contrariedad con el fin económico y social del derecho; 2.° contrariedad con los fines de su institución (receptado por la Ley 17.711); 3.° exceso sobre los límites de la moral, la buena fe y las buenas costumbres (criterio de Borda; ejemplo: intereses usurarios). \\[0.8em]

\textbf{¿Qué es la usura?} &
Una tasa de interés desmedida. Se acepta una tasa razonable, pero cuando esta se vuelve desmedida se exceden los límites de la moral. \\[0.8em]

\textbf{¿Cuál es el criterio mixto?} &
El juez pondera en cada caso concreto tanto los elementos subjetivos como los objetivos para determinar si el ejercicio fue abusivo. \\

\midrule
\multicolumn{2}{p{\dimexpr0.90\textwidth+2\tabcolsep\relax}}{%
\textbf{Resumen:} los criterios subjetivos atienden a la situación del agente (intención de dañar, negligencia, ausencia de interés); los objetivos, a factores externos a él (fin económico y social, fin de la institución, límites de la moral, la buena fe y las buenas costumbres). El criterio mixto encomienda al juez la ponderación de ambos en cada caso. Los criterios objetivos son los que terminaron siendo receptados por el derecho positivo argentino.
} \\
\end{longtable}

% ============================================================
\chapter{El abuso del derecho en el derecho argentino}
% ============================================================

Este capítulo analiza la evolución histórica del instituto en el ordenamiento jurídico argentino.

\section{El Código Civil de Vélez Sársfield (1871)}

El Código Civil de 1871 está imbuido de una filosofía liberal e individualista propia de la época. Vélez Sársfield entiende el derecho subjetivo como una potestad que el orden jurídico confiere para alcanzar la voluntad y, por lo tanto, \textbf{rechaza la noción de abuso del derecho}.

% TODO: contenido incompleto o ambiguo en las notas originales (el apunte atribuye la redacción originaria a "Decreto-Ley 340/1869")
\begin{formula}{art. 1071, redacción original del Código Civil (1871)}
``El ejercicio de un derecho propio o el cumplimiento de una obligación legal no puede constituirse como ilícito.''

\smallskip
\textit{Redacción originaria según Decreto-Ley 340/1869, en vigencia desde 1871. Esta norma rechazaba expresamente la posibilidad del abuso del derecho: ejercer un derecho nunca podía configurar un acto ilícito.}
\end{formula}

Según esta norma no podía haber abuso del derecho, porque ejercer un derecho no es un ilícito. Vélez Sársfield se aparta así de toda teoría del abuso del derecho y no recepta esa idea en ningún artículo.

\section{La jurisprudencia anterior a la reforma de 1968}

Antes de la Ley 17.711, los tribunales comenzaron a reconocer el instituto por vía pretoriana. A través de distintos fallos --- entre ellos casos de intereses excesivos --- la jurisprudencia entendió que podía haber abuso de derechos. El instituto se mezclaba con otros, pero estaba presente la idea de limitar la visión del derecho subjetivo como algo absoluto.

\section{La Constitución de 1949}

La Constitución de 1949, que luego fue derogada en 1956, constituyó un hito importante en la evolución constitucional del concepto.

\begin{formula}{art. 35, Constitución Nacional de 1949 (derogada en 1956)}
``Los derechos y garantías reconocidos no podrán ser alterados por las leyes que reglamenten su ejercicio; tampoco para amparar a ningún habitante de la nación en perjuicio, detrimento o menoscabo de otros. Los abusos de esos derechos que perjudiquen a la comunidad o lleguen a cualquier forma de explotación del hombre por el hombre configuran delitos que serán castigados por las leyes.''
\end{formula}

Se trata de una norma del campo constitucional, no específicamente del derecho civil, pero con una resonancia indudable en este. Corresponde a un momento de impronta social y de fuerte intervención estatal, que coloca límites a las potestades individuales, sobre todo a la propiedad: esta queda sometida también a una finalidad de bien común y a ciertos límites, de modo que no puede hacerse un uso absoluto del derecho.

\section{El fallo Reina (Corte Suprema, 1956)}

Días antes de la derogación de la Constitución de 1949, la Corte Suprema dictó el fallo \textbf{Reina}, en un caso de abuso del derecho. La Corte sostuvo que, aunque no estuviera en ninguna norma --- previendo que la Constitución de 1949 sería derogada ---, la teoría que ve el ejercicio abusivo del derecho es \textbf{inherente a la noción del derecho}, porque todo derecho es relativo y tiene una finalidad que debe cumplir.

Los derechos encuentran, entonces, un límite, y la teoría del abuso del derecho --- aunque no esté registrada en una norma --- corresponde a la concepción misma del derecho. El fallo buscaba dejar a salvo la institución más allá de la inminente derogación del artículo 35 de la Constitución de 1949.

\section{La Ley 17.711 (1968): reforma del artículo 1071}

En 1968, la Ley 17.711 produjo una reforma integral del Código Civil de Vélez Sársfield, que tocó lugares centrales; el artículo 1071 es uno de ellos.

% TODO: contenido incompleto o ambiguo en las notas originales (la referencia al "artículo referido al 954" figura así en la transcripción)
Esta reforma es parangonada con la modificación, en aquel momento, del artículo referido al 954, que incorporó la imprevisión y la lesión, así como el factor moral y el factor social, como limitación de una visión excesivamente individualista.

\begin{formula}{art. 1071, redacción según Ley 17.711 (1968)}
``El ejercicio regular de un derecho propio o el cumplimiento de una obligación legal no puede constituirse como ilícito. La ley no ampara el ejercicio abusivo de los derechos. Se considerará tal al que contraríe los fines que aquélla tuvo en mira al reconocerlos o al que exceda los límites impuestos por la buena fe, la moral y las buenas costumbres.''
\end{formula}

\subsection{Tres cambios fundamentales}

\begin{table}[H]
\centering
\begin{tabularx}{\textwidth}{
    >{\centering\arraybackslash}p{1cm}
    >{\raggedright\arraybackslash}p{4.2cm}
    >{\raggedright\arraybackslash}X
}
\toprule
\textbf{N.°} & \textbf{Cambio} & \textbf{Explicación} \\
\midrule
1.° & \textbf{Se agrega ``regular''} & El ejercicio de un derecho propio pasa a ser ``el ejercicio regular'' de un derecho propio. Lo regular es conforme a una cierta regla, medida o título. Si el ejercicio es irregular, se acerca al abuso del derecho. \\
\addlinespace
2.° & \textbf{Reconocimiento expreso} & Se agrega la frase ``La ley no ampara el ejercicio abusivo de los derechos''. Es una definición expresa de que no se admitirá el abuso del derecho. \\
\addlinespace
3.° & \textbf{Dos criterios objetivos} & (a) Contrario a los fines que tuvo en mira la ley al reconocer ese derecho (criterio de Josserand). (b) Exceso de los límites impuestos por la buena fe, la moral y las buenas costumbres (criterio de Borda). Se usa ``exceder'' y no ``contrario'' porque, si fuera contrario, habría un problema de compatibilidad con el objeto del acto. \\
\bottomrule
\end{tabularx}
\caption{Cambios introducidos por la Ley 17.711 al artículo 1071.}
\end{table}

\begin{definition}{ejercicio regular}
Es el ejercicio conforme a una cierta regla, medida o título. Si el ejercicio es irregular, se acerca al abuso del derecho.
\end{definition}

La palabra ``regular'' es muy importante: en el nuevo código se la encuentra en muchos otros artículos, como una manera de calificar el ejercicio de los derechos.

\section{El Código Civil y Comercial de la Nación (2015)}

El Código Civil y Comercial del año 2014, con vigencia desde 2015, ratifica la vigencia de la institución del abuso del derecho en el derecho argentino y la amplía.

\begin{formula}{art. 10, Código Civil y Comercial de la Nación (2015)}
``El ejercicio regular de un derecho propio o el cumplimiento de una obligación legal no puede constituirse como ilícito. La ley no ampara el ejercicio abusivo de los derechos. Se considera tal el que contraríe los fines del ordenamiento jurídico o el que exceda los límites impuestos por la buena fe, la moral y las buenas costumbres. El juez debe ordenar lo necesario para evitar los efectos del ejercicio abusivo o de la situación jurídica abusiva y, si correspondiere, procurar la reposición al estado de hecho anterior o fijar una indemnización.''
\end{formula}

\subsection{Tres cambios respecto de la Ley 17.711}

\begin{table}[H]
\centering
\begin{tabularx}{\textwidth}{
    >{\centering\arraybackslash}p{1cm}
    >{\raggedright\arraybackslash}p{4.2cm}
    >{\raggedright\arraybackslash}X
}
\toprule
\textbf{N.°} & \textbf{Cambio} & \textbf{Explicación} \\
\midrule
1.° & \textbf{Ubicación metodológica: Título Preliminar} & Ya no está en la parte de daños, sino en el Título Preliminar (art. 10). Esto le da fuerza expansiva a todo el derecho civil y comercial. \\
\addlinespace
2.° & \textbf{``Fines del ordenamiento jurídico''} & El art. 1071 hablaba de ``los fines que tuvo en mira la ley al reconocer ese derecho''; el art. 10 habla de ``los fines del ordenamiento jurídico''. Es más amplio. \\
\addlinespace
3.° & \textbf{Atribuciones expresas para el juez} & Se explicitan los poderes del juez: debe ordenar lo necesario para evitar los efectos del ejercicio abusivo; procurar la reposición al estado anterior; y fijar una indemnización por los daños si ya ocurrieron. \\
\bottomrule
\end{tabularx}
\caption{Cambios introducidos por el art. 10 del Código Civil y Comercial.}
\end{table}

\section{Otros artículos relacionados del Código Civil y Comercial}

El ejercicio regular y el abuso del derecho aparecen en múltiples artículos del Código, formando un sistema coherente.

\begin{table}[H]
\centering
\begin{tabularx}{\textwidth}{
    >{\raggedright\arraybackslash}p{2.3cm}
    >{\raggedright\arraybackslash}X
}
\toprule
\textbf{Artículo} & \textbf{Contenido / relación con el abuso del derecho} \\
\midrule
Art. 9 & Principio de buena fe. Base axiológica del artículo 10 sobre abuso del derecho. \\
\addlinespace
Art. 10 & Abuso del derecho. Norma principal del instituto, ubicada en el Título Preliminar. \\
\addlinespace
Art. 11 & Abuso de posición dominante. Se aplican los principios de los arts. 9 y 10 a quien abusa de su posición dominante en el mercado. \\
\addlinespace
Art. 14 & Derechos individuales y colectivos. La ley no ampara el ejercicio abusivo de los derechos individuales cuando pueda afectar al ambiente y a los derechos de incidencia colectiva. \\
\addlinespace
Art. 53 & Uso de la expresión ``ejercicio regular'' en materia de derechos de la personalidad. \\
\addlinespace
Art. 235 & Ejercicio regular de extraer aguas y otras disposiciones relativas al dominio público. \\
\addlinespace
Art. 240 & Límites al ejercicio de los derechos individuales sobre bienes en función de los derechos de incidencia colectiva. \\
\addlinespace
Art. 1.718 & Responsabilidad civil. El daño causado por el ejercicio regular de un derecho propio está justificado y no es antijurídico (inciso a). \\
\bottomrule
\end{tabularx}
\caption{Artículos del Código Civil y Comercial relacionados con el ejercicio regular y el abuso del derecho.}
\end{table}

La expresión ``ejercicio regular'' encontró mucho arraigo en el derecho civil y tiene su origen en la reforma de 1968.

\section{Cuadro Cornell: el abuso del derecho en el derecho argentino}

% TODO: contenido incompleto o ambiguo en las notas originales (la referencia al art. 279 / ex art. 953 figura solo en el cuadro Cornell del apunte)
\begin{longtable}{
    >{\raggedright\arraybackslash}p{0.27\textwidth}
    >{\raggedright\arraybackslash}p{0.63\textwidth}
}
\toprule
\textbf{Preguntas / palabras clave} & \textbf{Notas / respuestas} \\
\midrule
\endfirsthead
\toprule
\textbf{Preguntas / palabras clave} & \textbf{Notas / respuestas} \\
\midrule
\endhead
\bottomrule
\endfoot

\textbf{¿Cómo trató el tema el Código de Vélez?} &
El art. 1071 original decía que el ejercicio de un derecho propio no puede constituirse como ilícito. No recepcionó el abuso del derecho: para Vélez Sársfield, ejercer un derecho era siempre legítimo, por su concepción liberal e individualista del derecho subjetivo. \\[0.8em]

\textbf{¿Qué hizo la jurisprudencia antes de 1968?} &
Reconoció el instituto por vía pretoriana en distintos fallos, como casos de intereses excesivos, limitando la visión del derecho subjetivo como algo absoluto. \\[0.8em]

\textbf{¿Qué aportó la Constitución de 1949?} &
Su art. 35 establecía que los abusos de los derechos que perjudiquen a la comunidad o lleguen a cualquier forma de explotación del hombre por el hombre configuran delitos. Aunque fue derogada en 1956, marcó un punto de inflexión constitucional, con límites a las potestades individuales, sobre todo a la propiedad. \\[0.8em]

\textbf{¿Qué dijo la Corte en el fallo Reina (1956)?} &
Que la teoría del ejercicio abusivo del derecho es inherente a la noción del derecho mismo, porque todo derecho es relativo y tiene una finalidad. El instituto existe aunque no esté registrado en una norma. \\[0.8em]

\textbf{¿Cuáles son los 3 cambios de la Ley 17.711 (1968)?} &
1.° Se agrega la palabra ``regular'' al art. 1071. 2.° Se expresa que ``la ley no ampara el ejercicio abusivo de los derechos''. 3.° Se incorporan dos criterios objetivos: contrariedad con los fines que tuvo en mira la ley (Josserand), o exceso de los límites de la buena fe, la moral y las buenas costumbres (Borda). \\[0.8em]

\textbf{¿Por qué se usa ``exceder los límites'' y no ``contrario a''?} &
Porque si el derecho fuera directamente contrario, habría un problema de compatibilidad con el objeto del acto (art. 279 / ex art. 953 CC). ``Exceder los límites'' implica que el derecho en sí es válido, pero su ejercicio sobrepasa las fronteras de lo aceptable. \\[0.8em]

\textbf{¿Cuáles son los 3 cambios del Código Civil y Comercial (2015)?} &
1.° Se ubica en el Título Preliminar (art. 10), con alcance expansivo. 2.° Se amplía el criterio: de ``fines que tuvo la ley'' a ``fines del ordenamiento jurídico''. 3.° Se explicitan las atribuciones del juez: evitar los efectos del ejercicio abusivo, reponer al estado anterior e indemnizar. \\[0.8em]

\textbf{¿Qué otros artículos del CCC se relacionan con el ejercicio regular?} &
Arts. 9 (buena fe), 11 (abuso de posición dominante), 14 (derechos individuales y colectivos), 53, 235, 240 y 1.718 inc. a) (justificación del daño por ejercicio regular). \\

\midrule
\multicolumn{2}{p{\dimexpr0.90\textwidth+2\tabcolsep\relax}}{%
\textbf{Resumen:} el Código de Vélez Sársfield (1871) rechazaba el abuso del derecho al concebir el derecho subjetivo como potestad absoluta de la voluntad. La jurisprudencia lo incorporó por vía pretoriana; la Constitución de 1949 lo elevó al plano constitucional (aunque fue derogada en 1956); el fallo Reina lo consideró inherente a la noción del derecho; la Ley 17.711 (1968) lo positivizó en el art. 1071 con la palabra ``regular'' y dos criterios objetivos; y el Código Civil y Comercial (2015) lo amplió en el art. 10, ubicado en el Título Preliminar, extendiéndolo a los fines del ordenamiento jurídico y dotando al juez de atribuciones expresas para prevenir, reponer e indemnizar.
} \\
\end{longtable}
\partintro{Esta parte recorre la situación jurídica de la persona humana: las restricciones a su capacidad y el sistema de apoyos, los derechos personalísimos que protegen su dignidad, el fin de su existencia por la muerte y las figuras que regulan la incertidumbre sobre su vida (la ausencia y la presunción de fallecimiento).}
\part{El sujeto de la relación jurídica: la persona humana}
\chapter{Marco convencional y antecedentes}
\label{ch:bloque1}

\section{Presentación del tema}

Este capítulo y los tres siguientes reconstruyen el sistema de restricciones a la capacidad de las personas regulado en el Código Civil y Comercial argentino. Se aborda el contexto constitucional que ofrece la Convención sobre los Derechos de las Personas con Discapacidad (CDPD) y su artículo 12, los tres tipos de restricción que regula el código --capacidad restringida, incapacidad e inhabilitación por prodigalidad--, el procedimiento judicial que las rodea --legitimación, medidas cautelares, entrevista personal, dictamen interdisciplinario, sentencia, registro y revisión-- y el sistema de apoyos que reemplaza, como criterio general, a la antigua representación absoluta.

Quien ejerza la abogacía --en un estudio, en un juzgado, en el Ministerio Público o en cualquier área que involucre contratos, sucesiones, salud o familia-- tarde o temprano se encontrará con una persona cuya capacidad está restringida, o con la necesidad de evaluar si conviene promover esa restricción para proteger a alguien. Comprender este sistema permite redactar contratos que no terminen siendo nulos, asesorar correctamente a una familia que enfrenta el deterioro cognitivo de un mayor, intervenir con criterio en un proceso de determinación de la capacidad y, sobre todo, evitar dos errores opuestos: privar de autonomía a quien todavía puede decidir por sí mismo, o desproteger a quien necesita un apoyo real.

Estas reglas tampoco son ajenas a la vida cotidiana de las familias: un padre, una madre, un abuelo o un hermano pueden necesitar, en algún momento, ser acompañados con un apoyo para tomar decisiones. Entender este sistema ayuda a saber qué pedir, qué esperar de un juez y, fundamentalmente, cómo acompañar a esa persona sin sustituir su voluntad.

\begin{note}
El sistema puede organizarse en cuatro bloques que se relacionan entre sí. El primero explica por qué existe una convención internacional que exige no apartar a la persona de la toma de decisiones, sino acompañarla mediante apoyos. El segundo distingue qué tipo de restricción corresponde a cada situación: una restricción limitada para quien conserva gran parte de su autonomía (capacidad restringida), una restricción más intensa para quien no puede en absoluto interaccionar con su entorno (incapacidad), o una limitación específica para quien administra sus bienes de manera perjudicial para su familia (inhabilitación por prodigalidad). El tercero regula el procedimiento judicial que debe seguirse, con las garantías correspondientes, para decidir si corresponde imponer una restricción y de qué alcance. El cuarto explica qué sucede cuando una persona actúa en contra de lo dispuesto por la sentencia, cómo se recupera la plena capacidad y cómo funciona, en concreto, el sistema de apoyos que reemplaza a la antigua representación absoluta.
\end{note}

Al regular la capacidad jurídica se distingue entre capacidad de derecho y capacidad de ejercicio. Dentro de la capacidad de ejercicio, la regla siempre fue la capacidad, pero existen algunos casos de incapaces: las personas por nacer, las personas menores de edad --que van teniendo una capacidad progresiva-- y todo el universo de situaciones en que se presenta alguna discapacidad por razones de salud mental. Es en este último campo donde se ubica el tema desarrollado en esta sección: la Sección Tercera, Capítulo Segundo, Título Primero, Libro Primero del Código Civil y Comercial, que trata las restricciones a la capacidad.

\section{La Convención sobre los Derechos de las Personas con Discapacidad (CDPD)}

Para entender el tema es necesario tener presente el contexto de la Convención sobre los Derechos de las Personas con Discapacidad. Esta convención es del año 2006; la Argentina la ratifica en 2008 mediante la Ley 26.378.

\begin{important}
\textbf{Artículo 75, inciso 22 --- Constitución Nacional.} La jerarquía constitucional de ciertos tratados internacionales de derechos humanos se otorga mediante el procedimiento previsto en esta norma. En 2014, por este procedimiento, la CDPD obtuvo jerarquía constitucional, equiparándose a la primera parte de la Constitución Nacional (declaraciones, derechos y garantías).
\end{important}

En 2014 se le dio jerarquía constitucional: la Convención pasó a tener la misma jerarquía que la primera parte de la Constitución --las declaraciones, derechos y garantías--, incorporándose como uno de los tratados de derechos humanos con jerarquía constitucional.

\subsection{El artículo 12 de la CDPD: igual reconocimiento como persona ante la ley}

El artículo clave de la Convención para este tema es el \textbf{artículo 12}, titulado ``Igual reconocimiento como persona ante la ley''. La Convención parte de reafirmar que las personas con discapacidad son seres humanos que deben gozar de igualdad de condiciones y ser incorporadas e incluidas en la comunidad y en la sociedad, lo que supone superar barreras.

La discapacidad surge, según la Convención, de la interacción entre una deficiencia --que puede ser física, psíquica, sensorial o motriz-- y una barrera que impide la plena incorporación de la persona. Esta idea de barreras que deben eliminarse resulta especialmente clara en la dimensión arquitectónica: las rampas obligatorias en las construcciones o la adaptación de los colectivos con rampa son un modo de superar la barrera que representa, por ejemplo, una escalera.

Junto con la noción de barreras a eliminar, la Convención introduce dos nociones fundamentales: los \textbf{apoyos} y los \textbf{ajustes razonables}.

\begin{definition}{Apoyo (noción convencional)}
El apoyo consiste en brindar la ayuda necesaria para que la persona con discapacidad pueda integrarse por sí misma. Resulta fundamental en materia del ejercicio de la capacidad jurídica.
\end{definition}

\begin{definition}{Ajuste razonable}
Los ajustes razonables son las adaptaciones del procedimiento que permiten la plena integración y participación de la persona. Se encuentran, en el Código Civil y Comercial, en los artículos 35 y 36, referidos al proceso judicial vinculado con la capacidad. Por ejemplo, si se debe tomar una audiencia y la persona es sordomuda, corresponde disponer de un traductor de lenguaje de señas para que pueda entender lo que sucede y participar.
\end{definition}

\begin{important}
\textbf{Artículo 12, inciso 1 --- CDPD.} Las personas con discapacidad tienen derecho al reconocimiento de su personalidad jurídica.
\end{important}

\begin{important}
\textbf{Artículo 12, inciso 2 --- CDPD.} Las personas con discapacidad tienen capacidad jurídica en igualdad de condiciones con los demás en todos los aspectos de la vida.
\end{important}

Sobre el alcance de la expresión ``capacidad jurídica'' en este inciso existen dos posturas: quienes entienden que se refiere solo a la capacidad de derecho (ser titular de derechos) y quienes sostienen que comprende ambas dimensiones --capacidad de derecho y capacidad de ejercicio--. Agustina Palacios y Francisco Bariffi son referentes argentinos en el estudio de la Convención; Palacios sostiene, en un capítulo de libro sobre el tema, que el concepto de capacidad jurídica del artículo 12, inciso 2, comprende tanto la capacidad de derecho como la de ejercicio, es decir, que corresponde partir de un principio de capacidad. El propio Comité de los Derechos de las Personas con Discapacidad, creado por la Convención, sostiene la misma interpretación.

\begin{important}
\textbf{Artículo 12, inciso 3 --- CDPD.} Los Estados que han suscripto esta Convención deben garantizarle a las personas con discapacidad los apoyos necesarios para el ejercicio de su capacidad jurídica.
\end{important}

La palabra ``apoyo'' constituye un concepto clave del sistema del nuevo código. Se trata de un concepto amplio, que involucra un conjunto de medidas: por ejemplo, una persona con discapacidad visual puede utilizar apoyos informáticos o perros guía. En este contexto, la referencia es más específica: apoyos para el ejercicio de la capacidad jurídica, es decir, apoyos para la toma de decisiones.

\begin{example}{Apoyo para la venta de un inmueble}
Una persona con discapacidad mental que es dueña de un inmueble y quiere venderlo puede recibir del juez la designación de un apoyo que la asesore y preste su conformidad, a fin de evitar que alguien se aproveche de la situación y le compre el bien muy por debajo de su valor. La decisión de vender sigue siendo de la persona; el apoyo la asiste en la toma de esa decisión.
\end{example}

\begin{important}
\textbf{Artículo 12, inciso 4 --- CDPD (salvaguardias).} Los Estados deben asegurar que las medidas relativas al ejercicio de la capacidad jurídica respeten los derechos, la voluntad y las preferencias de la persona; que estén libres de conflicto de intereses e influencia indebida; que sean proporcionales y adaptadas a las circunstancias de la persona; que se apliquen en el plazo más corto posible; y que estén sujetas a exámenes periódicos por parte de una autoridad u órgano judicial competente, independiente e imparcial.
\end{important}

Las salvaguardias deben ser adecuadas y efectivas, y la medida que adopte un juez debe ajustarse a la persona concreta. En primer lugar, deben respetarse los derechos, la voluntad y las preferencias de la persona; a diferencia de la Convención de los Derechos del Niño, aquí no aparece la noción de ``interés superior'', sino la de respeto de la voluntad y las preferencias.

\begin{example}{Respeto de la voluntad y las preferencias}
Una persona dueña de un campo prefiere sembrar soja porque le gusta esa tarea. Si el apoyo designado considera más conveniente destinar el campo al ganado por rendir más, de todos modos debe respetar los derechos, la voluntad y las preferencias de la persona titular del campo.
\end{example}

Debe evitarse, además, todo conflicto de intereses o influencia indebida.

\begin{definition}{Conflicto de intereses e influencia indebida}
El conflicto de intereses es un concepto más objetivo: existe, por ejemplo, cuando quien es designado apoyo de una persona es, a la vez, quien le alquila un inmueble que administra o sobre el que debe decidir. La influencia indebida es más difícil de determinar, porque supone cierta manipulación de la voluntad de la persona, como puede ocurrir cuando quien ejerce una influencia decisiva convive con ella.
\end{definition}

\subsection{El Comité de la CDPD y el modelo de apoyo a la decisión}

Una observación general del Comité de Naciones Unidas sobre los Derechos de las Personas con Discapacidad, del año 2014, señala que los Estados deben reemplazar los regímenes basados en la sustitución de decisiones --lo que habitualmente se conoce como representación, en la cual otra persona sucede a la titular en su lugar-- por regímenes basados en el apoyo, de modo que la persona pueda decidir por sí misma. Esto es precisamente lo que instrumenta el nuevo código.

En definitiva, existe un delicado equilibrio entre autonomía y protección: la Convención tiende a enfatizar la autonomía, pero también prevé protecciones, de modo que corresponde articular ambos objetivos.

\section{El régimen del Código Civil derogado: incapaces e inhabilitados}

El código anterior, en su artículo 141, se refería a los incapaces. Eran declaradas incapaces de hecho las personas que, por demencia --también llamadas insanas--, padecieran una enfermedad mental que no les permitía administrar su persona ni sus bienes. Existían también los inhabilitados.

En su configuración posterior a la Ley 17.711, es decir, en la configuración inmediatamente anterior a la entrada en vigencia del nuevo código, el ordenamiento contemplaba dos figuras.

\subsection{Figura 1: los incapaces de hecho declarados judicialmente}

También llamados insanos o interdictos, eran personas declaradas incapaces de hecho por la concurrencia de dos factores: uno biológico (la presencia de una enfermedad mental) y uno jurídico (la ineptitud para administrar su persona o sus bienes). Se les designaba un curador y eran, básicamente, incapaces.

\subsection{Figura 2: los inhabilitados (artículo 152 bis)}

El inhabilitado, en cambio, partía de un principio de capacidad y se le nombraba, por sentencia judicial, un curador con funciones delimitadas por esa sentencia. Se preveían tres supuestos:

\begin{itemize}
\item \textbf{Embriaguez habitual o uso de estupefacientes}, cuando exponían al otorgamiento de actos perjudiciales para el patrimonio o la persona --por ejemplo, alguien con alcoholismo que puede aprovecharse y vender bienes a bajo precio o regalarlos.
\item \textbf{Disminución de facultades} que no llegara al supuesto de incapacidad del artículo 141, pero en la cual el juez estimara que podía producirse un daño a la persona o a su patrimonio.
\item \textbf{Prodigalidad}: quienes, por la prodigalidad en los actos de administración y disposición de sus bienes, exponían a su familia a la pérdida del patrimonio. Este supuesto solo procedía si la persona tenía familia y había dilapidado ya parte de sus bienes; es el caso típico de la persona con ludopatía.
\end{itemize}

\begin{note}
En el sistema del código anterior existían, en síntesis, dos grandes esquemas: el incapaz de hecho, totalmente incapaz y sujeto a curador, y el inhabilitado, capaz pero con restricciones puntuales, también sujeto a curador. En el inhabilitado la regla era la capacidad, con una excepción limitada a los tres supuestos mencionados.
\end{note}

\section{La Ley de Salud Mental 26.657 y el artículo 152 ter}

Dos años después de que Argentina aprobara la Convención, en 2010, se sancionó la Ley de Salud Mental (Ley 26.657), que incorporó al código anterior un artículo 152 ter con tres precisiones para el dictado de la sentencia de incapacidad o inhabilitación.

\begin{important}
\textbf{Artículo 152 ter --- Código Civil (texto incorporado por Ley 26.657).} Para declarar la incapacidad o la inhabilitación se requiere un examen de facultativos conformado por evaluaciones interdisciplinarias. La sentencia no podrá extenderse por más de tres años. La sentencia debe especificar las funciones y actos que se limitan.
\end{important}

Esta reforma constituyó un primer intento --parcial y no del todo coherente con el resto del sistema-- de recoger los principios de la Convención, ya que la figura de la incapacidad y su regla general seguían existiendo. Sin embargo, la exigencia de que la sentencia especificara qué funciones y actos se limitaban anticipaba una idea central del nuevo código: que la regla es la capacidad y que corresponde precisar qué es lo que no puede hacerse. Pese a ello, en la práctica muchos jueces continuaron dictando sentencias que declaraban a la persona ``incapaz para todos los actos de la vida civil''.

\section{Cuadro Cornell --- Bloque 1}

\begin{table}[H]
\centering
\begin{tabularx}{\textwidth}{
    >{\raggedright\arraybackslash}p{0.30\textwidth}
    >{\raggedright\arraybackslash}X
}
\cornellhead

\textbf{¿Por qué la CDPD cambió el paradigma de la capacidad jurídica?} &
Porque, a través de su artículo 12, reconoce la personalidad jurídica de las personas con discapacidad y su capacidad jurídica en igualdad de condiciones, reemplazando el modelo de sustitución de decisiones (representación) por un modelo de apoyo a la toma de decisiones, que promueve la autonomía y respeta la voluntad y las preferencias de la persona. \\

\textbf{¿Qué diferencia hay entre un apoyo y un ajuste razonable?} &
El apoyo es la ayuda que permite a la persona con discapacidad ejercer su capacidad jurídica y tomar decisiones por sí misma. El ajuste razonable es la adaptación del procedimiento judicial (por ejemplo, un intérprete de lengua de señas) que garantiza la accesibilidad y la participación efectiva de la persona en el proceso. \\

\textbf{¿Cómo regulaba la capacidad el Código Civil derogado?} &
Mediante dos figuras: los incapaces de hecho (declarados por un factor biológico y uno jurídico, con curador) y los inhabilitados del artículo 152 bis (embriaguez o toxicomanía, disminución de facultades, o prodigalidad), que partían de la capacidad con restricciones puntuales y también tenían curador. \\

\textbf{¿Qué cambios introdujo la Ley 26.657 antes del nuevo código?} &
Incorporó el artículo 152 ter, que exigió examen interdisciplinario para declarar incapacidad o inhabilitación, fijó un plazo máximo de tres años para la sentencia y obligó a especificar las funciones y actos limitados, anticipando el principio de que la regla es la capacidad. \\

\midrule
\multicolumn{2}{p{\textwidth}}{
\textbf{Resumen:} La CDPD, con jerarquía constitucional desde 2014, impone un modelo social de la discapacidad basado en apoyos y ajustes razonables, en reemplazo de la sustitución de la voluntad. El Código Civil derogado, incluso tras la reforma de la Ley 26.657, mantenía la incapacidad y la inhabilitación como ejes del sistema, con curadores en ambos casos; este es el punto de partida que el nuevo Código Civil y Comercial viene a transformar.
} \\
\bottomrule
\end{tabularx}
\end{table}

% =========================================================================
% CAPÍTULO 2
% =========================================================================
\chapter{Tipos de restricciones a la capacidad en el Código Civil y Comercial}
\label{ch:bloque2}

El nuevo Código Civil y Comercial regula el tema a partir de un concepto general de restricciones a la capacidad y de reglas generales aplicables a todos los supuestos.

\section{Reglas generales del artículo 31}

\begin{important}
\textbf{Artículo 31 --- Código Civil y Comercial (reglas generales).} La capacidad de ejercicio se presume, incluso cuando la persona se encuentra internada. Las limitaciones a la capacidad son de carácter excepcional y se imponen siempre en beneficio de la persona. La intervención estatal tiene carácter interdisciplinario, tanto en el tratamiento como en el proceso judicial. La persona tiene derecho a recibir información a través de medios y tecnologías adecuadas. La persona tiene derecho a participar en el proceso judicial con asistencia letrada, que debe ser proporcionada por el Estado si carece de ella. Deben priorizarse las alternativas terapéuticas menos restrictivas de los derechos y libertades.
\end{important}

La regla general es la capacidad: se presume incluso cuando la persona está internada. Por ello, la internación en un establecimiento psiquiátrico no implica necesariamente una restricción a la capacidad; puede haber internación sin restricción, restricción sin internación, o ambas situaciones combinadas.

Las limitaciones a la capacidad son excepcionales y se imponen en beneficio de la persona. En materia de salud mental se parte de que todas las personas son capaces, y la sentencia judicial cumple un rol central porque delimita con precisión qué actos y funciones se limitan.

La intervención estatal debe ser interdisciplinaria. Las disciplinas que habitualmente intervienen son la psiquiatría, la psicología y el trabajo social: la psiquiatría, por tratarse de una especialidad médica, permite diagnosticar y prescribir tratamiento; la psicología aporta la comprensión de los dinamismos psíquicos y el bienestar integral; el trabajo social aporta información sobre el entorno de la persona, dado que la intensidad de la restricción puede variar según cuente o no con una red de contención familiar y social.

\section{Personas con capacidad restringida (artículo 32, primer párrafo)}

\begin{important}
\textbf{Artículo 32, primer párrafo --- Código Civil y Comercial.} El juez puede restringir la capacidad para determinados actos de una persona mayor de trece años que padece una adicción o alteración mental permanente o prolongada de suficiente gravedad, siempre que estime que del ejercicio de su plena capacidad puede resultar un daño a su persona o a sus bienes.
\end{important}

La persona con capacidad restringida constituye el supuesto base del sistema. Se requiere, en primer lugar, sentencia judicial de restricción de la capacidad, en tanto la Convención establece que toda medida que afecte la capacidad de ejercicio debe ser dictada por una autoridad judicial independiente.

\subsection{Requisitos del artículo 32}

\begin{enumerate}
\item \textbf{Ser mayor de trece años.} Entre los 13 y los 18 años la regla es la incapacidad de ejercicio con capacidad progresiva; la sentencia dictada en esa etapa restringirá aquello vinculado con las atribuciones propias de un adolescente. Antes de los 13 años, el artículo 261 del código presume que la persona carece de discernimiento para los actos lícitos, por lo que no es necesaria una restricción judicial.
\item \textbf{Requisito intrínseco o biológico}: presencia de una adicción o de una alteración mental permanente o prolongada, de suficiente gravedad. El código no precisa el término ``adicción''; la doctrina tiende a interpretarlo a la luz del artículo 4 de la Ley de Salud Mental, que define la adicción como el consumo problemático de drogas, legales o ilegales. Una afición intrascendente --como cierta afición a una serie de televisión-- carece de la gravedad necesaria para justificar la limitación de la capacidad. La ``alteración mental permanente o prolongada'', también de suficiente gravedad, refiere a una situación en la que la persona está afectada en su funcionamiento o en sus capacidades mentales, a los fines de un proceso judicial civil, y no debe confundirse con el certificado de discapacidad, que es un concepto distinto.
\item \textbf{Requisito extrínseco}: que del ejercicio de la plena capacidad pueda resultar un daño a la persona o a sus bienes. El código no contempla el daño a terceros; si existiera riesgo de daño a terceros, correspondería otra regulación --la de la responsabilidad civil-- y, eventualmente, un supuesto de internación si el riesgo fuera cierto e inminente.
\end{enumerate}

\subsection{Régimen de las personas con capacidad restringida}

Las personas con capacidad restringida son, como regla, capaces; la sentencia limita únicamente los actos y funciones específicamente indicados, debe ser proporcional y ajustada a la situación concreta de la persona. Se les designa un \textbf{apoyo} --y no un curador--, cuya modalidad de actuación fija el juez: puede tratarse de un apoyo cuyo consentimiento deba concurrir junto con el de la persona, o de un apoyo de consejo. El artículo 101, inciso b, admite también el apoyo con funciones de representación, supuesto que la doctrina considera excepcional y que debe estar especialmente fundado, para no convertirse en una curatela con otro nombre.

\begin{example}{Apoyo con función de asistencia en la venta de un inmueble}
Si a una persona con una adicción se le designa un apoyo con funciones de asistencia para la venta de su inmueble, al momento de celebrarse la escritura concurren el escribano, la persona con capacidad restringida y el apoyo, y ambos firman en el ámbito de su respectiva intervención.
\end{example}

\section{Personas incapaces (artículo 32, último párrafo)}

\begin{important}
\textbf{Artículo 32, último párrafo --- Código Civil y Comercial.} Excepcionalmente, cuando la persona se encuentre absolutamente imposibilitada de interaccionar con su entorno y expresar su voluntad por cualquier modo, medio o formato adecuado, y el sistema de apoyos resulte ineficaz, el juez puede declarar la incapacidad y designar un curador.
\end{important}

La incapacidad constituye un supuesto excepcional, no la regla. Un análisis de las sentencias dictadas por la Cámara Nacional de Apelaciones en lo Civil durante los dos primeros años de vigencia del código mostró que los casos de capacidad restringida eran ampliamente mayoritarios frente a los de incapacidad, que resultaron poco frecuentes.

% TODO: contenido incompleto o ambiguo en las notas originales: no se identifican en el apunte los datos completos (autores, publicación) de la investigación referida sobre la Cámara Nacional de Apelaciones en lo Civil.

Se trata de casos muy severos: la persona no puede en absoluto entablar una relación con su entorno ni expresar su voluntad por ningún modo, medio o formato adecuado, y además el sistema de apoyos resulta ineficaz para asistirla. De no designarse un representante en estos supuestos extremos, la persona quedaría, como un ``muerto civil'', sin posibilidad de que se protejan sus derechos --por ejemplo, para participar en la asamblea de una sociedad anónima o para administrar bienes en caso de encontrarse en coma.

\subsection{Régimen de las personas incapaces}

La sentencia también debe establecer qué actos y funciones se limitan, aunque la limitación será mayor que en la capacidad restringida. Se designa un \textbf{curador} con funciones de representación, a diferencia del supuesto anterior, en el que se designaban apoyos. En el nuevo código, la figura del curador queda restringida a este supuesto excepcional del artículo 32, último párrafo.

\section{Inhabilitados por prodigalidad (artículo 48)}

\begin{important}
\textbf{Artículo 48 --- Código Civil y Comercial.} Pueden ser inhabilitados quienes por la prodigalidad en la gestión de sus bienes expongan a su cónyuge, conviviente o hijos menores de edad o con discapacidad a la pérdida del patrimonio. A estos inhabilitados se les designarán apoyos. Solo podrán otorgar los actos de disposición entre vivos previstos en la sentencia con la conformidad del apoyo. De ser necesario, el juez puede ampliar la nómina de actos que requieren la conformidad del apoyo.
\end{important}

La denominación de ``inhabilitados'' proviene del código anterior, pero en el nuevo código queda restringida exclusivamente al supuesto de prodigalidad; los otros dos supuestos que antes contemplaba el artículo 152 bis --embriaguez o toxicomanía, y disminución de facultades-- pasan a encuadrarse en la capacidad restringida.

\begin{definition}{Pródigo}
Pródigo es quien derrocha sus bienes --no debe confundirse con ``prodigio'', que designa a alguien extraordinario en sus capacidades--. La figura no protege solo a la persona pródiga, sino también a su familia frente al riesgo de pérdida del patrimonio.
\end{definition}

A diferencia del régimen anterior, que exigía la dilapidación efectiva de una parte del patrimonio, el artículo 48 no requiere que la dilapidación ya se haya producido: basta con que se ponga en riesgo la posibilidad de dilapidar el patrimonio.

\subsection{Régimen de los inhabilitados}

Las personas inhabilitadas son, como regla, capaces, pero tienen limitada la posibilidad de realizar actos de disposición de bienes --aquellos que impliquen enajenarlos-- y los demás actos que fije la sentencia. Se les designan apoyos.

\begin{example}{El pródigo y la protección del patrimonio familiar}
Ante una persona que dilapida bienes jugando en el casino, la familia puede iniciar la acción de inhabilitación; la sentencia puede disponer, por ejemplo, que para vender su automóvil o sus bienes registrables, o para hipotecar, necesite el consentimiento de su cónyuge. De ese modo se preserva el patrimonio familiar. Si la persona no tuviera familia, no podría ser declarada inhabilitada en estos términos.
\end{example}

\section{Cuadro Cornell --- Bloque 2}

\begin{table}[H]
\centering
\begin{tabularx}{\textwidth}{
    >{\raggedright\arraybackslash}p{0.30\textwidth}
    >{\raggedright\arraybackslash}X
}
\cornellhead

\textbf{¿Cuándo puede un juez restringir la capacidad de una persona?} &
Cuando la persona es mayor de trece años, padece una adicción o una alteración mental permanente o prolongada de suficiente gravedad, y del ejercicio de su plena capacidad puede resultar un daño a su persona o a sus bienes (art. 32, primer párrafo). \\

\textbf{¿Por qué la incapacidad es hoy un supuesto excepcional?} &
Porque solo procede cuando la persona está absolutamente imposibilitada de interaccionar con su entorno, no puede expresar su voluntad por ningún modo, medio o formato adecuado, y el sistema de apoyos resulta ineficaz; en la práctica, los casos de capacidad restringida resultan mucho más frecuentes que los de incapacidad. \\

\textbf{¿Qué protege realmente la inhabilitación por prodigalidad?} &
Protege al cónyuge, conviviente o hijos menores o con discapacidad frente al riesgo de pérdida del patrimonio familiar por la mala gestión de los bienes de la persona pródiga, limitando sus actos de disposición y exigiendo la conformidad de un apoyo. \\

\midrule
\multicolumn{2}{p{\textwidth}}{
\textbf{Resumen:} El Código Civil y Comercial distingue tres figuras según la intensidad de la limitación: la capacidad restringida (la regla, con apoyos), la incapacidad (excepcional, con curador) y la inhabilitación por prodigalidad (protección patrimonial familiar, con apoyos). En todos los casos rige el principio general de capacidad del artículo 31.
} \\
\bottomrule
\end{tabularx}
\end{table}

% =========================================================================
% CAPÍTULO 3
% =========================================================================
\chapter{Aspectos procesales}
\label{ch:bloque3}

El Código Civil y Comercial incorpora normas de tipo procesal. Aunque se trata de un código de fondo --y los códigos de procedimiento corresponden a las provincias y a la Nación, cada una para su propio fuero--, el código incluye normas procesales que apuntan a dar una regulación común a una materia que afecta directamente la capacidad de una persona.

\section{Juez competente y legitimados activos}

\subsection{Juez competente (artículo 36)}

El juez competente para intervenir en un proceso de determinación de la capacidad es el del domicilio o el del lugar de internación de la persona. La jurisprudencia tiende a aplicar un principio de inmediatez: si la persona tiene domicilio en un lugar pero está internada en otro, suele resultar competente el juez más cercano al lugar de la internación.

\subsection{Legitimados activos (artículo 33)}

\begin{important}
\textbf{Artículo 33 --- Código Civil y Comercial (legitimados para solicitar la declaración).} Están legitimados para solicitar la declaración de incapacidad y de capacidad restringida: a) el propio interesado; b) el cónyuge no separado de hecho y el conviviente mientras la convivencia no haya cesado; c) los parientes dentro del cuarto grado y, si fuera por afinidad, dentro del segundo grado; d) el Ministerio Público.
\end{important}

Una novedad del nuevo código es que el propio interesado puede iniciar el proceso; el código anterior no lo contemplaba. Esto cobra sentido, por ejemplo, cuando una persona a la que se le diagnosticó una enfermedad de deterioro cognitivo desea anticiparse y solicitar su propia restricción de la capacidad, proponiendo a alguien de su confianza como apoyo.

También puede solicitarlo el cónyuge, mientras no esté separado de hecho ni divorciado, y el conviviente. En torno al alcance del término ``conviviente'' existe cierta ambigüedad en el código, que la doctrina discute: si abarca a toda persona que convive o solo a quienes tienen una unión convivencial registrada; la interpretación parecería ser amplia.

Los parientes legitimados son los que se encuentran dentro del cuarto grado --quienes tendrían vocación hereditaria en caso de fallecimiento-- y, si el parentesco es por afinidad, hasta el segundo grado. El código anterior admitía como legitimado a cualquier vecino cuando la persona era considerada ``furiosa'' (peligrosa); esa posibilidad fue eliminada. Actualmente, quien no encuadre en los supuestos legitimados debe acudir al Ministerio Público (regulado en el artículo 103 del Código Civil y Comercial), que evaluará la situación y, en su caso, se presentará ante el juez para iniciar el proceso.

\section{Medidas cautelares y entrevista personal}

\subsection{Medidas cautelares (artículo 34)}

\begin{important}
\textbf{Artículo 34 --- Código Civil y Comercial (medidas cautelares).} Durante el proceso, el juez puede ordenar medidas de restricción al ejercicio de la capacidad, con el fin de garantizar los derechos personales y patrimoniales de la persona. El juez debe determinar con precisión para qué actos se disponen estas medidas y designar los apoyos necesarios, o excepcionalmente el curador para el caso de los incapaces.
\end{important}

Durante el trámite del proceso la persona es, en principio, capaz; sin embargo, el juez puede disponer restricciones cautelares, siempre de manera excepcional y justificada, respetando los principios generales del artículo 31 y de la Convención, para atender situaciones de urgencia --por ejemplo, la designación de un apoyo para decisiones referidas a un inmueble o a un fondo de inversión mientras la persona atraviesa un problema de salud mental.

\subsection{Entrevista personal (artículo 35)}

\begin{important}
\textbf{Artículo 35 --- Código Civil y Comercial (entrevista personal).} El juez debe garantizar la inmediatez con el interesado durante el proceso y entrevistarlo personalmente, antes de dictar resolución alguna. Debe asegurar la accesibilidad y los ajustes razonables del procedimiento de acuerdo a la situación de aquel. El Ministerio Público y, al menos, un letrado que preste asistencia al interesado durante el proceso deben estar presentes en las audiencias.
\end{important}

Si bien la entrevista personal ya era una práctica habitual en algunos fueros desde hace décadas, el nuevo código la establece como exigencia general: antes de adoptar una decisión tan significativa, el juez debe conocer personalmente a la persona involucrada, garantizando accesibilidad y ajustes razonables.

\section{La persona como parte y asistencia letrada}

\begin{important}
\textbf{Artículo 36 --- Código Civil y Comercial.} La persona en cuyo interés se lleva adelante el proceso es parte y puede aportar todas las pruebas que hacen a su defensa. Interpuesta la solicitud de declaración de incapacidad o de restricción a la capacidad ante el juez, éste debe nombrar un abogado que lo represente, si la persona afectada no tuviera letrado propio. En el caso de personas incapaces, también es necesaria la presencia del Ministerio Público.
\end{important}

El artículo 36 enfatiza que la persona interesada no puede quedar ausente de su propio proceso, situación que en ocasiones ocurría en el pasado. La persona reviste calidad de parte y tiene derecho a aportar pruebas; si carece de abogado propio, el Estado --normalmente a través del Ministerio Público-- debe proveerle uno, cuya función es específicamente la representación judicial, distinta de la del apoyo o de la del curador.

\section{Dictamen interdisciplinario y contenido de la sentencia}

\subsection{Artículo 37: el dictamen interdisciplinario}

\begin{important}
\textbf{Artículo 37 --- Código Civil y Comercial (dictamen e informe interdisciplinario).} Para la declaración de incapacidad o de capacidad restringida es necesario el dictamen de un equipo interdisciplinario. El dictamen debe contener: 1) diagnóstico y pronóstico; 2) época en que se manifestó la situación; 3) recursos personales, familiares y sociales existentes; 4) régimen aconsejable para la protección y asistencia del interesado, con la mayor autonomía posible. Estos puntos deben reflejarse en la sentencia.
\end{important}

La Convención no exige literalmente un análisis interdisciplinario, pero de su espíritu se desprende que el examen no debe limitarse a lo médico, sino integrar distintos saberes. La interdisciplina implica que las disciplinas involucradas --psiquiatría, psicología y trabajo social-- alcancen un conocimiento común, dentro de los límites de sus respectivas competencias profesionales.

% TODO: contenido incompleto o ambiguo en las notas originales: no queda resuelto en el apunte si el dictamen interdisciplinario debe ser conjunto (un único informe integrado) o secuencial (informes individuales de cada disciplina).

El equipo interdisciplinario debe expedirse sobre el diagnóstico y pronóstico --distinguiendo, por ejemplo, una crisis psiquiátrica aguda de una situación permanente o de un deterioro cognitivo severo--; la época en que se manifestó la situación, relevante para la eventual anulación de actos anteriores conforme al artículo 45; los recursos personales, familiares y sociales existentes, que permiten conocer la red de apoyo natural de la persona; y el régimen aconsejable para la protección, asistencia y promoción de la mayor autonomía posible.

\subsection{Artículo 38: el contenido de la sentencia}

\begin{important}
\textbf{Artículo 38 --- Código Civil y Comercial (alcance de la sentencia).} La sentencia debe determinar la extensión y alcance de la restricción y especificar las funciones y actos que se limitan, procurando que la afectación de la autonomía personal sea la menor posible. Asimismo, designa una o más personas de apoyo o curadores según el caso, y señala las condiciones de validez de los actos específicos sujetos a la restricción indicando si la persona debe actuar con la asistencia o representación del apoyo o curador.
\end{important}

El artículo 38 concreta la personalización que exige la Convención: a diferencia del artículo 152 ter del código anterior, que resultaba desajustado respecto del resto del ordenamiento, en el nuevo código todo el sistema parte de la capacidad, y la sentencia se limita a precisar qué actos y funciones se restringen y con qué modalidad de apoyo.

\section{Registro, revisión, internación y traslados}

\subsection{Registro (artículo 39)}

La sentencia debe registrarse en el Registro de Estado Civil, anotándose al margen del acta de nacimiento. A partir de la inscripción se producen efectos frente a terceros: los actos posteriores a la inscripción que sean contrarios a lo dispuesto por la sentencia son nulos. Cuando se levanta la restricción, corresponde cancelar la anotación mediante oficio al Registro Civil.

Este registro plantea una tensión entre la protección de la intimidad --por tratarse de un dato altamente sensible-- y la necesidad de dar publicidad suficiente, ya que quienes contratan deben poder conocer si su interlocutor tiene alguna restricción a la capacidad.

\subsection{Revisión (artículo 40)}

\begin{important}
\textbf{Artículo 40 --- Código Civil y Comercial (revisión).} La sentencia debe ser revisada por el juez en un plazo no superior a tres años, sobre la base de nuevos dictámenes interdisciplinarios y mediando la audiencia personal con el interesado. Puede ser revisada en cualquier momento a instancia del interesado. El Ministerio Público debe fiscalizar el cumplimiento efectivo de la revisión periódica establecida e instarla, si corresponde.
\end{important}

La revisión puede solicitarse en cualquier momento, pero no puede dejar de producirse por más de tres años. Corresponde un nuevo dictamen interdisciplinario, una nueva audiencia y una nueva sentencia; el Ministerio Público debe fiscalizar que el juez cumpla con esta obligación y, en su caso, instarla.

\subsection{Internación (artículo 41) y traslados (artículo 42)}

\begin{important}
\textbf{Artículo 41 --- Código Civil y Comercial (internación).} La internación sin consentimiento de una persona, tenga o no restringida su capacidad, procede solo si se cumplen los recaudos previstos en la legislación especial y se cumplen los requisitos de esta sección. En particular, la internación involuntaria solo es procedente cuando a criterio del equipo de salud mediare situación de riesgo cierto e inminente para la persona o para terceros.
\end{important}

No es posible internar a una persona contra su voluntad si no existe riesgo cierto e inminente de un daño de entidad para ella o para terceros. Este estándar, proveniente de la Ley de Salud Mental y ratificado por el nuevo código, busca evitar los abusos que en el pasado llevaron a internar a personas por carecer de vivienda o por resultar molestas; la internación involuntaria debe ser breve, acotada, orientada a la recuperación en un contexto de crisis y a la posterior reintegración, con supervisión periódica, debido proceso, control judicial y defensa --entre otros mecanismos, la Unidad de Letrados del artículo 22 de la Ley de Salud Mental, encargada de velar por los derechos de las personas internadas por razones de salud mental.

\section{Cuadro Cornell --- Bloque 3}

\begin{table}[H]
\centering
\begin{tabularx}{\textwidth}{
    >{\raggedright\arraybackslash}p{0.30\textwidth}
    >{\raggedright\arraybackslash}X
}
\cornellhead

\textbf{¿Quién puede iniciar un proceso de restricción de la capacidad?} &
El propio interesado, el cónyuge o conviviente no separado de hecho, los parientes dentro del cuarto grado (segundo grado por afinidad) y el Ministerio Público (art. 33). \\

\textbf{¿Por qué el juez debe entrevistar personalmente al interesado?} &
Porque el artículo 35 exige inmediatez y entrevista personal antes de dictar cualquier resolución, garantizando accesibilidad y ajustes razonables, de modo que el juez conozca directamente a la persona antes de restringir su capacidad. \\

\textbf{¿Qué garantiza que la persona no quede ausente de su propio proceso?} &
El artículo 36, que le reconoce calidad de parte, derecho a aportar pruebas y derecho a contar con un abogado --provisto por el Estado si no tuviera uno propio-- que la asista específicamente en el proceso judicial. \\

\textbf{¿Qué cuatro aspectos debe abordar toda sentencia?} &
Según el artículo 37, el dictamen y la sentencia deben contener: diagnóstico y pronóstico; época en que se manifestó la situación; recursos personales, familiares y sociales existentes; y régimen aconsejable para la protección y promoción de la mayor autonomía posible. \\

\textbf{¿Cuándo es legítimo internar a una persona contra su voluntad?} &
Solo ante riesgo cierto e inminente de daño para la persona o para terceros, fundado en evaluación interdisciplinaria, con control judicial inmediato, asistencia letrada y revisión periódica (art. 41). \\

\midrule
\multicolumn{2}{p{\textwidth}}{
\textbf{Resumen:} El proceso de determinación de la capacidad está atravesado por garantías procesales: legitimación amplia pero acotada, medidas cautelares excepcionales, entrevista personal obligatoria, calidad de parte del interesado con asistencia letrada, dictamen interdisciplinario con contenido mínimo obligatorio, registro con efectos frente a terceros, revisión periódica no mayor a tres años, e internación involuntaria limitada al riesgo cierto e inminente.
} \\
\bottomrule
\end{tabularx}
\end{table}

% =========================================================================
% CAPÍTULO 4
% =========================================================================
\chapter{Validez de los actos y sistema de apoyos}
\label{ch:bloque4}

\section{Nulidad de los actos posteriores y anteriores a la sentencia}

\subsection{Actos posteriores a la inscripción (artículo 44)}

\begin{important}
\textbf{Artículo 44 --- Código Civil y Comercial (actos posteriores a la sentencia).} Son nulos los actos de la persona incapaz o con capacidad restringida que contrarían lo dispuesto en la sentencia realizada con posterioridad a su inscripción en el Registro de Estado Civil y Capacidad de las Personas.
\end{important}

Si una persona sujeta a una sentencia que restringe su capacidad realiza un acto prohibido por esa sentencia, el acto es nulo. La nulidad es la sanción civil que deja sin efecto el acto y retrotrae la situación al estado anterior: si se entregó un inmueble, debe restituirse; si se entregó dinero, debe restituirse.

\begin{example}{Venta de un inmueble contraria a la sentencia}
Una persona con discapacidad mental que vive de la renta de varias cocheras tiene una sentencia que exige el apoyo de un hermano o hermana para vender inmuebles. Si, sin ese apoyo, vende un inmueble y el escribano no advierte la restricción, el acto es nulo: la venta no produce efectos y las partes deben restituirse lo entregado.
\end{example}

Se trata de una nulidad relativa. Para que proceda, el acto debe ser posterior a la inscripción de la sentencia y contrario a lo que ella dispone. Si la sentencia solo restringe, por ejemplo, la venta de inmuebles, el consentimiento que la misma persona preste para una intervención quirúrgica --materia no alcanzada por la sentencia-- es perfectamente válido.

\subsection{Actos anteriores a la sentencia (artículo 45)}

\begin{important}
\textbf{Artículo 45 --- Código Civil y Comercial (actos anteriores a la sentencia).} Los actos anteriores a la inscripción de la sentencia pueden ser declarados nulos, si perjudican a la persona incapaz o con capacidad restringida, y se cumple alguno de los siguientes requisitos: a) la enfermedad mental era ostensible al momento de la celebración del acto; b) quien contrató con ella era de mala fe; c) el acto fue a título gratuito.
\end{important}

Estos actos son, en principio, válidos, pero pueden anularse si concurre alguno de los requisitos del artículo 45. El dictamen interdisciplinario del artículo 37, al establecer la época en que se manifestó la situación, permite determinar si existían causales al momento de celebrarse el acto anterior.

\begin{example}{Contratante de mala fe}
Un profesional de la salud que atiende a una persona con capacidad restringida y, conociendo su situación, le compra un inmueble a bajo precio, actúa de mala fe. Si, en cambio, el acto fue a título gratuito, no existe perjuicio patrimonial porque la persona no recibió contraprestación alguna.
\end{example}

\section{Cese de las restricciones (artículo 47)}

\begin{important}
\textbf{Artículo 47 --- Código Civil y Comercial (cese de la incapacidad y de las restricciones a la capacidad).} El cese de la incapacidad o de la restricción a la capacidad debe decretarse por el juez que la declaró, previo examen de un equipo interdisciplinario integrado conforme a las pautas del artículo 37, que dictamine sobre el restablecimiento de la persona. Si el restablecimiento es total, el juez deja sin efecto la sentencia. Si el restablecimiento es parcial o el interesado solicita expresamente la revisión, el juez puede ampliar la nómina de actos que la persona puede realizar por sí o con la asistencia o representación de los apoyos designados.
\end{important}

Además de la revisión obligatoria del artículo 40, el interesado puede solicitar en cualquier momento el cese de la restricción, mediante el mismo procedimiento: un nuevo examen interdisciplinario conforme al artículo 37 (diagnóstico, pronóstico, época, recursos personales, familiares y sociales, y régimen aconsejable).

\section{El sistema de apoyos (artículo 43) y conclusiones}

\begin{important}
\textbf{Artículo 43 --- Código Civil y Comercial (concepto de apoyo).} Se entiende por apoyo cualquier medida de carácter judicial o extrajudicial que facilite a la persona la toma de decisiones para dirigir su persona, administrar sus bienes y celebrar actos jurídicos en general. Las medidas de apoyo tienen como función la de promover la autonomía y facilitar la comunicación, la comprensión y la manifestación de voluntad de la persona para el ejercicio de sus derechos. El interesado puede proponer al juez la designación de una o más personas de su confianza para que le presten apoyo. El juez debe evaluar los alcances de la designación y procurar la protección de la persona respecto de eventuales conflictos de intereses o influencia indebida. La resolución debe establecer la condición y la calidad de las medidas de apoyo y, de ser necesario, ser inscripta en el Registro de Estado Civil y Capacidad de las Personas.
\end{important}

El código toma de la Convención la noción de apoyo como cualquier medida judicial o extrajudicial que facilite la toma de decisiones. El apoyo judicial es el que se designa en la sentencia; el apoyo extrajudicial resulta más difícil de precisar, porque toda persona recibe apoyos informales de quienes la rodean, sin que ello implique control judicial. La cuestión se presenta cuando el apoyo adquiere relevancia jurídica.

\begin{example}{Apoyo extrajudicial sin relevancia jurídica}
Durante una situación de aislamiento, un vecino que hace las compras de una persona mayor que no puede salir --recibiendo el dinero, comprando y entregando la mercadería-- presta un apoyo de hecho que no requiere control judicial, en la medida en que no involucre actos como vender un inmueble, comprar un automóvil o hipotecar una propiedad.
\end{example}

\subsection{Salvaguardias: conflicto de intereses e influencia indebida}

Para designar apoyos, el juez debe valorar la eventual existencia de conflicto de intereses e influencia indebida, dos mecanismos de salvaguarda previstos por la Convención para evitar abusos. El conflicto de intereses es más objetivo --por ejemplo, cuando quien administra un inmueble como apoyo es, a la vez, quien lo alquila--; la influencia indebida es más difícil de probar, porque supone cierta manipulación de la voluntad, situación que puede presentarse, por ejemplo, cuando una persona mayor con discapacidad mental convive con uno de sus hijos y los demás hijos consideran que existe manipulación.

\subsection{Criterio de actuación de los apoyos}

El criterio de actuación de los apoyos es el respeto de la voluntad y las preferencias de la persona, principio que surge directamente del artículo 12 de la Convención. El apoyo puede ser propuesto por la propia persona interesada. La posibilidad de que una persona jurídica actúe como apoyo se encuentra en discusión doctrinaria y no consta que exista jurisprudencia que la haya admitido respecto de la institución donde se encuentra internada una persona.

% TODO: contenido incompleto o ambiguo en las notas originales: la posibilidad de que una persona jurídica sea designada apoyo queda planteada como debate abierto, sin resolución expresa en el material.

\subsection{Conclusiones generales}

El nuevo código se adapta a la Convención mediante conceptos y terminología propios de ella --apoyos, ajustes razonables, accesibilidad, autonomía, recursos personales y familiares--. Ratifica el criterio de la Ley de Salud Mental según el cual la regla es la capacidad y las limitaciones son excepcionales, por actos y funciones concretos. El sistema general es el de apoyos; se reconoce ampliamente el derecho de la persona a participar en el proceso y a ser escuchada; la sentencia debe revisarse periódicamente; y se consolida el criterio de intervención interdisciplinaria.

Subsiste, sin embargo, la cuestión de la incapacidad y de los supuestos de sustitución de la voluntad, en tanto se discute si esas formas de representación respetan plenamente la Convención. La doctrina mayoritaria entiende que no la vulneran, por tratarse de situaciones absolutamente excepcionales en las que la persona no puede expresarse de ningún modo y en las que, de no designarse representante, quedaría sin protección jurídica alguna. A ello se suma que la persona conserva el derecho a participar y puede dejar directivas anticipadas --instituto recibido de manera limitada en el artículo 139, que permite a una persona señalar anticipadamente quién podría ser su curador ante la eventual pérdida de conciencia.

\section{Cuadro comparativo: régimen anterior y régimen vigente}

\begin{longtable}{
    >{\raggedright\arraybackslash}p{0.22\textwidth}
    >{\raggedright\arraybackslash}p{0.36\textwidth}
    >{\raggedright\arraybackslash}p{0.36\textwidth}
}
\toprule
\textbf{Situación} & \textbf{Código Civil derogado} & \textbf{Código Civil y Comercial vigente} \\
\midrule
\endfirsthead
\toprule
\textbf{Situación} & \textbf{Código Civil derogado} & \textbf{Código Civil y Comercial vigente} \\
\midrule
\endhead

\textbf{Regla general} & Incapacidad como eje del sistema; el inhabilitado era la excepción acotada a tres casos. & La capacidad es la regla; toda restricción es excepcional y se impone en beneficio de la persona (art. 31). \\[0.4em]

\textbf{Figura de protección} & Curador con función de representación, para incapaces e inhabilitados. & Apoyos para personas con capacidad restringida e inhabilitados; curador solo para incapaces (art. 32, último párrafo). \\[0.4em]

\textbf{Intervención profesional} & Evaluación médica (luego interdisciplinaria, desde el art. 152 ter de 2010). & Dictamen interdisciplinario obligatorio: psiquiatría, psicología y trabajo social (art. 37). \\[0.4em]

\textbf{Duración de la sentencia} & Sin plazo, hasta el art. 152 ter (2010), que fijó un máximo de tres años. & Revisión obligatoria cada tres años como máximo, o antes a pedido del interesado (art. 40). \\[0.4em]

\textbf{Voz de la persona} & No estaba prevista la entrevista personal obligatoria ni la condición de parte. & Entrevista personal obligatoria (art. 35) y la persona es parte del proceso (art. 36). \\[0.4em]

\textbf{Sistema de apoyos} & No existía el sistema de apoyos; solo representación a través del curador. & Sistema de apoyos como regla; respeto de la voluntad y preferencias de la persona (art. 43). \\
\bottomrule
\end{longtable}

% =========================================================================
% CORNELL DE REPASO GENERAL (comprensivo, aportado por el propio apunte)
% =========================================================================

\section{Cuadro Cornell de repaso general}
\label{ch:cornell-general}

Este cuadro de repaso reúne las preguntas y conceptos clave desarrollados a lo largo de los cuatro bloques, pensado para el estudio activo previo a un examen.

\begin{longtable}{
    >{\raggedright\arraybackslash}p{0.30\textwidth}
    >{\raggedright\arraybackslash}p{0.62\textwidth}
}
\toprule
\textbf{Pregunta / concepto clave} & \textbf{Notas / respuesta / explicación} \\
\midrule
\endfirsthead
\toprule
\textbf{Pregunta / concepto clave} & \textbf{Notas / respuesta / explicación} \\
\midrule
\endhead

\textbf{¿Qué jerarquía tiene la CDPD en Argentina?} &
Fue ratificada en 2008 (Ley 26.378) y en 2014 obtuvo jerarquía constitucional mediante el procedimiento del artículo 75, inciso 22, de la Constitución Nacional, equiparándose a la primera parte de la Constitución. \\[0.4em]

\textbf{¿Qué establece el artículo 12 de la CDPD?} &
Reconoce la personalidad jurídica de las personas con discapacidad, su derecho a la capacidad jurídica en igualdad de condiciones, la obligación estatal de garantizar apoyos para el ejercicio de esa capacidad y la exigencia de salvaguardias que respeten la voluntad, evitando conflictos de intereses e influencia indebida. \\[0.4em]

\textbf{¿Cuáles son las reglas generales del artículo 31 del CCyC?} &
La capacidad se presume, incluso en caso de internación; las restricciones son excepcionales y en beneficio de la persona; la intervención es interdisciplinaria; hay derecho a información accesible y a asistencia letrada gratuita; y deben priorizarse las alternativas menos restrictivas. \\[0.4em]

\textbf{¿Cuáles son los tres tipos de restricción a la capacidad del CCyC?} &
Personas con capacidad restringida (art. 32, primer párrafo, la regla), personas incapaces (art. 32, último párrafo, supuesto excepcional) e inhabilitados por prodigalidad (art. 48). \\[0.4em]

\textbf{¿Qué requisitos exige el artículo 32 para restringir la capacidad?} &
Persona mayor de trece años; adicción o alteración mental permanente o prolongada de suficiente gravedad; y que el ejercicio de la plena capacidad pueda causarle un daño a su persona o a sus bienes. \\[0.4em]

\textbf{¿Cuándo corresponde declarar la incapacidad?} &
Solo cuando la persona esté absolutamente imposibilitada de interaccionar con su entorno, no pueda expresar su voluntad por ningún modo, medio o formato adecuado, y el sistema de apoyos resulte ineficaz. Se designa entonces un curador con funciones de representación. \\[0.4em]

\textbf{¿Qué protege la inhabilitación por prodigalidad?} &
Protege al cónyuge, conviviente o hijos menores o con discapacidad frente al riesgo de pérdida del patrimonio por la mala gestión de los bienes; limita los actos de disposición y exige la conformidad del apoyo para ciertos actos. \\[0.4em]

\textbf{¿Quiénes están legitimados para pedir la restricción de la capacidad?} &
El propio interesado, el cónyuge o conviviente no separado de hecho, los parientes dentro del cuarto grado (segundo grado por afinidad) y el Ministerio Público (art. 33). \\[0.4em]

\textbf{¿Qué exige la entrevista personal del artículo 35?} &
Que el juez garantice la inmediatez con el interesado y lo entreviste personalmente antes de dictar cualquier resolución, asegurando accesibilidad y ajustes razonables del procedimiento. \\[0.4em]

\textbf{¿Qué cuatro puntos debe contener el dictamen y reflejar la sentencia según el artículo 37?} &
1) Diagnóstico y pronóstico; 2) época en que se manifestó la situación; 3) recursos personales, familiares y sociales existentes; 4) régimen de protección y promoción de la mayor autonomía posible, sobre la base del dictamen interdisciplinario. \\[0.4em]

\textbf{¿Cuándo es válida una internación involuntaria?} &
Solo ante riesgo cierto e inminente de daño para la persona o para terceros, fundada en evaluación interdisciplinaria, con control judicial inmediato, asistencia letrada y revisión periódica (art. 41). \\[0.4em]

\textbf{¿Qué efectos produce un acto contrario a la sentencia?} &
Si es posterior a la inscripción de la sentencia, el acto es nulo (art. 44); si es anterior, puede anularse si la enfermedad era ostensible, hubo mala fe del contratante, o el acto fue a título gratuito (art. 45). \\[0.4em]

\textbf{¿Qué es un apoyo según el artículo 43?} &
Cualquier medida judicial o extrajudicial que facilite a la persona la toma de decisiones para dirigir su vida, administrar sus bienes y celebrar actos jurídicos, promoviendo su autonomía y respetando su voluntad y preferencias. \\[0.4em]

\textbf{¿Qué diferencia hay entre conflicto de intereses e influencia indebida?} &
El conflicto de intereses es objetivo y visible (el apoyo que administra un inmueble que él mismo alquila); la influencia indebida es más sutil y difícil de probar, porque supone una manipulación de la voluntad de la persona. \\
\midrule
\multicolumn{2}{p{0.94\textwidth}}{
\textbf{Síntesis final:} El Código Civil y Comercial argentino adapta el sistema de restricciones a la capacidad al modelo social de la discapacidad consagrado en la Convención sobre los Derechos de las Personas con Discapacidad (CDPD). Parte de un principio general de capacidad y reserva la restricción para casos excepcionales, debidamente fundados en un dictamen interdisciplinario y resueltos por sentencia judicial personalizada. Reconoce tres figuras --capacidad restringida, incapacidad e inhabilitación por prodigalidad-- que se diferencian por la intensidad de la limitación y por la naturaleza de la figura de protección: el apoyo, que asiste sin sustituir la voluntad, o el curador, reservado a los casos más extremos de imposibilidad absoluta de interacción. Todo el procedimiento está atravesado por garantías procesales --entrevista personal, calidad de parte, asistencia letrada, revisión periódica-- que buscan que la persona nunca quede ausente de las decisiones que la afectan. Finalmente, el sistema de apoyos del artículo 43, con sus salvaguardias contra el conflicto de intereses y la influencia indebida, condensa el espíritu de toda la reforma: acompañar para que la persona pueda decidir por sí misma, y proteger solo en la medida estrictamente necesaria.
} \\
\bottomrule
\end{longtable}
\chapter[Derechos personalísimos]{Derechos personalísimos: noción, caracteres y régimen}
\label{ch:personalisimos}

\section{Introducción}

Los derechos personalísimos son derechos subjetivos fundamentales que protegen la dignidad humana en sus manifestaciones físicas y espirituales. Estos derechos reconocen y garantizan a cada persona el respeto a su vida, integridad corporal, libertad, honor, intimidad e identidad.

En el ejercicio profesional de la abogacía, la comprensión de estos derechos resulta necesaria para:

\begin{itemize}
    \item Defender a clientes cuyos derechos personalísimos han sido vulnerados.
    \item Asesorar sobre el consentimiento informado en procedimientos médicos.
    \item Proteger la intimidad y la privacidad de datos en el entorno digital.
    \item Litigar casos de difamación, invasión de la privacidad e injurias.
    \item Orientar cuestiones relacionadas con la bioética y los experimentos humanos.
    \item Comprender la protección constitucional de la dignidad humana.
\end{itemize}

Estos derechos adquieren especial relevancia en el contexto de la tecnología digital avanzada, donde los datos personales, las imágenes y la información íntima pueden ser capturados, almacenados y difundidos sin consentimiento. El marco legal que se desarrolla a continuación permite proteger tanto los derechos propios como los ajenos frente a estos desafíos contemporáneos.

\section{Alcance y orden de exposición}

Los derechos personalísimos forman un sistema integrado en el que la dignidad humana constituye el principio rector. La exposición comienza con los conceptos generales y los antecedentes, continúa con las características y la clasificación, prosigue con la regulación en el Código Civil y Comercial y se detiene, finalmente, en el derecho a la libertad. Cada bloque se construye sobre el anterior, de forma análoga a las capas que envuelven un núcleo central de protección de la persona humana.

% TODO: verificar consistencia entre fuentes originales. El mapa de contenidos de la fuente anuncia el desarrollo posterior de los bloques de integridad física e integridad espiritual (por ejemplo, vida, integridad corporal, honor, imagen, intimidad e identidad), pero ese desarrollo no figura en los archivos recibidos; solo se desarrollan los fundamentos, las características, la clasificación, la regulación del Código y el derecho a la libertad.

\section{Noción de derechos personalísimos}

Los derechos personalísimos corresponden a la persona humana. Son derechos que pertenecen a todo ser humano por el solo hecho de ser tal.

\begin{definition}{Derechos personalísimos (Tobías)}
Los derechos personalísimos son derechos subjetivos que reconocen y garantizan a la persona humana el goce en su propia entidad, en su interioridad, en todas sus manifestaciones físicas y espirituales.
\end{definition}

Estos derechos no deben confundirse con los atributos de la personalidad, distinción que tampoco realiza el Código Civil y Comercial. Los atributos de la personalidad son cualidades esenciales que permiten que la persona despliegue su personalidad y entable relaciones jurídicas. Los derechos personalísimos, en cambio, son verdaderos derechos subjetivos que protegen dimensiones fundamentales de la persona.

Estos derechos protegen dos dimensiones fundamentales:

\begin{enumerate}
    \item La \textbf{dimensión de la integridad física}: la vida y la integridad corporal.
    \item La \textbf{dimensión espiritual}: la libertad, el honor, la imagen, la intimidad y la identidad de la persona.
\end{enumerate}

\section{Importancia y contexto actual}

A lo largo de la historia siempre ha existido una protección de los valores esenciales de la persona: la vida, la integridad física, la libertad, la imagen y el honor. Sin embargo, especialmente en la última mitad del siglo XX, se dieron circunstancias que motivaron a los juristas y al legislador a sancionar normas y a adoptar tratados internacionales orientados a proteger estas dimensiones de la persona humana.

Han surgido nuevas formas de amenaza que reclaman una respuesta jurídica. En el entorno de las biotecnologías, hoy es posible intervenir sobre el cuerpo humano de una manera mucho más poderosa: el conocimiento de la genética y de la biología ha dado lugar a enormes beneficios, pero también a nuevas formas de intervención que pueden no ser respetuosas de la integridad física.

\begin{note}
A lo largo de la historia se han realizado experimentos sobre personas humanas sin su consentimiento, aprovechando la condición de vulnerabilidad de determinados grupos, lo que ha dado lugar a fuertes reacciones éticas y jurídicas.
\end{note}

Las nuevas tecnologías inciden también en la transmisión de la vida humana, en los trasplantes de órganos y en las intervenciones biomédicas y biotecnológicas en general. Respecto de la investigación en seres humanos, la genética permite hoy intervenciones de gran alcance, como ocurrió en China en el año 2018, cuando se obtuvieron seres genéticamente modificados. Estas posibilidades tecnológicas son enormemente promisorias para el bienestar y el avance de las personas, pero presentan también aspectos problemáticos que requieren una reflexión cuidadosa.

\subsection{Datos personales y privacidad digital}

Otro ámbito relevante es el de la recolección, el uso, el tratamiento, la transmisión y el almacenamiento de los datos personales: datos vinculados con la identidad, con las características personales, con el consumo o con la salud. La cantidad de datos que hoy se almacena es enorme, fenómeno conocido como \textit{big data}: la tendencia a generar grandes colecciones de datos de las personas para elaborar perfiles, realizar seguimientos y dirigir ofertas publicitarias.

Los datos personales permiten tomar decisiones más racionales y conocer mejor a las personas, lo cual puede traer beneficios, pero también nuevas posibilidades de discriminación y de ataques a las personas en función de sus datos. Cuando se trata de datos sensibles existen nuevas posibilidades de manipulación; así, en determinadas situaciones políticas, las campañas se han dirigido a públicos específicos para persuadirlos e influir en su voto.

\begin{important}
El desarrollo de la informática ha transformado radicalmente la capacidad de procesamiento de datos personales —antes limitada a fichas escritas a mano— lo que exige nuevas formas de protección de la intimidad, la identidad y el honor de las personas. La posibilidad actual de captar la voz y la imagen de manera mucho más poderosa hace que estos elementos, una vez registrados, se conviertan también en datos personales sujetos a este mismo régimen de protección.
\end{important}

\section{Antecedentes en Argentina}

Entre los autores que en la Argentina han desarrollado el estudio de esta materia se encuentran Salvat, Cifuentes, Rivera, Navarro Floria y Tobías, entre otros. Cabe señalar que este tema desborda el contenido propio de la parte general del derecho civil y podría constituir por sí mismo una materia completa.

% TODO: contenido incompleto o ambiguo en las notas originales — el segundo autor mencionado aparece en el apunte como "el profesor Salto Cifuentes", lo que podría corresponder a dos autores distintos (Salvat y Cifuentes) mal unidos en la transcripción; se conserva la grafía más plausible sin poder confirmarlo con el material disponible.

\subsection{El Código de Vélez}

El Código de Vélez no contenía un tratamiento metódico y orgánico de los derechos personalísimos. No obstante, los autores reconocen la presencia del tema en algunos de sus artículos:

\begin{longtable}{p{0.30\textwidth} p{0.62\textwidth}}
\toprule
\textbf{Disposición} & \textbf{Contenido} \\
\midrule
\endfirsthead
\toprule
\textbf{Disposición} & \textbf{Contenido} \\
\midrule
\endhead
Artículos 428 y 2.126 & Hablan de derechos inherentes a la persona. \\
\addlinespace
Artículo 1.075 & Establece que toda afectación de un derecho que se confunda con la existencia de la persona puede ser materia de un delito. \\
\bottomrule
\end{longtable}

También en la nota al artículo 2.312, el Código de Vélez hacía referencia a ciertos derechos que tienen origen en la existencia del individuo mismo, como la libertad, el honor, el cuerpo y la patria potestad. Se trata de una mención presente en el código, aunque no sistemática ni ordenada.

\subsection{Evolución legislativa en Argentina}

\begin{longtable}{p{0.30\textwidth} p{0.62\textwidth}}
\toprule
\textbf{Norma} & \textbf{Contenido} \\
\midrule
\endfirsthead
\toprule
\textbf{Norma} & \textbf{Contenido} \\
\midrule
\endhead
Ley 11.723 (1933) — Ley de Propiedad Intelectual & Norma que continúa vigente. En sus artículos 31 a 35 regula la fotografía y el derecho a la imagen, incluyendo el consentimiento y los casos en que puede prescindirse de él. \\
\addlinespace
Ley 21.173 — incorporó el artículo 1.071 bis al código anterior & Protegía el derecho a la intimidad. Este artículo sería recogido posteriormente por el nuevo Código Civil y Comercial en su artículo 1.770. \\
% TODO: contenido incompleto o ambiguo en las notas originales — el apunte indica el año 1905 para esta incorporación, lo cual resulta inconsistente con la numeración de la ley citada; se conserva el dato tal como figura en el material original.
\addlinespace
Ley 24.193 (1992) — Ley de Trasplantes & Reemplazada en 2018 por la Ley 27.447. \\
\addlinespace
Ley 25.326 (2000) — Ley de Protección de Datos Personales & Aún vigente, con proyectos de reforma. \\
\addlinespace
Ley 26.529 (2009) — Ley de Derechos del Paciente & Reformada en 2012 por la Ley 26.742. \\
\bottomrule
\end{longtable}

Estas normas reflejan las dimensiones de los derechos personalísimos vinculadas con la integridad física, la integridad espiritual y la libertad.

\subsection{La Constitución Nacional de 1994}

La reforma constitucional de 1994 incorpora, en el artículo 75 inciso 22, diversos tratados de derechos humanos con jerarquía constitucional que consagran derechos correspondientes a la personalidad.

Se produce así una convergencia entre el derecho público y el derecho privado. Aunque los tratados de derechos humanos incorporados a la Constitución corresponden al derecho público, tienen resonancias directas en el derecho privado, en tanto corresponden a la persona humana, que se despliega fundamentalmente en el campo civil.

\begin{note}
La protección de estos derechos tiene también faceta penal: la protección del derecho al honor puede dar lugar al delito de calumnias e injurias, además de a la protección civil. Ambas protecciones suelen coexistir. El presente desarrollo se limita a las dimensiones civiles del asunto y a su incorporación al Código Civil y Comercial.
\end{note}

\section{Caracteres de los derechos personalísimos}

Según Tobías, los derechos personalísimos presentan los siguientes caracteres:

\begin{itemize}
    \item \textbf{Son derechos subjetivos privados:} corresponden a la esfera de las personas, dimensión donde operan el derecho privado y el derecho civil.
    \item \textbf{Son derechos innatos y originarios:} no son derechos concedidos por el legislador, sino reconocidos. Un ser humano, por el solo hecho de ser tal, tiene derecho a la vida, al honor, a la imagen y a la identidad.
    \item \textbf{Son vitalicios y necesarios:} acompañan a la persona durante toda su vida; una persona humana siempre debe tenerlos.
    \item \textbf{Son absolutos (o de exclusión):} tienen un sujeto pasivo universal. Frente a cualquier persona existe la obligación de respetarlos, sin que sea necesario haber celebrado un contrato ni haber tenido relación previa alguna.
    \item \textbf{Son inherentes a la persona:} no pueden separarse de la persona; forman parte de su ser mismo.
    \item \textbf{Son intransmisibles:} no es posible transmitir completamente, por ejemplo, el derecho a la integridad física, ni delegar en otro la toma de decisiones sobre estos derechos.
    \item \textbf{Son relativamente indisponibles:} en principio no se pueden disponer, salvo mediante disposiciones relativas, de carácter restrictivo y siempre con posibilidad de revocación.
    \item \textbf{Son irrenunciables:} no es posible renunciar a estos derechos. La manifestación de ceder o renunciar, por ejemplo, al derecho a la imagen, carece de validez.
\end{itemize}

\begin{important}
Estos caracteres protegen dimensiones fundamentales de la persona y, en última instancia, tienen como bien jurídico común la dignidad humana.
\end{important}

\section{Clasificación tripartita}

Se adopta aquí la clasificación tripartita habitualmente utilizada por el profesor Tobías, sin perjuicio de que otros autores propongan clasificaciones distintas.

\begin{table}[H]
\centering
\begin{tabular}{p{0.28\textwidth} p{0.28\textwidth} p{0.28\textwidth}}
\toprule
\textbf{Integridad física} & \textbf{Libertad} & \textbf{Integridad espiritual} \\
\midrule
Vida \newline Integridad corporal \newline Salud &
Concepto que comprende aspectos tanto corporales como espirituales &
Honor \newline Intimidad \newline Imagen \newline Identidad \\
\bottomrule
\end{tabular}
\end{table}

En torno a esta clasificación se organiza el desarrollo del capítulo, que aborda las disposiciones generales del Código Civil y Comercial y una primera aproximación al derecho a la libertad.

\section{Regulación en el Código Civil y Comercial}

Metodológicamente, el Código Civil y Comercial incorpora los derechos personalísimos en el Capítulo 3, Título Primero del Libro Primero. El Libro Primero —Parte General— se abre con el Título Primero, dedicado a la persona humana; dentro de este título, luego de regular el comienzo de la existencia de la persona y la capacidad de las personas, aparecen los derechos personalísimos en los artículos 51 a 62.

\begin{important}
La incorporación sistemática de los derechos personalísimos constituye una de las grandes novedades del Código Civil y Comercial, en la medida en que el Código de Vélez no traía un tratamiento sistemático de estos derechos. Esta incorporación otorga una mayor presencia e importancia a la dignidad humana en el tratamiento del Código Civil y Comercial.
\end{important}

\subsection{Disposiciones generales: artículo 51}

El capítulo se abre con el artículo 51, referido a la inviolabilidad de la persona humana:

\begin{formula}{Artículo 51 CCyC}
``La persona humana es inviolable y en cualquier circunstancia tiene derecho al reconocimiento y respeto de su dignidad.''
\end{formula}

El artículo no menciona explícitamente el derecho a la vida, como sí lo hacían proyectos anteriores, pero dicha protección puede reconocerse detrás de la noción de inviolabilidad de la persona: no es posible afectar su persona, ni su cuerpo, ni su integridad física, ni su integridad espiritual. La expresión ``reconocimiento'' señala que la dignidad no constituye una concesión del ordenamiento jurídico, sino algo preexistente que el ordenamiento se limita a reconocer.

\subsection{Medios de protección: prevención y reparación}

El artículo 52 establece que la enumeración de medios de protección no es taxativa, sin perjuicio de que exista una sistematización que distingue entre:

\begin{table}[H]
\centering
\begin{tabular}{p{0.30\textwidth} p{0.55\textwidth}}
\toprule
\textbf{Prevención} & \textbf{Reparación} \\
\midrule
Consiste en evitar que ocurra el hecho que menoscabaría la dignidad; se trata de impedir que se produzca el acto dañino. &
Si el hecho ocurriere, se trata de volver al estado anterior y, si ello no fuera posible, de indemnizar los daños producidos. \\
\bottomrule
\end{tabular}
\end{table}

El Código incorpora, además, nuevas formas de reparación, entre ellas la posibilidad de publicar la sentencia en la que se ha discutido la afectación de un derecho personalísimo. Esta forma de reparación, presente ya en artículos anteriores, se incorpora también en la parte de daños del Libro Tercero como posibilidad específica cuando se encuentra afectado el honor de la persona.

\subsection{Disponibilidad relativa: artículo 55}

El artículo 55 dispone que los derechos personalísimos no son disponibles, salvo mediante el consentimiento de la persona, consentimiento que se encuentra sujeto a una serie de límites.

\begin{formula}{Artículo 55 CCyC — Disponibilidad relativa}
El consentimiento no puede ser contrario a las buenas costumbres. El consentimiento no se presume; es de interpretación restrictiva y, ante la duda, se interpreta que no ha sido otorgado. Es libremente revocable en cualquier momento.
\end{formula}

\begin{important}
El consentimiento sobre los derechos personalísimos debe reunir las siguientes condiciones:
\begin{itemize}
    \item No puede ser contrario a la ley, la moral o las buenas costumbres.
    \item Debe ser expresado; no puede presumirse a partir de la mera falta de oposición. Puede manifestarse por escrito, en forma verbal o mediante un gesto comprensible, pero siempre debe existir una manifestación real.
    \item Es de interpretación restrictiva: ante la duda, se interpreta que no hubo consentimiento. Así, el consentimiento otorgado para el uso de la propia imagen en una campaña de bien común no se extiende, por interpretación restrictiva, a una campaña con fines comerciales.
    \item Es libremente revocable: la persona puede dar marcha atrás y revocar el consentimiento oportunamente otorgado.
\end{itemize}
\end{important}

\section{Derecho a la libertad}

El valor de la libertad constituye uno de los derechos personalísimos más importantes. Su fuente se ubica generalmente en el artículo 19 de la Constitución Nacional.

\begin{formula}{Artículo 19 de la Constitución Nacional}
``Las acciones privadas de los hombres que de ningún modo ofendan a la moral pública, ni perjudiquen a un tercero, están exentas de la autoridad de los magistrados. Ningún habitante de la Nación será obligado a hacer lo que la ley no manda, ni privado de lo que ella no prohíbe.''
\end{formula}

Este constituye un principio de libertad fundamental en el ordenamiento jurídico argentino. Desde la perspectiva civil, la protección de la libertad se vincula sobre todo con los supuestos en que una persona sufre violencia para celebrar un acto; en tal caso, la persona afectada tiene derecho a reclamar los daños y perjuicios correspondientes.

También existe protección de la libertad en el campo de los derechos constitucionales, a través del habeas corpus, de la acción de amparo y de la acción de inconstitucionalidad, instrumentos que permiten conocer la situación de una persona, su paradero y su situación respecto de su integridad y de su libertad.

\begin{note}
Al igual que en el resto de los derechos personalísimos, las formas de protección de la libertad se articulan a través del artículo 52, mediante la prevención y la reparación.
\end{note}

% ================================================================
% CUADRO CORNELL
% ================================================================
\section{Cuadro Cornell de repaso}

El siguiente cuadro permite repasar los conceptos centrales desarrollados en este capítulo mediante la técnica de recuperación activa propia del método Cornell.

\begin{longtable}{
    >{\raggedright\arraybackslash}p{0.30\textwidth}
    >{\raggedright\arraybackslash}p{0.60\textwidth}
}
\toprule
\textbf{Preguntas / palabras clave} & \textbf{Notas / respuestas} \\
\midrule
\endfirsthead
\toprule
\textbf{Preguntas / palabras clave} & \textbf{Notas / respuestas} \\
\midrule
\endhead

\textbf{¿Qué son los derechos personalísimos?} &
Son derechos subjetivos que reconocen y garantizan a la persona humana el goce en su propia entidad, en su interioridad, en todas sus manifestaciones físicas y espirituales. Protegen dimensiones fundamentales del ser humano inherentes a su naturaleza y dignidad. \\
\addlinespace

\textbf{¿Cuál es la diferencia entre derechos personalísimos y atributos de la personalidad?} &
Los atributos de la personalidad son cualidades esenciales que permiten que la persona despliegue su personalidad y entable relaciones jurídicas (como el nombre, el domicilio o la capacidad). Los derechos personalísimos son verdaderos derechos subjetivos que protegen dimensiones fundamentales de la persona (como la vida, el honor o la integridad física). \\
\addlinespace

\textbf{¿Por qué es importante estudiar los derechos personalísimos hoy?} &
Las circunstancias contemporáneas —biotecnologías, big data, privacidad digital— han generado nuevas amenazas a estos derechos que requieren una respuesta jurídica actualizada, orientada a proteger a las personas frente a formas de intromisión y vulneración que no existían en épocas anteriores. \\
\addlinespace

\textbf{¿Cuáles son los caracteres principales de los derechos personalísimos?} &
Son: subjetivos privados; innatos y originarios (reconocidos, no concedidos); vitalicios y necesarios; absolutos (sujeto pasivo universal); inherentes a la persona; intransmisibles; relativamente indisponibles; e irrenunciables. Todos estos caracteres buscan proteger la dignidad humana. \\
\addlinespace

\textbf{¿Cómo se clasifican los derechos personalísimos?} &
En forma tripartita: integridad física (vida, integridad corporal, salud); libertad (concepto que comprende aspectos corporales y espirituales); e integridad espiritual (honor, intimidad, imagen, identidad). \\
\addlinespace

\textbf{¿Dónde se regulan los derechos personalísimos en el Código Civil y Comercial?} &
En el Capítulo 3, Título Primero, Libro Primero (artículos 51 a 62). El artículo 51 abre el capítulo con el principio de inviolabilidad de la persona humana y su derecho al reconocimiento y respeto de su dignidad. \\
\addlinespace

\textbf{¿Cuáles son los medios de protección de los derechos personalísimos?} &
El artículo 52 establece: la prevención, que busca evitar que ocurra el acto dañino, y la reparación, que busca volver al estado anterior o indemnizar. También se prevé la publicación de la sentencia como forma específica de reparación en casos de daño al honor, la intimidad o la identidad. \\
\addlinespace

\textbf{¿Qué significa que los derechos personalísimos son relativamente indisponibles?} &
No pueden ser dispuestos en principio, salvo mediante consentimiento que: no sea contrario a las buenas costumbres; no se presuma (debe ser expresado); sea de interpretación restrictiva; y sea libremente revocable. \\
\addlinespace

\textbf{¿Cómo se protege la libertad en el derecho civil?} &
A través de acciones civiles cuando existe violencia o coacción para celebrar un acto (reclamo de daños y perjuicios). También existen protecciones constitucionales, como el habeas corpus, el amparo y la acción de inconstitucionalidad, más específicas respecto de la libertad física. \\
\addlinespace

\midrule
\multicolumn{2}{p{0.92\textwidth}}{
\textbf{Resumen:} Los derechos personalísimos constituyen un sistema integrado de protección de la dignidad humana en todas sus dimensiones. Fundados en el reconocimiento de que cada persona posee una dignidad inherente, protegen la vida, la integridad física y espiritual, la libertad, el honor, la intimidad y la identidad de todo ser humano. Aunque tienen raíces profundas en la historia del derecho, adquieren una importancia renovada en el contexto contemporáneo, en el que las biotecnologías, la digitalización y el almacenamiento masivo de datos han generado nuevas amenazas a estas esferas fundamentales de la persona. El Código Civil y Comercial supone un avance significativo al sistematizar estos derechos en su Parte General, reconociendo su carácter innato, vitalicio, necesario, absoluto, inherente, intransmisible e irrenunciable. Su protección se articula mediante dos vías principales: la prevención y la reparación.
} \\
\bottomrule
\end{longtable}
\chapter{La muerte: fin de la existencia de la persona humana}
\label{ch:muerte}

\section{Ubicación temática y relevancia práctica}
La regulación civil de la persona humana comienza con el primer momento de su existencia --que en el ordenamiento argentino, en consonancia con las normas constitucionales, se sitúa en el momento de la concepción-- y concluye con el fin de la existencia, es decir, con la muerte. El estudio de esta última etapa presenta dos grandes cuestiones: por un lado, la situación de la muerte propiamente dicha y su acreditación civil; por otro, el instituto de la presunción de fallecimiento, figura jurídica aplicable cuando una persona se encuentra ausente de su domicilio sin que se tengan noticias de ella durante el plazo fijado por la ley. En este capítulo se desarrolla exclusivamente el fin de la existencia por muerte comprobada; la presunción de fallecimiento se estudia en el capítulo~\ref{ch:ausencia}.

\begin{note}
Los distintos apartados de este capítulo se encuentran conectados entre sí. Ante el fallecimiento de una persona es necesario, en primer lugar, comprender qué es la muerte y cómo se comprueba. Cuando no se dispone de un cadáver, corresponde analizar cómo se acredita judicialmente el fallecimiento. Si dos o más personas fallecen simultáneamente, la conmoriencia determina las relaciones sucesorias entre ellas. Finalmente, la muerte produce consecuencias jurídicas concretas sobre el patrimonio y los derechos de la persona fallecida.
\end{note}

El fin de la existencia de la persona humana constituye uno de los pilares del derecho civil: el sistema de herencias, la donación de órganos, los registros civiles, los seguros de vida y las sucesiones dependen de determinar cuándo, cómo y con qué efectos jurídicos se produce la muerte. En el ejercicio profesional del derecho privado resulta indispensable gestionar partidas de defunción, iniciar sucesiones, asesorar a las familias en situaciones de presunción de fallecimiento, promover o impugnar declaraciones judiciales de muerte y conocer los límites de la transmisión hereditaria.

\section{La muerte: concepto, comprobación y normativa}

\subsection{Principio general}

El Código Civil y Comercial de la Nación (en adelante, CCyCN) establece el principio general que rige la finalización de la existencia de la persona humana:

\begin{formula}{Art. 93 CCyCN --- Principio general}
La existencia de la persona humana termina por su muerte.
\end{formula}

El Código Civil anterior contenía una disposición similar, aunque hacía referencia a la ``muerte natural''. El término ``natural'', en la redacción del Código de Vélez Sarsfield, no se oponía a ``provocada'' --como si una persona asesinada no muriera-- sino que aludía a la muerte biológica y, fundamentalmente, procuraba excluir cualquier forma de muerte civil.

\begin{note}
La muerte civil implicaba que el derecho consideraba muerta a una persona por el solo cumplimiento de determinada circunstancia --por ejemplo, una pena, o el ingreso a una comunidad religiosa, especialmente de votos solemnes, situación en la cual la persona ``moría al mundo'' para entrar en religión--. Esta institución nunca tuvo aplicación efectiva en el derecho argentino desde el Código Civil anterior, y el CCyCN simplifica la cuestión disponiendo, sin más, que la existencia de la persona termina con su muerte.
\end{note}

\subsection{¿Qué es la muerte?}

\begin{definition}{Muerte}
Cese irreversible de las funciones corporales; fin de la unidad de cuerpo y alma. El momento exacto en que la muerte ocurre permanece, en gran medida, en el plano del misterio: el derecho civil no define en qué instante preciso se produce, sino que regula el modo en que se \textbf{comprueba} su acaecimiento, comprobación que siempre resulta posterior al hecho biológico en sí.
\end{definition}

El momento exacto de la muerte de una persona resulta, por lo tanto, muy difícil de determinar. Aquello de lo que se ocupa el derecho civil es de su comprobación.

\subsection{Comprobación de la muerte}

\begin{formula}{Art. 94 CCyCN --- Comprobación de la muerte}
La comprobación de la muerte queda sujeta a los estándares médicos aceptados, aplicándose la legislación especial en el caso de ablación de órganos del cadáver.
\end{formula}

La comprobación de la muerte se realiza conforme a estándares médicos que van variando a lo largo del tiempo, a medida que avanza la ciencia, y que permiten detectar el cese irreversible de las funciones corporales --especialmente las cardiorrespiratorias y cerebrales, que son las que comandan el organismo humano.

Se habla de ``comprobación'' precisamente porque esta ocurre con posterioridad al hecho de la muerte. Lo que se obtiene es una certeza moral, no una certeza apodíctica: las certezas sobre la muerte se basan siempre en observaciones empíricas que, en principio, podrían tener alguna falla, pero que --con el acompañamiento de los estándares médicos-- se procuran hacer lo más rigurosas posible, de manera de no dejar duda alguna acerca del fallecimiento de la persona. De allí la confianza depositada en la intervención médica.

\subsection{Instrumentos para acreditar la muerte}

\begin{formula}{Ley 26.413 --- Registro del Estado Civil y Capacidad de las Personas}
Regula la forma en que deben ser expedidos el certificado médico de defunción y la partida de defunción. El certificado médico es expedido por un profesional de la salud siempre sobre la base de la presencia del cadáver. La partida de defunción es el instrumento público en el que consta el asiento en el registro civil, con la fecha, la hora y los demás elementos que rodean a la muerte.
\end{formula}

La normativa distingue con claridad dos instrumentos:

\begin{table}[H]
\centering
\caption{Certificado médico de defunción y partida de defunción}
\begin{tabularx}{\textwidth}{X X}
\toprule
\textbf{Certificado médico de defunción} & \textbf{Partida de defunción} \\
\midrule
Expedido por un profesional de la salud & Instrumento público del Registro Civil \\
Siempre requiere la presencia del cadáver & No requiere la presencia del cadáver (se basa en el certificado) \\
Documento de base para inscribir la muerte & Constancia del asiento en el Registro del Estado Civil \\
Contiene fecha, hora y elementos que rodean la muerte & Regulada también por la Ley 26.413 \\
\bottomrule
\end{tabularx}
\end{table}

\section{Donación de órganos y criterios de muerte encefálica}

\subsection{Consentimiento presunto}

Cuando una persona fallece sin haber expresado su voluntad respecto de la donación de órganos, la ley presume que es dadora.

\begin{formula}{Ley 27.447 (2018) --- Trasplante de Órganos y Tejidos}
Establece que cuando una persona fallece sin haber expresado su voluntad, se presume que es dadora de órganos. La negativa debe ser expresa, realizada ante los organismos señalados en la ley. El artículo 36 dispone que la muerte puede certificarse con la confirmación del cese irreversible de las funciones circulatorias o encefálicas. El artículo 37 remite a un protocolo fijado por el Ministerio de Salud con asesoramiento del INCUCAI (Resolución 716/2019), estableciendo que la certificación debe ser realizada por dos médicos, uno de los cuales debe ser neurólogo o neurocirujano y no puede integrar el equipo de ablación.
\end{formula}

\begin{note}
La Ley 27.447 presume que toda persona es donante de órganos, salvo que haya expresado directamente una negativa ante los organismos señalados en la norma. Esta solución ha generado controversias, tanto por suprimir el rol de la familia en la decisión como por presumir una voluntad que, para algunos, debería quedar estimulada en lugar de presumida.
\end{note}

\subsection{Antecedentes legislativos}

\begin{formula}{Ley 21.541 (1977) --- Primera ley de trasplantes}
Definió el fallecimiento como la disfunción total e irreversible de las funciones encefálicas, verificada por un equipo médico. La reglamentación establecía los signos mínimos que debían presentarse en su totalidad.
\end{formula}

\begin{formula}{Ley 24.193 (1993) --- Ley de trasplantes anterior a 2018}
Establecía que la muerte se constataba por: ausencia irreversible de respuesta cerebral con pérdida absoluta de conciencia; ausencia de respiración espontánea; ausencia de reflejos cefálicos; pupilas fijas no reactivas; e inactividad encefálica comprobada por medios técnicos o instrumentales (electroencefalograma plano y constante durante seis horas). No era necesario comprobar la inactividad encefálica en caso de paro cardiorrespiratorio total e irreversible.
\end{formula}

Esta definición de la Ley 24.193 estuvo vigente entre 1993 y 2018.

\subsection{Criterios actuales y Protocolo 2019}

A partir de 2018, la ley opta por no incorporar en su texto una descripción tan detallada de los signos, remitiéndola a un protocolo. Se dispone que el fallecimiento de una persona puede certificarse con la confirmación del cese irreversible de las funciones circulatorias o encefálicas --cualquiera de los dos supuestos, tanto el paro cardiorrespiratorio como el encefálico.

\begin{important}
La palabra clave del criterio vigente es la \textbf{irreversibilidad}. Si existe la posibilidad de reversión --como ocurre con un paro cardíaco que puede revertirse mediante reanimación-- no hay persona fallecida, sino una persona cuya función circulatoria puede ser terapéuticamente reanimada. El fallecimiento, en sentido propio, es un cese irreversible de las funciones.
\end{important}

Respecto de la certificación, el artículo 37 de la Ley 27.447 remite a un protocolo fijado por el Ministerio de Salud con asesoramiento del INCUCAI, concretado en la Resolución 716/2019. Dicho protocolo puede variar, en atención a la necesidad de actualización técnica de los medios biotecnológicos utilizados para el diagnóstico.

\subsection{Garantías y críticas al régimen vigente}

\begin{important}
Un sector de la doctrina ha criticado el sistema vigente por no establecer en el propio texto legal algunas garantías mínimas --como las que preveía la ley anterior--, delegando en una reglamentación un hecho tan decisivo como la certificación de la muerte. Esta cuestión ha merecido discusión doctrinaria; no obstante, la ley mantiene como criterio decisivo el del cese irreversible de las funciones.
\end{important}

Para la profundización de este debate cabe remitir a los trabajos del Dr.\ José Tobías y del Dr.\ Carlos Muñiz, quienes han escrito ampliamente sobre la determinación del momento de la muerte.

\section{Conmoriencia y prueba supletoria del fallecimiento}

\subsection{Conmoriencia}

\begin{formula}{Art. 95 CCyCN --- Conmoriencia}
Se presume que las personas que fallecen en un desastre común mueren al mismo tiempo. Salvo prueba en contrario, es una presunción iuris tantum.
\end{formula}

\begin{note}
A lo largo del tiempo han existido distintas tesis sobre la preeminencia --es decir, sobre quién moría primero, si los más jóvenes, los más ancianos, o de acuerdo con el sexo--. Los códigos civiles argentinos, ya desde la época de Vélez Sarsfield, establecen que todas las personas que se encuentran en esas circunstancias mueren al mismo tiempo, admitiéndose prueba en contrario.
\end{note}

\begin{example}{Accidente de tránsito}
Fallecen el padre y el hijo en un accidente de tránsito; sobrevive la madre. Se supone que no existen otros herederos.

Si el hijo fallece \textbf{después} del padre: hereda del padre y, al morir el propio hijo, lo hereda la madre.

Si el hijo fallece \textbf{antes} que el padre: el padre, al morir, no tendría herederos directos, y serían sus propios ascendientes quienes heredarían.

\textbf{Solución legal (Art. 95 CCyCN):} se presume que ambos fallecieron al mismo tiempo. No hay transmisión de derechos entre ellos, lo cual evita discusiones sobre quién murió primero y alteraciones en los órdenes sucesorios.
\end{example}

\subsection{Prueba supletoria del fallecimiento}

El segundo supuesto especial que corresponde analizar es el previsto en el artículo 98, segunda parte, del CCyCN, que no debe confundirse con la presunción de fallecimiento.

\begin{formula}{Art. 98 CCyCN --- Prueba supletoria del fallecimiento}
Cuando el cadáver no es hallado o no puede ser identificado, y la desaparición se produjo en circunstancias tales que la muerte debe tenerse por cierta, el juez puede tener por comprobada la muerte y disponer la correspondiente inscripción en el registro.
\end{formula}

Para inscribir una muerte se requiere, en principio, un certificado médico. Sin embargo, cuando el profesional de la salud no dispone de cadáver alguno, o no puede identificarlo, dicha vía resulta imposible. Ante la certeza de que la persona ha fallecido, se habilita entonces un trámite judicial.

El criterio decisivo para optar por este artículo, en lugar de la presunción de fallecimiento, es la certeza:

\begin{table}[H]
\centering
\caption{Prueba supletoria (Art. 98) frente a la presunción de fallecimiento}
\begin{tabularx}{\textwidth}{X X}
\toprule
\textbf{Prueba supletoria --- Art. 98} & \textbf{Presunción de fallecimiento} \\
\midrule
Hay \textbf{certeza} de que la persona falleció & Hay \textbf{duda o incertidumbre} sobre si falleció \\
Cadáver no hallado o no identificado & Persona ausente sin noticias durante el tiempo fijado por ley \\
Procedimiento judicial más breve & Procedimiento judicial mucho más largo \\
El juez ordena la inscripción de la muerte & Se declara presuntamente fallecida a la persona \\
\bottomrule
\end{tabularx}
\end{table}

\subsection{Ejemplos desarrollados}

\begin{example}{Avión que se desintegra en el aire (Art. 98)}
Un avión de Air France que se desintegró en el aire (vuelo AF 447), o el avión de LAPA en el Río Uruguay que se desintegró en el aire: no hubo sobrevivientes, y resulta imposible que los hubiera habido. Ante ese hecho, se aplica el artículo 98 --prueba supletoria-- porque existe certeza de que las personas fallecieron.
\end{example}

\begin{example}{Los Andes (presunción de fallecimiento)}
El caso del avión de la película ``Viven'', cuyos pasajeros --uruguayos-- caen en la cordillera de los Andes: no solo no se los encontraba, sino que tampoco se sabía si estaban muertos o no. No podía afirmarse que esas personas habían fallecido; correspondía, por lo tanto, aplicar la presunción de fallecimiento y esperar los plazos propios de ese instituto. De hecho, fueron encontrados a los tres meses.
\end{example}

\begin{example}{Alpinista en precipicio inaccesible (Art. 98)}
Un alpinista cae en un precipicio de acceso imposible, acompañado por compañeros de montañismo que certifican que la persona llegó a ese punto, tuvo un problema, se resbaló y cayó, resultando imposible que haya sobrevivido e imposible acceder al lugar. Con esos testimonios y las constancias del viaje puede realizarse la prueba supletoria: no corresponde permanecer en la incertidumbre, pues la persona falleció. El juez, teniendo por comprobada la muerte, dispone la inscripción en el registro.
\end{example}

\section{Efectos jurídicos de la muerte}

\subsection{Efectos extrapatrimoniales}

La muerte produce, en el campo extrapatrimonial, la extinción de los atributos de la persona: el nombre y la capacidad. Del domicilio subsiste el domicilio convencional --aquel que se elige para un contrato-- que se transmite a los sucesores conforme al artículo 1024 CCyCN. Se extingue el estado civil, pero subsisten los vínculos a los fines sucesorios: si existía cónyuge, y respecto del padre, la madre y los hijos, el parentesco se mantiene a esos efectos.

\subsection{Protección jurídica post mortem}

Se extinguen los derechos personalísimos, pero subsiste una protección jurídica sobre el honor de la persona fallecida, su intimidad y la inviolabilidad de sus papeles privados. Existe, asimismo, una protección jurídica del cadáver, orientada a que este sea tratado con dignidad.

\begin{note}
El cadáver posee un estatuto especial: no es la persona ni es el cuerpo, pero constituye una prolongación de la persona, y por ello existe un cuidado de esta y de su memoria. La protección de la memoria de los difuntos se vincula también con su intimidad y con la inviolabilidad de sus papeles privados.
\end{note}

\subsection{Efectos patrimoniales}

En el campo patrimonial, los derechos patrimoniales de la persona fallecida se transmiten a sus herederos conforme a las reglas dispuestas por el CCyCN, que pueden ser de dos tipos: testamentarias --cuando la propia persona ha determinado quiénes son sus herederos, existiendo porciones legítimas correspondientes a herederos forzosos sobre las cuales no puede disponerse, y una parte disponible según existan descendientes o solo ascendientes-- o, en su defecto, ab intestato.

\subsection{Derechos no transmisibles}

\begin{formula}{Art. 1024 CCyCN --- Efectos de los contratos: sucesores}
Las obligaciones que son inherentes a las personas no son transmisibles por causa de muerte.
\end{formula}

\begin{formula}{Art. 2280 CCyCN --- Situación de los herederos}
No se transmiten por sucesión los derechos que la ley dispone que no sean transmisibles.
\end{formula}

No se transmiten, por tanto, los derechos que la ley dispone como no transmisibles. Un ejemplo de ello son las obligaciones inherentes a las personas, que no son transmisibles por causa de muerte.

\begin{example}{El pintor famoso}
Se contrata a un pintor famoso para que realice el retrato de un familiar por cinco mil dólares. El pintor fallece. Su hija, también pintora, ofrece realizar el retrato en su lugar. La obligación no se transmite, porque era inherente a la persona --se habían considerado las cualidades específicas del pintor contratado-- y su transmisión resultaría incompatible con la naturaleza de la obligación pactada.
\end{example}

Los derechos reales de usufructo y de habitación constituyen otros ejemplos típicos de derechos que no se transmiten por causa de muerte: por su propia constitución, se extinguen con la muerte del beneficiario. La muerte produce, por regla general, la transmisión de los derechos patrimoniales, con excepción de aquellos que no se transmiten.

\section{Síntesis del capítulo}
El capítulo ha desarrollado las cuestiones relativas al fin de la existencia de la persona humana: la muerte, su comprobación, el caso de la conmoriencia, la prueba supletoria del fallecimiento y los efectos jurídicos de la muerte.

% ------------------------------------------------------------
% CUADRO CORNELL
% ------------------------------------------------------------
\section{Cuadro Cornell --- Repaso de la unidad}
\begin{longtable}{
    >{\raggedright\arraybackslash}p{0.30\textwidth}
    >{\raggedright\arraybackslash}p{0.60\textwidth}
}
\caption{Cuadro Cornell --- Fin de la existencia de la persona humana} \\
\toprule
\textbf{Preguntas / palabras clave} & \textbf{Notas / respuestas} \\
\midrule
\endfirsthead

\toprule
\textbf{Preguntas / palabras clave} & \textbf{Notas / respuestas} \\
\midrule
\endhead

\textbf{¿Cuándo termina la existencia de la persona? (Art. 93)} &
Con la muerte. El CCyCN simplifica la redacción anterior eliminando la referencia a la ``muerte natural'' y la figura de la muerte civil, ya inexistente en el derecho argentino desde el Código de Vélez. \\

\textbf{¿Qué es la muerte?} &
El cese irreversible de las funciones corporales. El momento exacto en que ocurre permanece en el plano del misterio; lo que el derecho regula es su comprobación posterior mediante estándares médicos. \\

\textbf{Art. 94 --- ¿Cómo se comprueba la muerte?} &
La comprobación queda sujeta a los estándares médicos aceptados, que evolucionan con la ciencia. Para la ablación de órganos rige la legislación especial (Ley 27.447 y Protocolo 2019). \\

\textbf{Certificado médico vs.\ partida de defunción} &
El certificado médico es expedido por un profesional de la salud, siempre ante la presencia del cadáver, y constituye el documento base. La partida de defunción es el instrumento público del Registro Civil donde consta la inscripción de la muerte (Ley 26.413). \\

\textbf{Ley 27.447 --- ¿Todos somos donantes de órganos?} &
Sí, por presunción legal. Quien no quiera serlo debe expresar una negativa ante los organismos designados. Se eliminó el rol de la familia que existía bajo la Ley 24.193 (Arts.\ 36 y 37). \\

\textbf{¿Qué exigía la Ley 24.193 (1993--2018) para la muerte encefálica?} &
Ausencia irreversible de respuesta cerebral; ausencia de respiración espontánea; ausencia de reflejos cefálicos; pupilas fijas no reactivas; encefalograma plano y constante durante seis horas. El paro cardiorrespiratorio total e irreversible eximía de la prueba encefálica. \\

\textbf{Ley 27.447 --- ¿cuál es el criterio actual de muerte?} &
Cese irreversible de las funciones circulatorias o encefálicas. Los signos específicos y los plazos se remiten al Protocolo del Ministerio de Salud (Res.\ 716/2019), elaborado con el INCUCAI. Debe certificarlo un equipo de dos médicos, uno de ellos neurólogo o neurocirujano ajeno al equipo de ablación. \\

\textbf{¿Qué es la conmoriencia? (Art. 95)} &
Presunción iuris tantum de que dos personas que fallecen en un desastre común mueren al mismo tiempo. Evita discusiones sobre quién murió primero y posibles alteraciones en los órdenes sucesorios; admite prueba en contrario. \\

\textbf{Prueba supletoria del fallecimiento (Art. 98, segunda parte)} &
Procede cuando el cadáver no es hallado o no puede identificarse, pero existe certeza de que la persona falleció. Un juez --no un médico-- ordena la inscripción de la muerte en el registro. Ejemplos: aviones desintegrados, alpinistas en precipicios inaccesibles. \\

\textbf{Diferencia entre el Art. 98 y la presunción de fallecimiento} &
Art.\ 98: hay certeza de la muerte, el procedimiento es más breve y el juez ordena la inscripción. Presunción de fallecimiento: hay duda o incertidumbre, el procedimiento es mucho más largo y se declara ``presuntamente fallecida'' a la persona. \\

\textbf{Efectos extrapatrimoniales de la muerte} &
Se extinguen el nombre, la capacidad, el estado civil y los derechos personalísimos. Subsisten el domicilio convencional (para contratos), los vínculos familiares a fines sucesorios y la protección jurídica del honor, la intimidad y la inviolabilidad del cadáver. \\

\textbf{Efectos patrimoniales de la muerte} &
Se transmiten a los herederos los derechos patrimoniales, conforme a las reglas testamentarias o ab intestato del CCyCN, respetando las porciones legítimas de los herederos forzosos. \\

\textbf{Derechos NO transmisibles por causa de muerte} &
Las obligaciones inherentes a la persona (intuitu personae) y los derechos reales de usufructo y de habitación (Arts.\ 1024 y 2280 CCyCN). Ejemplo: si se contrata a un pintor famoso y este fallece, su hija no está obligada a cumplir ni tiene derecho a exigir el pago. \\

\midrule
\multicolumn{2}{p{\dimexpr\textwidth-2\tabcolsep}}{
\textbf{Resumen:} El fin de la existencia de la persona humana se produce con la muerte, cuya única definición jurídica en el CCyCN es el cese irreversible de las funciones corporales (Art.\ 93). Como ese momento exacto es de difícil determinación, el derecho regula su comprobación: en el caso ordinario, mediante el certificado médico de defunción --que exige la presencia del cadáver-- y la posterior partida de defunción en el Registro Civil, ambos regulados por la Ley 26.413. Para la ablación de órganos, los criterios son más rigurosos y están regulados por la Ley 27.447 (2018) y el Protocolo del INCUCAI de 2019: basta el cese irreversible de las funciones circulatorias o encefálicas, certificado por dos médicos, uno de los cuales debe ser neurólogo o neurocirujano y ajeno al equipo de trasplante. Bajo esta ley, toda persona es donante presunta salvo negativa expresa. Cuando dos personas fallecen en un desastre común sin poder determinarse quién murió primero, la conmoriencia (Art.\ 95) presume que ambas murieron simultáneamente, evitando distorsiones en los órdenes sucesorios. Si el cadáver no es hallado pero hay certeza de la muerte, el juez puede tenerla por acreditada y ordenar su inscripción mediante la prueba supletoria del Art.\ 98, que se distingue de la presunción de fallecimiento por la ausencia de incertidumbre. Finalmente, la muerte produce efectos extrapatrimoniales --extinción del nombre, la capacidad y el estado civil, con subsistencia de los vínculos familiares a fines sucesorios y protección jurídica post mortem del honor y la intimidad-- y patrimoniales: transmisión del patrimonio a los herederos según las reglas del CCyCN, con excepción de los derechos inherentes a la persona (usufructo, habitación, obligaciones intuitu personae), que se extinguen con la muerte del titular.
} \\
\bottomrule
\end{longtable}
\chapter[Marco conceptual e histórico]{Ausencia y presunción de fallecimiento: marco conceptual e histórico}
\label{ch:ausencia}

El fin de la existencia de la persona admite, además de la muerte comprobada, una vertiente de incertidumbre sobre su vida o su muerte, regulada en el Código Civil y Comercial de la Nación (CCCN) mediante dos institutos diferenciados: la \textbf{ausencia} (arts.~79 a 84) y la \textbf{presunción de fallecimiento} (arts.~85 a 92). Los capítulos de esta sección tratan, sucesivamente, el marco conceptual e histórico, la ausencia, los tres casos de presunción de fallecimiento, el procedimiento y la sentencia, y las garantías y efectos de la reaparición.

\begin{note}
Tarde o temprano resulta habitual encontrarse con un cliente cuyo familiar desapareció: un accidente en ruta, una catástrofe natural, o una persona que simplemente dejó de dar señales de vida. Saber distinguir la ausencia de la presunción de fallecimiento no es un detalle académico: es la diferencia entre proteger el patrimonio familiar desde el primer mes o perderlo por falta de acción.

Estos institutos permiten: (1) asesorar a herederos sobre cuándo y cómo iniciar acciones judiciales; (2) gestionar la curatela de bienes de personas desaparecidas; (3) representar acreedores que necesitan cobrar deudas pendientes; (4) intervenir en procesos sucesorios abiertos tras una declaración de fallecimiento presunto; (5) proteger los derechos del ausente ante una eventual reaparición, gestionando la prenotación registral.

En síntesis, estas reglas existen para que el derecho privado pueda dar una respuesta ordenada, justa y temporalmente acotada a uno de los dramas más dolorosos de la vida real: no saber si un ser querido está vivo o muerto.
\end{note}

\textbf{¿Cómo se relacionan los temas?} Puede pensarse el proceso como una escalada progresiva: si no se tienen noticias de una persona y sus bienes quedan desprotegidos, corresponde solicitar primero la \textbf{declaración de ausencia} (arts.~79 a 84) para que un curador los cuide. Si transcurren tres años sin novedades, puede iniciarse la \textbf{presunción de fallecimiento} (arts.~85 a 92), que abre la sucesión pero con garantías por si la persona reaparece. Todo el proceso protege tanto a quienes esperan como a quien podría volver.

\section{Introducción: ausencia y presunción de fallecimiento}

El fin de la existencia de la persona está regulado en el Código Civil y
Comercial a través de dos figuras: la \textbf{ausencia}, por un lado, y la
\textbf{presunción de fallecimiento}, por el otro.

Ausencia y presunción de fallecimiento son dos figuras distintas. Es
habitual que se confundan bajo la expresión «ausencia con presunción de
fallecimiento», lo cual no resulta exacto: se trata de institutos
diferentes. Por supuesto, la presunción de fallecimiento supone que existe
una ausencia de la persona, pero lo que corresponde estudiar es,
específicamente, por un lado la \textbf{ausencia} y por el otro la
\textbf{presunción de fallecimiento}.

\section{El problema jurídico de fondo}

El problema jurídico de fondo puede formularse como la pregunta: ¿cuál es
la respuesta jurídica ante la ausencia prolongada de una persona de su
domicilio?

Esta cuestión ya encontraba explicación en el contexto de las guerras: la
persona que quedaba cautiva y de la que no se sabía nada. ¿Qué sucedía con
sus bienes? ¿Quién podía tomar medidas conservatorias? También pueden
pensarse situaciones catastróficas más recientes: la desaparición de un
avión, el tsunami que hubo en Tailandia a comienzos del año 2002, un
incendio, o cualquier catástrofe en la que una persona desaparece y no se
sabe si está viva o no.

Se produce una \textbf{incertidumbre}, y esa incertidumbre no puede durar
eternamente, porque para el derecho civil existe una cantidad de
relaciones y situaciones jurídicas que ameritan respuesta y toma de
decisiones.

No existe una institución estatal que actúe de oficio ante la mera falta
de contacto de una persona: este es un \textbf{sistema de derecho civil,
de derecho privado}, en el cual los particulares que tienen algún tipo de
interés en los bienes o en las relaciones jurídicas acuden a un proceso
judicial, ya sea para declarar la ausencia o para declarar la presunción
de fallecimiento.

Esto supone preguntarse: ¿cuándo se considera que una persona ha
fallecido? ¿Cuándo resulta razonable asumirlo así? ¿Quién puede pedirlo?
¿Qué sucede si la persona reaparece? ¿Qué resguardos jurídicos se toman en
el camino por si reaparece? Ese es el tema de fondo que estructura ambos
institutos.

\begin{important}
No debe confundirse la ausencia con la presunción de fallecimiento. Si
bien están ubicadas inmediatamente una al lado de la otra en el Código
Civil y Comercial, y tienen elementos comunes, son institutos distintos.
\end{important}

\section{Antecedentes históricos}

Los antecedentes de esta institución provienen del \textbf{Código de
Vélez Sársfield}, que regulaba la ausencia con presunción de fallecimiento
en los artículos 110 a 625. La doctrina señala con claridad cómo en Vélez
Sársfield se produce una \textbf{confusión de elementos}:

Por un lado, existen antecedentes del \textbf{derecho germánico}, donde
existía claramente el instituto de la declaración de muerte, que Vélez
recibe a través de Freitas. Por el otro lado, el \textbf{derecho
francés}, que no conocía esa declaración de muerte y simplemente regulaba
la ausencia, sin poner fin a la existencia de la persona. Vélez también
reconoce esta figura de la ausencia, de modo que el instituto queda
regulado con elementos mixtos.

Tal como fue regulado por Vélez, la institución finalmente recibió
críticas. En 1984 una ley ómnibus —una ley que reunía muchos temas en su
contenido—, la \textbf{Ley 14.394}, reguló por un lado la ausencia y por
otro lado la presunción de fallecimiento: dos instituciones distintas,
clarificando y resolviendo lo que Vélez había planteado. El régimen de
Vélez, entonces, quedó reemplazado por el sistema de la Ley 14.394.

\begin{important}
\textbf{CCCN — Estructura normativa actual.} El nuevo Código Civil y
Comercial deroga la Ley 14.394 en estas partes, pero sigue sus
lineamientos e incorpora tanto la ausencia como la presunción de
fallecimiento al articulado del Título I, Libro I:
\begin{itemize}
    \item \textbf{Ausencia} — Capítulo 6, Arts. 79-84 (no implica fin de la existencia de la persona).
    \item \textbf{Presunción de Fallecimiento} — Capítulo 7, Arts. 85-92 (implica inscripción en el Registro Civil y apertura de la sucesión).
\end{itemize}
Metodológicamente, el CCCN los incorpora con muy pocos cambios respecto de
la ley anterior.
\end{important}

% ======================================================================
% CAPÍTULO 2
% ======================================================================
\chapter{Ausencia (Arts. 79-84 CCCN)}

\section{Concepto de ausencia}

La ausencia es una institución relativamente sencilla. Si la persona
desapareció de su domicilio sin tenerse noticias de ella y sin dejar
apoderado, o si el apoderado no tiene atribuciones suficientes, puede
asignarse un \textbf{curador a los bienes}.

\begin{definition}{Ausencia}
La ausencia, en el proceso del Art. 79 CCCN, no busca declarar a la
persona presuntamente fallecida. Simplemente busca \textbf{garantizar la
adopción de medidas sobre los bienes} de una persona de la que no se
tienen noticias.
\end{definition}

\begin{example}{¿Qué NO es ausencia?}
Si alguien quedó varado en Qatar durante la pandemia, se tienen noticias
de esa persona: se sabe que está varada y que no puede regresar a su
domicilio. Eso no configura una ausencia. La ausencia es la situación de
la persona que desaparece, no se conecta a internet, no usa su tarjeta de
crédito, no pasa por migraciones, no tiene registro en su correo
electrónico, no tiene ningún tipo de actividad, y por lo tanto nadie de
las personas conocidas tiene noticias de ella.
\end{example}

\section{Legitimación y juez competente}

\begin{formula}{Art. 80 CCCN — Legitimación activa}
El Art. 80 establece quiénes pueden iniciar la acción: el Ministerio
Público y las personas que tengan interés legítimo en los bienes (por
ejemplo, los acreedores o los sucesores).
\end{formula}

Respecto del juez competente:

\begin{table}[H]
\centering
\begin{tabularx}{\textwidth}{>{\raggedright\arraybackslash}X >{\raggedright\arraybackslash}X}
\toprule
\textbf{Situación} & \textbf{Juez competente} \\
\midrule
Persona con domicilio en Argentina & Juez del domicilio \\
Sin domicilio en el país o lugar desconocido & Juez del lugar donde están los bienes que hay que cuidar \\
Bienes en distintas jurisdicciones & El primero que intervenga (principio de prevención) \\
\bottomrule
\end{tabularx}
\caption{Juez competente en el proceso de ausencia}
\end{table}

\section{Procedimiento y curatela}

Se cita al ausente por \textbf{cinco días por edictos}. Los edictos son
los avisos que se publican en el Boletín Oficial y en diarios de
circulación importante, mediante los cuales se convoca a la persona a que
se haga presente. Si no comparece, se le da intervención al
\textbf{defensor oficial}, que interviene en representación del ausente,
y pueden tomarse medidas precautorias para los bienes.

La prueba se orienta a demostrar que no hay noticias de esta persona: se
libran distintos oficios, se solicitan testigos, y se determina si existe
o no ausencia. El Ministerio Público procura defender al presunto
ausente, evaluando la verosimilitud de la situación alegada.

Determinada la ausencia, se nombra un curador. El procedimiento del
curador comprende actos de \textbf{conservación y administración de los
bienes}. No puede realizar actos extraordinarios: no puede vender ni
hipotecar un inmueble, porque en ese caso estaría comprometiendo el
patrimonio. Los actos extraordinarios que resulten de necesidad evidente e
impostergable requieren \textbf{autorización judicial}.

Los frutos —por ejemplo, la renta de un inmueble que se alquila— se
destinan al \textbf{sostenimiento de los familiares del ausente}.

\section{Conclusión de la curatela}

La curatela concluye:

\begin{enumerate}[label=(\alph*)]
    \item Porque reaparece el ausente, personalmente o por apoderado.
    \item Porque se constata su muerte: ya sea porque aparece el cadáver y se hace la inscripción del fallecimiento y la correspondiente partida, o bien porque se declara el fallecimiento presunto judicialmente, instituto que se analiza a continuación.
\end{enumerate}

% ======================================================================
% CAPÍTULO 3
% ======================================================================
\chapter{Presunción de fallecimiento (Arts. 85-92 CCCN)}

El sistema del Código es igual al sistema de la Ley 14.394 y se estructura
en torno a \textbf{tres casos}: un caso ordinario y dos casos
extraordinarios —el extraordinario genérico y el extraordinario
específico—.

\section{Caso ordinario (Art. 85 CCCN)}

\begin{formula}{Art. 85 CCCN — Caso ordinario}
La persona que se ausenta de su domicilio sin que se tengan noticias de
ella por el término de \textbf{tres (3) años}. Los elementos son: la
ausencia sin noticias durante tres años. No hay ningún elemento
particular adicional.
\end{formula}

El plazo es de \textbf{tres años}. Esto se vincula con situaciones como,
por ejemplo, una persona que desaparece —quizás víctima de un homicidio
no resuelto— y de la que no se sabe bien qué pasó. Al cabo de tres años
de incertidumbre, los familiares pueden solicitar su fallecimiento
presunto.

No es obligatorio: el cumplimiento del plazo de tres años no impone
iniciar el proceso de manera obligatoria. Esto se relaciona con la
certeza de las relaciones jurídicas. Dado que en la mayoría de los casos
se trata de un familiar y todo está rodeado de afectos y conmociones, hay
quienes siguen buscando y pueden dejar pasar cinco o seis años. Lo que
debe demostrarse en el proceso es que los supuestos legales se han
cumplido.

\begin{important}
\textbf{La última noticia.} El plazo de tres años —o el que corresponda
según el caso— se cuenta desde la \textbf{última noticia} que se tuvo de
la persona. Por ejemplo: si alguien cruzó por migraciones en una fecha
determinada por el paso fronterizo de Paso de los Libres, esa es la
última noticia que se tuvo. A partir de allí se cuentan los plazos.
\end{important}

\section{Caso extraordinario genérico (Art. 86, 1.\textsuperscript{er} párr.)}

El caso extraordinario se denomina así porque acontece algún hecho que
marca una \textbf{mayor probabilidad} de que haya habido un deceso de la
persona. El caso genérico comprende situaciones de incendio, terremoto,
acción de guerra u \textit{otro suceso semejante}, o bien la
participación de la persona en una actividad que implique el mismo
riesgo.

\begin{formula}{Art. 86, 1.\textsuperscript{er} párr. — Caso extraordinario genérico}
Si la persona hubiese estado en el lugar de un incendio, terremoto,
acción de guerra u otro suceso semejante susceptible de ocasionar la
muerte, o hubiese participado de una actividad que implique el mismo
riesgo, y no se tuviese noticias de ella durante \textbf{dos (2) años}
desde el día en que el suceso ocurrió o pudo haber ocurrido.
\end{formula}

\begin{example}{Casos de extraordinario genérico}
\begin{itemize}
    \item El alpinista que estaba escalando el Aconcagua o el Everest y dejó de dar señales.
    \item El tsunami de Tailandia.
    \item Un incendio, un terremoto, una acción bélica.
    \item Cualquier suceso que resulte muy probable que haya provocado la muerte.
\end{itemize}
\end{example}

En este caso, el plazo se adelanta a \textbf{dos años}, contados desde el
día en que el suceso ocurrió o pudo haber ocurrido. Por ejemplo, si se
trata de un alpinista y el último día que consta en un registro que
estaba subiendo fue el 10 de mayo de 2020, desde allí se cuentan los dos
años.

\section{Caso extraordinario específico (Art. 86, 2.\textsuperscript{do} párr.)}

El caso extraordinario específico es el supuesto más detallado por el
Código: el de un \textbf{buque naufragado o aeronave perdida}. Aquí el
acontecimiento resulta mucho más probable como causa de muerte, y el
plazo es de \textbf{seis meses} desde que el suceso ocurrió o pudo
ocurrir.

\begin{formula}{Art. 86, 2.\textsuperscript{do} párr. — Caso extraordinario específico}
Si el presunto fallecido estuvo en un buque naufragado o aeronave
siniestrada o perdida, y no se tuviese noticias de su suerte durante
\textbf{seis (6) meses} desde el día en que el suceso ocurrió o pudo
haber ocurrido.
\end{formula}

\section{Tabla comparativa de los tres casos}

\begin{table}[H]
\centering
\begin{tabularx}{\textwidth}{>{\raggedright\arraybackslash}X c c >{\raggedright\arraybackslash}X}
\toprule
\textbf{Tipo de caso} & \textbf{Artículo} & \textbf{Plazo} & \textbf{Circunstancia} \\
\midrule
Ordinario & Art. 85 & 3 años & Sin noticias \\
Extraordinario genérico & Art. 86, 1.\textsuperscript{er} párr. & 2 años & Catástrofe / actividad de riesgo \\
Extraordinario específico & Art. 86, 2.\textsuperscript{do} párr. & 6 meses & Buque naufragado / aeronave perdida \\
\bottomrule
\end{tabularx}
\caption{Comparación de los tres casos de presunción de fallecimiento}
\end{table}

\section{Casos reales ilustrativos}

\begin{example}{Caso ``Los Andes'' (1972) — Aeronave perdida}
En 1972, un avión militar uruguayo iba rumbo a Chile y cayó en la mitad
de la cordillera de los Andes. Durante 72 días se buscó a sus ocupantes
hasta que dos supervivientes lograron cruzar hacia el lado chileno y
dieron aviso. Esas personas estuvieron desaparecidas durante 72 días. Si
hubiera existido necesidad de declarar su fallecimiento, habría
correspondido el caso extraordinario específico (buque/aeronave), por lo
que hubiera sido necesario esperar seis meses.
% TODO: contenido incompleto o ambiguo en las notas originales — el apunte indica que "los encontraron antes de los 62 días", cifra que no coincide con los "72 días" mencionados previamente en el mismo párrafo; se conserva tal como figura en el material original.
Como los encontraron antes de los 62 días, ese plazo no llegó a cumplirse.
\end{example}

\begin{example}{Caso Malaysia Airlines MH370 (2014) — Aeronave perdida}
En 2014 desapareció un avión de Malaysia Airlines y nunca se supo con
certeza dónde terminó. El último contacto se toma como fecha para contar
los seis meses. Quien tuviera un familiar en ese vuelo debía esperar, como
mínimo, seis meses desde ese último contacto para poder iniciar el
proceso de presunción de fallecimiento bajo el caso extraordinario
específico.
\end{example}

% ======================================================================
% CAPÍTULO 4
% ======================================================================
\chapter{Procedimiento y sentencia}

\section{Legitimación y juez competente}

\begin{formula}{Art. 87 CCCN — Legitimación activa}
Cualquier persona que tenga un derecho subordinado a la muerte puede
pedir la declaración del fallecimiento presunto y debe justificar los
extremos legales: la ausencia, la falta de noticias y las diligencias
realizadas para averiguar la existencia del ausente.
\end{formula}

El \textbf{juez competente} es el juez del domicilio del ausente. Si bien
las normas procesales corresponden a cada provincia, en este caso el
Código Civil y Comercial las incorpora para dar una
\textbf{regulación común a todo el país}, evitando diferencias en el
tratamiento de una situación muy delicada, ya que se está declarando
fallecida a una persona para iniciar su sucesión.

\section{Diligencias probatorias}

Para justificar los extremos legales deben realizarse diligencias
tendientes a averiguar la existencia del ausente, entre ellas:

\begin{itemize}
    \item Oficios a distintas entidades.
    \item Pedido de informes a su tarjeta de crédito.
    \item Pedido de informes a su cuenta bancaria.
    \item Pedido a la empresa de telefonía celular.
    \item Pedido al club si es socio.
    \item Pedido a Migraciones para verificar entradas o salidas del país.
    \item Pedido de informe sobre actividad de su celular.
\end{itemize}

Se le nombra un \textbf{defensor al ausente} (o el defensor oficial). Se
lo cita por \textbf{edictos una vez por mes durante seis meses}: se
publican seis edictos, uno por mes. Se le nombra, asimismo, un curador de
los bienes si no hay quien los esté administrando.

\begin{important}
\textbf{Aclaración del Art. 79 (ausencia simple) frente a la presunción
de fallecimiento.} El Código aclara que el hecho de haber tramitado
previamente la simple ausencia (Art. 79) \textbf{no} es presupuesto
necesario ni suficiente para la presunción de fallecimiento. Es decir:
(a) puede iniciarse directamente la presunción de fallecimiento sin haber
tramitado la ausencia previa; y (b) haber tramitado la ausencia no exime
de tener que probar nuevamente, en el juicio de presunción de
fallecimiento, que se han realizado todas las diligencias para dar con el
ausente. Ambos procesos son independientes, aunque puedan darse en forma
sucesiva.
\end{important}

\section{Sentencia y día presuntivo de fallecimiento}

Pasados los seis meses de citación, recibida la prueba y oído el
defensor, el juez debe dictar sentencia. La sentencia contiene tres
elementos:

\begin{enumerate}
    \item Declara el fallecimiento presunto, comprobados que estén los extremos legales.
    \item Fija el día presuntivo de fallecimiento.
    \item Dispone la inscripción de la sentencia en el Registro Civil.
\end{enumerate}

\textbf{¿Cómo se fija el día presuntivo?}

\begin{table}[H]
\centering
\begin{tabularx}{\textwidth}{>{\raggedright\arraybackslash}X >{\raggedright\arraybackslash}X >{\raggedright\arraybackslash}X}
\toprule
\textbf{Caso} & \textbf{Día presuntivo} & \textbf{Ejemplo} \\
\midrule
Ordinario (Art. 85) & Último día del 1.\textsuperscript{er} año = mitad del plazo de 3 años & Última noticia: 31/12/2017 $\rightarrow$ día presuntivo: 30/06/2019 \\
Extraordinario genérico (Art. 86, 1.\textsuperscript{er} párr.) & El día en que ocurrió el suceso (o el término medio si se extendió en el tiempo) & Si el incendio duró varios días: el día del mediodía del período \\
Extraordinario específico (Art. 86, 2.\textsuperscript{do} párr.) & El último día en que se tuvo noticia del buque o aeronave & MH370: último contacto = 8 de marzo de 2014 $\rightarrow$ ese es el día presuntivo \\
\bottomrule
\end{tabularx}
\caption{Fijación del día presuntivo de fallecimiento según el caso}
\end{table}

El día presuntivo es de gran relevancia porque determina
\textbf{quiénes son los herederos} que entran en la sucesión del
presuntamente fallecido, declarado judicialmente. Si es posible, la
sentencia debe determinar en qué hora presuntiva se fijó el
fallecimiento; si no, se toma al final de ese día.

% ======================================================================
% CAPÍTULO 5
% ======================================================================
\chapter{Garantías y efectos de la reaparición}

\section{Las tres garantías adoptadas}

Con la declaración de fallecimiento presunto se abre la sucesión y los
herederos pueden tomar los bienes. Pero el Código adopta \textbf{tres
garantías} por si el ausente reapareciera:

\begin{table}[H]
\centering
\begin{tabularx}{\textwidth}{>{\raggedright\arraybackslash}p{0.28\textwidth} >{\raggedright\arraybackslash}X}
\toprule
\textbf{Garantía} & \textbf{Descripción} \\
\midrule
\textbf{Inventario} & Se hace un inventario detallado de todos los bienes del presuntamente fallecido (el acervo sucesorio / masa hereditaria). En base a ese inventario, si reaparece, los herederos deben rendir cuentas de todos esos bienes. \\
\textbf{Prenotación registral} & Al transferir los bienes a nombre de los herederos (por ejemplo, un departamento a nombre de dos hijos), se hace una prenotación: se deja constancia en el asiento registral de que el bien está sometido al régimen de declaración de fallecimiento presunto. \\
\textbf{Restricción para enajenar/gravar} & Para enajenar (vender) o gravar (hipotecar) los bienes heredados, los herederos necesitan autorización judicial. La prenotación avisa a terceros contratantes de que ese bien está sometido a este régimen. \\
\bottomrule
\end{tabularx}
\caption{Las tres garantías del sistema ante una eventual reaparición}
\end{table}

\section{Duración de la prenotación (Art. 91 CCCN)}

\begin{formula}{Art. 91 CCCN — Prenotación}
La prenotación caduca a los \textbf{cinco (5) años} desde el día
presuntivo del fallecimiento, o cuando la persona hubiera cumplido
\textbf{ochenta (80) años} de edad, lo que ocurra primero. Vencida la
prenotación, los herederos pueden enajenar libremente los bienes.
\end{formula}

\begin{important}
Son cinco años desde el \textbf{día presuntivo del fallecimiento}, no
desde el día en que la persona desapareció, y no desde la fecha de la
sentencia. Por eso el día presuntivo es tan relevante: es el que pone en
marcha el plazo de prenotación.
\end{important}

\section{Efectos de la reaparición (Art. 91 CCCN)}

Si el ausente reaparece \textbf{antes} de que venza la prenotación: se le
devuelven todos los bienes, porque el inventario asegura que se sabe
exactamente con qué se empezó, y la restricción de enajenación junto con
el requisito de autorización judicial protegieron el patrimonio.

Si el ausente reaparece \textbf{después} de vencida la prenotación
(cuando ya los bienes podían enajenarse libremente):

\begin{itemize}
    \item Se le entregan los bienes que aún existen.
    \item Se le entregan los bienes adquiridos con el valor de los que faltaron (subrogación real: por ejemplo, si se vendió un departamento y se compraron cuatro cocheras, se entregan las cuatro cocheras).
    \item Se le entrega el precio pendiente de cobro si se vendió un bien con saldo impago.
    \item Se le entregan los frutos no consumidos.
\end{itemize}

\begin{example}{Caso real: Oberá, Misiones (2011)}
Una persona se metió al río Uruguay y desapareció; la buscaron por tierra
y aire con helicópteros. Tiempo después, publicó en una red social una
foto desde Curitiba, en la playa. Se armó un escándalo: resultó que había
decidido hacerse ``desaparecer'' voluntariamente y fue descubierto con
una nueva familia. Si esa persona hubiera tenido bienes y se hubiera
iniciado la declaración, podría haber reaparecido y reclamado la
devolución de lo que le correspondía.
\end{example}

\begin{example}{Caso real: Inglaterra — reaparición tras cinco años}
Se conoció un caso en Inglaterra (noticias de diciembre de 2007) de una
persona que había desaparecido y reapareció cinco años después. Si
durante ese tiempo se hubiera declarado su fallecimiento presunto en
Argentina, al reaparecer tendría derecho a que se le restituyeran los
bienes —o su equivalente— conforme las reglas del Art. 91 CCCN.
\end{example}

% ======================================================================
% CUADRO CORNELL
% ======================================================================

\section{Cuadro de repaso Cornell}

El siguiente cuadro reúne, en formato Cornell, las preguntas y conceptos
centrales desarrollados a lo largo de los capítulos anteriores, con el
propósito de facilitar el repaso activo del material.

\begin{longtable}{
    >{\raggedright\arraybackslash}p{0.30\textwidth}
    >{\raggedright\arraybackslash}p{0.58\textwidth}
}
\toprule
\textbf{Preguntas / palabras clave} & \textbf{Notas / respuestas} \\
\midrule
\endfirsthead

\toprule
\textbf{Preguntas / palabras clave} & \textbf{Notas / respuestas} \\
\midrule
\endhead

\midrule
\multicolumn{2}{r}{\textit{Continúa en la página siguiente}} \\
\endfoot

\bottomrule
\endlastfoot

\textbf{¿Cuál es la diferencia entre ausencia y presunción de fallecimiento?} &
La ausencia (Arts. 79-84) es una medida conservatoria de bienes: nombra
un curador sin declarar fallecida a la persona. La presunción de
fallecimiento (Arts. 85-92) declara judicialmente el fallecimiento,
inscribe la sentencia en el Registro Civil y abre la sucesión. La
presunción supone ausencia, pero no es lo mismo. \\

\textbf{¿Quién puede pedir la declaración de ausencia?} &
El Ministerio Público y las personas que tengan interés legítimo en los
bienes (acreedores, sucesores). El juez competente es el del domicilio
del ausente; si no tiene domicilio en el país o es desconocido, el juez
del lugar donde están los bienes. \\

\textbf{¿Cuáles son los tres casos de presunción de fallecimiento?} &
1) Ordinario: ausencia sin noticias por 3 años (Art. 85). 2)
Extraordinario genérico: persona involucrada en catástrofe o actividad de
riesgo, plazo de 2 años (Art. 86, 1.\textsuperscript{er} párr.). 3)
Extraordinario específico: buque naufragado o aeronave perdida, plazo de
6 meses (Art. 86, 2.\textsuperscript{do} párr.). \\

\textbf{¿Cómo se fija el día presuntivo de fallecimiento?} &
En el ordinario: mitad del plazo (1 año y medio después de la última
noticia). En el extraordinario genérico: el día en que ocurrió el suceso
(o término medio si se extendió). En el extraordinario específico: el
último día en que se tuvo noticia del buque o aeronave. \\

\textbf{¿Cuáles son las tres garantías que protegen al ausente por si
reaparece?} &
1) Inventario detallado de todos los bienes. 2) Prenotación registral en
los asientos de los bienes. 3) Prohibición de enajenar o gravar sin
autorización judicial. La prenotación dura 5 años desde el día presuntivo
o hasta que la persona hubiera cumplido 80 años. \\

\textbf{¿Qué pasa si el ausente reaparece después de vencida la
prenotación?} &
Se le restituyen los bienes existentes, los adquiridos con el valor de
los que faltaron (subrogación real), el precio pendiente de cobro y los
frutos no consumidos. No puede recuperar bienes ya enajenados libremente
por los herederos una vez caduca la prenotación. \\

\textbf{¿Es necesario haber tramitado la ausencia antes de iniciar la
presunción de fallecimiento?} &
No. Según el Código, la declaración de ausencia (Art. 79) no es
presupuesto necesario ni suficiente para la presunción de fallecimiento.
Pueden tramitarse en forma sucesiva, pero en el juicio de presunción hay
que probar nuevamente las diligencias para dar con el ausente. \\

\textbf{¿Qué regula el CCCN sobre el procedimiento?} &
El CCCN incorpora normas procesales comunes para todo el país (citación
por edictos una vez por mes durante 6 meses, defensor oficial del
ausente, curador de bienes), para evitar diferencias entre jurisdicciones
en situaciones de gran sensibilidad jurídica y personal. \\

\midrule

\multicolumn{2}{p{\textwidth}}{
\textbf{Resumen:} El derecho civil argentino responde a la incertidumbre
generada por la desaparición de una persona mediante dos institutos
diferenciados que conviven en el CCCN (Arts. 79-92). La \textbf{ausencia}
(Arts. 79-84) es una herramienta de protección patrimonial preventiva:
permite nombrar un curador de bienes y adoptar medidas conservatorias sin
declarar fallecida a la persona. La \textbf{presunción de fallecimiento}
(Arts. 85-92) va más lejos: declara judicialmente el fallecimiento,
inscribe la sentencia en el Registro Civil y habilita la apertura de la
sucesión. Se estructura en tres casos según el grado de probabilidad de
muerte: el ordinario (3 años), el extraordinario genérico —catástrofes o
actividades de riesgo— (2 años) y el extraordinario específico —buque o
aeronave— (6 meses). La sentencia fija el día presuntivo, que determina
quiénes son los herederos y desde cuándo corre la prenotación registral
de cinco años. Durante ese período, los bienes no pueden enajenarse sin
autorización judicial, garantizando la restitución al ausente en caso de
reaparición. El sistema, con sus antecedentes en Vélez Sársfield y la Ley
14.394, busca equilibrar la certeza jurídica que exigen las relaciones
patrimoniales y sucesorias con la protección de los derechos del
ausente, a quien siempre se le reserva la posibilidad de recuperar su
patrimonio si regresa.
} \\

\bottomrule
\end{longtable}
\partintro{Junto a la persona humana, el ordenamiento reconoce como sujetos de derecho a las personas jurídicas. Esta parte estudia su concepto y naturaleza, el régimen de las personas jurídicas privadas, su responsabilidad civil, el fin de su existencia, las fundaciones y las asociaciones civiles y simples asociaciones.}
\part{El sujeto de la relación jurídica: las personas jurídicas}
\chapter{Introducción a las personas jurídicas}
\label{ch:pj}

Todo abogado, trabaje en derecho de familia, comercial, laboral, notarial o administrativo, se encuentra permanentemente con personas jurídicas. Saber identificar su tipo, leer su estatuto, comprender sus órganos de gobierno y conocer las reglas de responsabilidad es una habilidad transversal y cotidiana: desde redactar el estatuto de una ONG, hasta asesorar sobre la responsabilidad de los directivos de una sociedad anónima, gestionar la disolución de un consorcio o defender a una asociación civil en un proceso disciplinario. Este conocimiento también resulta útil en la vida cotidiana, dado que las personas se relacionan constantemente con personas jurídicas: el consorcio del edificio, el club al que se pertenece, la empresa en la que se trabaja, el Estado que regula.

\textbf{Ejemplo integrador que guía la exposición:} un grupo de vecinos forma un club de fútbol. Para ello necesitan constituir una persona jurídica. ¿Qué tipo eligen? ¿Qué nombre pueden usar? ¿Quién representa al club? ¿Pueden sancionar a un socio? ¿Qué pasa si el club se disuelve? Cada capítulo de esta parte responde a una de estas preguntas: el primero trata el concepto, la naturaleza y la clasificación; el segundo, el régimen de las personas jurídicas privadas; el tercero, su responsabilidad civil; el cuarto, el fin de su existencia y las fundaciones; y el quinto, las asociaciones civiles y simples asociaciones.

% ============================================================

\section{Presentación del tema}

El estudio del régimen de las personas jurídicas en el derecho civil y comercial se sitúa dentro de la Parte General del Código, que se ocupa de los elementos de la relación jurídica: el sujeto, el objeto y la causa. Habiendo estudiado a fondo la persona humana, corresponde ahora tratar la persona jurídica, es decir, los entes colectivos que entablan también relaciones jurídicas en el campo del derecho privado.

\section{El problema de fondo: la asociatividad humana}

El problema central que resuelve esta regulación es cómo dar tratamiento jurídico a la asociatividad, dado que el ser humano es un ser sociable por naturaleza que tiende a unirse con otros para fines útiles. Esa regulación jurídica es fundamental porque permite establecer con qué requisitos esos entes colectivos, morales o ideales pueden entablar relaciones jurídicas: qué tipos de personas jurídicas existen, cuándo hay personalidad, cuáles son sus atributos, cuándo empieza y cuándo finaliza su existencia, y qué consecuencias tiene.

Estas personas jurídicas apuntan muchas veces a facilitar la organización de la vida cotidiana: desde el Estado como organización política que trasciende a los ciudadanos y tiene una personalidad distinta a la de estos, hasta las sociedades, las mutuales, las cooperativas y las iglesias.

\begin{note}
Existe un tipo de persona jurídica de enorme trascendencia económica, las sociedades, que tiene su materia específica (el Derecho Societario); el Estado también tiene la suya (el Derecho Administrativo). En esta materia se estudian las grandes disposiciones generales del Código Civil y Comercial aplicables a todas las personas jurídicas.
\end{note}

\section{Regulación en el Código de Vélez (Ley 340, 1871)}

\begin{note}
\textbf{Código de Vélez --- artículos relevantes.}

\textbf{Art. 30:} ``Todos los entes susceptibles de adquirir derechos y contraer obligaciones son personas''.

\textbf{Art. 31:} Las personas se clasifican en personas físicas (o de existencia visible) y personas jurídicas (o de existencia ideal).

\textbf{Art. 32:} Las personas jurídicas son definidas por exclusión: son todos los entes susceptibles de adquirir derechos y contraer obligaciones que no son personas de existencia visible.

\textbf{Arts. 33 a 50:} Régimen general de las personas jurídicas.
\end{note}

\section{El nuevo Código Civil y Comercial: metodología}

El nuevo Código regula la materia de forma más organizada y sistemática. El Libro Primero, la Parte General, comienza con un Título Primero sobre la persona humana, y el Título Segundo trata de las personas jurídicas.

Dentro de los seis libros del Código, en el Libro Primero:
\begin{itemize}
    \item Persona humana: artículos 19 al 140.
    \item Personas jurídicas: a partir del artículo 141.
\end{itemize}

El Título Segundo sobre personas jurídicas se divide en:
\begin{itemize}
    \item Capítulo 1: Parte general (personalidad, composición y clasificación).
    \item Capítulo 2: Asociaciones civiles y simples asociaciones.
    \item Capítulo 3: Fundaciones.
\end{itemize}

\begin{note}
El nuevo Código adopta la denominación única de ``persona jurídica'', diferenciándose del Código anterior, que hablaba de ``persona de existencia ideal'' y ``persona jurídica''. Antes de la Ley 17.711 (1968) se discutía si el género era la persona de existencia ideal y la persona jurídica una especie; con dicha reforma se tendió a equipararlos. Otras denominaciones usadas en doctrina son ``personas morales'' o ``personas colectivas''.
\end{note}

\section{Definición del artículo 141 del CCyCN}

\begin{definition}{Persona jurídica (art. 141 CCyCN)}
Son personas jurídicas todos los entes a los cuales el ordenamiento jurídico les confiere aptitud para adquirir derechos y contraer obligaciones para el cumplimiento de su objeto y los fines de su creación.
\end{definition}

Esta definición presenta dos características destacables:

\begin{enumerate}
    \item \textbf{Se define por la positiva}: se incorpora como elemento al ordenamiento jurídico, que es quien confiere la aptitud. Esto marca la diferencia con las personas humanas, cuya personalidad es un dato de la naturaleza reconocido —no concedido— por el legislador.
    \item \textbf{Se incorpora el elemento finalista}: la aptitud está orientada al cumplimiento del objeto y los fines de la creación. Este es el fundamento del principio de especialidad.
\end{enumerate}

\section{Principio de especialidad}

A diferencia de la persona humana, que puede dedicarse a cualquier actividad lícita, las personas jurídicas tienen un objeto y un fin propios. El \textbf{principio de especialidad} significa que la persona jurídica, dentro de su posibilidad de adquirir derechos y contraer obligaciones, está acotada a sus fines, a su especialidad, a su objeto, y a aquello para lo cual fue constituida.

\begin{important}
\textbf{Artículo 156 CCyCN --- Objeto.} El objeto de la persona jurídica privada debe ser preciso y determinado.
\end{important}

Esto resulta especialmente relevante para la redacción de los estatutos: los constituyentes deben ser claros, explícitos y descriptivos sobre cómo piensan lograr sus fines y cuál es el objeto de la persona jurídica.

\section{Teorías sobre la naturaleza jurídica}

\subsection{Teoría de la ficción (Savigny)}

Savigny recoge elaboraciones que provenían del derecho canónico y desarrolla la idea de la ficción. Para esta teoría, la persona humana es la única dotada de voluntad; como para Savigny el derecho es una expresión de la voluntad, no existirían propiamente personas jurídicas, sino un artificio creado por el legislador, que puede ser ampliado o restringido, pero que carece de voluntad propia.

\begin{important}
Consecuencia fundamental de la teoría de la ficción: la persona jurídica no respondería por los daños causados —responderían solo las personas humanas—, y como se trataría de una representación con mandato, este mandato nunca podría ser conferido para el ilícito. Esto dio lugar a situaciones injustas y, con el tiempo, se modificó esta posición. El artículo 43 del Código de Vélez, en su redacción original, sostenía esta posición: las personas jurídicas no respondían por los daños causados por quienes las dirigían o administraban.
\end{important}

\subsection{Teorías negatorias de la personalidad}

Niegan que existan propiamente personas jurídicas. Sostienen que lo que existe son patrimonios de afectación o bienes colectivos, pero no algo que pueda llamarse persona jurídica.

\subsection{Teorías de la realidad}

Buscan encontrar un sustrato real detrás de lo que se llama persona jurídica:

\begin{itemize}
    \item \textbf{Teoría organicista}: la persona jurídica sería como un organismo vivo, con sus mecanismos de adopción de la voluntad (los órganos que la gobiernan). Existe una realidad que el legislador viene a tutelar, proteger y regular.
    \item \textbf{Teoría de la institución} (Hauriou, seguida en Argentina por Llambías): las personas jurídicas son instituciones. Una institución requiere tres elementos: (a) una idea, objeto o fines en torno a los cuales se nuclean personas; (b) un sustrato personal, es decir, las personas que se nuclean en torno a esa idea; y (c) un poder u organización que permite llevar adelante esos fines.
\end{itemize}

\subsection{Teorías normativistas (Kelsen)}

Las personas jurídicas son un centro de imputación de normas: una realidad a la que el derecho le asigna consecuencias jurídicas a través de normas dictadas por la autoridad competente.

\begin{note}
En el Código Civil y Comercial vigente se percibe una fuerte presencia de las teorías normativistas —el ordenamiento jurídico es quien confiere la aptitud—, pero también existe una dimensión institucional, en tanto se reconocen realidades organizativas como las simples asociaciones.
\end{note}

\section{Clasificación: personas jurídicas públicas y privadas}

El Código Civil y Comercial simplifica la clasificación entre públicas y privadas.

\subsection{Personas jurídicas públicas (arts. 146 y 147)}

\begin{definition}{Personas jurídicas públicas (art. 146 CCyCN)}
Son personas jurídicas públicas:
\begin{enumerate}[label=\alph*)]
    \item El Estado nacional, las provincias, la Ciudad Autónoma de Buenos Aires, los municipios, las entidades autárquicas y las demás organizaciones a las que el ordenamiento jurídico les atribuya ese carácter.
    \item Los estados extranjeros, las organizaciones a las que el derecho internacional público reconozca personalidad jurídica (como la ONU) y toda otra persona jurídica constituida en el extranjero cuyo carácter público resulte de su derecho aplicable.
    \item La Iglesia Católica.
\end{enumerate}
\end{definition}

La inclusión de la Iglesia Católica con personería jurídica pública se vincula con el artículo 2 de la Constitución Nacional, que establece que el gobierno federal sostiene el culto católico apostólico romano, en virtud de la preeminencia histórica de la Iglesia antes de la constitución del Estado argentino.

En cuanto a su reconocimiento, comienzo, capacidad, funcionamiento y organización, las personas jurídicas públicas se rigen por sus propias leyes y ordenamientos. Así, por ejemplo, si se litiga contra el Banco Central, habrá que revisar su Carta Orgánica y las leyes que lo regulan.

\subsection{Personas jurídicas privadas (art. 148)}

El artículo 148 enumera los tipos de personas jurídicas privadas.

\begin{table}[H]
\centering
\caption{Tipos de personas jurídicas privadas (art. 148 CCyCN)}
\begin{tabularx}{\textwidth}{
    >{\raggedright\arraybackslash}p{0.22\textwidth}
    >{\raggedright\arraybackslash}X
    >{\raggedright\arraybackslash}p{0.22\textwidth}
}
\toprule
\textbf{Tipo} & \textbf{Característica principal} & \textbf{Ley aplicable} \\
\midrule
Sociedades & Finalidad directa de lucro. Los socios participan en ganancias y pérdidas. & Ley 19.550 \\
Asociaciones civiles & No tienen fin de lucro. Unión de personas para un fin común (cultural, deportivo, social, etc.). & Arts. 168--186 CCyCN \\
Simples asociaciones & Similar a la asociación civil, pero sin autorización estatal para funcionar. & Arts. 187--192 CCyCN \\
Fundaciones & Patrimonio afectado por uno o varios fundadores a un fin de bien común. Sin miembros. & Arts. 193--224 CCyCN \\
Iglesias y comunidades religiosas (no católicas) & No tienen fin de lucro. Libertad de culto y conciencia. & Ley 21.745 \\
Mutuales & Ayuda mutua entre sus miembros. Socorros recíprocos. & Ley 20.321 \\
Cooperativas & Asociación de personas para fines comunes de producción, trabajo o consumo. & Ley 20.337 \\
Consorcios de propiedad horizontal & Conjunto de propietarios de unidades funcionales de un inmueble sometido al régimen de propiedad horizontal. & Art. 2044 CCyCN \\
\bottomrule
\end{tabularx}
\end{table}

\begin{note}
El artículo 149 establece que, si el Estado participa en una persona jurídica privada, ello no modifica el carácter privado de esta, aunque puede haber derechos y obligaciones diferenciadas en razón del interés público comprometido.
\end{note}

\section{Ley aplicable (art. 150)}

\begin{formula}{Ley aplicable a las personas jurídicas privadas (art. 150 CCyCN)}
Las personas jurídicas privadas se rigen:
\begin{enumerate}
    \item[1.°)] Por las normas imperativas de la ley especial o, en su defecto, de este Código;
    \item[2.°)] Por las normas del acto constitutivo con sus modificaciones y de los reglamentos, prevaleciendo las primeras en caso de divergencia;
    \item[3.°)] Por las normas supletorias de leyes especiales, o en su defecto, por las de este Título.
\end{enumerate}
\end{formula}

El estatuto tiene fuerza de norma jurídica. Por ello, siempre que se aborde una cuestión en la que interviene una persona jurídica, lo primero que corresponde hacer es leer su estatuto en su última versión, con todas las modificaciones, y también sus reglamentos.

\section{Personalidad diferenciada e inoponibilidad (art. 144)}

Un presupuesto fundamental es que la persona jurídica es distinta de sus miembros. Esto implica distintos patrimonios, distintas responsabilidades, y la posibilidad incluso de que el propio miembro demande a la persona jurídica o viceversa.

\begin{important}
\textbf{Regla general}: los miembros no responden por las obligaciones de la persona jurídica. Esta es una de las diferencias más importantes que debe considerar quien constituye una persona jurídica: responde con el capital que aporta, pero no con todo su patrimonio personal.
\end{important}

\begin{definition}{Inoponibilidad de la persona jurídica (art. 144 CCyCN)}
La actuación que esté destinada a la consecución de fines ajenos a la persona jurídica, que constituya un recurso para violar la ley, el orden público o la buena fe, o para frustrar derechos de cualquier persona, se imputa a quienes, a título de socios, asociados, miembros o controlantes directos o indirectos, la hicieron posible, quienes responderán solidaria e ilimitadamente por los perjuicios causados.
\end{definition}

Esta norma reconoce como antecedente el artículo 54 de la Ley de Sociedades (Ley 19.550), pero ahora se aplica a todo el sistema de personas jurídicas por estar ubicada en la Parte General.

La inoponibilidad significa que los efectos de la personería no pueden hacerse valer frente a quien ha sido dañado, quien puede accionar directamente contra los socios, asociados, miembros o controlantes que hicieron posible esa actuación dañina. Responden solidaria e ilimitadamente: todos responden por el todo y con todos sus patrimonios.

% ============================================================
\chapter{Régimen de las personas jurídicas privadas}
% ============================================================

El Código ofrece reglas generales sobre nombre, domicilio, capacidad, patrimonio, estatutos y comienzo de la existencia, que luego cada tipo de persona jurídica complementa con sus normas específicas.

\section{Nombre (art. 151)}

\begin{definition}{Nombre de la persona jurídica (art. 151 CCyCN)}
La persona jurídica debe tener un nombre que la identifique como tal, con el aditamento indicativo de la forma jurídica adoptada. La persona jurídica en liquidación debe aclarar esta circunstancia en la utilización de su nombre. El nombre debe satisfacer recaudos de veracidad, novedad y aptitud distintiva, tanto respecto de otros nombres como de marcas, nombres de fantasía u otras formas de referencia a bienes o servicios, se relacionen o no con el objeto de la persona jurídica.

No puede contener términos o expresiones contrarios a la ley, el orden público o las buenas costumbres, ni inducir a error sobre la clase u objeto de la persona jurídica.

Si se incluye el nombre de personas humanas, deben contar con la conformidad de estas, que se presume si son socios. Sus herederos pueden oponerse a la continuación del uso si acreditan perjuicios materiales o morales.
\end{definition}

Los requisitos pueden precisarse del siguiente modo:

\begin{itemize}
    \item \textbf{Indicación de la forma jurídica}: debe agregarse el tipo de persona jurídica (por ejemplo, ``El estudiante feliz S.A.''). Si está en liquidación, se agrega ``en liquidación''.
    \item \textbf{Veracidad}: concordancia entre el nombre y el objeto, para que no haya lugar a confusión. Esto se juega especialmente en ciertos ámbitos, como el de las universidades.
    \item \textbf{Novedad y aptitud distintiva}: el nombre debe distinguirse de otros nombres, marcas, nombres de fantasía u otras formas de referencia a bienes y servicios.
\end{itemize}

\begin{example}{Teatro Colón S.A.}
No se puede fundar una sociedad anónima con el nombre ``Teatro Colón S.A.'' porque se confundiría con el Teatro Colón preexistente. El nombre no tendría novedad ni aptitud distintiva.
\end{example}

\begin{example}{Fundación Lionel Messi / Fundación Borges}
Si se quiere usar el nombre de una persona humana, se necesita su conformidad. Si esa persona ya falleció, los herederos pueden oponerse si acreditan perjuicios materiales o morales. Así, si la Fundación Borges —supuestamente dedicada a las altas letras— organizara un concurso de letras de reggaetón, los herederos de Borges podrían oponerse alegando que ello afecta el buen nombre del escritor y resulta incompatible con los fines fundacionales.
\end{example}

\section{Domicilio y sede social (arts. 152 y 153)}

El domicilio es el asiento jurídico de la persona jurídica y permite ubicarla, determinar el juez competente y conocer su inscripción registral.

\begin{table}[H]
\centering
\caption{Domicilio y sede social}
\begin{tabularx}{\textwidth}{>{\raggedright\arraybackslash}X >{\raggedright\arraybackslash}X}
\toprule
\textbf{Domicilio} & \textbf{Sede social} \\
\midrule
Se fija en el estatuto. Puede expresarse simplemente con una ciudad (por ejemplo, ``Buenos Aires''). & Es la dirección física concreta donde tienen lugar la dirección y administración de la persona jurídica (por ejemplo, ``Av. Figueroa Alcorta 2263, CABA''). \\
Determina el juez competente y la jurisdicción registral. & Generalmente se inscribe ante el organismo de contralor (IGJ). Si está en el estatuto, cambiarla requiere modificar el estatuto. \\
\bottomrule
\end{tabularx}
\end{table}

\begin{important}
\textbf{Artículo 153 CCyCN --- Efectos de la sede social.} Se tienen por válidas y vinculantes para la persona jurídica todas las notificaciones efectuadas en la sede inscripta.
\end{important}

También existe el domicilio especial: cuando la persona jurídica posee establecimientos o sucursales, tiene su domicilio especial en el lugar de dichos establecimientos, pero solo para las obligaciones allí contraídas.

\begin{example}{Sucursal del Banco Santander en Mendoza}
Si un contratista de pintura firmó un contrato con la sucursal mendocina del Banco Santander y este no le paga, puede litigar en Mendoza. No necesita ir a Buenos Aires a litigar contra la casa matriz, porque la obligación se contrajo localmente.
\end{example}

\section{Capacidad y principio de especialidad (arts. 141, 156)}

La persona jurídica es capaz de derecho; la capacidad de ejercicio se realiza a través de sus órganos representativos. Pero esa capacidad tiene límites:

\begin{enumerate}
    \item \textbf{Principio de especialidad}: la capacidad está acotada al objeto y a los fines de la persona jurídica.
    \item \textbf{Limitaciones legales}: la ley prohíbe a las personas jurídicas adquirir ciertos derechos o los limita. Por ejemplo, el usufructo en favor de personas jurídicas tiene un plazo determinado, mientras que para personas humanas puede extenderse hasta su muerte.
    \item \textbf{Limitaciones por la naturaleza de las cosas}: las personas jurídicas no pueden adoptar niños, porque el derecho de familia se refiere a las personas humanas.
\end{enumerate}

\section{Patrimonio (arts. 143 y 154)}

La persona jurídica debe tener un patrimonio distinto del de sus miembros. Esto es una consecuencia directa de la personalidad diferenciada.

La persona jurídica tiene bienes a su nombre; si se trata de información registrable, se pueden inscribir preventivamente. Por ejemplo, si alguien dona un inmueble a una fundación y esta tarda en inscribirse, puede realizarse una anotación preventiva en el registro para proteger el bien mientras se completan los trámites.

\begin{note}
En cuanto a su duración, las personas jurídicas son en principio ilimitadas en el tiempo, excepto que el acto constitutivo disponga lo contrario. La prórroga está regulada en el artículo 165 y requiere decisión de los miembros o aprobación de la autoridad de contralor antes del vencimiento del plazo.
\end{note}

\section{Comienzo de la existencia (art. 142)}

\begin{definition}{Comienzo de la existencia (art. 142 CCyCN)}
La existencia de la persona jurídica privada comienza desde su constitución. No necesita autorización legal para funcionar, excepto disposición legal en contrario. En los casos en que se requiere autorización estatal, la persona jurídica no puede funcionar antes de obtenerla.
\end{definition}

En el Código anterior existía una discusión sobre si el comienzo de la existencia era el momento de la constitución o el de la autorización. El nuevo Código es claro: el comienzo es la constitución. Sin embargo, si la ley requiere autorización estatal para funcionar (como en el caso de las fundaciones), la persona jurídica no puede comenzar a operar hasta obtenerla.

\begin{example}{Fundación banco de alimentos}
Si hoy se constituye una fundación con objeto de banco de alimentos y se aporta un capital de \$100.000, la fundación existe desde ese momento. Pero no puede empezar a hacer compras, emitir facturas ni hacer donaciones hasta que se obtenga la autorización estatal. La constitución y la autorización son dos momentos distintos.
\end{example}

\section{Estatutos y gobierno (arts. 150, 158, 159, 160)}

El estatuto es una norma jurídica secundaria que emana de las partes que constituyen la persona jurídica y que la regula internamente. No debe confundirse con el acto constitutivo:

\begin{itemize}
    \item El \textbf{acto constitutivo} es el acto jurídico por el cual las partes deciden dar inicio a la persona jurídica. Generalmente, en ese acto se redacta el estatuto.
    \item El \textbf{estatuto} cobra independencia porque puede sufrir sucesivas modificaciones a lo largo del tiempo.
\end{itemize}

\begin{note}
Hay clubes que tienen más de cien años y cuyo estatuto ha cambiado muchísimas veces. Cuando se solicita el estatuto ante la Inspección de Justicia, se pide el actualizado con todos sus cambios.
\end{note}

\subsection{Contenido del estatuto}

El estatuto de una persona jurídica privada debe contener:

\begin{itemize}
    \item Nombre y domicilio.
    \item Objeto preciso y determinado (art. 156).
    \item Derechos y deberes de los miembros.
    \item Forma de gobierno: composición de los órganos, cómo se toman las decisiones, quórum y mayorías necesarias (incluidas las agravadas para decisiones trascendentes, como reforma del estatuto, expulsión de miembros o disolución).
    \item Atribuciones de cada órgano y sus miembros.
    \item Administración de bienes.
    \item Destino de los bienes en caso de disolución.
\end{itemize}

\subsection{Deber de lealtad y diligencia de los administradores (art. 159)}

\begin{important}
\textbf{Artículo 159 CCyCN --- Deber de lealtad y diligencia.} Los administradores de la persona jurídica deben obrar con lealtad y diligencia. No pueden perseguir ni favorecer intereses contrarios a los de la persona jurídica. Si en determinada operación tuvieran un interés contrario al de ella, deben hacerlo saber a los demás miembros del órgano de administración o, en su caso, al órgano de gobierno, y abstenerse de cualquier intervención relacionada con dicha operación.
\end{important}

Ser administrador de una persona jurídica implica una gran responsabilidad, porque se confían bienes económicos que no son propios.

% ============================================================
\chapter{Responsabilidad civil de las personas jurídicas}
% ============================================================

\section{Evolución histórica: del artículo 43 de Vélez al nuevo Código}

El Código de Vélez Sarsfield, adherente a la teoría de la ficción, establecía en el artículo 43 que las personas jurídicas no respondían por los daños que causaban quienes las dirigían o administraban: no se podía iniciar una acción de daños y perjuicios contra una persona jurídica.

La lógica era coherente con la teoría de la ficción: como la personería era solo ficción y el mandato nunca se confería para el ilícito, las actuaciones ilícitas se imputaban a las personas humanas, no a la persona jurídica. A lo largo del tiempo, esta solución fue generando la sensación de que constituía una injusticia.

% TODO: contenido incompleto o ambiguo en las notas originales — la referencia al profesor "Rivarola" y a la traducción del término utilizado por Freitas se transcribe tal como figura en el apunte, sin poder verificarse contra la fuente original.
Existieron distintas posturas doctrinarias para resolver el problema. Un profesor de apellido Rivarola acudió al diccionario que empleaba Vélez y encontró que, en el texto original de Freitas (en portugués), el término utilizado podía entenderse como ``cuando'' y no como ``aunque'', lo que cambiaría el sentido de la norma; sin embargo, esta interpretación no fue pacífica. El jurista Borda, en un fallo, llegó a declarar que el artículo 43 era de una injusticia intolerable.

\begin{important}
Todo cambió en 1968 con la Ley 17.711, que reformó el artículo 43 y estableció que la persona jurídica responde por los daños que causen quienes la dirigen o administran en ejercicio o con ocasión de sus funciones. El nuevo Código recoge esta solución en el artículo 1763.
\end{important}

\section{Tres supuestos de responsabilidad}

\subsection{Responsabilidad por actos de directivos y administradores (art. 1763)}

\begin{formula}{Responsabilidad de la persona jurídica (art. 1763 CCyCN)}
La persona jurídica responde por los daños que causen quienes las dirigen o administran en ejercicio o con ocasión de sus funciones.
\end{formula}

La expresión ``en ejercicio o con ocasión de sus funciones'' resulta determinante: si un presidente de la empresa está de vacaciones y comete un daño, la empresa no responde, porque el acto no ocurre en ejercicio ni con ocasión de sus funciones.

\subsection{Responsabilidad por cosas y actividades riesgosas (art. 1757)}

\begin{formula}{Hecho de las cosas y actividades riesgosas (art. 1757 CCyCN)}
Toda persona responde por el daño causado por el riesgo o vicio de las cosas, o de las actividades que sean riesgosas o peligrosas por su naturaleza, por los medios empleados o por las circunstancias de su realización.
\end{formula}

Esta norma se aplica a todas las personas, incluidas las jurídicas. Se trata de una responsabilidad objetiva: no depende de factores de atribución subjetivos (dolo o culpa); la responsabilidad se atribuye siempre.

\begin{example}{Sociedad anónima y accidente de tránsito}
Una sociedad anónima tiene un automóvil. Ese auto causa un daño. La persona jurídica responde, aunque no estuviera manejando: se le imputa la responsabilidad porque es dueña de la cosa riesgosa.
\end{example}

\subsection{Responsabilidad por los dependientes (art. 1753)}

\begin{formula}{Responsabilidad del principal por el hecho del dependiente (art. 1753 CCyCN)}
El principal responde objetivamente por los daños que causen los que están bajo su dependencia, o las personas de las cuales se sirve para el cumplimiento de sus obligaciones, cuando el hecho dañoso acaece en ejercicio o con ocasión de las funciones encomendadas.
\end{formula}

También esta norma engloba a las personas jurídicas. El término ``dependientes'' es genérico: incluye a cualquier persona de la que la persona jurídica se sirve para cumplir sus obligaciones, más allá de la relación de empleo formal.

\begin{example}{Personal de seguridad que golpea a alguien}
En una empresa, el personal de seguridad golpea a una persona de manera arbitraria y violenta. El daño causado se imputa tanto al empleado como a la persona jurídica. La empresa no puede escudarse en que ``fue el empleado''.
\end{example}

\section{Responsabilidad de los administradores (art. 160)}

\begin{formula}{Responsabilidad de los administradores (art. 160 CCyCN)}
Los administradores responden en forma ilimitada y solidaria frente a la persona jurídica, sus integrantes y terceros, por los daños causados por su culpa en el ejercicio o con ocasión de sus funciones, por acción u omisión.
\end{formula}

Esta disposición es la contracara del artículo 1763: si el presidente causó un daño, frente al tercero dañado responden tanto el presidente como la persona jurídica. Pero la persona jurídica, a su vez, puede iniciar acción contra el presidente por los daños que causó. Los administradores responden de forma ilimitada (con todos sus bienes) y solidaria (todos responden por el todo).

% ============================================================
\chapter{Fin de la existencia y fundaciones}
% ============================================================

\section{Causales de disolución (art. 163)}

\begin{formula}{Causales de disolución (art. 163 CCyCN)}
La persona jurídica se disuelve por:
\begin{enumerate}[label=\alph*)]
    \item La decisión de sus miembros adoptada por unanimidad o por la mayoría establecida por el estatuto o disposición especial;
    \item El cumplimiento de la condición resolutoria a la que el acto constitutivo subordinó su existencia;
    \item La consecución del objeto para el cual la persona jurídica se formó, o la imposibilidad sobreviniente de cumplirlo;
    \item El vencimiento del plazo;
    \item La declaración de quiebra;
    \item La fusión respecto de las personas jurídicas que se fusionan o la escisión respecto de la escindida;
    \item La reducción a uno del número de miembros, si la ley especial exige pluralidad de ellos y esta no es restablecida dentro de los tres meses;
    \item La denegatoria o revocación firmes de la autorización estatal para funcionar, cuando esta sea requerida;
    \item El agotamiento de los bienes destinados a sostenerla;
    \item Cualquier otra causa prevista en el estatuto o en disposiciones de este Título o de ley especial.
\end{enumerate}
\end{formula}

Entre las causales más relevantes cabe señalar:

\begin{itemize}
    \item \textbf{Inciso a) --- Decisión de los miembros}: por unanimidad si el estatuto no dice nada; si el estatuto establece una mayoría, se aplica esa mayoría.
    \item \textbf{Inciso b) --- Condición resolutoria}: cuando la existencia estaba subordinada a un hecho futuro e incierto que se cumple.
    \item \textbf{Inciso c) --- Consecución del objeto}: cuando la persona jurídica se creó para un fin específico y este se cumplió o se tornó imposible.
    \item \textbf{Inciso e) --- Quiebra}: cuando el pasivo supera al activo y no hay propuesta de acuerdo preventivo viable. La quiebra está regulada por ley especial y tiene por objeto liquidar el activo para pagar a los acreedores; excepcionalmente, puede haber avenimiento o conversión.
    \item \textbf{Inciso h) --- Revocación estatal}: el Estado puede revocar la autorización para funcionar, con las restricciones del artículo 164.
\end{itemize}

\begin{example}{Fundación condicionada a un programa estatal}
Una fundación se crea para ayudar a personas en situación de calle ``siempre y cuando no exista un programa especial del Gobierno de la Ciudad de Buenos Aires para esa problemática''. Si el gobierno implementa tal programa, se cumple la condición y la fundación se disuelve.
\end{example}

\begin{example}{Fundación para los cien años de un club}
Se constituye una fundación especial para organizar la celebración del centenario de un club. Pasada la celebración, el objeto se cumplió y la personería se disuelve.
\end{example}

\section{Revocación estatal --- garantías (art. 164)}

\begin{important}
\textbf{Artículo 164 CCyCN --- Revocación de la autorización estatal.} La revocación de la autorización estatal debe fundarse en la comisión de actos graves que violen la ley, el estatuto o el reglamento. La resolución que la disponga debe ser fundada y contener la indicación de los actos que la justifican. Se debe dar a los interesados la posibilidad de ser escuchados antes de que la resolución se dicte. Es apelable judicialmente. La apelación puede tener efectos suspensivos.
\end{important}

Las garantías que protegen a la persona jurídica frente a una revocación arbitraria pueden sintetizarse así:

\begin{enumerate}
    \item Solo procede por actos graves que violen la ley, el estatuto o el reglamento.
    \item La resolución debe ser fundada (no puede ser arbitraria).
    \item Se debe garantizar el derecho de defensa (descargo previo).
    \item Es apelable judicialmente, y la apelación puede tener efecto suspensivo (la persona jurídica sigue funcionando mientras se tramita).
\end{enumerate}

El fundamento constitucional de estas garantías es la libertad de asociación consagrada en el artículo 14 de la Constitución Nacional.

\section{Reconducción y liquidación}

\begin{itemize}
    \item \textbf{Reconducción}: si se decidió la disolución pero no terminó la liquidación, se puede revertir esa decisión siempre que la causa pueda ser removida. Por ejemplo, si se cumplió el plazo pero los miembros deciden ampliarlo y transformar la personería de tiempo limitado a indefinido.
    \item \textbf{Liquidación}: una vez dada la causal de disolución, se pagan las deudas y obligaciones pendientes con el activo, se abonan los gastos, y los bienes remanentes se destinan según lo establezca la ley o el estatuto. En asociaciones civiles y fundaciones generalmente se exige que los bienes vayan a otra institución de bien común, sin fin de lucro.
\end{itemize}

\section{Las fundaciones (arts. 193--224)}

\begin{definition}{Fundación (art. 193 CCyCN)}
Las fundaciones son personas jurídicas que se constituyen con una finalidad de bien común, sin propósito de lucro, mediante el aporte patrimonial de una o más personas, denominadas fundadores, destinando dicho aporte al cumplimiento de sus fines.
\end{definition}

Características distintivas de las fundaciones:

\begin{itemize}
    \item No tienen miembros: el fundador aporta un patrimonio, que queda afectado al cumplimiento de un fin de bien común.
    \item No tienen fin de lucro.
    \item Requieren autorización estatal para funcionar (art. 142).
    \item Pueden constituirse por acto entre vivos o por testamento.
\end{itemize}

\begin{example}{Fundación banco de alimentos (caso detallado)}
En un contexto de crisis, un fundador con capital decide constituir una fundación con \$100.000 para organizar un banco de alimentos. La fundación alquila un depósito, contrata empleados, establece convenios con instituciones que proveen alimentos y organiza campañas de donación. Todo eso es la fundación funcionando. No hay ``socios'': hay un fundador que dotó el patrimonio y un Consejo de Administración que lo gestiona.

También puede constituirse por testamento: un deportista que practique hockey puede establecer en su testamento que su parte disponible de la herencia vaya a una fundación para que personas sin recursos puedan practicar ese deporte.
\end{example}

\subsection{Patrimonio inicial suficiente}

El Código exige que la fundación cuente con un patrimonio inicial que posibilite razonablemente el cumplimiento de sus fines (recaudo no exigido para las asociaciones civiles). La Inspección de Personas Jurídicas no otorgará la autorización si los fondos son insuficientes.

El plan de acción a tres años (plan trienal) que se presenta ante el organismo de contralor puede incluir tanto los aportes iniciales como los ingresos futuros proyectados (eventos benéficos, campañas de donación, etc.).

\subsection{Órgano de gobierno: Consejo de Administración}

El único órgano de gobierno es el Consejo de Administración, dotado de todas las atribuciones para el gobierno y la administración. Debe tener al menos tres personas humanas como integrantes.

Puede delegar funciones en un Comité Ejecutivo (art. 205), que ejerce esas funciones entre las reuniones del Consejo. Los miembros del Comité Ejecutivo pueden ser remunerados, mientras que los del Consejo de Administración no.

\subsection{Control externo}

Las fundaciones no tienen órgano de fiscalización interna: el control es exclusivamente externo, a cargo de la autoridad de contralor (IGJ u organismo provincial equivalente). Entre las atribuciones del organismo de contralor (art. 222) se encuentran:

\begin{itemize}
    \item Aprobar estatutos y sus reformas.
    \item Fiscalizar el funcionamiento.
    \item Designar administradores interinos, suspender deliberaciones, remover administradores, convocar asambleas.
    \item Intervenir en la disolución y liquidación.
\end{itemize}

\subsection{Inmutabilidad del objeto (art. 216)}

No es posible reformar el objeto de la fundación, salvo que haya resultado de cumplimiento imposible. Esto se debe a que la fundación fue autorizada para ese objeto específico, y cambiarlo implicaría traicionar la voluntad inicial del fundador.

% ============================================================
\chapter{Asociaciones civiles y simples asociaciones}
% ============================================================

Las asociaciones civiles y las simples asociaciones son personas jurídicas privadas de enorme importancia social: clubes deportivos, ONG, asociaciones barriales, organizaciones de víctimas, entidades religiosas. La Constitución Nacional consagra la libertad de asociación como un derecho fundamental.

\section{Asociaciones civiles (arts. 168--186)}

\subsection{Concepto y objeto}

\begin{definition}{Objeto de las asociaciones civiles (art. 168 CCyCN)}
La asociación civil debe tener un objeto que no sea contrario al interés general o al bien común. El objeto no puede ser contrario a las disposiciones de este Código ni a la ley especial, ni referirse a cuestiones de interés particular de los asociados.
\end{definition}

La definición es por la negativa: cualquier finalidad que no afecte al interés general puede ser objeto de una asociación civil. La asociación civil es una corporación: unión de personas en torno a un fin común (cultural, artístico, deportivo, religioso, social).

\begin{important}
La asociación civil no tiene fin de lucro y no persigue el lucro como fin principal. Puede eventualmente generar ganancias si organiza un evento a beneficio, pero ese dinero se reinvierte en los fines de la asociación; no se distribuye entre los miembros.
\end{important}

\subsection{Diferencia fundamental con las sociedades}

En la sociedad, los socios aportan capital, participan de ganancias y pérdidas, y obtienen utilidades. En la asociación civil, los miembros aportan su cuota —fijada por el estatuto y las asambleas— que contribuye a los fines colectivos, sin distribuirse ganancias.

\begin{note}
La Resolución de la IGJ 7/2015 entiende por ``bien común'' a las finalidades que contribuyen al bien de la comunidad en general o a la mejora de las condiciones de vida, en contraposición al bien individual.
\end{note}

\subsection{Constitución y contenido del estatuto (art. 170)}

El instrumento constitutivo debe hacerse en escritura pública. El artículo 170 establece que debe contener:

\begin{itemize}
    \item Identificación de los constituyentes.
    \item Nombre, objeto y domicilio.
    \item Plazo de duración y causales de disolución.
    \item Composición del patrimonio inicial.
    \item Clases o categorías de socios (activos, cadetes, vitalicios, etc.).
    \item Régimen de admisión, renuncia y sanciones disciplinarias.
    \item Órganos de gobierno: composición, atribuciones, funcionamiento y duración de cargos.
    \item Cierre del ejercicio económico anual.
    \item Destino de los bienes en caso de liquidación: deben ir a otra entidad de bien común, pública o privada, sin fin de lucro, con domicilio en la República Argentina.
\end{itemize}

\section{Órganos de gobierno}

\subsection{Comisión Directiva (art. 171)}

Es el órgano de dirección y administración cotidiana de la asociación civil. Los cargos mínimos son Presidente, Secretario y Tesorero (los demás son Vocales). Cada estatuto puede agregar cargos adicionales: vicepresidentes primero y segundo, prosecretario, protesorero, etc.

Debe reunirse al menos una vez al mes. Sus atribuciones las fija el estatuto, y hay que distinguir entre las que tiene la Comisión como cuerpo colegiado y las que tiene el Presidente individualmente.

\begin{important}
No hay que confundir el hecho de que alguien sea presidente con que pueda resolver todo por sí mismo. El presidente es quien preside un órgano que es la Comisión Directiva.
\end{important}

\begin{example}{Venta de un inmueble por una asociación civil}
Si una asociación civil quiere vender un inmueble, hay que leer el estatuto para saber quién puede tomar esa decisión. A veces se exige que la asamblea autorice previamente la venta; otras veces basta con la Comisión Directiva. El abogado debe siempre revisar el estatuto antes de asesorar.
\end{example}

\subsection{Asamblea}

Es el órgano supremo: la reunión de todos los asociados. En asociaciones muy numerosas puede instrumentarse a través de representantes elegidos por los socios.

\begin{itemize}
    \item \textbf{Asamblea ordinaria}: generalmente una vez al año; aprueba la memoria y el balance presentado por la Comisión Directiva con el dictamen del órgano de fiscalización.
    \item \textbf{Asamblea extraordinaria}: se convoca cuando es necesario modificar el estatuto, elegir autoridades, aprobar reglamentos o tomar decisiones que requieran mayorías especiales.
\end{itemize}

\subsection{Órgano de fiscalización}

Puede recaer en personas no asociadas. Si la asociación cuenta con más de cien socios, debe tener una Comisión Revisora de Cuentas. Si la asociación exige una profesión determinada para ser socio, los integrantes del órgano de fiscalización no necesariamente deben contar con ese título, pero debe contratarse a profesionales independientes para el asesoramiento correspondiente.

Además del control interno, existe también la fiscalización estatal del artículo 174.

\section{Poder disciplinario}

La asociación civil puede sancionar a sus miembros. En esto se evidencia la distinción entre la persona jurídica y sus miembros.

\subsection{Principios que rigen el poder disciplinario}

\begin{enumerate}
    \item Debe estar regulado en el estatuto, con claridad sobre qué órgano interviene (asamblea, Comisión Directiva o Comisión de Disciplina).
    \item Se debe garantizar el derecho de defensa del afectado: debe darse vista al sumariado y permitirle hacer sus descargos.
    \item Gradualidad y proporcionalidad de la sanción: primero se amonesta, luego se suspende por períodos crecientes, y finalmente se expulsa. No se puede expulsar sin proceso.
    \item La decisión es apelable ante los tribunales judiciales si no se respetó el derecho de defensa, si fue desproporcionada o si no hubo debido proceso.
\end{enumerate}

\begin{example}{Destrucción de un vestuario}
Un socio muy enojado, luego de un partido, destruye el vestuario del club un domingo. El lunes, la Comisión Directiva puede suspenderlo transitoriamente como medida cautelar. Pero no puede expulsarlo sin proceso: debe abrirse un sumario, recopilarse pruebas y testimonios, notificarse al imputado, permitirle hacer su descargo y luego resolver. La sanción debe ser proporcional a la falta.
\end{example}

\subsection{Límites de la revisión judicial}

El juez puede revisar la forma del proceso disciplinario (si se respetó el derecho de defensa, si la sanción fue proporcional). Pero no puede revisar el fondo de la decisión —salvo grave violación de derechos humanos—, porque ello implicaría alterar el objeto y la autonomía de la persona jurídica.

\begin{example}{Club vegano y el vendedor de panchos}
Un club de personas veganas expulsa a un socio que insiste en vender panchos en las instalaciones. Si el juez le diera la razón al expulsado en cuanto al fondo, estaría alterando el objeto de la asociación y el derecho de todos sus miembros a mantener la identidad que los une. Por eso, el control judicial se centra en la forma, no en el fondo de la decisión disciplinaria.
\end{example}

\section{Simples asociaciones (arts. 187--192)}

Las simples asociaciones son personas jurídicas privadas que, a diferencia de las asociaciones civiles, no necesitan autorización estatal para funcionar. Existen muchísimas en la práctica.

\subsection{Diferencias con las asociaciones civiles}

\begin{table}[H]
\centering
\caption{Asociación civil frente a simple asociación}
\begin{tabularx}{\textwidth}{
    >{\raggedright\arraybackslash}p{0.24\textwidth}
    >{\raggedright\arraybackslash}X
    >{\raggedright\arraybackslash}X
}
\toprule
\textbf{Aspecto} & \textbf{Asociación civil} & \textbf{Simple asociación} \\
\midrule
Instrumento constitutivo & Escritura pública & Instrumento público o privado con firma certificada \\
Autorización estatal & Requerida & No requerida \\
Comienzo de existencia & Desde la constitución (previa autorización para funcionar) & Desde la constitución \\
Órgano de fiscalización & Obligatorio (Comisión Revisora de Cuentas si hay más de 100 socios) & Prescindible si hay menos de 20 socios; se certifica con contador externo \\
Responsabilidad en insolvencia & Separación plena entre persona jurídica y miembros & Administradores y miembros que participaron responden solidariamente con sus bienes propios \\
Régimen supletorio & No aplica & Se aplica el régimen de asociaciones civiles \\
\bottomrule
\end{tabularx}
\end{table}

\subsection{Responsabilidad en las simples asociaciones (art. 191)}

\begin{formula}{Responsabilidad en simples asociaciones (art. 191 CCyCN)}
Si la simple asociación no puede pagar a sus acreedores, los administradores y todos los miembros que hubieren participado de hecho en la administración responden solidariamente con sus propios bienes por las deudas de la simple asociación.

Los demás miembros que no participaron en la administración solo responden hasta el límite de las cuotas y contribuciones prometidas o debidas.
\end{formula}

En la simple asociación existe una menor diferencia entre la persona jurídica y sus miembros: quienes participaron de alguna forma en la administración y llevaron a un resultado de quebranto o insolvencia deben responder con sus propios bienes.

% ============================================================

\section{Cuadro Cornell de repaso general}
% ============================================================

El siguiente cuadro reúne, según el método Cornell, las preguntas centrales que atraviesan los cuatro bloques temáticos estudiados, con sus respuestas y una síntesis final para el repaso integral de la materia.

\begin{longtable}{
    >{\raggedright\arraybackslash}p{0.28\textwidth}
    >{\raggedright\arraybackslash}p{0.62\textwidth}
}
\cornellhead
\endfirsthead
\cornellhead
\endhead
\midrule
\multicolumn{2}{r}{\textit{continúa en la página siguiente}} \\
\midrule
\endfoot
\bottomrule
\endlastfoot

\textbf{¿Qué es una persona jurídica?} &
Son los entes a los cuales el ordenamiento jurídico les confiere aptitud para adquirir derechos y contraer obligaciones para el cumplimiento de su objeto y los fines de su creación (art. 141 CCyCN). A diferencia de la persona humana, cuya personalidad es reconocida por el legislador, la de las personas jurídicas es conferida por él. \\[0.4em]

\textbf{¿En qué consiste el principio de especialidad?} &
La capacidad de la persona jurídica está acotada a su objeto y sus fines específicos. No puede hacer cualquier cosa: solo aquello para lo que fue creada. De ahí la importancia de redactar el objeto con precisión y amplitud suficiente en el estatuto. \\[0.4em]

\textbf{¿Cuál es la diferencia entre la teoría de la ficción y las teorías de la realidad?} &
La teoría de la ficción (Savigny) entiende que la persona jurídica es un artificio del legislador, sin voluntad real. Las teorías de la realidad (organicista, de la institución) buscan un sustrato real. Las teorías normativistas las conciben como centros de imputación de normas. El CCyCN combina elementos normativistas e institucionales. \\[0.4em]

\textbf{¿Cómo clasifica el Código a las personas jurídicas?} &
Públicas: Estado nacional, provincias, municipios, entidades autárquicas, estados extranjeros, organismos internacionales, Iglesia Católica. Privadas: sociedades, asociaciones civiles, simples asociaciones, fundaciones, mutuales, cooperativas, consorcios de propiedad horizontal y otras previstas en leyes especiales. \\[0.4em]

\textbf{¿Qué norma prevalece según el artículo 150?} &
Orden de prevalencia: (1) normas imperativas de la ley especial; (2) normas imperativas del CCyCN; (3) estatuto y reglamentos; (4) normas supletorias. El estatuto tiene fuerza de norma jurídica en el ámbito de la persona jurídica. \\[0.4em]

\textbf{¿Cuándo se aplica la inoponibilidad del artículo 144?} &
Cuando la persona jurídica es usada para violar la ley, el orden público, la buena fe o para frustrar derechos de terceros, la personería se vuelve inoponible y se puede accionar directamente contra los socios, miembros o controlantes que hicieron posible esa actuación. Responden solidaria e ilimitadamente. \\[0.4em]

\textbf{¿Qué requisitos debe cumplir el nombre de una persona jurídica?} &
Debe indicar la forma jurídica adoptada; debe ser veraz, novedoso y con aptitud distintiva (diferenciarse de nombres, marcas y fantasías); no puede inducir a error sobre el objeto; no puede contener términos contrarios a la ley, el orden público o las buenas costumbres. Si incluye el nombre de una persona humana, se necesita su conformidad. \\[0.4em]

\textbf{¿Cuál es la diferencia entre domicilio y sede social?} &
Domicilio: ciudad indicada en el estatuto; determina jurisdicción y juez competente. Sede social: dirección física inscripta ante el organismo de contralor; las notificaciones allí enviadas se tienen por válidas. Domicilio especial: en el lugar de las sucursales, para las obligaciones allí contraídas. \\[0.4em]

\textbf{¿Cómo responde civilmente la persona jurídica?} &
Tres supuestos: (1) daños causados por directivos/administradores en ejercicio o con ocasión de sus funciones (art. 1763); (2) daños por cosas o actividades riesgosas, responsabilidad objetiva (art. 1757); (3) daños causados por dependientes en ejercicio o con ocasión de funciones (art. 1753). Los administradores también responden ilimitada y solidariamente ante la propia persona jurídica (art. 160). \\[0.4em]

\textbf{¿Cuáles son las causales de disolución del artículo 163?} &
Decisión de los miembros; condición resolutoria; consecución o imposibilidad del objeto; vencimiento del plazo; quiebra; fusión o escisión; reducción a uno del número de miembros; revocación o denegatoria de autorización estatal; agotamiento de bienes; otras causas estatutarias o legales. \\[0.4em]

\textbf{¿Qué garantías protegen frente a la revocación estatal del artículo 164?} &
Solo procede por actos graves que violen la ley, el estatuto o el reglamento. Requiere resolución fundada, garantía del derecho de defensa, y es apelable judicialmente con posible efecto suspensivo. El fundamento es la libertad de asociación constitucional. \\[0.4em]

\textbf{¿Qué caracteriza a las fundaciones?} &
Persona jurídica sin miembros, con fin de bien común y sin lucro, creada por el aporte patrimonial de uno o varios fundadores. Requiere autorización estatal para funcionar. Necesita patrimonio inicial suficiente (plan trienal). Órgano de gobierno: Consejo de Administración (mínimo tres personas). Solo fiscalización externa. El objeto es inmutable, salvo imposibilidad sobreviniente. \\[0.4em]

\textbf{¿Qué caracteriza a las asociaciones civiles?} &
Corporación de personas con fin común no lucrativo. Instrumento constitutivo: escritura pública. Tres órganos: Comisión Directiva, Asamblea (suprema) y órgano de fiscalización (Comisión Revisora de Cuentas si hay más de 100 socios). Poder disciplinario: proceso con derecho de defensa, gradualidad y proporcionalidad. Apelación judicial solo sobre la forma, salvo grave violación de derechos humanos. \\[0.4em]

\textbf{¿En qué se diferencian las simples asociaciones?} &
No requieren autorización estatal. Pueden constituirse por instrumento privado con firma certificada. Comienzan su existencia desde la constitución. Si hay menos de 20 socios, pueden prescindir del órgano de fiscalización. Principal diferencia: en caso de insolvencia, los administradores y miembros que participaron en la gestión responden solidariamente con sus propios bienes. \\[0.4em]

\midrule
\multicolumn{2}{p{0.94\textwidth}}{
\textbf{Resumen:}
Las personas jurídicas son entes colectivos a los que el ordenamiento jurídico --no la naturaleza-- otorga aptitud para ser titulares de derechos y obligaciones, con el fin de cumplir el objeto para el que fueron creadas (principio de especialidad). El Código Civil y Comercial las clasifica en públicas --el Estado y sus organismos, los estados extranjeros, los organismos internacionales y la Iglesia Católica-- y privadas --sociedades, asociaciones civiles, simples asociaciones, fundaciones, mutuales, cooperativas y consorcios--. Su existencia comienza desde la constitución, aunque algunas requieren autorización estatal para funcionar. Son entidades con personalidad diferenciada de sus miembros (patrimonios y responsabilidades distintas), salvo los supuestos de inoponibilidad del artículo 144, donde la personería cede frente al uso abusivo o fraudulento. Responden civilmente por los daños que causen sus directivos en ejercicio de sus funciones (art. 1763), por las cosas y actividades riesgosas que posean (art. 1757) y por los hechos dañosos de sus dependientes (art. 1753); además, los propios administradores responden ilimitada y solidariamente frente a la persona jurídica (art. 160). Su disolución puede deberse a la voluntad de los miembros, al cumplimiento o imposibilidad del objeto, al vencimiento del plazo, a la quiebra o a la revocación estatal --esta última con estrictas garantías de fundamentación y defensa--. Entre los tipos especiales, las fundaciones se distinguen por la ausencia de miembros, la afectación de un patrimonio a un fin altruista y el control externo exclusivo de la autoridad de contralor; las asociaciones civiles, por su carácter corporativo, sus tres órganos de gobierno y el poder disciplinario sobre sus miembros con plenas garantías procesales; y las simples asociaciones, por no requerir autorización estatal y por la responsabilidad solidaria que recae sobre los administradores y miembros activos ante la insolvencia.
} \\

\end{longtable}
\partintro{Esta parte conecta dos preguntas fundamentales del derecho privado: ¿quién responde? y ¿con qué responde? El patrimonio es el puente entre ambas. Se estudia el patrimonio, la clasificación de los derechos, los bienes y las cosas que lo componen y, finalmente, la garantía patrimonial de los acreedores, los bienes que quedan excluidos de ella y las acciones destinadas a conservarla.}
\part{El objeto de la relación jurídica: patrimonio, derechos y bienes}
\chapter{El patrimonio}

\section{Planteo general}

Esta parte conecta dos preguntas fundamentales del Derecho Privado: \textbf{¿quién responde?} y \textbf{¿con qué responde?} El patrimonio es el puente entre ambas: es el atributo que le da contenido económico a la persona jurídica y, al mismo tiempo, la garantía sobre la que los acreedores ejecutarán sus créditos. Entender qué son los bienes y las cosas ---y cómo se clasifican--- es entender de qué está hecho ese patrimonio.

\begin{note}
\textbf{Relevancia práctica.} Estas nociones constituyen la base de la práctica profesional jurídica cotidiana:
\begin{itemize}
    \item Al redactar una escritura de compraventa, es necesario distinguir si el bien es mueble o inmueble para determinar la forma requerida (escritura pública o simple contrato).
    \item Al asesorar a un cliente endeudado, permite determinar qué bienes puede conservar por ser inembargables.
    \item En un litigio, permite identificar qué acciones utilizar para recuperar el patrimonio de un deudor que intenta eludir a sus acreedores.
\end{itemize}
El patrimonio como garantía es el concepto central del crédito: no se presta dinero ni se celebra un contrato sin entender sobre qué responde el deudor.
\end{note}

\section{El patrimonio como atributo de la personalidad}

En este capítulo confluyen dos aspectos distintos y complementarios. El primero es el estudio del \textbf{patrimonio como atributo de la personalidad}. Existen distintas posturas: algunas señalan que el patrimonio es un atributo de la personalidad y otras no lo entienden así. Se adopta aquí ---junto con el profesor Lozano Roy--- la posición que entiende el patrimonio como atributo de la personalidad.

La materia aborda el patrimonio desde dos planos:

\begin{tabla}
\begin{tabularx}{\textwidth}{B{3.4cm}YY}
\toprule
Perspectiva & \textbf{Qué estudia} & \textbf{Dónde se ubica} \\
\midrule
Atributo de la personalidad & El patrimonio como cualidad inherente al sujeto de derecho & Dentro del estudio del sujeto de la relación jurídica \\
Objeto de la relación jurídica & Los bienes y cosas que conforman el patrimonio & Dentro del estudio del objeto de la relación jurídica \\
\bottomrule
\end{tabularx}
\end{tabla}

La materia recorre tres grandes ejes: el \textbf{sujeto} (persona humana y jurídica), el \textbf{objeto} (bienes y cosas) y la \textbf{causa} de la relación jurídica. El patrimonio aparece en los dos primeros ejes: como atributo del sujeto y como continente de los objetos.

\begin{norma}{Artículo 15 CCCN --- Titularidad de derechos individuales}
Las personas son titulares de los derechos individuales sobre los bienes que integran su patrimonio conforme a lo que se establece en este Código.
\end{norma}

\section{El patrimonio como garantía común de los acreedores}

El artículo 242 constituye el núcleo funcional del patrimonio:

\begin{norma}{Artículo 242 CCCN --- Garantía común}
Todos los bienes del deudor están afectados al cumplimiento de sus obligaciones y constituyen la garantía común de sus acreedores, con excepción de aquellos que este Código o leyes especiales declaran inembargables e inejecutables.
\end{norma}

\begin{example}{La deuda de un millón de pesos}
Supóngase que una persona debe un millón de pesos. El acreedor la intima a pagar y el deudor no paga. El acreedor no puede limitarse a solicitar cortésmente el pago: en algún momento deberá ejecutar compulsivamente al deudor.

Si el deudor no entrega voluntariamente el dinero, no realiza una transferencia bancaria ni cumple de ninguna manera la obligación, el acreedor se cobra \textbf{ejecutando su patrimonio}. Si el deudor tiene un auto a su nombre y un departamento, el acreedor traba embargo sobre esos bienes y, si el deudor no paga, los ejecuta judicialmente; con el producido de los bienes se cobra.
\end{example}

\section{El patrimonio como universalidad de derecho}

\begin{definition}{Patrimonio (universalidad de derecho)}
Según la teoría clásica ---Llambías y otros autores---, el patrimonio es una \textbf{universalidad de derecho}: un conjunto de bienes y deudas que conforman una unidad porque la ley se la otorga, con independencia de su contenido concreto en cada momento.
\end{definition}

Se destacan dos notas esenciales:

\begin{tabla}
\begin{tabularx}{\textwidth}{B{3.0cm}YY}
\toprule
Nota & \textbf{Descripción} & \textbf{Consecuencia práctica} \\
\midrule
Unicidad & Cada persona tiene un solo patrimonio & No puede fragmentarlo a voluntad, salvo casos legales \\
Mutabilidad del contenido & Los bienes entran y salen, pero el patrimonio subsiste & Los acreedores se cobran sobre los bienes actuales, no sobre los que existían al contraer la deuda \\
\bottomrule
\end{tabularx}
\end{tabla}

Los acreedores pueden cobrarse con los bienes actuales del deudor. El hecho de que el deudor haya vendido un departamento no implica que los acreedores hayan perdido su garantía: la garantía es el \textbf{patrimonio como tal}, con independencia de su contenido. Esto indica también que, si el contenido del patrimonio de la persona va cambiando con el tiempo, sigue siendo su patrimonio.

\subsection{Los patrimonios especiales}

El Código Civil y Comercial reconoce la existencia de patrimonios especiales, disipando la ambigüedad que existía en el código anterior. Son universalidades autónomas que conviven con el patrimonio general:

\begin{tabla}
\begin{tabularx}{\textwidth}{B{4.6cm}Y}
\toprule
Tipo de patrimonio especial & \textbf{Descripción} \\
\midrule
Fondo de comercio & Conjunto de bienes afectados a una explotación comercial \\
Masa de la quiebra & Bienes del fallido afectados al proceso concursal \\
Masa hereditaria & Bienes del causante antes de la partición \\
Herencia del presunto fallecido & Patrimonio sujeto a prenotación durante el período legal \\
\bottomrule
\end{tabularx}
\end{tabla}

\section{Cuadro Cornell: El patrimonio}

\cornellinicio
\begin{longtable}{|Q|Z|}
\hline
\rowcolor{gray!15}\textbf{Preguntas / palabras clave} & \textbf{Notas / respuestas} \\ \hline
\endfirsthead
\hline
\rowcolor{gray!15}\textbf{Preguntas / palabras clave} & \textbf{Notas / respuestas} \\ \hline
\endhead

\cornellpregunta{¿Qué es el patrimonio y por qué se dice que es una universalidad de derecho?}{%
Es una universalidad de derecho: el conjunto de bienes (activo) y deudas (pasivo) de una persona, al que el ordenamiento jurídico le otorga unidad. Es una unidad porque la ley se la otorga, con independencia de su contenido concreto en cada momento (teoría clásica: Llambías y otros autores). En la materia se lo entiende, junto con el profesor Lozano Roy, como atributo de la personalidad.}

\cornellpregunta{¿Cuáles son las notas esenciales del patrimonio?}{%
\textbf{Unicidad:} cada persona tiene un solo patrimonio, que no puede fragmentar a voluntad salvo casos legales.

\textbf{Mutabilidad del contenido:} los bienes entran y salen, pero el patrimonio subsiste; por eso los acreedores se cobran sobre los bienes actuales y no sobre los que existían al contraer la deuda.}

\cornellpregunta{¿Puede existir más de un patrimonio para una misma persona?}{%
Cada persona tiene un solo patrimonio (unicidad), salvo casos legales. El CCCN reconoce patrimonios especiales, universalidades autónomas que conviven con el patrimonio general: fondo de comercio, masa de la quiebra, masa hereditaria y herencia del presunto fallecido.}

\cornellpregunta{¿Por qué se dice que el patrimonio es la garantía común de los acreedores?}{%
Porque todos los bienes del deudor responden por sus obligaciones (Arts. 242 y 743 CCCN). Si el deudor no paga voluntariamente, el acreedor puede ejecutar compulsivamente los bienes que integran su patrimonio, con excepción de los declarados inembargables.}

\cornellpregunta{¿Desde qué perspectivas se estudia el patrimonio?}{%
Como atributo de la personalidad (dentro del estudio del sujeto de la relación jurídica) y como conjunto de los bienes y cosas que lo conforman (dentro del estudio del objeto de la relación jurídica). Los tres ejes de la materia son el sujeto, el objeto y la causa; el patrimonio aparece en los dos primeros.}

\cornellresumen{El patrimonio es una universalidad de derecho ---conjunto de bienes y deudas--- y un atributo de la personalidad. Es único y de contenido cambiante, aunque el CCCN admite patrimonios especiales. Es la garantía común de los acreedores (Art. 242 CCCN): todos los bienes del deudor responden por sus obligaciones, salvo los declarados inembargables e inejecutables, y el acreedor puede cobrarse ejecutándolos cuando el deudor no paga voluntariamente.}

\end{longtable}
\cornellfin

% =====================================================================
% CAPÍTULO 2 — CLASIFICACIÓN DE LOS DERECHOS
% =====================================================================
\chapter{Clasificación de los derechos}

\section{Derechos patrimoniales y extrapatrimoniales}

La gran clasificación de los derechos distingue los que componen el patrimonio de los que quedan fuera de él:

\begin{tabla}
\begin{tabularx}{\textwidth}{B{4.0cm}YY}
\toprule
Categoría & \textbf{Subcategorías} & \textbf{Caracteres} \\
\midrule
Extrapatrimoniales & Derechos personalísimos; derechos de familia & Sin contenido económico directo; ligados a la persona o al vínculo familiar \\
Patrimoniales --- Reales & Dominio, condominio, usufructo, uso, habitación, servidumbre, hipoteca, prenda & Relación directa entre la persona y la cosa; oponibles \textit{erga omnes} \\
Patrimoniales --- Personales & Obligaciones y contratos & Relación entre acreedor y deudor respecto de una prestación \\
Patrimoniales --- Intelectuales & Obras científicas, literarias, artísticas, invenciones, software & Recaen sobre creaciones inmateriales \\
\bottomrule
\end{tabularx}
\end{tabla}

\section{Los derechos reales}

\begin{definition}{Derechos reales}
Los derechos reales son aquellos que conceden un \textbf{señorío inmediato sobre una cosa}: existe una relación directa de la persona con la cosa.
\end{definition}

El derecho real por antonomasia es el \textbf{dominio}, que es el derecho más perfecto y confiere el uso y goce de una cosa. Técnicamente, en el Código Civil y Comercial ---cuya tradición proviene de los romanos---, el \textit{dominus} es el señor de una \textit{res}; por eso estos derechos son «reales»: el término viene del latín \textit{res}, que significa «cosa».

\begin{tabla}
\begin{tabularx}{\textwidth}{B{4.4cm}P{4.2cm}Y}
\toprule
Tipo & \textbf{Ejemplos} & \textbf{Nota} \\
\midrule
Sobre cosa propia --- plenos & Dominio & El más amplio: uso, goce y disposición \\
Sobre cosa propia --- limitados & Condominio & Dos o más titulares sobre la misma cosa \\
Sobre cosa ajena --- de disfrute & Usufructo, uso, habitación, servidumbre & El titular goza de la cosa de otro \\
Sobre cosa ajena --- de garantía & Hipoteca (inmuebles), prenda (muebles) & Garantizan el cumplimiento de una obligación \\
\bottomrule
\end{tabularx}
\end{tabla}

\section{Los derechos personales}

\begin{important}
No deben confundirse los derechos \textbf{personales} con los derechos \textbf{personalísimos}.
\end{important}

Los derechos se llaman «personales» porque confieren prerrogativas para exigir de \textbf{otro} una determinada prestación: la relación no es directa entre una cosa y su dueño, sino entre un acreedor y un deudor en torno a una prestación. Estos derechos se plasman fundamentalmente en obligaciones y contratos.
% TODO: contenido incompleto o ambiguo en las notas originales (la frase original sobre el sujeto que «se plasma fundamentalmente en obligaciones y contratos» resulta ininteligible; se conservó únicamente la referencia a obligaciones y contratos, que coincide con el cuadro de clasificación)

\section{Los derechos intelectuales}

\begin{definition}{Derechos intelectuales}
Los derechos intelectuales confieren un \textbf{derecho de explotación, de uso y de goce sobre una obra intelectual}: una obra científica, la creación de una invención, la obra literaria o artística, el software.
\end{definition}

Se trata de distintos tipos de creaciones que tienen un campo de inmaterialidad, porque no son una casa, un auto o una computadora, sino obras intelectuales.

\section{Cuadro Cornell: Clasificación de los derechos}

\cornellinicio
\begin{longtable}{|Q|Z|}
\hline
\rowcolor{gray!15}\textbf{Preguntas / palabras clave} & \textbf{Notas / respuestas} \\ \hline
\endfirsthead
\hline
\rowcolor{gray!15}\textbf{Preguntas / palabras clave} & \textbf{Notas / respuestas} \\ \hline
\endhead

\cornellpregunta{¿Cómo se clasifican los derechos según su relación con el patrimonio?}{%
\textbf{Extrapatrimoniales} (personalísimos y de familia): sin contenido económico directo, ligados a la persona o al vínculo familiar; quedan fuera del patrimonio.

\textbf{Patrimoniales:} reales, personales e intelectuales.}

\cornellpregunta{¿Qué son los derechos reales y cuál es el derecho real por antonomasia?}{%
Conceden un señorío inmediato sobre una cosa: hay una relación directa de la persona con la cosa, y son oponibles \textit{erga omnes}. El derecho real por antonomasia es el dominio, el derecho más perfecto, que confiere el uso y goce de una cosa. Se los llama «reales» por el latín \textit{res} (cosa).}

\cornellpregunta{¿Qué diferencia hay entre ser dueño de una cosa y tener un derecho sobre ella?}{%
El dominio es un derecho real sobre cosa propia, pleno y el más amplio (uso, goce y disposición). Los derechos reales sobre cosa ajena son de disfrute (usufructo, uso, habitación, servidumbre: el titular goza de la cosa de otro) o de garantía (hipoteca sobre inmuebles, prenda sobre muebles). El condominio es un derecho sobre cosa propia limitado: dos o más titulares sobre la misma cosa.}

\cornellpregunta{¿Por qué se llaman «personales» los derechos personales si no son personalísimos?}{%
Porque confieren prerrogativas para exigir de otro una determinada prestación: la relación no es entre una cosa y su dueño, sino entre un acreedor y un deudor en torno a una prestación. Se plasman en obligaciones y contratos. No deben confundirse con los derechos personalísimos, que son extrapatrimoniales.}

\cornellpregunta{¿Qué son los derechos intelectuales?}{%
Derechos de explotación, de uso y de goce sobre una obra intelectual (obra científica, invención, obra literaria o artística, software). Recaen sobre creaciones con un campo de inmaterialidad.}

\cornellresumen{Los derechos se clasifican en extrapatrimoniales (personalísimos y de familia, sin contenido económico directo) y patrimoniales (reales, personales e intelectuales). Los derechos reales otorgan un señorío inmediato sobre una cosa y son oponibles \textit{erga omnes}; el dominio es el derecho real por antonomasia. Los derechos personales confieren la facultad de exigir de otro una prestación y se plasman en obligaciones y contratos; no deben confundirse con los personalísimos. Los derechos intelectuales recaen sobre creaciones inmateriales.}

\end{longtable}
\cornellfin

% =====================================================================
% CAPÍTULO 3 — BIENES Y COSAS
% =====================================================================
\chapter{Bienes y cosas}

\section{Distinción entre bienes y cosas}

\begin{norma}{Artículo 16 CCCN --- Bienes y cosas}
Los derechos referidos en el primer párrafo del artículo 15 pueden recaer sobre bienes susceptibles de valor económico. Los bienes materiales se llaman cosas. Las disposiciones referentes a las cosas son aplicables a la energía y a las fuerzas naturales susceptibles de ser puestas al servicio del hombre.
\end{norma}

\begin{tabla}
\begin{tabularx}{\textwidth}{B{3.4cm}YY}
\toprule
Concepto & \textbf{Definición} & \textbf{Ejemplos} \\
\midrule
Bien (sentido amplio) & Todo objeto susceptible de valor económico, material o inmaterial & Un crédito, un derecho de autor, una casa, dinero \\
Cosa & Objeto material susceptible de valor económico & Casa, auto, computadora, ganado \\
Bien inmaterial & Bien sin corporalidad física & Un crédito, una marca, un derecho de usufructo \\
\bottomrule
\end{tabularx}
\end{tabla}

\section{Muebles e inmuebles}

La distinción entre cosas inmuebles y cosas muebles es la clasificación más importante de las cosas en el Código.

\subsection{Inmuebles}

\begin{tabla}
\begin{tabularx}{\textwidth}{B{2.8cm}YY}
\toprule
Tipo de inmueble & \textbf{Descripción} & \textbf{Ejemplos} \\
\midrule
Por naturaleza & Fijados al suelo de manera natural o que se encuentran bajo él sin hecho del hombre & El suelo, los minerales, los árboles arraigados \\
Por accesión & Cosas muebles inmovilizadas por su adhesión física al suelo con carácter perdurable & Paredes, pisos, ventanas, puertas, techos de un edificio \\
\bottomrule
\end{tabularx}
\end{tabla}

\begin{example}{Límite de la accesión}
Los pisos, las paredes, las ventanas, las puertas y los techos de una casa están compuestos por cosas que eran muebles pero que fueron incorporadas y fijadas al piso con carácter perdurable, por lo que se convierten en inmuebles. La ventana es separable, pero es parte del inmueble.

En cambio, una biblioteca colocada en un departamento no es parte del inmueble aunque esté puesta con carácter permanente, porque está afectada a la actividad de su titular: si se vende el departamento, se lo vende sin la biblioteca, salvo que se trate de un placar construido directamente e incorporado con las paredes.
\end{example}

\subsection{Cosas muebles}

\begin{tabla}
\begin{tabularx}{\textwidth}{B{3.8cm}YY}
\toprule
Tipo de cosa mueble & \textbf{Descripción} & \textbf{Ejemplos} \\
\midrule
Por desplazamiento propio (semovientes) & Se mueven por sí mismas & Ganado \\
Por fuerza externa & Se desplazan mediante fuerza humana o mecánica & Muebles del hogar, electrodomésticos \\
Locomóviles & Vehículos con propulsión propia & Automóviles, motos, camiones, aeronaves \\
Muebles registrables & Muebles que por su importancia la ley manda registrar & Automotores, buques, aeronaves \\
Muebles no registrables & Sin obligación de registro & La mayoría de los objetos cotidianos \\
\bottomrule
\end{tabularx}
\end{tabla}

\subsection{Importancia práctica de la distinción entre muebles e inmuebles}

\begin{tabla}
\begin{tabularx}{\textwidth}{B{3.4cm}YY}
\toprule
Aspecto jurídico & \textbf{Cosas inmuebles} & \textbf{Cosas muebles} \\
\midrule
Forma de transmisión & Escritura pública obligatoria & Mayor libertad de formas \\
Ley aplicable (DIPr) & Ley del territorio donde está el bien & Ley del domicilio del propietario \\
Prescripción adquisitiva & 20 años (10 con justo título y buena fe) & Plazo más corto; sin registro: posesión vale título \\
Garantías reales & Hipoteca & Prenda \\
\bottomrule
\end{tabularx}
\end{tabla}

\section{Cosas divisibles e indivisibles}

\begin{definition}{Cosas divisibles}
Son divisibles las cosas que pueden dividirse en porciones reales sin ser destruidas, y cada una de las cuales forma un todo homogéneo y análogo tanto a las otras partes como a la cosa original.
\end{definition}

\begin{example}{División de una cosa}
Si deben hacerse, a partir de diez quintales de soja, dos entregas de cinco quintales, la cosa es perfectamente divisible: cada parte es un todo homogéneo. En cambio, si se divide un auto, no quedan dos autos: queda un auto destruido.
\end{example}

La divisibilidad tiene un límite legal: la división no puede tornar antieconómico el uso y aprovechamiento de la cosa.

\begin{example}{División de un lote}
Dividir un lote de diez metros por diez en pequeños lotes de un metro cuadrado sería absolutamente inconveniente y antieconómico.
\end{example}

Por eso la reglamentación en materia de inmuebles fija unidades mínimas ---en algunos lugares, de 866 metros cuadrados---. Las autoridades locales (municipios o provincias) tienen las potestades para reglamentar las divisiones en materia de inmuebles.

\section{Cosas principales y accesorias}

\begin{definition}{Cosas principales y accesorias}
Son \textbf{principales} las cosas que pueden existir por sí mismas; \textbf{accesorias}, aquellas cuya existencia y naturaleza son determinadas por otra de la cual dependen o están adheridas.
\end{definition}

El régimen jurídico de la cosa accesoria es el de la principal.

\begin{example}{El cuadro y su marco}
Sea un cuadro compuesto por la tela, la pintura y un marco: el marco es accesorio con relación a la tela y a la pintura, y el régimen jurídico aplicable es el de la principal. Pero si se trata de una lámina que se compra en cualquier kiosco y se la coloca en un marco enormemente valioso de oro, lo principal pasa a ser el marco, y la lámina colocada en él resulta absolutamente accesoria.
\end{example}

\begin{norma}{Artículo 230 CCCN --- Criterio residual}
Si no es posible determinar la relación de principalidad y accesoriedad entre las cosas, se considera principal la de mayor valor. Si los valores son iguales, no hay cosa principal ni accesoria.
\end{norma}

\section{Cosas consumibles y no consumibles; fungibles y no fungibles}

\begin{tabla}
\begin{tabularx}{\textwidth}{B{3.0cm}YYY}
\toprule
Clasificación & \textbf{Definición} & \textbf{Ejemplos} & \textbf{Importancia} \\
\midrule
Consumibles & Su existencia termina con el primer uso & Alimentos, combustible & Relevante en obligaciones de restitución \\
No consumibles & No dejan de existir por el uso habitual & Muebles, ropa (se deterioran, pero subsisten) & Base de contratos de uso y préstamo \\
Fungibles & Un individuo equivale a otro de la misma especie; pueden sustituirse & Toneladas de trigo de igual calidad, billetes & Admite sustitución en el cumplimiento \\
No fungibles & No pueden sustituirse por otro individuo & Cuadro de artista reconocido, botella de vino numerada & Incumplimiento si se entrega cosa distinta \\
\bottomrule
\end{tabularx}
\end{tabla}

Las obras artísticas no son fungibles; las cosas producidas industrialmente, por lo general, sí lo son.

\section{Frutos y productos}

\begin{tabla}
\begin{tabularx}{\textwidth}{B{3.2cm}YYY}
\toprule
Categoría & \textbf{Definición} & \textbf{Renovación} & \textbf{Ejemplos} \\
\midrule
Frutos naturales & Producidos espontáneamente por la naturaleza & Se renuevan periódicamente & Manzanas de un árbol \\
Frutos industriales & Producidos con intervención humana (agricultura, industria) & Se renuevan con la actividad & Cosechas, producción fabril \\
Frutos civiles & Rentas producidas por el bien a través de actos jurídicos & Se renuevan periódicamente & Alquileres, intereses del dinero prestado \\
Productos & Objetos no renovables que al separarse alteran o disminuyen la sustancia & No se renuevan & Minerales extraídos de una cantera \\
\bottomrule
\end{tabularx}
\end{tabla}

\begin{example}{Producto y fruto}
Si se extrae cobre de una cantera, se trata de un \textbf{producto}: la cantera no volverá a producir ese mineral sino en dos millones de años, de modo que se extraen productos y se disminuye el valor de la cosa. En cambio, las manzanas de un árbol se renuevan: son un \textbf{fruto}.
\end{example}

Esta distinción adquiere importancia cuando debe restituirse algo con sus frutos y sus productos.

\begin{example}{Restitución de una casa por un poseedor de mala fe}
Quien recibe una casa y es de mala fe debe restituirla con las rentas: si la casa producía un alquiler mensual de veinte mil pesos, debe devolver la casa con las rentas que hubiese producido durante esos años.
\end{example}

Los artículos 1935 y 1936 CCCN (y el artículo 234, sobre bienes fuera del comercio) regulan los deberes de restitución de frutos y productos según la buena o mala fe del poseedor.
% TODO: contenido incompleto o ambiguo en las notas originales (no se explica por qué el Art. 234 se menciona junto a los Arts. 1935 y 1936; se conservó la referencia tal como figura)

\section{Bienes según las personas}

\begin{tabla}
\begin{tabularx}{\textwidth}{B{2.9cm}P{1.7cm}YY}
\toprule
Categoría & \textbf{Artículo CCCN} & \textbf{Caracteres} & \textbf{Ejemplos} \\
\midrule
Dominio público del Estado & Art. 235 & Inalienables, inembargables, imprescriptibles; uso de todos & Ríos, lagos, playas, plazas, calles \\
Dominio privado del Estado & Art. 236 & El Estado es titular como particular; pueden enajenarse & Inmueble fiscal no afectado a uso público \\
Bienes de los particulares & Art. 238 & Todos los bienes no enumerados en los artículos anteriores & Cualquier bien en manos de un privado \\
Bienes fuera del comercio & Art. 234 & Prohibida su enajenación por ley o por acto jurídico & Bienes de uso cultual, bienes embargados \\
Bienes de incidencia colectiva & Art. 240 & Interés colectivo en su protección; no apropiables individualmente & Ecosistema, fauna, biodiversidad, agua, valores culturales \\
\bottomrule
\end{tabularx}
\end{tabla}

\subsection{Bienes de incidencia colectiva}

\begin{definition}{Bienes de incidencia colectiva}
Son bienes que no son estrictamente individuales, sino que corresponden de forma homogénea a un conjunto de personas o a la comunidad como tal. Por ejemplo: el ecosistema, la flora, la fauna, la biodiversidad, el agua y los valores culturales.
\end{definition}

\section{Cuadro Cornell: Bienes y cosas}

\cornellinicio
\begin{longtable}{|Q|Z|}
\hline
\rowcolor{gray!15}\textbf{Preguntas / palabras clave} & \textbf{Notas / respuestas} \\ \hline
\endfirsthead
\hline
\rowcolor{gray!15}\textbf{Preguntas / palabras clave} & \textbf{Notas / respuestas} \\ \hline
\endhead

\cornellpregunta{¿Qué diferencia hay entre un bien y una cosa?}{%
Todo bien es un objeto susceptible de valor económico (material o inmaterial). Una cosa es un bien material (objeto corporal). Todos los objetos inmateriales con valor económico son bienes, pero no cosas (Art. 16 CCCN).}

\cornellpregunta{¿Cuál es la clasificación más importante de las cosas?}{%
Muebles e inmuebles. Los inmuebles pueden ser por naturaleza (el suelo y lo adherido orgánicamente) o por accesión (cosas muebles incorporadas con carácter perdurable). Los muebles pueden desplazarse por sí mismos (semovientes) o por fuerza externa; algunos son registrables (automotores).}

\cornellpregunta{¿Por qué importa distinguir muebles de inmuebles?}{%
Define la ley aplicable (en DIPr), el plazo de prescripción adquisitiva (20/10 años para inmuebles \textit{vs.} plazos menores para muebles), la forma de transmisión (escritura pública para inmuebles) y el tipo de garantía real (hipoteca \textit{vs.} prenda).}

\cornellpregunta{¿Cuándo una cosa es divisible?}{%
Cuando puede dividirse en porciones reales sin ser destruida y cada porción forma un todo homogéneo y análogo a las otras partes y a la cosa original (diez quintales de soja en dos entregas de cinco; un auto no es divisible: dividido queda destruido). Límite: la división no puede tornar antieconómico el uso y aprovechamiento; en inmuebles, las autoridades locales fijan unidades mínimas.}

\cornellpregunta{¿Cuándo una cosa es principal o accesoria?}{%
Es principal la que puede existir por sí sola; accesoria, la que depende de otra. El régimen jurídico sigue a lo principal. Si no se puede determinar cuál es principal, prevalece la de mayor valor (Art. 230 CCCN); si los valores son iguales, no hay cosa principal ni accesoria.}

\cornellpregunta{¿Qué diferencia hay entre cosas consumibles y no consumibles, y entre fungibles y no fungibles?}{%
\textbf{Consumibles:} su existencia termina con el primer uso (alimentos, combustible). \textbf{No consumibles:} no dejan de existir por el uso habitual.

\textbf{Fungibles:} un individuo equivale a otro de la misma especie y pueden sustituirse, lo que admite sustitución en el cumplimiento (trigo de igual calidad, billetes). \textbf{No fungibles:} no pueden sustituirse; entregar una cosa distinta es incumplimiento (cuadro de artista reconocido).}

\cornellpregunta{¿Qué diferencia a un fruto de un producto?}{%
El fruto se produce de modo renovable sin alterar ni disminuir la sustancia de la cosa (manzanas, intereses, alquileres). El producto es no renovable y su separación altera o disminuye la sustancia (minerales extraídos de una cantera).}

\cornellpregunta{¿Por qué importa distinguir entre frutos y productos?}{%
Interesa cuando debe restituirse algo con sus frutos y productos: los Arts. 1935 y 1936 CCCN regulan los deberes de restitución según la buena o mala fe del poseedor. Quien recibe una casa de mala fe debe restituirla con las rentas que hubiese producido.}

\cornellpregunta{¿Cómo se clasifican los bienes según las personas?}{%
\begin{itemize}[nosep,leftmargin=1.2em]
\item \textbf{Dominio público del Estado} (Art. 235): inalienables, inembargables, imprescriptibles; uso de todos.
\item \textbf{Dominio privado del Estado} (Art. 236): el Estado es titular como particular; pueden enajenarse.
\item \textbf{Bienes de los particulares} (Art. 238): los no enumerados en los artículos anteriores.
\item \textbf{Bienes fuera del comercio} (Art. 234): prohibida su enajenación por ley o por acto jurídico.
\item \textbf{Bienes de incidencia colectiva} (Art. 240): interés colectivo en su protección; no apropiables individualmente.
\end{itemize}}

\cornellpregunta{¿Qué son los bienes de incidencia colectiva?}{%
Bienes no estrictamente individuales, que corresponden de forma homogénea a un conjunto de personas o a la comunidad como tal (ecosistema, flora, fauna, biodiversidad, agua, valores culturales).}

\cornellresumen{Todo bien es un objeto susceptible de valor económico; las cosas son los bienes materiales (Art. 16 CCCN). Las cosas se clasifican en muebles e inmuebles (la distinción más importante), divisibles e indivisibles, principales y accesorias, consumibles y no consumibles, fungibles y no fungibles, y se distingue entre frutos y productos. Según las personas, los bienes pueden ser del dominio público o privado del Estado, de los particulares, fuera del comercio o de incidencia colectiva. Cada clasificación tiene consecuencias jurídicas precisas.}

\end{longtable}
\cornellfin

% =====================================================================
% CAPÍTULO 4 — GARANTÍA PATRIMONIAL, ACREEDORES Y ACCIONES
% =====================================================================
\chapter{Garantía patrimonial, acreedores y acciones}
\label{ch:garantia}

\section{El patrimonio como garantía: repaso y profundización}

El artículo 743 precisa cuáles son los bienes que constituyen la garantía:

\begin{norma}{Artículo 743 CCCN --- Bienes que constituyen la garantía}
Son bienes del deudor sujetos a ejecución todos los que integran su patrimonio a la fecha de la sentencia o en el que se hubiese determinado, y los que adquiera hasta la extinción de la obligación. El acreedor puede exigir la venta judicial de los bienes del deudor, pero sólo en la medida necesaria para satisfacer su crédito. Todos los acreedores pueden ejecutar estos bienes en posición igualitaria, excepto que exista una causa legal de preferencia.
\end{norma}

\section{Tipos de acreedores y preferencias de cobro}

Cuando hay concurrencia de acreedores ---especialmente en situaciones de insolvencia--- el orden de cobro no es necesariamente igualitario.

\begin{tabla}
\begin{tabularx}{\textwidth}{B{2.7cm}YYY}
\toprule
Tipo de acreedor & \textbf{Fuente de preferencia} & \textbf{Alcance} & \textbf{Ejemplo} \\
\midrule
Con privilegio especial & La ley & Sobre determinados bienes & Acreedor laboral sobre maquinarias del empleador en quiebra \\
Con privilegio general & La ley & Sobre la generalidad del patrimonio & Créditos del fisco en algunos supuestos \\
Con garantía real & Autonomía de la voluntad (contrato) & Sobre el bien gravado & Hipotecario que ejecuta el inmueble hipotecado \\
Quirografario (común) & Ninguna preferencia & Sobre el remanente, en igualdad & Acreedor sin ninguna garantía \\
\bottomrule
\end{tabularx}
\end{tabla}

\begin{example}{Concurrencia de acreedores}
Supóngase que el deudor debe 10.000 dólares a un acreedor, 20.000 a otro y 30.000 a un tercero. Tiene un único bien: un departamento valuado en 50.000 dólares.

\textbf{Escenario A (todos quirografarios).} El pasivo total es de 60.000 y el activo de 50.000: hay insolvencia. Cada acreedor cobrará proporcionalmente sobre los 50.000.

\textbf{Escenario B (uno tiene hipoteca por 50.000).} El acreedor hipotecario ejecuta primero el departamento, cobra sus 50.000 y los otros dos no cobran nada, porque el bien con garantía real tiene preferencia.
\end{example}
% TODO: contenido incompleto o ambiguo en las notas originales (en el escenario B la hipoteca es «por 50.000», pero las deudas individuales planteadas son de 10.000, 20.000 y 30.000; se conservó el planteo tal como figura)

\section{Bienes inembargables (Art. 744 CCCN)}

Razones humanitarias justifican la exclusión de ciertos bienes de la garantía común.

\begin{tabla}
\begin{tabularx}{\textwidth}{P{8.4cm}Y}
\toprule
\textbf{Bien inembargable} & \textbf{Fundamento} \\
\midrule
Ropa, muebles e inmuebles de uso indispensable del deudor, cónyuge, conviviente e hijos & Necesidades básicas del hogar \\
Instrumentos necesarios para ejercer la profesión, arte u oficio & Sin ellos, la persona no tendría manera de sobrevivir \\
Sepulcros afectados a su destino (salvo que la deuda sea por su construcción o reparación) & Respeto a la memoria de los difuntos \\
Bienes afectados a cualquier religión reconocida por el Estado (templos, objetos de culto) & Libertad religiosa \\
Derechos de usufructo, uso y habitación (en los términos del Código) & Carácter personalísimo de estos derechos \\
Indemnizaciones por daño moral o lesiones a la integridad psicofísica & Compensación de daños personalísimos \\
Indemnizaciones por alimentos al cónyuge, conviviente e hijos en caso de homicidio & Protección de la familia del fallecido \\
Otros declarados inembargables por leyes especiales (por ejemplo, el salario mínimo vital y móvil) & Sustento mínimo para la vida \\
\bottomrule
\end{tabularx}
\end{tabla}

\begin{important}
El telón de fondo de todas estas exclusiones es proteger la dignidad humana en sus dimensiones más fundamentales.
\end{important}

\section{La protección de la vivienda (Arts. 244 y ss. CCCN)}

El Código Civil y Comercial reemplazó la denominación «bien de familia» (Ley 14.394) por el concepto de «protección de la vivienda».

\begin{norma}{Artículo 244 CCCN --- Afectación}
Puede afectarse al régimen previsto en este Capítulo, un inmueble destinado a vivienda, por su totalidad o hasta una parte de su valor. Esta protección no excluye la concurrencia de otros acreedores.
\end{norma}

\begin{tabla}
\begin{tabularx}{\textwidth}{B{3.4cm}Y}
\toprule
Aspecto & \textbf{Régimen} \\
\midrule
Requisito & Inscripción en el Registro de la Propiedad Inmueble (trámite sencillo) \\
Efecto & El inmueble queda inembargable frente a las deudas posteriores a la inscripción \\
Deudas anteriores & Los acreedores con deudas anteriores no se ven afectados \\
Límite cuantitativo & Puede afectarse la totalidad o solo parte del valor \\
Límite de unidades & No puede afectarse más de un inmueble \\
\bottomrule
\end{tabularx}
\end{tabla}

\section{Acciones de los acreedores}

\subsection{Medidas preventivas}

Antes de llegar a la ejecución, los acreedores disponen de medidas preventivas para conservar la garantía:

\begin{tabla}
\begin{tabularx}{\textwidth}{B{3.4cm}YY}
\toprule
Medida & \textbf{Descripción} & \textbf{Ejemplo} \\
\midrule
Embargo preventivo & Afecta un bien determinado para que no salga del patrimonio del deudor & Trabar embargo sobre la casa del deudor que está por venderla antes de obtener sentencia \\
Inhibición general de bienes & Impide al deudor enajenar cualquier bien cuando no se conocen bienes concretos & Anotar en el registro una inhibición general cuando el deudor no tiene bienes identificables \\
\bottomrule
\end{tabularx}
\end{tabla}

Los deudores pueden plantearse la estrategia de sacar bienes de su patrimonio y transformarlos en dinero en efectivo, de tal manera que los acreedores, cuando vengan a cobrar, no encuentren bienes a su nombre. Estas medidas sirven para contrarrestar esa estrategia.

\subsection{Acción directa (Arts. 736 y 737 CCCN)}
% TODO: contenido incompleto o ambiguo en las notas originales (el índice del apunte menciona «Art. 737 y 739 CCCN» para las acciones directa y subrogatoria; el desarrollo cita «Art. 736 y ss.» para la acción directa y «Art. 739» para la subrogatoria)

\begin{definition}{Acción directa}
Es la acción que se le da al acreedor para \textbf{cobrar lo que un tercero debe a su deudor}, hasta el importe del propio crédito.
\end{definition}

Sus caracteres son los siguientes:

\begin{tabla}
\begin{tabularx}{\textwidth}{B{3.6cm}Y}
\toprule
Carácter & \textbf{Descripción} \\
\midrule
Propia & El acreedor la ejerce por derecho propio \\
Exclusiva & En beneficio exclusivo del acreedor que la ejerce, no de todos \\
Excepcional & Solo procede en los casos expresamente previstos por la ley \\
Restrictiva & De interpretación restrictiva \\
Presupone homogeneidad & Los créditos deben ser de la misma naturaleza (por ejemplo, dinerarios) \\
\bottomrule
\end{tabularx}
\end{tabla}

\begin{example}{Acción directa con créditos homogéneos}
Supóngase que se prestó dinero a una persona y que esta, a su vez, se lo prestó a una tercera. Los créditos son homogéneos, porque el propio dinero del acreedor pasó de uno a otro. El acreedor va directamente contra el deudor de su deudor y le ejecuta la deuda, pero solo hasta el importe de su crédito: si prestó doscientos mil pesos a su deudor y este le prestó a otro trescientos mil, puede reclamarle al tercero doscientos mil; los cien mil restantes se los reclamará su deudor.
\end{example}

Esta figura aparece también denominada en las fuentes originales «acción de delegación directa».
% TODO: verificar referencia presente en las fuentes originales (denominación «acción de delegación directa»).

\begin{formula}{Art.~737 CCCN --- Requisitos de la acción directa}
\begin{itemize}
    \item Existencia de un crédito exigible del acreedor contra su propio deudor.
    \item Existencia de una deuda correlativa exigible del tercero demandado a favor del deudor.
    \item Homogeneidad de ambos créditos entre sí.
    \item Ninguno de los dos créditos debe estar embargado ni ser indisponible.
    \item El deudor debe ser notificado de la demanda.
\end{itemize}
\end{formula}

\subsection{Acción subrogatoria (Arts. 739 y 740 CCCN)}

\begin{definition}{Acción subrogatoria}
El acreedor ocupa el lugar de su deudor y ejerce un derecho que le es propio a este, cuando el deudor está inactivo y esa omisión perjudica el cobro del crédito.
\end{definition}

\begin{example}{La cochera heredada}
Supóngase que se prestaron doscientos mil pesos a una persona. Esta tiene un familiar que fallece y le deja en herencia una cochera que vale cuatrocientos mil pesos. El deudor, considerando que el acreedor se quedaría con la cochera apenas la tuviera, decide no iniciar la sucesión y permanece inactivo.

El acreedor, que advierte que su deudor está inactivo y tiene la posibilidad de exigir judicialmente un crédito, y que esa inactividad lo perjudica, se subroga: ocupa su lugar e inicia la sucesión. La cochera ingresa al patrimonio del deudor y el acreedor traba embargo sobre ese bien para cobrarse. Pero aquí no tiene ningún privilegio especial: cualquier otro acreedor embargante tendrá también preferencia.
\end{example}

\begin{formula}{Arts.~739 y 740 CCCN --- Acción subrogatoria}
\textbf{Art.~739.} El acreedor de un crédito cierto, exigible o no, puede ejercer judicialmente los derechos patrimoniales de su deudor si este es remiso en hacerlo y esa omisión afecta el cobro de su acreencia. El acreedor no goza de preferencia alguna sobre los bienes obtenidos por ese medio.

\smallskip
\textbf{Art.~740.} Los derechos y acciones del deudor excluidos de la subrogación son los inherentes a su persona, los no patrimoniales, los patrimoniales no susceptibles de ejecución forzada y los excluidos por disposición legal.
\end{formula}

El crédito o la acción recuperados ingresan al patrimonio del deudor; por eso la acción subrogatoria beneficia a todos los acreedores, mientras que en la acción directa resulta beneficiado únicamente el acreedor que actúa.

\subsubsection{Comparación entre la acción directa y la acción subrogatoria}

\begin{tabla}
\begin{tabularx}{\textwidth}{B{3.4cm}YY}
\toprule
Criterio & \textbf{Acción directa} & \textbf{Acción subrogatoria} \\
\midrule
Nombre en el que se actúa & Nombre propio del acreedor & Nombre del deudor (subrogación) \\
Beneficiario & Solo el acreedor que la ejerce & Todos los acreedores del deudor \\
Privilegio & Sí: cobra directamente hasta su crédito & No: el bien ingresa al patrimonio; todos pueden ejecutarlo \\
Presupuesto & Créditos homogéneos y caso legal expreso & Inactividad del deudor que perjudica la acreencia \\
Límite & Hasta el importe del propio crédito & Todo lo recuperado ingresa al patrimonio del deudor \\
Normas & Arts. 736 y 737 CCCN & Arts. 739 y 740 CCCN \\
\bottomrule
\end{tabularx}
\end{tabla}

\begin{important}
La acción subrogatoria, a diferencia de la acción directa, no otorga ningún privilegio y beneficia a todos los acreedores.
\end{important}

\subsubsection{Derechos excluidos de la acción subrogatoria}

No se pueden subrogar derechos vinculados con la personalidad ni de ejercicio absolutamente discrecional por el titular (Arts. 741--742 CCCN). Como ejemplo se menciona el reconocimiento de un hijo o las cuestiones de puro ejercicio de la voluntad familiar. Según otra de las fuentes, los derechos excluidos son los inherentes a la persona del deudor, los no patrimoniales, los patrimoniales no susceptibles de ejecución forzada y los excluidos por disposición legal (art. 740 CCCN).
% TODO: verificar consistencia entre fuentes originales: una fuente cita los arts. 741 y 742 y otra el art. 740 para los derechos excluidos de la subrogación.


\subsection{Acción de simulación y acción revocatoria (inoponibilidad)}
% TODO: contenido incompleto o ambiguo en las notas originales (el índice del apunte asocia esta sección con «nulidad relativa», concepto que no se desarrolla en el cuerpo de las notas)

\begin{tabla}
\begin{tabularx}{\textwidth}{B{2.6cm}P{3.2cm}YY}
\toprule
Acción & \textbf{Denominación en el CCCN} & \textbf{Presupuesto} & \textbf{Efecto} \\
\midrule
De simulación & Nulidad por simulación & El deudor aparenta sacar bienes pero los conserva (por ejemplo, venta ficticia al hermano/a sin contraprestación) & Nulidad del acto; la cosa vuelve al patrimonio del deudor \\
Revocatoria / pauliana & Inoponibilidad (CCCN) & El deudor vende realmente su bien a un cómplice para defraudar a sus acreedores & El acto es válido entre partes pero inoponible al acreedor; puede ejecutar el bien \\
\bottomrule
\end{tabularx}
\end{tabla}

Ambas acciones se distinguen mediante casos paralelos.

\begin{example}{Simulación}
El deudor llama a su hermana y le dice «simulemos que te vendo mi casa». La hermana no tiene dinero ni bienes, pero aparece como compradora. Cuando el acreedor va a ejecutar al deudor, la casa no está a su nombre. El acreedor alega que el acto fue simulado; el juez decreta la nulidad, las cosas vuelven al estado anterior y el acreedor puede cobrarse.
\end{example}

\begin{example}{Inoponibilidad}
El deudor vende efectivamente su casa ---el dinero ingresa a su patrimonio---, pero el acreedor demuestra que el comprador era cómplice del deudor. Al revocar la venta, esta no se anula: es inoponible al acreedor defraudado, que ejecuta el bien como si la venta no hubiera existido. El adquirente cómplice pierde el bien, pero tiene acción para cobrarse contra el deudor originario.
\end{example}

\section{Cuadro Cornell: Garantía, acreedores y acciones}

\cornellinicio
\begin{longtable}{|Q|Z|}
\hline
\rowcolor{gray!15}\textbf{Preguntas / palabras clave} & \textbf{Notas / respuestas} \\ \hline
\endfirsthead
\hline
\rowcolor{gray!15}\textbf{Preguntas / palabras clave} & \textbf{Notas / respuestas} \\ \hline
\endhead

\cornellpregunta{¿Qué bienes del deudor están sujetos a ejecución?}{%
Según el Art. 743 CCCN, todos los que integran su patrimonio a la fecha de la sentencia (o en el que se hubiese determinado) y los que adquiera hasta la extinción de la obligación. El acreedor puede exigir la venta judicial solo en la medida necesaria para satisfacer su crédito, y todos los acreedores ejecutan en posición igualitaria salvo causa legal de preferencia.}

\cornellpregunta{¿Qué tipos de acreedores existen y cómo afecta eso al cobro?}{%
Con privilegio especial (sobre determinados bienes, por ejemplo el laboral), con privilegio general (sobre todo el patrimonio), con garantía real (hipoteca o prenda, elegida voluntariamente) y quirografarios (sin preferencia, cobran en igualdad sobre el remanente). En caso de insolvencia, cobran en ese orden.
% TODO: contenido incompleto o ambiguo en las notas originales (el orden de cobro «en ese orden» figura solo en el cuadro Cornell del apunte; el desarrollo no precisa la prelación entre las distintas preferencias)
}

\cornellpregunta{¿Qué bienes son inembargables? (Art. 744 CCCN)}{%
Ropa y muebles indispensables del hogar, instrumentos del oficio o profesión, sepulcros (salvo deuda de construcción), bienes de culto, ciertos derechos de usufructo, uso y habitación, indemnizaciones por lesiones o daño moral, alimentos por homicidio, y los declarados así por leyes especiales (salario mínimo). El fundamento común es proteger la dignidad humana en sus dimensiones más fundamentales.}

\cornellpregunta{¿Cómo se protege la vivienda frente a las deudas?}{%
Inscribiendo el inmueble en el Registro de la Propiedad Inmueble (Arts. 244 y ss. CCCN). Desde la inscripción, las deudas posteriores no pueden ejecutarlo. Solo se protege un inmueble, en su totalidad o hasta una parte de su valor. Las deudas anteriores a la inscripción no se ven afectadas.}

\cornellpregunta{¿Qué medidas preventivas tienen los acreedores antes de la ejecución?}{%
\textbf{Embargo preventivo:} afecta un bien determinado para que no salga del patrimonio del deudor. \textbf{Inhibición general de bienes:} impide al deudor enajenar cualquier bien cuando no se conocen bienes concretos. Sirven para contrarrestar la estrategia del deudor de convertir sus bienes en dinero en efectivo para que los acreedores no encuentren bienes a su nombre.}

\cornellpregunta{¿Qué es la acción directa?}{%
Permite al acreedor cobrar directamente lo que un tercero le debe a su deudor (Art. 736 CCCN). La ejerce en nombre propio, solo lo beneficia a él y opera hasta el límite de su crédito. Requiere homogeneidad de créditos y un caso legal expreso.}

\cornellpregunta{¿Qué es la acción subrogatoria?}{%
El acreedor ocupa el lugar de su deudor inactivo y ejerce los derechos patrimoniales de este en su nombre (Art. 739 CCCN). El bien recuperado ingresa al patrimonio del deudor y beneficia a todos los acreedores, no solo al subrogante.}

\cornellpregunta{¿Qué diferencias hay entre la acción directa y la acción subrogatoria?}{%
\begin{itemize}[nosep,leftmargin=1.2em]
\item \textbf{Directa:} se actúa en nombre propio; beneficia solo al acreedor que la ejerce; con privilegio (cobra directamente hasta su crédito); requiere créditos homogéneos y caso legal expreso.
\item \textbf{Subrogatoria:} se actúa en nombre del deudor; beneficia a todos los acreedores; sin privilegio (el bien ingresa al patrimonio y todos pueden ejecutarlo); requiere inactividad del deudor que perjudique la acreencia.
\end{itemize}}

\cornellpregunta{¿Qué derechos no pueden ejercerse mediante la acción subrogatoria?}{%
Los derechos vinculados con la personalidad y los de ejercicio absolutamente discrecional por el titular (Arts. 741--742 CCCN); por ejemplo, el reconocimiento de un hijo o cuestiones de puro ejercicio de la voluntad familiar.}

\cornellpregunta{¿Qué diferencia hay entre la acción de simulación y la de inoponibilidad?}{%
En la simulación el acto es ficticio (el deudor aparenta vender pero conserva el bien); el efecto es la nulidad. En la inoponibilidad (acción revocatoria o pauliana) el acto es real pero fraudulento: el acreedor lo hace inoponible a sí mismo y puede ejecutar el bien aunque el tercero lo haya adquirido realmente.}

\cornellresumen{Los bienes del deudor responden por sus obligaciones y los acreedores los ejecutan en posición igualitaria salvo causa legal de preferencia (Art. 743 CCCN). Los acreedores pueden tener privilegio especial, privilegio general, garantía real o ser quirografarios. La garantía tiene límites: los bienes inembargables (Art. 744) y la protección de la vivienda (Art. 244). Para conservar y recuperar la garantía, los acreedores cuentan con medidas preventivas (embargo preventivo, inhibición general de bienes) y con las acciones directa, subrogatoria, de simulación y de inoponibilidad.}

\end{longtable}
\cornellfin

% =====================================================================
% CAPÍTULO 5 — SÍNTESIS GENERAL DE LA UNIDAD
% =====================================================================

\section{Síntesis de la parte}
\subsection{Recorrido de la parte}

Esta parte recorrió el patrimonio, el patrimonio como garantía, los distintos tipos de acreedores y los distintos tipos de acciones de los acreedores. Desde el patrimonio como garantía existe la posibilidad de ejercer acciones; asimismo, hay bienes que quedan excluidos de ella, y que el Código enuncia en los artículos 744 y 244.

\subsection{Síntesis final}

Esta parte construyó la arquitectura conceptual sobre la que descansa buena parte del Derecho Privado argentino. El patrimonio es el puente entre la persona y el mercado: como atributo de la personalidad, le da contenido económico al sujeto; como garantía común (Art. 242 CCCN), le da sustento a todo el sistema crediticio.

Para entender de qué está hecho ese patrimonio, el Código clasifica los bienes y cosas según múltiples criterios ---muebles e inmuebles, fungibles y no fungibles, divisibles e indivisibles, principales y accesorias, frutos y productos---, cada uno con consecuencias jurídicas precisas que atraviesan el Derecho Real, Obligacional, Contractual, de Familia y de Sucesiones.

El sistema de protección articula dos tensiones: la del acreedor, que necesita ejecutar el patrimonio para cobrar, y la del deudor, que no puede ser reducido a la indigencia. Esta tensión se resuelve mediante los bienes inembargables y la protección de la vivienda.

Finalmente, las acciones directa, subrogatoria, de simulación y de inoponibilidad son los instrumentos procesales que permiten al acreedor conservar y recuperar la garantía frente a estrategias de vaciamiento patrimonial.
\partintro{Si el sujeto es quien entabla la relación y el objeto es la materia sobre la que recae, la causa explica cómo surgen las relaciones jurídicas. Esta parte estudia los hechos jurídicos y el acto voluntario, la imputabilidad, el acto jurídico y sus elementos (objeto, causa y forma), los instrumentos, la clasificación y las modalidades de los actos y, por último, los efectos de los actos jurídicos respecto de las partes, los representantes, los sucesores y los terceros.}
\part{La causa de la relación jurídica: hechos y actos jurídicos}
\chapter{Fundamentos de los Hechos y Actos Jurídicos}

\begin{note}
Los primeros capítulos de esta parte abordan el tercer elemento de la relación jurídica: la causa, es decir, el origen del cual emergen las relaciones jurídicas. El recorrido conceptual parte de la clasificación general de los hechos jurídicos, avanza hacia la distinción entre hechos voluntarios e involuntarios (discernimiento, intención y libertad) y culmina con el estudio de la imputabilidad, es decir, de las consecuencias que se atribuyen a quien realiza un acto voluntario, así como del tratamiento excepcional que reciben los actos involuntarios a través del principio de equidad.
\end{note}

\section{Introducción: la causa como tercer elemento de la relación jurídica}

La teoría de los hechos y actos jurídicos constituye el tercer elemento de la relación jurídica, luego de haberse estudiado el sujeto y el objeto. Mientras el sujeto es quien entabla la relación y el objeto es la materia sobre la que ésta recae, la causa explica cómo surgen las relaciones jurídicas.

\begin{table}[H]
\centering
\caption{Elementos de la relación jurídica}
\begin{tabularx}{\textwidth}{l l X}
\toprule
\textbf{Elemento} & \textbf{Definición} & \textbf{Contenido} \\
\midrule
Sujeto & Quien entabla la relación & Persona humana o jurídica \\
Objeto & Materia sobre la que recae la relación & Bienes y cosas \\
Causa & De dónde surge la relación jurídica & \textbf{Hechos y actos jurídicos} \\
\bottomrule
\end{tabularx}
\end{table}

Esta teoría proyecta efectos sobre el conjunto del derecho civil y comercial, en la medida en que toda disciplina jurídica se ocupa, en definitiva, de hechos que dan origen a relaciones y situaciones jurídicas.

\section{Ubicación metodológica en el Código Civil y Comercial}

En el Código de Vélez, los hechos y actos jurídicos se regulaban en el Libro Segundo, Sección Segunda, dedicada a los derechos personales. El artículo 896 marcaba el inicio de la regulación de los hechos jurídicos, distanciada de los derechos reales y de otras áreas del código.

\begin{table}[H]
\centering
\caption{Cambio metodológico: Código de Vélez frente al CCCN}
\begin{tabularx}{\textwidth}{X X}
\toprule
\textbf{Código de Vélez (anterior)} & \textbf{CCCN actual} \\
\midrule
Libro Segundo -- Derechos personales & Libro Primero -- Parte General (metodológicamente más correcto) \\
Art. 896: iniciaba los hechos jurídicos & Desarrolla la teoría en nueve capítulos integrados a la parte general \\
Distanciado de los derechos reales y otras áreas & Integra la Parte General y fue bien recibido por la doctrina \\
\bottomrule
\end{tabularx}
\end{table}

El Código Civil y Comercial de la Nación (CCCN) desarrolla la teoría en nueve capítulos: disposiciones generales, los vicios (error, dolo, violencia), los actos jurídicos propiamente dichos, los vicios de los actos, las modalidades, la representación y la ineficacia. Este cambio metodológico constituye un acierto del nuevo Código, en la medida en que integra la parte general del derecho civil. Concretamente, la regulación de los hechos y actos jurídicos se ubica en el Libro Primero, Título Cuarto del Código.

\section{Definición de hecho jurídico en el CCCN}

\begin{definition}{Art. 257 CCCN -- Hecho jurídico}
El hecho jurídico es el acontecimiento que, conforme al ordenamiento jurídico, produce el nacimiento, modificación y extinción de relaciones o situaciones jurídicas.
\end{definition}

\begin{table}[H]
\centering
\caption{Comparación entre el art. 257 CCCN y el art. 896 del Código de Vélez}
\begin{tabularx}{\textwidth}{l X X}
\toprule
\textbf{Aspecto} & \textbf{CCCN (Art. 257)} & \textbf{Código de Vélez (Art. 896)} \\
\midrule
Término clave & Produce (certeza) & Susceptible de producir (posibilidad) \\
Referencia al orden jurídico & Sí: ``conforme al ordenamiento jurídico'' & No hacía esta referencia expresa \\
Efectos abarcados & Relaciones o situaciones jurídicas & Derechos u obligaciones solamente \\
\bottomrule
\end{tabularx}
\end{table}

El CCCN incorpora la expresión ``conforme al ordenamiento jurídico'', señalando que es el ordenamiento jurídico el que genera estas consecuencias. Esta afirmación no debe entenderse exclusivamente desde una visión positivista, como si únicamente el ordenamiento jurídico las generase, pues existe al respecto un debate doctrinario abierto.

El reemplazo del término ``susceptible de producir'' por ``produce'' resulta más preciso, en tanto al derecho le interesan los acontecimientos que efectivamente producen consecuencias jurídicas.

\begin{example}{El hecho jurídico frente al acontecimiento cotidiano}
Tomar un vaso de agua no constituye un hecho jurídico. Arrojarlo por la ventana y dañar a un transeúnte sí lo es, en tanto genera una consecuencia jurídica relevante para el ordenamiento.
\end{example}

\section{Clasificación general de los hechos jurídicos}
\label{sec:clasifhechos}

El cuadro general de clasificación de los hechos jurídicos resulta central para el desarrollo de la materia.

\begin{table}[H]
\centering
\caption{Hechos jurídicos: clasificación general}
\begin{tabularx}{\textwidth}{X X}
\toprule
\textbf{Hechos de la naturaleza} & \textbf{Hechos humanos} \\
\midrule
Granizo, terremoto, aluvión. Relevan por sus efectos, no por su origen. &
\textbf{Involuntarios}: sin discernimiento, intención o libertad. \\
\cmidrule(l){2-2}
 & \textbf{Voluntarios} (con discernimiento, intención y libertad): \\
 & \quad -- \textbf{Lícitos}: simples actos lícitos / actos jurídicos \\
 & \quad -- \textbf{Ilícitos}: delitos (con dolo) / cuasidelitos (con culpa) \\
\bottomrule
\end{tabularx}
\end{table}

\section{Hechos voluntarios e involuntarios: importancia de la distinción}

\begin{example}{Caso del ataque epiléptico en la estación de subte}
Una persona sufre un ataque de epilepsia en una estación de subte y se desvanece, golpeando a otra persona, que a su vez golpea a una tercera que casi cae a las vías del tren.

Conclusión jurídica: el hecho se origina de manera involuntaria. No hubo voluntad de empujar; la persona sufrió un ataque de epilepsia y se desvaneció. Desde el punto de vista jurídico, este hecho es tratado de manera diferente al supuesto en que existe voluntad, pues quienes tienen dominio de sus actos se hacen cargo de ellos, existiendo mayor responsabilidad cuando media voluntad.
\end{example}

\section{Actos ilícitos: dolo y culpa}

Los actos ilícitos son los prohibidos por la ley que generan daños. A diferencia del derecho penal, donde los delitos se encuentran tipificados, en el derecho civil los ilícitos comprenden todos los actos que causan un daño, presumiéndose que dicho daño es siempre antijurídico.

\begin{table}[H]
\centering
\caption{Delito civil y cuasidelito civil}
\begin{tabularx}{\textwidth}{l X}
\toprule
\textbf{Figura} & \textbf{Definición} \\
\midrule
Delito civil & Acto ilícito cometido con dolo (intención de dañar o manifiesta indiferencia por los intereses ajenos) \\
Cuasidelito civil & Acto ilícito cometido con culpa (negligencia, imprudencia o impericia, sin intención de dañar) \\
\bottomrule
\end{tabularx}
\end{table}

\begin{note}
Los términos ``delito civil'' y ``cuasidelito civil'' no aparecen literalmente en el texto del CCCN. Se trata de una clasificación histórica elaborada por la doctrina.

Cabe señalar, asimismo, una crítica metodológica: los actos ilícitos se encuentran regulados en el Libro Tercero del CCCN (Derechos Personales -- Responsabilidad Civil), cuando, al menos en sus elementos centrales, deberían haberse tratado en el Libro Primero, junto a la teoría general de los hechos jurídicos.
\end{note}

\begin{definition}{Art. 1724 CCCN -- Factores subjetivos de atribución (dolo y culpa)}
El dolo se configura por la producción de un daño de manera intencional o con manifiesta indiferencia por los intereses ajenos.

La culpa consiste en la omisión de la diligencia debida según la naturaleza de la obligación y las circunstancias de las personas, el tiempo y el lugar. Comprende la imprudencia, la negligencia y la impericia en el arte o profesión.
\end{definition}

El CCCN amplía la noción clásica romanista del dolo (intención de dañar), incorporando también la ``manifiesta indiferencia''.

\begin{example}{Caso del conductor que ignora a la anciana}
Un conductor circula a 120 km/h y advierte que una anciana cruza correctamente con el semáforo en verde para peatones. Pese a verla, no realiza ningún intento de esquivarla ni de frenar, y la mata.

Análisis: no puede determinarse con certeza que el conductor haya querido matarla con premeditación, pero sí mostró una ``manifiesta indiferencia'' ante los intereses ajenos. Por lo tanto, se configura el dolo por indiferencia, que el CCCN incorpora expresamente.
\end{example}

\begin{table}[H]
\centering
\caption{Formas de culpa}
\begin{tabularx}{\textwidth}{l X}
\toprule
\textbf{Forma de culpa} & \textbf{Definición} \\
\midrule
Negligencia & Torpeza, falta de cuidado o diligencia debida. Ej.: un médico que no se lava las manos antes de una intervención. \\
Imprudencia & Realizar algo osado que no correspondía, asumiendo un riesgo innecesario sin intención de dañar. \\
Impericia & Falta de conocimientos técnicos propios de la profesión u oficio. Ej.: un médico que interviene sin los conocimientos necesarios. \\
\bottomrule
\end{tabularx}
\end{table}

\section{Actos lícitos: simples actos lícitos y actos jurídicos}

\begin{definition}{Art. 258 CCCN -- Simple acto lícito}
La acción voluntaria no prohibida por la ley, de la que resulta alguna adquisición, modificación o extinción de relaciones o situaciones jurídicas.
\end{definition}

\begin{definition}{Art. 259 CCCN -- Acto jurídico}
El acto jurídico es el acto voluntario lícito que tiene por fin inmediato la adquisición, modificación o extinción de relaciones o situaciones jurídicas.
\end{definition}

\begin{table}[H]
\centering
\caption{Simple acto lícito frente a acto jurídico}
\begin{tabularx}{\textwidth}{l X X}
\toprule
\textbf{Criterio} & \textbf{Simple acto lícito (Art. 258)} & \textbf{Acto jurídico (Art. 259)} \\
\midrule
Voluntad & Sí, voluntario & Sí, voluntario \\
Licitud & Sí, no prohibido por la ley & Sí, lícito \\
Fin inmediato & No busca producir consecuencias jurídicas directamente & Sí: fin inmediato de producir consecuencias jurídicas \\
Ejemplo & Pasear por la calle (puede derivar en consecuencias jurídicas imprevistas) & Contrato de mutuo (préstamo de dinero con plazo e interés pactados) \\
\bottomrule
\end{tabularx}
\end{table}

\begin{example}{El contrato de mutuo como acto jurídico}
En un contrato de mutuo, una persona decide prestarle dinero a otra, fijando el plazo y, eventualmente, una tasa de interés. Todo ello se establece en un contrato firmado por personas humanas que resuelven voluntariamente realizar un acto lícito --prestarse dinero-- con el fin inmediato de generar una relación jurídica.
\end{example}

% =====================================================================
% CAPÍTULO 2
% =====================================================================
\chapter{Teoría del Acto Voluntario}

\section{Las condiciones del acto voluntario}

\begin{definition}{Art. 260 CCCN -- Acto voluntario}
El acto voluntario es el ejecutado con discernimiento, intención y libertad, que se manifiesta por un hecho exterior.
\end{definition}

\begin{table}[H]
\centering
\caption{Condiciones del acto voluntario}
\begin{tabularx}{\textwidth}{l l X}
\toprule
\textbf{Condición} & \textbf{Tipo} & \textbf{Descripción} \\
\midrule
Discernimiento & Interna & Aptitud para diferenciar lo bueno de lo malo, lo verdadero de lo falso, lo justo de lo injusto y sus consecuencias \\
Intención & Interna & Propósito de llevar a cabo un acto determinado, de manera consciente y dirigida a ese fin \\
Libertad & Interna & Posibilidad de actuar de acuerdo a la propia conveniencia, deseo y convicción, sin coacción \\
Exteriorización & Externa & Manifestación visible del acto: expresa (oral, escrita, signos inequívocos) o tácita \\
\bottomrule
\end{tabularx}
\end{table}

\section{El discernimiento}

El discernimiento es la aptitud para diferenciar lo bueno de lo malo, lo verdadero de lo falso, lo justo de lo injusto y sus consecuencias. Se trata de una comprensión genérica, una capacidad o aptitud para comprender lo que está sucediendo.

\begin{table}[H]
\centering
\caption{Causas obstativas del discernimiento (Art. 261 CCCN)}
\begin{tabularx}{\textwidth}{l X}
\toprule
\textbf{Causa} & \textbf{Detalle} \\
\midrule
Privación de razón (permanente o transitoria) & Ej.: golpe en la cabeza que genera un estado de inconsciencia durante el cual se realizó un acto perjudicial. No requiere sentencia judicial de incapacidad. \\
Menor de diez (10) años -- para actos ilícitos & El código presume que antes de esa edad no hay discernimiento para comprender lo injusto y sus consecuencias. \\
Menor de trece (13) años -- para actos lícitos & Para los actos lícitos se exige mayor madurez cognitiva; la edad mínima es superior. \\
\bottomrule
\end{tabularx}
\end{table}

\begin{important}
El discernimiento no debe confundirse con la capacidad. La capacidad es un concepto jurídico integral que se refiere a la aptitud para entablar relaciones jurídicas, mientras que el discernimiento es un componente del acto voluntario. Puede haber discernimiento sin capacidad: es el caso de las personas menores de edad.
\end{important}

\begin{example}{Discernimiento sin capacidad jurídica}
Un niño de diez años puede tener discernimiento suficiente para un acto lícito, pero no necesariamente capacidad jurídica para celebrar ese acto.
\end{example}

\begin{table}[H]
\centering
\caption{Discernimiento frente a capacidad jurídica}
\begin{tabularx}{\textwidth}{X X}
\toprule
\textbf{Discernimiento (Art. 261 CCCN)} & \textbf{Capacidad jurídica} \\
\midrule
Elemento del acto voluntario. Se evalúa en cada caso concreto. No depende de sentencia judicial. & Concepto jurídico integral. Puede requerir sentencia judicial. Se adquiere con la mayoría de edad (18 años). \\
\bottomrule
\end{tabularx}
\end{table}

\section{La intención y sus vicios: error y dolo}

La intención es el propósito de llevar a cabo un acto determinado con consciencia. Mientras el discernimiento es la característica general que puede estar presente o no en cada acto, la intención es el propósito de realizar ese acto concreto, preciso y específico.

\begin{table}[H]
\centering
\caption{Vicios de la intención}
\begin{tabularx}{\textwidth}{l X X}
\toprule
\textbf{Vicio} & \textbf{Definición} & \textbf{Ejemplo} \\
\midrule
Error & Ignorancia o falso conocimiento de la realidad que vicia la voluntad & Creer que se vende una casa cuando la otra parte entiende que es una donación. Error sobre la naturaleza del acto. \\
Dolo (vicio de la voluntad) & Maniobra engañosa por la que se induce a la otra parte a celebrar el acto & Alquilar una propiedad mostrando fotos falsas; al llegar, la casa está en mal estado y lejos de la playa. \\
\bottomrule
\end{tabularx}
\end{table}

La palabra ``dolo'' posee al menos dos significados en el derecho civil: (1) como intención de dañar, elemento subjetivo de la responsabilidad civil, y (2) como maniobra engañosa que vicia la voluntad.

El error puede recaer sobre los objetos, las personas, las causas o las cualidades de la cosa, dando lugar a distintas circunstancias que se estudian en el capítulo~\ref{ch:error}.

Respecto del dolo como vicio, la cuestión central consiste en determinar cuándo el engaño tiene entidad suficiente para afectar la intención, es decir, cuándo puede afirmarse que el acto no se habría realizado de no mediar el engaño.

\section{La libertad y el vicio de violencia}

La libertad es la posibilidad de llevar a cabo el acto de acuerdo con la propia conveniencia, deseo y convicción.

\begin{table}[H]
\centering
\caption{Formas de violencia como causa obstativa de la libertad}
\begin{tabularx}{\textwidth}{l X}
\toprule
\textbf{Forma de violencia} & \textbf{Definición} \\
\midrule
Fuerza física irresistible & Uso de la fuerza física de manera irresistible sobre el sujeto para que realice el acto. Ej.: torturar a alguien para que firme. \\
Intimidación (violencia moral) & Amenaza de sufrir un mal grave e inminente que no puede contrarrestarse, y que lleva a la persona a realizar el acto sin libertad real. \\
\bottomrule
\end{tabularx}
\end{table}

Es posible comprender, querer y realizar un acto con esa intención, y sin embargo no ser suficientemente libre por encontrarse la persona bajo amenaza. Un acto realizado bajo amenaza no es, por tanto, un acto voluntario.

\section{La exteriorización de la voluntad}

La voluntad debe exteriorizarse para ser considerada tal. Esto puede ocurrir de manera expresa o tácita.

\begin{table}[H]
\centering
\caption{Modalidades de exteriorización de la voluntad}
\begin{tabularx}{\textwidth}{l X X}
\toprule
\textbf{Modalidad} & \textbf{Descripción} & \textbf{Ejemplo} \\
\midrule
Expresa oral & La voluntad se comunica verbalmente & Pactar verbalmente un préstamo \\
Expresa escrita & Instrumentada en documento público o privado & Contrato firmado; escritura pública \\
Signos inequívocos & Gestos que no admiten otra interpretación & Hacer el gesto de pedir un café en un bar \\
Ejecución de hecho material & Realización de un acto que revela la voluntad & Girar dinero como pago de una deuda \\
Tácita & Se infiere de conductas que permiten conocer con certidumbre la voluntad & Serie de actos que evidencian aceptación de una oferta \\
\bottomrule
\end{tabularx}
\end{table}

\begin{definition}{Art. 264 CCCN -- Límite a la voluntad tácita}
Si la ley o la convención exigen que la voluntad sea expresa, no puede suplirse con una expresión tácita.
\end{definition}

Un ejemplo de esta regla es el consentimiento exigido para los actos de disposición de derechos, respecto de los cuales la ley exige que sea expreso, no pudiendo interpretarse un consentimiento tácito en ese contexto.

\section{El silencio como (no) manifestación de voluntad}

\begin{definition}{Art. 263 CCCN -- El silencio como manifestación de voluntad}
El silencio opuesto a actos o a una interrogación no es considerado como una manifestación de voluntad conforme al acto o la interrogación, excepto en los casos en que haya un deber de expedirse que puede resultar de la ley o de la voluntad de las partes.
\end{definition}

En el derecho civil y comercial, el silencio no constituye, en principio, una manifestación de voluntad: quien calla no otorga.

\begin{example}{El silencio frente a una oferta no solicitada}
Si una persona recibe un mensaje de texto que indica ``te has ganado un auto; a partir de mañana se descuentan \$20.000 de tu tarjeta de crédito durante 60 meses, salvo que contestes que no lo querés en 24 horas'', y guarda silencio, ese silencio no puede considerarse aceptación. De lo contrario, sería necesario responder permanentemente a las publicidades para aclarar que no se las desea.
\end{example}

\begin{table}[H]
\centering
\caption{Excepciones: cuándo el silencio sí es manifestación de voluntad}
\begin{tabularx}{\textwidth}{l X}
\toprule
\textbf{Excepción} & \textbf{Descripción} \\
\midrule
Deber legal de expedirse & La ley obliga a pronunciarse. Ej.: intimación para reconocer o desconocer la firma de un instrumento. Si se guarda silencio, se tiene por reconocida. \\
Deber convencional de expedirse & Las partes acordaron que el silencio implica determinada manifestación. Ej.: contrato que establece que se renueva automáticamente si la parte no comunica lo contrario con un mes de anticipación. \\
Relaciones entre declaraciones precedentes & Cuando una serie de declaraciones anteriores llevaron razonablemente a la otra parte a esperar una determinada respuesta, y el silencio posterior confirma esa dirección. \\
\bottomrule
\end{tabularx}
\end{table}

% =====================================================================
% CAPÍTULO 3
% =====================================================================
\chapter{Imputabilidad de los Actos Voluntarios}

\section{El deber de reparar: fundamento y presupuestos}

Quienes realizan un acto voluntario deben hacerse cargo de sus consecuencias, en virtud de la relación de causalidad existente entre la acción y el resultado. La cuestión central consiste en determinar hasta dónde se extiende esa responsabilidad.

\begin{definition}{Art. 1716 CCCN -- Deber de reparar}
La violación del deber de no dañar a otro, o el incumplimiento de una obligación, da lugar a la reparación del daño causado, conforme con las disposiciones de este Código.
\end{definition}

\begin{definition}{Art. 1717 CCCN -- Antijuridicidad}
Cualquier acción u omisión que causa un daño a otro es antijurídica si no está justificada.
\end{definition}

\begin{table}[H]
\centering
\caption{Presupuestos de la responsabilidad civil}
\begin{tabularx}{\textwidth}{c l X}
\toprule
\textbf{N.\textdegree{}} & \textbf{Presupuesto} & \textbf{Descripción} \\
\midrule
1 & Daño & Debe existir un daño efectivo a una persona o a su patrimonio. \\
2 & Antijuridicidad & El daño debe ser antijurídico (se presume siempre, salvo causa de justificación -- Art. 1718). \\
3 & Factor de atribución & Subjetivo: dolo o culpa (Art. 1724). Objetivo: la ley establece que se responde independientemente de dolo o culpa (ej.: art. 1757 -- responsabilidad por cosas y actividades riesgosas). \\
4 & Relación de causalidad & El acto voluntario debe ser la causa adecuada del daño producido (nexo causal adecuado). \\
\bottomrule
\end{tabularx}
\end{table}

\begin{example}{Concausalidad: auto que choca a motociclista sin casco}
Un auto choca a una moto cuyo conductor no lleva casco. El accidente le provoca una lesión cerebral severa, agravada precisamente por la falta de casco.

Análisis: quien chocó a la moto tiene una parte fundamental en la causación del daño, pues sin el choque el motociclista no se habría caído. Pero la negligencia del conductor que no llevaba casco también incide en el resultado dañoso, ya que con casco se habrían evitado o reducido las consecuencias. Esto no exime al conductor del auto de responder, pero obliga a analizar cómo incidió cada causa (concausalidad).
\end{example}

\section{Tipos de consecuencias: inmediatas, mediatas y casuales}

\begin{table}[H]
\centering
\caption{Clasificación de las consecuencias del acto}
\begin{tabularx}{\textwidth}{l X X}
\toprule
\textbf{Tipo} & \textbf{Definición} & \textbf{Ejemplo (choque de auto)} \\
\midrule
Inmediatas & Las que acostumbran suceder según el curso normal y ordinario de las cosas (Art. 1726 CCCN). & Auto choca a persona $\rightarrow$ fractura de cadera. Resultado directo del impacto. \\
Mediatas & Resultan de la conexión del hecho con un acontecimiento distinto (Art. 1726 CCCN). & Internación $\rightarrow$ pérdida de ingresos como carpintero durante 20 días. \\
Casuales & No pueden ser previstas, o resultan de circunstancias extraordinarias (Art. 1727 CCCN). & Pérdida de ingresos $\rightarrow$ depresión $\rightarrow$ divorcio $\rightarrow$ hijos con adicciones. Imprevisibles. \\
\bottomrule
\end{tabularx}
\end{table}

\begin{formula}{Arts. 1726 y 1727 CCCN -- ¿Por cuáles consecuencias se responde?}
Se reparan las consecuencias inmediatas y las mediatas previsibles. No se responde por las consecuencias casuales (mediatas no previsibles).
\end{formula}

\section{El criterio de previsibilidad y la graduación de la responsabilidad}

\begin{definition}{Art. 1725 CCCN -- Valoración de la conducta}
Cuanto mayor sea el deber de obrar con prudencia y pleno conocimiento de las cosas, mayor es la diligencia exigible al agente y la valoración de la previsibilidad de las consecuencias.
\end{definition}

Cuanto mayor es lo que una persona conoce, y cuanto mayor es su deber de conocer y obrar con prudencia, mayor es la diligencia que se le exige y, en consecuencia, mayor es también la valoración de la previsibilidad de las consecuencias por las que deberá responder.

\begin{table}[H]
\centering
\caption{Graduación de la diligencia exigible según la situación}
\begin{tabularx}{\textwidth}{X X}
\toprule
\textbf{Situación} & \textbf{Exigencia de diligencia} \\
\midrule
Profesional de la salud tratando un paciente & Alta: se le exige la diligencia propia de su especialidad y el estándar de cuidado de la comunidad médica. \\
Trabajador en planta con material nuclear o contaminante & Muy alta: el riesgo elevado implica mayor previsibilidad exigida. \\
Trabajador en tareas de bajo riesgo & Estándar ordinario de diligencia. \\
\bottomrule
\end{tabularx}
\end{table}

La condición especial o la facultad intelectual particular del agente no se toma en cuenta a los fines de esta graduación, a menos que haya sido específicamente contemplada en un contrato.

\section{El caso fortuito y la fuerza mayor como eximentes}

\begin{definition}{Art. 1730 CCCN -- Caso fortuito y fuerza mayor}
Se considera caso fortuito o fuerza mayor al hecho que no ha podido ser previsto, o que, habiendo sido previsto, no ha podido ser evitado. El caso fortuito o fuerza mayor exime de responsabilidad, excepto disposición en contrario.
\end{definition}

\begin{important}
Cuando algo no puede ser previsto, el supuesto excede el ámbito de lo voluntario y se enmarca dentro del caso fortuito o la fuerza mayor. Se trata de un principio clásico del derecho civil: cuando algo no es voluntario, no se responde.
\end{important}

% =====================================================================
% CAPÍTULO 4
% =====================================================================
\chapter{Responsabilidad por Actos Involuntarios}

\section{Regla general: no responsabilidad por actos involuntarios}

Los actos involuntarios, a diferencia de los voluntarios, no generan responsabilidad en principio. Ello obedece a que, si el agente no obró con discernimiento, intención y libertad, no puede imputársele el deber de responder por las consecuencias del hecho.

\begin{important}
En principio, los actos involuntarios no generan responsabilidad, salvo las excepciones expresamente previstas por el ordenamiento.
\end{important}

El antiguo artículo 907 del Código de Vélez establecía que en ningún caso se respondía por los actos involuntarios.

\section{La evolución: el principio de equidad}

A partir, especialmente, de la jurisprudencia francesa, comenzó a abrirse camino la idea de que la regla de no responsabilidad por actos involuntarios podía resultar injusta en determinados casos.

\begin{example}{El caso que motivó la incorporación de la equidad}
Una persona privada de la razón por una enfermedad mental, con un gran patrimonio, causa un daño enorme a una persona muy pobre con familia numerosa.

Dilema: desde el punto de vista teórico, el hecho es involuntario y no puede imputarse. Pero cabe preguntarse si es justo que una persona muy rica, cuya conducta involuntaria arruina a una familia pobre, no repare absolutamente nada. Aquí aparece el concepto de equidad.
\end{example}

\begin{table}[H]
\centering
\caption{Dos concepciones de equidad}
\begin{tabularx}{\textwidth}{l X}
\toprule
\textbf{Concepción} & \textbf{Contenido} \\
\midrule
Equidad romana / igualitaria & Buscar que las personas tengan lo indispensable para la vida; igualdad material o distributiva. \\
Epieikeia aristotélica (equidad correctiva) & Rectificación de lo justo legal en el caso concreto. Cuando la aplicación estricta de la ley conduce a un resultado manifiestamente injusto, el juez puede apartarse de ella y disponer algo distinto para ese caso. Esta es la concepción que opera en el artículo 1750 CCCN. \\
\bottomrule
\end{tabularx}
\end{table}

\section{El artículo 1750 del CCCN: régimen vigente}

\begin{definition}{Art. 1750 CCCN -- Daños causados por actos involuntarios}
El autor de un daño causado por un acto involuntario responde por razones de equidad. Se aplica lo dispuesto en el artículo 1742. El acto realizado por quien sufre fuerza irresistible no genera responsabilidad para su autor, sin perjuicio de la que corresponde a título personal a quien ejerce esa fuerza.
\end{definition}

\begin{table}[H]
\centering
\caption{Evolución normativa: del art. 907 (mod. 1968) al art. 1750 CCCN}
\begin{tabularx}{\textwidth}{l X X}
\toprule
\textbf{Aspecto} & \textbf{Anterior (Art. 907 mod. 1968)} & \textbf{CCCN (Art. 1750)} \\
\midrule
Regla general & Nunca se respondía & El autor responde por razones de equidad \\
Facultad judicial & ``Los jueces podrán disponer'' & Texto más imperativo; el autor ``responde'' \\
Quantum & Patrimonio del autor y situación de la víctima & Se remite al art. 1742 (atenuación) \\
Debate doctrinario & Incorporado en 1968, debatido desde entonces & Jornadas Nacionales de Derecho Civil (Santa Fe): ¿es posible una indemnización de cero? \\
\bottomrule
\end{tabularx}
\end{table}

\begin{note}
Constituye una cuestión doctrinaria debatida si el juez podría considerar que la equidad implica no establecer ninguna indemnización, frente a la posición de otros autores que sostienen que siempre debe establecerse alguna. La mayoría de la doctrina se inclina por esta última posición, aunque se trata de un debate abierto en el que se ha señalado que podría el juez, por razones de equidad, no fijar indemnización alguna.

Se menciona al Dr. Juan Carlos Palmero, jurista cordobés, como el autor que más ha profundizado este tema en Argentina, con una publicación específica sobre el daño involuntario.
\end{note}

\section{Criterios de la equidad para la atenuación (Art. 1742 CCCN)}

\begin{definition}{Art. 1742 CCCN -- Atenuación de la responsabilidad}
El juez, al fijar la indemnización, puede atenuarla si es equitativo en función del patrimonio del deudor, la situación personal de la víctima y las circunstancias del hecho. Esta facultad no es aplicable en caso de dolo del responsable.
\end{definition}

\begin{table}[H]
\centering
\caption{Factores para la determinación de la indemnización equitativa}
\begin{tabularx}{\textwidth}{l X}
\toprule
\textbf{Factor} & \textbf{Descripción} \\
\midrule
Patrimonio del causante del daño & ¿Cuánta capacidad económica tiene quien causó el daño involuntario? \\
Situación personal de la víctima & ¿Cuáles son las circunstancias de la persona que sufrió el daño? \\
Circunstancias del hecho & El contexto particular del caso concreto que rodea al evento dañoso. \\
\bottomrule
\end{tabularx}
\end{table}

El juez debe considerar estos factores, pudiendo apartarse de la regla general por razones de equidad y disponer el monto de la indemnización, llegando incluso a la posibilidad de que no exista indemnización, según una posición doctrinaria minoritaria; la mayoría de los autores sostiene, en cambio, que siempre corresponde otorgar algún tipo de indemnización.

Con el desarrollo de la imputabilidad de los actos voluntarios y del régimen de los actos involuntarios queda cerrada la presentación general de la teoría de los hechos y actos jurídicos.

% =====================================================================
% CORNELL
% =====================================================================

\section{Cuadro Cornell de Repaso General}
\textit{Hechos y Actos Jurídicos · Actos Voluntarios · Imputabilidad}

\begin{longtable}{
    >{\raggedright\arraybackslash}p{0.30\textwidth}
    >{\raggedright\arraybackslash}p{0.60\textwidth}
}
\toprule
\textbf{Preguntas / palabras clave} &
\textbf{Notas / respuestas} \\
\midrule
\endfirsthead
\multicolumn{2}{l}{\textit{Cuadro Cornell de Repaso General (continuación)}} \\
\toprule
\textbf{Preguntas / palabras clave} &
\textbf{Notas / respuestas} \\
\midrule
\endhead
\midrule
\multicolumn{2}{r}{\textit{Continúa en la página siguiente}} \\
\endfoot
\bottomrule
\endlastfoot

\textbf{¿Qué es un hecho jurídico según el CCCN?} &
El hecho jurídico es el acontecimiento que, conforme al ordenamiento jurídico, produce el nacimiento, modificación y extinción de relaciones o situaciones jurídicas (Art. 257 CCCN). El CCCN mejoró la redacción anterior: usa ``produce'' en lugar de ``susceptible de producir'', y habla de ``relaciones o situaciones jurídicas'' en lugar de solo ``derechos u obligaciones''. \\

\textbf{¿Cómo se clasifican los hechos jurídicos?} &
Hechos de la naturaleza (sin intervención humana: granizo, terremoto, aluvión) y hechos humanos (con intervención de la voluntad). Los humanos se subclasifican en voluntarios (con discernimiento, intención y libertad) e involuntarios. \\

\textbf{¿Cuándo un acto humano es voluntario?} &
Cuando reúne tres condiciones internas --discernimiento (aptitud para distinguir lo bueno de lo malo), intención (propósito dirigido al acto concreto) y libertad (ausencia de coacción)-- y se exterioriza (Art. 260 CCCN). \\

\textbf{¿Qué son las causas obstativas y cuáles corresponden a cada condición?} &
Son los vicios o impedimentos que eliminan cada condición. Discernimiento: privación de razón, minoría de edad (menor de 10 años para actos ilícitos; menor de 13 para actos lícitos). Intención: error y dolo. Libertad: violencia (fuerza o intimidación). \\

\textbf{¿En qué se diferencia el acto jurídico del simple acto lícito?} &
Ambos son voluntarios y lícitos. La diferencia radica en el fin inmediato: el acto jurídico tiene por fin inmediato producir consecuencias jurídicas (Art. 259 CCCN); el simple acto lícito no busca ese resultado directamente, aunque puede producirlo (Art. 258 CCCN). \\

\textbf{¿Cuáles son los dos significados del ``dolo'' en el Derecho Civil?} &
(1) Como factor de atribución: intención de dañar o manifiesta indiferencia por los intereses ajenos (Art. 1724 CCCN). (2) Como vicio de la voluntad: maniobra engañosa que induce a la otra parte a celebrar un acto que de otro modo no habría realizado. \\

\textbf{¿Qué diferencia al dolo de la culpa como factores de atribución?} &
Dolo: el agente produce el daño de manera intencional o con manifiesta indiferencia. Culpa: omisión de la diligencia debida, sin intención de dañar. Se subdivide en negligencia (torpeza), imprudencia (temeridad) e impericia (falta de conocimiento técnico). \\

\textbf{¿Cuándo el silencio es manifestación de voluntad?} &
En principio no lo es (Art. 263 CCCN). Excepcionalmente sí, cuando: (a) la ley obliga a expedirse; (b) las partes pactaron que el silencio implica determinada consecuencia; o (c) existe una serie de declaraciones precedentes que permiten inferir el sentido del silencio. \\

\textbf{¿Cuáles son los cuatro presupuestos de la responsabilidad civil?} &
1) Daño (efectivo). 2) Antijuridicidad (se presume, salvo causa de justificación -- Art. 1717). 3) Factor de atribución (subjetivo: dolo/culpa; objetivo: la ley). 4) Relación de causalidad adecuada entre el acto y el daño. \\

\textbf{¿Por cuáles consecuencias responde el agente?} &
Por las consecuencias inmediatas (curso natural y ordinario -- Art. 1726) y las mediatas previsibles. No se responde por las casuales (mediatas imprevisibles -- Art. 1727). El caso fortuito --imprevisto o inevitable-- exime de responsabilidad (Art. 1730). \\

\textbf{¿Qué pasa con los actos involuntarios y el deber de reparar?} &
En principio no generan responsabilidad. Pero el Art. 1750 CCCN establece que el autor de un daño involuntario responde por razones de equidad. El juez pondera el patrimonio del causante, la situación de la víctima y las circunstancias del hecho (remisión al Art. 1742). Quien actúa bajo fuerza irresistible no responde; sí responde quien ejerció esa fuerza. \\

\textbf{¿Qué es la equidad en este contexto?} &
En su sentido aristotélico (epieikeia): la rectificación de lo justo legal cuando la aplicación estricta de la norma conduce a un resultado manifiestamente injusto en el caso concreto. Permite al juez apartarse de la regla general de no responsabilidad por actos involuntarios para establecer (o no) una indemnización justa. \\

\midrule

\multicolumn{2}{p{0.94\textwidth}}{
\vspace{0.2em}
\textbf{Resumen:}
La teoría de los hechos y actos jurídicos constituye el tercer elemento de la relación jurídica, es decir, la fuente de donde emergen derechos y obligaciones. El punto de partida es la clasificación general: hechos de la naturaleza y hechos humanos; dentro de los humanos, voluntarios e involuntarios según concurran discernimiento, intención y libertad. Los actos voluntarios se dividen en lícitos --simples actos lícitos y actos jurídicos (Art. 259 CCCN), cuyo elemento diferenciador es el fin inmediato de producir consecuencias jurídicas-- e ilícitos --delitos (con dolo) y cuasidelitos (con culpa)--. La exteriorización de la voluntad puede ser expresa o tácita; el silencio, como regla general, no es manifestación de voluntad salvo excepciones legales o convencionales. La segunda parte abordó la imputabilidad: el deber de reparar (Arts. 1716-1717 CCCN) nace cuando confluyen cuatro presupuestos --daño, antijuridicidad, factor de atribución y nexo causal adecuado--. Las consecuencias inmediatas y mediatas previsibles son reparables; las casuales no. El caso fortuito exime de responsabilidad (Art. 1730). Finalmente, los actos involuntarios, si bien no generan responsabilidad en principio, pueden dar lugar a una indemnización por equidad en virtud del Art. 1750 CCCN, aplicando los criterios de atenuación del Art. 1742: patrimonio del causante, situación personal de la víctima y circunstancias del hecho. Esta tensión entre la regla de la no responsabilidad y la equidad como válvula correctora constituye uno de los debates doctrinarios más actuales del derecho civil argentino.
\vspace{0.2em}
} \\

\end{longtable}
\chapter[El acto jurídico]{El acto jurídico: concepto y elementos}
\label{ch:actojuridico}

El acto jurídico es, dentro de la clasificación general de los hechos jurídicos (sección~\ref{sec:clasifhechos}), un hecho humano voluntario lícito que tiene por fin inmediato producir la adquisición, modificación o extinción de derechos, relaciones o situaciones jurídicas. Es la categoría central de esta parte.

\section{Acto jurídico y contrato: alcance de la teoría general}

El Código Civil y Comercial argentino contiene un sistema general de los actos jurídicos, a diferencia del Código francés, que no trata los actos jurídicos en general sino que se limita al contrato. \textbf{El acto jurídico es un concepto más amplio que el contrato}; el contrato es una especie dentro del género acto jurídico.

\begin{definition}{Contrato}
El contrato es el acto jurídico en el cual hay acuerdo de voluntades para el surgimiento de obligaciones en el campo patrimonial.
\end{definition}

\begin{example}{Contratos como especie de acto jurídico}
\begin{itemize}
    \item \textbf{Contrato de compraventa}: acuerdo de voluntades para transferir la propiedad de una cosa contra el pago de un precio.
    \item \textbf{Contrato de locación}: acuerdo para transferir la tenencia de un bien —mueble o inmueble— contra la entrega de una renta mensual (canon).
\end{itemize}
\end{example}

No todos los actos jurídicos son contratos. Existen actos jurídicos que no responden a esta estructura de acuerdo patrimonial de voluntades:

\begin{itemize}
    \item \textbf{Testamento}: acto que tiene por finalidad disponer de los bienes para el supuesto de la propia muerte, asignándolos según la voluntad del testador. Produce efectos luego de la muerte, pero no constituye un contrato.
    \item \textbf{Matrimonio}: acto jurídico de carácter especial que crea relaciones jurídicas de naturaleza personal.
\end{itemize}

\section{Definición y características del acto jurídico}

\begin{definition}{Acto jurídico (Artículo 259, Código Civil y Comercial de la Nación)}
Acto jurídico es el acto voluntario lícito que tiene por fin inmediato la adquisición, modificación o extinción de derechos, relaciones o situaciones jurídicas.
\end{definition}

De esta definición se desprenden tres elementos fundamentales: que el acto sea \textbf{voluntario}, que sea \textbf{lícito}, y que tenga por \textbf{fin inmediato} producir consecuencias jurídicas.

\begin{important}
La idea de ``fin inmediato'' es crucial: significa que el acto jurídico busca de manera directa la producción de consecuencias jurídicas. Esto es lo que lo diferencia del simple acto lícito, que puede producir consecuencias jurídicas, pero sin que ello constituya su finalidad.
\end{important}

\section{Acto jurídico frente al simple acto lícito}

La distinción entre el acto jurídico y el simple acto lícito reside en la intención del sujeto respecto del efecto jurídico producido.

\begin{example}{El pescador}
Un pescador que sale a pescar, al apropiarse de lo pescado, realiza un simple acto lícito. Se produce un efecto jurídico —el pescador pasa a ser dueño de lo que pescó—, pero su intención original al pescar no es adquirir propiedad, sino pescar. El efecto jurídico es una consecuencia accesoria o secundaria del acto, no su fin inmediato.

Por el contrario, cuando se celebra un contrato de compraventa, el fin inmediato del sujeto es adquirir la propiedad de una cosa. El efecto jurídico es precisamente el objetivo buscado, la intención directa. Por ello, se trata de un acto jurídico.
\end{example}

\section{Los elementos del acto jurídico}

El acto jurídico se compone de los siguientes elementos esenciales:

\begin{enumerate}
    \item \textbf{Sujeto}: la persona que concurre al acto para su celebración.
    \item \textbf{Objeto}: la materia sobre la que recae el acto, es decir, los bienes o hechos que constituyen su contenido.
    \item \textbf{Forma}: la manera en que se exterioriza el acto.
    \item \textbf{Causa}: la finalidad perseguida por las partes al realizar el acto.
\end{enumerate}

% TODO: contenido incompleto o ambiguo en las notas originales — el tratamiento detallado del elemento "sujeto" y del elemento "forma" no fue desarrollado en esta clase; el apunte indica expresamente que la forma se abordará en una clase posterior, y el sujeto no vuelve a ser tratado en el material disponible.

El sujeto de los actos jurídicos se estudió en las Partes~II y III. Los capítulos siguientes desarrollan el \textbf{objeto} y la \textbf{causa} del acto jurídico; la \textbf{forma} se estudia en el capítulo~\ref{ch:forma}, y los elementos accidentales o modalidades (condición, plazo y cargo), que pueden estar o no presentes en un acto determinado, en el capítulo~\ref{ch:modalidades}.
\chapter{El objeto del acto jurídico}

\section{Fundamento normativo}

El artículo 279 del Código Civil y Comercial es central en materia de objeto de los actos jurídicos. Este artículo reconoce como antecedente el artículo 953 del Código Civil de Vélez, disposición que tuvo enorme influencia a lo largo de toda la vigencia del código anterior, en tanto marcaba el contenido moral del acto.

\begin{definition}{Objeto del acto jurídico (Artículo 279, Código Civil y Comercial de la Nación)}
El objeto del acto jurídico debe ser un hecho o un bien. Los hechos no pueden ser imposibles (ni jurídica ni prácticamente), prohibidos por la ley, o contrarios a la moral y a las buenas costumbres. Los bienes que son objeto del acto no pueden ser prohibidos, ni ser de carácter extrapatrimonial a menos que exista una autorización legal.
\end{definition}

\section{Contenido moral y límite de la autonomía de la voluntad}

Un acto jurídico debe ser conforme a la ley, a la moral y al orden público. No existe una libertad absoluta en la autonomía de la voluntad: esta se encuentra limitada y moldeada por razones de distinto tipo —la ley, el orden público, la moral y la dignidad de la persona.

\begin{note}
A lo largo de la historia, la disposición del artículo 953 del Código de Vélez originó numerosas limitaciones y aplicaciones jurisprudenciales, entre ellas:
\begin{itemize}
    \item Antes de la sanción de la Ley 17.711, dicha disposición dio lugar a la aplicación del vicio de \textbf{lesión}, que permitía anular un acto cuando existía una notable desproporción entre las prestaciones.
    \item Limitación en la aplicación de \textbf{cláusulas penales}, entendidas como la cláusula que se impone a quien incumple un contrato, obligándolo a pagar, por ejemplo, intereses punitorios.
    \item Limitación de los intereses cuando resultan \textbf{usurarios}: cobrar una tasa de interés excesiva o desproporcionada constituye usura, considerada inmoral. Aunque ello depende de la situación y el contexto (por ejemplo, en escenarios inflacionarios), cobrar tasas del 200 al 300 por ciento de interés es claramente inmoral.
\end{itemize}
\end{note}

Existe una tensión ideológica sobre este punto: una mirada más liberal sostiene que cada parte es libre de pactar las tasas de interés que desee. Sin embargo, en Argentina —y en la mayoría de las legislaciones modernas— aun cuando no exista un límite normativo expreso a los intereses, el objeto del acto jurídico debe ser conforme a la moral y a las buenas costumbres, principio que conserva plena vigencia en la actualidad.

\section{Requisitos del objeto}

El artículo 279 divide el objeto del acto jurídico en dos grandes categorías: los hechos y los bienes.

\subsection{Los hechos como objeto del acto}

Los hechos, en tanto acontecimientos, no pueden ser:

\begin{enumerate}
    \item \textbf{Imposibles} (ni jurídica ni prácticamente): no puede pactarse lo imposible. Cruzar el océano Atlántico nadando, por ejemplo, no puede ser objeto de un acto jurídico por tratarse de un hecho prácticamente imposible.
    \item \textbf{Prohibidos por la ley}: si la ley prohíbe determinado hecho, este no puede ser objeto de un acto jurídico.
    \item \textbf{Contrarios a la moral y a las buenas costumbres}: constituye una exigencia fundamental que el acto jurídico se adecue a los principios morales básicos.
\end{enumerate}

\subsection{Los bienes como objeto del acto}

El artículo 279 hace referencia genérica a prohibiciones que pueden surgir de la ley. Así, por ejemplo, ciertos bienes no pueden ser objeto de determinados contratos: los bienes de dominio público no pueden ser objeto de un acto jurídico de disposición.

\section{Objeto ilícito: el caso de la venta de influencias}

\begin{example}{Venta de influencias y acto de corrupción}
Una persona se compromete a realizar tratativas para obtener un privilegio en sede administrativa. En el fondo, se trata de una venta de influencias, un acto de corrupción.

Este acto tiene un objeto ilícito porque vender influencias es contrario a la ética pública y constituye un acto de corrupción que afecta el orden público.

Si posteriormente esa persona pretende hacer valer la obligación —por ejemplo, cobrar lo que le fue prometido—, la consecuencia es la \textbf{nulidad del acto}. La nulidad surge precisamente porque el objeto no responde a los requisitos del artículo 279. Se trata de una nulidad de tipo \textbf{absoluta}, regulada en el artículo 386 del Código Civil y Comercial.
\end{example}
\chapter{La causa del acto jurídico}
\label{ch:causa}

\section{Causa fuente y causa fin}

El concepto de causa admite varias acepciones en el derecho civil. En particular, se distingue entre la causa fuente y la causa fin.

\begin{definition}{Causa fuente}
Es la fuente de las obligaciones, es decir, el origen del cual surgen. Las obligaciones surgen de los contratos, en general de los actos jurídicos, y también de los actos ilícitos. No puede haber una obligación sin una causa, sin un elemento que le dé origen.
\end{definition}

\begin{definition}{Causa fin}
Es el fin perseguido por las partes al momento de celebrar un acto: qué se propusieron las partes al celebrarlo y cuál es la finalidad que persiguen.
\end{definition}

El desarrollo que sigue se concentra en el problema de la \textbf{causa fin}, y no en la causa fuente.

\section{Evolución histórica de la teoría de la causa}

El concepto de causa ha sido objeto de debate a lo largo de la historia del derecho, dando lugar a distintas posturas doctrinarias. Elementos vinculados a la valoración de los actos aparecen desde la antigüedad, especialmente con los canonistas y los glosadores.

\subsection{Teoría causalista}

Esta teoría, atribuida a un autor francés, sostiene que la causa fin de un contrato es la prestación de la otra parte. Se presenta como algo evidente que, al celebrar un contrato sinalagmático —perfectamente bilateral, en el que una parte tiene una prestación que la obliga y la otra, a su vez, otra prestación que también la obliga—, cada parte persigue como causa la prestación de la contraria.

\begin{example}{Compraventa bajo la teoría causalista}
En una compraventa, la causa por la cual una parte celebra el contrato es adquirir la cosa, y para la otra parte, la causa es recibir el precio.
\end{example}

\subsection{Crítica anticausalista}

Un autor belga —mencionado posteriormente por otros autores, entre ellos Planiol— critica la tesis causalista. El argumento principal sostiene que, si la causa fin es la prestación de la otra parte, entonces la causa se identifica con el objeto mismo del contrato, sin ofrecer ninguna novedad explicativa.

La crítica anticausalista señala que la idea de causa fin no agrega nada que no esté ya detrás del objeto: si el objeto es la transferencia de la propiedad (en una compraventa) y la causa es adquirir esa propiedad, causa y objeto se superponen. La postura anticausalista afirma, en consecuencia, que la causa no constituye un elemento independiente.

\subsection{Teoría moderna de la causa}

Enrique Capitant, hacia comienzos del siglo XX, retoma la idea de la causa y la profundiza. Para Capitant, la causa relevante es la finalidad perseguida por las partes, una finalidad que estas incorporan y que buscan en el momento de celebrar el acto.

\section{La causa en el Código de Vélez}

El Código Civil de Vélez no menciona la causa entre los elementos del acto de manera explícita en la sección específica de actos jurídicos. Sin embargo, el tema de la causa aparecía en algunos artículos ubicados en la parte de obligaciones (artículos 499 a 502), en el Libro Segundo, que se abría con la sección de obligaciones.

\begin{definition}{Artículo 499, Código Civil de Vélez}
No hay obligación sin causa. Se entiende por tal el hecho o acto lícito o ilícito en la relación de familia, o entre extraños, que es la causa generadora del derecho que faculta al acreedor a exigir del deudor la prestación debida.
\end{definition}

Este artículo se refiere claramente a la \textbf{causa fuente} (el hecho o acto que genera la obligación). Sin embargo, en los artículos 500 y 502, y especialmente en las notas que Vélez Sarsfield colocaba al pie de la ley —donde se critica la confusión entre causa fuente y causa fin—, aparece la idea de causa fin vinculada a la finalidad perseguida.

\begin{definition}{Artículo 502, Código Civil de Vélez (principio de la causa ilícita)}
De ningún efecto es el acto constituido con una causa ilícita, si esa causa es contraria a las leyes o al orden público.
\end{definition}

Aunque la causa no estaba expresamente mencionada como elemento en el Código de Vélez, existía la presunción de que, si las partes no expresaban la causa, esta debía existir igualmente, y la obligación se mantenía válida aun cuando la causa expresada fuera falsa.

\begin{important}
Una obligación con causa ilícita carece de todo efecto. Si la causa es ilícita, si contraviene las leyes o el orden público, el acto es nulo.
\end{important}

% TODO: contenido incompleto o ambiguo en las notas originales — el apunte menciona una observación del docente sobre el "artículo 429 del Código de Vélez", indicando que, cuando dicho artículo se refiere a una causa ilícita, en realidad remite a la causa fuente y no a la causa fin (ejemplo dado: quien causa un daño a otro debe repararlo, lo cual es causa fuente de la obligación de reparar). El apunte no precisa con exactitud el contenido ni la ubicación de dicho artículo dentro del sistema de Vélez, por lo que se conserva la referencia tal como fue consignada en clase.

\section{La causa en el Código Civil y Comercial}

El nuevo Código Civil y Comercial incorpora expresamente la causa como elemento del acto jurídico. El artículo 281 es la disposición central de esta regulación.

\begin{definition}{Artículo 281, Código Civil y Comercial}
La causa es el fin inmediato perseguido por las partes que ha sido determinante de la voluntad. Cuando las partes no la expresan, se presume que existe. Se presume también su licitud, a menos que resulte ilícita del acto mismo. También integran la causa los motivos exteriorizados cuando sean lícitos y hayan sido incorporados al acto en forma expresa o tácita, y sean esenciales para ambas partes.
\end{definition}

En rigor, lo que resulta relevante de esta definición es que la causa es el fin perseguido por las partes que ha sido determinante de la voluntad.

\subsection{Dimensiones objetiva y subjetiva de la causa}

El artículo 281 incorpora dos dimensiones de la causa:

\begin{enumerate}
    \item \textbf{Dimensión objetiva}: el fin inmediato autorizado por el ordenamiento jurídico. Es la finalidad típica que el derecho reconoce como válida para ese tipo de acto. En una compraventa, la causa objetiva es transferir la propiedad de una cosa a cambio de un precio.
    \item \textbf{Dimensión subjetiva}: los motivos o intenciones de las partes. Estos integran la causa cuando son exteriorizados (expresados o manifestados), lícitos, y esenciales para ambas partes; es decir, cuando ambas partes conocen y aceptan esos motivos como parte de la razón por la cual celebran el acto.
\end{enumerate}

\begin{example}{El caso de la profesora de francés}
Una profesora de francés contrata a un escritor para que redacte un ensayo que ella presentará, traducido, en un concurso destinado a obtener una beca para viajar a Francia. La profesora presentará el ensayo como propio.

En este caso se distingue con claridad entre objeto y causa:
\begin{itemize}
    \item El \textbf{objeto} del acto es contratar a alguien para que realice un ensayo o un trabajo. Esto es, en sí mismo, perfectamente válido y lícito.
    \item La \textbf{causa} (finalidad perseguida) es presentar el ensayo como propio para engañar a un concurso y obtener una beca. Esto es ilícito.
\end{itemize}

Si una parte persigue una causa ilícita pero la otra parte lo ignora, esta última no resulta afectada: el motivo ilícito no integra el acto. Pero si se interpreta que ambas partes conocen la causa y esta es esencial para ambas, entonces dicha causa integra el acto y lo vicia.
\end{example}

\subsection{Síntesis de las características de la causa}

\begin{enumerate}
    \item El Código Civil y Comercial incorpora la causa como elemento del acto jurídico en el artículo 281.
    \item La causa comprende dos dimensiones: una \textbf{objetiva} (el fin inmediato autorizado por el ordenamiento, determinante de la voluntad) y una \textbf{subjetiva} (los motivos o intenciones de las partes).
    \item Solo integran la causa los motivos o intenciones cuando: (a) son exteriorizados, (b) son lícitos, y (c) son esenciales para ambas partes.
    \item Quedan excluidos de la causa los puros motivos o intenciones personales que no estén exteriorizados ni sean esenciales para ambas partes.
\end{enumerate}

\section{Régimen de la causa}

\begin{definition}{Artículo 282, Código Civil y Comercial}
Cuando la causa no está expresada, se presume que existe. Se presume también su licitud, a menos que resulte ilícita del acto mismo.
\end{definition}

El artículo 282 retoma los principios ya contenidos en los artículos 500, 501 y 502 del Código de Vélez:

\begin{itemize}
    \item Si la causa no está expresada, se presume que existe (presunción de causa).
    \item El acto vale aunque la causa expresada sea falsa. Si lo que se aparenta es falso pero existe una causa verdadera, se estará ante una simulación, aunque el acto siga siendo válido.
    \item Una causa ilícita vicia el acto: si la causa es ilícita —contraria a las leyes y al orden público—, el acto es nulo o de ningún efecto.
\end{itemize}

\section{Actos abstractos}

\begin{definition}{Artículo 1014, Código Civil y Comercial (en materia de contratos)}
La inexistencia, falsedad o ilicitud de la causa no son discutibles en el acto abstracto.
\end{definition}

% TODO: contenido incompleto o ambiguo en las notas originales — el apunte atribuye esta regla al "artículo 283" en el desarrollo oral de la clase, mientras que el recuadro normativo transcripto en el propio material la identifica como "artículo 1014". Se conserva la referencia tal como figura en las notas originales, sin resolver la discrepancia mediante conocimiento externo.

Un \textbf{acto abstracto} es aquel que se independiza de su causa y recibe un tratamiento autónomo. El caso típico es el pagaré.

\begin{example}{El pagaré}
Una persona firma un pagaré porque contrae una deuda con un hospital. Al ingresar al hospital, el pagaré adquiere autonomía propia: puede circular y transferirse a terceros. Si posteriormente esa persona ya pagó al hospital y alguien pretende ejecutar el pagaré, no puede discutirse la causa (el hecho de que ya se pagó). La causa por la cual se emitió el pagaré no se discute en los títulos abstractos.
\end{example}

\section{Repaso mediante el método Cornell}

El siguiente cuadro reproduce, en formato de tabla Cornell, las preguntas y respuestas centrales para el repaso activo de los contenidos desarrollados en estos capítulos.

\begin{center}
\begin{longtable}{
    >{\raggedright\arraybackslash}p{0.28\textwidth}
    >{\raggedright\arraybackslash}p{0.62\textwidth}
}
\toprule
\textbf{Preguntas / palabras clave} & \textbf{Notas / respuestas} \\
\midrule
\endfirsthead
\toprule
\textbf{Preguntas / palabras clave} & \textbf{Notas / respuestas} \\
\midrule
\endhead
\midrule
\multicolumn{2}{r}{\textit{Continúa en la página siguiente}} \\
\midrule
\endfoot
\bottomrule
\endlastfoot

\textbf{¿Qué es un acto jurídico?} &
Es el acto voluntario lícito con fin inmediato de producir la adquisición, modificación o extinción de derechos, relaciones o situaciones jurídicas (art.\ 259, Código Civil y Comercial). Lo fundamental es que la intención directa de generar consecuencias jurídicas distingue al acto jurídico del simple acto lícito. \\

\addlinespace

\textbf{¿Cuál es la diferencia entre acto jurídico y simple acto lícito?} &
En el acto jurídico, la producción de consecuencias jurídicas es el fin inmediato e intencional. En el simple acto lícito, las consecuencias jurídicas se producen, pero no son el fin buscado (por ejemplo, el pescador que se apropia de lo que pesca). Lo determinante es la intención o finalidad del sujeto respecto del efecto jurídico. \\

\addlinespace

\textbf{¿Cuáles son los elementos del acto jurídico?} &
1.\ Sujeto: la persona que realiza el acto. 2.\ Objeto: la materia o los bienes sobre los que recae el acto. 3.\ Forma: la manera de exteriorizar el acto. 4.\ Causa: la finalidad perseguida por las partes. \\

\addlinespace

\textbf{¿Qué requisitos debe cumplir el objeto?} &
Debe ser posible (física y jurídicamente); no debe estar prohibido por la ley; debe ser lícito, es decir, conforme a la moral, a las buenas costumbres y al orden público; y no debe ser de carácter extrapatrimonial, salvo autorización legal. \\

\addlinespace

\textbf{¿Qué es la causa del acto?} &
Es el fin inmediato perseguido por las partes que ha sido determinante de la voluntad (art.\ 281, Código Civil y Comercial). Incorpora una dimensión objetiva (el fin autorizado por el derecho) y una dimensión subjetiva (los motivos lícitos, exteriorizados y esenciales para ambas partes). Se distingue de la causa fuente (origen de la obligación), que resulta menos relevante a estos efectos. \\

\addlinespace

\textbf{¿Qué diferencia hay entre objeto ilícito y causa ilícita?} &
En el objeto ilícito, la materia misma sobre la que recae el acto es ilícita (por ejemplo, vender un bien prohibido o contratar un acto ilegal). En la causa ilícita, el objeto es lícito, pero la finalidad o el motivo que persiguen las partes es ilícito (por ejemplo, contratar la redacción de un ensayo cuyo fin es el fraude académico). Ambas vician el acto, aunque por razones distintas. \\

\addlinespace

\textbf{¿Qué es un acto abstracto?} &
Es el acto que se independiza de su causa y recibe un tratamiento autónomo. El ejemplo típico es el pagaré: una vez emitido, puede circular sin vincularse a la causa original, y dicha causa no puede discutirse al momento de su ejecución. La inexistencia, falsedad o ilicitud de la causa no son discutibles en el acto abstracto. \\

\midrule
\multicolumn{2}{p{0.90\textwidth}}{
\textbf{Resumen:} Los actos jurídicos constituyen el fenómeno jurídico fundamental mediante el cual las personas, voluntariamente, crean, modifican o extinguen derechos y obligaciones. Su estudio requiere comprender que un acto jurídico no es simplemente cualquier conducta lícita: su característica esencial es que su fin inmediato es la producción intencional de consecuencias jurídicas. Para que un acto sea válido debe reunir cuatro elementos: un sujeto (persona con capacidad), un objeto (materia lícita, posible y conforme a la moral), una forma (manera de exteriorización) y una causa (finalidad perseguida). El objeto y la causa constituyen los límites más importantes a la autonomía de la voluntad: no puede haber objeto imposible, prohibido o contrario a la moral, ni causa ilícita. La distinción entre objeto ilícito y causa ilícita es crucial, porque permite comprender que un acto puede tener un objeto perfectamente válido pero una finalidad viciada. El Código Civil y Comercial introduce innovaciones relevantes: incorpora la causa expresamente como elemento (art.\ 281), reconoce sus dimensiones objetiva y subjetiva, y establece presunciones de existencia y licitud de la causa. Salvo en los actos abstractos (como el pagaré), la causa es siempre discutible y determinante: una causa ilícita torna nulo el acto, sin importar cuán lícito sea su objeto.
} \\

\end{longtable}
\end{center}
\chapter{La forma del acto jurídico}
\label{ch:forma}
% ============================================================

\section{Concepto de forma}

\begin{definition}{Forma del acto jurídico}
La forma es la exteriorización de la voluntad de los sujetos respecto del objeto, en orden a conseguir el fin jurídico y su protección al celebrar el acto. Ningún acto puede carecer de forma; en este sentido, la forma es un elemento esencial de los actos jurídicos.
\end{definition}

La forma del acto jurídico puede ser muy variada. El ordenamiento contempla desde formas simples —como un gesto que expresa la voluntad de adquirir algo— hasta actos de gran solemnidad, como el matrimonio y el testamento, para los cuales la ley establece formas determinadas. Entre estos extremos, existen los contratos verbales, los contratos electrónicos celebrados mediante formularios por internet y los instrumentos privados y públicos.

La forma se vincula con dos cuestiones centrales del derecho civil:

\begin{itemize}
    \item La \textbf{validez} del acto: en ciertos casos, la forma es requisito de existencia del acto y su incumplimiento genera la nulidad.
    \item La \textbf{prueba} del acto: en otros casos, la forma es exigida por la ley con el único fin de facilitar la demostración de que el acto fue celebrado.
\end{itemize}

\section{Regulación en el Código Civil y Comercial}

El Código Civil y Comercial de la Nación (CCyCN) regula la forma del acto jurídico dentro del Capítulo Quinto dedicado a los actos jurídicos, en las siguientes secciones del Libro Primero, Título Cuarto:

\begin{itemize}
    \item \textbf{Sección Tercera:} Forma y prueba del acto jurídico.
    \item \textbf{Sección Cuarta:} Instrumento público.
    \item \textbf{Sección Quinta:} Escritura pública.
    \item \textbf{Sección Sexta:} Instrumentos privados y particulares.
\end{itemize}

\section{El formalismo en el derecho: ventajas y desventajas}

La regulación jurídica de las formas no es neutral. Exige ponderarse en función de sus ventajas y desventajas para el tráfico jurídico.

\subsection{Ventajas del formalismo}

\begin{table}[H]
\centering
\begin{tabularx}{\textwidth}{>{\bfseries\raggedright\arraybackslash}p{0.35\textwidth}
                              >{\raggedright\arraybackslash}X}
\toprule
Ventaja & Explicación \\
\midrule
Facilita la prueba del acto & El acto queda documentado y puede acreditarse con mayor certeza. \\[4pt]
Publicita el acto & El acto es conocido por los demás; ingresa al conocimiento público. \\[4pt]
Colabora en la percepción de impuestos & Como ocurre con los contratos de alquiler que deben informarse a la AFIP. \\[4pt]
Permite distinguir el acto concluido de los preparatorios & La forma cristaliza el acto jurídico. Los borradores de un testamento no son un testamento; éste queda concluido con la forma mandada por la ley. \\[4pt]
Protege a terceros & Los derechos adquiridos respecto de terceros resultan tutelados cuando se cumple la forma. El tercero conoce su situación jurídica; esto es especialmente relevante en materia de inmuebles. \\
\bottomrule
\end{tabularx}
\end{table}

\begin{example}{Testamento — cristalización del acto}
Una persona puede estar pensando en hacer un testamento, consultar a distintas personas y elaborar borradores. En un momento determinado, la forma cristaliza ese acto jurídico. Los borradores no sirven: no son un testamento. El testamento queda concluido con la forma mandada por la ley: por escritura pública o por instrumento ológrafo.
\end{example}

\subsection{Desventajas del formalismo}

Las formas más estrictas presentan también inconvenientes:

\begin{itemize}
    \item Las formas son más costosas y engorrosas.
    \item Existe mayor riesgo de que el acto resulte inválido por algún defecto formal.
    \item Le dan \textbf{menor celeridad} a las transacciones.
    \item Generan incomodidades indudables en el tráfico jurídico cotidiano.
\end{itemize}

Por estas razones, el ordenamiento jurídico impone formas más exigentes únicamente a los actos más importantes o trascendentes, donde la mayor seguridad justifica la mayor incomodidad.

% ============================================================
\chapter{Clasificación de los Actos Jurídicos según la Forma}
% ============================================================

\section{Forma esencial y forma instrumental}

Es necesario distinguir dos sentidos del concepto de forma:

\begin{itemize}
    \item \textbf{Forma esencial:} es la manifestación de la voluntad que realizan las partes para celebrar el acto jurídico.
    \item \textbf{Forma instrumental:} es cuando la ley impone que algunos actos deben celebrarse mediante documentos escritos, ya sean públicos o privados.
\end{itemize}

El Código Civil y Comercial se concentra en la forma instrumental. Los instrumentos podrán ser públicos o privados.

\section{Primera clasificación: actos formales y no formales}

\begin{definition}{Actos formales}
Son aquellos respecto a los cuales la ley impone el cumplimiento de ciertas formas o solemnidades que deben cumplirse para la validez del acto. Ejemplo: el testamento, el matrimonio.
\end{definition}

\begin{definition}{Actos no formales}
Son aquellos que no tienen ninguna solemnidad prescripta por la ley; las partes son libres de realizarlos con libertad de formas.
\end{definition}

\section{Segunda clasificación: actos solemnes y no solemnes}

Los actos formales, a su vez, se clasifican en solemnes y no solemnes:

\begin{definition}{Actos formales solemnes}
Son aquellos en los que el rigor formal impuesto por la ley se vincula con la validez misma del acto. Si no se cumple con la forma, el acto es inválido.
\end{definition}

\begin{definition}{Actos formales no solemnes (ad probationem)}
Son aquellos en los que la forma es prescripta por la ley únicamente para facilitar la prueba del acto, sin afectar su validez.
\end{definition}

\section{Solemnidad absoluta y solemnidad relativa}

Dentro de los actos formales solemnes, el Código establece una distinción fundamental:

\begin{table}[H]
\centering
\begin{tabularx}{\textwidth}{>{\raggedright\arraybackslash}X
                              >{\raggedright\arraybackslash}X}
\toprule
\textbf{Solemnidad absoluta} & \textbf{Solemnidad relativa} \\
\midrule
La solemnidad es parte esencial del acto. Si no se cumple, el acto es ineficaz por nulidad de forma y no produce ningún efecto. & El acto no es válido (no genera los efectos propios del acto), pero genera la obligación de otorgar la forma prescripta. \\[6pt]
\textit{Ejemplo:} El testamento ológrafo que no está escrito y firmado de puño y letra por el testador. & \textit{Ejemplo:} El boleto de compraventa de un inmueble. No produce el efecto de la transferencia de la propiedad, pero sí genera la obligación de escriturar. \\
\bottomrule
\end{tabularx}
\end{table}

\subsection{Conversión del acto jurídico}

\begin{formula}{Art. 285 CCyCN — Conversión del acto jurídico}
\textit{El acto que no se otorga en la forma exigida por la ley no queda concluido como tal, pero vale como acto en el que las partes se obligan a cumplir con la expresada formalidad, salvo que ella se exija bajo sanción de nulidad.}
\end{formula}

La \textbf{conversión del acto jurídico} opera en los supuestos de solemnidad relativa: el acto no produce los efectos propios del negocio jurídico que se procuraba celebrar, pero se convierte en otro acto —la obligación de otorgar el instrumento proyectado—. Cuando la forma es exigida bajo sanción de nulidad (solemnidad absoluta), el incumplimiento genera la nulidad y no produce ningún otro efecto.

\begin{example}{Boleto de compraventa de inmueble}
Una compraventa de inmueble celebrada mediante boleto de compraventa (instrumento privado) no produce el efecto traslativo de dominio, ya que la ley exige escritura pública para ello. Sin embargo, en virtud de la conversión del artículo 285, genera la obligación de las partes de escriturar, es decir, de otorgar el instrumento proyectado.
\end{example}

% ============================================================

\section{Principio de Libertad de Formas}
% ============================================================

\begin{formula}{Art. 284 CCyCN — Libertad de formas}
\textit{Si la ley no designa una forma determinada para la exteriorización de la voluntad, las partes pueden utilizar la que estimen conveniente. Las partes pueden convenir una forma más exigente que la impuesta por la ley.}
\end{formula}

El principio de libertad de formas es la regla general del sistema. Si la ley no designa una forma determinada, las partes son libres de elegir la que consideren conveniente. Esta libertad tiene dos dimensiones:

\begin{itemize}
    \item \textbf{Libertad de elección:} ante el silencio de la ley sobre la forma, las partes determinan libremente el modo de exteriorizar su voluntad.
    \item \textbf{Autonomía para elevar la forma:} las partes pueden convenir una forma más exigente que la impuesta por la ley. Por ejemplo, celebrar un contrato de locación por escritura pública, aunque la ley no lo requiera.
\end{itemize}

\begin{example}{Locación por escritura pública}
Las partes de un contrato de alquiler pueden convenir celebrarlo por escritura pública, aunque la ley no exija esa forma para la locación. Ello implica costos adicionales (honorarios del escribano), pero resulta perfectamente válido.
\end{example}

Por supuesto, cuando la ley manda una forma determinada, las partes deben partir de la forma legalmente impuesta. El artículo 284 es el punto de partida para aquellos actos respecto de los cuales la ley no designa forma especial.

% ============================================================
\chapter{Los Instrumentos — Clasificación General}
% ============================================================

\section{Concepto de instrumento}

El concepto de instrumento refiere a la forma escrita. Es la expresión de la voluntad a través de un documento que da mayor seguridad y certeza acerca de los derechos y obligaciones asumidos por las partes.

\begin{formula}{Art. 286 CCyCN — Expresión escrita}
\textit{La expresión escrita puede tener lugar por instrumentos públicos, o por instrumentos particulares firmados o no firmados. Puede hacerse constar en cualquier soporte, siempre que su contenido sea representado con texto inteligible, aunque su lectura exija medios técnicos.}
\end{formula}

\begin{important}
La referencia a medios técnicos contempla los soportes informáticos y electrónicos. El Código extiende la categoría de instrumento a todo soporte inteligible, independientemente del medio requerido para su lectura.
\end{important}

\section{Clasificación de los instrumentos}

\begin{table}[H]
\centering
\begin{tabularx}{\textwidth}{>{\raggedright\arraybackslash}p{0.30\textwidth}
                              >{\raggedright\arraybackslash}p{0.30\textwidth}
                              >{\raggedright\arraybackslash}X}
\toprule
\textbf{Instrumentos públicos} & \textbf{Instrumentos privados} & \textbf{Particulares no firmados} \\
\midrule
Interviene un oficial público. Tienen validez y fuerza probatoria por sí mismos con los alcances que les concede el Código. & Surgen de las partes. Son instrumentos particulares firmados. La \textbf{firma} es un requisito esencial. & Todo escrito no firmado: impresos, registros visuales o auditivos de cosas o hechos, registros de la palabra y de la información, cualquiera sea el medio empleado. \\
\bottomrule
\end{tabularx}
\end{table}

\begin{formula}{Art. 287 CCyCN — Instrumentos particulares}
\textit{Los instrumentos particulares pueden estar firmados o no. Si lo están, se llaman instrumentos privados. Si no lo están, se los denomina instrumentos particulares no firmados; esta categoría comprende todo escrito no firmado, entre otros, los impresos, los registros visuales o auditivos de cosas o hechos y, cualquiera que sea el medio empleado, los registros de la palabra y de la información.}
\end{formula}

El Código Civil y Comercial amplió el tratamiento jurídico de esta materia respecto del Código anterior, incorporando la categoría de instrumentos particulares no firmados para dar cobertura a la realidad del tráfico contemporáneo, caracterizado por la informatización: formularios web, constancias electrónicas, mensajes de texto, registros digitales.

\begin{example}{Constancia de compra por internet}
Una constancia generada por una plataforma de comercio electrónico al completar un formulario de compra en línea —que registra los datos del contrato, el precio y el método de pago— es un instrumento particular no firmado. No tiene firma en sentido estricto, por lo que no es instrumento privado con su fuerza probatoria plena, pero puede tener valor como prueba en caso de conflicto.
\end{example}

\section{La firma como disposición general}

\begin{formula}{Art. 288 CCyCN — Firma}
\textit{La firma prueba la autoría de la declaración de voluntad expresada en el texto al cual corresponde. La firma consiste en el nombre del firmante o en un signo; en ambos casos se realiza con el trazo habitual y de tal manera que identifique a la persona.}
\end{formula}

En el Código Civil y Comercial, la regulación de la firma fue trasladada de la sección de instrumentos privados —donde se encontraba en el Código de Vélez— a las disposiciones generales sobre la forma, lo que significa que rige tanto para los instrumentos públicos (incluyendo la escritura pública) como para los privados. La firma:

\begin{itemize}
    \item Prueba la \textbf{autoría} de la declaración de voluntad.
    \item Expresa la \textbf{conformidad} con el contenido del documento.
    \item Consiste en el nombre o en un signo habitual e identificatorio de la persona.
\end{itemize}

\subsection{Firma digital y firma electrónica}

El artículo 288 contempla también la cuestión de los instrumentos generados por medios electrónicos y remite a la normativa específica sobre firma digital.

\begin{important}
La \textbf{firma digital} no debe confundirse con el agregado de datos personales al pie de un correo electrónico (nombre, cargo, teléfono). Este último no constituye firma digital en sentido jurídico.
\end{important}

\begin{definition}{Firma digital (Ley 25.506 — año 2001)}
Se entiende por firma digital al resultado de aplicar a un documento digital un procedimiento matemático (algoritmo) que requiere información de exclusivo conocimiento del firmante y que se encuentra bajo su absoluto control. La firma digital debe haber sido: (a) creada durante el período de vigencia del certificado digital válido del firmante; (b) verificada por referencia a los datos de verificación de firma; y (c) generada por un certificado emitido por un certificador licenciado.
\end{definition}

\begin{definition}{Firma electrónica}
Es un conjunto de datos electrónicos, ligados o asociados de manera lógica a otros datos electrónicos, que el signatario utiliza con fines de identificación, pero que no alcanza el estándar de la firma digital. No cuenta con certificación indubitable de autoría e integridad.
\end{definition}

La diferencia entre ambas figuras es fundamental en materia probatoria:

\begin{table}[H]
\centering
\begin{tabularx}{\textwidth}{>{\raggedright\arraybackslash}X
                              >{\raggedright\arraybackslash}X}
\toprule
\textbf{Firma digital} & \textbf{Firma electrónica} \\
\midrule
Certificación digital indubitable de \textbf{autoría} e \textbf{integridad} mediante tercero certificador licenciado. El documento firmado digitalmente no puede ser modificado con posterioridad sin que ello resulte detectable. & Identificación electrónica sin certificación indubitable. Puede ser desconocida; quien la invoca debe acreditar su validez mediante prueba. \\[6pt]
\textbf{Mayor valor probatorio.} Garantizadas la autoría y la integridad. & \textbf{Menor valor probatorio.} La autoría y la integridad no están garantizadas por estas certificaciones. \\
\bottomrule
\end{tabularx}
\end{table}

% ============================================================
\chapter{Instrumentos Públicos}
% ============================================================

\section{Concepto y caracteres}

\begin{definition}{Instrumento público}
Son los documentos a los cuales la ley les asigna un valor de autenticidad por sí mismos y que, como tales, no requieren de un reconocimiento previo de la firma para hacer plena fe de su contenido. La autenticidad del instrumento público proviene de la ley, no de la voluntad de las partes.
\end{definition}

En la mayoría de los instrumentos públicos interviene un oficial público —funcionario público o escribano— que certifica la validez del acto. Sin embargo, lo que otorga la fuerza probatoria propia al instrumento público es la ley: es ella la que, ante la intervención de determinado funcionario revestido de esa credibilidad social e institucional, concede al instrumento el carácter de auténtico.

Los caracteres esenciales del instrumento público son:

\begin{itemize}
    \item \textbf{Autenticidad:} hace plena fe de su contenido sin necesidad de reconocimiento previo.
    \item \textbf{Fuerza probatoria propia:} su valor probatorio proviene de la ley.
    \item \textbf{Oponibilidad erga omnes:} tiene valor frente a todos, no sólo entre las partes.
\end{itemize}

\section{Enumeración legal}

\begin{formula}{Art. 289 CCyCN — Enumeración de instrumentos públicos}
\textit{Son instrumentos públicos: a) Las escrituras públicas y sus copias o testimonios; b) Los instrumentos que extienden los escribanos o los funcionarios públicos con los requisitos que establecen las leyes; c) Los títulos emitidos por el Estado nacional, provincial o la Ciudad Autónoma de Buenos Aires, conforme a las leyes que autorizan su emisión.}
\end{formula}

Cada categoría comprende:

\begin{enumerate}
    \item \textbf{Escrituras públicas y sus copias o testimonios:} instrumento matriz del protocolo notarial y sus reproducciones autenticadas.

    \item \textbf{Instrumentos de escribanos o funcionarios públicos:} comprende, entre otros, las actas notariales, las sentencias judiciales homologadas, las actas celebradas en sede judicial, y los convenios celebrados ante funcionario competente (como un convenio de alimentos labrado en sede judicial).

    \item \textbf{Títulos emitidos por el Estado:} títulos de deuda pública nacional, provincial o de la Ciudad Autónoma de Buenos Aires, emitidos conforme a las leyes que autorizan su emisión.
\end{enumerate}

\section{Requisitos de validez}

\begin{formula}{Art. 290 CCyCN — Requisitos de validez del instrumento público}
\textit{Son requisitos de validez del instrumento público: a) La actuación del oficial público en los límites de sus atribuciones y de su competencia territorial, excepto que el lugar sea generalmente tenido como comprendido en ella; b) Las firmas del oficial público, de las partes, y en su caso, de sus representantes; si alguno de ellos no firma por sí o no puede hacerlo, otro de los intervinientes lo hará a su ruego, haciéndose constar esta circunstancia.}
\end{formula}

Los dos requisitos esenciales son:

\begin{itemize}
    \item \textbf{Competencia del oficial público:} comprende la competencia \emph{funcional} (atribuciones para celebrar ese tipo de acto) y la competencia \emph{territorial} (actuación dentro de la jurisdicción asignada). La actuación fuera de estos límites invalida el instrumento.
    \item \textbf{Firmas:} deben firmar el oficial público, las partes y, en su caso, los representantes. La falta de firma de cualquiera de los intervinientes invalida el instrumento para todos.
\end{itemize}

\begin{important}
Es presupuesto de validez que el oficial público se encuentre efectivamente en funciones. Los actos otorgados antes de la notificación de la suspensión o cesación de funciones son válidos. Si la persona ejercía efectivamente un cargo existente bajo apariencia de legitimidad, la falta de investidura no afecta al acto ni al instrumento dentro de los límites de la buena fe.
\end{important}

\section{Defectos de forma}

Los defectos de forma en el instrumento público —enmiendas, agregados, borraduras, entrelineados, alteraciones— deben ser salvados al final del instrumento antes de las firmas. Si no se salvan, el instrumento carece de validez en esa parte.

\begin{example}{Enmienda en número de DNI}
Si en una escritura pública se consignó erróneamente el número del documento nacional de identidad de uno de los comparecientes y se corrigió mediante una enmienda manuscrita sobre el texto impreso, esa corrección debe ser salvada al final del instrumento antes de las firmas. De lo contrario, la enmienda no salvada afecta la validez del instrumento.
\end{example}

\begin{note}
Si el instrumento contiene las firmas de las partes pero adolece de un defecto que le impide valer como instrumento público (por ejemplo, por incompetencia del escribano), puede igualmente valer como \textbf{instrumento privado}, en la medida en que el acto de que se trate no exija forma pública bajo sanción de nulidad.
\end{note}

\section{Eficacia probatoria del instrumento público}

Este es uno de los aspectos de mayor relevancia práctica. El instrumento público hace plena fe, pero el Código distingue dos situaciones según el tipo de contenido:

\begin{formula}{Art. 296 CCyCN — Eficacia probatoria del instrumento público}
\textit{El instrumento público hace plena fe: a) En cuanto a que se ha realizado el acto, la fecha, el lugar y los hechos que el oficial público enuncia como cumplidos por él o ante él hasta que sea declarado falso en juicio civil o criminal; b) En cuanto al contenido de las declaraciones sobre convenciones, disposiciones, pagos, reconocimientos y enunciaciones de hechos directamente relacionados con el objeto principal del acto instrumentado, hasta que se produzca prueba en contrario.}
\end{formula}

\subsection{Hechos percibidos directamente por el oficial público}

Los \textbf{hechos que el oficial público enuncia como cumplidos por él o ante él} —la fecha, el lugar, la celebración del acto, los hechos ocurridos en su presencia— hacen plena fe y sólo pueden ser cuestionados mediante un \textbf{juicio de redargución de falsedad} (acción civil o criminal tendiente a demostrar la falsedad del instrumento). Quien pretenda impugnarlos debe acusar al oficial público de haber faltado a la verdad, lo que supone una acusación de suma gravedad dada la credibilidad que la sociedad deposita en la función notarial.

\begin{example}{Acta notarial — constancia de hechos ante el escribano}
Si un escribano labra un acta en la que consigna: ``En el día de la fecha me constituí en el domicilio de la calle X y constaté que se realizaba la actividad Y, en presencia de las personas Z'', quien pretenda desacreditar ese hecho debe iniciar un juicio de redargución de falsedad, dado que estaría afirmando que el escribano mintió en su función oficial.
\end{example}

\subsection{Declaraciones de las partes}

En cuanto al \textbf{contenido de las declaraciones de las partes} —convenciones, disposiciones, pagos, reconocimientos y enunciaciones de hechos directamente relacionados con el objeto del acto—, el instrumento hace plena fe pero \textbf{admite prueba en contrario}. En este caso, no es necesario iniciar un juicio de redargución de falsedad: basta con demostrar por otros medios que el contenido de esas declaraciones no es veraz.

\begin{example}{Reconocimiento de deuda simulado}
Si en una escritura pública una de las partes reconoce adeudar una suma de dinero, ese reconocimiento es el contenido de su declaración de voluntad. Si posteriormente se pretende demostrar que ese reconocimiento fue simulado, no es necesaria la redargución de falsedad —el escribano no mintió: certificó que el compareciente hizo esa declaración—, sino que puede acreditarse la simulación por prueba en contrario.
\end{example}

% ============================================================
\chapter{La escritura pública y las actas notariales}
% ============================================================

\section{Concepto y protocolo notarial}

\begin{formula}{Art. 299 CCyCN — Escritura pública. Definición}
\textit{La escritura pública es el instrumento matriz extendido en el protocolo de un escribano público o de otro funcionario autorizado para ejercer las mismas funciones, que contiene uno o más actos jurídicos.}
\end{formula}

La escritura pública es el tipo de instrumento público más importante en la regulación civil y comercial. Sus características esenciales son:

\begin{itemize}
    \item \textbf{Instrumento matriz:} la escritura original es aquella que queda integrada al protocolo del escribano. Lo que las partes conservan en sus domicilios es una \textbf{copia} o \textbf{testimonio}, que también tiene carácter de instrumento público y hace plena fe tanto como la escritura matriz.
    \item \textbf{Contenido:} la escritura contiene \textbf{actos jurídicos}. Esto la distingue del acta notarial, cuyo objeto es la comprobación de \textbf{hechos}.
    \item \textbf{Función notarial:} los escribanos son regulados en cada provincia y en la Ciudad Autónoma de Buenos Aires. El Código Civil y Comercial establece reglas comunes para este instrumento con validez en todo el país.
\end{itemize}

\begin{important}
Si existe discordancia entre el texto de la copia o testimonio y el de la escritura matriz, prevalece el contenido de la escritura matriz.
\end{important}

\subsection{El protocolo notarial}

El protocolo está conformado por folios numerados correlativamente. No pueden quedar hojas en blanco —lo que implicaría la posibilidad de antedatar documentos. Los asientos deben realizarse en orden cronológico. La reglamentación del libro del protocolo es de competencia local (provincial o de la Ciudad Autónoma de Buenos Aires).

\section{Requisitos de la escritura pública}

La escritura pública debe ser recibida personalmente por el escribano, quien tiene el \textbf{deber de calificación} de todos los presupuestos y elementos del acto: verifica la identidad de los comparecientes, su capacidad, la legitimidad del acto y la regularidad de los títulos.

\subsection{Estructura: encabezado, cuerpo y pie}

El contenido de la escritura pública (artículo 305 CCyCN) se organiza en tres grandes partes:

\begin{table}[H]
\centering
\begin{tabularx}{\textwidth}{>{\bfseries\raggedright\arraybackslash}p{0.20\textwidth}
                              >{\raggedright\arraybackslash}X}
\toprule
Parte & Contenido \\
\midrule
Encabezado & Lugar y fecha de otorgamiento (con posibilidad de incluir la hora). Datos de los comparecientes: nombre, apellido, documento nacional de identidad, domicilio real, estado civil. En el caso de segundas nupcias, se hace constar. Nombre del cónyuge, relevante por cuestiones vinculadas a la vivienda y al régimen patrimonial. Si comparece una persona jurídica: denominación completa, domicilio y datos del representante y de su mandato. \\[4pt]
Cuerpo & Naturaleza del acto e individualización de los bienes que constituyen su objeto. En las transmisiones de derechos reales sobre inmuebles, incluye el \textbf{``corresponde''}: la relación de antecedentes de titularidad del bien transmitido. \\[4pt]
Pie / Cierre & Constancia de la lectura del instrumento por el escribano ante los comparecientes. Salvado de enmiendas, borrones o entrelineados realizados antes de las firmas. Firma de los otorgantes, del escribano y de los testigos, si los hubiera. Si algún compareciente no puede firmar, otro firmará a su ruego, dejándose constancia del impedimento y de la impresión digital del otorgante. \\
\bottomrule
\end{tabularx}
\end{table}

\subsection{El ``corresponde'' en las transmisiones de inmuebles}

El \textbf{``corresponde''} es una cláusula obligatoria en las escrituras que instrumentan transmisiones de derechos reales sobre inmuebles. Contiene la relación de todos los antecedentes en virtud de los cuales el transmitente llegó a ser titular del derecho que transmite.

\begin{example}{Cadena de antecedentes}
Si el transmitente adquirió el inmueble por herencia, el corresponde consigna: ``Le corresponde por declaratoria de herederos de fecha X y partición de fecha Y. A su vez, el causante lo había adquirido por escritura de fecha Z, que obra en el folio N del Registro de Propiedad Inmueble.'' Esta cláusula permite encadenar las distintas escrituras y verificar la titularidad del bien a lo largo del tiempo.
\end{example}

La omisión del ``corresponde'' no genera la invalidez de la escritura, pero puede acarrear sanciones para el escribano en el ejercicio de su función profesional.

\section{Copias y testimonios}

Las copias o testimonios de la escritura matriz deben entregarse a las partes. Pueden reproducirse por cualquier medio que asegure su permanencia. Si alguna de las partes requiere una nueva copia, debe entregársele, salvo que exista una obligación de dar dinero pendiente de cumplimiento.

% ============================================================

\section{Actas Notariales}
% ============================================================

\begin{formula}{Art. 310 CCyCN — Actas notariales}
\textit{Se denominan actas los documentos notariales que tienen por objeto la comprobación de hechos.}
\end{formula}

\begin{definition}{Acta notarial}
Es el documento expedido por un escribano a requerimiento de parte, cuyo objeto es la comprobación de \textbf{hechos}. A diferencia de la escritura pública —que contiene actos jurídicos (Art. 299)—, el acta sólo constata hechos, sin crear, modificar ni extinguir relaciones jurídicas.
\end{definition}

El acta notarial implica la concurrencia del escribano a un lugar determinado para constatar un hecho y dejar constancia de él. Entre sus aplicaciones habituales se encuentran:

\begin{itemize}
    \item Constatación del estado de un bien al momento de su restitución (locaciones).
    \item Intimaciones de pago al deudor.
    \item Comprobación de que un acreedor se negó a recibir un pago.
    \item Constancia de hechos relevantes para el ejercicio de derechos.
\end{itemize}

\begin{example}{Acta de constatación al recuperar un local comercial}
Al recuperar la posesión de un local comercial después de una locación, el locador puede requerir al escribano que concurra al acto de entrega de llaves y labre un acta notarial que constate el estado en que se encuentra el inmueble. El escribano describe en el acta todo lo que constata, y ese documento tendrá valor probatorio en caso de reclamarse daños o el incumplimiento de las obligaciones de restitución.
\end{example}

Los requisitos del acta notarial están regulados en el artículo 311 CCyCN. El acta tiene valor probatorio respecto de:

\begin{itemize}
    \item Los \textbf{hechos} que el notario tuvo a la vista: su existencia, características y estado.
    \item Las \textbf{personas} que concurrieron, si quisieron identificarse.
    \item Las \textbf{declaraciones} efectuadas por los presentes.
\end{itemize}

\begin{note}
La Ciudad Autónoma de Buenos Aires regula el ejercicio de la función notarial mediante la Ley 404 (año 2000). El escribano puede ser titular o adscripto. El ingreso al notariado requiere la aprobación de un procedimiento de selección competitivo y una investidura formal.
\end{note}

% ============================================================
\chapter{Instrumentos particulares y fecha cierta}
% ============================================================

\section{Clasificación y jerarquía}

El Código Civil y Comercial sistematiza los instrumentos particulares en dos categorías:

\begin{itemize}
    \item \textbf{Instrumentos particulares firmados} (instrumentos privados): requieren firma de las partes como requisito esencial.
    \item \textbf{Instrumentos particulares no firmados}: todo escrito no firmado, registro visual o auditivo, o cualquier otro soporte inteligible sin firma.
\end{itemize}

La categoría más importante es la de los instrumentos privados (particulares firmados). Los instrumentos particulares no firmados constituyen una categoría incorporada por el nuevo Código para dar cobertura jurídica a la realidad del tráfico contemporáneo de bienes y servicios.

\section{Instrumentos privados: el requisito de la firma}

\begin{important}
El \textbf{requisito esencial} de los instrumentos privados es la \textbf{firma} de las partes (Art. 288). Sin firma, el documento es un instrumento particular no firmado, con menor fuerza probatoria.
\end{important}

El Código Civil y Comercial eliminó el requisito del \textbf{doble ejemplar} que imponía el Código anterior (Código de Vélez). En virtud de ello, ya no es jurídicamente exigible que los instrumentos se celebren en tantos ejemplares como partes intervinientes. No obstante, la buena práctica profesional recomienda mantener este recaudo, especialmente cuando una de las partes asume obligaciones relevantes.

\section{Reconocimiento de la firma}

El instrumento privado tiene una firma que prueba la declaración de voluntad, pero para que se produzca la plenitud de la fuerza probatoria entre las partes, la firma debe ser \textbf{reconocida}.

\begin{formula}{Art. 314 CCyCN — Reconocimiento de la firma}
\textit{Todo aquel contra quien se presente un instrumento cuya firma se le atribuye debe manifestarse al respecto. El silencio se tendrá por reconocimiento de la firma.}
\end{formula}

\begin{important}
El silencio frente a la intimación de reconocimiento de firma equivale al reconocimiento. Esto constituye un supuesto de silencio como manifestación de la voluntad (Art. 263 CCyCN): ante la obligación legal de expedirse, el silencio importa aceptación tácita.
\end{important}

El reconocimiento de la firma puede ser:

\begin{itemize}
    \item \textbf{Expreso:} la persona declara que la firma le pertenece.
    \item \textbf{Tácito:} guarda silencio ante la intimación, o realiza gestiones sin desconocer la autoría de la firma.
    \item \textbf{Forzoso:} se acredita mediante prueba, especialmente a través de la \textbf{pericia caligráfica} (cotejo de firmas realizado por un perito calígrafo designado judicialmente).
\end{itemize}

\begin{example}{Negación de contrato de locación}
Si el locatario niega haber suscripto el contrato de alquiler, el locador puede exhibirle el instrumento en juicio e intimarlo a que reconozca la firma. Si la niega, podrá acreditarse la autoría mediante pericia caligráfica, que establecerá con alta probabilidad la autenticidad de la firma atribuida al demandado.
\end{example}

El reconocimiento de la firma importa el reconocimiento de todo el cuerpo del instrumento: las cláusulas, disposiciones, convenciones y reconocimientos que contiene. Lo que el instrumento público tiene por sí, el instrumento privado lo adquiere con el reconocimiento de firma.

\section{Firma a ruego e impresión digital}

\begin{formula}{Art. 313 CCyCN — Firma a ruego e impresión digital}
\textit{Si alguno de los firmantes no sabe o no puede firmar, puede dejarse constancia de la impresión digital o bien de la firma de otra persona a ruego, en presencia de dos testigos que también lo suscriban.}
\end{formula}

Cuando una parte no sabe o no puede firmar, el Código prevé dos alternativas:

\begin{enumerate}
    \item \textbf{Firma a ruego:} otra persona firma en nombre del que no puede hacerlo, en presencia de dos testigos que también suscriben el instrumento. Esta modalidad tiene pleno valor probatorio, similar a un mandato.

    \item \textbf{Impresión digital:} la estampación de la huella dactilar del signatario. Sin embargo, conforme al artículo 314 CCyCN, el documento signado con impresión digital vale únicamente como \textbf{principio de prueba por escrito} y puede ser impugnado. Su menor valor probatorio obedece a que la impresión digital puede obtenerse sin el pleno consentimiento de la persona.
\end{enumerate}

\section{Documento firmado en blanco}

\begin{formula}{Art. 315 CCyCN — Documento firmado en blanco}
\textit{El firmante de un documento en blanco puede impugnar su contenido mediante la prueba de que no responde a sus instrucciones, pero no puede valerse de testigos si no existe principio de prueba por escrito. El desconocimiento del documento no puede invocarse para afectar al tercero que de buena fe celebró un acto jurídico basándose en él.}
\end{formula}

El documento firmado en blanco implica una suerte de mandato implícito: quien firma en blanco autoriza al receptor del documento a completarlo dentro de las instrucciones acordadas. Si el documento fue completado en forma contraria a esas instrucciones, el firmante puede impugnar su contenido, pero con importantes limitaciones probatorias:

\begin{itemize}
    \item No puede valerse exclusivamente de testigos si no existe principio de prueba por escrito.
    \item Si el instrumento fue utilizado de buena fe por un tercero que celebró un acto jurídico basándose en él, ese tercero queda protegido: el firmante no puede invocar el desconocimiento del documento en su perjuicio.
\end{itemize}

\begin{important}
La firma de documentos en blanco es una práctica de alto riesgo, ya que quien tiene el documento puede completarlo con obligaciones enormemente perjudiciales para el patrimonio del firmante.
\end{important}

\section{Enmiendas en instrumentos privados}

Las enmiendas, raspaduras o entrelineados que afecten partes esenciales del instrumento privado deben ser salvadas al final, con la firma de las partes. Diferencia con la escritura pública: en los instrumentos privados, las enmiendas no salvadas no generan la nulidad del instrumento, sino que el juez determinará en qué medida el defecto excluye o reduce su fuerza probatoria.

% ============================================================

\section{Fecha Cierta y Oponibilidad a Terceros}
% ============================================================

\subsection{El instrumento privado frente a terceros}

El instrumento privado tiene fuerza probatoria entre las partes desde el reconocimiento de la firma. Frente a terceros, sin embargo, rige una regla diferente: el instrumento privado \textbf{no es oponible} a terceros mientras no adquiera \textbf{fecha cierta}.

\begin{formula}{Art. 317 CCyCN — Fecha cierta}
\textit{La eficacia probatoria de los instrumentos privados reconocidos se extiende a los terceros desde su fecha cierta. Adquieren fecha cierta el día en que acontece un hecho del que resulta como consecuencia ineludible que el documento ya estaba firmado o no pudo ser firmado después.}
\end{formula}

\begin{definition}{Fecha cierta}
No es la fecha consignada en el instrumento, sino el momento a partir del cual existe certeza objetiva de que ese instrumento ya estaba firmado o no pudo ser firmado con posterioridad. Sirve para hacer oponible el instrumento privado a terceros.
\end{definition}

\subsection{Medios para adquirir fecha cierta}

El Código de Vélez (artículo 1035) enumeraba taxativamente los casos de adquisición de fecha cierta:

\begin{itemize}
    \item Exhibición del instrumento en juicio o ante una repartición pública donde quedara archivado.
    \item Reconocimiento ante un escribano y testigos.
    \item Transcripción en cualquier registro público.
    \item Fallecimiento de la parte que lo firmó.
\end{itemize}

El Código Civil y Comercial \textbf{abandona la enumeración taxativa} y establece que la fecha cierta puede acreditarse por cualquier medio de prueba, que el juez apreciará de manera rigurosa.

\begin{example}{Fecha cierta por exhibición en juicio}
Cuando un instrumento privado se incorpora a un expediente judicial, la mesa de entradas de tribunales le estampa un sello que acredita la fecha de presentación. Desde ese momento existe certeza de que el instrumento ya existía: puede oponerse a terceros como instrumento firmado con anterioridad a esa fecha.
\end{example}

\begin{example}{Fecha cierta por fallecimiento del firmante}
Si uno de los firmantes de un instrumento privado fallece, desde ese momento se tiene certeza de que el instrumento fue firmado antes de la muerte: con posterioridad al fallecimiento, esa persona ya no pudo haberlo suscripto. El instrumento adquiere fecha cierta desde el día del fallecimiento.
\end{example}

% ============================================================
\chapter{Instrumentos Particulares no Firmados y Correspondencia}
% ============================================================

\section{Concepto y amplitud}

Los instrumentos particulares no firmados comprenden todos los escritos no firmados, impresos, registros visuales o auditivos de cosas o hechos, cualquiera sea el medio empleado, y los registros de la palabra y de la información. Esta amplísima categoría abarca tanto los soportes en papel como los medios electrónicos: mensajes de texto, correos electrónicos, publicaciones en redes sociales, tickets electrónicos, registros de plataformas web.

\begin{note}
A diferencia de los instrumentos privados, los instrumentos particulares no firmados no cuentan con firma. La firma prueba la autoría de la declaración de voluntad; su ausencia implica que el instrumento tiene menor fuerza probatoria y que la valoración de su alcance queda librada a la apreciación judicial.
\end{note}

\section{Fuerza probatoria — criterios del juez}

\begin{formula}{Art. 319 CCyCN — Valor probatorio del instrumento particular}
\textit{El valor probatorio de los instrumentos particulares debe ser apreciado por el juez ponderando, entre otras pautas, la congruencia entre lo sucedido y narrado, la precisión y claridad técnica del texto, los usos y prácticas del tráfico, las relaciones precedentes y la confiabilidad de los soportes utilizados y de los procedimientos técnicos que se apliquen.}
\end{formula}

Las pautas que el juez debe ponderar son:

\begin{enumerate}
    \item \textbf{Congruencia entre lo sucedido y lo narrado:} correspondencia entre lo que dice el instrumento y los hechos ocurridos y acreditados por otros medios.

    \item \textbf{Precisión y claridad técnica del texto:} un formulario web completo y preciso tendrá mayor valor que un mensaje informal con redacción ambigua.

    \item \textbf{Usos y prácticas del tráfico:} la forma habitual en que se concluyen los actos jurídicos en ese sector de actividad o negocio.

    \item \textbf{Relaciones precedentes entre las partes:} el historial de intercambios previos que contextualice el instrumento en cuestión.

    \item \textbf{Confiabilidad de los soportes y procedimientos técnicos empleados:} la seguridad e integridad de los sistemas informáticos utilizados para generar o transmitir el instrumento.
\end{enumerate}

\begin{example}{Ticket electrónico como prueba}
En un conflicto sobre el pago del precio en una compraventa de bienes, el comprador presenta el ticket generado por la plataforma de comercio electrónico y el resumen de tarjeta de crédito que registra el débito correspondiente. La congruencia entre ambos elementos es un indicio de peso que el juez puede ponderar para tener por acreditado el pago.
\end{example}

\section{La correspondencia}

\begin{formula}{Art. 318 CCyCN — Correspondencia}
\textit{La correspondencia, cualquiera sea el medio empleado para crearla o transmitirla, puede presentarse como prueba por el destinatario, pero la que es confidencial no puede ser utilizada sin consentimiento del remitente.}
\end{formula}

La correspondencia —cartas, correos electrónicos, mensajes— pertenece, una vez recibida, al destinatario. Éste puede utilizarla como prueba en juicio contra el remitente, con la excepción de la \textbf{correspondencia confidencial}: si la comunicación tiene carácter confidencial, no puede ser utilizada sin el consentimiento del remitente.

La confidencialidad puede surgir:
\begin{itemize}
    \item De la manifestación expresa de las partes (por ejemplo, que el remitente la haya designado así).
    \item De la naturaleza misma de la comunicación (una carta dirigida a un amigo de confianza sobre asuntos personales).
\end{itemize}

Esta protección tiene fundamento constitucional: la \textbf{inviolabilidad de la correspondencia y los papeles privados} reconocida por la Constitución Nacional como garantía vinculada a la intimidad y privacidad de las personas.

\begin{example}{Uso de correspondencia confidencial}
Si alguien envía una carta a un amigo en la que reconoce adeudar sumas de dinero a terceros, y ese amigo entrega la carta a los acreedores para que la usen en un juicio, el uso de esa correspondencia puede ser cuestionado por su carácter confidencial. Sin el consentimiento del remitente, los destinatarios originales no pueden utilizarla.
\end{example}

% ============================================================
% CUADRO CORNELL
% ============================================================
\section{Cuadro Cornell — Repaso Integral}

\begin{table}[H]
\centering
\begin{tabularx}{\textwidth}{
    >{\raggedright\arraybackslash}p{0.30\textwidth}
    >{\raggedright\arraybackslash}X
}
\toprule
\textbf{Pregunta / Concepto clave} &
\textbf{Notas / Respuesta} \\
\midrule

\textbf{¿Qué es la forma del acto jurídico? ¿Es un elemento esencial?} &
Es la exteriorización de la voluntad de los sujetos respecto del objeto, en orden a conseguir el fin jurídico. Ningún acto puede carecer de forma: en ese sentido, la forma es un elemento esencial. \\[6pt]

\textbf{¿Cuáles son las dos dimensiones de la forma?} &
La forma tiene incidencia sobre la \textbf{validez} del acto (en ciertos casos, su incumplimiento produce nulidad) y sobre la \textbf{prueba} del acto (en otros casos, la forma es exigida sólo para facilitar su demostración). \\[6pt]

\textbf{¿Cuáles son las ventajas y desventajas del formalismo?} &
\textit{Ventajas:} facilita la prueba, publicita el acto, colabora en la percepción impositiva, permite distinguir el acto concluido de los preparatorios y protege a los terceros. \textit{Desventajas:} mayor costo, mayor riesgo de invalidez por defecto formal, menor celeridad en las transacciones, incomodidades prácticas. \\[6pt]

\textbf{¿Cuál es el principio general que rige la forma? (Art. 284)} &
El \textbf{principio de libertad de formas}: si la ley no designa una forma determinada, las partes pueden utilizar la que estimen conveniente. Las propias partes pueden pactar formas más exigentes que las legales. \\[6pt]

\textbf{¿Cómo se clasifican los actos según la forma?} &
Actos \textbf{formales} (con forma legal impuesta) y \textbf{no formales} (libertad de formas). Los formales se subdividen en \textbf{solemnes} (la forma es requisito de validez) y \textbf{no solemnes} o \emph{ad probationem} (la forma es exigida sólo para la prueba). \\[6pt]

\textbf{¿Qué diferencia hay entre solemnidad absoluta y solemnidad relativa?} &
En la \textbf{solemnidad absoluta}, el incumplimiento produce la nulidad del acto, sin ningún otro efecto. En la \textbf{solemnidad relativa}, el acto no produce los efectos propios del negocio, pero genera la obligación de otorgar la forma prescripta (conversión del acto jurídico, Art. 285). \\[6pt]

\textbf{¿Qué es la conversión del acto jurídico? (Art. 285)} &
Es el efecto propio de la solemnidad relativa: el acto no concluido en la forma requerida por la ley no queda concluido como tal, pero vale como acto por el cual las partes se obligan a cumplir con la formalidad, siempre que ésta no sea exigida bajo sanción de nulidad. Ejemplo: el boleto de compraventa genera la obligación de escriturar. \\[6pt]

\textbf{¿Cuáles son las categorías de instrumentos en el CCyCN? (Arts. 286--287)} &
(1) Instrumentos \textbf{públicos}: con autenticidad propia por ley. (2) Instrumentos \textbf{privados}: particulares firmados; requisito esencial: firma (Art. 288). (3) Instrumentos \textbf{particulares no firmados}: todo escrito, registro visual, auditivo o electrónico sin firma. \\[6pt]

\textbf{¿Qué prueba la firma? ¿Qué diferencia hay entre firma digital y firma electrónica? (Art. 288, Ley 25.506)} &
La firma prueba la \textbf{autoría} de la declaración de voluntad. La \textbf{firma digital} garantiza autoría e integridad mediante algoritmo matemático y certificado emitido por certificador licenciado. La \textbf{firma electrónica} es sólo un conjunto de datos identificatorios sin esa certificación; puede ser desconocida y su validez debe acreditarse. \\[6pt]

\textbf{¿Qué son los instrumentos públicos y cuáles son sus caracteres? (Art. 289)} &
Son los documentos a los que la ley les atribuye autenticidad propia, sin necesidad de reconocimiento previo. Hacen plena fe de su contenido. Enumeración: escrituras públicas y copias, instrumentos de escribanos o funcionarios públicos, y títulos emitidos por el Estado. \\[6pt]

\textbf{¿Cuáles son los requisitos de validez del instrumento público? (Art. 290)} &
(a) Actuación del oficial público dentro de sus \textbf{atribuciones} (competencia funcional) y \textbf{competencia territorial}. (b) \textbf{Firmas} del oficial, las partes y sus representantes. La falta de cualquier firma invalida el instrumento para todos. \\[6pt]

\textbf{¿Cómo funciona la eficacia probatoria del instrumento público? (Art. 296)} &
Distingue dos regímenes: (a) Los \textbf{hechos percibidos por el oficial} (fecha, lugar, actos cumplidos ante él) sólo pueden cuestionarse mediante \textbf{juicio de redargución de falsedad}. (b) El \textbf{contenido de las declaraciones de las partes} hace plena fe pero admite \textbf{prueba en contrario} sin necesidad de redargución. \\[6pt]

\textbf{¿Qué es la escritura pública y en qué se diferencia del acta notarial? (Arts. 299 y 310)} &
La \textbf{escritura pública} es el instrumento matriz extendido en el protocolo notarial que contiene \textbf{actos jurídicos}. El \textbf{acta notarial} es el documento notarial que tiene por objeto la comprobación de \textbf{hechos}. La distinción es esencial: la escritura crea, modifica o extingue derechos; el acta sólo constata hechos. \\[6pt]

\textbf{¿Qué estructura tiene la escritura pública? (Art. 305)} &
Tres partes: (1) \textbf{Encabezado}: lugar, fecha y datos de los comparecientes. (2) \textbf{Cuerpo}: naturaleza del acto, individualización de bienes y, en transmisiones de inmuebles, el ``corresponde''. (3) \textbf{Pie}: constancia de lectura, salvado de enmiendas, firmas. \\[6pt]

\textbf{¿Qué es el ``corresponde'' y para qué sirve?} &
Es la cláusula que, en las escrituras de transmisión de derechos reales sobre inmuebles, relaciona todos los antecedentes de titularidad del bien. Permite encadenar las sucesivas escrituras y verificar la historia dominial. Su omisión no invalida la escritura pero puede generar sanciones para el escribano. \\[6pt]

\textbf{¿Cuál es el requisito esencial del instrumento privado y cómo se reconoce la firma? (Arts. 288 y 314)} &
El requisito esencial es la \textbf{firma}. Para que el instrumento tenga plena fuerza probatoria entre las partes, la firma debe ser \textbf{reconocida}: expresa, tácita (el silencio frente a la intimación equivale a reconocimiento, Art. 314) o forzosa (pericia caligráfica). El reconocimiento de firma implica el reconocimiento de todo el contenido del instrumento. \\[6pt]

\textbf{¿Qué es la firma a ruego y cuál es la diferencia con la impresión digital? (Art. 313)} &
La \textbf{firma a ruego} es la firma de un tercero en nombre de quien no sabe o no puede firmar, en presencia de dos testigos. Equivale a la firma y otorga plena fuerza probatoria. La \textbf{impresión digital} sólo vale como \emph{principio de prueba por escrito} (Art. 314) y puede ser impugnada, por su menor seguridad. \\[6pt]

\textbf{¿Qué es la fecha cierta y para qué sirve? (Art. 317)} &
Es el momento a partir del cual existe certeza objetiva de que el instrumento ya estaba firmado o no pudo firmarse después. Desde la fecha cierta, el instrumento privado es \textbf{oponible a terceros}. El CCyCN abandona la enumeración taxativa del Código anterior: la fecha cierta puede acreditarse por cualquier medio, que el juez apreciará rigurosamente. \\[6pt]

\textbf{¿Cómo valora el juez un instrumento particular no firmado? (Art. 319)} &
Debe ponderar: congruencia entre lo narrado y lo ocurrido; precisión y claridad técnica del texto; usos y prácticas del tráfico; relaciones precedentes entre las partes; y confiabilidad de los soportes y procedimientos técnicos empleados. \\[6pt]

\textbf{¿Puede utilizarse la correspondencia como prueba en juicio? (Art. 318)} &
Sí: el destinatario puede presentarla como prueba contra el remitente. Sin embargo, si la correspondencia es \textbf{confidencial} —por manifestación expresa o por la naturaleza de la comunicación—, no puede utilizarse sin consentimiento del remitente. Esta protección tiene fundamento constitucional (inviolabilidad de la correspondencia). \\

\midrule

\multicolumn{2}{p{\textwidth}}{%
\textbf{Resumen:} La forma del acto jurídico es el modo en que la voluntad de las partes se exterioriza y produce efectos jurídicos. El CCyCN parte del principio de libertad de formas (Art. 284), pero impone formas determinadas para los actos más trascendentes, con distintas consecuencias según se trate de solemnidad absoluta (nulidad del acto) o relativa (conversión: obligación de otorgar la forma prescripta, Art. 285). En materia de instrumentos, el Código distingue tres grandes categorías: los \textbf{instrumentos públicos}, que gozan de autenticidad propia y hacen plena fe, con dos regímenes de impugnación según se trate de hechos percibidos por el oficial (redargución de falsedad) o de declaraciones de las partes (prueba en contrario); la \textbf{escritura pública}, instrumento matriz del protocolo notarial que contiene actos jurídicos y requiere una estructura formal precisa; y los \textbf{instrumentos particulares}, que incluyen los privados (requisito esencial: firma; fuerza probatoria plena tras el reconocimiento) y los particulares no firmados (fuerza probatoria sujeta a la apreciación judicial según las pautas del Art. 319). La \textbf{firma digital} (Ley 25.506) garantiza autoría e integridad en el mundo digital, con mayor valor probatorio que la firma electrónica. La \textbf{fecha cierta} (Art. 317) determina desde cuándo el instrumento privado puede oponerse a terceros, y puede acreditarse por cualquier medio apreciado rigurosamente por el juez. Estas reglas configuran un sistema orientado a dar seguridad jurídica al tráfico de bienes y servicios, equilibrando la certeza del formalismo con la agilidad que demanda la vida contemporánea.
}\\

\bottomrule
\end{tabularx}
\end{table}

% ============================================================

% ============================================================
\chapter{Clasificación de los Actos Jurídicos}
% =====================================================

En el régimen del Código de Vélez se enunciaba expresamente una serie de clasificaciones de los actos jurídicos --unilaterales y bilaterales, entre otras--. El Código Civil y Comercial de la Nación (CCyC) vigente, en cambio, optó metodológicamente por no incluir una clasificación sistematizada de los actos jurídicos: no existen artículos dedicados a definir de manera unificada las distintas clases de actos. Sin embargo, estas clasificaciones subsisten y se manifiestan a lo largo de todo el articulado, en la medida en que el Código atribuye consecuencias jurídicas distintas según el acto sea bilateral, entre vivos o de última voluntad, oneroso o gratuito, entre otras categorías.

\begin{note}
Estas clasificaciones no son puramente teóricas, sino que tienen consecuencias prácticas concretas y se encuentran presentes en distintos lugares del Código. A lo largo de las distintas asignaturas del plan de estudios de abogacía es habitual encontrar aplicaciones de estas categorías.
\end{note}

Su relevancia práctica es amplia: saber si un acto es formal o no formal determina si se requiere escritura pública; saber si es oneroso o gratuito modifica las reglas de buena fe aplicables (art. 392, CCyC); y saber si el acto está sujeto a condición, plazo o cargo indica en qué momento nace, se suspende o se extingue la obligación asumida. Estas categorías atraviesan los contratos, las sucesiones, el derecho de familia, las obligaciones y los derechos reales, y constituyen el marco conceptual necesario para interpretar el Código en su conjunto.

\begin{example}{Aplicación integradora de las clasificaciones}
En la redacción de una compraventa de un inmueble, dicho acto resulta simultáneamente bilateral, oneroso, formal (requiere escritura pública), patrimonial y de disposición. Si además las partes acuerdan que la entrega quede sujeta a que el banco apruebe el crédito del comprador, se configura una condición suspensiva; si se pacta que el contrato se resuelve en caso de que el banco no se expida dentro de los sesenta días, se configura un plazo resolutorio; y si el vendedor se compromete a pintar el departamento antes de la entrega, se configura un cargo. De este modo, las clasificaciones y las modalidades estudiadas en este capítulo convergen en un único documento cotidiano de la práctica profesional.
\end{example}

\section{Actos Unilaterales y Bilaterales}

El criterio de esta clasificación es el número de partes que concurren a la celebración del acto.

\begin{definition}{Actos unilaterales y bilaterales}
El acto es \textbf{unilateral} cuando una sola parte concurre a su celebración. Es \textbf{bilateral} cuando dos o más partes concurren a su celebración.
\end{definition}

\begin{table}[H]
\centering
\begin{tabularx}{\textwidth}{l X l}
\toprule
\textbf{Tipo de acto} & \textbf{Definición} & \textbf{Ejemplo típico} \\
\midrule
Unilateral & Una sola parte concurre a la celebración del acto. & El testamento \\
Bilateral & Dos o más partes concurren a la celebración del acto. & Un contrato; el matrimonio \\
\bottomrule
\end{tabularx}
\caption{Actos unilaterales y bilaterales}
\end{table}

\section{Actos entre Vivos y de Última Voluntad}

\begin{definition}{Actos entre vivos y de última voluntad}
Los actos \textbf{entre vivos} son otorgados por personas vivas para producir efectos en vida de los otorgantes. Los actos \textbf{de última voluntad} son aquellos que producen sus efectos luego de la muerte de quien los otorga.
\end{definition}

\begin{table}[H]
\centering
\begin{tabularx}{\textwidth}{l X X}
\toprule
\textbf{Tipo de acto} & \textbf{Definición} & \textbf{Cuándo produce efectos} \\
\midrule
Entre vivos & Son otorgados por personas vivas para producir efectos en vida. & Desde el momento de su celebración \\
De última voluntad & Producen sus efectos luego de la muerte del otorgante. & Desde la muerte del otorgante \\
\bottomrule
\end{tabularx}
\caption{Actos entre vivos y de última voluntad}
\end{table}

\section{Actos a Título Oneroso y a Título Gratuito}

Esta clasificación resulta especialmente relevante a lo largo de todo el Código, dado que el CCyC tiende a proteger de manera diferenciada a quienes adquieren a título oneroso y de buena fe frente a quienes adquieren a título gratuito.

\begin{definition}{Actos onerosos y gratuitos}
El acto es \textbf{oneroso} cuando existen prestaciones recíprocas: cada parte entrega algo a cambio de lo que recibe. El acto es \textbf{gratuito} (o a título de liberalidad) cuando una de las partes recibe una prestación sin dar nada a cambio, lo cual enriquece a esa parte.
\end{definition}

\begin{table}[H]
\centering
\begin{tabularx}{\textwidth}{l X X}
\toprule
\textbf{Tipo de acto} & \textbf{Definición} & \textbf{Ejemplo} \\
\midrule
Oneroso & Hay prestaciones recíprocas: cada parte da algo a cambio. & Compraventa \\
Gratuito (liberalidad) & Una de las partes recibe sin dar nada a cambio; enriquece a esa parte. & Donación; renuncia a un derecho \\
\bottomrule
\end{tabularx}
\caption{Actos a título oneroso y a título gratuito}
\end{table}

El caso típico del título gratuito es la donación, pero también puede configurarse a título gratuito la renuncia a un derecho: quien tiene derecho a cobrar y renuncia a ese cobro está beneficiando a la otra parte sin recibir nada a cambio.

\begin{formula}{Art. 392, CCyC --- Efectos respecto de terceros en cosas registrables}
El Código diferencia la protección al subadquirente según haya adquirido a título oneroso o gratuito. El subadquirente de buena fe y a título oneroso goza de mayor protección.
\end{formula}

\section{Actos Principales y Accesorios}

\begin{definition}{Actos principales y accesorios}
Los actos \textbf{principales} son aquellos que no dependen, para su existencia, de otro acto. Los actos \textbf{accesorios} son aquellos que siguen la suerte del acto principal, del cual dependen para existir.
\end{definition}

\begin{table}[H]
\centering
\begin{tabularx}{\textwidth}{l X X}
\toprule
\textbf{Tipo de acto} & \textbf{Definición} & \textbf{Ejemplo} \\
\midrule
Principal & No depende de otro acto para existir. & Contrato de alquiler \\
Accesorio & Sigue la suerte del principal; depende de él para existir. & La garantía del contrato de alquiler \\
\bottomrule
\end{tabularx}
\caption{Actos principales y accesorios}
\end{table}

A modo de ejemplo, una garantía que acompaña a un contrato de alquiler constituye un acto accesorio respecto del principal y sigue su suerte, sin perjuicio de los matices particulares que se estudian en cada contrato en concreto.

\section{Actos Formales y No Formales}

Esta es una clasificación de particular importancia, que se desarrolla en profundidad en el capítulo~\ref{ch:forma}.

\begin{definition}{Actos formales y no formales}
Los actos \textbf{formales} son aquellos en los que la ley prescribe cierta solemnidad. Dentro de estos, la solemnidad puede hacer a la validez del acto (\textbf{formal solemne} o absolutamente formal) o hacer únicamente a su prueba (\textbf{formal no solemne} o relativamente formal). Los actos \textbf{no formales} son aquellos respecto de los cuales la ley no exige una forma determinada, quedando las partes en libertad de elegir la que prefieran.
\end{definition}

\begin{table}[H]
\centering
\begin{tabularx}{\textwidth}{l X X}
\toprule
\textbf{Tipo de acto} & \textbf{Definición} & \textbf{Consecuencia del incumplimiento de la forma} \\
\midrule
Formal solemne (absolutamente formal) & La ley prescribe una solemnidad cuyo incumplimiento invalida el acto. & Nulidad del acto \\
Formal no solemne (relativamente formal) & La ley prescribe una forma, pero su incumplimiento afecta solo la prueba, no la validez. & Puede probarse por otro medio \\
No formal & La ley no manda ninguna forma determinada; las partes pueden elegir la que prefieran. & No hay consecuencia: el acto es válido \\
\bottomrule
\end{tabularx}
\caption{Actos formales y no formales}
\end{table}

\begin{example}{Acto formal solemne}
El matrimonio es un acto formal solemne: se encuentra rodeado de una serie de solemnidades marcadas por el Código Civil y no puede celebrarse sino de acuerdo con los rituales que este establece. Cualquier otra forma determina que el matrimonio no sea válido.
\end{example}

\begin{example}{Acto no formal}
El alquiler de un inmueble es, en principio, un acto no formal: la ley no exige que se celebre por escritura pública. No obstante, las partes pueden optar por una forma más solemne que la exigida legalmente --por ejemplo, contratando a un escribano-- por razones de seguridad jurídica, sin que ello sea un requisito legal.
\end{example}

\section{Actos Patrimoniales y Extrapatrimoniales}

\begin{definition}{Actos patrimoniales y extrapatrimoniales}
Los actos \textbf{patrimoniales} tienen contenido económico y están vinculados con el patrimonio. Los actos \textbf{extrapatrimoniales} no tienen contenido económico directo y se vinculan con derechos personalísimos o con el derecho de familia.
\end{definition}

\begin{table}[H]
\centering
\begin{tabularx}{\textwidth}{l X X}
\toprule
\textbf{Tipo de acto} & \textbf{Definición} & \textbf{Ejemplo} \\
\midrule
Patrimonial & Tiene contenido económico; está vinculado con el patrimonio. & Un contrato de compraventa \\
Extrapatrimonial & No tiene contenido económico directo; vinculado con derechos personalísimos o de familia. & El reconocimiento de un hijo \\
\bottomrule
\end{tabularx}
\caption{Actos patrimoniales y extrapatrimoniales}
\end{table}

El reconocimiento de un hijo es un acto jurídico unilateral mediante el cual una persona expresa su voluntad de reconocerlo. Se trata de un acto típico del derecho de familia, extrapatrimonial en sí mismo, aun cuando pueda tener consecuencias patrimoniales indirectas, como el derecho a alimentos. Los actos vinculados con derechos personalísimos también son extrapatrimoniales.

\section{Actos de Disposición, de Administración y de Conservación}

\begin{definition}{Actos de disposición, administración y conservación}
Los actos de \textbf{disposición} producen una variación sustancial en la composición del patrimonio. Los actos de \textbf{administración} no producen un cambio sustancial en el patrimonio. Los actos de \textbf{conservación} están dedicados al mantenimiento de los bienes.
\end{definition}

\begin{table}[H]
\centering
\begin{tabularx}{\textwidth}{l X X}
\toprule
\textbf{Tipo de acto} & \textbf{Definición} & \textbf{Ejemplo} \\
\midrule
Disposición & Se produce una variación sustancial en la composición del patrimonio. & Venta de un inmueble: sale el bien, entra dinero \\
Administración & No produce un cambio sustancial en el patrimonio. & Cobrar un alquiler, recibir una renta \\
Conservación & Dedicados al mantenimiento de los bienes. & Reparaciones, actos de preservación \\
\bottomrule
\end{tabularx}
\caption{Actos de disposición, administración y conservación}
\end{table}

\begin{note}
Esta clasificación resulta especialmente relevante en el derecho de familia y de las personas, en materia de tutela, curatela y régimen de poderes. El representante legal de una persona incapaz, por ejemplo, necesita autorización judicial para los actos de disposición, pero no para los de administración ordinaria.
\end{note}

\section{Actos Puros y Simples, y Actos Modales}

\begin{definition}{Actos puros y simples, y actos modales}
Los actos \textbf{puros y simples} no están sujetos a ninguna modalidad y producen sus efectos de inmediato. Los actos \textbf{modales} son aquellos en los cuales las partes eligen sujetarlo a alguno de los tres modos posibles: condición, plazo o cargo.
\end{definition}

\begin{table}[H]
\centering
\begin{tabularx}{\textwidth}{l X}
\toprule
\textbf{Tipo de acto} & \textbf{Definición} \\
\midrule
Puro y simple & No está sujeto a ninguna modalidad. Produce efectos de inmediato. \\
Modal & Las partes eligen sujetarlo a alguno de los tres modos: condición, plazo o cargo. \\
\bottomrule
\end{tabularx}
\caption{Actos puros y simples, y actos modales}
\end{table}

\section{Actos Causales y Abstractos}

La noción de acto abstracto se desarrolla también al estudiar el régimen de la causa (capítulo~\ref{ch:causa}).


\begin{definition}{Actos causales y abstractos}
El acto \textbf{causal} permanece vinculado a su causa, la cual influye en su validez y en sus efectos. El acto \textbf{abstracto} se independiza de su causa.
\end{definition}

\begin{table}[H]
\centering
\begin{tabularx}{\textwidth}{l X}
\toprule
\textbf{Tipo de acto} & \textbf{Definición} \\
\midrule
Causal & El acto permanece vinculado a su causa; la causa influye en su validez y efectos. \\
Abstracto & El acto se independiza de su causa. \\
\bottomrule
\end{tabularx}
\caption{Actos causales y abstractos}
\end{table}

% TODO: contenido incompleto o ambiguo en las notas originales -- el apunte remite a un desarrollo previo de los actos abstractos ("se estudiaron previamente en la clase") que no forma parte de este documento.

\section{Actos Positivos y Negativos}

\begin{definition}{Actos positivos y negativos}
El acto \textbf{positivo} involucra una acción, es decir, un hacer. El acto \textbf{negativo} involucra una omisión, es decir, un no hacer.
\end{definition}

\begin{table}[H]
\centering
\begin{tabularx}{\textwidth}{l X X}
\toprule
\textbf{Tipo de acto} & \textbf{Definición} & \textbf{Ejemplo} \\
\midrule
Positivo & Involucra una acción (un hacer). & Entregar una cosa, pagar una suma \\
Negativo & Involucra una omisión (un no hacer). & Comprometerse a no competir en un rubro comercial \\
\bottomrule
\end{tabularx}
\caption{Actos positivos y negativos}
\end{table}

\begin{example}{Acto negativo: cláusula de no competencia}
Es habitual que quienes se dedican a una determinada actividad --por ejemplo, la producción de helados-- al vender la marca y la empresa se comprometan a no generar una nueva empresa ni emprendimientos vinculados con ese rubro durante los próximos cinco años. Este acto jurídico se encuadra dentro de los actos negativos, en tanto las partes se obligan a un ``no hacer'', es decir, a una omisión.
\end{example}

% =====================================================
\chapter{Modalidades del acto jurídico}
\label{ch:modalidades}
% =====================================================

Las modalidades del acto jurídico son la condición, el plazo y el cargo. En el régimen del Código de Vélez, las modalidades estaban reguladas dentro de la sección de obligaciones del Libro Segundo. En el Código Civil y Comercial vigente, en cambio, se las incorporó al Libro Primero, Título Cuarto, dentro del capítulo de hechos y actos jurídicos, en un capítulo específico dedicado a las modalidades del acto jurídico.

\section{La Condición}

\begin{formula}{Art. 343, CCyC --- Alcances}
Se denomina condición a la cláusula de los actos jurídicos por la cual las partes subordinan su plena eficacia o resolución a un hecho futuro e incierto.
\end{formula}

\begin{definition}{Condición}
La condición es una cláusula por la cual se subordina la adquisición o la resolución de un acto a un hecho futuro e incierto. Lo esencial de la noción es la incertidumbre: el hecho condicionante puede o no llegar a suceder.
\end{definition}

\begin{example}{Condición suspensiva (1)}
``Te voy a regalar un auto si Argentina sale campeón del mundo en el próximo mundial.'' Se trata de un hecho futuro e incierto: puede ocurrir o no. El efecto --la donación del auto-- queda suspendido hasta que ocurra, o no, el hecho condicionante.
\end{example}

\begin{example}{Condición suspensiva (2)}
``Te voy a regalar mi biblioteca jurídica si te recibís de abogado/a.'' El hecho es incierto: puede ocurrir o no. La obligación existe desde la constitución del acto, pero no es exigible --no tiene plena eficacia-- hasta el momento en que ocurre el hecho condicionante.
\end{example}

\begin{example}{Condición resolutoria}
``Te dono mi auto, pero siempre y cuando mi hermano no vuelva de España. Si mi hermano vuelve, tenés que devolverme el auto.'' Aquí el acto ya produce efectos desde el principio; el hecho condicionante produce la resolución, es decir, la devolución de los efectos hacia atrás.
\end{example}

\begin{table}[H]
\centering
\begin{tabularx}{\textwidth}{l X X}
\toprule
\textbf{Tipo de condición} & \textbf{Punto de partida} & \textbf{Efecto del hecho condicionante} \\
\midrule
Suspensiva & El acto existe pero NO produce efectos todavía. & Si ocurre el hecho, el acto produce sus efectos; si no ocurre, el acto no produce efectos. \\
Resolutoria & El acto YA produce efectos desde su celebración. & Si ocurre el hecho, los efectos cesan (se resuelven), con retroactividad. \\
\bottomrule
\end{tabularx}
\caption{Diferencia entre condición suspensiva y resolutoria}
\end{table}

\subsection{La Condición Impropia o Suposición}

El nuevo Código incorporó también la llamada ``condición impropia'' o ``suposición'': aquella en la que las partes condicionan el acto a un hecho presente o pasado que, al momento del acto, es ignorado por ellas.

\begin{formula}{Art. 343, CCyC (segundo párrafo)}
Se equipara a la condición la cláusula por la cual las partes sujetan la eficacia del acto a la existencia de cierto estado de cosas al tiempo de celebrarlo.
\end{formula}

\begin{example}{Condición impropia o suposición}
Una persona ignora la situación patrimonial de un familiar que vive lejos y, sin saberlo, celebra un acto sujeto a que ese familiar no haya caído en quiebra: por ejemplo, ``te presto la plata en la medida en que tal familiar mío no haya terminado en quiebra''. Aunque el hecho ya ocurrió o no ocurrió al momento de celebrar el acto, la incertidumbre subjetiva de las partes --que lo ignoran-- hace que este supuesto también valga como condición.
\end{example}

\subsection{Clasificación de las Condiciones}

\begin{table}[H]
\centering
\begin{tabularx}{\textwidth}{l l X}
\toprule
\textbf{Criterio} & \textbf{Tipo} & \textbf{Descripción} \\
\midrule
Según el hecho & Positiva & Implica una acción (que algo ocurra). \\
Según el hecho & Negativa & Implica una omisión (que algo no ocurra). \\
Según el efecto & Suspensiva & Subordina la eficacia del acto a un hecho futuro e incierto. \\
Según el efecto & Resolutoria & Subordina la extinción del acto a un hecho futuro e incierto. \\
Según la fuente & Causal & El hecho no depende de las partes; depende del azar o de un tercero. \\
Según la fuente & Potestativa & Depende de la voluntad del deudor o del acreedor. \\
Según la fuente & Mixta & Depende en parte del deudor y en parte de un tercero o del azar. \\
\bottomrule
\end{tabularx}
\caption{Clasificación de las condiciones}
\end{table}

La condición \textbf{causal} no depende de las partes, sino del azar o de un tercero --por ejemplo, que la selección nacional resulte campeona--. La condición \textbf{potestativa} depende del acreedor o del deudor --por ejemplo, ``te voy a regalar mi biblioteca si te recibís de abogado/a'' es potestativa del acreedor--. La condición \textbf{mixta} depende en parte del deudor y en parte de un hecho ajeno a él.

\begin{important}
La condición puramente potestativa del deudor está prohibida: ``te comprometo a entregarte algo si yo decido hacerlo'' no puede depender completamente de la voluntad del obligado.
\end{important}

\subsection{Condiciones Prohibidas}

\begin{formula}{Art. 344, CCyC --- Condiciones prohibidas}
Es nula la condición que: (1) sea un hecho imposible, (2) sea contraria a las buenas costumbres, (3) esté prohibida por el ordenamiento jurídico, (4) dependa exclusivamente de la voluntad del obligado (puramente potestativa). La condición de no hacer algo imposible no perjudica la validez de la obligación.
\end{formula}

\begin{formula}{Art. 345, CCyC --- Condiciones que gravan libertades fundamentales}
Las condiciones cuyo cumplimiento depende de que el beneficiario cambie de domicilio, religión u otras libertades fundamentales se tienen por no escritas.
\end{formula}

\begin{example}{Condición que grava una libertad fundamental}
``Te dono la casa si cambiás de religión'' constituye una condición prohibida, en tanto grava la libertad fundamental de profesar la religión que la persona elija.
\end{example}

\begin{table}[H]
\centering
\begin{tabularx}{\textwidth}{X X}
\toprule
\textbf{Tipo de condición prohibida} & \textbf{Consecuencia jurídica} \\
\midrule
Hecho imposible / contraria a buenas costumbres / puramente potestativa (art. 344, párr. 1) & Nulidad absoluta del acto \\
Que grava libertades fundamentales (art. 345, párr. 3) & Se tiene por no escrita (la obligación principal podría subsistir) \\
Imposible en testamentos (excepción) & La condición se tiene por no puesta; la disposición vale \\
\bottomrule
\end{tabularx}
\caption{Condiciones prohibidas y sus consecuencias}
\end{table}

Existe una diferencia entre el primer y el tercer párrafo del artículo 344: las condiciones ilícitas del primer párrafo producen la nulidad absoluta del acto, mientras que las condiciones comprendidas en el tercer párrafo se tienen por no escritas, discutiéndose en la doctrina si la obligación principal subsiste o no. La excepción se presenta en los testamentos: las condiciones prohibidas incorporadas en un testamento no afectan la validez de las disposiciones testamentarias de las que forman parte.

\section{El Plazo}

A diferencia de la condición, el plazo es la modalidad que posterga o limita temporalmente los efectos del acto. En el plazo hay un hecho futuro pero \textbf{cierto}: se sabe que va a ocurrir, aunque no siempre se sepa exactamente cuándo.

\begin{example}{Plazo suspensivo y plazo resolutorio}
Plazo suspensivo: ``Te voy a donar estas casas dentro de un año.'' No hay incertidumbre sobre si llegará el día, aunque pueda haberla sobre el momento exacto en el caso del plazo incierto. Plazo resolutorio: ``Te presto el uso de la casa por un año.'' Los efectos comienzan de inmediato, pero cesan cuando vence el plazo.
\end{example}

\begin{table}[H]
\centering
\begin{tabularx}{\textwidth}{l X X}
\toprule
\textbf{Tipo de plazo} & \textbf{Definición} & \textbf{Ejemplo} \\
\midrule
Suspensivo & Posterga los efectos del acto hasta el vencimiento del plazo. & ``Te dono la casa dentro de un año.'' \\
Resolutorio & Los efectos del acto cesan al vencimiento del plazo. & ``Te presto la casa por un año.'' \\
Cierto & Se conoce la fecha exacta del vencimiento. & ``En 30 días'', ``el 31 de diciembre'' \\
Incierto & No se sabe exactamente cuándo vencerá, pero es seguro que ocurrirá. & ``Cuando muera tal persona'' \\
Determinado & Las partes, la ley o el juez fijaron la fecha o el momento. & La ley o el contrato establece una fecha \\
Indeterminado & No se precisó fecha, pero existe un plazo posible. & Cláusula ``a mejor fortuna'' (art. 889, CCyC) \\
Esencial & El cumplimiento fuera de término no satisface el interés del acreedor. & Catering para una fiesta: si llega dos semanas tarde, ya no sirve \\
Accidental & El cumplimiento tardío puede igualmente satisfacer el interés del acreedor. & Entrega de materiales de construcción con demora \\
\bottomrule
\end{tabularx}
\caption{Clasificación del plazo}
\end{table}

\begin{example}{Plazo esencial}
Una persona contrata un servicio de catering para una fiesta y, por un error, el servicio se entrega dos semanas después del evento. El plazo era esencial --la fiesta misma-- por lo que la prestación tardía no satisface el interés del acreedor y el incumplimiento es total.
\end{example}

El plazo puede surgir de la convención de las partes, de la ley, puede ser judicial, o bien indeterminado, cuando se establece la fecha de inicio pero no la de finalización.

\begin{formula}{Art. 889, CCyC --- Cláusula ``a mejor fortuna''}
El Código admite la cláusula ``a mejor fortuna'': ``yo te presto plata y me devolvés cuando estés bien económicamente.'' Se trata de un plazo indeterminado. El juez, a solicitud del acreedor, puede fijar un plazo y hacer cumplir la obligación.
\end{formula}

\section{El Cargo}

\begin{definition}{Cargo}
El cargo es una obligación accesoria impuesta al adquirente de un derecho, que no impide los efectos del acto pero que, sin embargo, puede exigirse.
\end{definition}

\begin{example}{Cargo simple}
Una persona dona un campo a otra con el cargo de que mantenga en la casa a los caseros. Si el donatario no cumple, los caseros podrán exigir el cumplimiento del cargo o reclamar la indemnización de los daños y perjuicios ocasionados por el incumplimiento; sin embargo, en principio, el acto de donación no se cae.
\end{example}

\begin{example}{Cargo con finalidad benéfica}
Se dona a una fundación un terreno con el cargo de que el producto de la explotación económica de ese terreno se destine a obras para los barrios vulnerables. La fundación que recibe la donación debe administrar el bien conforme a ese cargo.
\end{example}

\begin{table}[H]
\centering
\begin{tabularx}{\textwidth}{X X X}
\toprule
\textbf{Característica} & \textbf{Condición} & \textbf{Cargo} \\
\midrule
Suspende efectos del acto & SÍ (si es suspensiva) & NO \\
Coercible (exigible) & Cuando se cumple el hecho: sí & Sí, siempre \\
Si falla en cumplirlo, revoca el acto & Sí (si es resolutoria) & Solo si fue impuesto como cargo condicional \\
Quién puede exigirlo & Las partes & Los beneficiarios y quien constituyó el derecho \\
\bottomrule
\end{tabularx}
\caption{Diferencias entre condición y cargo}
\end{table}

El cargo es coercitivo pero no suspensivo, a diferencia de la condición suspensiva, que sí suspende los efectos. El cargo siempre se impone al adquirente como una obligación que debe cumplir y, en tanto obligación, debe reunir los requisitos propios de las obligaciones: ser lícito, posible, determinado y de contenido patrimonial. Surge de la autonomía de la voluntad, es excepcional y accidental --puede o no estar presente--, constituye una obligación accesoria del derecho transmitido, debe ser expreso y es restrictivo del derecho adquirido, en la medida en que recibir algo con el cargo de cumplir una obligación restringe el derecho de quien lo recibe.

\subsection{Cargo Simple y Cargo Condicional}

\begin{table}[H]
\centering
\begin{tabularx}{\textwidth}{l X X}
\toprule
\textbf{Tipo de cargo} & \textbf{Definición} & \textbf{Consecuencia del incumplimiento} \\
\midrule
Cargo simple & Obligación accesoria que no condiciona la adquisición del derecho. & Se exige el cumplimiento o se reclaman daños. El acto no se revoca. \\
Cargo condicional & El cumplimiento del cargo fue puesto como condición del acto. & Si no se cumple, se puede revertir la adquisición del derecho. Opera como condición resolutoria. \\
\bottomrule
\end{tabularx}
\caption{Cargo simple y cargo condicional}
\end{table}

El cargo simple no tiene efecto resolutorio sobre el derecho principal. El cargo condicional, en cambio, opera como una condición resolutoria (o suspensiva), de modo que su incumplimiento afecta a la propia adquisición del derecho: si el cargo fue puesto como condicional, la adquisición puede revertirse en caso de incumplimiento.

\subsection{Cargos Prohibidos}

\begin{formula}{Art. 357, CCyC --- Cargos prohibidos}
Son de aplicación a los cargos las normas establecidas para las condiciones prohibidas (art. 344, CCyC). Sin embargo, cuando se impone un cargo prohibido, no se produce la nulidad del acto, sino que el cargo se tiene por no escrito.
\end{formula}

Los cargos prohibidos son, en su contenido, los mismos supuestos previstos para las condiciones prohibidas: el artículo 357 remite al artículo 344. La diferencia radica en la consecuencia: cuando se impone un cargo prohibido, este se tiene por no escrito, sin que ello acarree la nulidad del acto principal.

% =====================================================

\section{Cierre y Síntesis}
% =====================================================

Al finalizar el desarrollo de las clasificaciones y las modalidades del acto jurídico --condición, plazo y cargo-- corresponde una aclaración fundamental sobre la naturaleza de estas últimas.

\begin{important}
Las modalidades --condición, plazo y cargo-- \textbf{no} son una clasificación de los actos jurídicos. Son modos, cláusulas, elementos accidentales de un acto jurídico: pueden o no estar presentes. Si las partes las pactan, rigen las reglas de su regulación civil; si no las incorporan, el acto es puro y simple y produce efectos de inmediato.
\end{important}

Los actos jurídicos son manifestaciones de voluntad orientadas a producir efectos jurídicos. El Código Civil y Comercial no los clasifica en un artículo único, sino que sus categorías atraviesan todo el texto legal. Estas clasificaciones --unilateral/bilateral, entre vivos/última voluntad, oneroso/gratuito, principal/accesorio, formal/no formal, patrimonial/extrapatrimonial, disposición/administración/conservación, puro/modal, causal/abstracto, positivo/negativo-- no son puramente teóricas: determinan las reglas aplicables en cada situación jurídica concreta. Sobre esa base, el acto puede quedar sujeto a modalidades accidentales: la condición (hecho futuro e incierto que suspende o resuelve los efectos), el plazo (hecho futuro cierto que posterga o limita los efectos en el tiempo) y el cargo (obligación accesoria que no suspende el acto, pero que es coercitiva). Cada modalidad presenta variantes propias, reglas de validez específicas y artículos particulares dentro del CCyC (arts. 343 a 357).

% =====================================================

\section{Cuadro Cornell de Repaso}
% =====================================================

El siguiente cuadro reúne, según el método Cornell, las preguntas centrales que permiten repasar activamente el contenido desarrollado en los capítulos anteriores, junto con sus respuestas y una síntesis final integradora.

\begin{longtable}{
    >{\raggedright\arraybackslash}p{0.30\textwidth}
    >{\raggedright\arraybackslash}p{0.62\textwidth}
}
\toprule
\textbf{Preguntas / palabras clave} & \textbf{Notas / respuestas} \\
\midrule
\endfirsthead
\toprule
\textbf{Preguntas / palabras clave} & \textbf{Notas / respuestas} \\
\midrule
\endhead

\textbf{¿Qué son las clasificaciones de los actos jurídicos y para qué sirven?} &
Son categorías doctrinarias --y en gran parte también legales-- que agrupan los actos según criterios como el número de partes, el momento de producción de efectos, la existencia de prestaciones recíprocas, la dependencia respecto de otro acto, la forma, el contenido patrimonial o el nivel de variación del patrimonio. No son puramente teóricas: determinan las reglas aplicables al acto en cada situación concreta del Código y de la práctica profesional. \\

\textbf{Diferencia entre acto unilateral y bilateral. Ejemplos.} &
Unilateral: una sola parte (el testamento). Bilateral: dos o más partes (el contrato, el matrimonio). Esta diferencia resulta clave porque el régimen de perfeccionamiento, revocación y efectos es distinto en cada caso. \\

\textbf{¿Por qué importa si un acto es oneroso o gratuito?} &
El CCyC protege de manera distinta a los subadquirentes según hayan adquirido a título oneroso o gratuito. Conforme el art. 392, CCyC, quien adquiere de buena fe y a título oneroso goza de mayor protección frente a la evicción y otras acciones de terceros que quien adquirió gratuitamente. \\

\textbf{Acto formal solemne vs.\ formal no solemne vs.\ no formal.} &
Formal solemne: la forma es requerida para la validez; si no se cumple, el acto es nulo (testamento, matrimonio). Formal no solemne: la forma prescrita afecta solo la prueba; el acto igual es válido. No formal: la ley no exige ninguna forma; las partes eligen cómo exteriorizar su voluntad. \\

\textbf{¿Qué es la condición? Definición legal.} &
Cláusula por la cual se subordina la adquisición o la resolución de un acto a un hecho futuro e incierto (art. 343, CCyC). Lo esencial es la incertidumbre: el hecho puede o no ocurrir. Si el hecho es cierto, se trata de un plazo, no de una condición. \\

\textbf{Condición suspensiva vs.\ resolutoria.} &
Suspensiva: los efectos del acto no se producen hasta que ocurra el hecho condicionante; el acto existe, pero no es exigible. Resolutoria: el acto ya produce efectos; si ocurre el hecho, los efectos cesan y se resuelven retroactivamente. Ejemplo de suspensiva: ``te dono la biblioteca si te recibís''. Ejemplo de resolutoria: ``te dono el auto, pero si vuelve mi hermano, me lo devolvés''. \\

\textbf{¿Qué condiciones están prohibidas? Consecuencias.} &
Son nulas las condiciones que: (1) son imposibles, (2) son contrarias a las buenas costumbres, (3) son puramente potestativas del deudor (art. 344, CCyC). Se tienen por no escritas las que gravan libertades fundamentales, como la religión o el domicilio (art. 345). En los testamentos, las condiciones prohibidas no afectan la validez de la disposición testamentaria. \\

\textbf{Diferencia entre condición y plazo.} &
En ambos casos el hecho es futuro. La diferencia clave es la certeza: en la condición el hecho es incierto (puede no ocurrir); en el plazo el hecho es cierto (se sabe que ocurrirá, aunque no siempre se sepa exactamente cuándo, como en el plazo incierto: ``cuando muera tal persona''). \\

\textbf{Tipos de plazo según su certeza y determinación.} &
Cierto: se conoce la fecha exacta. Incierto: no se sabe cuándo, pero es seguro que ocurrirá. Determinado: fijado por las partes, la ley o el juez. Indeterminado: sin fecha precisada (por ejemplo, la cláusula ``a mejor fortuna'', art. 889, CCyC). Esencial: el cumplimiento tardío no satisface el interés del acreedor (catering para una fiesta). Accidental: el cumplimiento tardío igualmente puede satisfacer el interés del acreedor. \\

\textbf{¿Qué es el cargo? ¿Suspende los efectos del acto?} &
El cargo es una obligación accesoria impuesta al adquirente de un derecho que NO suspende los efectos del acto: el acto produce efectos de inmediato. Sin embargo, el cargo es coercitivo: puede exigirse su cumplimiento. Si fue impuesto como condición (cargo condicional), su incumplimiento puede revertir la adquisición del derecho. Los cargos prohibidos se tienen por no escritos (art. 357, que remite al art. 344, CCyC), sin que ello invalide el acto principal. \\

\textbf{¿Son las modalidades una clasificación de los actos jurídicos?} &
No. Las modalidades (condición, plazo, cargo) son elementos accidentales del acto jurídico: pueden o no estar presentes. Si las partes las incorporan, el acto queda sujeto a esas reglas; si no las incorporan, el acto es puro y simple y produce efectos de inmediato, sin condicionamiento alguno. \\

\midrule
\multicolumn{2}{p{\dimexpr\textwidth-2\tabcolsep}}{
\textbf{Resumen:} Los actos jurídicos admiten múltiples clasificaciones --según el número de partes, el momento de sus efectos, la existencia de contraprestación, su dependencia de otro acto, su forma, su contenido patrimonial y su incidencia en el patrimonio-- que determinan las reglas aplicables en cada caso concreto del Código Civil y Comercial. Además de estas clasificaciones, el acto jurídico puede quedar sujeto a modalidades accidentales: la condición (hecho futuro e incierto), el plazo (hecho futuro cierto) y el cargo (obligación accesoria y coercitiva que no suspende el acto). Cada modalidad tiene reglas propias de validez, variantes específicas y condiciones o cargos expresamente prohibidos por la ley (arts. 343 a 357, CCyC). Comprender estas categorías es la base del razonamiento jurídico: permiten leer el Código, redactar documentos, asesorar a clientes y resolver casos con precisión técnica en cualquier rama del derecho privado.
} \\
\bottomrule
\end{longtable}
\chapter[Principios generales]{Principios generales de los efectos del acto jurídico}

\section{Introducción: los efectos del acto jurídico}

El estudio de los efectos del acto jurídico se organiza como un sistema concéntrico. El núcleo es el principio \textbf{\textit{nemo plus iuris}} (nadie transmite más de lo que tiene), que luego se proyecta hacia las partes, sus representantes, los sucesores y, finalmente, los terceros. Puede pensarse como las ondas que produce una piedra en el agua: el acto jurídico es la piedra, y cada onda alcanza a un grupo diferente de personas con reglas propias.

\begin{definition}{Acto jurídico}
Acto \textbf{voluntario lícito} que tiene por fin inmediato producir el \textbf{nacimiento, la modificación, la transferencia o la extinción} de relaciones e instituciones jurídicas.
\end{definition}

Por definición, entonces, el acto jurídico busca producir efectos jurídicos. La cuestión central consiste en determinar \textbf{a quiénes abarcan esos efectos}. Los interrogantes que estructuran el tema son los siguientes:

\begin{itemize}
    \item ¿Los efectos abarcan a las partes?
    \item ¿Cómo están abarcados los sucesores?
    \item ¿Cómo están abarcados los acreedores?
    \item ¿Quiénes son representantes?
    \item ¿Qué reglas rigen estos efectos de los actos jurídicos?
\end{itemize}

\subsection{Importancia práctica}

Estos conceptos constituyen el fundamento de cada operación jurídica que realiza un abogado o una abogada. Se proyectan, por ejemplo, en las siguientes situaciones:

\begin{itemize}
    \item \textbf{Redacción de un contrato de compraventa:} es necesario saber si el cliente puede transmitir válidamente el derecho que pretende vender (\textit{nemo plus iuris}).
    \item \textbf{Otorgamiento o recepción de un poder notarial:} debe determinarse si el representante actuó dentro de sus facultades.
    \item \textbf{Asesoramiento a herederos:} se distingue qué obligaciones del causante los alcanzan y cuáles se extinguen.
    \item \textbf{Cobro de un crédito impago:} se elige entre la acción directa —para cobrar directamente— y la acción subrogatoria —para ingresar bienes al patrimonio del deudor—.
\end{itemize}

No se trata de temas abstractos: son herramientas cotidianas de la abogacía en materia contractual, sucesoria y ejecutiva.

\section{Nemo plus iuris (Art.~399 CCCN)}

\subsection{La transmisibilidad de los derechos como regla general}

Los derechos son \textbf{transmisibles}, salvo estipulación válida de las partes o prohibición legal. Sobre esta base se establece la regla central que organiza todo el tema: \textbf{nadie puede transmitir un derecho mejor que el que tiene}.

Esta regla ya se encontraba en el Código anterior, en el art.~3270: nadie puede transmitir a otro un derecho mejor o más extenso que el que tiene, salvo las excepciones legalmente dispuestas. Proviene del Derecho Romano y en algunos libros se la designa como \textit{nemo plus iuris}, por sus primeras palabras. Ya estaba presente en el Digesto.

\begin{formula}{Art.~399 CCCN --- Regla general (\textit{nemo plus iuris})}
Nadie puede transmitir a otro un derecho mejor o más extenso que el que tiene, sin perjuicio de las excepciones legalmente dispuestas.
\end{formula}

% TODO: contenido incompleto o ambiguo en las notas originales
% (las notas indican "en 398 estamos en el Título Quinto del Libro Primero" al ubicar el art. 399; se conserva la ubicación por Título y Libro, sin precisar el número de artículo)
\begin{note}
El \textit{nemo plus iuris} se encuentra en el art.~399, hacia el final del Libro Primero del Código Civil y Comercial de la Nación, dentro del Título Quinto, denominado «Transmisión de los derechos». Esto se relaciona con el estudio de los efectos del acto: quien transmite derechos no puede transmitirlos de un modo mejor que el que tiene.
\end{note}

\begin{example}{el inquilino que no puede vender}
Quien es inquilino de un inmueble no puede venderlo, porque no puede transmitir un derecho mejor que el que tiene. Su derecho sobre el inmueble es simplemente el uso y goce en los términos de un contrato de locación, y no el derecho en los términos del dominio del dueño de la cosa. Esto constituye claramente un \textbf{límite a la transmisibilidad de los derechos}.
\end{example}

\subsection{Caso práctico: Juan, Verónica y Daniela}

El principio se ilustra con un caso que involucra a tres personas.

\begin{example}{transmisión de propiedad viciada por intimidación}
Juan ejerce violencia (intimidación) sobre Verónica y logra que ella le transfiera su propiedad bajo amenazas. Juan aparece como propietario, pero su título está viciado por intimidación. Luego, Juan le vende el inmueble a Daniela.

\smallskip
\textbf{Pregunta:} ¿Daniela recibió la propiedad con el vicio o sin él?

\smallskip
Conforme al \textit{nemo plus iuris}, Juan transmite el derecho con el mismo alcance que lo tiene, es decir, con el vicio. Por lo tanto, Daniela, como sucesora, recibiría un derecho afectado por ese vicio de intimidación.
\end{example}

\section[Excepciones al nemo plus iuris]{Excepciones al \textit{nemo plus iuris}}

El mismo caso plantea una \textbf{excepción crucial} al principio general: es necesario determinar si Daniela es o no de buena fe. Pueden presentarse, así, excepciones al \textit{nemo plus iuris}.

\subsection{Adquirente de buena fe y a título oneroso (Art.~392 CCCN)}

Una de esas excepciones se encuentra en el art.~392. Si Daniela es de \textbf{buena fe} y a \textbf{título oneroso}, adquiere la propiedad, y quien sale perdiendo es Verónica —quien transmitió bajo intimidación—, porque Daniela se convierte en adquirente de buena fe a título oneroso. \textbf{Quien es de buena fe queda protegido.}

\begin{formula}{Art.~392 CCCN --- Efectos respecto de terceros en casos de nulidad}
El subadquirente de buena fe y a título oneroso queda a salvo de cualquier acción de nulidad que afecte el antecedente de su título, así como de los actos de los que fueran parte terceros de buena fe, con anterioridad a la inscripción de la sentencia que no fueren supuestos de restricciones a la capacidad, simulación o fraude.
\end{formula}

\subsection{Otros supuestos de protección de terceros}

Los supuestos en que los terceros quedan protegidos frente al principio general son los siguientes:

% TODO: contenido incompleto o ambiguo en las notas originales
% (la enumeración original menciona dos veces los "cambios en los poderes que no son oponibles" y los formula de manera imprecisa;
%  además, el texto del art. 392 transcripto excluye supuestos de simulación o fraude, mientras que la enumeración los incluye entre los supuestos de protección de terceros)
\begin{itemize}
    \item los \textbf{cambios en los poderes que no son oponibles}, según dependa de la posibilidad o no de hacerlos conocer a los terceros (cuestión vinculada con la protección de los terceros y la confirmación);
    \item el \textbf{poseedor de buena fe de cosas muebles}, que queda protegido;
    \item los supuestos de \textbf{simulación} (Art.~337);
    \item los supuestos de \textbf{fraude} (Art.~340).
\end{itemize}

\section{Cuadro Cornell: principios generales}

\begin{table}[H]
\centering
\small
\renewcommand{\arraystretch}{1.3}
\begin{tabularx}{\textwidth}{
    >{\raggedright\arraybackslash}p{0.28\textwidth}
    >{\raggedright\arraybackslash}X
}
\toprule
\textbf{Preguntas / palabras clave} &
\textbf{Notas / respuestas} \\
\midrule

\textbf{¿Qué busca producir el acto jurídico y cuál es la cuestión central sobre sus efectos?} &
El acto jurídico es un acto voluntario lícito cuyo fin inmediato es producir el nacimiento, la modificación, la transferencia o la extinción de relaciones e instituciones jurídicas; por definición, busca producir efectos jurídicos. La cuestión central es a quiénes abarcan esos efectos: las partes, los sucesores, los acreedores y los representantes, y qué reglas los rigen. \\

\midrule

\textbf{¿Qué es el \textit{nemo plus iuris}?} &
Regla del art.~399 CCCN: nadie puede transmitir a otro un derecho mejor o más extenso que el que tiene, sin perjuicio de las excepciones legalmente dispuestas. Proviene del Derecho Romano (estaba ya en el Digesto) y se encontraba en el Código anterior (art.~3270). Se ubica en el Título Quinto del Libro Primero, «Transmisión de los derechos». Es un principio cardinal del derecho privado. Los derechos son transmisibles, salvo estipulación válida de las partes o prohibición legal. \\

\midrule

\textbf{¿Cómo opera el \textit{nemo plus iuris} en el caso de Juan, Verónica y Daniela?} &
Juan obtuvo la propiedad de Verónica mediante intimidación y luego se la vendió a Daniela. Juan transmite el derecho con el mismo alcance que lo tiene, es decir, con el vicio; por lo tanto, Daniela, como sucesora, recibe un derecho afectado por ese vicio, salvo que sea adquirente de buena fe y a título oneroso (art.~392 CCCN). \\

\midrule

\textbf{¿Cuándo queda protegido el adquirente frente al \textit{nemo plus iuris}?} &
Según el art.~392 CCCN, el subadquirente de buena fe y a título oneroso queda a salvo de cualquier acción de nulidad que afecte el antecedente de su título; en el caso, Daniela adquiere y quien pierde es Verónica. También quedan protegidos el poseedor de buena fe de cosas muebles, los terceros en los supuestos de simulación (art.~337) y de fraude (art.~340), y los cambios en los poderes que no son oponibles a terceros. \\

\midrule

\textbf{¿Por qué un inquilino no puede vender el inmueble que alquila?} &
Porque no puede transmitir un derecho mejor que el que tiene: su derecho es solo al uso y goce en los términos de un contrato de locación, y no en los términos del dominio del dueño de la cosa. Es un límite a la transmisibilidad de los derechos. \\

\midrule

\multicolumn{2}{p{\dimexpr\textwidth-2\tabcolsep\relax}}{
\textbf{Resumen:}
El acto jurídico busca producir efectos, y el primer principio para determinar su alcance es el \textit{nemo plus iuris} (art.~399 CCCN): nadie transmite un derecho mejor o más extenso que el que tiene, de modo que el adquirente recibe el derecho con sus vicios. La excepción estudiada es la del adquirente de buena fe y a título oneroso (art.~392 CCCN), a la que se suman otros supuestos de protección de terceros.
} \\

\bottomrule
\end{tabularx}
\end{table}

% ============================================================
% CAPÍTULO 2
% ============================================================
\chapter[Partes y efecto relativo]{Las partes y el efecto relativo de los actos jurídicos}

\section{Efecto relativo de los actos jurídicos (Art.~1021 CCCN)}

El segundo gran marco teórico es el efecto relativo. Los actos jurídicos producen efectos entre las partes que los celebran; es decir, tienen un \textbf{efecto relativo}: sus efectos afectan a las partes que celebran el acto.

\begin{formula}{Art.~1021 CCCN --- Regla general (efecto relativo)}
El contrato solo tiene efectos entre las partes contratantes; no lo tiene con respecto a terceros, excepto en los casos previstos por la ley.
\end{formula}

Nadie puede verse afectado ni beneficiado por un acto del que no fue parte. Los actos jurídicos tienen valor entre las partes; como regla general, así lo establece el art.~1021.

\section{Definición de parte (Art.~1023 CCCN)}

La noción de «parte» resulta decisiva para determinar quién queda alcanzado por los efectos del acto. Esto conduce a considerar cuál es la definición de «parte», que se encuentra en el art.~1023.

\begin{definition}{Parte}
\textbf{Parte} es quien concurre a un acto ejerciendo prerrogativas jurídicas propias.

\smallskip
\textbf{Art.~1023 CCCN --- Parte del contrato.} Se considera parte del contrato a quien:
\begin{enumerate}[label=\alph*)]
    \item lo otorga a nombre propio, aunque lo haga en interés ajeno;
    \item es representado por un otorgante que actúa en su nombre o en su interés;
    \item manifiesta la voluntad contractual, aunque sea transmitida por un corredor o por un agente sin representación.
\end{enumerate}
\end{definition}

\subsection{Distinción entre parte y representante}

Quien actúa en nombre propio, o por medio de un representante que actúa en su nombre, es la parte, aunque en el acto conste que concurre el señor Fulanito de Tal en representación de Menganito de Tal. En ese caso, la parte es el \textbf{representado} (Menganito de Tal). El \textbf{representante} no ejerce una prerrogativa jurídica propia, sino la de otro.

\section{Cuadro Cornell: partes y efecto relativo}

\begin{table}[H]
\centering
\small
\renewcommand{\arraystretch}{1.3}
\begin{tabularx}{\textwidth}{
    >{\raggedright\arraybackslash}p{0.28\textwidth}
    >{\raggedright\arraybackslash}X
}
\toprule
\textbf{Preguntas / palabras clave} &
\textbf{Notas / respuestas} \\
\midrule

\textbf{¿Cuál es el efecto relativo de los actos jurídicos?} &
Los contratos solo producen efectos entre las partes (art.~1021 CCCN); no los producen respecto de terceros, excepto en los casos previstos por la ley. Un tercero ni puede beneficiarse ni perjudicarse por un acto en el que no intervino, salvo excepciones legales (estipulación a favor de tercero, acciones directa y subrogatoria). \\

\midrule

\textbf{¿Quién es «parte» en un acto jurídico?} &
Quien concurre al acto ejerciendo prerrogativas jurídicas propias (art.~1023 CCCN). Se considera parte a quien: a) otorga el contrato a nombre propio, aunque lo haga en interés ajeno; b) es representado por un otorgante que actúa en su nombre o en su interés; c) manifiesta la voluntad contractual, aunque sea transmitida por un corredor o por un agente sin representación. \\

\midrule

\textbf{¿Es parte el representante?} &
No. El representante no ejerce una prerrogativa jurídica propia, sino la de otro: actúa en nombre del representado, quien sí es la parte, aunque en el acto conste que concurre otra persona en su representación. El mandante de un poder, por ejemplo, es la parte aunque no esté presente. \\

\midrule

\multicolumn{2}{p{\dimexpr\textwidth-2\tabcolsep\relax}}{
\textbf{Resumen:}
Los actos jurídicos tienen efecto relativo: producen efectos entre las partes y, como regla, no alcanzan a terceros (art.~1021 CCCN). Es parte quien concurre ejerciendo prerrogativas jurídicas propias (art.~1023 CCCN); por eso, cuando hay representación, la parte es el representado y no el representante.
} \\

\bottomrule
\end{tabularx}
\end{table}

% ============================================================
% CAPÍTULO 3
% ============================================================
\chapter{Representación}

El Código regula, en el Título Cuarto del Libro Primero, todo lo relativo a la representación. Esta materia es muy importante para la vida de los abogados, ya que generalmente quienes ejercen la profesión cumplen funciones de representación, dadas por poderes o por distintos mecanismos.

\section{Tipos de representación (Art.~358 CCCN)}

El art.~358 del CCCN clasifica la representación en tres tipos:

\begin{center}
\renewcommand{\arraystretch}{1.3}
\begin{tabularx}{\textwidth}{
    >{\raggedright\arraybackslash\bfseries}p{2.4cm}
    >{\raggedright\arraybackslash}p{4.4cm}
    >{\raggedright\arraybackslash}X
}
\toprule
\textbf{Tipo} & \textbf{Origen} & \textbf{Ejemplo} \\
\midrule
Voluntaria & Acto jurídico (poder notarial) & Juana le otorga poder a Daniela para celebrar un contrato de mutuo en su nombre. \\
\addlinespace
Legal & Regla de derecho (la ley) & Padres que representan a sus hijos menores (Art.~101 CCCN); tutores; curadores. \\
\addlinespace
Orgánica & Estatuto de persona jurídica & El presidente de una sociedad que firma contratos porque el estatuto le atribuye esa función. \\
\bottomrule
\end{tabularx}
\end{center}

\section{Representación voluntaria (Arts.~362--381 CCCN)}

Los arts.~362 a 381 regulan la forma, la capacidad, los vicios, la actuación y la extinción de la representación voluntaria. En cuanto a la capacidad:

\begin{itemize}
    \item \textbf{Representante:} basta con el \textbf{discernimiento}. Puede serlo desde que tiene discernimiento para los actos lícitos, a los trece años.
    \item \textbf{Representado:} se requiere capacidad para otorgar el poder; debe ser una persona capaz.
\end{itemize}

\section{Facultades expresas (Art.~375 CCCN)}

El art.~375 es de máxima relevancia práctica: regula los casos en que se requieren \textbf{facultades expresas}. Quien, como abogado, vaya a realizar alguno de estos actos de mayor relevancia necesita lo que se denomina un \textbf{poder especial} de su representado.

\begin{formula}{Art.~375 CCCN --- Facultades expresas}
Las facultades para ciertos actos deben ser expresas. Entre ellos:
\begin{itemize}
    \item reconocer hijos;
    \item aceptar herencias;
    \item contraer matrimonio;
    \item hacer donaciones;
    \item avalar;
    \item constituir hipotecas;
    \item otorgar testamentos;
    \item transar;
    \item comprometer en árbitros;
    \item y otros actos de especial trascendencia.
\end{itemize}
\end{formula}

\section{Cuadro Cornell: representación}

\begin{table}[H]
\centering
\small
\renewcommand{\arraystretch}{1.3}
\begin{tabularx}{\textwidth}{
    >{\raggedright\arraybackslash}p{0.28\textwidth}
    >{\raggedright\arraybackslash}X
}
\toprule
\textbf{Preguntas / palabras clave} &
\textbf{Notas / respuestas} \\
\midrule

\textbf{¿Qué tipos de representación existen y cuál es el origen de cada una?} &
Arts.~358 y ss. CCCN. 1) \textbf{Voluntaria}: surge de un acto jurídico (poder notarial). 2) \textbf{Legal}: surge de la ley (padres respecto de hijos menores, tutores, curadores). 3) \textbf{Orgánica}: surge del estatuto de una persona jurídica (presidente de una sociedad). \\

\midrule

\textbf{¿Qué capacidad se requiere en la representación voluntaria?} &
Arts.~362 a 381 CCCN. Para el representante basta con el discernimiento, que se tiene desde los trece años para los actos lícitos. Para otorgar el poder se requiere capacidad en el representado, que debe ser una persona capaz. \\

\midrule

\textbf{¿Cuándo se requieren facultades expresas y qué se necesita en esos casos?} &
Art.~375 CCCN: las facultades para ciertos actos deben ser expresas, y para ellos se necesita un poder especial del representado. Entre esos actos: reconocer hijos, aceptar herencias, contraer matrimonio, hacer donaciones, avalar, constituir hipotecas, otorgar testamentos, transar, comprometer en árbitros y otros actos de especial trascendencia. \\

\midrule

\multicolumn{2}{p{\dimexpr\textwidth-2\tabcolsep\relax}}{
\textbf{Resumen:}
La representación puede ser voluntaria (nace de un acto jurídico), legal (nace de la ley) u orgánica (nace del estatuto de una persona jurídica). En la representación voluntaria basta el discernimiento para el representante, mientras que el representado debe ser capaz para otorgar el poder. Para los actos de mayor relevancia enumerados en el art.~375 CCCN se requieren facultades expresas, es decir, un poder especial.
} \\

\bottomrule
\end{tabularx}
\end{table}

% ============================================================
% CAPÍTULO 4
% ============================================================
\chapter{Sucesores}

A un acto concurren las partes con sus representantes. Después se encuentran los \textbf{sucesores}, que son aquellos que reciben un derecho. El sucesor puede ser \textbf{universal} o \textbf{singular}.

\section{Sucesor universal}

\begin{definition}{Sucesor universal}
Es quien recibe \textbf{todo o una parte indivisa del patrimonio} de otro.
\end{definition}

Quien sucede a su padre porque este falleció se convierte en su sucesor universal. Y aunque sean dos herederos, ambos son universales porque tienen \textbf{vocación al todo}.

\begin{example}{sucesión entre dos hermanos}
Supongamos que el padre tiene un departamento, una cochera, un auto y dinero en efectivo. Los dos hermanos tienen vocación al todo: sin testamento, cada uno al 50\,\%. No es que cada uno tenga el 50\,\% del dinero más el 50\,\% del auto: cada uno tiene vocación al todo.

\smallskip
Después, en la partición, si uno se queda con el departamento y el otro con todos los demás bienes, en ese momento cada uno se convierte en sucesor singular de lo que recibió.
\end{example}

\begin{formula}{Art.~1024 CCCN --- Efectos respecto de los sucesores universales}
Los efectos del contrato se extienden activa y pasivamente a los sucesores universales, a no ser que las obligaciones que de él nacen sean inherentes a la persona, o que la transmisión sea incompatible con la naturaleza de la obligación, o esté prohibida por una cláusula del contrato o la ley.
\end{formula}

Los límites de la transmisión a los sucesores universales se ilustran con el siguiente ejemplo.

\begin{example}{obligación inherente a la persona (pintor famoso)}
Supongamos que se contrata a un pintor famoso —con nombre propio, por sus condiciones artísticas— para que realice un retrato. Si el pintor fallece, su sucesor no puede venir a cumplir el contrato, porque se trata de una obligación inherente a la persona del artista.
\end{example}

Tampoco se transmiten las obligaciones de derecho de familia, ni se transmite el usufructo o el contrato de uso y habitación.

\section{Sucesor singular}

\begin{definition}{Sucesor singular}
Es quien recibe un \textbf{derecho determinado} (por ejemplo, un inmueble o un auto); es decir, quien sigue un bien determinado.
\end{definition}

\begin{important}
Los sucesores singulares, por regla, \textbf{no se ven afectados} por los actos jurídicos de la parte de quien reciben el bien, y se asimilan a los terceros completamente extraños. La excepción está dada por:
\begin{itemize}
    \item los actos jurídicos que son \textbf{antecedentes del derecho transmitido};
    \item las \textbf{obligaciones que pesan sobre la cosa};
    \item los \textbf{derechos accesorios del objeto}.
\end{itemize}
\end{important}

\begin{example}{inmueble con contrato de locación}
Quien hereda un inmueble sujeto a un contrato de locación recibe el inmueble con el contrato de locación, aunque sea sucesor singular, porque sigue un bien determinado.
\end{example}

\begin{example}{auto con prenda}
Quien hereda un auto sujeto a una prenda hereda el auto con la prenda: la prenda lo afecta. Esa es la consecuencia de ser sucesor singular.
\end{example}

\section{Cuadro Cornell: sucesores}

\begin{table}[H]
\centering
\small
\renewcommand{\arraystretch}{1.3}
\begin{tabularx}{\textwidth}{
    >{\raggedright\arraybackslash}p{0.28\textwidth}
    >{\raggedright\arraybackslash}X
}
\toprule
\textbf{Preguntas / palabras clave} &
\textbf{Notas / respuestas} \\
\midrule

\textbf{¿Qué son los sucesores universales?} &
Quienes reciben la totalidad o una parte indivisa del patrimonio de otra persona (art.~1024 CCCN): herederos \textit{ab intestato} o testamentarios. Aunque sean varios herederos, todos son universales porque tienen vocación al todo. \\

\midrule

\textbf{¿Se extienden los efectos del contrato a los sucesores universales y cuáles son los límites?} &
Sí: los efectos del contrato se extienden activa y pasivamente a los sucesores universales (art.~1024 CCCN), salvo que las obligaciones sean inherentes a la persona, que la transmisión sea incompatible con la naturaleza de la obligación, o que esté prohibida por una cláusula del contrato o la ley. \\

\midrule

\textbf{¿Qué ejemplos de obligaciones o derechos no se transmiten a los sucesores universales?} &
La obligación de un pintor famoso contratado por sus condiciones artísticas para realizar un retrato (obligación inherente a la persona: si fallece, su sucesor no puede cumplir el contrato); las obligaciones de derecho de familia; el usufructo o el contrato de uso y habitación. \\

\midrule

\textbf{¿Qué son los sucesores singulares y qué los afecta?} &
Quienes reciben un derecho determinado (un inmueble, un auto). Se asimilan a terceros respecto de los actos del causante, pero los afectan: los antecedentes del derecho transmitido, las obligaciones y cargas que pesan sobre la cosa y los derechos accesorios del objeto (por ejemplo, la prenda sobre un auto heredado, o la locación de un inmueble heredado). \\

\midrule

\textbf{¿Cuándo un sucesor universal se convierte en singular?} &
En la partición: si, por ejemplo, de los bienes del padre un hermano se queda con el departamento y el otro con los demás bienes, cada uno se convierte en sucesor singular de lo que recibió. \\

\midrule

\multicolumn{2}{p{\dimexpr\textwidth-2\tabcolsep\relax}}{
\textbf{Resumen:}
Los sucesores universales reciben todo o una parte indivisa del patrimonio y ocupan el lugar del causante: los efectos del contrato se les extienden activa y pasivamente, salvo obligaciones inherentes a la persona, incompatibles con la naturaleza de la obligación o vedadas por el contrato o la ley (art.~1024 CCCN). Los sucesores singulares siguen un bien determinado: por regla no se ven afectados por los actos de la parte de quien lo reciben, salvo por los antecedentes del derecho transmitido, las obligaciones que pesan sobre la cosa y los derechos accesorios del objeto.
} \\

\bottomrule
\end{tabularx}
\end{table}

% ============================================================
% CAPÍTULO 5
% ============================================================
\chapter[Terceros y acciones]{Los terceros y las acciones a su disposición}

\section{Situación de los terceros}

Los terceros no pueden ni beneficiarse ni perjudicarse por los contratos, salvo que exista disposición legal o que así se haya convenido entre las partes. Por regla, los terceros son \textbf{ajenos} y no pueden verse afectados por las disposiciones de un acto en el que no fueron parte.

% TODO: contenido incompleto o ambiguo en las notas originales
% (la expresión "tanto en tales asistencias" del ejemplo original es imprecisa; se conserva tal como figura)
\begin{example}{estipulación a favor de un tercero}
Supongamos que se establece una estipulación a favor de un tercero: «yo le voy a dar plata para que vos le des a Fulanito de Tal tanto en tales asistencias». Esto es una estipulación a favor de un tercero y es perfectamente válido.
\end{example}

\section{Las acciones de los acreedores frente a terceros}

Cuando un acreedor debe actuar frente a quien no es parte en el vínculo que lo une con su deudor, el ordenamiento le concede, excepcionalmente, la \textbf{acción directa} (arts.~736 y 737 CCCN) contra el deudor de su deudor y la \textbf{acción subrogatoria} (arts.~739 y 740 CCCN) para ejercer los derechos de un deudor inactivo. Ambas se desarrollan y se comparan en el capítulo~\ref{ch:garantia}, junto con la acción de simulación y la acción de inoponibilidad, estudiadas en la Parte~VI.

\section{Cuadro Cornell: terceros y acciones}

\begin{table}[H]
\centering
\small
\renewcommand{\arraystretch}{1.3}
\begin{tabularx}{\textwidth}{
    >{\raggedright\arraybackslash}p{0.28\textwidth}
    >{\raggedright\arraybackslash}X
}
\toprule
\textbf{Preguntas / palabras clave} &
\textbf{Notas / respuestas} \\
\midrule

\textbf{¿Puede un contrato obligar o beneficiar a un tercero?} &
Por regla, no: los terceros son ajenos y no pueden ni beneficiarse ni perjudicarse por los contratos, salvo disposición legal o que se haya convenido entre las partes. Por ejemplo, es perfectamente válida una estipulación a favor de un tercero. \\

\midrule

\textbf{¿Qué es la acción directa y qué requisitos tiene?} &
Permite al acreedor cobrar lo que un tercero debe a su deudor, actuando directamente contra el deudor de su deudor, hasta el importe de su propio crédito (arts.~736--737 CCCN). La ejerce por derecho propio y en su exclusivo beneficio; es excepcional y de interpretación restrictiva. Requisitos: crédito exigible del acreedor contra su deudor; deuda correlativa exigible del tercero a favor del deudor; homogeneidad de ambos créditos; ningún crédito embargado ni indisponible; notificación de la demanda al deudor. \\

\midrule

\textbf{¿Qué es la acción subrogatoria y qué derechos quedan excluidos?} &
El acreedor de un crédito cierto, exigible o no, ocupa el lugar de su deudor remiso y ejerce judicialmente los derechos patrimoniales que este no ejerce, cuando esa omisión afecta el cobro de su acreencia (arts.~739--740 CCCN). El acreedor no goza de preferencia sobre los bienes obtenidos. Quedan excluidos los derechos inherentes a la persona del deudor, los no patrimoniales, los patrimoniales no susceptibles de ejecución forzada y los excluidos por disposición legal. \\

\midrule

\textbf{¿En qué se diferencian la acción directa y la subrogatoria?} &
Directa: exige créditos homogéneos; el acreedor actúa por derecho propio, en su nombre; beneficia solo al acreedor que la ejerce; se limita al importe de su crédito. Subrogatoria: exige la inacción del deudor que perjudica al acreedor; el acreedor actúa en nombre del deudor; lo recuperado ingresa al patrimonio del deudor y beneficia a todos los acreedores. \\

\midrule

\multicolumn{2}{p{\dimexpr\textwidth-2\tabcolsep\relax}}{
\textbf{Resumen:}
Los terceros, por regla, son ajenos al acto y no pueden beneficiarse ni perjudicarse por él, salvo disposición legal o convenio. Excepcionalmente, un acreedor puede actuar mediante la acción directa, que cobra contra el deudor del deudor y solo beneficia a quien la ejerce (arts.~736--737 CCCN), o mediante la acción subrogatoria, que ejerce en nombre del deudor inactivo y beneficia a todos los acreedores (arts.~739--740 CCCN).
} \\

\bottomrule
\end{tabularx}
\end{table}

% ============================================================
% CAPÍTULO 6
% ============================================================

\section{Síntesis general de los efectos del acto jurídico}

Los efectos del acto jurídico se rigen por dos principios estructurantes: el \textbf{\textit{nemo plus iuris}} (nadie puede transmitir un derecho mejor que el que tiene, Art.~399 CCCN) y el \textbf{efecto relativo} (los actos solo producen efectos entre las partes, Art.~1021 CCCN).

A partir de estos ejes, el sistema regula quiénes resultan alcanzados:

\begin{itemize}
    \item las \textbf{partes};
    \item sus \textbf{representantes} (voluntarios, legales u orgánicos);
    \item los \textbf{sucesores universales}, que ocupan el mismo lugar que el causante en todos los derechos y obligaciones transmisibles;
    \item los \textbf{sucesores singulares}, que solo reciben los antecedentes del derecho transmitido, las cargas sobre la cosa y sus accesorios.
\end{itemize}

Los \textbf{terceros}, en principio ajenos, pueden excepcionalmente actuar mediante la \textbf{acción directa} (Arts.~736--737 CCCN) o la \textbf{acción subrogatoria} (Arts.~739--740 CCCN).

Estas instituciones son herramientas cotidianas de la práctica jurídica: delimitan quiénes quedan obligados por los contratos, cómo se transmiten los bienes en una sucesión y de qué modo un acreedor puede cobrar su crédito cuando el deudor es negligente.
\partintro{Un acto jurídico puede existir y, sin embargo, presentar defectos que impiden que produzca todos sus efectos normales. Esta parte estudia los vicios de la voluntad (error, dolo y violencia), el vicio de lesión, que afecta el equilibrio de las prestaciones, y los vicios de simulación y de fraude, que afectan la sinceridad del acto y la garantía de los acreedores.}
\part{Vicios de los actos jurídicos: voluntad, equilibrio y sinceridad}
\chapter{Fundamentos y método del Código Civil y Comercial}
% ============================================================

\section{Introducción: ¿vicios de la voluntad o del consentimiento?}

La doctrina ha discutido si estos institutos debían denominarse \textbf{vicios de la voluntad} o \textbf{vicios del consentimiento}. La mayoría de los juristas argentinos, al igual que el propio Código, se inclina por la primera denominación, señalando que es conceptualmente más amplia que la segunda.

La razón es de fondo: hablar de vicios del consentimiento los reduciría a defectos propios de los contratos ---un tipo específico de acto jurídico---, mientras que los vicios de la voluntad afectan a todo acto voluntario.

\begin{articulo}{Art. 260 CCyCN --- Acto voluntario}
Son requisitos del acto voluntario: discernimiento, intención y libertad (condiciones internas) y manifestación de la voluntad (condición externa). Los vicios afectan a alguno de esos elementos.
\end{articulo}

\begin{articulo}{Art. 261 CCyCN --- Acto involuntario}
Determina cuándo falta el discernimiento como condición interna del acto voluntario.
\end{articulo}

Los actos voluntarios son aquellos que presentan tres condiciones internas: \textbf{discernimiento, intención y libertad}, y una condición externa que es la \textbf{manifestación de la voluntad}. Los vicios afectan alguno de estos elementos. El estudio de los vicios constituye una introducción al régimen de las \textbf{ineficacias}: la posibilidad de que el acto no produzca efectos jurídicos por la presencia de un defecto en su celebración.

Las preguntas centrales que guían este estudio son:
\begin{itemize}
    \item Bajo qué circunstancias un acto voluntario puede ser anulado por la presencia de un defecto en su celebración.
    \item Cómo se afecta la intención y cómo se afecta la libertad como elementos internos.
    \item Qué requisitos deben reunir los vicios para que pueda declararse la nulidad.
    \item Qué consecuencias jurídicas se siguen de la presencia de un vicio.
\end{itemize}

\section{Método del Código Civil y Comercial}

Desde el punto de vista metodológico, los capítulos 2 a 4 del Título Cuarto del CCyCN organizan los vicios comenzando en el artículo 265 con el error. El Código distingue dos grandes grupos bajo el término común de «vicios»:

\begin{table}[H]
\centering
\begin{tabularx}{\textwidth}{>{\bfseries\raggedright\arraybackslash}X >{\raggedright\arraybackslash}X}
\toprule
\textbf{Vicios de la voluntad} & \textbf{Vicios propios de los actos jurídicos} \\
\midrule
Error (Arts. 265--270) & Lesión \\
Dolo (Arts. 271--275) & Simulación \\
Violencia (Arts. 276--278) & Fraude \\
\bottomrule
\end{tabularx}
\caption{Clasificación de los vicios en el CCyCN}
\end{table}

Los primeros ---error, dolo y violencia--- afectan los elementos internos del acto voluntario. Los segundos ---lesión, simulación y fraude--- tienen un tratamiento posterior en el Código y son propios de los actos jurídicos.

% ============================================================
\chapter{El error como vicio de la voluntad (Arts.~265--270)}
\label{ch:error}
% ============================================================

\section{Concepto y clasificación del error de hecho}

El error afecta la \textbf{intención} como elemento interno del acto voluntario. La doctrina trata conjuntamente dos supuestos: la \textbf{ignorancia} ---la falta de conocimiento sobre algún elemento del acto--- y el \textbf{error propiamente dicho} ---un falso conocimiento---. Jurídicamente ambos se tratan como un único supuesto.

La primera gran distinción es entre el error de derecho y el error de hecho:

\begin{table}[H]
\centering
\begin{tabularx}{\textwidth}{>{\raggedright\arraybackslash}X >{\raggedright\arraybackslash}X}
\toprule
\textbf{Error de derecho} & \textbf{Error de hecho} \\
\midrule
Falso conocimiento de la normativa legal aplicable a una relación o situación jurídica. & Recae sobre los elementos fácticos del acto o circunstancias del mismo. \\[6pt]
Regulado en el Art. 8 CCyCN: la ignorancia de la ley no sirve de excusa para su incumplimiento. & Regulado en los Arts. 265 a 270 CCyCN. \\[6pt]
No puede invocarse como vicio. La ley se presume conocida por todos. & Es el que puede dar lugar a la nulidad del acto, con los requisitos del caso. \\
\bottomrule
\end{tabularx}
\caption{Error de derecho versus error de hecho}
\end{table}

\begin{articulo}{Art. 8 CCyCN --- Principio de inexcusabilidad}
La ignorancia de las leyes no sirve de excusa para su cumplimiento, si la excepción no está autorizada por el ordenamiento jurídico.
\end{articulo}

Dentro del error de hecho, la clasificación fundamental es:

\begin{itemize}
    \item \textbf{Error esencial}: recae sobre un elemento esencial del acto. Puede causar la nulidad si además es reconocible.
    \item \textbf{Error incidental}: recae sobre elementos secundarios que no son determinantes de la voluntad. No causa nulidad.
\end{itemize}

La fórmula es: cuando el error es \textbf{esencial Y reconocible}, causa la nulidad del acto (si es bilateral o unilateral de aceptación).

\section{El error reconocible --- Art. 266}

El CCyCN introduce un requisito que no estaba en el Código de Vélez: la \textbf{reconocibilidad}. En el Código anterior se exigía que el error fuera esencial y \textbf{excusable}. El nuevo Código reemplaza la excusabilidad por la reconocibilidad.

\begin{articulo}{Art. 266 CCyCN --- Error reconocible}
El error es reconocible cuando el destinatario de la declaración lo pudo conocer según la naturaleza del acto, las circunstancias de persona, tiempo y lugar.
\end{articulo}

\begin{table}[H]
\centering
\begin{tabularx}{\textwidth}{>{\raggedright\arraybackslash}X >{\raggedright\arraybackslash}X}
\toprule
\textbf{Código de Vélez} & \textbf{CCyCN (vigente)} \\
\midrule
Error esencial + excusable. & Error esencial + reconocible. \\[4pt]
Miraba al sujeto que erraba: ¿tenía razón para equivocarse? & Mira al destinatario: ¿pudo darse cuenta del error ajeno? \\[4pt]
Requisito subjetivo centrado en quien declara. & Requisito objetivo centrado en quien recibe la declaración. \\
\bottomrule
\end{tabularx}
\caption{Cambio de requisito: Código de Vélez versus CCyCN}
\end{table}

La excusabilidad del Código de Vélez miraba al sujeto que emitía la declaración: que hubiera tenido alguna razón objetiva para errar. Ese criterio fue abandonado por el nuevo Código, siguiendo el modelo del derecho italiano.

El error es reconocible cuando el destinatario de la declaración no pudo \textit{no conocer} el error, según la naturaleza del acto o las circunstancias de persona, tiempo y lugar. Existe así una expectativa de que la declaración viciada por un error esencial debió haber sido advertida por la parte destinataria.

Si el destinatario no advierte el error siendo que cualquiera lo hubiera advertido ---pero tampoco actúa con malicia (porque si actúa con malicia incurriría en dolo)---, debe admitir que el acto sea anulado. La consecuencia de que el error sea reconocible es la nulidad relativa del acto.

\section{El error esencial --- supuestos del Art. 267}

\begin{articulo}{Art. 267 CCyCN --- Supuestos de error esencial}
El error de hecho es esencial cuando recae sobre: (a) la naturaleza del acto; (b) el objeto sobre un bien o un hecho diverso o de distinta especie, o una calidad, extensión o suma diversa; (c) la cualidad sustancial del bien que haya sido determinante de la voluntad; (d) los motivos personales relevantes que hayan sido incorporados expresa o tácitamente al acto; (e) la persona con la cual se celebró o a la cual se refiere el acto, si ella fue determinante para su celebración.
\end{articulo}

\subsection{Inciso a) --- Error sobre la naturaleza del acto}

Es el error sobre el \textbf{tipo de acto} que, conforme al ordenamiento jurídico, cada parte cree celebrar.

\begin{example}{Error sobre la naturaleza del acto}
A le dice a B: «Me das ese cuadro». A interpreta que B está haciendo una donación; B interpreta que está realizando un comodato (préstamo gratuito). Hay un error sobre la naturaleza del acto: uno pensó que era una donación y el otro que era un préstamo. Consecuencia: nulidad relativa, si el error es reconocible y esencial.
\end{example}

\subsection{Inciso b) --- Error sobre el objeto}

Es el error sobre un \textbf{bien o un hecho diverso o de distinta especie}, o sobre una calidad, extensión o suma diversa.

\begin{example}{Error sobre el objeto --- distintas especies}
A cree estar comprando un caballo pura sangre para competencias; B interpreta que A está comprando un caballo de carga para trabajo en el campo. El objeto del acto ---el caballo--- no es el mismo en la mente de ambas partes. Si el error es reconocible, da lugar a la nulidad.
\end{example}

\begin{example}{Error sobre el objeto --- calidad}
A cree estar comprando un vino de excelente calidad y en realidad es un vino de calidad inferior. Si la otra parte debería haberse dado cuenta del error (reconocibilidad), el acto puede ser anulado.
\end{example}

\subsection{Inciso c) --- Error sobre la cualidad sustancial del bien}

Este supuesto presenta un factor \textbf{objetivo} y uno \textbf{subjetivo}: por un lado, la cualidad debe ser sustancial al bien (factor objetivo); por otro lado, debe haber sido determinante de la voluntad según la apreciación común a las circunstancias del caso (factor subjetivo).

\begin{table}[H]
\centering
\begin{tabularx}{\textwidth}{>{\raggedright\arraybackslash}X >{\raggedright\arraybackslash}X}
\toprule
\textbf{Código de Vélez} & \textbf{CCyCN (vigente)} \\
\midrule
Utilizaba el criterio de la cualidad sustancial sin la palabra «sustancial» en ciertos pasajes, generando debate. & Incorpora la palabra «sustancial» y la conjuga con el requisito de que haya sido «determinante de la voluntad». \\[4pt]
Discusión entre criterio objetivo y subjetivo de la cualidad. & Combina ambos factores: objetivo (cualidad sustancial del bien) + subjetivo (determinante de la voluntad según apreciación común). \\
\bottomrule
\end{tabularx}
\caption{Debate histórico sobre la cualidad sustancial}
\end{table}

\begin{example}{Error sobre la cualidad sustancial}
A está comprando una notebook y necesita que tenga buena memoria RAM porque va a trabajar con edición de videos. Esto es una cualidad sustancial de la computadora (factor objetivo), y además ha sido determinante para su voluntad de contratar (factor subjetivo). Si se le entrega una computadora sin esa capacidad y el error era reconocible, puede dar lugar a la anulabilidad del acto.
\end{example}

\subsection{Inciso d) --- Error sobre los motivos}

Este inciso se vincula con la causa de los actos jurídicos (Art. 281 CCyCN). El error puede recaer sobre los \textbf{motivos personales relevantes} que hayan sido incorporados expresa o tácitamente al acto.

\begin{example}{Error sobre los motivos}
A quiere hacer una donación a una persona pensando que esa persona le salvó la vida. En realidad no es la misma persona. El motivo ---la gratitud hacia quien le salvó la vida--- fue incorporado al acto y es esencial. Si el error es reconocible, puede anularse la donación.
\end{example}

\subsection{Inciso e) --- Error sobre la persona}

El error sobre la persona solo causa nulidad cuando la identidad de la persona fue \textbf{determinante para la celebración del acto}. Si no es determinante, es un error incidental y no produce nulidad.

\begin{example}{Error sobre la persona}
A quiere contratar a un pintor famoso para que le haga un retrato de su madre. Existe un homónimo ---alguien con el mismo nombre--- y A contrata al homónimo creyendo que es el pintor famoso. Hay un error esencial sobre la persona: era determinante quién era el pintor, el homónimo debería haberse dado cuenta de que no era el artista buscado, y el acto puede anularse.
\end{example}

\section{Error de cálculo --- Art. 268}

\begin{articulo}{Art. 268 CCyCN --- Error de cálculo}
El error de cálculo no da lugar a la nulidad del acto, sino solamente a su rectificación, excepto que sea determinante del consentimiento.
\end{articulo}

El error de cálculo no da lugar a la nulidad sino únicamente a una \textbf{rectificación}, dado que se trata de un error aritmético que no fue determinante del consentimiento.

\begin{example}{Error de cálculo}
Las partes pactan vender 40 kg de mercadería a \$100 por kg. Queda claro el valor unitario y la cantidad. Al realizar la multiplicación en el contrato, escriben mal el precio total. Eso es simplemente un error aritmético que se corrige (rectifica) pero no da lugar a la nulidad.
\end{example}

\section{Subsistencia del acto y nulidad relativa --- Art. 269}

\begin{articulo}{Art. 269 CCyCN --- Subsistencia del acto}
La parte que incurre en error no puede solicitar la nulidad del acto si la otra ofrece ejecutarlo con las modalidades y el contenido que aquella entendió celebrar.
\end{articulo}

El artículo 269 le concede a la parte destinataria de la declaración la posibilidad de \textbf{ofrecer ejecutar el acto con las modalidades y el contenido que la parte en error pretendió} que se celebrara.

\begin{example}{Subsistencia del acto}
A le dice a B: «Me das tu auto». B lo interpreta como préstamo; A como donación. A alega error y busca anular el acto. B puede decir: «Bueno, celebremos el acto bajo la modalidad del préstamo (comodato): te presto el auto». Si eso le conviene a A, se salva la nulidad. Esto responde al principio general de conservación de los actos.
\end{example}

\begin{important}
El principio de conservación de los actos es un principio general del derecho civil: se tiende a interpretar que los actos deben ser conservados, preservando el valor de los negocios jurídicos.
\end{important}

\section{Error en la declaración --- Art. 270}

Las disposiciones sobre el error se aplican tanto a la declaración de voluntad en sí como a su \textbf{transmisión}. Pueden presentarse problemas cuando quien transmite la voluntad es un \textbf{mandatario} (con representación) o simplemente un \textbf{mensajero}. En general, se aplican todas las disposiciones sobre error esencial y reconocible.

\section{Consecuencias del error: nulidad relativa y prescripción}

La sanción que se sigue de la presencia del error es la \textbf{nulidad relativa} del acto, regulada en el Art. 388 CCyCN y susceptible de confirmación.

\begin{articulo}{Arts. 2562 y 2563 CCyCN --- Prescripción}
El plazo de prescripción para la acción de nulidad por vicios del consentimiento es de \textbf{dos años}. Comienza a correr desde que el error se conoció o pudo ser conocido.
\end{articulo}

% ============================================================
\chapter{El dolo como vicio de la voluntad (Arts.~271--275)}
% ============================================================

\section{Las acepciones del dolo en el derecho civil}

La palabra «dolo» tiene múltiples significados en el CCyCN. Es fundamental no confundirlos:

\begin{table}[H]
\centering
\begin{tabularx}{\textwidth}{>{\raggedright\arraybackslash}p{0.30\textwidth} >{\raggedright\arraybackslash}X}
\toprule
\textbf{Acepción} & \textbf{Definición y referencia} \\
\midrule
1) Elemento del acto ilícito & Intención de dañar o manifiesta indiferencia por los intereses ajenos (Art. 1724 CCyCN). Componente del factor subjetivo de atribución de responsabilidad. \\[4pt]
2) Inejecución deliberada de obligación & Incumplimiento voluntario de una obligación. Figuraba en el Código de Vélez; el nuevo CCyCN no lo define con la misma fórmula. \\[4pt]
3) Vicio de la voluntad \emph{(el que se estudia aquí)} & Maniobra engañosa realizada para determinar a una persona a celebrar un acto que no tenía intención de celebrar, o a celebrarlo en condiciones más desventajosas. \\
\bottomrule
\end{tabularx}
\caption{Las tres acepciones del dolo en el CCyCN}
\end{table}

Lo que tienen en común todas estas acepciones es una suerte de \textbf{conciencia de antijuridicidad}: el dolo siempre expresa una cierta conciencia de que se está actuando de modo contrario al derecho, intencionadamente causando un perjuicio o una situación prohibida.

\section{Concepto: acción y omisión dolosa --- Art. 271}

\begin{articulo}{Art. 271 CCyCN --- Acción y omisión dolosa}
Acción dolosa es toda aserción de lo falso, disimulación de lo verdadero, cualquier artificio, astucia o maquinación que se emplee para la celebración del acto. La omisión dolosa causa los mismos efectos que la acción dolosa, cuando el acto no se habría realizado sin la reticencia u ocultación.
\end{articulo}

La incorporación de la \textbf{omisión dolosa} es una novedad bien recibida por la doctrina. Significa que la maniobra engañosa puede realizarse tanto a través de acciones como de omisiones: ocultar información o ser reticente al proporcionarla puede inducir a la otra parte a celebrar un acto que afecta sus intereses y configura, por tanto, el concepto de engaño.

Las formas que puede tomar la maniobra engañosa son:
\begin{itemize}
    \item \textbf{Aseverar lo falso}: afirmar algo que se sabe incorrecto.
    \item \textbf{Disimular lo verdadero}: ocultar la realidad.
    \item \textbf{Cualquier artificio, astucia o maquinación}: no es una mentira casual; tiene que ser una maniobra con entidad engañosa.
    \item \textbf{Omisión dolosa}: silencio o reticencia que induce al otro a celebrar el acto en condiciones perjudiciales.
\end{itemize}

\begin{note}
En el fondo de estas conductas hay un \textbf{deber de buena fe} entre las partes: un deber de lealtad y de informar. En el derecho del consumidor estos principios son especialmente relevantes: la publicidad engañosa y la omisión de información relevante son manifestaciones colaterales del problema del dolo que pueden tener regulación específica.
\end{note}

\section{Dolo esencial --- requisitos del Art. 272}

\begin{articulo}{Art. 272 CCyCN --- Dolo esencial}
El dolo es esencial y causa la nulidad del acto si es grave, si es determinante de la voluntad, si causa un daño importante y si no ha habido dolo por ambas partes.
\end{articulo}

Para que el dolo cause la nulidad del acto deben reunirse los cuatro requisitos siguientes:

\subsection*{a) Gravedad}

La gravedad es una característica de la maniobra engañosa en sí misma. Para que el dolo sea grave debe tener \textbf{entidad para engañar}: no debe ser una maniobra burda que cualquier persona advertiría a primera vista. El parámetro es el de la persona razonablemente precavida que toma las precauciones comunes. La gravedad de la maniobra es un requisito distinto del daño importante.

\subsection*{b) Determinante de la voluntad}

El dolo debe haber sido \textbf{determinante de la voluntad}. Este requisito es fundamental para distinguirlo del dolo incidental: si sin el engaño la persona igualmente hubiera celebrado el acto, no se configura dolo esencial.

\subsection*{c) Daño importante}

Debe haberse causado un daño: una afectación de un interés o de un derecho de la parte que es víctima del acto.

\subsection*{d) No dolo de ambas partes}

No debe haber habido dolo recíproco, es decir, de ambas partes. Un principio fundamental del derecho es que \textbf{nadie puede alegar su propia torpeza}: si ambas partes se engañaban mutuamente, ninguna puede pedir la nulidad.

\section{Dolo incidental --- Art. 273}

\begin{articulo}{Art. 273 CCyCN --- Dolo incidental}
El dolo incidental no es determinante de la voluntad: la parte igualmente habría celebrado el acto, aunque en condiciones distintas. No causa la nulidad del acto, pero sí puede generar una indemnización de daños y perjuicios.
\end{articulo}

\begin{example}{Dolo incidental}
A está comprando un terreno. La otra parte realiza una maniobra para hacerle creer que el terreno tiene cierta habilitación que en realidad no tiene. Pero A de todos modos iba a comprar ese terreno. El engaño le ha provocado un perjuicio ---por ejemplo, pagó más de lo que debía---, pero no fue determinante de que celebrara el acto. En ese caso, el dolo incidental solo puede dar lugar a una acción de daños y perjuicios, pero no a la nulidad.
\end{example}

\begin{table}[H]
\centering
\begin{tabularx}{\textwidth}{>{\raggedright\arraybackslash}X >{\raggedright\arraybackslash}X}
\toprule
\textbf{Dolo esencial} & \textbf{Dolo incidental} \\
\midrule
Es determinante de la voluntad: sin el engaño, el acto no se hubiera celebrado. & No es determinante: la parte igualmente hubiera celebrado el acto. \\[4pt]
Reúne los requisitos del Art. 272: grave, determinante, daño importante, sin dolo recíproco. & El engaño existe pero no alcanza los requisitos del Art. 272 en cuanto a su incidencia en la voluntad. \\[4pt]
Consecuencia: nulidad relativa del acto + daños y perjuicios. & Consecuencia: solo daños y perjuicios (si se prueban). Sin nulidad. \\
\bottomrule
\end{tabularx}
\caption{Dolo esencial versus dolo incidental}
\end{table}

\section{Dolo de tercero --- Art. 274}

\begin{articulo}{Art. 274 CCyCN --- Dolo de tercero}
El dolo de un tercero causa la nulidad del acto si es esencial. Si el dolo de un tercero es incidental, solo da lugar a daños y perjuicios. La parte que conoció o debió conocer el dolo del tercero responde solidariamente con éste por los daños causados.
\end{articulo}

El dolo de tercero es aquel cometido por alguien distinto de la parte con la cual se está contratando. Las reglas son simétricas al dolo de la propia parte:

\begin{itemize}
    \item Si el tercero reúne los requisitos del Art. 272 (dolo esencial): causa la nulidad.
    \item Si el dolo del tercero reúne los requisitos del Art. 273 (dolo incidental): solo genera responsabilidad por daños y perjuicios.
\end{itemize}

Si la parte contratante tuvo \textbf{conocimiento del dolo del tercero} y no lo comunicó, responde \textbf{solidariamente} por todos los daños causados. La responsabilidad solidaria implica que responde por el total de la deuda, no solo por su parte proporcional.

\section{Consecuencias del dolo: nulidad y daños --- Art. 275}

\begin{articulo}{Art. 275 CCyCN --- Responsabilidad por daños}
El autor del dolo esencial o incidental debe reparar el daño causado. Cuando el dolo sea del tercero, también responde la parte que al tiempo de la celebración del acto tuvo conocimiento del dolo del tercero.
\end{articulo}

Las consecuencias del dolo esencial son:

\begin{enumerate}
    \item \textbf{Nulidad relativa del acto} (Art. 388 CCyCN). Al ser relativa, está dispuesta en beneficio de la parte afectada y es susceptible de confirmación.
    \item \textbf{Acción de daños y perjuicios} en virtud del Art. 275.
    \item \textbf{Prescripción de dos años} desde el momento en que se conoció o pudo ser conocido el dolo (Arts. 2562 y 2563 CCyCN).
\end{enumerate}

% ============================================================
\chapter{La violencia como vicio de la voluntad (Arts.~276--278)}
% ============================================================

\section{Concepto de violencia: fuerza e intimidación}

La violencia es una \textbf{coerción grave de tipo físico o moral} que se ejerce sobre una persona para determinarla a realizar un acto en contra de su voluntad. Es un vicio que afecta la \textbf{libertad} como elemento interno del acto voluntario: la posibilidad de autodeterminarse y tomar decisiones sin coerciones injustas.

La violencia adopta dos grandes formas:

\begin{table}[H]
\centering
\begin{tabularx}{\textwidth}{>{\raggedright\arraybackslash}X >{\raggedright\arraybackslash}X}
\toprule
\textbf{Fuerza (Art. 276, primera parte)} & \textbf{Intimidación (Art. 276, segunda parte)} \\
\midrule
Violencia física irresistible. & Amenazas que generan el temor de sufrir un mal grave e inminente. \\[4pt]
Ejemplo: torturas físicas para que alguien firme un contrato. & Ejemplo: «Si no me vendés tu casa, voy a hacerte daño a vos o a tu familia». \\[4pt]
La coerción es de tipo corporal y directa. & La coerción opera a través del temor generado por las amenazas. \\
\bottomrule
\end{tabularx}
\caption{Formas de la violencia}
\end{table}

\begin{example}{Fuerza física irresistible}
A es torturado físicamente para que firme la escritura de venta de su casa. Aparece todo en regla: hay una escritura firmada, inscripta en el Registro de la Propiedad. Pero ese acto está viciado por fuerza física irresistible. Se inicia un juicio de nulidad para volver la situación al estado anterior a la celebración del acto. Este hecho también puede configurar delitos penales, pero aquí se analiza la dimensión civil.
\end{example}

\section{Requisitos de la intimidación --- Art. 276}

\begin{articulo}{Art. 276 CCyCN --- Fuerza e intimidación}
La fuerza irresistible y las amenazas que generan el temor de sufrir un mal grave e inminente que no se puedan contrarrestar o evitar en la persona o bienes de la parte o de un tercero, causan la nulidad del acto. La relevancia de las amenazas debe ser juzgada teniendo en cuenta la situación del amenazado y las demás circunstancias del caso.
\end{articulo}

\subsection*{a) Las amenazas deben ser injustas}

Las amenazas que configuran el vicio son \textbf{amenazas injustas}. El ejercicio legítimo de un derecho no constituye una amenaza en este sentido. Si una persona le debe dinero a otra y le dice «si no me pagás, voy a iniciar un juicio», no le está amenazando en el sentido jurídicamente relevante: está ejerciendo su derecho. Una amenaza ilícita sería, en cambio: «si no me vendés tu casa, te voy a romper la tuya».

\subsection*{b) El temor se mide según las circunstancias del amenazado}

La relevancia de las amenazas ha de ser juzgada teniendo en cuenta la \textbf{situación del amenazado y las demás circunstancias del caso}. No hay un estándar único: la misma amenaza puede configurar intimidación para una persona pero no para otra, según sus capacidades y contexto.

\begin{example}{Medición del temor según las circunstancias}
Si alguien con millones de seguidores en redes sociales amenaza con publicar que una persona es delincuente, eso puede generar un temor fundado a sufrir un daño grave a la reputación. El contexto (quién amenaza, con qué medios, a quién) determina la relevancia de la amenaza.
\end{example}

\subsection*{c) El mal debe ser grave e inminente}

\textbf{Grave}: el mal debe tener entidad suficiente. Puede ser un mal sobre la persona (integridad física, moral o espiritual) o sobre los bienes.

\textbf{Inminente}: significa que el mal va a ocurrir y que no puede ser contrarrestado ni evitado. El término «evitar» fue agregado por el nuevo Código; no figuraba en el anterior.

\subsection*{d) El mal puede recaer sobre la parte o un tercero}

En el Código de Vélez se enumeraban taxativamente las personas sobre quienes podía recaer la amenaza: cónyuge, ascendientes o descendientes. Esta enumeración generó una disputa: ¿era taxativa o ejemplificativa? La doctrina mayoritaria la interpretaba como ejemplificativa.

El nuevo Código resuelve el problema directamente indicando: \textbf{sobre la parte o un tercero}. Puede ser un familiar, un amigo, o incluso alguien que no se conoce tanto, pero cuya situación genera el temor suficiente.

\section{Violencia de tercero y responsabilidad solidaria --- Art. 278}

Al igual que en el dolo, la violencia puede provenir de la propia parte o de un tercero. En ambos casos da lugar a la nulidad del acto. La diferencia relevante es la \textbf{responsabilidad de la parte cómplice}.

\begin{articulo}{Art. 278 CCyCN --- Responsabilidad por daños causados por violencia}
El autor de la fuerza irresistible y de las amenazas debe reparar los daños causados. La parte que al tiempo de la celebración del acto tuvo conocimiento de la fuerza o de las amenazas responde solidariamente con el autor por los daños.
\end{articulo}

\begin{example}{Violencia de tercero --- complicidad de la parte}
Un tercero amenaza a A para que le venda su casa a B. A le vende la casa a B (que es la parte contratante), pero la amenaza vino del tercero. Si B tuvo conocimiento de la fuerza o de la intimidación ejercida por el tercero, responde solidariamente por todos los daños y perjuicios: responde por el total, no solo por su parte proporcional.
\end{example}

\section{El temor reverencial}

El \textbf{temor reverencial} es el respeto o sumisión de una persona hacia sus superiores (por ejemplo, de un militar hacia un oficial superior). El temor reverencial \textbf{no es causa de nulidad}.

Para configurar el vicio de violencia deben existir amenazas concretas: la sola presencia de personas con autoridad no produce una situación de afectación de la libertad. El temor a una figura de autoridad, por sí solo, es temor reverencial y no causa nulidad. Este criterio se mantiene tanto en el Código de Vélez como en el nuevo CCyCN.

\begin{note}
Si el superior jerárquico realiza amenazas concretas que reúnen los requisitos de los Arts. 276 y siguientes, sí puede configurarse el vicio de violencia y dar lugar tanto a la nulidad como a los daños y perjuicios.
\end{note}

\section{Consecuencias de la violencia: nulidad y prescripción}

Las consecuencias de la violencia son análogas a las del dolo y el error:

\begin{enumerate}
    \item \textbf{Nulidad relativa del acto}. Al ser relativa, la parte afectada puede confirmar el acto si así lo desea.
    \item \textbf{Acción de daños y perjuicios} (Art. 278 CCyCN).
\end{enumerate}

La prescripción también es de \textbf{dos años} (Arts. 2562 y 2563 CCyCN), pero existe una diferencia crucial respecto del error y el dolo: el plazo comienza a correr desde el momento en que \textbf{cesa la violencia}, no desde que se conoció el vicio.

\begin{articulo}{Diferencia en el inicio del plazo de prescripción}
Error: desde que se conoció o pudo ser conocido el error.\\
Dolo: desde que se conoció o pudo ser conocido el dolo.\\
Violencia: desde que cesa la violencia (Arts. 2562 y 2563 CCyCN).
\end{articulo}

\begin{example}{Inicio del plazo desde el cese de la violencia}
A es secuestrado durante un año y medio. Durante ese tiempo, bajo torturas, firma un contrato vendiendo su casa. Después lo liberan. Desde el momento en que cesa la violencia ---desde que recobra la libertad--- comienza a correr el plazo de dos años para iniciar la acción de nulidad. Mientras dura la coerción, la persona no puede libremente ejercer su derecho a pedir la nulidad.
\end{example}

% ============================================================

\section{Cuadro Cornell --- Vicios de la Voluntad}
% ============================================================

\begin{table}[H]
\centering
\begin{tabularx}{\textwidth}{
    >{\raggedright\arraybackslash}p{0.30\textwidth}
    >{\raggedright\arraybackslash}X
}
\toprule
\textbf{Preguntas / palabras clave} & \textbf{Notas / respuestas} \\
\midrule

\textbf{¿Por qué se prefiere hablar de «vicios de la voluntad» y no de «vicios del consentimiento»?} &
Porque hablar de vicios del consentimiento los limitaría a los contratos (un tipo específico de acto jurídico). Los vicios de la voluntad, en cambio, afectan a todo acto voluntario. El CCyCN adopta esta denominación más amplia. \\[6pt]

\textbf{¿Cuáles son los elementos internos del acto voluntario y qué vicio afecta a cada uno?} &
Los tres elementos internos son discernimiento, intención y libertad (Art. 260 CCyCN). La \textbf{intención} es afectada por el error y el dolo; la \textbf{libertad} es afectada por la violencia (fuerza o intimidación). \\[6pt]

\textbf{¿Qué diferencia hay entre error esencial y error incidental?} &
El error esencial recae sobre un elemento esencial del acto (naturaleza, objeto, cualidad sustancial, motivos determinantes, persona). El error incidental recae sobre elementos secundarios no determinantes de la voluntad. Solo el error esencial ---y reconocible--- puede causar la nulidad. \\[6pt]

\textbf{¿Qué significa que un error sea «reconocible» (Art. 266)?} &
Que el destinatario de la declaración pudo conocer el error según la naturaleza del acto y las circunstancias de persona, tiempo y lugar. Es un criterio objetivo centrado en el destinatario, que reemplaza la «excusabilidad» del Código de Vélez (criterio subjetivo centrado en quien erraba). \\[6pt]

\textbf{¿Qué sucede si el destinatario ofrece cumplir el acto según lo que el otro quiso? (Art. 269)} &
Se evita la nulidad del acto. La parte en error no puede solicitar la nulidad si el destinatario ofrece ejecutar el acto con las modalidades y el contenido que aquella pretendió celebrar. Esto responde al principio de conservación de los actos. \\[6pt]

\textbf{¿Qué es el dolo como vicio de la voluntad? (Art. 271)} &
Una maniobra engañosa: toda aserción de lo falso, disimulación de lo verdadero, cualquier artificio, astucia o maquinación para inducir a una persona a celebrar un acto. La omisión dolosa (reticencia u ocultación) tiene los mismos efectos que la acción dolosa cuando el acto no se habría realizado sin ella. \\[6pt]

\textbf{¿Cuáles son los requisitos del dolo esencial? (Art. 272)} &
Para causar la nulidad, el dolo debe ser: (1) \textbf{grave} ---con entidad para engañar---; (2) \textbf{determinante de la voluntad} ---sin el engaño, el acto no se hubiera celebrado---; (3) causante de un \textbf{daño importante}; (4) \textbf{unilateral} ---no haber habido dolo de ambas partes---. \\[6pt]

\textbf{¿Qué diferencia al dolo incidental del dolo esencial? (Art. 273)} &
El dolo incidental no es determinante de la voluntad: la parte igualmente hubiera celebrado el acto. No causa la nulidad, pero sí puede dar lugar a una indemnización por daños y perjuicios si se prueban. \\[6pt]

\textbf{¿Puede el dolo cometido por un tercero causar la nulidad? (Art. 274)} &
Sí. El dolo de tercero causa la nulidad si es esencial, y solo daños y perjuicios si es incidental. Si la parte contratante tuvo conocimiento del dolo del tercero, responde solidariamente por todos los daños. \\[6pt]

\textbf{¿Cuáles son las dos formas de la violencia? (Art. 276)} &
\textbf{Fuerza}: coerción física irresistible. \textbf{Intimidación}: amenazas injustas que generan el temor de sufrir un mal grave e inminente que no puede contrarrestarse ni evitarse, sobre la persona o bienes de la parte o de un tercero. \\[6pt]

\textbf{¿Cómo se mide si las amenazas son relevantes como intimidación?} &
Se juzga teniendo en cuenta la situación del amenazado y las demás circunstancias del caso. No hay un estándar único: la misma amenaza puede configurar intimidación para una persona pero no para otra, según sus capacidades y contexto. \\[6pt]

\textbf{¿Qué es el temor reverencial y tiene efectos jurídicos?} &
Es el respeto o sumisión de una persona hacia su superior jerárquico (por ejemplo, un militar ante un oficial superior). El temor reverencial solo, sin amenazas concretas, no es causa de nulidad. Solo cuando se suman amenazas reales que reúnen los requisitos del Art. 276 puede configurarse el vicio de violencia. \\[6pt]

\textbf{¿Desde cuándo corre el plazo de prescripción para cada vicio?} &
En todos los casos el plazo es de dos años (Arts. 2562 y 2563 CCyCN). La diferencia está en el punto de inicio: \textbf{error y dolo}: desde que se conoció o pudo ser conocido el vicio. \textbf{Violencia}: desde que cesa la violencia. \\[6pt]

\textbf{¿Cuál es el tipo de nulidad que producen los vicios de la voluntad?} &
Todos los vicios de la voluntad (error, dolo, violencia) producen \textbf{nulidad relativa} (Art. 388 CCyCN). Eso significa que está dispuesta en beneficio de la parte afectada, puede ser confirmada por ella y prescribe a los dos años. \\

\midrule

\multicolumn{2}{p{\textwidth}}{%
\textbf{Resumen:} Los vicios de la voluntad son defectos que afectan los elementos internos del acto voluntario ---intención y libertad--- y que, cuando reúnen los requisitos que establece el CCyCN, producen la nulidad relativa del acto. El \textbf{error} (Arts. 265--270) afecta la intención por una falsa o errónea representación de la realidad: para causar la nulidad debe ser esencial ---recaer sobre la naturaleza del acto, el objeto, la cualidad sustancial, los motivos determinantes o la persona--- y reconocible por el destinatario, criterio que el nuevo Código tomó del derecho italiano en reemplazo de la excusabilidad del Código de Vélez. El \textbf{dolo} (Arts. 271--275) afecta la intención mediante una maniobra engañosa ---acción u omisión--- que para producir la nulidad debe ser grave, determinante de la voluntad, causante de un daño importante y unilateral; si no es determinante, solo genera daños y perjuicios (dolo incidental). La \textbf{violencia} (Arts. 276--278) afecta la libertad mediante fuerza física irresistible o intimidación a través de amenazas injustas que generan el temor de sufrir un mal grave e inminente; la relevancia de las amenazas se mide según las circunstancias del amenazado. En los tres vicios, la consecuencia es la nulidad relativa del acto ---confirmable por la parte afectada--- y la posibilidad de reclamar daños y perjuicios. El plazo de prescripción es de dos años en todos los casos: para el error y el dolo, desde que se conoció o pudo conocerse el vicio; para la violencia, desde que cesa la coerción.%
} \\

\bottomrule
\end{tabularx}
\caption{Cuadro Cornell --- Vicios de la Voluntad (Error, Dolo y Violencia)}
\end{table}
\chapter[Antecedentes históricos y derecho comparado]{Antecedentes históricos y derecho comparado de la lesión}
% ===========================================================================

\section{Los vicios de los actos jurídicos según el Código Civil y Comercial}

El estudio de la lesión se inserta dentro del análisis general de los vicios propios de los actos jurídicos. En el ordenamiento argentino vigente, estos vicios son tres: la \textbf{lesión}, la \textbf{simulación} y el \textbf{fraude}.

A diferencia del tratamiento que recibían los vicios de la voluntad —cada uno regulado en un capítulo propio—, en el Código Civil y Comercial (CCC) la lesión, la simulación y el fraude están tratados dentro del capítulo sexto del título cuarto, distribuidos en tres secciones.

\begin{note}
La lesión no estaba incluida originariamente en el Código de Vélez Sarsfield. Se trata de un instituto que resulta central para comprender el funcionamiento concreto de la justicia conmutativa en el derecho privado patrimonial, razón por la cual merece un desarrollo extenso y diferenciado, con antecedentes históricos propios.
\end{note}

\section{Concepto de lesión: una distinción necesaria}

\begin{definition}{Lesión como vicio de los actos jurídicos}
Cuando se habla de lesión como vicio de los actos jurídicos, \textbf{no} se hace referencia a la lesión física —como en el derecho penal, donde el delito de lesiones se clasifica en graves, levísimas y leves—. Se trata de un concepto jurídico distinto, vinculado a la existencia de una \textbf{desproporción entre las prestaciones} de un acto jurídico.
\end{definition}

En todo acto jurídico —especialmente en la compraventa— existe la expectativa de que las prestaciones sean equivalentes. Este principio se vincula con la denominada \textbf{justicia conmutativa}, entendida como la justicia entre particulares que busca una igualdad aritmética: que lo que una parte da sea igual a lo que recibe.

La pregunta central que atraviesa la construcción histórica del instituto es la siguiente: en principio, se supone que las partes son libres y pactan voluntariamente los términos de su acuerdo, de modo que no habría razón para invalidarlo si no existió error, dolo ni violencia. Sin embargo, cuando aparece una desproporción notable entre las prestaciones, surge el interrogante de en qué punto esa desproporción se vuelve tan grave como para dar lugar a la lesión.

A este elemento objetivo —la desproporción— se suma un segundo elemento, de naturaleza subjetiva: la explotación, por una de las partes, de la inferioridad de la otra, con el fin de obtener una ventaja patrimonial evidentemente desproporcionada.

\section{El derecho romano y el Código Justinianeo}

En el derecho romano existen dos rescriptos atribuidos a los emperadores Diocleciano y Maximino que habrían instituido la figura, aunque se ha discutido si se trata de textos insertados con posterioridad. Lo indudable es que el instituto de la lesión aparece recogido en el Código Justinianeo, la recopilación del ordenamiento jurídico romano elaborada bajo el Emperador Justiniano.

\begin{formula}{Código Justinianeo (Corpus Iuris Civilis)}
``Si tú o tu padre hubiera vendido por menor precio una cosa, es humano que, restituyendo todo el precio a los compradores, recobres el fondo vendido mediando la autoridad del juez, aunque si el comprador lo prefiere, reciban lo que les falta al justo precio. Pero se considera que el precio es menor si no se hubiera recibido ni la mitad del verdadero precio.''
\end{formula}

Este instituto recibió el nombre de \textbf{lesión enorme}. Era concedida al vendedor en la compraventa de inmuebles, y exigía que la desproporción alcanzara una medida determinada: más de la mitad del valor real de la cosa.

La lesión enorme daba lugar a dos acciones posibles:

\begin{itemize}
    \item una acción para recuperar el inmueble (el ``fundo''); o bien
    \item una acción para que el comprador tuviera la opción de pagar el complemento del precio faltante, es decir, un reajuste.
\end{itemize}

\begin{example}{Reajuste del precio en el derecho romano}
Si la cosa vale 100 y se pagó 40, el comprador —a quien le interesa conservar la cosa— puede optar por pagar los 60 que faltan, en lugar de restituir el fondo.
\end{example}

\section{La Edad Media y el precio justo}

Los glosadores, que continuaron comentando y ampliando el derecho romano, extendieron la aplicación de la lesión a la mayoría de los contratos durante la Edad Media. Los canonistas, por su parte, pusieron especial énfasis en la inmoralidad del desequilibrio entre las prestaciones.

Este factor moral resulta históricamente relevante, porque en torno a la lesión siempre ha oscilado la tensión entre una mirada objetiva —centrada en la desproporción del objeto— y un reproche dirigido a quien explota una situación ajena.

En este contexto adquiere especial relevancia el aporte de \textbf{Tomás de Aquino} y su idea del \textbf{precio justo}, un concepto de difícil determinación pero de gran influencia doctrinal: cuando el precio resulta manifiestamente injusto, se configura una posible lesión.

\begin{formula}{Tomás de Aquino — El precio justo}
Aun la ley humana —en referencia al derecho romano— considera ilícito que el vendedor venda una cosa por más de lo que vale, o que el comprador la adquiera por menos de su valor, aunque no lo tiene por ilícito si la diferencia no es excesiva. Se reputa ilícito cuando el engaño supera la mitad del valor de la cosa, en clara referencia a la métrica del Código Justinianeo.
\end{formula}

\section{La legislación española como antecedente del derecho argentino}

La evolución de la figura en las fuentes españolas —antecedente directo del derecho argentino— puede sintetizarse del siguiente modo:

\begin{table}[H]
\centering
\caption{Tratamiento de la lesión en la legislación española}
\begin{tabularx}{\textwidth}{>{\bfseries\raggedright\arraybackslash}p{0.28\textwidth} X}
\toprule
\textbf{Fuente española} & \textbf{Tratamiento de la lesión} \\
\midrule
Fuero Juzgo & Rigió durante la dominación visigoda. No admite la lesión, pero combate seriamente la usura, instituto muy cercano. \\
Fuero Real & Concede la acción al vendedor cuando lo engañan en más de la mitad del valor, si el comprador no acepta pagar la diferencia. \\
Las Leyes de Partidas & Consolidan la acción tanto para el vendedor como para el comprador. \\
Novísima Recopilación & Acepta la acción en compraventa, permuta, rentas y otros contratos semejantes. \\
\bottomrule
\end{tabularx}
\end{table}

\section{El derecho comparado moderno}

\subsection{El Código Civil Francés (1804): la lesión objetiva}

El Código Napoleónico constituye el ejemplo por excelencia de lo que se ha denominado \textbf{lesión objetiva}: aquella que contempla únicamente los elementos objetivos —es decir, la desproporción— sin tomar en consideración, en la configuración de la norma, el elemento subjetivo de la explotación de la inferioridad de la otra parte.

\begin{formula}{Código Civil Francés — art. 1.674 (1804)}
Cuando el vendedor resulte lesionado en más de siete doceavos del precio del inmueble, tendrá derecho a solicitar la rescisión de la venta, aunque en el contrato haya renunciado expresamente a solicitar dicha rescisión y aunque haya declarado donar las plusvalías.
\end{formula}

La medida establecida —más de 7/12— representa un poco más que la mitad (que serían 6/12). El código francés sigue así la métrica del derecho romano, pero la precisa mediante un porcentaje específico, tomando el valor al momento de la venta.

\begin{example}{Complemento del precio en el código francés}
Se concede al comprador —lesionante— la posibilidad de ofrecer el complemento del precio con la deducción de la décima parte. Si la cosa vale 100 y se pagaron 40, el comprador puede ofrecer el complemento hasta 90, pagando los 50 restantes, y de esa manera conservar la propiedad.
\end{example}

En el derecho argentino, esta acción de complemento recibe el nombre de \textbf{reajuste equitativo} y es fijada por el juez. En el código francés, la prescripción de la acción era de dos años.

\subsection{El Código Civil Alemán — BGB (1900): la lesión subjetiva}

En 1900, la reforma del código civil alemán incorpora por primera vez la fórmula de la denominada \textbf{lesión subjetiva}, estrechamente vinculada al combate de la usura. La doctrina coincide en señalar este momento como el surgimiento de dicha categoría.

\begin{formula}{Código Civil Alemán — BGB (síntesis, 1900)}
Un acto jurídico contrario a las buenas costumbres por el cual alguien, explotando la necesidad, la ligereza o la inexperiencia de otro, obtiene para sí o para un tercero que a cambio de una prestación le prometa o entregue ventajas patrimoniales que exceden de tal forma el valor de la prestación que, teniendo en cuenta las circunstancias, existe una desproporción chocante con ella.
\end{formula}

Aparecen aquí los elementos subjetivos del lesionante —que explota la necesidad, la ligereza o la inexperiencia— y se hace prometer o entregar ventajas patrimoniales. Existe, entonces, una acción activa: el lesionante toma la iniciativa para explotar una inferioridad. Del lado del lesionado, la inferioridad se define como necesidad, ligereza o inexperiencia, términos que luego reproducirá casi textualmente el art.\ 954 argentino y, con alguna variante, el código vigente.

\begin{important}
El código alemán declaraba el acto nulo en su totalidad, sin admitir nulidades parciales ni reajuste de precios, por tratarse de un acto contrario a las buenas costumbres. Existía un deber de restituir; la acción no era confirmable, era imprescriptible y podía ser declarada de oficio.
\end{important}

\subsection{El Código Suizo (1912) y el Código Italiano (1942)}

En 1912, el Código Suizo de las Obligaciones sigue al modelo alemán con variantes propias: coloca la prescripción en un año, torna la acción anulable y permite la confirmación del acto, acercándose así a la idea de un vicio que da lugar a una nulidad relativa susceptible de confirmación.

En 1942 se sanciona el Código Civil Italiano, que ejercerá una influencia decisiva sobre la reforma del art.\ 954 argentino en 1968.

\begin{formula}{Código Civil Italiano — art. 1.448 (1942, síntesis)}
Si hubiera desproporción en las prestaciones de una de las partes, y esa desproporción dependiese del estado de una de ellas, del cual se aprovechara la otra para obtener ventajas, la parte damnificada podrá demandar la rescisión del contrato. La acción no es admisible si la lesión es inferior a la mitad del valor que la prestación ejecutada o prometida por la parte perjudicada tenía en el momento del contrato. La lesión debe perdurar hasta el momento en que se proponga la demanda.
\end{formula}

\begin{note}
La exigencia de que la lesión perdure hasta el momento de la demanda no estaba presente en los códigos anteriores, y será recogida tanto por la ley 17.711 como por el CCC vigente. En el código italiano, además, los contratos aleatorios no pueden ser rescindidos por esta vía.
\end{note}

% ---------------------------------------------------------
% CORNELL — CAPÍTULO 1
% ---------------------------------------------------------
\section{Cuadro de repaso — Método Cornell}
\begin{table}[H]
\centering
\begin{tabularx}{\textwidth}{
    >{\raggedright\arraybackslash}p{0.30\textwidth}
    >{\raggedright\arraybackslash}X
}
\toprule
\textbf{Preguntas / palabras clave} & \textbf{Notas / respuestas} \\
\midrule

\textbf{¿Qué son los vicios propios de los actos jurídicos en el CCC?} &
Son tres: lesión, simulación y fraude. Están regulados en el capítulo sexto del título cuarto del Código Civil y Comercial, distribuidos en tres secciones. \\

\textbf{¿Qué es la lesión como vicio jurídico?} &
No es la lesión física, sino una desproporción entre las prestaciones de un acto jurídico, combinada con la explotación, por una de las partes, de la inferioridad de la otra. Se vincula con la idea de justicia conmutativa (igualdad aritmética entre lo que se da y lo que se recibe). \\

\textbf{¿Cuál es el origen histórico del instituto?} &
El Código Justinianeo (siglo VI) consagró la ``lesión enorme'' en la compraventa de inmuebles: si el precio era inferior a la mitad del valor real, el vendedor podía pedir la nulidad o el complemento del precio. \\

\textbf{¿Qué distinción clave introduce el BGB alemán de 1900?} &
Introduce la lesión subjetiva: ya no basta la desproporción objetiva, sino que se requiere además la explotación de la necesidad, ligereza o inexperiencia de la otra parte. Esta fórmula fue luego adoptada por Argentina en 1968 y rige, con variantes, en el art.\ 332 del CCC. \\

\midrule

\multicolumn{2}{p{\dimexpr\textwidth-2\tabcolsep\relax}}{
\textbf{Resumen:} La lesión, desde la lesión enorme del Código Justinianeo (más de la mitad del valor en la compraventa de inmuebles), pasando por el énfasis moral de la Edad Media y el precio justo de Tomás de Aquino, evolucionó hacia dos grandes modelos de derecho comparado: la \textbf{lesión objetiva} del código francés (que exige únicamente la desproporción) y la \textbf{lesión subjetiva} del BGB alemán (que agrega la explotación de la inferioridad de la otra parte), modelo este último que —junto con el código italiano de 1942— influirá decisivamente en la incorporación de la figura al derecho argentino.
} \\

\bottomrule
\end{tabularx}
\end{table}

% ===========================================================================
\chapter{La lesión en el derecho argentino}
% ===========================================================================

\section{El Código de Vélez (1869): la exclusión de la lesión}

Vélez Sarsfield no incluyó la lesión entre los vicios de los actos jurídicos. Al tratar los actos jurídicos en el art.\ 944, y luego de referirse a la simulación, el codificador dedicó una extensa nota —ubicada luego del art.\ 943— para explicar por qué no adoptaba la figura.

\begin{note}
Técnicamente, la nota no se refiere al contenido del art.\ 943 —que trataba sobre la violencia—, sino que Vélez la ubicó allí, antes del comienzo del tratamiento de los actos jurídicos, con el propósito de explicar en algún lugar del código por qué no incorporaba la lesión. Años más tarde, Borda señaló en un fallo de 1958 que una nota al pie no constituye ley.
\end{note}

Vélez ofrece dos grandes razones para no reconocer la lesión:

\begin{formula}{Nota de Vélez al art. 943 — Primer argumento}
``Para sostener nosotros que la lesión enorme, menor o mínima vicia los actos y obtenernos por tanto de proyectar disposiciones de la materia, basta comparar las diversas legislaciones y de las diferencias entre ellas resultados que no han tenido un principio uniforme al establecer esta ley.''
\end{formula}

El codificador observa que la legislación comparada presenta numerosas diferencias: unas la conceden al vendedor, otras al comprador; unas la aplican a la compraventa, otras a distintos actos; unas imponen una medida de más de la mitad, otras no. Ante esta disparidad de criterios, Vélez concluye que no encuentra un principio uniforme y que, por lo tanto, no se siente obligado a incorporar la figura.

\begin{formula}{Nota de Vélez al art. 943 — Segundo argumento}
``Dejaríamos de ser responsables de nuestras acciones si la ley nos permitiera enmendar todos nuestros errores o todas las consecuencias de ellos. El consentimiento libre prestado sin error, dolo y violencia y con la solemnidad requerida por las leyes debe hacer irrevocables los contratos.''
\end{formula}

Este segundo argumento expresa que lo pactado sin error, dolo ni violencia es válido, y que quien se equivoca debe hacerse cargo de las consecuencias de su propio error. Se advierte aquí una visión marcadamente individualista y liberal, coincidente con otras opciones adoptadas por Vélez a lo largo de todo el código.

\section{La jurisprudencia anterior a la ley 17.711}

Durante el siglo XX, distintos fallos judiciales fueron invalidando actos manifiestamente desproporcionados, invocando generalmente el art.\ 953 —referido al objeto del acto—, que disponía que el objeto no podía consistir en cosas o hechos contrarios a la moral. Cuando se pactaba algo contrario a la moral, ello podía dar lugar a la nulidad, de modo que el vicio de lesión se introducía por vía del art.\ 953.

La nota de Vélez al art.\ 943 seguía vigente, pero algunos jueces —entre ellos, Borda, en un fallo de la Cámara Nacional de Apelaciones de 1958 conocido como el caso ``Orlando''— sostuvieron que las notas no constituyen texto legal y que la lesión resulta incompatible con las buenas costumbres.

\begin{example}{Casos de aplicación jurisprudencial de la lesión (antes de 1968)}
Préstamos usurarios, honorarios exorbitantes del administrador de una sucesión, venta a precio vil (muy bajo) de un terreno, y moderación de cláusulas penales excesivas.
\end{example}

En 1961, el Tercer Congreso de Derecho Civil tuvo una importancia decisiva para la reforma de la ley 17.711, al formular recomendaciones que constituyeron, en gran medida, la base de la reforma de 1968.

\section{La ley 17.711 y la reforma al art. 954 (1968)}

La ley 17.711 incorporó diversas figuras al Código Civil —entre ellas, el abuso del derecho y modificaciones en los efectos de la declaración de voluntad en el tiempo— y también la lesión, incorporada en el art.\ 954.

\begin{formula}{Art. 954 — Ley 17.711 (1968, texto incorporado)}
``También podrá demandarse la nulidad o modificación de los actos jurídicos cuando una de las partes, explotando la necesidad, ligereza o inexperiencia de la otra, obtuviera por medio de ellos una ventaja patrimonial evidentemente desproporcionada y sin justificación. Se presume, salvo prueba en contrario, que existe tal explotación en caso de notable desproporción de las prestaciones. Los cálculos deben hacerse según valores al tiempo del acto y la desproporción deberá subsistir al momento de la demanda. Solo el lesionado o sus herederos podrán ejercer la acción cuya prescripción operará a los cinco años. El accionante tiene opción para demandar la nulidad o el reajuste equitativo del convenio, pero la primera de estas acciones se transformará en acción de reajuste si éste fuera ofrecido por el demandado.''
\end{formula}

En este texto se incorporan claramente un elemento objetivo —la desproporción— y un elemento subjetivo —la explotación de la necesidad, ligereza o inexperiencia—. El nuevo código sustituirá luego la palabra ``ligereza'' por ``debilidad psíquica''. Cabe destacar que el texto del art.\ 954 no establece una métrica fija para medir el elemento objetivo.

\section{El art. 332 del CCC: texto y diferencias con el art. 954}

El Código Civil y Comercial vigente sigue, en lo sustancial, al art.\ 954 para regular la lesión, en un único artículo ubicado en la sección primera del capítulo sexto del título cuarto.

\begin{important}
Aunque se trate de un único artículo, el art.\ 332 contiene un desarrollo normativo extenso, que exige un análisis detenido de cada uno de sus elementos.
\end{important}

\begin{formula}{Art. 332 — Código Civil y Comercial de la Nación (vigente)}
``Puede demandarse la nulidad o modificación de los actos jurídicos cuando una de las partes, explotando la necesidad, debilidad psíquica o inexperiencia de la otra, obtuviera por medio de ellos una ventaja patrimonial evidentemente desproporcionada y sin justificación. Se presume, salvo prueba en contrario, que existe tal explotación en casos de notable desproporción de las prestaciones. Los cálculos deben hacerse según valores al tiempo del acto y la desproporción debe subsistir al momento de la demanda. El afectado tiene opción para demandar la nulidad o un reajuste equitativo del convenio, pero la primera de estas acciones se transforma en acción de reajuste si éste es ofrecido por el demandado al contestar la demanda. Solo el lesionado o sus herederos podrán ejercer la acción.''
\end{formula}

Las diferencias del CCC respecto del art.\ 954 pueden sintetizarse así:

\begin{table}[H]
\centering
\caption{Diferencias entre el art. 954 (ley 17.711) y el art. 332 (CCC)}
\begin{tabularx}{\textwidth}{>{\bfseries\raggedright\arraybackslash}p{0.20\textwidth} X}
\toprule
\textbf{Modificación} & \textbf{Explicación} \\
\midrule
Cambio 1 & La palabra ``ligereza'' es reemplazada por ``debilidad psíquica''. El nuevo código clarifica así un punto que había dado lugar a un problema interpretativo, entre quienes ubicaban la ligereza como un obrar irreflexivo sin mayores requisitos y quienes la entendían como expresión de una cierta debilidad psíquica; el CCC se pronuncia por esta última postura. \\
Cambio 2 & El plazo de prescripción se reduce de cinco a dos años. Metodológicamente resulta correcto tratar el plazo en el libro dedicado a las prescripciones, y no dentro del propio artículo. \\
\bottomrule
\end{tabularx}
\end{table}

\section{Ámbito de aplicación del art. 332}

El art.\ 332 recoge, en principio, todos los actos jurídicos en general. Sin embargo, el acto debe ser bilateral y no a título gratuito, puesto que en los actos a título gratuito siempre existe una desproporción entre las prestaciones que resulta ajena a la lógica de la lesión.

\section{El elemento objetivo: la desproporción de las prestaciones}

El primer elemento exigido por la norma es el elemento objetivo: la desproporción de las prestaciones. A diferencia del código francés, que establece un criterio fijo (7/12 partes), el CCC habla de ``ventaja evidentemente desproporcionada'' y de ``notable desproporción'', sin fijar un criterio numérico. No obstante, se exige que la diferencia sea muy notable: no bastaría, por ejemplo, una diferencia del veinte por ciento.

Un requisito adicional es que la desproporción debe ser \textbf{sin justificación}.

\begin{example}{Desproporción con justificación}
Un tío que va a asistir al casamiento de un sobrino le vende un departamento que vale 100.000 dólares en 40.000, como regalo de casamiento. Tras su muerte, otros herederos —primos del beneficiado— demandan por lesión, invocando la diferencia entre el valor real y el precio pagado. El demandado se defenderá alegando que existió una justificación: la intención de realizar una liberalidad. Podrá discutirse si ese regalo excede o no la porción disponible —cuestión de legítima—, pero no podrá hablarse de lesión en sentido propio, porque no hubo explotación sino una justificación para la desproporción.
\end{example}

Los cálculos deben realizarse tomando en cuenta los valores al momento del acto, puesto que todo vicio afecta al acto desde su origen. Sin embargo, la desproporción debe además subsistir al momento de la demanda, exigencia que proviene del código italiano.

\begin{example}{Desproporción que no subsiste}
Una persona compra muy barato un inmueble en un barrio cotizado. La desproporción existe al momento del acto, pero luego el valor de la propiedad desciende notablemente porque el barrio se deteriora y se vuelve inhabitable. Como la desproporción no subsiste al momento de la demanda, no existe acción por lesión, aun cuando el vicio se hubiera configurado al momento del acto.
\end{example}

\section{Los elementos subjetivos}

El art.\ 332 exige, además del elemento objetivo, la \textbf{explotación} por parte del lesionante. Se trata de algo más intenso que un simple ``aprovechar'' —entendido como que la otra parte no advierte la desproporción y el lesionante simplemente continúa adelante—: existe aquí una acción más activa, consistente en ``hacerse dar o prometer'' aprovechando la inferioridad de la otra parte. Este elemento resulta, en la práctica, de difícil prueba.

Del lado del lesionado, el CCC exige que se configure alguno de los siguientes tres estados de inferioridad:

\begin{table}[H]
\centering
\caption{Estados de inferioridad del lesionado (art. 332 CCC)}
\begin{tabularx}{\textwidth}{>{\bfseries\raggedright\arraybackslash}p{0.22\textwidth} X}
\toprule
\textbf{Estado de inferioridad} & \textbf{Definición y aclaración} \\
\midrule
Necesidad & Falta de lo indispensable para la vida. No basta con alegar, por ejemplo, que se tenía necesidad de asistir a la final de un mundial y por ello se vendió una propiedad muy barata. La necesidad exige una situación urgente, vital e indispensable —como cubrir una operación médica de un familiar—, no un deseo o conveniencia. \\
Debilidad psíquica & Presencia de alguna alteración mental. Reemplaza a la palabra ``ligereza'' del código anterior, que había generado dos posturas doctrinarias: obrar irreflexivo sin mayores requisitos, o expresión de una cierta debilidad psíquica. El nuevo código se pronuncia por esta última. \\
Inexperiencia & Falta del conocimiento que otorgan la vida y las relaciones. Quien se dedica profesionalmente al mercado inmobiliario durante treinta años no puede alegar inexperiencia en esa materia; distinta es la situación de una persona que se traslada de otro lugar, desconoce el ambiente y carece de formación en la materia del acto. \\
\bottomrule
\end{tabularx}
\end{table}

\begin{important}
En general se interpreta que estos tres supuestos de inferioridad constituyen una enumeración taxativa: no puede invocarse otro distinto. El proyecto de reforma del CCC del año 2018 proponía ampliar estos supuestos, pero en el derecho vigente son únicamente tres: necesidad, debilidad psíquica e inexperiencia.
\end{important}

\section{La presunción legal del art. 332}

El art.\ 332, al igual que su antecedente el art.\ 954, establece que se presume —salvo prueba en contrario— la existencia de explotación en caso de notable desproporción de las prestaciones.

\begin{important}
Esta presunción constituye una inversión de la carga de la prueba: si el lesionado prueba la notable desproporción de las prestaciones, se presume que hubo explotación, y corresponderá entonces al lesionante —demandado— probar que no existió tal explotación. Existe así una preponderancia práctica del elemento objetivo dentro de una norma que, en su configuración, es de lesión subjetiva.
\end{important}

Parte de la doctrina discute si el término ``notable'' equivale a ``evidentemente desproporcionada''. En general, se considera que lo notable es aquello que salta a la vista sin necesidad de una investigación mayor.

Existe además una discusión acerca del alcance de la presunción: una postura sostiene que solo alcanza al elemento subjetivo del lesionante, debiendo probarse por separado la necesidad, debilidad psíquica o inexperiencia del lesionado; otra postura entiende que, probada la desproporción, ambos elementos subjetivos se presumen.

\begin{note}
Esta presunción legal constituye una novedad propia del derecho argentino, sin equivalente en el código italiano, el código suizo ni el código alemán, y es la que otorga al elemento objetivo su especial preponderancia práctica.
\end{note}

\section{Las acciones, la legitimación y la prescripción}

La lesión da lugar a dos acciones:

\begin{enumerate}
    \item \textbf{Acción de nulidad}, de carácter relativo —susceptible de confirmación—, en tanto está establecida para la protección de la parte lesionada, al igual que ocurre con los restantes vicios.
    \item \textbf{Reajuste equitativo}, mediante el cual el juez corrige la desproporción en el caso concreto, determinando lo que resulte justo. A diferencia del esquema del código francés, que deducía una décima parte para el pago del complemento del precio, en el derecho argentino el reajuste queda librado a la valoración judicial, sin que la ley fije una pauta.
\end{enumerate}

\begin{important}
El reajuste equitativo prima sobre la acción de nulidad cuando es ofrecido por el demandado al contestar la demanda, en virtud del principio de conservación de los actos jurídicos.
\end{important}

La legitimación para iniciar la acción corresponde exclusivamente al lesionado o a sus herederos, tal como establecía ya el art.\ 954. El plazo de prescripción es de \textbf{dos años}, contado desde la fecha en que la obligación debía cumplirse —plazo que la ley 17.711 fijaba en cinco años.

\section{Figuras afines: cláusula penal y usura}

El art.\ 794, relativo a la cláusula penal, guarda semejanza con la lesión: los jueces pueden reducir las penas establecidas en una cláusula penal —cláusula contractual que fija sanciones para el caso de incumplimiento— cuando resulten excesivas. Se trata de una figura próxima a la lesión, en cuanto también existe desproporción susceptible de reducción judicial, aunque con una diferencia relevante: no produce nulidad, sino únicamente reajuste.

La \textbf{usura}, en cambio, constituye un delito en el derecho penal, regulado por el art.\ 175 bis del Código Penal.

\begin{formula}{Art. 175 bis — Código Penal (síntesis) — Usura}
Quien, aprovechando la necesidad, ligereza o inexperiencia de otro, se hiciera dar o prometer, en cualquier forma, para sí o para un tercero, intereses u otras ventajas pecuniarias evidentemente desproporcionadas con su prestación, u otorgare a sabiendas un crédito usurario, será reprimido.
\end{formula}

En la figura penal se requiere dolo. Debe notarse que esta norma no utiliza el término ``explotación'' sino ``aprovechar'', y exige que el autor se haga dar o prometer las ventajas: existe una acción activa dirigida a lograr esa entrega o promesa, que constituye el núcleo del delito.

\section{Debate doctrinario: la esencia de la lesión}

Existe un debate doctrinario acerca de cuál es la esencia de la lesión. Tobías, en su Tratado de Derecho Civil, identifica tres posturas posibles:

\begin{table}[H]
\centering
\caption{Posturas doctrinarias sobre la naturaleza jurídica de la lesión}
\begin{tabularx}{\textwidth}{>{\bfseries\raggedright\arraybackslash}p{0.28\textwidth} X}
\toprule
\textbf{Postura} & \textbf{Contenido} \\
\midrule
Vicio del objeto & Enfatiza el elemento objetivo: existe un problema en el objeto del acto, que resulta desproporcionado. Para Tobías, esta postura ignora la presencia de la explotación de la inferioridad. \\
Vicio de la voluntad & Enfatiza el elemento subjetivo del lesionado: por su inferioridad —necesidad, debilidad psíquica o inexperiencia— estaría viciada su voluntad. Tobías objeta que, de seguirse este criterio, resultaría irrelevante la existencia de desproporción. \\
Comportamiento reprochable (Tobías) & Es una respuesta del ordenamiento frente a un comportamiento reprochable, vinculado a la buena fe. Se reprocha la explotación. Esta postura se encuentra más cerca del antecedente alemán. \\
\bottomrule
\end{tabularx}
\end{table}

En definitiva, los tres elementos —objeto, voluntad y conducta— están en juego a la hora de determinar la esencia de la lesión, y de fondo siempre se encuentra la justicia conmutativa. El elemento objetivo resulta indudable: el objeto del acto debe ser proporcionado, y si existe desproporción, el acto no realiza plenamente la justicia. Al mismo tiempo, existe una valoración desde las buenas costumbres, que fundamenta el reproche jurídico cuya consecuencia es la nulidad.

\begin{note}
En síntesis: el vicio de lesión en el ordenamiento argentino se recorrió desde el Código de Vélez —que no lo contempló, por las razones expuestas en su nota al art.\ 943— pasando por la jurisprudencia posterior y la reforma introducida por la ley 17.711 en el art.\ 954, cuyo texto es sustancialmente reproducido por el art.\ 332 del CCC vigente. Dentro de esta figura se analizaron su fuente, su método, su ámbito de aplicación, sus elementos, las acciones a las que da lugar, la legitimación activa, el plazo de prescripción, su comparación con otras figuras afines y el problema de su esencia jurídica.
\end{note}

% ---------------------------------------------------------
% CORNELL — CAPÍTULO 2
% ---------------------------------------------------------
\section{Cuadro de repaso — Método Cornell}
\begin{longtable}{
    >{\raggedright\arraybackslash}p{0.30\textwidth}
    >{\raggedright\arraybackslash}p{0.58\textwidth}
}
\toprule
\textbf{Preguntas / palabras clave} & \textbf{Notas / respuestas} \\
\midrule
\endhead

\textbf{¿Por qué Vélez no incluyó la lesión?} &
Por dos razones: (1) la disparidad de criterios en el derecho comparado, que le impedía identificar un principio uniforme; (2) su visión liberal-individualista, según la cual, si no hubo error, dolo ni violencia, el contrato es válido y cada parte debe hacerse cargo de sus propias decisiones. \\

\textbf{¿Cómo incorporó Argentina la lesión antes de 1968?} &
La jurisprudencia —especialmente el fallo de Borda en el caso ``Orlando'' (1958)— la aplicó por vía del art.\ 953 (objeto contrario a la moral). El Tercer Congreso de Derecho Civil (1961) recomendó luego su incorporación expresa. \\

\textbf{¿Qué exige el art. 332 del CCC para configurar la lesión?} &
Un elemento objetivo: desproporción evidente y sin justificación, que exista al momento del acto y subsista al momento de la demanda. Dos elementos subjetivos: explotación por parte del lesionante, y necesidad, debilidad psíquica o inexperiencia del lesionado. \\

\textbf{¿Cómo opera la presunción del art. 332?} &
Si se prueba la notable desproporción de las prestaciones, se presume —salvo prueba en contrario— la explotación. Esto invierte la carga probatoria: es el demandado quien debe demostrar que no explotó. \\

\textbf{¿Qué acciones otorga la lesión?} &
Nulidad relativa (confirmable) o reajuste equitativo (fijado por el juez). El reajuste prima sobre la nulidad cuando es ofrecido por el demandado al contestar la demanda, en virtud del principio de conservación del acto. \\

\textbf{¿Cuál es el plazo de prescripción?} &
Dos años, contados desde que la obligación debía ser cumplida. El CCC redujo el plazo de cinco años que establecía la ley 17.711. \\

\textbf{¿Cuál es la naturaleza jurídica de la lesión según la doctrina?} &
Existe debate entre tres posturas: (1) vicio del objeto; (2) vicio de la voluntad; (3) conducta reprochable, contraria a la buena fe y a las buenas costumbres. Tobías se inclina por esta tercera postura, más próxima al modelo alemán. \\

\midrule

\multicolumn{2}{p{0.88\textwidth}}{
\textbf{Resumen:} Argentina, luego de que Vélez la excluyera por razones liberales y por la falta de un criterio uniforme en el derecho comparado, incorporó la lesión en 1968 mediante la ley 17.711, siguiendo el modelo alemán e italiano. El CCC vigente la regula en el art.\ 332 con tres elementos: uno objetivo (desproporción evidente, sin justificación, que subsista al tiempo de la demanda), uno subjetivo del lesionante (explotación activa) y uno subjetivo del lesionado (necesidad, debilidad psíquica o inexperiencia). La presunción legal de explotación ante la notable desproporción —novedad argentina sin equivalente en los códigos extranjeros— otorga preponderancia práctica al elemento objetivo. Las acciones disponibles son la nulidad relativa y el reajuste equitativo, prevaleciendo este último por el principio de conservación del acto. El plazo de prescripción es de dos años.
} \\

\bottomrule
\end{longtable}
\chapter[La simulación]{La simulación como vicio de los actos jurídicos}
\label{ch:simulacion}

Tanto la simulación como el fraude son vicios de los actos jurídicos que hacen que un acto, pese a existir, no produzca todos sus efectos normales. La diferencia entre ambos puede describirse como la que existe entre una fachada falsa y una casa vaciada por dentro: en la simulación se construye una fachada --un acto aparente-- que oculta o no oculta nada real, y por eso se ataca con la nulidad; en el fraude la casa es real, pero el deudor la vació de bienes para no pagarles a sus acreedores, y por eso se ataca con la inoponibilidad y no con la nulidad. Este capítulo trata la simulación; el siguiente, el fraude, y al final se comparan ambas figuras.

\begin{note}
Estos dos vicios explican por qué un acto que parece perfecto en el papel puede terminar siendo anulado o simplemente ignorado por la ley frente a ciertas personas. Para cualquier abogado --litigante, asesor de empresas, escribano o juez-- identificar una simulación (por ejemplo, una venta que en realidad esconde una donación para perjudicar herederos) o un fraude (un deudor que vacía su patrimonio para no pagar) constituye una herramienta central de la práctica profesional: de estos institutos depende poder recuperar un bien familiar, cobrar una deuda, defender a un cliente acusado de haber ``prestado su nombre'', o asesorar correctamente antes de firmar un contrato para no quedar involucrado, sin saberlo, en un acto simulado o fraudulento.

También resulta útil en la vida personal: son temas que aparecen en sucesiones familiares, en préstamos entre parientes, en la compra de bienes a nombre de terceros o en cualquier situación donde alguien pueda intentar aprovecharse de un acuerdo informal. Entender la diferencia entre lo lícito y lo ilícito, y saber qué prueba se necesita para desarmar un acto aparente, protege el propio patrimonio y el de la familia.

\end{note}

\section{Concepto y noción general}

La simulación es un vicio especial de los actos jurídicos, porque supone una situación en la que las partes acuerdan una declaración de voluntad o un acto jurídico que, o bien no es real, o bien es aparente y encubre otro acto verdadero.

La simulación puede tener fines lícitos o ilícitos, según las circunstancias en las que las personas involucradas la lleven a cabo.

\begin{example}{Simulación ilícita entre herederos}
Un padre o una madre --viudo o viuda-- tiene tres hijos y quiere beneficiar en realidad a uno de ellos, su hijo favorito, simulando venderle su única casa. Los otros dos hijos no se enteran del verdadero propósito del acto.

Cuando fallece la persona y los herederos inician la sucesión, se encuentran con que la casa figura vendida, pero el hijo que aparece como comprador en realidad no tenía fondos, pues atravesaba una situación económica muy mala. Ante la pregunta de los demás hermanos, queda claro que el acto tuvo que haber sido simulado: en apariencia se le vendió la casa, pero en realidad se la estaban regalando, de modo que se trató de una donación.

Frente a esto, se ataca el acto para que se declare su nulidad y se lo considere una donación, de modo que la porción de la legítima hereditaria que corresponde a los demás herederos reingrese al acervo hereditario. Este es un ejemplo de simulación \textbf{ilícita}, porque perjudica a terceros --en este caso, a los otros coherederos.
\end{example}

\begin{example}{Simulación lícita: el auto a nombre del hermano}
Una persona que se muda al exterior por tiempo indeterminado acuerda con su hermano poner el automóvil a su nombre, simulando una venta, para que este lo utilice mientras dure la ausencia, con el compromiso de restituirlo a su regreso.

La titularidad del auto queda simulada: el hermano actúa todo el tiempo como si el auto fuera suyo, y cuando la persona regresa y lo reclama, se produce la necesidad de volver las cosas al estado anterior.
\end{example}

\section{Definición doctrinaria y ubicación metodológica}

La doctrina define la simulación como la situación en la que dos personas, o una de ellas y un tercero que aparentemente no interviene en el acto, se ponen de acuerdo para declarar que celebran actos jurídicos que en realidad no se han propuesto llevar a cabo, o para aparentar un acto distinto, ocultando el acto verdadero bajo la apariencia de otro que declaran haber realizado.

\begin{definition}{Simulación}
Habrá simulación cuando dos personas, o una de ellas y un tercero que aparentemente no interviene en el acto, se ponen de acuerdo para declarar que celebran actos jurídicos que en realidad no se han propuesto llevar a cabo, o para aparentar un acto distinto, ocultando el acto verdadero bajo la apariencia de otro que declaran haber realizado.
\end{definition}

A partir de esta definición surgen los dos grandes tipos de simulación: la \textbf{absoluta}, cuando no hay nada de real, y la \textbf{relativa}, cuando la simulación encubre otro acto.

Metodológicamente, el tema se ubica en la Sección Segunda del Capítulo 6, ``De los vicios de los actos jurídicos'', del Título Cuarto (``Hechos y actos jurídicos'') del Libro Primero, Parte General, del Código Civil y Comercial (CCyC).

\section{Definición legal de la simulación}

La definición legal que trae el Código Civil y Comercial viene prácticamente igual a la que traía el código anterior en su artículo 955.

\begin{formula}{Definición legal de la simulación}
``La simulación tiene lugar cuando se encubre el carácter jurídico de un acto con la apariencia de otro, o cuando el acto contiene cláusulas que no son sinceras, o fechas que no son verdaderas, o cuando por el acto se constituyen o transmiten derechos a personas interpuestas que no son aquellas para quienes en realidad se constituyen o transmiten.''
\end{formula}

\section{Simulación absoluta y simulación relativa}

Dentro de esta clasificación se distingue, por un lado, la simulación absoluta y, por otro, la simulación relativa.

\subsection{Simulación absoluta}

Es absoluta cuando el acto nada tiene de verdadero.

\begin{example}{Simulación absoluta}
Una madre simula venderle un bien a un hijo para eludir a sus acreedores. Lo ``vende'', pero sigue viviendo allí: no hay ninguna transferencia de dinero, no hay nada real. El bien sigue siendo de ella; lo único que hace es aparentar haberlo puesto a nombre de su hijo. Esto constituye una simulación absoluta.
\end{example}

\subsection{Simulación relativa}

Es relativa aquella en que el acto aparente encubre a otro real y verdadero.

\begin{formula}{Art. 334 CCyC}
``Si el acto simulado encubre a otro real, este es plenamente eficaz, siempre que concurran los requisitos propios de esa categoría de acto, y no sea ilícito ni perjudique a un tercero.''
\end{formula}

\begin{example}{Simulación relativa}
Se simula una compraventa, pero en realidad se trata de una donación, sin que esto perjudique a terceros ni sea ilícito. Será una donación hecha por determinados fines, que aparentemente se hizo como compraventa, pero en realidad es una donación.
\end{example}

El artículo 334 deja a salvo el acto encubierto: el acto aparente se anula y queda subsistente el acto encubierto, plenamente eficaz.

\section{Simulación lícita e ilícita}

La otra clasificación relevante distingue entre simulación lícita e ilícita: la simulación \textbf{lícita} es aquella que no perjudica a terceros ni viola la ley; la simulación \textbf{ilícita} es aquella que perjudica a terceros o viola la ley.

Retomando los ejemplos ya expuestos: el caso del deudor que simula vender un bien para que sus acreedores no puedan ejecutarlo es ilícito, porque los perjudica. En cambio, el caso de la persona que pone el auto a nombre de su hermano, sin perjudicar a nadie ni hacer nada contrario a la ley --simplemente porque resulta más práctico, ya que el hermano lo va a usar y asume todas las responsabilidades del caso, con el compromiso de devolverlo--, constituye una simulación lícita.

\section{Naturaleza de la sanción: nulidad relativa}

La simulación da lugar a la nulidad, y se trata de una nulidad de tipo \textbf{relativa}, porque está establecida en beneficio de las partes.

Entre los vicios que se tratan en el Título Cuarto del Código --error, dolo, violencia (vicios de la voluntad), vicios de los actos voluntarios, lesión y simulación-- todos dan lugar a la nulidad. El fraude es el único vicio de ese título que da lugar a la inoponibilidad.

Algunos autores sostenían que, ante una simulación, se estaba frente a un acto inexistente, pero esta postura no encuentra eco en el Código Civil y Comercial, que siempre refiere a la nulidad del acto.

\begin{important}
Si al leer algún texto o algún autor se encuentra una referencia a la inexistencia del acto simulado, en realidad esto no tiene asidero: la inexistencia es una categoría de ineficacia reconocida en otros países, pero en el derecho argentino no está reconocida, salvo en algún supuesto muy específico. Por lo tanto, al hablar de simulación siempre corresponde hablar de nulidad.
\end{important}

\section{La acción de simulación entre las partes}

La acción de simulación puede ser ejercida por las partes entre sí, o por terceros.

\begin{example}{Acción entre las partes}
Una persona regresa de un viaje y reclama a su hermano la restitución del auto que le había puesto a su nombre de manera simulada, comprometiéndose este a devolvérselo a su regreso. El hermano desconoce que el acto haya sido simulado y sostiene que el auto es suyo.

Frente a ese desconocimiento, la persona debe iniciar una acción de simulación contra su hermano ante el juez, solicitando que se declare la nulidad del acto simulado, para que las cosas vuelvan al estado anterior. Este es un ejemplo de acción entre las propias partes.
\end{example}

\subsection{Procedencia de la acción entre partes}

\begin{formula}{Art. 335 CCyC (primer párrafo)}
``Los que otorgan un acto simulado ilícito, o que perjudica a terceros, no pueden ejercer acción alguna el uno contra el otro sobre la simulación, excepto que las partes no puedan obtener beneficio alguno de las resultas del acto practicado con la simulación.''
\end{formula}

Si la simulación es ilícita, la acción entre las partes no procede, por regla, salvo que ninguna de ellas pueda obtener un beneficio de la simulación. En cambio, en la simulación lícita no está presente este requisito, y la acción procede siempre.

\begin{example}{Excepción a la regla}
Una persona simuló vender su casa para eludir a sus acreedores y, luego, gracias a una herencia de un pariente lejano, obtiene dinero suficiente para pagarles. Con esa disponibilidad, inicia una acción de simulación contra el adquirente simulado de su casa, para recuperar la propiedad. Aunque la simulación fue ilícita, al dejar sin efecto el acto simulado no obtiene ningún beneficio indebido, porque ya pagó a los acreedores: estos no perdieron nada, no fueron perjudicados. Por eso, en este caso, la acción entre partes puede proceder.
\end{example}

\subsection{La prueba entre las partes: el contradocumento}

Entre las partes, la simulación debe probarse a través del \textbf{contradocumento}. El contradocumento es el instrumento donde consta la verdadera intención de las partes: si la simulación es una declaración de voluntad aparente que encubre otro acto real, o que no tenía nada de real, en el contradocumento se plasma la voluntad real, y las partes dejan en claro que se trató de un acto simulado. Durante la simulación se redacta también un contradocumento, justamente para que, ante la necesidad de probar que el acto es simulado, la parte interesada pueda exhibirlo.

En la redacción originaria de Vélez, de 1871 (vigente desde 1869), no había posibilidad de eximirse de presentar el contradocumento. Luego, la jurisprudencia fue incorporando excepciones, y en el período 1968/1969 la Ley 17.711 incorporó un párrafo al código anterior que estableció la posibilidad de prescindir del contradocumento. Esta solución fue receptada por el nuevo Código en el artículo 335, que agrega, respecto del código anterior, la exigencia de justificar las razones por las cuales el contradocumento no existe o no puede ser presentado.

\begin{formula}{Art. 335 CCyC (segundo párrafo)}
``La simulación alegada por las partes debe probarse mediante el respectivo contradocumento. Puede prescindirse de él cuando la parte justifica las razones por las cuales no existe o no puede ser presentado y median circunstancias que hacen inequívoca la simulación.''
\end{formula}

\begin{example}{Prescindencia del contradocumento}
Dos hermanos, por la cercanía y confianza entre ellos, no consideraron necesario escribir un contradocumento, o bien este se destruyó --por ejemplo, en un incendio-- y no puede ser presentado. En ese caso deben acreditarse circunstancias que hagan inequívoca la simulación: por ejemplo, demostrar que el hermano titular registral tenía bienes que en realidad seguía usando el otro hermano --el departamento, el automóvil--, de modo de dar al juez la convicción inequívoca de que efectivamente se está ante un acto simulado.
\end{example}

Complementariamente, el artículo 1028 del código anterior establecía que el contradocumento particular que altere lo expresado en un instrumento público puede invocarse entre las partes, pero no es oponible respecto de terceros.

\begin{important}
Si el acto simulado se instrumentó por escritura pública, puede utilizarse un contradocumento simplemente particular, es decir, que no sea escritura pública. Aunque la escritura pública tiene mayor fuerza probatoria que el instrumento particular, el contradocumento particular puede ser invocado entre las partes si tiene firma, ya que esta le da mayor certeza. Frente a terceros, en cambio, es inoponible. Si no tiene firma, habrá que probar las circunstancias que hacen inequívoca la simulación.
\end{important}

\section{La acción de simulación ejercida por terceros}

\subsection{Procedencia}

La acción ejercida por terceros procede cuando la simulación es ilícita. El tercero no tendría ningún interés en atacar una simulación lícita; la razón de su acción reside, precisamente, en la existencia de una simulación ilícita.

\begin{formula}{Art. 336 CCyC}
``Los terceros cuyos derechos o intereses legítimos son afectados por el acto simulado pueden demandar su nulidad.''
\end{formula}

Esta norma habilita a los terceros a iniciar una demanda de nulidad; se trata, nuevamente, de una demanda de nulidad, y no de inexistencia ni de inoponibilidad.

\subsection{La prueba por indicios}

Según el artículo 336, el tercero puede acreditar la simulación por cualquier medio de prueba: no se le exige el contradocumento, ya que este lo tienen las partes que hicieron el acto simulado, y no el tercero.

Ante la imposibilidad de recurrir al contradocumento, el tercero recurre a indicios. La doctrina, con aportes de distintos autores especializados en el tema, reúne una enumeración de indicios:

\begin{itemize}
    \item Carencia de medios económicos del adquirente para haber podido pagar el precio.
    \item El parentesco muy próximo, o la amistad íntima, entre las partes del acto.
    \item La ausencia de ejecución material del acto o contrato (por ejemplo, el vendedor que sigue viviendo en el inmueble ``vendido'').
    \item La venta de la cosa a una sociedad anónima en la que el propio enajenante o sus parientes son accionistas.
    \item La venta de un bien que el vendedor necesita para su sustento (por ejemplo, el auto con el cual trabaja y se traslada).
    \item El precio vil, es decir, un precio muy bajo en relación con el valor real de la cosa.
    \item La ausencia de incorporación del bien en las declaraciones impositivas (no haberlo denunciado ante la AFIP).
    \item Los pagos anticipados: pagar por anticipado todo el precio antes de la escrituración resulta sospechoso.
    \item La retroventa: es poco habitual que alguien compre algo con la posibilidad de que se lo revendan a él mismo.
    \item Que la enajenación no sea útil ni necesaria: no había ninguna necesidad real de vender.
\end{itemize}

Estos son indicios surgidos de la jurisprudencia, que el tercero que alega una simulación debe reunir para convencer a un juez de que lo sucedido no es real, sino que encubre otro acto, o que directamente no tiene nada de real.

\section{Efectos de la simulación frente a terceros}

La simulación ilícita que perjudica a un tercero provoca la nulidad del acto; si el acto simulado encubre a otro real, este es plenamente eficaz si concurren los requisitos propios de su categoría, y no es ilícito ni perjudica a terceros. En las simulaciones relativas por cláusulas simuladas, si es posible distinguir y separar la cláusula falsa de las demás, esta pierde valor, pero el acto sigue siendo válido en lo restante, por el principio de conservación de los actos.

\begin{formula}{Art. 337 CCyC}
``La simulación no puede oponerse a los acreedores del adquirente simulado que de buena fe hayan ejecutado los bienes comprendidos en el acto. La acción del tercero contra los subadquirentes solo procede si adquirió por título gratuito, o si es cómplice en la simulación. El subadquirente de mala fe y quien contrató de mala fe con el deudor responden solidariamente por los daños causados al tercero que ejerció la acción, si los derechos se transmitieron a un adquirente de buena fe y a título oneroso, o de otro modo se hizo imposible su recaudo.''
\end{formula}

\begin{example}{Efectos frente a terceros y subadquirentes}
Una persona simula venderle el auto a su hermano; este, a su vez, tiene sus propios acreedores, que ejecutan el auto de buena fe, sin saber nada de la simulación. En ese caso quedan protegidos, pues la lógica siempre es proteger al de buena fe y a título oneroso.

Si el hermano, a su vez, transmitió el auto a un subadquirente, la acción contra este solo procede si adquirió a título gratuito, o si fue cómplice de la simulación. Si el subadquirente es de buena fe y a título oneroso --por ejemplo, compró el vehículo mediante un aviso, sin conocer la simulación-- queda protegido. Si, en cambio, el hermano se puso de acuerdo con esa otra persona y le vendió el auto sabiendo todo, la acción procede, aunque la complicidad resulta muy difícil de demostrar.

El subadquirente de mala fe y el deudor responden por los daños: procede una acción de daños junto con la nulidad, dirigida contra el subadquirente de mala fe. En cambio, si el subadquirente es de buena fe pero a título gratuito --por ejemplo, si el bien fue donado a una institución de beneficencia-- no responde por los daños, sino en la medida en que se haya enriquecido, debiendo restituir en esa medida.
\end{example}

La regla general que atraviesa este tema es que siempre se protege a quien actuó de buena fe y a título oneroso; quien actuó de mala fe responde por los daños, y quien actuó de buena fe pero a título gratuito responde solo en la medida de su enriquecimiento. Se trata de un principio fundamental en todo el tema de los actos jurídicos dentro del sistema de derecho civil argentino.

\section{Prescripción de la acción de simulación}

\begin{formula}{Art. 2563, inc. b, CCyC}
El plazo de prescripción de dos años se computa, en la acción de simulación ejercida por parte, ``desde que, requerida una de ellas, se negó a dejar sin efecto el acto simulado''; y en la acción de simulación que ejerce un tercero, desde que conoció o pudo conocer el vicio del acto jurídico.
\end{formula}

La acción entre partes no comienza a correr desde que se celebró el acto simulado, sino desde que una de ellas desconoce, en los hechos, la simulación.

\begin{example}{Cómputo del plazo entre partes}
Una persona regresa de España, reclama el auto y su hermano le dice que el auto es suyo. Aunque hayan pasado varios años desde su partida, el plazo de prescripción empieza a correr recién cuando regresa y el hermano se niega a devolvérselo, y no desde el momento de la partida.
\end{example}

La acción ejercida por un tercero prescribe a los dos años desde que este conoció o pudo conocer el vicio del acto jurídico.

% -----------------------------------------------------------
% CORNELL --- CAPÍTULO 1
% -----------------------------------------------------------
\section{Repaso Cornell: la simulación}
\begin{table}[H]
\centering
\begin{tabularx}{\textwidth}{
    >{\raggedright\arraybackslash}p{0.28\textwidth}
    >{\raggedright\arraybackslash}X
}
\toprule
\textbf{Preguntas / palabras clave} &
\textbf{Notas / respuestas} \\
\midrule

\textbf{¿Qué es la simulación?} &
Un vicio en el que las partes acuerdan una declaración de voluntad que no es real, o que es aparente y encubre otro acto. Puede tener fines lícitos (por ejemplo, poner un auto a nombre de un hermano) o ilícitos (por ejemplo, simular una venta para perjudicar a coherederos). \\

\textbf{¿Simulación absoluta o relativa?} &
Absoluta: el acto no tiene nada de real (por ejemplo, una madre que ``vende'' un bien pero sigue siendo suya). Relativa: el acto aparente encubre a otro real y válido (art. 334 CCyC), que subsiste plenamente eficaz si no es ilícito ni perjudica a terceros. \\

\textbf{¿Simulación lícita o ilícita?} &
Lícita: no perjudica a terceros ni viola la ley. Ilícita: perjudica a terceros o viola la ley (por ejemplo, vender bienes para eludir acreedores). \\

\textbf{¿Qué sanción tiene la simulación?} &
La nulidad, de carácter relativo (no la inexistencia, categoría que el Código Civil y Comercial no reconoce para estos casos). \\

\textbf{¿Cómo se prueba entre las partes?} &
Mediante el contradocumento (art. 335 CCyC); puede prescindirse de él si se justifican las razones de su ausencia y median circunstancias que hacen inequívoca la simulación. \\

\textbf{¿Cómo prueba un tercero?} &
Por cualquier medio de prueba (art. 336 CCyC), típicamente a través de indicios: falta de medios económicos del comprador, parentesco cercano, precio vil, ausencia de ejecución material del contrato, pagos anticipados, retroventa, entre otros. \\

\textbf{¿Quién queda protegido frente a terceros?} &
Los acreedores de buena fe del adquirente simulado que hayan ejecutado bienes (art. 337 CCyC). Los subadquirentes de mala fe o a título gratuito responden; los de buena fe y título oneroso quedan protegidos. \\

\textbf{¿Cuándo prescribe la acción de simulación?} &
A los dos años (art. 2563 inc. b CCyC): entre partes, desde que una de ellas se negó a dejar sin efecto el acto; por terceros, desde que conocieron o pudieron conocer el vicio. \\

\midrule

\multicolumn{2}{p{\dimexpr\textwidth-2\tabcolsep\relax}}{
\textbf{Resumen:} La simulación y el fraude son los dos últimos vicios de los actos jurídicos estudiados en esta parte, y ambos responden a la misma preocupación: qué hacer cuando un acto jurídico, sin ser lo que aparenta, afecta a personas ajenas a él. En la simulación, las partes se ponen de acuerdo para declarar algo que no es real o que oculta otro acto; según perjudique o no a terceros o viole la ley, será ilícita o lícita, y siempre se sanciona con la nulidad relativa del acto ostensible, probándose entre las partes mediante contradocumento y, por los terceros, mediante cualquier indicio idóneo.
} \\

\bottomrule
\end{tabularx}
\end{table}

% ===========================================================
% CAPÍTULO 2 --- EL FRAUDE
% ===========================================================
\chapter[El fraude]{El fraude como vicio de los actos jurídicos}

\section{Introducción y distinción con el fraude penal}

La palabra ``fraude'' no debe confundirse con el delito de fraude, si bien puede darse que la misma conducta configure, a la vez, un fraude civil y un delito. En este capítulo se consideran los aspectos civiles del fraude, vinculados con la garantía del patrimonio a favor de los acreedores.

En la concepción integral del derecho civil y del Código Civil y Comercial, el patrimonio opera como garantía de los acreedores. Por eso, si un deudor enajena bienes para sustraerlos de sus acreedores, comete un fraude: elude, por así decirlo, la acción de sus acreedores. El código da una respuesta a esta conducta a través de determinadas disposiciones.

\section{Definición doctrinaria}

\begin{definition}{Fraude}
El fraude tiene lugar cuando una persona, en forma intencional, se desapodera de bienes o poderes para colocarse en una situación de insolvencia y eludir de esta manera la acción de sus acreedores.
\end{definition}

La insolvencia expresa una situación en la que el pasivo supera al activo, y en la que, por lo tanto, el deudor no tiene bienes suficientes para hacer frente a sus deudas.

\begin{example}{Insolvencia y dinero en efectivo}
Un deudor vende su casa para evitar que se la ejecuten. Si conserva el dinero de la venta en efectivo --en una caja fuerte, en un bolso, o de cualquier otra forma-- resulta muy difícil de embargar, salvo que en algún mandamiento se logre identificarlo. Así, la persona aparece como sin bienes, pero los acreedores pueden descubrir que hubo un acto de enajenación y atacarlo.
\end{example}

\section{La inoponibilidad como consecuencia del fraude}

El fraude no da lugar a una nulidad, sino a la \textbf{inoponibilidad}. La acción que corresponde al fraude se denomina acción de inoponibilidad en el nuevo Código; históricamente se la conocía como acción pauliana o acción revocatoria, nombre que usaba el código de Vélez. El nuevo Código, siguiendo lo que ya aceptaba la mayoría de la doctrina antes de su entrada en vigencia, habla directamente de inoponibilidad.

El acto es válido entre las partes, pero inoponible: tiene una ineficacia relativa. Esta ineficacia consiste en que el acto no produce efectos respecto de ciertas personas determinadas, que son los acreedores defraudados que inician la acción. Frente a ellos, el acto se comporta como si no se hubiera celebrado, es decir, es inoponible; pero frente a las propias partes, y frente a cualquier otra persona que no sea el acreedor accionante, el acto sigue siendo válido.

\begin{important}
El concepto de inoponibilidad --y no de nulidad-- es fundamental para diferenciar al fraude de los demás vicios de los actos jurídicos.
\end{important}

\section{La acción de inoponibilidad}

Metodológicamente, el fraude está tratado en la Sección Tercera del Capítulo 6 (``De los vicios de los actos jurídicos''), del Título Cuarto, del Libro Primero, Parte General.

\begin{formula}{Art. 338 CCyC}
``Todo acreedor puede solicitar la declaración de inoponibilidad de los actos celebrados por su deudor en fraude de sus derechos, y de las renuncias al ejercicio de derechos o facultades con los que hubiese podido mejorar o evitado empeorar su estado de fortuna.''
\end{formula}

El acreedor puede pedir la declaración de inoponibilidad tanto de los actos hechos en fraude, como de las renuncias hechas para eludir la acción de los acreedores.

\begin{example}{Renuncia en fraude a los acreedores}
Quien renuncia a cobrar una indemnización o un crédito perjudica a sus acreedores, porque estos se ven afectados por esa renuncia. Los acreedores pueden alegar que esa renuncia fue hecha en fraude a sus derechos.
\end{example}

\begin{note}
Esta inoponibilidad, como supuesto de ineficacia relativa, debe complementarse con los artículos 396 y 397 del Código, que regulan la inoponibilidad en general y que se estudian al abordar las ineficacias.
\end{note}

\section{Requisitos de procedencia de la acción}

\subsection{Crédito de causa anterior al acto impugnado}

El primer requisito es que el crédito sea de causa anterior al acto impugnado, salvo que el deudor haya actuado con el propósito de defraudar a futuros acreedores.

\begin{example}{Crédito anterior al acto impugnado}
Un deudor tiene una deuda desde el 1 de enero de 2018; a partir de esa fecha, su acreedor tiene la expectativa de que todo su patrimonio constituye la garantía del crédito. El acreedor no podría impugnar un acto realizado en 2017, porque en ese momento no existía la deuda. Si al 1 de enero de 2018 el deudor tenía un departamento y un auto, y durante 2018 enajena el departamento y se vuelve insolvente --porque además el auto se destruye en un choque--, el acreedor puede atacar la venta del departamento, porque su crédito es anterior (enero de 2018) y la venta es posterior (mayo). En cambio, si antes tenía además una cochera que vendió durante 2017, el acreedor no puede atacar esa venta, porque su crédito es posterior a ella, salvo que demuestre que la venta se hizo a sabiendas, para defraudar a los futuros acreedores.

Este sería el caso de un deudor que tiene varios inmuebles y, mientras un acreedor tramita un crédito que tarda varios meses en salir, vende todos sus inmuebles; cuando el crédito finalmente se otorga, el deudor ya no tiene nada, pero lo hizo a sabiendas, en fraude de los futuros acreedores.
\end{example}

\subsection{Que el acto haya causado o agravado la insolvencia}

El segundo requisito es que el acto haya causado o agravado la insolvencia.

\begin{example}{Insolvencia causada o agravada}
Un deudor que tiene cinco departamentos y vende uno, conservando otros cuatro con los cuales puede responder perfectamente, no incurre en fraude: este supone una situación de incapacidad e impotencia para atender los créditos.
\end{example}

\subsection{La complicidad del adquirente a título oneroso}

El tercer requisito exige distinguir según el acto haya sido a título oneroso o a título gratuito. Si es a título oneroso, debe probarse que quien contrató con el deudor conocía, o debía conocer, que el acto provocaba o agravaba la insolvencia, es decir, que fue cómplice porque conocía la situación. Si el acto fue a título gratuito, no se exige este requisito de complicidad, y la acción procede siempre, porque no hay nadie que se haya enriquecido en perjuicio de acreedores que quedaron desprovistos del bien.

\begin{example}{Complicidad del adquirente}
Un contador que conoce perfectamente la situación contable de su cliente le compra un departamento: en ese caso corresponde analizar si conocía o debía conocer el estado de insolvencia del vendedor. La complicidad es, en la práctica, difícil de probar.
\end{example}

\section{Efectos de la acción de inoponibilidad}

El acto es válido entre las partes, pero sus efectos son inoponibles al acreedor defraudado que obtiene la sentencia a su favor. A diferencia de la nulidad, que vuelve las cosas al estado anterior, en el fraude el acto sigue siendo válido entre las partes; lo único que ocurre es que, al ejecutarse la cosa, el acreedor cobra su crédito, y el excedente queda para quien había adquirido el bien.

\begin{example}{Efectos de la inoponibilidad}
Una persona vende un departamento a un pariente cercano en fraude de sus acreedores. Se prueba el fraude; la deuda era de 40.000 dólares y el departamento vale 100.000. Se ejecuta el bien, se obtienen 100.000 dólares: los acreedores cobran su parte, que son 40.000, y los 60.000 restantes no son de ellos, sino de quien había adquirido el inmueble, que se queda con ese excedente. El adquirente tiene, a su vez, un reclamo contra el deudor por lo que pagó, aunque probablemente no tenga de dónde cobrarse.

Esta es una diferencia importante entre el fraude y la simulación: en la simulación, si se declara la nulidad, las cosas vuelven al estado anterior; en el fraude, si al ejecutar la cosa hay un excedente, este se entrega a quien había adquirido la cosa, porque el acto vale entre las partes.
\end{example}

\section{El desinterés del adquirente}

\begin{formula}{Art. 341 CCyC}
El adquirente puede desinteresar a los acreedores que promovieron la acción, ofreciendo una garantía suficiente, o pagando directamente el crédito de quienes se hubiesen presentado.
\end{formula}

\begin{example}{Desinterés del adquirente}
Un adquirente que se encuentra muy conforme con el departamento adquirido --por su cercanía al trabajo y al colegio de sus hijos-- y a quien se le inicia la acción por haber sido cómplice, ofrece una garantía o paga directamente a los acreedores. En ese caso, la acción de inoponibilidad se detiene, porque los acreedores quedan desinteresados.
\end{example}

\section{Extensión de la inoponibilidad}

\begin{formula}{Art. 342 CCyC}
La declaración de inoponibilidad beneficia solo a los acreedores que promovieron la acción, y hasta el importe de sus respectivos créditos.
\end{formula}

Solo el acreedor que promueve la acción se beneficia con ella; si existen otros acreedores que no la iniciaron, no pueden beneficiarse. Retomando el ejemplo anterior: si el acreedor accionante tiene un crédito de 40.000 dólares y el departamento vale 100.000, se beneficia solo con los 40.000; los 60.000 restantes quedan para quien hubiera adquirido la cosa.

\section{Efectos frente a subadquirentes}

Este punto resulta muy similar al analizado para la simulación. Si el adquirente tiene, a su vez, sus propios acreedores, que embargan y ejecutan de buena fe, estos quedan protegidos: en el fraude, tampoco puede oponérsele al acreedor de buena fe que haya ejecutado los bienes comprendidos en el acto.

Si existe un subadquirente, corresponde analizar si es o no cómplice del fraude, y si adquirió a título oneroso o gratuito. Si es de mala fe, o si adquirió a título gratuito, en cualquiera de los dos casos procede la acción. Si el subadquirente es de buena fe y a título oneroso, queda protegido, y la acción se dirige, en forma de daños y perjuicios, contra el cómplice y el deudor que celebró el acto en fraude.

El criterio para determinar la complicidad es si el sujeto conocía el estado de insolvencia y que el acto lo agravaba. Se aplica también la regla de la responsabilidad solidaria por daños y perjuicios para quien actuó de mala fe, mientras que quien actuó de buena fe y a título gratuito responde solo en la medida de su enriquecimiento, es decir, de lo que deba restituir al acreedor defraudado.

\section{Prescripción de la acción de inoponibilidad}

El plazo de prescripción es de dos años, y se cuenta desde que el acreedor conoció o pudo conocer el vicio del acto jurídico.

\begin{important}
Con el Código de Vélez, la lesión prescribía a los cinco años y el fraude al año; el nuevo Código bajó el plazo de la lesión a dos años y elevó el del fraude también a dos años, unificando así el plazo de prescripción de todos los vicios: error, dolo, violencia, lesión, simulación y fraude prescriben, todos, a los dos años, aunque cambia el momento desde el cual comienza a correr ese plazo en cada caso.
\end{important}

A diferencia de la simulación y de los demás vicios, en el caso del fraude no se trata de una acción de nulidad, sino de una acción de inoponibilidad.

% -----------------------------------------------------------
% CORNELL --- CAPÍTULO 2
% -----------------------------------------------------------
\section{Repaso Cornell: el fraude}
\begin{table}[H]
\centering
\begin{tabularx}{\textwidth}{
    >{\raggedright\arraybackslash}p{0.28\textwidth}
    >{\raggedright\arraybackslash}X
}
\toprule
\textbf{Preguntas / palabras clave} &
\textbf{Notas / respuestas} \\
\midrule

\textbf{¿Qué es el fraude como vicio civil?} &
El acto por el cual el deudor, en forma intencional, se desapodera de bienes para colocarse en insolvencia y eludir a sus acreedores. No debe confundirse con el delito penal de fraude. \\

\textbf{¿Qué sanción tiene el fraude?} &
La inoponibilidad (ineficacia relativa), no la nulidad: el acto es válido entre las partes, pero no produce efectos frente al acreedor defraudado que obtuvo la sentencia a su favor (art. 338 CCyC). \\

\textbf{¿Qué requisitos exige la acción de inoponibilidad?} &
1) Crédito de causa anterior al acto (salvo fraude a futuros acreedores); 2) que el acto haya causado o agravado la insolvencia; 3) si es a título oneroso, que el adquirente conociera o debiera conocer ese estado (complicidad). \\

\textbf{¿Qué ocurre si hay un excedente al ejecutar el bien?} &
El excedente queda para el adquirente, no para el acreedor: a diferencia de la nulidad, el fraude no retrotrae el acto por completo, solo lo hace inoponible en la medida del crédito (arts. 341 y 342 CCyC). \\

\textbf{¿Cómo se desinteresa el adquirente?} &
Ofreciendo garantía suficiente o pagando directamente a los acreedores que promovieron la acción (art. 341 CCyC), lo que detiene la acción. \\

\textbf{¿Cuándo prescribe la acción de inoponibilidad?} &
A los dos años desde que el acreedor conoció o pudo conocer el vicio del acto (antes, con el Código de Vélez, era un año). \\

\midrule

\multicolumn{2}{p{\dimexpr\textwidth-2\tabcolsep\relax}}{
\textbf{Resumen:} En el fraude, a diferencia de la simulación, el acto es plenamente real: el deudor efectivamente enajena un bien, pero lo hace para colocarse en insolvencia y eludir a sus acreedores; por eso la sanción no es la nulidad, sino la inoponibilidad, que deja vigente el acto entre las partes y solo lo desconoce frente al acreedor que promovió la acción, protegiendo siempre al adquirente de buena fe y a título oneroso.
} \\

\bottomrule
\end{tabularx}
\end{table}

% ===========================================================
% CAPÍTULO 3 --- SÍNTESIS COMPARATIVA
% ===========================================================

\section[Síntesis comparativa]{Comparación entre simulación, fraude y acción subrogatoria}

Para cerrar el estudio de los vicios de los actos jurídicos, corresponde comparar la acción de simulación, la acción por fraude y la acción subrogatoria (esta última no es, en rigor, un vicio, pero guarda cierto parentesco con las anteriores).

\begin{table}[H]
\centering
\caption{Simulación, fraude y acción subrogatoria}
\begin{tabularx}{\textwidth}{
    >{\raggedright\arraybackslash}p{0.20\textwidth}
    >{\raggedright\arraybackslash}X
    >{\raggedright\arraybackslash}X
    >{\raggedright\arraybackslash}X
}
\toprule
\textbf{Criterio} & \textbf{Simulación} & \textbf{Fraude} & \textbf{Acción subrogatoria} \\
\midrule
¿A quién beneficia? & A todos los acreedores (vuelve el bien al patrimonio del deudor) & Solo al acreedor que la ejerce & A todos los acreedores \\
\addlinespace
¿En nombre de quién se ejerce? & En derecho propio del acreedor & En derecho propio del acreedor & Por el derecho del deudor (por ejemplo, cobrar una herencia) \\
\addlinespace
¿Requiere insolvencia del deudor? & No la requiere & Sí la requiere & No la requiere; requiere inacción del deudor \\
\addlinespace
Acción que persigue & La nulidad del acto & La inoponibilidad del acto & El fin propio del derecho que se ejerce (por ejemplo, la herencia) \\
\addlinespace
¿Quién puede ejercerla? & Cualquier persona con interés (por ejemplo, herederos entre sí) & El acreedor con interés & El acreedor con interés \\
\bottomrule
\end{tabularx}
\end{table}

Con esta comparación se cierra el estudio de los vicios de los actos jurídicos. Lo que sigue en las próximas unidades son los supuestos de ineficacia y, sobre todo, el estudio de la nulidad como el gran supuesto de ineficacia civil.

% -----------------------------------------------------------
% CORNELL --- CAPÍTULO 3
% -----------------------------------------------------------
\subsection*{Repaso Cornell: síntesis comparativa}
\begin{table}[H]
\centering
\begin{tabularx}{\textwidth}{
    >{\raggedright\arraybackslash}p{0.28\textwidth}
    >{\raggedright\arraybackslash}X
}
\toprule
\textbf{Preguntas / palabras clave} &
\textbf{Notas / respuestas} \\
\midrule

\textbf{¿En qué se diferencian simulación, fraude y acción subrogatoria?} &
La simulación persigue la nulidad y beneficia a todos los acreedores; el fraude persigue la inoponibilidad, beneficia solo al acreedor accionante y requiere insolvencia; la acción subrogatoria no es un vicio, se ejerce por el derecho del deudor, no requiere insolvencia sino inacción, y beneficia a todos los acreedores. \\

\midrule

\multicolumn{2}{p{\dimexpr\textwidth-2\tabcolsep\relax}}{
\textbf{Resumen:} Ambas figuras, junto con la acción subrogatoria, integran el sistema de protección del patrimonio como garantía común de los acreedores, y sientan las bases conceptuales --nulidad, inoponibilidad, buena y mala fe, título oneroso y gratuito-- aplicables a la ineficacia de los actos jurídicos.
% TODO: verificar consistencia entre fuentes originales: el material anuncia el estudio posterior de la nulidad como supuesto general de ineficacia, pero ese desarrollo no figura en los archivos recibidos.

} \\

\bottomrule
\end{tabularx}
\end{table}
\partintro{Así como existen hechos y actos que dan nacimiento a una relación jurídica, existen también hechos y actos que la extinguen. Esta parte estudia el cierre del ciclo de la relación jurídica por tres vías: los hechos ajenos a la voluntad de las partes, los actos jurídicos extintivos y el transcurso del tiempo (prescripción y caducidad).}
\part{La extinción de las relaciones jurídicas}
\chapter[Introducción a la extinción de las relaciones jurídicas]{Introducción: el cierre del ciclo de la relación jurídica}
%==============================================================================

\section{Objeto del tema}

El estudio de la parte general del derecho civil comprende el sujeto, el objeto y la causa de las relaciones jurídicas. Dentro de la causa se ubican los \textbf{hechos jurídicos}, que dan inicio a las relaciones y producen consecuencias jurídicas. Del mismo modo en que existen hechos y actos que dan nacimiento a una relación jurídica, existen también hechos y actos que, en algún momento, la \textbf{extinguen}. Ese es el objeto de esta parte.

\begin{note}
El desarrollo del tema es introductorio: la extinción de las obligaciones se trata con mayor profundidad en \textbf{Obligaciones}, y la extinción de los contratos, en \textbf{Contratos}.
\end{note}

En materia de hechos jurídicos extintivos se sigue la doctrina de Antonio Aguiar Anzorena.

\section{Relevancia práctica}

La prescripción y la caducidad suelen tratarse al final del estudio de la materia y no siempre se estudian con detenimiento. Sin embargo, tienen una relevancia decisiva en el ejercicio profesional: todo abogado recibirá clientes con una deuda para cobrar, y lo primero que deberá verificar es si la acción sigue vigente o ya prescribió, porque de ello depende toda la estrategia del caso.

Conocer cuándo un contrato puede resolverse, rescindirse o revocarse permite asesorar con precisión en negociaciones, redactar cláusulas de salida seguras y anticipar riesgos. Esta lógica es útil también en la vida cotidiana: saber que las deudas y los reclamos tienen plazos permite tomar decisiones a tiempo (enviar un reclamo antes de que se pierda el derecho, comprender por qué una obligación antigua ya no puede exigirse judicialmente, o entender los efectos de un arrepentimiento en una seña).

\section{Las tres vías de extinción}

Un contrato de alquiler ilustra los tres bloques temáticos.

\begin{example}{El contrato de alquiler}
El contrato nace como relación jurídica. Mientras existe, puede extinguirse:
\begin{itemize}
    \item por un \textbf{hecho ajeno a la voluntad} de las partes, por ejemplo si fallece el inquilino y era el único con derecho a habitar el inmueble (\textbf{hechos extintivos});
    \item por un \textbf{acto de voluntad} de las partes, como cuando el inquilino paga todos los meses hasta el final o decide rescindir el contrato antes de tiempo (\textbf{actos extintivos});
    \item por el \textbf{transcurso del tiempo}: si el propietario nunca reclama una deuda de alquileres atrasados, la acción para cobrarla puede prescribir, y si no ejerce ciertos derechos puntuales dentro del plazo que fija la ley, esos derechos caducan (\textbf{prescripción y caducidad}).
\end{itemize}
\end{example}

Los tres bloques constituyen las tres formas en que el derecho explica el final de una misma historia: la de una relación jurídica que nace, vive y en algún momento se apaga.

\section{Enumeración de la doctrina seguida}

\textbf{Hechos jurídicos extintivos:}
\begin{itemize}
    \item la muerte;
    \item la imposibilidad (sobrevenida);
    \item la confusión;
    \item la compensación;
    \item la caducidad.
\end{itemize}

\textbf{Actos jurídicos extintivos:}
\begin{itemize}
    \item el pago;
    \item la dación en pago;
    \item la novación;
    \item la transacción;
    \item la renuncia a los derechos;
    \item la remisión de deuda;
    \item y, como los tres principales: la resolución, la rescisión y la revocación.
\end{itemize}
\chapter[Hechos jurídicos extintivos]{Hechos jurídicos extintivos}
%==============================================================================

%------------------------------------------------------------------------------
\section{La muerte como hecho extintivo}

Los efectos jurídicos de la muerte se desarrollan en el capítulo~\ref{ch:muerte}; aquí interesa su carácter de hecho que extingue relaciones jurídicas.

%------------------------------------------------------------------------------

La muerte produce la extinción de la persona: finaliza la existencia de la persona humana. Como consecuencia, muchas relaciones jurídicas se extinguen, especialmente los \textbf{derechos extrapatrimoniales}.

También se extinguen con la muerte algunos derechos patrimoniales, como el \textbf{usufructo}: es un derecho que se constituye de por vida y, por lo tanto, se extingue con la muerte de su titular, sin transmitirse a los herederos.

\begin{important}
\textbf{El usufructo como excepción a la regla.} En materia de derechos patrimoniales, la regla general es que, a la muerte de una persona, su patrimonio se transmite a los herederos, quienes pueden ser instituidos por testamento o por disposición de la ley, existiendo algunos herederos forzosos con derecho a una porción legítima. Sin embargo, esta regla de la transmisión no rige en todos los casos: en algunos supuestos la muerte extingue directamente la relación jurídica, sin transmitirla a nadie. Uno de ellos es el usufructo.
\end{important}

%------------------------------------------------------------------------------
\section{La imposibilidad sobrevenida}
%------------------------------------------------------------------------------

El segundo hecho extintivo es la \textbf{imposibilidad}. Debe tratarse de una imposibilidad \textbf{sobreviniente} y no inicial: si la imposibilidad existiera desde el comienzo, el acto jurídico no podría formarse válidamente, porque su objeto sería imposible.

\begin{articulo}{Artículo 279 del Código Civil y Comercial}
Los actos jurídicos no pueden tener por objeto un hecho imposible.
\end{articulo}

El fundamento de esta regla se expresa en el adagio latino según el cual \textit{a lo imposible nadie está obligado}: lo imposible no puede exigirse.

La imposibilidad sobrevenida, en cambio, se produce una vez que la relación jurídica ya comenzó. Se ubica en el título de las obligaciones del Código Civil y Comercial, en el Libro Tercero.

\begin{articulo}{Artículo 955 del Código Civil y Comercial}
La imposibilidad sobrevenida, objetiva, absoluta y definitiva de la prestación, producida por caso fortuito o fuerza mayor, extingue la obligación, sin responsabilidad.
\end{articulo}

Se trata, entonces, de una extinción \textbf{sin responsabilidad}, siempre que la imposibilidad sea absolutamente imprevisible (caso fortuito o fuerza mayor). En cambio, si la obligación se vuelve imposible de cumplir por causas imputables al deudor (por ejemplo, si el deudor no pagó por su propia culpa), el deudor debe responder igualmente.

%------------------------------------------------------------------------------
\section{La confusión}
%------------------------------------------------------------------------------

\begin{definition}{Confusión}
Hecho extintivo que consiste en que la calidad de \textbf{acreedor} y la de \textbf{deudor} se reúnen en una misma persona y en un mismo patrimonio.
\end{definition}

\begin{articulo}{Artículos 931 y 932 del Código Civil y Comercial}
Regulan la extinción de la obligación por confusión, cuando las calidades de acreedor y de deudor se reúnen en una misma persona y en un mismo patrimonio.
\end{articulo}

\begin{example}{Juana y su tía}
Juana le presta 20.000 dólares a su tía. La tía fallece y Juana es su única heredera. Al heredar, Juana ocupa el lugar de su tía y se convierte, al mismo tiempo, en acreedora y deudora de la misma deuda, porque la tía le debía a ella. La obligación queda confundida y se extingue por confusión.
\end{example}

La confusión extingue la obligación en la proporción en que efectivamente se produce.

%------------------------------------------------------------------------------
\section{La compensación}
%------------------------------------------------------------------------------

\begin{definition}{Compensación}
Tiene lugar cuando dos personas reúnen, recíprocamente, la calidad de acreedor y de deudor la una de la otra. Sus créditos se compensan hasta el monto de la menor deuda.
\end{definition}

\begin{example}{Compensación de créditos recíprocos}
Si yo te debo 10.000 y vos me debés 20.000, a partir de ese momento se compensan: vos me debés 10.000, porque 20.000 menos 10.000 es 10.000. Se extinguen ambas deudas, con fuerza de pago, hasta el monto de la menor, y desde el momento en que ambas obligaciones comenzaron a coexistir en condiciones de ser compensables.
\end{example}

Se trata de un \textit{hecho} extintivo y no de un acto porque esta situación puede producirse automáticamente, sin que medie una declaración de voluntad expresa de las partes: en ese caso, la compensación opera \textbf{de pleno derecho}.

%------------------------------------------------------------------------------
\section{La caducidad (introducción)}
%------------------------------------------------------------------------------

\begin{definition}{Caducidad}
Pérdida de un derecho por no ejercitarlo en el plazo fijado por la ley. Siempre tiene \textbf{fuente legal}.
\end{definition}

Este instituto se desarrolla en profundidad junto con la prescripción (véase el capítulo~\ref{cap:prescripcion}). Desde ahora cabe señalar que la caducidad extingue el derecho no ejercido y que, a diferencia de la prescripción, \textbf{no se suspende ni se interrumpe}. Existen algunos supuestos, establecidos por la ley, en los que no corre la caducidad.

\begin{important}
No debe confundirse la \textbf{caducidad}, que extingue el derecho no ejercido, con la \textbf{prescripción liberatoria o extintiva}, que extingue la acción.
\end{important}

La caducidad rige para algunos derechos muy puntuales respecto de los cuales el Código o las leyes especiales establecen que, transcurrido cierto tiempo (generalmente plazos cortos) sin que se ejerza el derecho, este caduca.

\begin{example}{La compensación económica en el divorcio}
El derecho a exigir la compensación económica luego de un divorcio está regulado en los artículos 441 y 442 del Código Civil y Comercial. Conforme al artículo 442, la acción para reclamarla caduca a los seis meses de haberse dictado la sentencia de divorcio.
\end{example}

\begin{articulo}{Artículo 442 del Código Civil y Comercial}
La acción para reclamar la compensación económica caduca a los seis meses de haberse dictado la sentencia de divorcio.
\end{articulo}

%------------------------------------------------------------------------------
\section{Cuadro Cornell: hechos jurídicos extintivos}
%------------------------------------------------------------------------------

{\renewcommand{\arraystretch}{1.2}
\begin{longtable}{
    >{\raggedright\arraybackslash}p{\cornellL}
    >{\raggedright\arraybackslash}p{\cornellR}
}
\toprule
\textbf{Preguntas / palabras clave} &
\textbf{Notas / respuestas} \\
\midrule
\endfirsthead

\toprule
\textbf{Preguntas / palabras clave} &
\textbf{Notas / respuestas} \\
\midrule
\endhead

\textbf{¿Qué hechos extinguen relaciones jurídicas sin mediar un acto de voluntad?} &
La muerte, la imposibilidad sobrevenida, la confusión, la compensación y la caducidad. \\
\addlinespace

\textbf{¿Por qué el usufructo se extingue con la muerte de su titular?} &
Porque es un derecho constituido de por vida y no se transmite a los herederos, aunque la regla general sea que el patrimonio sí se transmite. \\
\addlinespace

\textbf{¿Qué distingue a la imposibilidad sobrevenida de la imposibilidad inicial?} &
La inicial impide que el acto se forme válidamente, porque el objeto sería imposible (art.~279 CCyC). La sobrevenida ocurre luego de nacida la obligación; si es objetiva, absoluta y definitiva, producida por caso fortuito o fuerza mayor, extingue la obligación sin responsabilidad del deudor (art.~955 CCyC). Si se debe a causas imputables al deudor, este responde igualmente. \\
\addlinespace

\textbf{¿Qué es la confusión?} &
La reunión de las calidades de acreedor y deudor en una misma persona y en un mismo patrimonio (arts.~931 y 932 CCyC), como en el caso de Juana, que hereda a su tía acreedora. \\
\addlinespace

\textbf{¿Qué es la compensación y por qué es un hecho extintivo?} &
La extinción recíproca de créditos entre dos personas que son acreedoras y deudoras entre sí, hasta el monto de la menor deuda. Es un hecho porque puede operar de pleno derecho, sin declaración de voluntad de las partes. \\
\addlinespace

\textbf{¿Qué es la caducidad y en qué se diferencia de la prescripción?} &
La pérdida de un derecho por no ejercerlo dentro del plazo fijado por la ley; siempre tiene fuente legal. Extingue el derecho no ejercido, mientras que la prescripción extintiva extingue la acción. No se suspende ni se interrumpe. Ejemplo: art.~442 CCyC (seis meses desde la sentencia de divorcio). \\

\midrule
\multicolumn{2}{p{\dimexpr\cornellL+\cornellR+2\tabcolsep\relax}}{%
\textbf{Resumen:} los hechos extintivos operan sin necesidad de un acto de voluntad: la muerte (que extingue derechos extrapatrimoniales y algunos patrimoniales, como el usufructo), la imposibilidad sobrevenida (sin responsabilidad si obedece a caso fortuito o fuerza mayor), la confusión (acreedor y deudor en una misma persona), la compensación (créditos recíprocos hasta el monto de la menor deuda) y la caducidad (pérdida del derecho no ejercido en el plazo legal).} \\
\bottomrule
\end{longtable}}

%==============================================================================
\chapter[Actos jurídicos extintivos]{Actos jurídicos extintivos}
%==============================================================================

Los actos jurídicos extintivos son aquellos que suponen una \textbf{declaración de voluntad} dirigida a poner fin a la relación jurídica.

%------------------------------------------------------------------------------
\section{El pago}
%------------------------------------------------------------------------------

El acto extintivo de una obligación por excelencia es el pago.

\begin{definition}{Pago}
Cumplimiento de la prestación que constituye el objeto de la obligación.
\end{definition}

\begin{important}
No debe asociarse el pago sólo con la entrega de dinero en efectivo: \textbf{pagar es cumplir con la obligación}.
\end{important}

\begin{example}{El pago no es sólo entregar dinero}
Si yo te presto una cosa mueble y vos te comprometés a darme plata a cambio, cuando yo te entrego la cosa pactada estoy pagando, porque cumplo mi obligación; y cuando vos me entregás el dinero, vos también estás pagando. Del mismo modo, si a una persona la contratan para hacer una asesoría de servicios jurídicos, esa persona paga cuando cumple con la prestación de hacer esa asesoría.
\end{example}

%------------------------------------------------------------------------------
\section{La dación en pago}
%------------------------------------------------------------------------------

\begin{definition}{Dación en pago}
Acto extintivo en el cual el acreedor acepta voluntariamente, en pago, una prestación diversa de la debida.
\end{definition}

\begin{example}{Cambiar un servicio por dinero}
Me comprometo a prestar servicios jurídicos pero, por un problema, no puedo hacerlo. En lugar de esa prestación, ofrezco dinero al acreedor y este lo acepta en pago de la deuda original. Eso es una dación en pago.
\end{example}

%------------------------------------------------------------------------------
\section{La novación}
%------------------------------------------------------------------------------

\begin{definition}{Novación}
Extinción de una obligación por la creación de otra nueva, destinada a reemplazarla. Aparece una cierta modificación y la nueva obligación sustituye a la anterior.
\end{definition}

\begin{example}{Reemplazo del objeto de la obligación}
Yo debía entregar una cosa determinada y, mediante novación, pactamos que ahora debo entregar otra cosa distinta. Se extingue la primera obligación y aparece la segunda en su reemplazo.
\end{example}

%------------------------------------------------------------------------------
\section{La transacción}
%------------------------------------------------------------------------------

La transacción es un \textbf{contrato}. Con ella ya se sale de los artículos ubicados en la parte de las obligaciones (en torno al número 900) para pasar a la parte de contratos especiales, dentro del mismo Libro Tercero, en torno al artículo 1641. Su ubicación como contrato es una novedad del Código actual.

La palabra «transar» tiene en derecho civil y comercial un alcance muy específico: significa poner fin a obligaciones dudosas o litigiosas.

\begin{definition}{Transacción}
Contrato por el cual las partes, haciéndose concesiones recíprocas, extinguen obligaciones dudosas o litigiosas.
\end{definition}

\begin{example}{El accidente de tránsito}
Ante un litigio (por ejemplo, un accidente de tránsito) una de las partes reclama una indemnización de 200.000 pesos y la otra ofrece una suma sensiblemente menor. Quien le pone fin a ese litigio, habitualmente, es un juez, que determina qué es lo justo. Pero si las partes desean ponerle fin anticipadamente y llegan a un acuerdo (una de ellas dice «yo reconozco mi responsabilidad, te voy a pagar cien mil pesos» y la otra acepta), están haciendo una transacción. La transacción implica ceder en parte el propio derecho: quien reclamaba 200.000 termina recibiendo 100.000.
\end{example}

En la práctica profesional, un abogado necesita contar con \textbf{poder especial} de su cliente para hacer una transacción en su nombre, y existen ciertos contratos que no admiten transacción. Estas cuestiones están reguladas en detalle en el Código y se profundizan en Obligaciones y en Contratos.

%------------------------------------------------------------------------------
\section{La renuncia}
%------------------------------------------------------------------------------

\begin{definition}{Renuncia}
Acto jurídico extintivo que consiste en la abdicación de un derecho.
\end{definition}

\begin{articulo}{Artículo 944 del Código Civil y Comercial}
Toda persona puede renunciar a los derechos conferidos por la ley cuando la renuncia no está prohibida y sólo afecta intereses privados. No se admite la renuncia anticipada de las defensas que pueden hacerse valer en juicio.
\end{articulo}

\begin{example}{Renunciar a cobrar una deuda}
«Te renuncio a cobrarte esa deuda que tengo». Es una liberalidad: en el fondo se está regalando algo, y sigue siendo un acto a título gratuito. Por ello puede ser eventualmente impugnado si hubiera acreedores interesados en que esa deuda sí se cobre.
\end{example}

%------------------------------------------------------------------------------
\section{La remisión de deuda}
%------------------------------------------------------------------------------

\begin{definition}{Remisión de deuda}
Acto jurídico que se produce cuando el acreedor entrega voluntariamente al deudor el título en el que consta la deuda (por ejemplo, un pagaré).
\end{definition}

%------------------------------------------------------------------------------
\section{La resolución}
%------------------------------------------------------------------------------

La resolución, la rescisión y la revocación son los tres institutos principales con los que se cierra el estudio de los actos jurídicos extintivos.

\begin{definition}{Resolución}
Extinción de los efectos de la relación jurídica que se produce cuando ocurre un acontecimiento que, conforme a previsiones legales precisas o a pactos de las partes, provoca dicha extinción.
\end{definition}

\begin{important}
En derecho civil, «resolución» tiene un alcance técnico: no significa resolver un problema en el sentido de encontrarle una solución, sino \textbf{ponerle fin}. En general, la resolución pone fin de manera \textbf{retroactiva}: opera hacia atrás, hacia el momento de origen de las actuaciones. Este efecto se denomina, en latín, \textit{ex tunc}.
\end{important}

\subsection{La cláusula resolutoria (o pacto comisorio)}

Es una cláusula de un contrato por la cual las partes tienen la potestad de resolverlo, es decir, de ponerle fin, en caso de que la otra parte no cumpla. Puede ser \textbf{expresa} (pactada explícitamente) o \textbf{implícita} (presente en todos los contratos bilaterales, aunque nada se haya pactado). También se la conoce como \textbf{pacto comisorio}, expreso o implícito.

\begin{articulo}{Artículo 1087 del Código Civil y Comercial}
Todo contrato bilateral tiene implícita la facultad de resolver las obligaciones emergentes de él en caso de incumplimiento de una de las partes.
\end{articulo}

Los artículos 1086, 1088 y 1089 del Código precisan los requisitos y el modo en que debe ejercerse esta facultad.
% TODO: contenido incompleto en las fuentes originales: no se profundiza en los requisitos ni en el modo de ejercicio de esta facultad.


\begin{example}{El pintor que no cumplió el día pactado}
Se pacta que un pintor venga a pintar una casa el primero de enero. El pintor se toma diez días de vacaciones y recién el diez de enero pretende que se lo acepte igual. El dueño de la casa le dice «usted incumplió porque no vino el día que habíamos pactado» y resuelve el contrato, debiendo notificarlo conforme a los requisitos correspondientes.
\end{example}

\begin{example}{El catering que llega tarde}
Se contrata un servicio de catering para un casamiento el día viernes y el proveedor pretende entregar las cosas recién el domingo. El cliente responde: «usted incumplió, yo no le voy a pagar». El proveedor incumplió su parte y, por lo tanto, el cliente tampoco cumple la suya: se resuelve el contrato.
\end{example}

\subsection{La señal o arras}

La señal (también llamada arras) es una suma de dinero que se entrega al inicio de un contrato y que puede tener dos alcances distintos:

\begin{itemize}
    \item \textbf{Facultad de arrepentirse:} debe pactarse expresamente. Quien entregó la señal puede luego arrepentirse y poner fin al contrato (cláusula de arrepentimiento).
    \item \textbf{Función confirmatoria del acto:} es el alcance que se le atribuye si no se pactó nada distinto. Quien entregó la señal debe cumplir todo el contrato igualmente.
\end{itemize}

\begin{important}
Es muy importante ver cómo se pactó la entrega de la seña, porque de ello depende si el contrato queda firme o si existe la posibilidad de arrepentirse.
\end{important}

\begin{articulo}{Artículo 1059 del Código Civil y Comercial}
Regula la señal o arras: su función confirmatoria y, si así fue pactado expresamente, la facultad de arrepentimiento.
\end{articulo}

\subsection{La condición resolutoria}

Es una cláusula que subordina la resolución (es decir, la extinción) de una relación jurídica a un \textbf{hecho futuro e incierto}.

\subsection{La frustración de la finalidad}

El Código Civil y Comercial incorpora, como causal de resolución, la frustración definitiva de la finalidad del contrato: cuando el contrato se frustra en su finalidad, la parte perjudicada puede declarar su resolución, siempre que medie una causa extraordinaria en las circunstancias.

\begin{articulo}{Artículo 1090 del Código Civil y Comercial}
Regula la frustración definitiva de la finalidad del contrato como causal de resolución, cuando obedece a una alteración extraordinaria de las circunstancias existentes al momento de su celebración.
\end{articulo}

\begin{example}{La pandemia y el casamiento frustrado}
Alguien debía entregar algo el diez de abril porque iba a celebrarse un casamiento, y este finalmente no se realizó porque se prohibieron todas las reuniones de personas. Se frustró la finalidad del contrato y la parte perjudicada puede poner fin al acto invocando esa frustración.
\end{example}

\subsection{La teoría de la imprevisión}

La teoría de la imprevisión tiene como antecedente el artículo 1198 del Código de Vélez, que no contenía originalmente este instituto: fue incorporado por la reforma de la ley 17.711. Se aplica cuando la prestación a cargo de una de las partes se torna \textbf{excesivamente onerosa} por una alteración extraordinaria de las circunstancias existentes al momento de la celebración del contrato, producida por una causa sobreviniente.

Esta figura se compara con la lesión, que también supone una desproporción entre las prestaciones, aunque se trata de institutos distintos: en la imprevisión, la alteración de las bases del contrato produce una desproporción que torna excesivamente onerosa la posición de una de las partes.

\begin{articulo}{Artículo 1091 del Código Civil y Comercial}
Autoriza, ante la excesiva onerosidad sobreviniente de la prestación, la resolución total o parcial del contrato, o su adecuación.
\end{articulo}

\begin{note}
La figura tuvo su apogeo periódicamente en Argentina con los cambios en el dólar y la hiperinflación. Sin embargo, se advierte que la inflación en Argentina «ya no es algo previsible», por lo que el instituto va teniendo, con el tiempo, situaciones de aplicación cada vez más acotadas, porque siempre debe tratarse de algo realmente imprevisible.
\end{note}

\subsection{Efectos generales de la resolución}

La resolución opera, generalmente, en forma \textbf{retroactiva} (\textit{ex tunc}).

\begin{articulo}{Artículo 1079 del Código Civil y Comercial}
Salvo disposición legal en contrario, la resolución produce efectos retroactivos, con excepción de las prestaciones cumplidas que quedan firmes en cuanto resulten equivalentes.
\end{articulo}

Si ya se han cumplido determinadas prestaciones y quedaron firmes, la resolución no las alcanza. En esto se diferencia de las dos figuras siguientes (rescisión y revocación), que operan hacia el futuro.

%------------------------------------------------------------------------------
\section{La rescisión}
%------------------------------------------------------------------------------

\begin{definition}{Rescisión}
Forma de extinción que opera a futuro: no tiene efectos retroactivos. Puede ser bilateral o unilateral, en los casos autorizados por la ley.
\end{definition}

% TODO: contenido incompleto o ambiguo en las notas originales (el apunte cita nuevamente el artículo 1079 para los efectos de la rescisión, tras haberlo citado para los efectos de la resolución)
\begin{articulo}{Artículo 1079 del Código Civil y Comercial (efectos de la extinción por declaración de una de las partes)}
La rescisión opera a futuro. Excepto estipulación en contrario, sólo produce efectos para el futuro y no afecta el derecho de terceros.
\end{articulo}

\begin{example}{Rescisión bilateral: un contrato de servicios jurídicos}
Se trata, generalmente, de contratos de tracto sucesivo: por ejemplo, un contrato de servicios jurídicos que se viene pagando todos los meses. Un día, ambas partes deciden rescindirlo de común acuerdo. Eso es una rescisión bilateral.
\end{example}

\begin{example}{Rescisión unilateral: el contrato de alquiler}
El contrato puede extinguirse total o parcialmente por la sola declaración de una de las partes, en los casos en que el propio contrato o la ley le otorgan esa facultad. El ejemplo típico es el contrato de alquiler, que puede ser rescindido por el inquilino con cierto preaviso, a partir de los seis meses de celebrado el contrato: la ley concede unilateralmente al locatario la potestad de rescindir.
\end{example}

%------------------------------------------------------------------------------
\section{La revocación}
%------------------------------------------------------------------------------

\begin{definition}{Revocación}
Acto jurídico unilateral mediante el cual una de las partes deja sin efecto la relación por su exclusiva voluntad. Al igual que la rescisión, es unilateral y produce efectos a futuro: no es retroactiva.
\end{definition}

\begin{example}{El mandato revocado}
Se contrata a un abogado para que actúe en nombre de una persona (mandato). En un momento se produce un quiebre de la confianza, o simplemente existe la voluntad de no continuar con esa representación, y se revoca el mandato. Todo lo que el abogado hizo en nombre del mandante hasta ese momento es válido, porque tenía mandato vigente; pero desde el día en que se revoca, ya no tiene mandato para actuar en su nombre.
\end{example}

\begin{example}{El testamento revocado}
Una persona otorga un testamento válido. Mientras vive, puede revocarlo y otorgar otro nuevo, quedando revocado el anterior.
\end{example}

\subsection{Cuadro comparativo: resolución, rescisión y revocación}

\begin{table}[H]
\centering
\begin{tabularx}{\textwidth}{
    >{\raggedright\arraybackslash}p{3.0cm}
    >{\raggedright\arraybackslash}X
    >{\raggedright\arraybackslash}X
}
\toprule
\textbf{Criterio} & \textbf{Resolución} & \textbf{Rescisión / Revocación} \\
\midrule
Efectos en el tiempo & Retroactivos (\textit{ex tunc}) & Hacia el futuro (\textit{ex nunc}) \\
\addlinespace
Origen típico & Incumplimiento, condición, frustración de la finalidad, imprevisión & Acuerdo de partes o facultad legal / contractual unilateral \\
\addlinespace
Partes involucradas & Puede requerir declaración de la parte cumplidora & Rescisión: bilateral o unilateral. Revocación: unilateral \\
\addlinespace
Prestaciones ya cumplidas & En principio quedan sin efecto, salvo equivalencia & Quedan firmes; no se retrotraen \\
\bottomrule
\end{tabularx}
\end{table}

%------------------------------------------------------------------------------
\section{Cuadro Cornell: actos jurídicos extintivos}
%------------------------------------------------------------------------------

{\renewcommand{\arraystretch}{1.2}
\begin{longtable}{
    >{\raggedright\arraybackslash}p{\cornellL}
    >{\raggedright\arraybackslash}p{\cornellR}
}
\toprule
\textbf{Preguntas / palabras clave} &
\textbf{Notas / respuestas} \\
\midrule
\endfirsthead

\toprule
\textbf{Preguntas / palabras clave} &
\textbf{Notas / respuestas} \\
\midrule
\endhead

\textbf{¿Cuál es el acto extintivo por excelencia?} &
El pago: el cumplimiento de la prestación que constituye el objeto de la obligación. No se limita a la entrega de dinero: pagar es cumplir con la obligación. \\
\addlinespace

\textbf{¿Qué diferencia a la dación en pago de la novación?} &
En la dación en pago el acreedor acepta voluntariamente una prestación distinta de la debida. En la novación nace una obligación nueva que reemplaza y extingue a la anterior. \\
\addlinespace

\textbf{¿Qué es la transacción?} &
Un contrato (ubicado en torno al art.~1641 CCyC) por el cual las partes, mediante concesiones recíprocas, extinguen obligaciones dudosas o litigiosas, evitando o poniendo fin a un litigio. El abogado necesita poder especial de su cliente para transar. \\
\addlinespace

\textbf{¿Qué diferencia a la renuncia de la remisión de deuda?} &
La renuncia es la abdicación de un derecho conferido por la ley, cuando no está prohibida y sólo afecta intereses privados (art.~944 CCyC); puede ser una liberalidad a título gratuito. La remisión de deuda es la entrega voluntaria, por el acreedor, del título en el que consta la deuda. \\
\addlinespace

\textbf{¿Qué efecto tiene la resolución sobre la relación jurídica?} &
Un efecto retroactivo (\textit{ex tunc}), salvo las prestaciones ya cumplidas que resulten equivalentes y queden firmes (art.~1079 CCyC). \\
\addlinespace

\textbf{¿Qué formas o causas puede adoptar la resolución?} &
La cláusula resolutoria o pacto comisorio (expreso o implícito, art.~1087 CCyC), la señal o arras (art.~1059 CCyC), la condición resolutoria, la frustración de la finalidad (art.~1090 CCyC) y la teoría de la imprevisión (art.~1091 CCyC). \\
\addlinespace

\textbf{¿Cómo influye la forma en que se pactó la señal?} &
Si se pactó expresamente como facultad de arrepentirse, quien la entregó puede poner fin al contrato. Si no se pactó nada distinto, es confirmatoria del acto y quien la entregó debe cumplir todo el contrato. \\
\addlinespace

\textbf{¿Cuándo se aplica la teoría de la imprevisión?} &
Cuando la prestación de una de las partes se torna excesivamente onerosa por una alteración extraordinaria de las circunstancias existentes al celebrarse el contrato, producida por una causa sobreviniente e imprevisible. \\
\addlinespace

\textbf{¿Qué diferencia a la rescisión y a la revocación de la resolución?} &
Ambas operan hacia el futuro (\textit{ex nunc}), sin efecto retroactivo. La rescisión puede ser bilateral o unilateral; la revocación es unilateral (ejemplos: mandato, testamento). \\

\midrule
\multicolumn{2}{p{\dimexpr\cornellL+\cornellR+2\tabcolsep\relax}}{%
\textbf{Resumen:} los actos extintivos requieren una declaración de voluntad. El pago cumple la obligación; la dación en pago la satisface con una prestación distinta; la novación la reemplaza por otra nueva; la transacción pone fin a obligaciones dudosas o litigiosas mediante concesiones recíprocas; la renuncia abdica un derecho y la remisión entrega el título de la deuda. La resolución opera hacia atrás (\textit{ex tunc}); la rescisión y la revocación, hacia el futuro.} \\
\bottomrule
\end{longtable}}

%==============================================================================
\chapter[Prescripción y caducidad]{Prescripción y caducidad}\label{cap:prescripcion}
%==============================================================================

%------------------------------------------------------------------------------
\section{Importancia de la prescripción en la práctica}
%------------------------------------------------------------------------------

La prescripción es un instituto de enorme importancia que, por tratarse generalmente al final de la materia, no siempre se estudia con detenimiento; sin embargo, tiene una relevancia decisiva en la vida profesional del abogado.

\begin{example}{El cliente que llega con una deuda vieja}
Un cliente llega con una deuda y pregunta: «¿yo tengo derecho a cobrar esto? Hace tiempo que pasó, alguna vez le mandé un mensajito para reclamarla». El abogado pregunta cuánto tiempo pasó y a veces la respuesta es «como cinco años». Ahí hay que advertir que pueden quedar apenas tres días antes de que la deuda prescriba, y actuar con urgencia. Por eso, lo primero que hay que revisar al recibir a un cliente son las \textbf{fechas}.
\end{example}

El fundamento de la prescripción es que existe la expectativa de que los derechos se ejerzan en un tiempo oportuno: no se puede vivir con la indeterminación de que en cualquier momento alguien venga a reclamar algo pasado un tiempo indefinido. La inacción hace presumir, por ley, que decayó el interés en ejercer el derecho. Si una persona tiene una deuda a su favor y pasan veinte años sin reclamarla, ya no parece tener tanto interés en cobrarla, y esto hace que la deuda (más precisamente, la acción para reclamarla) prescriba.

%------------------------------------------------------------------------------
\section{Prescripción adquisitiva y prescripción extintiva}
%------------------------------------------------------------------------------

\begin{definition}{Prescripción}
Instituto jurídico en virtud del cual el transcurso del tiempo opera la adquisición de un derecho, o la extinción de la acción por la inacción del titular del derecho.
\end{definition}

La prescripción está tratada en el Libro Sexto, Título Primero del Código Civil y Comercial, que reúne las disposiciones comunes a los derechos reales y personales. Tiene dos grandes formas, la \textbf{adquisitiva} y la \textbf{extintiva o liberatoria}, que comparten algunos elementos comunes (sobre todo, la idea del transcurso del tiempo unida a una cierta inacción) pero presentan muchas diferencias entre sí y persiguen finalidades distintas. El énfasis de esta parte está en la prescripción extintiva o liberatoria.

\subsection{La prescripción adquisitiva (usucapión)}

\begin{definition}{Prescripción adquisitiva (usucapión)}
Modo por el cual, mediante la posesión de una cosa reunidas ciertas condiciones, se puede adquirir el dominio, ante el manifiesto desinterés del propietario en la atención, cuidado y vigilancia del bien.
\end{definition}

Está regulada en el Libro Cuarto, sobre Derechos Reales, y se profundiza en esa materia.

\begin{definition}{Posesión}
Relación real: la que se tiene directamente con la cosa en los hechos, es decir, el \textit{corpus} (la detentación material de algo), unida al \textit{animus domini} (el ánimo de comportarse como dueño de la cosa, sin reconocer en ningún otro un derecho superior).
\end{definition}

\begin{articulo}{Artículos 1898 y 1899 del Código Civil y Comercial}
Para la prescripción adquisitiva de inmuebles, el plazo común es de veinte años de posesión continua e ininterrumpida, sea de buena o de mala fe. Existe un plazo más corto de diez años cuando el poseedor cuenta con justo título y buena fe.
\end{articulo}

Para las cosas muebles rigen plazos distintos: si la cosa fue hurtada o perdida, el plazo de prescripción adquisitiva es de dos años; si se trata de una cosa mueble registrable, el plazo es de diez años.

\begin{example}{El automóvil poseído por más de diez años}
Una persona recibió un automóvil del titular registral o de un funcionario autorizado, de buena fe, y lo posee durante más de diez años sin que conste a su nombre en el registro de la propiedad. Esa persona pasa a tener la propiedad del bien; para ello debe promover un juicio de prescripción adquisitiva (usucapión), que es un modo de adquirir el dominio.
\end{example}

\subsection{La prescripción extintiva o liberatoria}

\begin{definition}{Prescripción extintiva o liberatoria}
Extingue la acción por la cual una persona puede reclamar un derecho, por el transcurso del tiempo fijado por la ley y la inacción del titular. Consiste en la falta de ejercicio, en el tiempo, de las acciones tendientes a reclamar los derechos, lo que determina la imposibilidad práctica de ejercerlos en el futuro por el transcurso de los plazos legales.
\end{definition}

\begin{important}
La prescripción extintiva extingue la \textbf{acción}, y no el derecho, porque puede ocurrir que finalmente se pague una deuda ya prescripta.
\end{important}

\begin{example}{El amigo que paga una deuda de veinte años}
Una persona le prestó dinero a un amigo hace mucho tiempo y nunca se lo reclamó; pasaron veinte años. Un día se encuentran y el amigo le cuenta que está pasando una mala situación económica. El otro le dice: «hace veinte años me prestaste esta plata, tomá, te pago lo que te debía», y le paga. Tiempo después, ese amigo fallece y sus herederos preguntan si, por haber pagado una deuda prescripta, corresponde la devolución. La respuesta es que no: no se pagó algo que no se debía. La deuda existía; lo que se había extinguido era solamente la acción para reclamarla judicialmente.
\end{example}

\begin{articulo}{Artículo 2538 del Código Civil y Comercial}
El pago espontáneo de una obligación prescripta no es repetible.
\end{articulo}

«Repetible» significa que puede pedirse el reembolso de lo pagado; por lo tanto, lo pagado espontáneamente no puede reclamarse. En el Código Civil anterior, esta situación se llamaba «obligación natural». El nuevo Código evita expresamente ese nombre y utiliza la expresión «pago espontáneo», pero en los hechos sigue siendo la misma idea: una obligación que tiene fundamento en la idea natural de justicia (dar a cada uno lo suyo). La deuda existe; lo que ocurre es que prescribió y ya no es exigible judicialmente. La acción que acompaña al derecho se extingue y el derecho queda desprovisto de acción, pero sigue siendo un derecho que existía.

Consecuencia práctica: cuando la contraparte intenta hacer valer un derecho ya prescripto, el abogado debe oponer la \textbf{excepción de prescripción}, defensa procesal por la cual se señala al juez que esa acción está prescripta y por qué motivos.

\subsection{Reglas generales sobre la prescripción}

Las normas de prescripción son \textbf{imperativas}: no pueden ser modificadas por las partes. No se puede pactar, por ejemplo, que una deuda sea imprescriptible.

\begin{articulo}{Artículo 2533 del Código Civil y Comercial}
Las normas relativas a la prescripción no pueden ser modificadas por convención de las partes.
\end{articulo}

La prescripción opera a favor y en contra de todas las personas, salvo disposición legal en contrario. En las obligaciones con sujetos múltiples (varios acreedores o varios deudores) cualquier acreedor puede oponer la prescripción, aunque el obligado no la invoque, y esta puede también renunciarse. % TODO: contenido incompleto o ambiguo en las notas originales (no se precisa quién puede oponer la prescripción en las obligaciones con sujetos múltiples)

La renuncia a la prescripción sólo puede hacerse una vez que la prescripción ya se ganó: la persona a cuyo favor corrió el plazo puede renunciar a invocarla, de modo similar a la renuncia como acto extintivo. La prescripción puede ser invocada en todos los casos, salvo en los supuestos en que la ley indique expresamente lo contrario.

%------------------------------------------------------------------------------
\section{Suspensión e interrupción del cómputo de la prescripción}
%------------------------------------------------------------------------------

Es muy importante distinguir dos categorías de hechos que pueden incidir sobre el cómputo del plazo de prescripción: los \textbf{hechos suspensivos} y los \textbf{hechos interruptivos}.

\subsection{La suspensión}

\begin{definition}{Suspensión del cómputo}
Detiene el conteo del tiempo mientras dura la causa suspensiva, pero aprovecha (conserva) el período ya transcurrido antes de que la suspensión comenzara.
\end{definition}

\begin{articulo}{Artículo 2539 del Código Civil y Comercial}
La suspensión de la prescripción detiene el cómputo del tiempo por el lapso que dura, pero aprovecha el período transcurrido hasta que ella comenzó.
\end{articulo}

\begin{example}{Cómo se cuenta un plazo suspendido}
El plazo de prescripción empieza a correr el 1 de enero de 2017 y el 30 de junio de 2017 ocurre un hecho suspensivo que dura un año: el cómputo se detiene ese día. El tramo ya corrido (del 1 de enero al 30 de junio de 2017) es válido y se conserva. Cuando, un año después, el 30 de junio de 2018, se reanuda el cómputo, se sigue contando desde el sexto mes en adelante, sin perder el tiempo ya corrido.
\end{example}

\textbf{Supuestos de suspensión}

\begin{itemize}
    \item \textbf{Interpelación fehaciente.} Es el acto típico del abogado que recibe un caso y reclama la deuda al deudor (o al poseedor) de manera extrajudicial, mediante carta documento o acta notarial, con fuerza probatoria hacia terceros. Suspende el cómputo \textbf{por una sola vez y por seis meses}.
\end{itemize}

\begin{articulo}{Artículo 2541 del Código Civil y Comercial}
La interpelación fehaciente hecha por el titular del derecho contra el deudor o el poseedor suspende la prescripción por seis meses, por una sola vez.
\end{articulo}

\begin{itemize}
    \item \textbf{Pedido de mediación.} Suspende el cómputo durante todo el tiempo que dura el trámite de mediación y hasta veinte días después de que se labra el acta de cierre, momento en el cual se reanuda el cómputo del plazo.
\end{itemize}

\begin{articulo}{Artículo 2542 del Código Civil y Comercial}
El curso de la prescripción se suspende desde la expedición por medio fehaciente de la solicitud de mediación, hasta veinte días contados desde el momento en que el acta que da por concluido el procedimiento se encuentre a disposición de las partes.
\end{articulo}

\begin{itemize}
    \item \textbf{Otros supuestos especiales.} Se trata de vínculos de confianza en los que los conflictos de interés dificultan ejercer acciones entre sí. Están suspendidas las acciones:
    \begin{itemize}
        \item entre cónyuges, mientras dura el matrimonio, para no sumar un motivo más a las disputas conyugales;
        \item entre convivientes en unión convivencial;
        \item entre las personas con capacidad restringida y sus representantes, apoyos o tutores, mientras dura la responsabilidad parental, la tutela o la curatela;
        \item entre las personas jurídicas y sus administradores;
        \item a favor y en contra de los herederos, mientras dure la respectiva situación.
    \end{itemize}
\end{itemize}

\begin{articulo}{Artículo 2543 del Código Civil y Comercial}
Enumera los supuestos especiales de suspensión de la prescripción entre cónyuges, convivientes, personas con capacidad restringida y sus representantes o apoyos, personas jurídicas y sus administradores, y a favor y en contra de los herederos.
\end{articulo}

\subsection{La interrupción}

\begin{definition}{Interrupción del cómputo}
A diferencia de la suspensión, tiene por no sucedido el lapso que la precede e inicia un nuevo plazo. Es decir, invalida todo el tiempo ya transcurrido.
\end{definition}

\begin{example}{Cómo se cuenta un plazo interrumpido}
La prescripción empezó a correr el 1 de enero de 2017 y el 30 de junio de 2017 ocurre un hecho interruptivo que dura un año. Cuando finaliza ese hecho, el cómputo empieza de cero: no se aprovecha nada del tiempo ya transcurrido. Esto es lo que distingue, en la práctica, a la interrupción de la suspensión.
\end{example}

\textbf{Supuestos de interrupción}

\begin{itemize}
    \item \textbf{Reconocimiento de deuda.} Cuando la persona obligada firma un documento en el que reconoce la deuda.
\end{itemize}

\begin{articulo}{Artículo 2545 del Código Civil y Comercial}
El curso de la prescripción se interrumpe por el reconocimiento que el deudor o poseedor efectúa del derecho de aquel contra quien prescribe.
\end{articulo}

\begin{itemize}
    \item \textbf{Petición judicial.} Se interrumpe la prescripción por toda petición del titular del derecho ante la autoridad judicial que traduzca la intención de no abandonar el derecho.
\end{itemize}

\begin{articulo}{Artículo 2546 del Código Civil y Comercial}
La petición judicial interrumpe la prescripción, aunque sea defectuosa, realizada por quien no es capaz o ante un tribunal incompetente, y aunque se presente dentro del plazo de gracia previsto en el ordenamiento procesal.
\end{articulo}

La redacción del nuevo Código es muy amplia para describir este supuesto: \textbf{no se exige necesariamente una demanda} en sentido estricto, sino cualquier petición que traduzca la intención de no abandonar el derecho; puede bastar, por ejemplo, una verificación de crédito en un proceso sucesorio. Interrumpir resulta más beneficioso para el acreedor que suspender, porque invalida todo el tiempo transcurrido y hace que se cuente todo de nuevo.

Sobre el plazo de gracia: los plazos procesales vencen a las 24 horas, pero como los tribunales dejan de funcionar antes, se habilitan las dos primeras horas hábiles del día siguiente para presentar el escrito dentro de ese plazo, lo que se conoce habitualmente como «las dos primeras».

\begin{itemize}
    \item \textbf{Solicitud de arbitraje.} Equivale, en los hechos, a interponer una demanda o a un reconocimiento de deuda, y también interrumpe el cómputo de la prescripción.
\end{itemize}

\begin{articulo}{Artículo 2548 del Código Civil y Comercial}
La solicitud de arbitraje interrumpe la prescripción.
\end{articulo}

\subsection{Cuadro comparativo: suspensión e interrupción}

\begin{table}[H]
\centering
\begin{tabularx}{\textwidth}{
    >{\raggedright\arraybackslash}p{3.3cm}
    >{\raggedright\arraybackslash}X
    >{\raggedright\arraybackslash}X
}
\toprule
\textbf{Criterio} & \textbf{Suspensión} & \textbf{Interrupción} \\
\midrule
Tiempo ya transcurrido & Se conserva (se aprovecha) & Se invalida por completo \\
\addlinespace
Reinicio del cómputo & Continúa sumando desde donde se detuvo & Empieza de cero \\
\addlinespace
Supuestos & Interpelación fehaciente, mediación, causas entre cónyuges y convivientes, entre otras & Reconocimiento de deuda, petición judicial, solicitud de arbitraje \\
\bottomrule
\end{tabularx}
\end{table}

%------------------------------------------------------------------------------
\section{La dispensa y los plazos de prescripción}
%------------------------------------------------------------------------------

\subsection{La dispensa}

\begin{definition}{Dispensa}
Tercer supuesto, distinto de la suspensión y de la interrupción. A diferencia de ambas, ocurre una vez que la prescripción ya se cumplió, cuando por alguna circunstancia extraordinaria era imposible ejercer la acción en el momento en que finalizaba el plazo.
\end{definition}

En el Código anterior este supuesto se trataba como un caso de suspensión. El nuevo Código lo clarifica y lo llama \textbf{dispensa}: la doctrina ya interpretaba que se trataba de un supuesto de dispensa, y ahora se lo ubica claramente como tal.

\begin{articulo}{Artículo 2550 del Código Civil y Comercial}
El juez puede dispensar de la prescripción ya cumplida al titular de la acción, si dificultades de hecho o maniobras dolosas le obstaculizaron temporalmente el ejercicio de la acción, y si el titular hace valer sus derechos dentro de los seis meses siguientes a la cesación de los obstáculos.
\end{articulo}

\begin{example}{La persona postrada en un hospital}
Se debe recurrir ante un juez y explicarle, por ejemplo: «yo estaba en cama, totalmente postrado, al momento en que prescribía esta acción, y me era imposible ejercerla». Si el juez entiende que existieron esas circunstancias extraordinarias, concede una dispensa que permite iniciar la acción dentro de los seis meses siguientes. Apenas la persona sale del hospital, inicia la demanda y pide la dispensa de la prescripción.
\end{example}

La dispensa se aplica también a personas incapaces sin representante: por ejemplo, un niño de cinco años, sin representantes, que tenía que cobrar ciertas deudas; en ese caso se espera a que existan representantes que puedan actuar en su nombre.

\begin{important}
La diferencia central con la suspensión y la interrupción es que la dispensa se \textbf{tramita ante un juez} y recae sobre una prescripción ya cumplida.
\end{important}

\subsection{Los plazos de prescripción}

Se han visto distintos plazos de prescripción a lo largo de las clases. Por ejemplo, la acción por vicios tiene un plazo de prescripción de dos años. El plazo genérico es de cinco años.

\begin{articulo}{Artículo 2560 del Código Civil y Comercial}
El plazo de la prescripción es de cinco años, excepto que esté previsto uno diferente en la legislación local.
\end{articulo}

A partir de los artículos siguientes, el Código regula distintos plazos especiales, desde más largos hasta más cortos, según el tipo de acción de que se trate.

En general, el cómputo del plazo corre desde el día en que la prestación es \textbf{exigible}, salvo distintas excepciones:
\begin{itemize}
    \item en el caso de la violencia, el plazo comienza a correr desde que cesa la violencia;
    \item en la simulación entre las partes, el plazo corre desde que una de ellas desconoce el carácter simulado del acto.
\end{itemize}

\begin{important}
Es un tema importantísimo, porque es la fecha a partir de la cual se cuenta todo.
\end{important}

%------------------------------------------------------------------------------
\section{La caducidad en profundidad}
%------------------------------------------------------------------------------

La caducidad es también una forma de operar del transcurso del tiempo, pero produce la extinción de un \textbf{derecho no ejercido}, no de la acción, como ocurre en la prescripción.

Los plazos de caducidad están fijados específicamente por la ley: no todos los derechos tienen un plazo de caducidad, sino solamente aquellos respecto de los cuales la ley expresamente lo determina. Un ejemplo es el del artículo 442 sobre la compensación económica en materia de divorcio, cuya acción caduca a los seis meses de dictada la sentencia.

\begin{articulo}{Artículo 2567 del Código Civil y Comercial (régimen general de caducidad)}
Los plazos de caducidad no se suspenden ni se interrumpen, excepto disposición legal en contrario.
\end{articulo}

La caducidad no admite suspensión ni interrupción y, en general, sus plazos son más cortos que los de la prescripción.

Hay actos que impiden que se produzca la caducidad: el cumplimiento del acto previsto por la ley, la realización del acto jurídico correspondiente, o el reconocimiento del derecho por parte de quien podría beneficiarse de la caducidad.

Al igual que la prescripción, las normas sobre caducidad establecidas por la ley son, en general, \textbf{imperativas}: las partes no pueden renunciarla ni alterarla por disposición contractual cuando la caducidad proviene de la ley (sí es disponible cuando surge de la voluntad de las partes). % TODO: contenido incompleto o ambiguo en las notas originales (la caducidad se presenta como de fuente siempre legal en la introducción, pero aquí se menciona una caducidad que surge de la voluntad de las partes)

Cuando concurren, sobre un mismo supuesto, normas de caducidad y de prescripción, \textbf{priman las de caducidad}, y para lo no regulado se aplican, supletoriamente, las normas relativas a la prescripción.

%------------------------------------------------------------------------------
\section{Cuadro Cornell: prescripción y caducidad}
%------------------------------------------------------------------------------

{\renewcommand{\arraystretch}{1.2}
\begin{longtable}{
    >{\raggedright\arraybackslash}p{\cornellL}
    >{\raggedright\arraybackslash}p{\cornellR}
}
\toprule
\textbf{Preguntas / palabras clave} &
\textbf{Notas / respuestas} \\
\midrule
\endfirsthead

\toprule
\textbf{Preguntas / palabras clave} &
\textbf{Notas / respuestas} \\
\midrule
\endhead

\textbf{¿Qué es la prescripción?} &
El instituto por el cual el transcurso del tiempo opera la adquisición de un derecho (prescripción adquisitiva) o la extinción de la acción por inacción del titular (prescripción extintiva o liberatoria). Se trata en el Libro Sexto, Título Primero del CCyC. \\
\addlinespace

\textbf{¿Qué diferencia hay entre prescripción adquisitiva y extintiva?} &
La adquisitiva otorga el dominio por la posesión de una cosa (para inmuebles: veinte años, o diez con justo título y buena fe). La extintiva extingue la acción para reclamar un derecho, no el derecho en sí mismo. \\
\addlinespace

\textbf{¿Puede pagarse válidamente una deuda ya prescripta?} &
Sí. El pago espontáneo de una obligación prescripta no es repetible (art.~2538 CCyC): el derecho existía y sólo se había extinguido la acción para exigirlo judicialmente. Si se intenta reclamar un derecho prescripto, debe oponerse la excepción de prescripción. \\
\addlinespace

\textbf{¿Pueden las partes modificar las normas de prescripción?} &
No. Son imperativas y no pueden modificarse por convención de las partes (art.~2533 CCyC). La renuncia a la prescripción sólo puede hacerse una vez que esta ya se ganó. \\
\addlinespace

\textbf{¿Qué diferencia a la suspensión de la interrupción del cómputo?} &
La suspensión detiene el cómputo pero conserva el tiempo ya transcurrido (art.~2539 CCyC); la interrupción invalida todo lo transcurrido y reinicia el plazo desde cero. \\
\addlinespace

\textbf{¿Qué hechos suspenden la prescripción?} &
La interpelación fehaciente (seis meses, una sola vez, art.~2541 CCyC), el pedido de mediación (hasta veinte días después del acta de cierre, art.~2542 CCyC) y ciertas situaciones entre cónyuges, convivientes, personas con capacidad restringida y sus representantes, personas jurídicas y sus administradores, y herederos (art.~2543 CCyC). \\
\addlinespace

\textbf{¿Qué hechos interrumpen la prescripción?} &
El reconocimiento de deuda (art.~2545 CCyC), la petición judicial, interpretada en sentido amplio y aun si es defectuosa (art.~2546 CCyC), y la solicitud de arbitraje (art.~2548 CCyC). \\
\addlinespace

\textbf{¿Qué es la dispensa de la prescripción?} &
Una autorización judicial para ejercer la acción dentro de los seis meses siguientes a la cesación de los obstáculos, pese a que el plazo de prescripción ya se cumplió, cuando existió una imposibilidad de hecho justificada (por ejemplo, una internación) (art.~2550 CCyC). \\
\addlinespace

\textbf{¿Cuál es el plazo genérico de prescripción y desde cuándo se cuenta?} &
Cinco años (art.~2560 CCyC), salvo que esté previsto uno diferente; existen plazos especiales, más largos y más cortos. En general se cuenta desde que la prestación es exigible, con excepciones (violencia, simulación entre las partes). \\
\addlinespace

\textbf{¿Por qué la caducidad es más rígida que la prescripción?} &
Porque no admite suspensión ni interrupción (art.~2567 CCyC), sus plazos suelen ser más breves y extingue el derecho no ejercido, no sólo la acción. Cuando concurren normas de caducidad y de prescripción, priman las de caducidad. \\

\midrule
\multicolumn{2}{p{\dimexpr\cornellL+\cornellR+2\tabcolsep\relax}}{%
\textbf{Resumen:} el transcurso del tiempo puede hacer adquirir un derecho (prescripción adquisitiva) o extinguir la acción por la inacción del titular (prescripción extintiva), que deja subsistente un derecho sin acción. El cómputo puede detenerse (suspensión, que conserva el tiempo corrido), reiniciarse (interrupción, que lo invalida) o ser habilitado judicialmente una vez cumplido el plazo (dispensa). La caducidad extingue el derecho no ejercido en el plazo legal, de manera más rígida e imperativa.} \\
\bottomrule
\end{longtable}}

%==============================================================================

\section{Síntesis de la extinción de las relaciones jurídicas}
%==============================================================================

El tema recorre el cierre del ciclo de vida de las relaciones jurídicas: cómo terminan, después de haber nacido y desarrollado sus efectos. Ese final puede producirse:

\begin{itemize}
    \item por \textbf{hechos ajenos a la voluntad} de las partes: la muerte, la imposibilidad sobrevenida, la confusión o la compensación;
    \item por \textbf{actos jurídicos} mediante los cuales las partes deciden ponerle fin a la relación: el pago y la dación en pago, que la cumplen; la novación, que la reemplaza; la transacción y la renuncia, que ceden derechos; la remisión, que perdona una deuda; y la resolución, la rescisión y la revocación, que la dejan sin efecto, hacia atrás o hacia adelante según el caso;
    \item por el simple \textbf{transcurso del tiempo}, que castiga la inacción del titular de un derecho a través de la prescripción (que extingue la acción, pero no siempre el derecho) y de la caducidad (que extingue directamente el derecho no ejercido, de manera más rígida e imperativa).
\end{itemize}

El hilo conductor es la idea de que ningún derecho puede sostenerse en una incertidumbre indefinida: el ordenamiento jurídico necesita fijar, tarde o temprano, un punto final, ya sea porque los hechos de la vida lo imponen, porque las partes lo deciden o porque el tiempo, por sí solo, se encarga de hacerlo. Comprender estas reglas, y sobre todo los plazos y las causales de suspensión, interrupción y dispensa de la prescripción, es una de las herramientas de mayor relevancia práctica en el ejercicio cotidiano de la abogacía.

\partintro{Esta parte recorre el tramo del proceso civil que va desde la producción de los medios de prueba hasta el momento en que el juez se dispone a sentenciar. Los dos primeros capítulos abordan dos medios de prueba fundamentales del proceso civil y comercial argentino: la \textbf{prueba testimonial} (declaración de testigos) y la \textbf{prueba de confesión} (absolución de posiciones). En ambos casos se analiza el régimen legal aplicable conforme al Código Procesal Civil y Comercial de la Nación (CPCCN), los deberes y derechos de los declarantes, las formas válidas de ofrecimiento, producción y valoración de la prueba, y los supuestos de caducidad que operan objetivamente ante la inactividad procesal de las partes. Se incluyen referencias al Código Civil y Comercial de la Nación (CCyC), jurisprudencia de la Corte Suprema de Justicia de la Nación (CSJN) y ejemplos prácticos del ejercicio forense.

El proceso judicial puede compararse con una investigación periodística: primero hay que conseguir los testimonios de quienes presenciaron los hechos (prueba testimonial) y luego obtener reconocimientos de los propios involucrados (prueba de confesión). Cada medio de prueba sigue tres etapas: \textbf{ofrecimiento}, \textbf{producción} y \textbf{valoración}.

Los tres capítulos siguientes se ocupan de la prueba pericial, la prueba indiciaria y el reconocimiento judicial. Luego se estudia cómo concluye la etapa probatoria, lo que conduce a los alegatos y a la etapa conclusional, y se cierra con los sistemas de valoración de la prueba, es decir, con el modo en que el juez aprecia la prueba producida al dictar sentencia.}
\part{La prueba y la conclusión de la causa}

\chapter[Prueba testimonial]{Prueba testimonial (arts.~426 a 456 CPCCN)}\label{cap:testimonial}

% =====================================================

La prueba testimonial comienza a regularse a partir del artículo 426 del CPCCN.

\section{Concepto de testigo}

\begin{definition}{Testigo}
El \textbf{testigo} es un tercero que no es parte en el proceso, convocado para que preste declaración sobre hechos de su conocimiento. Es convocado durante la etapa probatoria para que relate lo que sabe sobre los hechos controvertidos o ventilados en la causa.
\end{definition}

\section{Objeto del testimonio y valor según el tipo de percepción}

El objeto del testimonio son los \textbf{hechos alegados por las partes} en sus escritos postulatorios. El testigo puede conocer esos hechos a través de distintas vías:
\begin{itemize}
    \item haberlos visto directamente;
    \item haberlos escuchado;
    \item haberlos conocido a través del relato de otra persona.
\end{itemize}

No todos los testimonios tienen la misma eficacia probatoria: no es lo mismo quien vio un hecho directamente que quien lo conoce de oídas. La valoración del juez al dictar sentencia será distinta en cada caso. Sin embargo, todos esos tipos de testimonio son admisibles.

\begin{norma}{Artículo 426 CPCCN --- Procedencia de la prueba testimonial}
\textit{Toda persona mayor de 14 años podrá ser propuesta como testigo y tendrá el deber de comparecer y declarar, salvo las excepciones establecidas por la ley.}
\end{norma}

\section{Deberes del testigo}

Sobre todo testigo recaen tres deberes:

\begin{center}
\renewcommand{\arraystretch}{1.3}
\begin{tabularx}{\textwidth}{
    >{\raggedright\arraybackslash}p{0.17\textwidth}
    >{\raggedright\arraybackslash}X
    >{\raggedright\arraybackslash}X}
\toprule
\textbf{Deber} & \textbf{Alcance} & \textbf{Excepciones} \\
\midrule
\textbf{Comparecer} & Es una carga pública. Toda persona convocada como testigo debe presentarse ante el juzgado. & Cargo que implique declarar por escrito; causa justificada debidamente acreditada; irregularidad en la notificación (menos de 3 días de anticipación). \\
\addlinespace
\textbf{Declarar} & El testigo tiene la obligación de prestar declaración testimonial. & Secreto profesional (militar, científico, artístico); declaración que perjudique el honor del testigo o lo exponga a enjuiciamiento penal. \\
\addlinespace
\textbf{Decir verdad} & Antes de declarar, el testigo presta juramento o promesa de decir verdad. Se le informan las consecuencias penales del falso testimonio. & No hay excepciones al juramento, pero si el testigo miente incurre en el delito de falso testimonio. \\
\bottomrule
\end{tabularx}
\end{center}

\subsection*{Regla de los 70 kilómetros}

Respecto del deber de comparecer, si el domicilio del testigo se encuentra dentro de un radio de \textbf{70 kilómetros} del tribunal, tiene la obligación de concurrir personalmente. Si su domicilio está fuera de ese radio, se lo convoca mediante la \textbf{Ley 22.172}, a través de un oficio dirigido al juzgado de turno en la sede de su domicilio. Ese oficio debe contener los datos personales del testigo y el interrogatorio completo, ya que la parte contraria también tiene derecho a controlar este medio de prueba y a formular sus propias preguntas.

\section{Testigos excluidos: art. 427 CPCCN y vinculación con el CCyC}

El artículo 427 establece que no pueden ser ofrecidos como testigos los consanguíneos o afines en línea directa de las partes, ni el cónyuge aunque estuviere separado legalmente, salvo en los casos de reconocimiento de firmas. Los parientes colaterales (hermanos, tíos, sobrinos) sí pueden ser testigos.

\begin{norma}{Artículo 427 CPCCN --- Testigos excluidos}
\textit{No podrán ser ofrecidos como testigos los consanguíneos o afines en línea directa de las partes, ni el cónyuge aunque estuviere separado legalmente, salvo si se trata de reconocimiento de firmas.}
\end{norma}

\begin{norma}{Artículo 711 del Código Civil y Comercial de la Nación --- Testigos en juicios de familia}
\textit{Los parientes y allegados a las partes pueden ser ofrecidos como testigos en juicios de familia.}

El fundamento es que, en los conflictos familiares, muchas veces quienes mejor conocen la problemática son precisamente los parientes o allegados, y excluirlos generaría una imposibilidad de probar los hechos alegados (por ejemplo, en casos de violencia familiar).
\end{norma}

\section{Los tres momentos de la prueba testimonial}

Es importante distinguir tres momentos en todo medio probatorio:
\begin{itemize}
    \item \textbf{Ofrecimiento}: se realiza en los escritos postulatorios (demanda y contestación de demanda).
    \item \textbf{Producción}: se lleva a cabo durante la etapa probatoria, una vez abierta la causa a prueba.
    \item \textbf{Valoración}: las partes valoran la prueba en los alegatos; el juez la valora al dictar sentencia.
\end{itemize}

\section{Ofrecimiento de la prueba testimonial}

Para ofrecer un testigo en el escrito postulatorio, el código exige informar:
\begin{itemize}
    \item nombre completo del testigo;
    \item ocupación;
    \item domicilio.
\end{itemize}

El \textbf{artículo 333 CPCCN} agrega un elemento adicional: argumentar el \textbf{objeto de la citación}, es decir, explicar sintéticamente por qué se convoca a esa persona y sobre qué hechos declarará.

\begin{example}{Objeto de la citación}
``Se ofrece a Juan Pérez porque estuvo presente en el momento del accidente'', o ``porque es el médico que estuvo a cargo del tratamiento de la parte actora''.
\end{example}

Se pueden ofrecer \textbf{hasta ocho testigos}. Si la complejidad del caso lo amerita, pueden admitirse más, previa decisión judicial sobre la necesidad de la ampliación.

\section{Producción de la prueba testimonial: citación a audiencia}

Una vez que la causa está abierta a prueba y el juez ha fijado las audiencias de testigos (en la etapa probatoria, luego de la audiencia preliminar), se debe notificar a los testigos. Los medios de notificación admitidos son:
\begin{itemize}
    \item cédula de notificación;
    \item telegrama colacionado;
    \item acta notarial;
    \item carta documento.
\end{itemize}

La notificación debe cursarse con al menos \textbf{tres días} de anticipación a la audiencia. La cédula debe informar tanto la \textbf{audiencia principal} (primera) como la \textbf{audiencia supletoria} (segunda), con la advertencia de que si el testigo no concurre a la primera sin causa justificada, será convocado a la segunda con citación de la fuerza pública (es decir, a través de la policía).

\section{Desarrollo de la audiencia testimonial e interrogatorio preliminar}

Al comenzar la audiencia, el testigo se identifica con su documento nacional de identidad y brinda sus datos personales. Luego se le toma juramento o promesa de decir verdad, explicándole las consecuencias penales del falso testimonio. Inmediatamente después, se procede al interrogatorio preliminar.

\begin{norma}{Artículo 441 CPCCN --- Interrogatorio preliminar (generales de la ley)}
\textit{Aunque las partes no lo pidan, los testigos serán siempre preguntados por: 1) Su nombre, edad, estado, profesión y domicilio; 2) Si es pariente por consanguinidad o afinidad de alguna de las partes y en qué grado; 3) Si tiene interés directo o indirecto en el pleito; 4) Si es amigo íntimo o enemigo de alguna de las partes; 5) Si es dependiente, acreedor o deudor de algunos de los litigantes, o si tiene algún otro género de relación con ellos o algún interés en el pleito.}
\end{norma}

Estas preguntas se denominan, en el lenguaje forense, \textbf{``generales de la ley''}. Su propósito es evaluar si el testigo puede estar afectado por alguna de estas circunstancias, lo que podría incidir en la valoración de su testimonio e incluso en su idoneidad (conforme al art.~456 CPCCN).

\section{Forma de las preguntas al testigo: art. 443}

Luego del interrogatorio preliminar, cada parte puede formular su interrogatorio al testigo, que puede presentarse por escrito o realizarse a viva voz. El artículo 443 establece los requisitos que deben cumplir las preguntas:

\begin{norma}{Artículo 443 CPCCN --- Forma de las preguntas}
\textit{Las preguntas no contendrán más de un hecho, serán claras y concretas. No se formularán las que estén concebidas en términos afirmativos, sugieran la respuesta o sean ofensivas o vejatorias. No podrán contener referencias de carácter técnico, salvo si fuesen dirigidas a personas especializadas.}
\end{norma}

En la práctica forense, la pregunta se formula con la siguiente fórmula ritual: \textit{``Para que jure, diga el testigo si sabe y cómo le consta\ldots{} [el hecho sobre el que se pregunta]''}.

La pregunta debe ser abierta y no sugestiva porque el objetivo es conocer genuinamente lo que el testigo sabe y cómo lo sabe, no confirmar una versión ya armada. Por ejemplo, si se pregunta si el testigo sabe qué sucedió tal día en tal lugar, se busca que el testigo relate libremente, no que responda ``sí'' o ``no'' a una afirmación del abogado.

Si una pregunta viola los requisitos del artículo 443, la parte contraria puede oponerse. Esa oposición genera una \textbf{incidencia} dentro de la audiencia: si el juez hace lugar a la oposición, la pregunta queda rechazada y las costas de la incidencia se imponen a quien formuló la pregunta defectuosa.

\section{Testigo técnico y perito}

La mención a las referencias técnicas del artículo 443 permite distinguir dos figuras que no deben confundirse:

\begin{center}
\renewcommand{\arraystretch}{1.3}
\begin{tabularx}{\textwidth}{
    >{\raggedright\arraybackslash}p{0.2\textwidth}
    >{\raggedright\arraybackslash}X
    >{\raggedright\arraybackslash}X}
\toprule
 & \textbf{Testigo técnico} & \textbf{Perito} \\
\midrule
\textbf{Fungibilidad} & No es fungible. Es una persona específica que conoció los hechos. & Es fungible. Se designa por sorteo de oficio por el juez. \\
\addlinespace
\textbf{Conocimiento del hecho} & Conoció los hechos en el pasado, cuando ocurrieron. & Toma conocimiento del caso en el presente, al momento de su designación. \\
\addlinespace
\textbf{Ejemplo} & El médico que trató al paciente (parte actora) durante su enfermedad. & El perito médico designado por sorteo para realizar la pericia médica en el proceso. \\
\bottomrule
\end{tabularx}
\end{center}

Al testigo técnico, que es una persona especializada, sí se le pueden hacer preguntas de carácter técnico. Al perito, en cambio, se le formulan puntos de pericia en los escritos postulatorios.

\section{Negativa a responder: art. 444}

\begin{norma}{Artículo 444 CPCCN --- Negativa a responder}
\textit{El testigo podrá rehusarse a contestar las preguntas si la respuesta lo expusiera a enjuiciamiento penal o comprometiera su honor, o si no pudiera responder sin revelar un secreto profesional, militar, científico, artístico o industrial.}
\end{norma}

\section{Caducidad de la prueba testimonial: arts. 432 y 437}

Debe distinguirse la negligencia de la caducidad en la producción de la prueba. La \textbf{negligencia} se basa en un elemento subjetivo: el desinterés de la parte en producir su prueba a lo largo del período probatorio. La \textbf{caducidad}, en cambio, opera de forma objetiva: el código establece supuestos específicos que, al configurarse, producen automáticamente la pérdida del derecho a esa prueba.

\begin{norma}{Artículo 432 CPCCN --- Caducidad de la prueba testimonial}
\textit{A pedido de parte y sin sustanciación, se tendrá por desistido del testigo a la parte que lo propuso, en los siguientes supuestos: a) Si no hubiera activado la citación del testigo y este no hubiera comparecido por esta razón; b) Si no habiendo comparecido el testigo a la primera audiencia sin invocar causa justificada, no requiriera oportunamente las medidas de compulsión necesarias; c) Si fracasara la segunda audiencia por motivos no imputables a la parte y no solicitara nueva audiencia dentro del quinto día.}
\end{norma}

\begin{norma}{Artículo 437 CPCCN --- Cuarto supuesto de caducidad}
\textit{Si la parte que ofreció al testigo no concurriera a la audiencia por sí o por apoderado y no hubiere dejado interrogatorio, se la tendrá por desistida de aquel sin sustanciación.}
\end{norma}

\begin{center}
\renewcommand{\arraystretch}{1.3}
\begin{tabularx}{\textwidth}{
    >{\centering\arraybackslash}p{0.06\textwidth}
    >{\raggedright\arraybackslash}X
    >{\raggedright\arraybackslash}X
    >{\raggedright\arraybackslash}p{0.17\textwidth}}
\toprule
\textbf{N.º} & \textbf{Supuesto} & \textbf{Descripción} & \textbf{Norma} \\
\midrule
1 & No se notificó al testigo y no compareció. & La parte no libró la cédula ni ningún medio de notificación al testigo. & Art. 432 inc.~a) \\
\addlinespace
2 & El testigo no compareció a la 1.ª audiencia y no se solicitaron medidas de compulsión. & El testigo faltó sin causa y la parte no gestionó la segunda audiencia con fuerza pública. & Art. 432 inc.~b) \\
\addlinespace
3 & Fracasó la 2.ª audiencia y no se pide una nueva en 5 días. & La segunda audiencia fracasó por causas no imputables a la parte, pero esta no solicitó nueva audiencia dentro del quinto día. & Art. 432 inc.~c) \\
\addlinespace
4 & La parte oferente no concurrió y no dejó interrogatorio. & El testigo sí compareció, pero la parte que lo ofreció no se presentó ni dejó el pliego de preguntas. & Art. 437 \\
\bottomrule
\end{tabularx}
\end{center}

El desistimiento del testigo por quien lo ofreció es válido, sin necesidad de consentimiento de la contraparte, siempre que sea previo a la celebración de la audiencia.

\section{Idoneidad del testigo: art. 456}

\begin{norma}{Artículo 456 CPCCN --- Idoneidad de los testigos}
\textit{Dentro del plazo de prueba, las partes podrán alegar y probar acerca de la idoneidad de los testigos. El juez apreciará, según las reglas de la sana crítica y en oportunidad de dictar sentencia definitiva, las circunstancias y motivos que corroboren o disminuyan la fuerza de las declaraciones.}
\end{norma}

Este artículo permite a las partes generar un \textbf{incidente} por separado del proceso principal para demostrar que un testigo no era idóneo para declarar. Se puede ofrecer en ese incidente toda la prueba que se considere necesaria. Es un mecanismo de defensa ante una declaración que, si bien era formalmente admisible, reveló en su contenido elementos que cuestionan la imparcialidad o credibilidad del testigo.

\cornelltitle{prueba testimonial}
\cornellblock{
\cqrow{¿Quién es un testigo?}{Es un tercero ajeno al proceso, convocado en la etapa probatoria para declarar sobre hechos de su conocimiento, ya sea porque los vio, los escuchó o los conoció a través del relato de terceros. No es parte del proceso.
}
\cqrow{¿Cuáles son los deberes del testigo?}{Tres deberes: 1) comparecer (carga pública; se lo puede citar con fuerza pública en caso de incumplimiento); 2) declarar (salvo secreto profesional o riesgo de autoincriminación); 3) decir verdad (juramento previo; el falso testimonio es delito penal).
}
\cqrow{¿Quiénes están excluidos como testigos?}{Conforme al art.~427, los consanguíneos o afines en línea directa y el cónyuge. Excepción: en juicios de familia (art.~711 CCyC) se admiten parientes y allegados, porque son quienes mejor conocen la problemática familiar.
}
\cqrow{¿Cómo se ofrece la prueba testimonial?}{En los escritos postulatorios, indicando nombre completo, ocupación y domicilio del testigo, más el objeto de la citación (art.~333). Se pueden ofrecer hasta 8 testigos; el juez puede admitir más si es necesario.
}
\cqrow{¿Qué son las ``generales de la ley''?}{Son las preguntas del interrogatorio preliminar (art.~441), que se formulan a todo testigo antes del interrogatorio de fondo: datos personales, parentesco con las partes, interés en el pleito, amistad o enemistad, relación de dependencia. Permiten evaluar la idoneidad e imparcialidad del testigo.
}
\cqrow{¿Cómo deben formularse las preguntas a un testigo?}{Deben contener un solo hecho y ser claras y concretas; NO estar formuladas en términos afirmativos, NO sugerir la respuesta, NO ser vejatorias u ofensivas y NO contener referencias técnicas (salvo al testigo técnico). Se formulan así: ``Para que jure, diga el testigo si sabe y cómo le consta\ldots{}''.
}
\cqrow{¿Qué diferencia hay entre testigo técnico y perito?}{El testigo técnico conoció los hechos en el pasado (por ejemplo, el médico tratante) y no es fungible. El perito es fungible, es designado por sorteo de oficio, toma conocimiento en el presente y dictamina sobre puntos de pericia formulados en los escritos postulatorios. Al primero se le pueden hacer preguntas técnicas; al segundo se le formulan puntos de pericia.
}
}
\medskip
\cornellblocksum{
\cqrow{¿Cuáles son los supuestos de caducidad de la prueba testimonial?}{1) No se notificó al testigo (art.~432 a); 2) no se gestionaron medidas de compulsión tras la 1.ª audiencia fallida (art.~432 b); 3) no se pidió nueva audiencia en 5 días tras fracasar la 2.ª (art.~432 c); 4) la parte oferente no concurrió ni dejó interrogatorio (art.~437). Los 3 primeros son objetivos; todos operan a pedido de parte.
% TODO: contenido incompleto o ambiguo en las notas originales (alcance de "objetivos" y "a pedido de parte" respecto del art. 437)
}
}{La prueba testimonial convoca a terceros ajenos al proceso para que declaren sobre hechos de su conocimiento, imponiéndoles los deberes de comparecer, declarar y decir verdad bajo juramento. Se ofrece en los escritos postulatorios, se produce en la etapa probatoria y se valora en los alegatos y en la sentencia. El CPCCN regula minuciosamente el modo de formular las preguntas (art.~443), el interrogatorio preliminar (art.~441) y los supuestos de caducidad (arts.~432 y 437).
}

\chapter[Prueba de confesión]{Prueba de confesión (arts.~404 a 424 CPCCN)}\label{cap:confesion}

% =====================================================

\section{Concepto de confesión: confesión espontánea y provocada}

\begin{definition}{Confesión}
La confesión como prueba consiste en una \textbf{declaración formulada por quien es parte en el proceso}, sobre hechos personales o de su conocimiento, que produce un efecto desfavorable al confesante y favorable para la otra parte.
\end{definition}

Se distingue entre:
\begin{itemize}
    \item \textbf{Confesión espontánea}: ocurre cuando la parte demandada reconoce un hecho en la contestación de demanda. Al reconocer ese hecho, genera plena prueba sobre él, sin necesidad de ningún otro medio probatorio.
    \item \textbf{Confesión provocada (absolución de posiciones)}: es el medio de prueba propiamente dicho. Se produce en audiencia, donde la parte absuelve las posiciones formuladas por la contraria.
\end{itemize}

La confesión en sí no es el medio de prueba: el medio de prueba es la absolución de posiciones, o confesión provocada en juicio.

\section{Momento del ofrecimiento y de la producción}

Al igual que la prueba testimonial, la confesión se ofrece en los escritos postulatorios. Su producción --- la absolución de posiciones --- se lleva a cabo en la audiencia preliminar, conforme al artículo 360 inciso 4 del CPCCN. En la práctica, muchos jueces fijan una audiencia posterior específicamente para este fin, aunque la norma establece la audiencia preliminar como el momento de producción.

\section{Quiénes pueden ser citados a absolver posiciones: art. 405}

\begin{norma}{Artículo 405 CPCCN --- Personas que pueden ser citadas}
\textit{Podrán asimismo ser citados a absolver posiciones, además del actor y el demandado: los representantes de los incapaces por los hechos en que hayan intervenido personalmente en ese carácter; los apoderados por hechos realizados en nombre de sus mandantes, estando vigente el mandato, y por hechos anteriores cuando estuvieran sus representados fuera del lugar en que se sigue el juicio, siempre que el apoderado tuviese facultades para ello y la parte contraria lo consienta; y los representantes legales de las personas jurídicas, sociedades o entidades colectivas que tuviesen facultad para obligarlas.}
\end{norma}

\section{Elección del absolvente en personas jurídicas: art. 406}

Para las empresas y personas jurídicas, el artículo 406 permite oponerse, dentro del quinto día de notificada la audiencia, a que absuelva el representante elegido por el ponente, siempre que:
\begin{enumerate}
    \item se alegue que ese representante no intervino personalmente ni tuvo conocimiento directo de los hechos;
    \item se indique en el mismo escrito el nombre del representante que absolverá las posiciones;
    \item se deje constancia de que esa persona fue notificada de la audiencia y presta su conformidad.
\end{enumerate}

\section{Citación: domicilio constituido y domicilio real}

La notificación de la audiencia de absolución de posiciones debe hacerse por cédula. La notificación varía según la forma de actuación de la parte en el proceso:

\begin{center}
\renewcommand{\arraystretch}{1.3}
\begin{tabularx}{\textwidth}{
    >{\raggedright\arraybackslash}X
    >{\raggedright\arraybackslash}X
    >{\raggedright\arraybackslash}X}
\toprule
\textbf{Modalidad} & \textbf{Tipo de actuación} & \textbf{Domicilio donde se notifica} \\
\midrule
Parte con letrado patrocinante (actúa por derecho propio) & Firma sus presentaciones junto con su abogado patrocinante. & Domicilio constituido (generalmente el estudio del abogado). \\
\addlinespace
Parte representada por apoderado & El apoderado actúa en nombre de la parte. & Domicilio real de la parte (no el domicilio procesal del apoderado). \\
\bottomrule
\end{tabularx}
\end{center}

La cédula debe entregarse con al menos \textbf{tres días} de anticipación. En caso de urgencia debidamente justificada, el juez puede reducir ese plazo mediante resolución, pero la anticipación no puede ser inferior a un día.

\section{Reserva del pliego de posiciones: art. 410}

\begin{norma}{Artículo 410 CPCCN --- Reserva del pliego de posiciones}
\textit{El pliego deberá ser entregado en secretaría hasta media hora antes de la hora fijada para la audiencia, en sobre cerrado, al que se le pondrá cargo. Si la parte que pidió las posiciones no comparece sin justa causa a la audiencia, no hubiese dejado pliego, y comparece el citado, perderá el derecho a exigirlas.}
\end{norma}

El sobre cerrado se presenta en la mesa de entradas del juzgado hasta media hora antes del inicio de la audiencia. El fundamento de este límite temporal es evitar que la parte ponente prepare dos pliegos distintos (uno para el caso de que el absolvente comparezca y otro para el caso de incomparecencia) y elija el que más le convenga según las circunstancias.

\section{Forma de redactar las posiciones: art. 411}

\begin{norma}{Artículo 411 CPCCN --- Forma de las posiciones}
\textit{Las posiciones serán claras y concretas, no contendrán más de un hecho, serán redactadas en forma afirmativa y deberán versar sobre puntos controvertidos que se refieran a la actuación personal del absolvente. El juez podrá modificar de oficio, sin recurso alguno, el orden y los términos de las posiciones propuestas por las partes sin alterar su sentido, y podrá asimismo eliminar las que fueran manifiestamente inútiles o superfluas.}
\end{norma}

La diferencia fundamental con el interrogatorio de testigos es que, mientras las preguntas a los testigos deben ser abiertas y no sugestivas, las posiciones de la confesión deben formularse en términos afirmativos, de manera que el absolvente pueda responder simplemente ``sí'' o ``no'' (aunque luego puede agregar las aclaraciones que estime pertinentes).

La fórmula es: \textit{``Para que jure como que es cierto que usted se encontraba en tal lugar tal día''}. Eso es la posición.

Un aspecto crucial es que el ponente, al formular la posición, \textbf{está reconociendo ese hecho}. Por lo tanto:
\begin{itemize}
    \item si el absolvente contesta ``sí, es cierto'': hay plena prueba sobre ese hecho (ambas partes lo reconocen);
    \item si el absolvente contesta ``no es cierto'': ese hecho deberá probarse por otro medio.
\end{itemize}

\section{Posición impertinente: art. 414}

\begin{norma}{Artículo 414 CPCCN --- Posición impertinente}
\textit{La parte podrá negarse a contestar la posición si esta fuera impertinente. El juez valorará esa oposición al dictar sentencia y, si considera que la posición era pertinente, podrá tener por confeso al absolvente en ese punto.}
\end{norma}

A diferencia de lo que ocurre en la prueba testimonial (donde la oposición a una pregunta genera una incidencia con fundamentos), en la confesión la oposición a una posición no requiere fundamentación: el abogado del absolvente simplemente manifiesta que se opone en los términos del artículo 414 y eso queda asentado en el acta. Sin embargo, si el juez al dictar sentencia considera que la posición era pertinente, tendrá por confeso al absolvente, lo que puede tener consecuencias muy significativas para el resultado del proceso.

\section{Interrogatorio de las partes: art. 415}

\begin{norma}{Artículo 415 CPCCN --- Interrogatorio de oficio}
\textit{El juez podrá interrogar de oficio a las partes en cualquier estado del proceso, y estas podrán hacerse recíprocamente las preguntas y observaciones que juzgaren convenientes en la audiencia que corresponda, siempre que el juez no las declare superfluas o improcedentes.}
\end{norma}

La posibilidad de que las partes se formulen preguntas recíprocas en cualquier audiencia, o de que el juez interrogue de oficio, es un instrumento muy valioso para esclarecer la verdad de los hechos, pero en la práctica se utiliza con poca frecuencia. Con la reforma procesal y el proceso por audiencias este panorama debería cambiar.

\section{Confesión ficta: art. 417}

\begin{norma}{Artículo 417 CPCCN --- Confesión ficta}
\textit{Si el citado no compareciera a declarar dentro de la media hora de fijada para la audiencia, o si habiendo comparecido rehusara a responder o respondiera de una manera evasiva, el juez al sentenciar lo tendrá por confeso sobre los hechos personales, teniendo en cuenta las circunstancias de la causa y las demás pruebas producidas. En caso de incomparecencia del absolvente, aunque no se hubiera extendido acta, se aplicará lo establecido en este párrafo, si el ponente hubiera presentado oportunamente el pliego de posiciones y el absolvente estuviere debidamente notificado.}
\end{norma}

Para que se configure la confesión ficta deben reunirse los siguientes requisitos:
\begin{itemize}
    \item el absolvente debe estar debidamente notificado de la audiencia (con todos los requisitos de forma ya explicados);
    \item el ponente debe haber presentado el pliego en sobre cerrado hasta media hora antes;
    \item el absolvente no comparece dentro de la media hora de la hora fijada, o se niega a contestar, o responde de manera evasiva.
\end{itemize}

Configurados esos supuestos, el juez al dictar sentencia puede tener por confeso al absolvente. La norma aclara que el juez tomará en cuenta ``las circunstancias de la causa y las demás pruebas producidas'': no es una consecuencia automática absoluta, sino que debe apreciarse en el contexto del proceso.

\section{Valor probatorio de la confesión ficta}

Existen distintas interpretaciones doctrinarias y jurisprudenciales sobre el valor de la confesión ficta:

\begin{center}
\renewcommand{\arraystretch}{1.3}
\begin{tabularx}{\textwidth}{
    >{\raggedright\arraybackslash}p{0.25\textwidth}
    >{\raggedright\arraybackslash}X}
\toprule
\textbf{Posición doctrinaria} & \textbf{Contenido} \\
\midrule
\textbf{Primer sector} & La confesión ficta importa plena prueba siempre que no haya otro elemento de juicio que la contraríe. \\
\addlinespace
\textbf{Segundo sector} & La confesión ficta es plena prueba siempre que haya otros medios de prueba que colaboren y la corroboren. \\
\addlinespace
\textbf{Tercer sector} & No tiene valor absoluto en ningún supuesto: siempre debe conjugarse la confesión ficta con el resto de las constancias y medios de prueba del expediente. \\
\bottomrule
\end{tabularx}
\end{center}

\begin{norma}{Fallo de la CSJN --- 6 de febrero de 2001 (LL 2001-B, pág. 959)}
La CSJN habilitó excepcionalmente el recurso extraordinario en una cuestión de prueba --- materia normalmente ajena a la instancia extraordinaria por ser de derecho común --- y revocó la sentencia de cámara que había desestimado el valor de la confesión ficta. El caso involucraba a una empresa de transportes que había sido tenida por confesa en audiencia de absolución de posiciones, y existían además otros elementos de juicio que corroboraban esa conclusión.
\end{norma}

Desde esta perspectiva, no habría razón para que una ficción legal (la configuración de la confesión ficta) tenga más valor que la realidad que surge del resto de las constancias de la causa. Desde un enfoque sistémico del proceso, el juez al dictar sentencia debe valorar la totalidad de las constancias de la causa, sin que la confesión ficta sea automáticamente superior a la realidad probatoria del expediente.

\section{Confesión expresa: art. 423}

\begin{norma}{Artículo 423 CPCCN --- Efectos de la confesión judicial expresa}
\textit{La confesión judicial expresa constituirá plena prueba, salvo cuando: 1) dicho medio de prueba estuviere excluido por la ley respecto de los hechos que constituyen el objeto del juicio, o sobre derechos que el confesante no puede renunciar o transigir válidamente; 2) recayere sobre hechos cuya investigación prohíbe la ley; 3) se opusiere a las constancias de instrumentos fehacientes de fecha anterior agregados al expediente.}
\end{norma}

La confesión judicial expresa constituye plena prueba, salvo en esos tres supuestos. Es decir, cuando hay confesión expresa no es necesario recurrir a la valoración del resto de la prueba (salvo las excepciones). Esta es la diferencia fundamental con la confesión ficta: la ficta es una creación legal que opera ante la inactividad del absolvente; la expresa es una declaración voluntaria y directa del confesante sobre un hecho.

\section{Confesión extrajudicial: art. 424}

La confesión extrajudicial es aquella que se produce fuera del proceso judicial propiamente dicho. Se mencionan dos ejemplos:
\begin{itemize}
    \item \textbf{Declaración ante autoridad policial}: tiene valor y puede hacerse valer en el proceso.
    \item \textbf{Manifestaciones ante el perito designado judicialmente}: si la parte, durante la realización de una pericia (por ejemplo, una pericia médica), le hace saber al perito algún hecho relevante, esas manifestaciones también pueden tener valor confesional.
\end{itemize}

Esto fue objeto de debate, pero se considera que no afecta el derecho de defensa porque ocurre en el marco del proceso y todas las partes tienen posibilidad de controlarlo. En todo caso, el juez lo valorará en la sentencia junto con el resto de los medios de prueba.

\begin{norma}{Artículo 424 CPCCN --- Alcance de la confesión}
\textit{En caso de duda, la confesión deberá interpretarse en favor de quien la hace. La confesión es indivisible, salvo cuando el confesante invoca hechos impeditivos, modificativos o extintivos, o absolutamente separables e independientes; o cuando las circunstancias calificativas expuestas por quien confiesa fueren contrarias a una presunción legal o inverosímiles, y las modalidades del caso hicieran procedente la divisibilidad.}
\end{norma}

\section{Ausencia del absolvente por enfermedad y regla de los 300 kilómetros}

Si el absolvente no puede concurrir a la audiencia por enfermedad, el artículo 418 permite que el juez o un funcionario del juzgado investido de autoridad se traslade al domicilio o lugar donde se encuentre el absolvente para tomar la declaración. La enfermedad debe justificarse con la debida anticipación y por escrito (art.~419).

En cuanto a la distancia, existe una regla análoga a la de los 70 kilómetros para testigos: si el absolvente domicilia dentro de los \textbf{300 kilómetros} del juzgado, debe concurrir ante el juez de la causa. Si su domicilio está a más de 300 kilómetros, se aplica la Ley 22.172 y se libra un oficio al juzgado competente en la zona del domicilio del absolvente, adjuntando el pliego de posiciones.

\section{Desarrollo de la audiencia de absolución de posiciones}

Al momento de celebrarse la audiencia, se procede a la apertura del sobre que contiene el pliego de posiciones. Si ambas partes ofrecieron el medio de prueba, se absolverán primero las posiciones de una de ellas y luego las de la otra. La audiencia puede ser grabada o puede levantarse un acta escrita. Las posiciones (que ya estaban en el sobre cerrado) se leen al absolvente, quien las contesta por sí o por no, pudiendo agregar las aclaraciones que estime pertinentes.

\cornelltitle{prueba de confesión}
\cornellblocksum{
\cqrow{¿En qué consiste la confesión como prueba?}{En una declaración de quien es parte sobre hechos personales o de su conocimiento, con efecto desfavorable al confesante y favorable a la contraria. Puede ser espontánea (en la contestación de demanda) o provocada (absolución de posiciones en audiencia).
}
\cqrow{¿Cómo se redactan las posiciones?}{En forma afirmativa (``Para que jure como que es cierto que usted se encontraba en\ldots{}''), con un solo hecho, sobre la actuación personal del absolvente y de manera que pueda contestarse por sí o por no. El ponente, al formularlas, reconoce el hecho.
}
\cqrow{¿Qué es la confesión ficta?}{Es la ficción legal del art.~417: si el absolvente, debidamente notificado, no comparece, se niega a contestar o responde evasivamente, el juez al dictar sentencia puede tenerlo por confeso en los hechos de las posiciones. No es automática: el juez evalúa las demás constancias de la causa.
}
\cqrow{¿Qué valor probatorio tiene la confesión ficta frente a la confesión expresa?}{La confesión judicial EXPRESA constituye plena prueba (art.~423), salvo tres excepciones legales. La confesión FICTA tiene valor discutido: para un sector doctrinal es plena prueba si no hay elementos en contra; para otro, requiere corroboración de otros medios; para un tercero, siempre debe cotejarse con el resto de la prueba. La CSJN (fallo del 6/2/2001, LL 2001-B, p.~959) revocó una sentencia que había desestimado la confesión ficta cuando había otros elementos corroborantes.
}
\cqrow{¿Cuándo opera la caducidad en la prueba de confesión?}{Cuando el ponente no comparece a la audiencia, no deja el pliego y el absolvente sí comparece: el ponente pierde el derecho a exigir la absolución de posiciones (art.~410 in fine).
}
}{La prueba de confesión provoca una declaración de las propias partes sobre hechos personales: puede ser espontánea (reconocimiento en la contestación de demanda, que genera plena prueba de inmediato) o provocada a través de la absolución de posiciones, cuya forma afirmativa y específica la distingue de la testimonial. La confesión ficta (art.~417) opera ante la incomparecencia o evasiva del absolvente notificado y, si bien genera una presunción legal desfavorable, su valoración debe conjugarse con el resto de las constancias de la causa, tal como lo señaló la Corte Suprema en el fallo del 6 de febrero de 2001. En conjunto, ambos medios de prueba ilustran la tensión permanente del proceso civil entre el formalismo procedimental y la búsqueda de la verdad objetiva de los hechos.
}

\chapter[La prueba pericial]{La prueba pericial: concepto, sujetos y procedimiento}\label{cap:pericial}

Este capítulo y los dos siguientes abordan tres medios de prueba del proceso civil argentino: la \textbf{prueba pericial}, la \textbf{prueba indiciaria} (con las presunciones legales y las presunciones \textit{iuris hominis}) y el \textbf{reconocimiento judicial}.

Respecto de la prueba pericial se estudian su concepto, los sujetos involucrados (el perito designado de oficio y el consultor técnico de parte), el procedimiento completo desde el ofrecimiento hasta la valoración del dictamen, las obligaciones y derechos del perito, las causales de recusación y remoción, el anticipo de gastos y las consecuencias del incumplimiento. Luego se analiza la prueba indiciaria, es decir, cómo el juez puede formar su convicción a partir de indicios acreditados en el expediente cuando no existe prueba directa. Finalmente se estudia el reconocimiento judicial como la prueba más directa: su concepto, su procedimiento y su vinculación con las facultades del artículo 475 del CPCCN.

\textbf{Normas principales:} arts.\ 459, 163 inc.\ 5, 386, 476, 477 y 479 del CPCCN; art.\ 243 del Código Penal.

\begin{example}{Cómo se conectan los tres medios de prueba}
En un juicio por accidente de tránsito, el juez no sabe si el auto tenía los frenos en mal estado, lo que requiere un perito mecánico. Tampoco hay testigos directos del momento de la colisión, por lo que se necesitan indicios: marcas de frenado, velocidad, comportamiento previo del conductor. Finalmente, el juez visita el lugar del accidente en persona (reconocimiento judicial).

Estos tres instrumentos forman el sistema que permite al juez conocer hechos técnicos o no presenciados y dictar sentencia con convicción fundada.
\end{example}

\section{Finalidad de la prueba en el proceso}

Las pruebas tienen por finalidad \textbf{acreditar} la ocurrencia de un hecho de interés para la litis: un hecho que es \textbf{controvertido}, que debe ser \textbf{conducente} y que fue fijado por el juez en la audiencia preliminar. Quien tiene la carga de acreditar ese hecho debe hacerlo a través de los medios de prueba disponibles.

Se sabe que un hecho existió porque deja un registro, una huella en las cosas o en las personas. Para reproducirlo en el proceso es necesario consultar la fuente y hacerlo a través de los medios correspondientes.

Las pruebas estudiadas hasta este punto (documental, informativa, testimonial y confesional) son medios que tienden a reconstruir hechos para formar la convicción del juez. El juez no puede ver todo ni sabe cómo sucedieron las cosas, por lo que debe fijar los hechos a través de la prueba que las partes ofrezcan y produzcan en la etapa probatoria.

\section{Cuándo es necesaria la prueba pericial}

Cuando para conocer los hechos controvertidos no bastan los conocimientos del común de la gente (ni los del juez, incluso el juez laboral) y hace falta un conocimiento especial en alguna ciencia, arte, industria o actividad técnica especializada, las partes deben valerse de la \textbf{prueba pericial}. Para ello se propone una prueba pericial, a fin de que se realice una peritación.

\subsection{Características de la peritación}

\begin{definition}{Peritación}
Actividad que se desarrolla por encargo judicial, realizada por una persona ajena a las partes y al juez, para hacer conocer al juez determinadas cuestiones técnicas o científicas que escapan a su conocimiento común.
\end{definition}

La peritación no se realiza por iniciativa propia: se ofrece como prueba en los escritos constitutivos del proceso (demanda, contestación, reconvención, contestación de excepciones previas) y es el juez quien la ordena.

\section{Quiénes son los peritos}

Si las profesiones están reglamentadas, son peritos quienes tienen \textbf{título habilitante} para ejercer la profesión. Si la profesión no está reglamentada, se debe buscar a personas que conozcan el tema discutido y tengan amplia experiencia en él, que pueden ser nombradas a falta de profesión reglamentada y de título habilitante.

\subsection{Dónde se ubica a los peritos}

Una vez ofrecida la prueba pericial y admitida por el juez, este designa y sortea un perito dentro de la lista de peritos que llevan las cámaras de cada fuero. Quienes quieren actuar como peritos (por ejemplo ingenieros, arquitectos, contadores, agrimensores, médicos de distintas especialidades) se anotan en las listas de peritos judiciales y, en el momento que corresponda, deben ser sorteados.

\begin{note}
La expresión «desinsaculado» proviene de la práctica antigua en que el sorteo del perito se hacía extrayendo un papel o una bolilla de un saco (\textit{insacular}). Por eso se sigue diciendo que «ha sido desinsaculado el perito».
\end{note}

Generalmente interviene un \textbf{perito único designado de oficio} por el juez, cuando existen razones necesarias para que intervengan personas con conocimientos especiales respecto de determinados hechos.

\begin{definition}{Perito}
Tercero ajeno al proceso, con conocimientos especiales, idóneo y, en su caso, con título habilitante, que esclarece determinadas cuestiones técnicas, científicas, artísticas o industriales.
\end{definition}

El perito es un tercero auxiliar del juez. Tiene el deber de ser \textbf{imparcial} y no puede tener otro interés que esclarecer los hechos, de modo que el juez pueda entenderlos y aplicar la norma que rige. La materia escapa al conocimiento del juez, aunque este deba fijar los hechos.

\section{El consultor técnico (perito de parte)}

La misma complejidad que requiere la intervención del perito para ayudar al juez la enfrenta el abogado que prepara la demanda. Así como el abogado es el asistente técnico en el aspecto jurídico (patrocinio jurídico), la parte puede contar con un asistente en materias científicas o técnicas.

\begin{definition}{Consultor técnico (perito de parte)}
Asistente de la parte en materias científicas o técnicas. Tanto el actor como el demandado tienen derecho a designarlo, y puede designarse también en los escritos constitutivos del proceso.
\end{definition}

\subsection{Diferencias entre perito designado de oficio y consultor técnico}

\begin{table}[H]
\centering
\small
\begin{tabularx}{\textwidth}{
    >{\raggedright\arraybackslash}p{0.19\textwidth}
    >{\raggedright\arraybackslash}X
    >{\raggedright\arraybackslash}X
}
\toprule
\textbf{Criterio} & \textbf{Perito (de oficio)} & \textbf{Consultor técnico} \\
\midrule
Designación & Por el juez, de oficio, por sorteo & Por la parte (actor o demandado) \\
Imparcialidad & Debe ser absolutamente imparcial y objetivo & Puede y debe ser parcial; actúa en favor de su cliente \\
% TODO: contenido incompleto o ambiguo en las notas originales (la tabla dice "con matices"; más adelante se afirma que la aceptación del cargo es facultativa)
Obligatoriedad & Obligado a aceptar el cargo (con matices) & Facultativo; la parte puede removerlo en cualquier momento \\
Dictamen / presentación & Obligado a presentar dictamen en tiempo y forma & No está obligado a presentar nada \\
Recusación & Puede ser recusado con causa (art.\ 17 CPCCN) & No puede ser recusado (es de parte, se asume parcialidad) \\
Honorarios & A cargo de la condena en costas & A cargo de la condena en costas \\
Base del dictamen & Expediente + estudios propios + apoyaturas técnicas & Expediente + asesoría al abogado y a la parte \\
\bottomrule
\end{tabularx}
\end{table}

Una diferencia clave es la imparcialidad. El consultor técnico es un perito de parte: puede ser parcial y, de hecho, debería serlo, porque es contratado por su cliente para que lo asesore y fundamente científico-técnicamente la cuestión de modo que lo favorezca. El perito, en cambio, debe ser absolutamente imparcial y objetivo.

El perito puede basarse en estudios que realiza en instituciones especializadas por fuera del propio expediente judicial. Puede valerse incluso, si la cuestión es compleja, de la ayuda de otros profesionales, pero no puede remitirse en su dictamen únicamente a esas otras cuestiones: el dictamen debe ser hecho por el perito, sin perjuicio de que se valga de otras apoyaturas técnicas o de los exámenes y experimentos necesarios.

\section{Actividad de la contraparte al ofrecer la prueba pericial}

Quien ofrece la prueba pericial debe indicar claramente la \textbf{especialidad} del perito y los \textbf{puntos de pericia}, y puede designar en ese mismo acto al consultor técnico.

Cuando el actor ofrece la prueba pericial, se corre traslado al demandado junto con el traslado de la demanda para que ejerza las facultades que le permite el Código Procesal. Conforme al artículo 468 (mencionado también, en las fuentes originales, como art.~478), la contraparte puede:

% TODO: contenido incompleto o ambiguo en las notas originales (numeración del artículo: 468 o 478)
\begin{enumerate}[label=\alph*)]
    \item Impugnar la procedencia de la prueba pericial por no corresponder conforme al artículo 459, es decir, por no requerirse conocimientos especiales técnicos, científicos, etc.\ para acreditar los hechos que se pretenden a través de los puntos de pericia indicados.
    \item Manifestar desinterés en la pericia, es decir, sostener que se va a abstener de participar en ella.
    \item Impugnar la especialidad del perito propuesto.
    \item Impugnar los puntos de pericia por sostener que están errados, equivocados o que no son procedentes.
    \item Proponer otros puntos de pericia distintos.
    \item Designar su propio consultor técnico.
\end{enumerate}

\begin{formula}{Art. 459 CPCCN --- Procedencia de la prueba pericial}
Procederá la prueba pericial cuando la apreciación de los hechos controvertidos requiere conocimientos especiales en alguna ciencia, arte, industria o actividad técnica especializada.
\end{formula}

\subsection{Costas, traslados y simetría}

% TODO: contenido incompleto o ambiguo en las notas originales (el texto fuente dice "no se hace mérito de la policía"; se interpreta como "pericia")
Si la contraria impugna la procedencia (o manifiesta desinterés) y el juez hace lugar a la pericia, pero en la sentencia esta no resulta un elemento determinante (no se hace mérito de la pericia, no fue necesaria), quien impugnó o manifestó desinterés no queda obligado al pago de las costas de la pericia.

Si la contraparte impugna los puntos de pericia y ofrece otros distintos, se debe correr traslado de vuelta a quien ofreció la prueba. Este traslado va \textbf{por nota}, no por cédula.

Todo lo anterior se aplica en forma simétrica: si el demandado ofrece prueba pericial, se corre traslado al actor para que pueda ejercer la misma actividad. La diferencia es que el traslado al demandado (de la pericial ofrecida por el actor) se corre \textbf{por cédula}, mientras que el traslado al actor (de la pericial ofrecida por el demandado) se corre \textbf{por nota}.

\section{La audiencia preliminar y la designación del perito}

Una vez presentadas todas las impugnaciones y traslados, se llega a la audiencia preliminar, donde el juez debe resolver:

\begin{itemize}
    \item Todo planteo relativo a la especialidad del perito.
    \item Las impugnaciones respecto de los puntos de pericia.
    \item La especialidad del perito, si las partes no se pusieron de acuerdo.
\end{itemize}

En la audiencia preliminar, si las partes se ponen de acuerdo, pueden designar ellas mismas al perito. Sin perjuicio de ello, dado que la calidad del perito se orienta a formar la convicción del juez, este podría cuestionar al perito designado si no es el especialista adecuado o si carece de título habilitante.

El juez tiene también la facultad de agregar otros puntos de pericia o modificar los propuestos, de modo que sean adecuados para acreditar los hechos discutidos.

Si las partes no se ponen de acuerdo, el juez determina la especialidad del perito, fija los puntos de pericia, puede agregar otros si lo considera necesario y puede modificar los existentes. El perito debe ser designado de oficio, sorteado de entre los integrantes de la lista de las cámaras del fuero para la especialidad que resulte.

\section{Recusación del perito}

A partir de la designación del perito en la audiencia preliminar corre un plazo de \textbf{cinco días} para que las partes lo recusen si lo estiman necesario.

Las causales de recusación son las mismas que corresponden a la recusación con causa del juez, contenidas en el artículo 17 del Código Procesal.

\begin{formula}{Art. 17 CPCCN --- Causales de recusación con causa (aplicables al perito)}
Como ejemplos se mencionan: que el perito tenga manifiesto interés en el pleito, ser amigo íntimo o enemigo de alguna de las partes, haber sido denunciado por alguna de las partes, entre otras causales equivalentes a las de recusación judicial.
\end{formula}

Además existen causales específicas de la cuestión pericial:

\begin{itemize}
    \item \textbf{Falta de título habilitante:} el perito carece del título y la profesión está reglamentada.
    \item \textbf{Incompetencia del perito:} tiene un título de profesión en general pero se necesita un especialista en determinada materia.
\end{itemize}

\begin{example}{Incompetencia del perito}
Si se designó un médico clínico y hace falta un médico traumatólogo o un médico psiquiatra; o si se designó un ingeniero cuando en realidad corresponde un agrimensor, como en una pericia sobre inmuebles rurales en la que hay que determinar su superficie.
\end{example}

Si el perito recusado acepta las causales, el juez lo remueve directamente y designa otro perito de oficio. Si el perito rechaza las causales, se forma un \textbf{incidente} que resuelve el juez, y la resolución es irrecurrible en cuanto a si corresponde o no la recusación.

\begin{important}
La notificación de la designación del perito va por nota, por lo que puede sorprender a la parte que deba recusar dentro del plazo por alguna de estas causales. Es necesario estar atento.
\end{important}

\section{Aceptación del cargo, obligaciones y sanciones del perito}

Para intervenir válidamente, el perito no solo debe ser designado por el juez sino que debe presentarse y \textbf{aceptar efectivamente el cargo} ante el Oficial Primero.

La aceptación del cargo es facultativa: no es una carga pública, por lo que el perito puede no aceptarlo. Sin embargo, si reiteradamente y sin justificación un perito no acepta el cargo, puede ser removido de la lista de la cámara que corresponde.

Una vez aceptado el cargo, el perito tiene el deber de:

\begin{itemize}
    \item comparecer todas las veces que sea citado al proceso;
    \item presentar el dictamen en tiempo y forma;
    \item concurrir a las audiencias y dar las explicaciones que correspondan.
\end{itemize}

\begin{formula}{Art. 243 del Código Penal --- Sanción penal al perito}
Establece la sanción para el perito que, legalmente citado, no compareciere o rehusare dar las explicaciones que le sean requeridas: pena de prisión de 15 días a un mes e inhabilitación especial de un mes a un año.
\end{formula}

Además de las sanciones penales, el perito puede ser responsabilizado por los daños y perjuicios que ocasione su incumplimiento, lo que se reclamará por vía incidental.

\section{Anticipo para gastos}

Dentro de los \textbf{tres días} de haber aceptado el cargo, el perito puede solicitar un anticipo para gastos cuando la pericia requiere determinados materiales, insumos o gastos de traslado importantes.

\begin{important}
El anticipo para gastos no debe ser un solapado aumento de honorarios. El perito tiene derecho a recibirlo, pero también el deber de rendir cuenta de los gastos realizados y del destino de los fondos en el dictamen correspondiente, de modo que las partes puedan controlar que los gastos eran tales y fueron destinados a la pericia.
\end{important}

La resolución que fija el anticipo para gastos solamente es susceptible del recurso de \textbf{reposición}; no se puede apelar a la cámara. Una vez notificada y firme, el anticipo debe depositarse dentro de los \textbf{cinco días} de ser notificado por cédula. Si la parte obligada no lo deposita oportunamente, se la tiene por desistida de la prueba pericial y pierde la prueba.

\section{Remoción del perito}

Son causales de remoción del perito:

\begin{itemize}
    \item la renuncia del perito durante el proceso;
    \item que, después de haber aceptado el cargo, rehúse presentarse a dar el dictamen;
    \item que no presente el dictamen oportunamente (luego de haber pedido ampliación del plazo y, en su caso, si no está justificada la demora).
\end{itemize}

En caso de remoción, el perito debe, si corresponde, pagar los daños y es pasible de las sanciones civiles y penales que puedan corresponder.

El plazo para presentar el dictamen, salvo disposición en contrario, es de \textbf{15 días}. Corresponde a la parte que ofreció la prueba pericial realizar las diligencias necesarias para que la prueba no caduque, siendo fundamental notificar al perito de su designación para que se presente en el expediente y acepte el cargo.

\cornelltitle{la prueba pericial, sujetos y procedimiento inicial}
\cornellblocksum{
\cqrow{¿Cuándo procede la prueba pericial?}{Cuando la apreciación de los hechos controvertidos requiere conocimientos especiales en alguna ciencia, arte, industria o actividad técnica especializada que exceden la cultura media del juez (art.\ 459 CPCCN).
}
\cqrow{¿Qué es la peritación?}{Es una actividad desarrollada por encargo judicial, realizada por un tercero ajeno a las partes y al juez, con conocimientos especiales, para hacer conocer al juez cuestiones técnicas o científicas que escapan a su conocimiento común.
}
\cqrow{¿Cómo se designa el perito?}{Por sorteo (desinsaculación) de entre los integrantes de la lista de peritos de la cámara del fuero correspondiente. La especialidad y los puntos de pericia se fijan en la audiencia preliminar.
}
\cqrow{¿Cuál es la diferencia clave entre perito y consultor técnico?}{El perito es designado de oficio, debe ser imparcial y está obligado a presentar el dictamen. El consultor técnico es de parte, puede y debe ser parcial, y su intervención es facultativa; puede ser removido por la parte que lo designó en cualquier momento.
}
\cqrow{¿Cuáles son las causales de recusación del perito?}{Las mismas del art.\ 17 CPCCN para la recusación del juez (interés en el pleito, amistad íntima, enemistad, etc.), más las específicas: falta de título habilitante y falta de la especialidad requerida para los puntos de pericia. El plazo es de cinco días desde la designación.
}
\cqrow{¿Qué pasa si el perito no acepta el cargo reiteradamente?}{Puede ser removido de la lista de la cámara correspondiente, para desalentar la conducta de elegir casos en función de los honorarios esperados.
}
\cqrow{¿Qué consecuencias tiene no pagar el anticipo para gastos?}{Si la parte obligada no deposita el anticipo dentro de los 5 días de notificada por cédula, se la tiene por desistida de la prueba pericial y pierde la prueba (la resolución solo admite recurso de reposición).
}
}{La prueba pericial se ofrece con indicación de especialidad y puntos de pericia, se corre traslado a la contraparte y, en la audiencia preliminar, el juez fija la especialidad y los puntos. El perito, tercero imparcial designado de oficio por sorteo, puede ser recusado en cinco días, debe aceptar el cargo, puede pedir un anticipo para gastos (solo reposición; depósito en cinco días bajo pena de pérdida de la prueba) y puede ser removido por renuncia o incumplimiento. El consultor técnico, en cambio, es de parte y puede ser parcial.
}

\chapter[Práctica de la pericia]{Práctica de la pericia: pasos, dictamen y valoración}\label{cap:practica-pericia}

Establecidos los sujetos de la prueba pericial y su procedimiento inicial, corresponde analizar cómo se realiza la pericia, qué ocurre con el dictamen y de qué modo lo valora el juez.

%==================================================================

\section{Presencia de las partes en la realización de la pericia}

Si las partes quieren presenciar la realización de la prueba pericial, deben hacérselo saber al perito en el expediente. En tal caso, el perito tiene la obligación de informar, con antelación suficiente, el expediente, el lugar, el día y la hora en que llevará a cabo la pericia, de modo que las partes, sobre todo los consultores técnicos, puedan concurrir y hacer las observaciones que consideren pertinentes.

\begin{important}
Si las partes solicitan presenciar la pericia y el perito no informa el día en que realizará su tarea, la pericia es \textbf{nula}, porque se viola el principio de bilateralidad de la audiencia, que tiene raigambre constitucional y no puede ser dejado de lado por ningún sistema que cree el legislador para llevar adelante el proceso.
\end{important}

\section{Los pasos para realizar la prueba pericial}

\subsection{Paso 1: determinación del tema pericial}

El perito debe tener muy en claro qué cuestionario se le pide que conteste y qué cuestiones debe resolver, lo que surge de los puntos de pericia. Puede requerir explicaciones al juez y a las partes respecto de esos puntos e, incluso, puede advertir que exceden su especialidad, en cuyo caso deberá indicar cuál sería el perito de la especialidad correcta.

\subsection{Paso 2: observación del campo de trabajo}

Como experto, el perito debe determinar cuál es su campo de trabajo, qué elementos necesita para llevar a cabo la pericia y dónde se pueden encontrar. Tiene facultad de concurrir a los lugares pertinentes y de recurrir a entidades especializadas para satisfacer las dudas y requerimientos que puedan surgir. Debe recolectar los elementos que luego tendrá que acompañar en cuanto resulten necesarios para el dictamen.

\subsection{Paso 3: elección de los medios}

El perito puede valerse de exámenes de personas, de cosas y de estudios especializados.

\begin{important}
La elección de los medios (la selección de determinados medios importantes para el caso) puede ocasionar después una impugnación o, en su caso, la nulidad de la pericia, si esto pone en riesgo la entidad científica o técnica de la pericia.
\end{important}

\begin{example}{Pericia médica}
El médico legista designado para llevar a cabo la pericia necesita valerse de tomografías, resonancias y estudios complejos (de biología molecular, etc.) que escapan incluso al saber del experto. Se está dando así un paso casi fuera de la prueba pericial e invadiendo la esfera de las pruebas científicas.
\end{example}

El autor Falcón propone establecer la \textbf{prueba científica} como medio probatorio autónomo, porque ciertas cuestiones que encuadrarían dentro de ella requieren estudios especializados, equipos especiales, conjuntos de personas preparadas e instituciones de la materia. Sin embargo, el perito puede integrar su pericia con estos exámenes requeridos a otras instituciones, lo que la convierte en una prueba compleja que excede el marco pericial y traspasa quizás la frontera entre la prueba pericial y la científica.

Esta frontera es de dudosa delimitación, porque el conocimiento avanza paulatinamente: tanto el que tiene el perito respecto de cuestiones que antes no podía alcanzar ni interpretar aun siendo experto, como el de la propia ciencia, que avanza al mismo tiempo en cuestiones más novedosas, complejas y dificultosas que requieren interpretación.

\begin{formula}{Art. 476 CPCCN --- Opinión de instituciones}
Permite que se soliciten opiniones a universidades, instituciones o entidades públicas o privadas respecto de cuestiones técnicas o científicas que resulten complejas. Esto se engloba en lo que puede considerarse prueba científica, que para el autor Falcón debería ser considerada un medio autónomo.
\end{formula}

\subsection{Paso 4: agregación y preservación de los elementos}

El perito no solo debe llevar adelante las operaciones que correspondan sino también preservar los medios de los que se valió, para acompañarlos al expediente.

\begin{example}{Prueba informática}
En una pericia sobre cuestiones electrónicas, el perito debe investigar si en determinada computadora existen ciertos mensajes u otros archivos digitales. Debe preservar esa prueba adecuadamente y rescatarla para asegurar su integridad, trasladarla después y acompañarla como fundamento de sus conclusiones.
% TODO: contenido incompleto o ambiguo en las notas originales (el texto fuente dice "desde el peligro [riesgo] que tiene que ir a investigar")
\end{example}

\subsection{Paso 5: realización de las operaciones técnicas y emisión del dictamen}

El perito debe realizar todas las operaciones técnicas que correspondan para evacuar todos los puntos de pericia y sacar las conclusiones correctas en función de lo que se le pide resolver.

Finalmente debe emitir su \textbf{dictamen}, que deberá ser técnico y científico y valerse de las pruebas necesarias. El dictamen tiene que ser acompañado de las pruebas que sean necesarias.

\section{La prueba documental que acompaña el perito}

Existe un punto delicado: el acompañamiento por el perito de prueba documental que \textquotedblleft descubre\textquotedblright{} puede importar, en muchos casos, una agregación extemporánea de esa prueba al proceso.

La prueba documental debe acompañarse junto con la demanda, la contestación, la reconvención y la contestación de excepciones, salvo que se trate de documentos nuevos o recién conocidos. Los documentos posteriores a esas oportunidades o recién conocidos pueden acompañarse en el proceso hasta el llamamiento de autos para sentencia en primera instancia, sin perjuicio de hacerlo en segunda instancia.

% TODO: contenido incompleto o ambiguo en las notas originales (el pasaje sobre la agregación en segunda instancia -"debe apelarse en la sentencia definitiva y fundarse en forma libre"- no resulta claro)

\begin{important}
El perito no puede acompañar prueba documental supliendo la negligencia de la parte que la tenía a su disposición. Esto rige cuando el perito, al realizar la pericia, concurre a bancos, entidades públicas, etc., y solicita documentación que le será necesaria y que luego acompañará como fundamento de su dictamen al presentarlo en el expediente.
\end{important}

\section{Notificación del dictamen y actividad de las partes}

El perito debe presentar el dictamen \textbf{por escrito, con copias para las partes}. Las partes son notificadas por cédula del dictamen pericial y, a partir de ahí, tienen un plazo de \textbf{cinco días} para realizar distintas actividades.

Dentro de los cinco días de notificadas por cédula, las partes pueden:

\begin{enumerate}[label=\alph*)]
    \item Impugnar la pericia por contener falsedades.
    \item Plantear la nulidad de la prueba pericial (no expresamente contemplada en el código, pero disponible por vía incidental).
    \item Pedir una ampliación de la pericia, dado que el perito no contestó todos los puntos solicitados.
    \item Pedir una nueva pericia por ser inconsistente la anterior y no poder ser corregida por ampliación.
    \item Pedir explicaciones del perito.
\end{enumerate}

La nulidad de la pericia puede darse, por ejemplo, cuando se ha ejercido violencia sobre el perito, cuando el perito perdió sus facultades mentales de modo que la pericia no puede considerarse válida, o cuando tiene groseros errores de fundamentación que ya no admiten corrección a través de un pedido de explicaciones. Si prospera la nulidad, se debe realizar una nueva pericia por otros expertos, no por el mismo perito que ya participó.

\begin{important}
Si no se impugna la pericia por inconsistencias o falsedades dentro de los cinco días de notificado el dictamen por cédula, la impugnación puede hacerse hasta el momento de los alegatos. Si tampoco se hizo en esa oportunidad y el juez, que es quien valora la prueba pericial, la utiliza como elemento de convicción para dictar sentencia, es posible cuestionarla ante la cámara, al expresar agravios en el recurso de apelación contra la sentencia definitiva, por las inconsistencias o falsedades de que adolezca, aunque no se haya hecho al correrse el traslado de la pericia.
\end{important}

\section{El artículo 475 CPCCN: medidas complementarias}

El artículo 475 del Código Procesal permite que las partes ofrezcan, o que el juez ordene de oficio, la ejecución de las siguientes medidas:

\begin{formula}{Art. 475 CPCCN --- Medidas complementarias de la prueba pericial}
\begin{itemize}
    \item \textbf{Inc.\ 1:} planos, relevamientos, tomas fotográficas y filmaciones cinematográficas de objetos, documentos o lugares, con el empleo de medios o instrumentos técnicos.
    \item \textbf{Inc.\ 2:} exámenes científicos necesarios para el mejor esclarecimiento de los hechos controvertidos.
    \item \textbf{Inc.\ 3:} reconstrucción de hechos para comprobar si se han producido o pudieron producirse de una determinada manera.
\end{itemize}
\end{formula}

Las medidas del inciso 1 (planos, fotografías, películas) son en realidad \textbf{prueba documental}, dado que la concepción de prueba documental consiste en la representación del pensamiento humano en algún tipo de objeto. El reconocimiento de hechos tampoco es prueba pericial en sí misma: es otro medio de prueba. Estas facultades del artículo 475 se vinculan con el reconocimiento judicial del artículo 479.

\section{Valoración del dictamen pericial (art.\ 477 CPCCN)}

El dictamen pericial \textbf{no es vinculante} para el juez salvo que la ley disponga lo contrario.

\begin{formula}{Art. 477 CPCCN --- Valoración del dictamen pericial}
El juez valorará el dictamen pericial de acuerdo con las reglas de la sana crítica. Podrá apartarse de las conclusiones del perito cuando existan otras pruebas o razones que justifiquen tal apartamiento.
\end{formula}

El código procesal exige al juez que valore el dictamen en función de las reglas de la sana crítica: la prueba pericial no escapa a la valoración que, para todos los medios de prueba, debe hacerse conforme a esas reglas.

\begin{example}{Juicio de demencia}
La ley dispone lo contrario cuando, en el juicio de demencia, establece que el presunto insano debe ser examinado por tres médicos psiquiatras y que, si lo encuentran en sano juicio, el juez no podría declararlo demente. Esta pericia, por la ley y por su resultado, es absolutamente vinculante para el juez.

No sucede lo mismo a la inversa: aun cuando los médicos psiquiatras consideren demente a la persona, ello no vincula al juez. Fundado en otras razones que surjan del expediente y en lo que valore por la regla de la sana crítica, puede apartarse del dictamen y no declarar demente a esa persona.
\end{example}

Para valorar el dictamen, el juez debe ponderar:

\begin{itemize}
    \item cuán completo sea el dictamen pericial;
    \item cuán competentes e idóneos sean los expertos que lo elaboraron;
    \item cuán coherente y lógico sea el dictamen en sí mismo;
    \item cuán fundado esté en cuestiones teóricas, científicas y en estudios;
    \item cuán completo sea en cuanto a que se acompañaron todos los elementos utilizados al elaborar la pericia.
\end{itemize}

\section{El perito árbitro: cuando el dictamen sí es vinculante}

En el proceso arbitral (título tercero del código procesal), lo que deciden los peritos árbitros en el ámbito de su competencia respecto de cuestiones de hecho es \textbf{vinculante} para el juez.

El artículo 516 del código procesal permite que, en la etapa de ejecución de sentencia, cuando existen cuestiones complicadas o de difícil dilucidación (por ejemplo, liquidaciones de sociedades o de sociedad conyugal), se tome como puerta de salida del proceso, como método alternativo de resolución de conflictos, dejar esas cuestiones a lo que resuelvan estos peritos árbitros, que utilizan la misma regla que los amigables componedores.

El perito árbitro vincula al juez en cuanto a la fijación del hecho a través del laudo emitido en el ámbito de la cuestión que se le dejó para resolver, pero no lo limita en cuanto a la aplicación del derecho en ese proceso.

\cornelltitle{práctica de la pericia, dictamen y valoración}
\cornellblocksum{
\cqrow{¿Qué es el principio de bilateralidad en la pericia?}{Si las partes solicitaron presenciar la pericia, el perito tiene la obligación de informarles con antelación suficiente el lugar, día y hora. Si no lo hace, la pericia es nula, porque se viola el principio de bilateralidad de audiencia, de raigambre constitucional.
}
\cqrow{¿Cuáles son los pasos del perito para realizar la pericia?}{1) Determinación del tema pericial (interpretar los puntos de pericia). 2) Observación del campo de trabajo y recolección de elementos. 3) Elección de los medios (con cautela, pues puede dar lugar a impugnación o nulidad). 4) Agregación y preservación de los elementos. 5) Realización de las operaciones técnicas y emisión del dictamen.
}
\cqrow{¿Qué puede hacer la parte en los 5 días posteriores al dictamen?}{Impugnar por falsedades o inconsistencias, plantear la nulidad, pedir ampliación, pedir nueva pericia o pedir explicaciones al perito. Vencido ese plazo precluyen algunas facultades (la impugnación puede hacerse hasta los alegatos).
}
\cqrow{¿Qué medidas complementarias prevé el art.\ 475 CPCCN?}{Planos, relevamientos, fotografías y filmaciones (inc.\ 1, que constituyen en realidad prueba documental); exámenes científicos (inc.\ 2); y reconstrucción de hechos (inc.\ 3).
}
\cqrow{¿Es vinculante el dictamen pericial para el juez?}{No, salvo que la ley disponga lo contrario (por ejemplo, en el juicio de demencia, si los tres médicos psiquiatras dictaminan que la persona está en sano juicio, el juez no puede declararla demente). El juez valora el dictamen según las reglas de la sana crítica y puede apartarse fundadamente.
}
\cqrow{¿Qué pondera el juez al valorar el dictamen?}{Su completitud, la competencia e idoneidad de los expertos, su coherencia lógica, su fundamento teórico, científico y en estudios, y que se hayan acompañado todos los elementos utilizados.
}
\cqrow{¿Cuándo sí es vinculante el dictamen?}{Cuando interviene un perito árbitro (art.\ 516 CPCCN, en etapa de ejecución de sentencia): su fijación de hechos es vinculante para el juez, aunque este conserva la aplicación del derecho.
}
}{La pericia se realiza en cinco pasos: determinación del tema, observación del campo de trabajo, elección de los medios, preservación de los elementos y emisión del dictamen, con aviso previo a las partes que pidieron presenciarla (bajo pena de nulidad por violación de la bilateralidad). Presentado el dictamen y notificado por cédula, las partes tienen cinco días para impugnar, plantear nulidad, pedir ampliación, nueva pericia o explicaciones. El juez valora el dictamen según la sana crítica (art.\ 477) y puede apartarse fundadamente; solo excepcionalmente la ley lo torna vinculante (juicio de demencia), y el perito árbitro sí vincula al juez en cuestiones de hecho.
}

\chapter[Prueba indiciaria y reconocimiento judicial]{Prueba indiciaria y reconocimiento judicial}\label{cap:indiciaria}

%==================================================================

\section{La prueba indiciaria: concepto y necesidad}

La prueba pericial es un medio complejo no solo por la materia que requiere la intervención del perito, sino por la práctica de la pericia, que tiene distintos momentos y puede requerir la participación de terceras personas o de determinadas instituciones.

Se plantean entonces las siguientes cuestiones: ¿en qué condiciones puede el juez apartarse de lo dispuesto por estos expertos? ¿Qué eficacia probatoria tienen estos dictámenes y los razonamientos que hacen los expertos para ilustrar el conocimiento?

La prueba indiciaria es necesaria precisamente cuando falta la demostración acabada, a través de los otros medios probatorios (documental, informativa, pericial, reconocimiento judicial), de la certeza respecto de la existencia del hecho de importancia para la litis. Si los hechos están acreditados categóricamente por los otros medios establecidos en el código, no es necesario recurrir a estas pruebas llamadas indiciarias.

\section{Presunciones: concepto y tipos}

Al hablar de presunciones existe confusión con los términos que utilizan la ley o la jurisprudencia, que hablan de indicios o los asocian a la conclusión que se extrae de ellos.

Sobre la naturaleza de la prueba indiciaria hay dos posturas. Para algunos autores no es un medio de prueba sino que se limita a ser un razonamiento, una forma de razonamiento jurídico. Otros sostienen que, aunque considerada de forma indirecta y a partir de la utilización de principios lógicos, atiende a la demostración de un hecho controvertido de interés para el proceso y, de esta forma, es un medio de prueba.

\begin{table}[H]
\centering
\small
\begin{tabularx}{\textwidth}{
    >{\raggedright\arraybackslash}X
    >{\raggedright\arraybackslash}X
    >{\raggedright\arraybackslash}X
}
\toprule
\textbf{Tipo de presunción} & \textbf{Quién la establece} & \textbf{Admite prueba en contrario} \\
\midrule
Legal -- \textit{iuris et de iure} (de pleno derecho) & El legislador (la ley) & No: es absoluta \\
Legal -- \textit{iuris tantum} (relativa) & El legislador (la ley) & Sí: puede ser desvirtuada \\
Judicial -- \textit{hominis} & El juez & Depende de la valoración del juez \\
\bottomrule
\end{tabularx}
\end{table}

% TODO: contenido incompleto o ambiguo en las notas originales (el texto fuente dice "como dice la solar, 'tiene por cierto'")
En la presunción legal, el legislador le asigna peso a la prueba: acreditados determinados extremos (los hechos que tienen que actuar), el juez \textit{tiene por cierto} el hecho y extrae su consecuencia. En las presunciones \textit{hominis}, en cambio, es el juez quien asigna peso a los indicios, a diferencia de la presunción legal, en la que el peso lo asigna el legislador.

\section[Presunciones judiciales (hominis)]{Las presunciones judiciales (hominis): requisitos del art.\ 163 inc.\ 5}

Las presunciones judiciales (\textit{hominis}) no están establecidas en el código procesal como un medio particular de prueba en el capítulo correspondiente, sino en el artículo 163, referido a la sentencia, particularmente en su inciso 5.

\begin{formula}{Art. 163 inc.\ 5 CPCCN --- Presunciones judiciales}
Las presunciones no establecidas por ley constituirán prueba cuando se funden en hechos reales y probados y cuando, por su número, precisión, gravedad y concordancia, produzcan convicción según las reglas de la sana crítica, de acuerdo con la naturaleza del juicio.
\end{formula}

Los requisitos de los indicios son los siguientes:

\begin{itemize}
    \item \textbf{Hechos reales y probados:} los indicios deben estar acreditados en el expediente. El indicio no puede ser una conjetura: tiene que ser un hecho concreto, real y probado, que por sí solo no basta para tener por acreditado el otro hecho presumido, pero que sumado a otros indicios sí puede lograrlo.
    \item \textbf{Numerosos:} debe existir pluralidad de indicios.
    \item \textbf{Precisos:} deben tener una interpretación en un solo sentido, sin admitir otra.
    \item \textbf{Graves:} tienen peso por sí mismos para determinar, con cierto grado de realidad, la dirección hacia determinado hecho.
    \item \textbf{Concordantes:} son convergentes, en el sentido de que todos apuntan hacia el mismo lado y forman una cadena sin lagunas y sin lugar a dudas respecto de lo que se puede concluir a partir del conjunto.
\end{itemize}

\section{La conducta procesal de las partes como indicio}

El artículo 163 dispone también que la conducta de las partes puede asumir la entidad de indicio. Por ejemplo: una conducta omisiva, la conducta de la parte que destruye documentación, la parte que no contesta la demanda o la que no actúa con buena fe a lo largo del proceso. Esa conducta, sumada a otros elementos de prueba, puede ser un elemento de convicción que complemente las conclusiones del proceso.

\begin{important}
El juez no podría valerse únicamente de la conducta de las partes para condenar a una de ellas o hacer lugar a la pretensión de la contraria. Pero, reunidos determinados elementos probatorios y verificada una conducta desleal, omisiva o inequitativa de la contraria, esa conducta puede ser un indicio del cual el juez se valga para fallar en contra de quien la asume y perjudica a la contraparte.
\end{important}

Asimismo, al dictar sentencia el juez puede considerar la \textbf{conducta desplegada por las partes durante el proceso}. Si una parte está en condiciones de aportar determinado medio de prueba y omite hacerlo sin explicar por qué, lo más razonable es pensar que ese medio de prueba le era desfavorable. No puede premiarse una conducta que obstaculiza la finalidad del proceso, que es averiguar cómo sucedieron los hechos para determinar si corresponde condenar o absolver al demandado.

Las reglas de la sana crítica, a las que remite el artículo 163 inciso 5 del CPCCN para valorar las presunciones, se desarrollan en el capítulo~\ref{cap:valoracion}.

\section{El reconocimiento judicial (art.\ 479 CPCCN)}

\begin{definition}{Reconocimiento judicial}
Prueba más directa y más fiable, en la que el juez se apersona y percibe directamente los hechos, sin que intermedie para su representación ninguna cosa ni persona. No interviene un testigo que declare lo que percibió por sus sentidos ni un dictamen pericial que saque conclusiones: es el propio juez quien toma conocimiento directo en el lugar de los hechos.
\end{definition}

\begin{formula}{Art. 479 CPCCN --- Reconocimiento judicial}
El juez podrá ordenar, de oficio o a petición de parte, el reconocimiento judicial de lugares o de cosas. Se interpreta que también comprende personas, como en los procesos de restricción de la capacidad, donde el juez tiene entrevista con el presunto incapaz. Si el conocimiento excede el propio del juez, puede valerse de un perito para que interprete y señale las distintas circunstancias en juego.
\end{formula}

Es una prueba \textbf{privativa} del juez. Esto no quiere decir que no pueda ser ofrecida por una parte (puede ser ofrecida por una parte o bien ordenada de oficio), sino que el juez es quien puede decidir la utilidad que le significa esa prueba, en función de lo que podrá recabar y del conocimiento que necesita para tener certeza de cómo sucedieron los hechos de interés para el proceso.

En el acto del reconocimiento judicial, el juez puede hacer comparecer a quienes considere pertinentes, puede pedir la concurrencia de testigos y puede valerse también de las medidas previstas en el artículo 475 (planos, croquis, fotografías, filmaciones, reconstrucción de hechos) para aprehender adecuadamente los hechos que deberá fijar después en el proceso judicial.

\section{Forma de la diligencia del reconocimiento judicial}

El juez debe comparecer personalmente, con las personas del tribunal que considere pertinentes. Se debe levantar un \textbf{acta} dejando constancia de las distintas circunstancias percibidas.

\begin{important}
La narración que haga el juez de sus percepciones en el acta debe ser lo más objetiva posible, porque la etapa probatoria no es el momento de valorar la prueba y podría incurrir en prejuzgamiento. Debe dejar constancia objetiva, además, de las aclaraciones y observaciones de las partes, de los consultores técnicos, de los peritos, de las declaraciones eventuales de los testigos y de todas las dificultades que se presenten.
\end{important}
% TODO: contenido incompleto o ambiguo en las notas originales (el texto fuente alude a "autores técnicos [consultores técnicos]" y a "añadir presentadas por el juez" en la enumeración del contenido del acta)

\cornelltitle{prueba indiciaria y reconocimiento judicial}
\cornellblocksum{
\cqrow{¿Qué son las presunciones hominis?}{Son presunciones judiciales: el juez infiere un hecho no acreditado directamente a partir de indicios que son hechos reales, probados en el expediente, numerosos, precisos, graves y concordantes (art.\ 163 inc.\ 5 CPCCN).
}
\cqrow{¿Cuál es la diferencia entre presunción \textit{iuris et de iure} e \textit{iuris tantum}?}{La \textit{iuris et de iure} (de pleno derecho) no admite prueba en contrario: la conclusión es inexorable. La \textit{iuris tantum} (relativa) admite prueba en contrario: la presunción cede si se prueba lo contrario.
}
\cqrow{¿Qué son las reglas de la sana crítica?}{Un razonamiento libre de error y de vicio, sistematizado por Falcón en ocho reglas, con el que el juez determina cuál de las posiciones de las partes es la correcta y puede demostrar, de manera controlable por las partes, cómo formó su convicción (art.\ 386 CPCCN).
}
\cqrow{¿Qué es el reconocimiento judicial?}{La prueba más directa y fiable: el juez percibe personalmente y sin intermediarios los hechos en el lugar donde ocurrieron (lugares, cosas y personas). Se labra un acta objetiva. El juez puede valerse de peritos para interpretar lo observado y de las facultades del art.\ 475 CPCCN.
}
\cqrow{¿Por qué el acta del reconocimiento judicial debe ser objetiva?}{Porque en la etapa probatoria el juez no debe valorar todavía la prueba, ya que incurriría en prejuzgamiento. El acta debe dejar constancia objetiva de lo percibido y de las aclaraciones de las partes, consultores técnicos, peritos y testigos.
}
}{Cuando no hay prueba directa, el juez puede inferir un hecho controvertido a partir de indicios reales, probados, numerosos, precisos, graves y concordantes (presunciones \textit{hominis}, art.\ 163 inc.\ 5), a diferencia de las presunciones legales (\textit{iuris et de iure}, que no admiten prueba en contrario, e \textit{iuris tantum}, que sí la admiten). La convicción se forma según las reglas de la sana crítica (art.\ 386). El reconocimiento judicial (art.\ 479) es la prueba más directa y fiable: el juez percibe lugares, cosas y personas sin intermediarios, puede valerse de peritos y de las medidas del art.\ 475. Todo el andamiaje probatorio sirve a un único fin: formar la convicción del juez para que pueda dictar una sentencia fundada y justa.
}

\chapter{Los alegatos}\label{cap:alegatos}

Concluida la producción de la prueba, la etapa final del proceso civil ordinario atraviesa tres momentos: los alegatos, la etapa conclusional (``autos para sentencia'') y la valoración de la prueba por el juez al dictar sentencia. Este capítulo se ocupa del primero; los dos siguientes, de los restantes.

Los \textbf{alegatos} son el acto procesal escrito (o, en el futuro, oral) mediante el cual las partes analizan la prueba producida y procuran formar la convicción del juez antes del dictado de la sentencia. La \textbf{etapa conclusional} es el momento en que el juez cierra el debate, con sus distintas vías de acceso, sus efectos y la regulación de la morosidad judicial. Los \textbf{sistemas de valoración de la prueba} explican cómo aprecia el juez la prueba al dictar sentencia, con especial atención al sistema de la sana crítica (art.~386 CPCCN) y sus reglas según Falcón, ilustrado con ejemplos de sentencias reales.

\begin{note}
\textbf{Analogía para comprender el recorrido.} Puede imaginarse un partido de fútbol: la etapa probatoria es el partido en sí (los jugadores ---las partes--- muestran sus jugadas ---las pruebas---). Los alegatos son el análisis posterior al partido: cada equipo explica al árbitro ---el juez--- por qué ganó. Luego el árbitro se retira (autos para sentencia) y delibera utilizando las reglas del reglamento ---la sana crítica--- para dictar el resultado final: la sentencia.
\end{note}

\section{Concepto de alegato}

\begin{definition}{Alegato}
Acto procesal en virtud del cual se concluye con la etapa probatoria y mediante el cual las partes se expresan respecto del mérito de la prueba que se produjo en el expediente.
\end{definition}

En este acto procesal, cada parte analiza las pruebas producidas para arribar a una conclusión y trata de formar la convicción del juez:

\begin{itemize}
    \item Si es el \textbf{actor}: para demostrar por qué es viable hacer lugar a su pretensión.
    \item Si es el \textbf{demandado}: para demostrar por qué debe rechazarse la demanda y la pretensión invocada por el actor.
\end{itemize}

Esto muestra que todas las etapas del proceso están vinculadas entre sí, más allá de que rija un sistema de \textbf{preclusión}, en virtud del cual se avanza por etapas y no pueden realizarse los actos procesales que no se ejercieron en la etapa para la cual estaban destinados.

\section{Naturaleza jurídica: acto escrito y facultativo}

El alegato presenta dos características fundamentales:

\begin{enumerate}[label=\Alph*)]
    \item \textbf{Es escrito.} El Código Procesal Civil y Comercial de la Nación (CPCCN) es fundamentalmente un código escrito. La mayoría de los actos son escritos, con algunas implementaciones de oralidad en las audiencias (preliminar, confesional, testimonial) y eventualmente en audiencias conciliatorias (art.~36 CPCCN). Por eso, el alegato es un acto procesal escrito.
    \item \textbf{Es facultativo.} No es una carga ni un deber, y no trae aparejada sanción ni consecuencia alguna para la parte que no lo presente.
\end{enumerate}

Debe recordarse la distinción entre ambas categorías: el \textbf{deber} es aquel que trae aparejada una sanción, mientras que la \textbf{carga} ---imperativo del propio interés--- es aquella cuya inobservancia, si bien no trae aparejada una sanción, puede traer aparejada una consecuencia (por ejemplo, quien no contesta la demanda puede ser declarado rebelde). En el caso de los alegatos no existe consecuencia alguna.

\section{Utilidad del alegato}

Aunque su presentación sea facultativa, el alegato tiene una doble utilidad.

\subsection{Formar la convicción del juez}

La finalidad inmediata es demostrarle al juez que lo invocado en los escritos introductorios quedó acreditado de determinado modo.

\subsection{Base para la expresión de agravios en apelación}

Cuando se interpone un recurso de apelación contra la sentencia de primera instancia, debe fundarse mediante un escrito llamado \textbf{expresión de agravios}, que consiste en una crítica concreta y razonada de la sentencia impugnada.

Si la parte ya realizó un alegato en el que efectuó su propia valoración de la prueba a la luz de los hechos y arribó a su conclusión, ese alegato puede servirle de base para atacar la sentencia en la que considera que se incurrió en un error, ya sea porque se fijaron mal los hechos o porque se valoró mal la prueba.

\section{Procesos en los que el alegato está excluido}

El legislador excluyó la facultad de alegar en los siguientes procesos, por razones de lógica vinculadas a su extensión o características:

\begin{table}[H]
\centering
\begin{tabularx}{\textwidth}{>{\raggedright\arraybackslash}p{4.2cm} >{\raggedright\arraybackslash}X}
\toprule
\textbf{Proceso} & \textbf{Razón de la exclusión} \\
\midrule
Proceso sumarísimo & Es mucho menos extenso que el ordinario; no tiene sentido alegar. \\
Juicio ejecutivo & No hay valoración de prueba general; la única prueba se vincula a las excepciones. \\
Proceso sucesorio & Es un proceso voluntario; no corresponde el alegato. \\
Proceso voluntario en general & Por sus particularidades, el legislador lo excluyó expresamente. \\
\bottomrule
\end{tabularx}
\end{table}

\section{Regulación legal: artículo 482 CPCCN}

El alegato se encuentra regulado en el artículo 482 del CPCCN, denominado ``Agregación de pruebas -- Alegatos''. Producida la prueba, el expediente se pone a disposición de las partes para alegar. El artículo regula la agregación de los cuadernos de prueba al expediente principal y establece los plazos para retirar el expediente en préstamo y para presentar los alegatos.

Actualmente el sistema de cuadernos de prueba prácticamente no existe debido a la digitalización del expediente: toda la prueba se produce dentro del expediente digital.

\subsection{Procedimiento para llegar a la etapa de alegatos}

Una vez producida toda la prueba (o decretada la negligencia o caducidad de alguna), cualquiera de las partes puede peticionar la clausura de la etapa probatoria y que se pongan los autos para alegar. El juez dicta entonces una resolución cuyo texto típico es:

\begin{quote}
``Buenos Aires, 10 de junio de 2020. Siendo exacto lo expuesto y no existiendo medidas pendientes de producción, decreto clausurada la etapa probatoria y pongo los autos en Secretaría para alegar en los términos del artículo 482 del Código Procesal. El orden establecido para retirar en préstamo el expediente será el siguiente: primero la actora, luego el demandado, luego la citada en garantía. Notifíquese.''
\end{quote}

\subsection{Plazos: individual y común}

El artículo 482 establece dos plazos distintos:

\begin{table}[H]
\centering
\begin{tabularx}{\textwidth}{>{\raggedright\arraybackslash}p{2.3cm} >{\raggedright\arraybackslash}p{2.6cm} >{\raggedright\arraybackslash}p{2.6cm} >{\raggedright\arraybackslash}X}
\toprule
\textbf{Tipo} & \textbf{Finalidad} & \textbf{Duración} & \textbf{Particularidades} \\
\midrule
Plazo individual & Retirar en préstamo el expediente & 6 días por parte & Comienza a correr desde que la resolución queda firme y notificada. El orden lo fija el juez: actor $\rightarrow$ demandado $\rightarrow$ citada en garantía. \\
\addlinespace
Plazo común & Presentar el escrito del alegato & 6 días $\times$ n.º de partes & Vence el mismo día para todos. Se computa multiplicando 6 días por el número de partes (si actúan con la misma representación, cuentan como una). \\
\bottomrule
\end{tabularx}
\end{table}

\begin{formula}{Plazo común para alegar}
\[
\text{Plazo común} = 6\ \text{días} \times \text{n.º de partes}
\]
\end{formula}

\begin{example}{Cómputo del plazo común}
Actor y demandado: 2 partes $\times$ 6 días $=$ 12 días de plazo común para alegar. Si hay 3 partes: $3 \times 6 = 18$ días. Si el demandado y la citada en garantía tienen el mismo abogado, cuentan como una sola parte.
\end{example}

\subsection{Consecuencia por no devolver el expediente en término}

Si la parte que retiró el expediente en préstamo no lo devuelve dentro de los 6 días, sin necesidad de intimación previa, se tiene por no presentado el alegato.

\begin{important}
Para la devolución del expediente no rige el plazo de gracia de las dos primeras horas, porque no se trata de presentar un escrito sino de devolver el expediente. Por ello es fundamental que el expediente sea devuelto dentro de los seis días, para no quedar sujeto a la interpretación del tribunal.
\end{important}

\subsection{Cómputo del plazo: firmeza de la resolución}

Los plazos comienzan a correr luego de que la resolución de clausura quede firme (notificada y no impugnada). El plazo varía según quién la haya firmado:

\begin{itemize}
    \item Si la firma el \textbf{juez}: la resolución queda firme a los \textbf{5 días} de notificada.
    \item Si la firma el \textbf{secretario o prosecretario}: el mecanismo de impugnación es el del art.~38 ter CPCCN, y el plazo es de \textbf{3 días}.
\end{itemize}

Recién una vez firme la resolución corren simultáneamente el plazo individual para retirar el expediente y el plazo común para alegar.

\section{Contenido estructural del alegato}

No existe en la ley un artículo que regule expresamente el contenido del alegato (a diferencia de la demanda ---art.~330--- o la contestación ---art.~356---). No obstante, se recomienda el libro \textit{Cómo se hace un alegato}, de Rojas y Falcón (Ed. Abeledo Perrot), que enseña de manera teórico-práctica la estructura correcta.

La estructura recomendada es:

\begin{enumerate}
    \item \textbf{Suma o sumario} (``Presenta alegato''), encabezado y datos del expediente.
    \item \textbf{Descripción y relato de los hechos:} cuáles invocó el actor en la demanda, cuál versión dio el demandado al contestar y qué hechos quedaron controvertidos (por ejemplo, cuál de los rodados violó la luz del semáforo).
    \item \textbf{Análisis estático de las pruebas:} se analiza una por una la prueba producida, marcando lo favorable y atacando la prueba de la contraria (por ejemplo, cuestionando la idoneidad de testigos que están dentro de las generales de la ley por ser amigos o empleados de la otra parte).
    \item \textbf{Análisis dinámico:} síntesis de qué hechos quedaron acreditados con qué prueba, para formar convicción en el juez.
    \item \textbf{Conclusión y firma.}
\end{enumerate}

En cuanto al contenido de derecho, existe una corriente \textbf{restrictiva} (solo hechos y prueba) y una corriente \textbf{amplia} (puede incluir argumentaciones jurídicas), dado el principio \textit{iura novit curia}, en virtud del cual el juez conoce el derecho y puede apartarse del invocado por las partes.

\section{El alegato y la implementación de la oralidad}

Existe un anteproyecto de nuevo Código Procesal en el Congreso que regula el proceso ordinario como un \textbf{proceso por audiencias}, con mayor oralidad:

\begin{enumerate}
    \item \textbf{Primera etapa escrita:} demanda y contestación.
    \item \textbf{Audiencia preliminar:} el juez ordena la producción de la prueba escrita (pericial, informes) y fija la audiencia de vista de causa.
    \item \textbf{Audiencia de vista de causa:} se produce toda la prueba no escrita (testigos, partes, peritos que explican su dictamen). Una vez terminada la prueba oral, cada parte tiene \textbf{20 minutos} para alegar oralmente.
\end{enumerate}

El alegato oral requiere otra preparación de los operadores jurídicos: técnicas de litigación oral y un estudio mucho más detenido del expediente antes de la audiencia, ya que habrá que tener analizado todo lo ocurrido hasta ese momento y, además, hablar en el alegato de la prueba producida en esa misma audiencia.

La finalidad del alegato es la misma en ambos sistemas: comunicarse con el juez, demostrar cómo la prueba acredita los hechos y persuadirlo para que, al dictar sentencia, haga viable la pretensión. Lo que cambia es la forma (de escrita pasa a oral), lo que permite una mayor inmediación con la jurisdicción.

\cornelltitle{los alegatos}
\cornellblock{
\cqrow{¿Qué es el alegato y cuál es su función principal?}{Es el acto procesal (escrito u oral) mediante el cual las partes analizan la prueba producida en la etapa probatoria y procuran formar la convicción del juez. Tiene dos utilidades centrales: (1) demostrar al juez que los hechos invocados quedaron acreditados, para obtener una sentencia favorable; y (2) sentar las bases para la expresión de agravios en el recurso de apelación, permitiendo atacar los errores en los que pueda incurrir el juez al valorar la prueba.
}
\cqrow{¿Es obligatorio presentar el alegato?}{No. Es un acto facultativo: no constituye ni una carga ni un deber. Quien no lo presenta no sufre consecuencia alguna (no hay rebeldía ni sanción). Su omisión solo perjudica en la medida en que el juez carezca de los argumentos de esa parte al momento de valorar la prueba.
}
\cqrow{¿En qué procesos no se puede alegar?}{El alegato está excluido en: (a) proceso sumarísimo (por su menor extensión); (b) juicio ejecutivo (no hay prueba general que valorar); (c) proceso sucesorio; (d) procesos voluntarios. El legislador los excluye por razones lógicas vinculadas a las particularidades de cada uno.
}
\cqrow{Art.~482 CPCCN: ¿cuáles son los dos plazos y cómo se computan?}{
}
\cqrow{¿Cuándo comienzan a correr los plazos del art.~482?}{Luego de que la resolución de clausura quede firme: si la firmó el juez, a los 5 días de la notificación; si la firmó el secretario o prosecretario, a los 3 días (art.~38 ter CPCCN). Recién una vez firme la resolución, corren simultáneamente el plazo individual para retirar el expediente y el plazo común para alegar.
}
\cqrow{¿Cuál es la estructura recomendada de un alegato?}{1.º Suma o encabezado. 2.º Relato de los hechos invocados y de cuáles quedaron controvertidos. 3.º Análisis estático de cada prueba producida (propia y de la contraria), con posibilidad de atacar la idoneidad de testigos. 4.º Análisis dinámico: qué hechos quedaron acreditados con qué prueba. 5.º Conclusión y firma.
}
}
\medskip
\cornellblocksum{
\cqrow{¿Qué implica la oralidad en los alegatos según el anteproyecto?}{En el proceso por audiencias propuesto, luego de la audiencia de vista de causa (donde se produce toda la prueba oral), cada parte tiene 20 minutos para alegar verbalmente. Esto requiere otra preparación: dominio de técnicas de litigación oral y análisis detenido del expediente antes de la audiencia, ya que deberá mencionarse también la prueba producida ese mismo día.
}
}{el alegato ---acto procesal escrito y facultativo regulado en el art.~482 CPCCN--- cierra la etapa probatoria y le permite a cada parte presentar al juez su análisis de la prueba producida, con la doble utilidad de persuadirlo en el caso concreto y de sentar las bases para una eventual expresión de agravios si la sentencia resulta desfavorable. Sus plazos son individuales para retirar el expediente en préstamo (6 días por parte, en el orden fijado por el juez) y comunes para presentar el escrito (6 días por el número de partes), y comienzan a correr luego de que quede firme la resolución de clausura de la etapa probatoria.
}

\chapter[Etapa conclusional]{La etapa conclusional: autos para sentencia}\label{cap:conclusional}

%==============================================================

\section{Concepto}

Así como el proceso se inicia con la promoción de la demanda, en algún momento debe finalizar. En líneas generales, finaliza con el dictado de la sentencia en la \textbf{etapa conclusional} del proceso.

Este es el momento de máxima expresión de la jurisdicción: el juez crea la \textbf{norma individual} para el caso, que no es otra cosa que la sentencia, en función de los hechos que tuvo por ciertos y de la aplicación de las normas abstractas al caso concreto.

No obstante, puede suceder que no se llegue a esta etapa si el proceso termina antes por \textbf{modos anormales}: acuerdo conciliatorio, transacción, desistimiento del proceso por el actor, allanamiento del demandado o caducidad de instancia (art.~310 CPCCN). Estos modos anormales se tratan por separado.

\section{Vías de acceso a la etapa conclusional}

No siempre se llega a la etapa conclusional por el mismo camino. Hay distintas vías.

\subsection{Vía normal: luego de la etapa probatoria y los alegatos}

Producida la prueba, clausurada la etapa probatoria y vencido el plazo para alegar (hayan presentado o no los alegatos las partes), se entra directamente en la etapa conclusional.

\subsection{Cuestión de puro derecho (art.~359, primer párrafo, CPCCN)}

El primer párrafo del art.~359 CPCCN establece que, contestado el traslado de la demanda o la reconvención, vencidos los plazos para hacerlo y resueltas las excepciones, si la cuestión pudiera ser resuelta como de puro derecho así se decidirá, y firme que se encuentre la providencia se llamará autos para sentencia.

Esto ocurre cuando no hay hechos controvertidos, sino que el debate se refiere únicamente a cuáles son las normas aplicables al caso. En ese supuesto, finalizada la etapa introductoria y antes de la audiencia preliminar, el juez declara que la cuestión es de puro derecho y llama autos para sentencia.

\subsection{Supuestos en la audiencia preliminar (art.~360 inc.~6 y arts.~361--362 CPCCN)}

En la audiencia preliminar también puede resolverse la cuestión como de puro derecho en los siguientes casos:

\begin{itemize}
    \item Las partes reconocen recíprocamente, en el acto de la audiencia, los hechos que antes eran controvertidos según los escritos de demanda y contestación, con lo que ya no hay hechos controvertidos que requieran prueba.
    \item Una de las partes se opuso a la apertura a prueba y el juez hace lugar a esa oposición.
    \item Las partes desisten de todas las pruebas ofrecidas (por ejemplo, porque entienden que con la prueba documental acompañada es suficiente).
    \item Ya se realizó prueba anticipada (art.~326 CPCCN) y el juez entiende que con esa prueba es suficiente.
\end{itemize}

\section{El auto para sentencia}

\begin{definition}{Auto para sentencia}
Resolución que da inicio a la etapa conclusional. Puede dictarla el juez de oficio o a petición de parte, mediante un escrito de dos líneas solicitando que el expediente pase a dictado de sentencia.
\end{definition}

Su texto típico es:

\begin{quote}
``Buenos Aires, desde junio de 2020. Autos para sentencia.''
\end{quote}
% TODO: contenido incompleto o ambiguo en las notas originales (la fórmula "desde junio de 2020" figura así en el apunte).

\subsection{Efectos del auto para sentencia}

\begin{enumerate}
    \item \textbf{Cierre del debate:} las partes ya no pueden realizar ningún acto procesal, presentar escritos ni introducir nuevos documentos. Toda la actividad queda clausurada.
    \item \textbf{Límite para documentos nuevos:} se vincula con el art.~335 CPCCN (documentos de fecha posterior). Solo podían acompañarse hasta que el expediente tuviera auto para sentencia. Una segunda oportunidad se presenta ante la Cámara (art.~260 CPCCN), cuando el expediente se eleva para resolver el recurso de apelación.
    \item \textbf{Improcedencia de la caducidad de instancia:} una vez dictado el auto para sentencia, ya no corre el plazo de caducidad de instancia (art.~310 CPCCN), porque las partes no pueden realizar ninguna actividad procesal.
\end{enumerate}

El art.~313 inc.~4 CPCCN establece que es improcedente la caducidad de instancia si se hubiese llamado autos para sentencia, salvo si se dispusiese prueba de oficio, cuando la producción dependiera de la actividad de las partes.

\subsection{Excepción: medidas para mejor proveer (art.~36 inc.~4 CPCCN)}

\begin{definition}{Medidas para mejor proveer}
Diligencias necesarias para esclarecer la verdad de los hechos controvertidos, que el juez puede ordenar de oficio dejando sin efecto el auto para sentencia, siempre que no afecten el derecho de defensa de las partes y no suplan la negligencia u omisiones de las partes.
\end{definition}

\begin{example}{Límite de las medidas para mejor proveer}
En un juicio de mala praxis en el que la parte actora no ofreció una prueba pericial médica, el juez no puede, como medida para mejor proveer, ordenar la realización de un dictamen pericial médico. Sí puede, en cambio, en un accidente de tránsito en el que un dictamen pericial médico resulta oscuro en algunas cuestiones, dejar sin efecto el auto para sentencia y fijar una audiencia convocando a las partes y al perito para que brinden aclaraciones.
\end{example}

\begin{example}{Oficio a la municipalidad}
En un accidente de tránsito donde se discute el funcionamiento de un semáforo, el juez puede librar un oficio a la municipalidad correspondiente para que informe si el semáforo de determinada intersección el día del accidente funcionaba o estaba fuera de servicio.
\end{example}

\begin{important}
Si el juez ordena la medida pero deja la diligencia a cargo de las partes (por ejemplo, que la parte impulse el diligenciamiento del oficio), vuelve a computarse el plazo de caducidad de instancia mientras no se conteste el oficio y la parte no peticione que el expediente regrese a sentencia.
\end{important}

\subsection{Excepción: cambio de juez}

Si el juez que intervenía se jubila o toma licencia y quien dictará sentencia es otro magistrado, se deja sin efecto el auto para sentencia para notificar a las partes quién es el nuevo juez y permitir que ejerzan, si corresponde, la recusación con causa. Una vez vencido el plazo sin cuestionamientos, el expediente vuelve a estar en estado de sentencia.

\section{Plazos para dictar sentencia y notificación}

Una vez dictado el auto para sentencia, comienza a correr para el juez el plazo para dictar sentencia:

\begin{itemize}
    \item \textbf{Primera instancia} (proceso ordinario): 40 días.
    \item \textbf{Cámara} (tribunal colegiado): 60 días desde que la sala dicta su propio auto para sentencia (por ser tres jueces que deben votar).
\end{itemize}

La sentencia debe ser notificada \textbf{de oficio} por el juzgado o tribunal a todas las partes del proceso (actor, demandado, terceros, citadas en garantía) y a los profesionales a quienes se hayan regulado honorarios. Se notifica la parte dispositiva (la decisión final, también llamada ``fallo'').

\textbf{Excepción: demandado rebelde.} Aunque durante el proceso las notificaciones al rebelde se hacen por ministerio de ley, la sentencia sí debe notificársele. Esta cédula se dirige al domicilio real denunciado y es en formato papel, no electrónico. Su confección suele quedar a cargo de la parte actora.

\section{Morosidad judicial (arts.~167 y 168 CPCCN)}

El código prevé consecuencias para el juez que no dicta sentencia en los plazos establecidos. Se distinguen dos supuestos:

\begin{table}[H]
\centering
\begin{tabularx}{\textwidth}{>{\raggedright\arraybackslash}p{5cm} >{\raggedright\arraybackslash}X}
\toprule
\textbf{Supuesto de morosidad} & \textbf{Consecuencia} \\
\midrule
Demora en providencias simples o interlocutorias & Pesará negativamente en la calificación del juez (mal desempeño de funciones). No hay multa. \\
\addlinespace
Demora en el dictado de la sentencia definitiva & El juez puede ser sancionado con una multa de hasta el 15\% de su remuneración básica, y/o el expediente puede ser remitido a otro juez para que dicte sentencia. \\
\bottomrule
\end{tabularx}
\end{table}

Si el juez advierte que no va a poder dictar sentencia en término, debe presentar una nota a su superior (la Cámara, si es juez de primera instancia; la Corte, si es integrante de una sala de Cámara) con \textbf{diez días de anticipación}, explicando los motivos. Si esos motivos son atendibles, se concederá una prórroga o se evitará la multa y, de corresponder, se determinará si es más idóneo que la dicte otro juez.

\cornelltitle{la etapa conclusional}
\cornellblock{
\cqrow{¿Qué es la etapa conclusional y a través de qué vías se accede a ella?}{Es la etapa en que el juez dicta la sentencia (norma individual). Se accede por: (1) la vía normal: luego de la etapa probatoria y vencido el plazo para alegar; (2) la cuestión de puro derecho (art.~359, párr.~1.º): cuando no hay hechos controvertidos; (3) los supuestos del art.~360 inc.~6 y de los arts.~361 y 362: reconocimiento recíproco en la audiencia, oposición a la apertura a prueba, desistimiento de prueba, prueba anticipada suficiente.
}
\cqrow{¿Qué es el auto para sentencia y qué efectos produce?}{Es la resolución que da inicio a la etapa conclusional. Produce: (a) el cierre del debate (no más escritos ni documentos nuevos); (b) la improcedencia de la caducidad de instancia (art.~313 inc.~4); (c) el inicio del plazo para que el juez dicte sentencia (40 días en primera instancia; 60 días en Cámara). El juez puede dejarlo sin efecto para dictar medidas para mejor proveer (art.~36 inc.~4) o ante el cambio de magistrado.
}
\cqrow{¿Qué son las medidas para mejor proveer?}{Diligencias ordenadas de oficio por el juez (art.~36 inc.~4 CPCCN) para esclarecer la verdad de los hechos controvertidos, siempre que no afecten el derecho de defensa ni suplan la negligencia de las partes. Ejemplos: convocar al perito para aclarar un dictamen oscuro; librar oficio a la municipalidad para informar si un semáforo funcionaba. Si la diligencia queda a cargo de las partes, vuelve a correr la caducidad de instancia.
}
\cqrow{¿Cómo se notifica la sentencia?}{De oficio, por el juzgado o tribunal, a todas las partes (actor, demandado, terceros, citadas en garantía) y a los profesionales a quienes se hayan regulado honorarios; se notifica la parte dispositiva. Al demandado rebelde se le notifica por cédula en papel dirigida al domicilio real denunciado, cuya confección suele quedar a cargo de la parte actora.
}
}
\medskip
\cornellblocksum{
\cqrow{¿Qué consecuencias tiene la morosidad judicial en el dictado de sentencia?}{Si el juez no dicta sentencia en el plazo establecido (40 días en primera instancia, 60 en Cámara), puede ser sancionado con una multa de hasta el 15\% de su remuneración básica, y/o el expediente puede ser remitido a otro juez. Si advierte que no podrá cumplir el plazo, debe notificar a su superior con 10 días de anticipación; si los motivos son atendibles, puede obtenerse una prórroga.
}
}{una vez vencido el plazo para alegar, el proceso entra en la etapa conclusional, que se materializa mediante el auto para sentencia: resolución que cierra el debate, impide nuevas presentaciones de las partes, hace improcedente la caducidad de instancia y activa el plazo del juez para dictar sentencia (40 días en primera instancia, 60 en Cámara). A esta etapa se puede arribar también directamente desde la etapa introductoria si la cuestión es de puro derecho ---sin hechos controvertidos--- o por distintos supuestos que pueden darse en o en torno a la audiencia preliminar.
}

\chapter[Valoración de la prueba]{Sistemas de valoración de la prueba}\label{cap:valoracion}

%==============================================================

\section{Concepto de valoración de la prueba}

\begin{definition}{Valoración de la prueba}
Valorar la prueba significa apreciarla; en sentido figurado, ponerle un precio. Llegado el momento de la sentencia, el juez valora la prueba producida en la etapa probatoria.
\end{definition}

En este punto cobra especial importancia la \textbf{inmediación}.

\begin{definition}{Inmediación}
Contacto directo del juez con la producción de la prueba.
\end{definition}

No es lo mismo que el testigo haya declarado ante un funcionario del juzgado que transcribió su declaración en un acta, a que el testigo haya declarado en presencia del juez, con una declaración filmada que el juez puede reproducir al dictar sentencia. La inmediación permite una sentencia más eficaz, porque el juez aprecia de mejor modo la prueba.

\section{Sistemas de apreciación de la prueba (según Falcón)}

\subsection{Sistema de apreciación prohibida o excluida}

La \textbf{apreciación prohibida} refiere a los casos en que el legislador estableció normativamente la prohibición de ciertos medios de prueba. Por ejemplo, el art.~427 CPCCN establece quiénes no pueden ser testigos (parientes ascendientes, descendientes). Si declararan, el juez no puede apreciar esa declaración.

La \textbf{apreciación excluida} refiere a la facultad del juez (art.~386, párr.~2.º, CPCCN) de no estar obligado a ponderar una por una todas las pruebas: puede excluir aquellas que no considere trascendentes y detenerse en las que resulten más eficientes para formar su convicción.

\subsection{Sistema de la prueba tasada o prueba legal}

En este sistema, el legislador establece de antemano qué valor debe asignársele a cada prueba. Tiene un origen histórico, cuando se consideraba que tal vez el juez no estaba lo suficientemente formado para apreciar la prueba.

\begin{example}{Regla tasada}
Una regla tasada podría ser: ``un hecho solo puede probarse con dos o más testigos; un solo testigo constituye semiplena prueba''.
\end{example}

Es un sistema que puede dar más certezas, pero también mayor injusticia: muchas veces los testigos no se consiguen, y tasar las pruebas puede resultar injusto porque estas pueden ser más o menos valiosas en función de los hechos del caso concreto. Aplicado de modo único y aislado, puede no ser el más eficaz; prácticamente ya no se utiliza de modo único.

\subsection{Sistema de apreciación no tasada o libre convicción}

Dentro de este sistema existe un espectro:

\begin{itemize}
    \item En un extremo, la \textbf{apreciación arbitraria}, donde el juez valora la prueba sin seguir ningún tipo de regla.
    \item En el otro extremo, la \textbf{libre convicción}, donde el juez valora la prueba siguiendo las reglas de la lógica, la experiencia y la intuición (entendida como sentido común).
\end{itemize}

\subsection{Sistema de apreciación conjunta}

No es un sistema puro. En él se aplican, según la prueba de que se trate, criterios de los distintos sistemas anteriores: a veces la prueba tasada, a veces la prohibida o excluida, a veces la libre convicción.

\section{La sana crítica: el sistema del Código (art.~386 CPCCN)}

\begin{formula}{Art. 386 CPCCN --- Regla de valoración de la prueba}
Salvo disposición legal en contrario, los jueces formarán su convicción respecto de la prueba de conformidad con las reglas de la sana crítica. No tendrán el deber de expresar en la sentencia la valoración de todas las pruebas producidas, sino únicamente de las que fueren esenciales y decisivas para el fallo de la causa.
\end{formula}


\begin{definition}{Sana crítica}
Sistema de apreciación de la prueba que es científico y que se basa en las reglas de la experiencia y la lógica. También se lo denomina conjunto de reglas que permiten esclarecer la verdad al momento de dictar sentencia, siempre librado de todo vicio y error. Es el sistema adoptado por el legislador en el CPCCN.
\end{definition}

La sana crítica es, además, el razonamiento que sigue el juez, libre de error y de vicio, para determinar cuál de las dos posiciones asumidas por las partes en el proceso es la correcta. Está sistematizada de modo que el juez pueda demostrar cómo formó su convencimiento de manera concreta, lógica y razonada, y que las partes puedan controlar la construcción de la sana crítica y, si necesitan impugnarla, tengan los fundamentos para hacerlo.

\begin{note}
La expresión «salvo disposición legal en contrario» del artículo 386 remite a la prueba legal tasada: si el legislador le asignó valor a la prueba, el juez no puede apartarse de ella y debe comenzar por esas pruebas legales tasadas. Fuera de ese supuesto, el juez forma su convicción a través de las reglas de la sana crítica.
\end{note}

Si el juez comete un error al valorar la prueba o al fijar los hechos, la sentencia podrá ser objeto de recurso de apelación, para que la Cámara subsane ese error.

\section{Las reglas de la sana crítica según Falcón}

Falcón enumera las siguientes reglas que el juez debe seguir al momento de dictar sentencia:

\begin{longtable}{>{\raggedright\arraybackslash}p{1.6cm} >{\raggedright\arraybackslash}p{3.8cm} >{\raggedright\arraybackslash}p{7.6cm}}
\toprule
\textbf{Regla} & \textbf{Denominación} & \textbf{Desarrollo} \\
\midrule
\endfirsthead
\toprule
\textbf{Regla} & \textbf{Denominación} & \textbf{Desarrollo} \\
\midrule
\endhead

\textbf{1} & \textbf{Solo se prueban los hechos afirmados} & El juez, al dictar sentencia, no puede tener en cuenta otros hechos que no sean los afirmados por las partes. Los hechos no afirmados no existen. Está vinculado con el principio de congruencia. \\
\addlinespace
\textbf{2} & \textbf{Los hechos a probar deben ser controvertidos} & Solo los hechos afirmados y además controvertidos determinan el ámbito probatorio. El juez valorará la prueba producida para acreditar esos hechos. \\
\addlinespace
\textbf{3} & \textbf{Primero debe aplicarse la prueba tasada (si existe)} & Si hay alguna norma que le da cierto valor a determinada prueba, el juez debe aplicarla primero. Ejemplo: art.~579 CCyCN; en un juicio de filiación, la negativa del supuesto progenitor a realizarse la prueba genética se considera indicio grave en su contra. \\
\addlinespace
\textbf{4} & \textbf{Análisis estático de cada medio probatorio} & El juez analiza individualmente todos los medios de prueba producidos para ver cuál es más fiable. El código no establece un orden de valor, pero la doctrina entiende que en principio el orden es: 1.º documental, 2.º informativa, 3.º confesional, 4.º pericial, 5.º testimonial. \\
\addlinespace
\textbf{5} & \textbf{Análisis dinámico y en conjunto de la prueba} & Una vez analizada cada prueba estáticamente, el juez hace un análisis conjunto a la luz de los hechos del caso. Este análisis puede llevar a dar mayor valor a un medio probatorio distinto del orden estático. Ejemplo: en un accidente de tránsito, la declaración de un testigo presencial puede tener más valor que una exposición policial que no aclara cómo ocurrió el hecho. \\
\addlinespace
\textbf{6} & \textbf{Presunciones judiciales y conducta de las partes} & Para obtener la amalgama que estructura el relato de manera coordinada, el juez puede utilizar presunciones judiciales y considerar la conducta desarrollada por las partes a lo largo del proceso (art.~163 inc.~5 CPCCN). \\
\addlinespace
\textbf{7} & \textbf{En caso de duda: carga probatoria} & En caso de duda al momento de valorar la prueba, el juez debe tener por no acreditado el hecho y aplicar las reglas de la carga de la prueba (art.~377 CPCCN): quién tenía la carga de probar ese hecho y si produjo o no la prueba. \\
\addlinespace
\textbf{8} & \textbf{Todo debe volcarse por escrito en la sentencia} & Toda la actividad valorativa debe expresarse por escrito en la sentencia, explicando los fundamentos, argumentos y pasos seguidos para fijar ciertos hechos y para obtener la certeza en el pronunciamiento. La convicción no solo tiene que ser del magistrado, sino también de terceros; de allí la importancia de la construcción de la sentencia. \\
\bottomrule
\end{longtable}

\section{Ejemplos de sentencias reales (accidentes de tránsito)}

Se presentan dos sentencias reales de expedientes de accidentes de tránsito para ejemplificar las reglas de la sana crítica.

\subsection{Ejemplo 1: sentencia con valoración de testigos presenciales}

La sentencia se estructura en tres partes.

\textbf{A) Resultandos (descripción lineal y objetiva).} El juez relata cronológicamente todo lo que consta en el expediente:

\begin{quote}
``A fojas 1/40 se presenta la actora con el patrocinio letrado de tales abogados, promueve demanda contra el demandado Juárez y, si está en garantía, a determinada compañía de seguros. La pretensión tiene por objeto que se condene al demandado a pagar la suma de \$100.000 o lo que más o menos resulte de la prueba del expediente, más intereses y costas. Relata que el día 1.º de diciembre de 2020, cuando era tal hora, iba en su vehículo y [relata la ocurrencia del accidente]. A fojas 75/80 se presenta el demandado, contesta demanda y niega todos los hechos invocados. A fojas 100/120 se presenta la aseguradora, reconoce que existía una póliza de seguros vigente, pero se adhiere a la contestación del demandado. A fojas 98 se celebra la audiencia preliminar. A fojas 157 se clausura el período probatorio. A fojas 160/161 obra el alegato de la parte actora; las restantes partes no presentaron alegatos. A fojas 163 se llama autos para sentencia.''
\end{quote}

\textbf{B) Considerandos (análisis y valoración de la prueba).} El juez aplica las reglas de la sana crítica. En primer lugar, determina qué debe probarse y a quién correspondía la carga:

\begin{quote}
``Efectuada la síntesis de lo actuado, lo que corresponde es determinar si se probó la ocurrencia del hecho, cuya carga recaía sobre la parte actora (artículo 377 del Código Procesal).''
\end{quote}

Luego analiza la prueba testimonial:

\begin{quote}
``A tales fines, la actora ofreció y produjo las testimoniales de los señores Fernández y Palombo, quienes declararon en mi presencia en el marco de un proyecto de procesos orales firmados [indica ver el sobre reservado en el expediente a fojas 159]. Ambos testigos dieron versión concreta de lo acontecido, señalaron dónde se encontraban, qué fue lo que observaron. Dijeron que estaban fumando en el balcón que da a la calle Laprida cuando se produjo el accidente, que vieron cómo fue la colisión [transcribe las partes trascendentes de los testimonios].''
\end{quote}
% TODO: contenido incompleto o ambiguo en las notas originales (las expresiones "si está en garantía" y "proyecto de procesos orales firmados" figuran así en el apunte).

Y valora la credibilidad de los testigos:

\begin{quote}
``No detecté en aquel momento ningún elemento que me hiciera dudar de su veracidad en el momento en que declararon ante él. Todo lo contrario, los recuerdo como testigos confiables que relataron lo que vieron. Ahora, al repasar sus testimonios gracias a las posibilidades que ofrece el registro fílmico, tengo la misma impresión. A ello corresponde añadir que nadie cuestionó el contenido de sus declaraciones ni si no fueron impugnadas por la otra parte por considerar que no eran idóneos o que tenían contradicciones o que habían faltado a la verdad.''
\end{quote}

\textbf{C) Parte dispositiva o fallo.} El juez concluye:

\begin{quote}
``Por ende, tendré por acreditada la materialidad del hecho.''
\end{quote}

Tenido por acreditado el hecho, el juez determina quién es el responsable y pasa a analizar los rubros indemnizatorios reclamados a la luz de la prueba producida, para fijar la procedencia de cada uno y los montos.

\subsection{Ejemplo 2: sentencia con impugnación de testigo (concubina del demandado)}

En otra sentencia, el juez explicita que aplica la sana crítica y explica por qué excluye un testimonio:

\begin{quote}
``Las pruebas han de apreciarse de conformidad a las reglas de la sana crítica (artículo 386 del Código Procesal). Cabe aclarar que los jueces no estamos obligados a ponderar una por una [las pruebas] ---lo que establece el artículo 386---.''
\end{quote}

El juez menciona luego la declaración de una testigo ofrecida por el demandado:

\begin{quote}
``A fojas 210 obra la declaración testimonial de la señora González, quien brindó la siguiente versión de los hechos [la transcribe]. Sin embargo, su declaración fue impugnada por la parte actora a fojas [tanto], ya que la testigo dijo ser pareja/concubina del demandado, afirmando también que al momento de prestar declaración estaba esperando un hijo del demandado.''
\end{quote}

El CPCCN excluye expresamente como testigo al cónyuge, no a la concubina. Sin embargo, en este caso, durante la audiencia se le hicieron preguntas para acreditar sus datos personales y la testigo reconoció ser concubina del demandado y estar esperando un hijo suyo. Por eso la actora impugnó esa declaración.

El juez resuelve:

\begin{quote}
``Lo importante de un testigo es que la declaración sea idónea para crear su convicción respecto de lo que sucedió en la realidad de los hechos, y que no rige aquí la regla de que el testigo único es un testigo nulo. Pero dice: el hecho de que la testigo sea concubina del demandado obliga a examinarla con el mayor rigor, por lo que, ante la ausencia de otros elementos probatorios que precisen las declaraciones que ella dijo y que se habían atacado a través de la impugnación, no puede tenerse en cuenta su declaración.''
\end{quote}

El juez realiza así un \textbf{análisis dinámico}: se aparta de esa declaración no solo por la relación de la testigo con el demandado, sino también porque no existen otros elementos probatorios que corroboren lo que ella dijo. Por eso no resulta suficiente para formar convicción sobre lo ocurrido.

\cornelltitle{valoración de la prueba}
\cornellblock{
\cqrow{¿Cuáles son los 4 sistemas de valoración de la prueba?}{Según Falcón: (1) prohibida o excluida (el legislador prohíbe ciertos medios o el juez puede excluir los no trascendentes ---art.~386, párr.~2.º---); (2) tasada o legal (la ley predetermina el valor de cada prueba; hoy casi no se usa en forma pura); (3) libre convicción (sin reglas ---arbitraria--- o con reglas de lógica y experiencia); (4) conjunta (mezcla de los anteriores según la prueba de que se trate).
}
\cqrow{¿Qué es la sana crítica y cuál es su base legal?}{Es el sistema adoptado por el CPCCN (art.~386). Es un sistema científico de apreciación de la prueba basado en las reglas de la lógica, la experiencia y el sentido común (intuición). Permite al juez esclarecer la verdad al dictar sentencia, libre de todo vicio y error. Los jueces no están obligados a valorar todas las pruebas producidas, solo las esenciales y decisivas.
}
\cqrow{¿Cuáles son las 8 reglas de la sana crítica según Falcón?}{(1) Solo se prueban los hechos afirmados. (2) Esos hechos deben ser además controvertidos. (3) Primero se aplica la prueba tasada, si existe. (4) Análisis estático de cada medio probatorio (orden de fiabilidad doctrinario: documental, informativa, confesional, pericial, testimonial). (5) Análisis dinámico y conjunto (puede variar el orden según el caso). (6) Uso de presunciones judiciales y de la conducta de las partes (art.~163 inc.~5). (7) En caso de duda, aplicar la carga probatoria (art.~377). (8) Todo debe volcarse por escrito en la sentencia, con fundamentos.
}
\cqrow{¿Por qué el análisis dinámico puede alterar el orden de fiabilidad de la prueba?}{Porque el valor de una prueba depende de los hechos del caso concreto. Ejemplo: aunque doctrinariamente la documental está por encima de la testimonial, en un accidente de tránsito donde lo controvertido es cómo ocurrió el hecho, la declaración de un testigo presencial que vio la colisión tendrá más valor que una simple exposición policial que no aclara los detalles de la ocurrencia del accidente.
}
}
\medskip
\cornellblocksum{
\cqrow{¿Cómo se aplicaron las reglas de la sana crítica en las sentencias reales analizadas?}{En la primera sentencia, el juez acreditó el hecho con la testimonial de testigos presenciales, que declararon en su presencia, cuyo contenido nadie cuestionó. En la segunda, descartó la declaración de la concubina del demandado por la falta de elementos que la corroboraran, aplicando el rigor que impone su relación con una de las partes.
}
}{al momento de dictar sentencia, el juez valora la prueba conforme al sistema de la sana crítica (art.~386 CPCCN): un sistema científico basado en la lógica, la experiencia y el sentido común que, según Falcón, se articula en ocho reglas, que van desde tomar en cuenta solo los hechos afirmados y controvertidos, pasando por el análisis estático y dinámico de la prueba, hasta la exigencia de volcar por escrito y con fundamentos toda la actividad valorativa en la sentencia. Dos sentencias reales de accidentes de tránsito ilustran cómo estas reglas se aplican: el hecho se acredita con la testimonial de testigos presenciales (dando mayor valor dinámico a esa prueba sobre la exposición policial) y se descarta la declaración de la concubina del demandado por falta de corroboración.
}

\partintro{Esta parte estudia las distintas formas en que concluye un proceso civil y las resoluciones con que el tribunal se pronuncia. Primero se examinan los \textbf{modos anormales de terminación} ---es decir, aquellos que no son la sentencia definitiva--- y la caducidad de la instancia. Luego se analiza la clasificación de las \textbf{resoluciones judiciales}, el contenido de la sentencia definitiva, la actuación del juez con posterioridad a la sentencia, la \textbf{cosa juzgada} y la relación entre la sentencia civil y la sentencia penal.


\begin{note}
Para comprender cómo se relacionan estos temas puede imaginarse el proceso judicial como un partido de fútbol: las resoluciones judiciales son las decisiones del árbitro, desde una simple advertencia hasta la tarjeta roja o el gol válido. Cada decisión tiene su nombre, su nivel de complejidad y el recurso (la herramienta) que el litigante puede usar para impugnarla. Cuando el partido termina y el resultado es definitivo, sin posibilidad de revancha, se está ante la cosa juzgada. Y si además hubo una investigación disciplinaria paralela (proceso penal), ambas decisiones deben coordinar sus conclusiones para evitar un <<escándalo lógico>>.
\end{note}}
\part{La terminación del proceso y la sentencia}

\chapter[Modos anormales de terminación]{Modos anormales de terminación del proceso}\label{cap:modos}

\section{Encuadre: modo normal y modos anormales de terminación}

Este capítulo y el siguiente abordan los \textbf{modos anormales de terminación del proceso civil}, regulados en los artículos 304 a 318 del Código Procesal Civil y Comercial de la Nación (CPCCN). A diferencia del modo normal de terminación ---la sentencia definitiva---, estos institutos permiten que el proceso concluya de manera anticipada o por circunstancias distintas al dictado de una sentencia.

Las etapas del proceso ---la etapa introductoria, la etapa probatoria y la presentación de los alegatos--- conducen a la sentencia, que constituye el \textbf{modo normal} de terminar el proceso (sin perjuicio de los recursos de apelación): se trata de la \textbf{sentencia definitiva}. Los modos anormales, en cambio, definen otras formas en que un proceso puede finalizar.

\begin{note}
\textbf{Aclaración terminológica.} La expresión <<modo anormal>> es la denominación utilizada por el CPCCN; se trata de un término arbitrario adoptado por el Código. Algunos autores prefieren hablar de <<modos extraordinarios>> de terminación del proceso, por oposición al modo <<ordinario>> o <<normal>> (la sentencia). La diferencia es meramente semántica.
\end{note}

El CPCCN prevé \textbf{cinco} modos anormales de terminación del proceso, regulados a partir del artículo 304, que se estudian en el siguiente orden:

\begin{enumerate}
    \item El \textbf{desistimiento} ---del proceso o del derecho---, acto unilateral o bilateral del actor o reconviniente.
    \item El \textbf{allanamiento} ---sometimiento total o parcial del demandado a las pretensiones del actor---.
    \item La \textbf{transacción} ---contrato bilateral con concesiones recíprocas que extingue obligaciones litigiosas---.
    \item La \textbf{conciliación} ---acuerdo intraprocesal celebrado ante el juez---.
    \item La \textbf{caducidad de la instancia} ---extinción del proceso por inactividad de la parte que tiene la carga de impulsarlo, operada en los plazos previstos en el artículo 310 del CPCCN---.
\end{enumerate}

Cada instituto se analiza en sus requisitos, efectos procesales y diferencias con figuras afines, con especial énfasis en la jurisprudencia y la práctica profesional.

\begin{example}{Cómo se relacionan los modos anormales}
Una persona inicia una demanda por cobro de honorarios. Lo normal sería que el proceso avance hasta la sentencia. Pero el actor puede \textbf{desistir} antes de que se notifique al demandado; el deudor puede \textbf{allanarse} y pagar; ambos pueden llegar a una \textbf{transacción}, o el juez puede convocarlos a \textbf{conciliación}. Y si nadie actúa durante 6 meses, opera la \textbf{caducidad de la instancia}. Todos estos son modos en que el proceso puede terminar sin sentencia.
\end{example}

% ----------------------------------------------------------
\section{El desistimiento (art. 304 y ss., CPCCN)}

Desistir, en términos procesales, significa \textbf{renunciar} o \textbf{abandonar}. % TODO: contenido incompleto o ambiguo en las notas originales (las notas originales utilizan también el verbo <<aplicar>> como sinónimo de desistir)
El desistimiento es el primer modo anormal de terminación del proceso previsto en el CPCCN y su nota característica es la renuncia o el abandono de una posición procesal. Puede ser un acto \textbf{unilateral} o \textbf{bilateral}, según la etapa procesal en que se produzca.

\subsection{¿Quién puede desistir?}

El desistimiento corresponde a quien \textbf{articula una pretensión} y no a quien plantea una defensa. Pueden desistir:

\begin{itemize}
    \item el \textbf{actor}, en la demanda principal;
    \item el \textbf{reconviniente}, en caso de existir reconvención.
\end{itemize}

El demandado que se limita a defenderse no puede desistir, ya que no articula una pretensión propia: el desistimiento es un acto de quien plantea la pretensión y no de quien se defiende.

\subsection{Tipos de desistimiento: del proceso y del derecho}

Como modo anormal de terminar el proceso deben distinguirse dos supuestos: el \textbf{desistimiento del proceso} y el \textbf{desistimiento del derecho}. Este último se encuentra en algunos libros como \textbf{desistimiento de la acción}. Ambas modalidades tienen efectos radicalmente distintos.

\begin{table}[H]
\centering
\caption{Desistimiento del proceso y desistimiento del derecho}
\small
\begin{tabularx}{\textwidth}{
    >{\raggedright\arraybackslash}p{0.22\textwidth}
    >{\raggedright\arraybackslash}X
    >{\raggedright\arraybackslash}X
}
\toprule
\textbf{Aspecto} & \textbf{Desistimiento del proceso} & \textbf{Desistimiento del derecho / acción} \\
\midrule
¿Quién puede hacerlo? & Solo el actor (o reconviniente) & Solo el actor (o reconviniente) \\
Conformidad de la contraria & Necesaria si ya se notificó el traslado de demanda & No se requiere (es unilateral siempre) \\
Efecto sobre el derecho & No lo extingue; puede replantearse en otro proceso & Lo extingue definitivamente \\
Efecto de cosa juzgada & No genera cosa juzgada material sobre el derecho & Opera con efecto de cosa juzgada material \\
Posibilidad de nuevo proceso & Sí, siempre que no esté prescripto el derecho & No; el derecho fue renunciado \\
\bottomrule
\end{tabularx}
\end{table}

\subsubsection{Desistimiento del derecho}

Si el actor desiste del \textbf{derecho}, es decir, de la pretensión articulada, \textbf{no se requiere la conformidad de la otra parte}. El efecto que provoca es el de la \textbf{cosa juzgada}: el actor no podrá reclamar ese derecho en otro proceso, porque lo desistió. El desistimiento opera con calidad de cosa juzgada y esa parte no podrá hacer valer ese derecho en un proceso distinto.

Se trata de un acto procesal unilateral. Su fundamento radica en que el actor, al renunciar a su propio derecho sustancial, no perjudica al demandado, quien se ve beneficiado, ya que no podrá ser perseguido nuevamente por esa pretensión. El efecto es análogo al de la cosa juzgada material.

\subsubsection{Desistimiento del proceso}

El desistimiento del \textbf{proceso} puede realizarse de forma unilateral y sin requerir la conformidad de la otra parte \textbf{únicamente antes de notificar la demanda} al demandado. Si ya se notificó el traslado de la demanda, se requiere la \textbf{conformidad del demandado}.

La razón es que este desistimiento no genera el mismo efecto respecto de la cosa juzgada: el derecho no se desistió, sino solamente se renunció a continuar con esa instancia. Por lo tanto, el derecho puede hacerse valer en otro proceso y el actor podría iniciar otro proceso y formular el mismo reclamo. Por esta circunstancia, una vez notificado el traslado de la demanda, el actor solo puede desistir si cuenta con la conformidad del demandado.

\subsection{Legitimación para desistir}

Para desistir se necesita no solo la legitimación sustancial en términos de capacidad, sino también la \textbf{legitimación procesal}. Si quien desiste es un abogado con un poder que no es suficiente para desistir de la acción o del derecho, el desistimiento no es válido: el representante debe estar \textbf{expresamente facultado} para desistir, ya sea del proceso o del derecho.

\begin{important}
\textbf{Requisito de legitimación doble.} Para desistir válidamente se requiere: (1) \textbf{legitimación sustancial} ---capacidad de derecho sobre la pretensión--- y (2) \textbf{legitimación procesal} ---el representante o apoderado debe tener poder especial y expreso que lo autorice a desistir---. El poder general de representación no es suficiente.
\end{important}

\subsection{Oportunidad para desistir}

El desistimiento puede llevarse a cabo \textbf{en cualquier instancia del proceso antes del dictado de la sentencia}. Una vez dictada la sentencia, el proceso terminó, por lo que no se podría desistir del proceso. Además, si la sentencia reconoció un derecho, hizo lugar a la pretensión del actor o la rechazó, tampoco se podría desistir. Una vez dictada la sentencia definitiva, el proceso ha concluido por su modo normal y el desistimiento ya no tiene objeto.

% ----------------------------------------------------------
\section{El allanamiento (art. 307, CPCCN)}

El allanamiento es una de las actitudes posibles del demandado al contestar la demanda y se encuentra contemplado en el artículo 307 del CPCCN.

\begin{articulo}{Art. 307 --- CPCCN}
Regula el allanamiento como actitud del demandado frente a las pretensiones del actor. Prevé los efectos según si el allanamiento es total o parcial, y si viene acompañado o no del cumplimiento de la pretensión.
\end{articulo}

\begin{definition}{Allanamiento}
El allanamiento es un acto jurídico procesal que consiste en el \textbf{sometimiento total o parcial del demandado a las pretensiones del actor}.
\end{definition}

Por oposición al desistimiento, el allanamiento corresponde al \textbf{demandado}. Puede allanarse a las pretensiones del actor el demandado, o bien el \textbf{reconvenido}, que se allana a las pretensiones del reconviniente.

\subsection{Naturaleza: ¿reconocimiento de los hechos o del derecho?}

\begin{important}
Allanarse \textbf{no constituye el reconocimiento ni de la verdad de los hechos ni del derecho}. Es un acto voluntario del demandado que, a través de la actitud que asume, permite el progreso de la pretensión articulada por el actor.
\end{important}

Esta es una distinción esencial: el allanamiento no implica que el demandado reconozca que los hechos alegados por el actor sean verdaderos ni que el derecho invocado sea correcto. Simplemente expresa su voluntad de no oponerse al progreso de la pretensión.

\subsection{Características}

El allanamiento es un \textbf{acto unilateral}: no requiere la conformidad de la otra parte. Si el demandado quiere allanarse, no necesita la conformidad del actor.

\subsection{Modalidades: expreso y tácito}

Se trata de un acto unilateral que puede ser \textbf{expreso} o \textbf{tácito}.

\begin{itemize}
    \item \textbf{Expreso:} debe ser contundente y no generar ninguna duda; la parte demandada manifiesta expresamente que se allana, es decir, que se somete total o parcialmente a las pretensiones de la otra parte.
    \item \textbf{Tácito:} siempre genera mayores dudas. Por ello, el allanamiento tácito exige una actitud del demandado que no genere ninguna duda, que sea \textbf{incontrovertible}, de la que surja que se está sometiendo a las pretensiones del actor aunque no lo manifieste expresamente.
\end{itemize}

\begin{example}{Allanamiento tácito claro}
El demandado, sin oponerse a la pretensión del actor, se presenta y paga, satisfaciendo así la totalidad de la pretensión. Aunque no manifestó expresamente <<me allano>>, el efecto es incontrovertiblemente ese.
\end{example}

\begin{example}{Allanamiento tácito dudoso (demanda de desalojo)}
En una demanda de desalojo, la parte demandada se presenta y deposita en el tribunal las llaves del inmueble. ¿Es un allanamiento tácito? Puede haber otras cuestiones en discusión simultáneamente, por lo que la actitud puede no ser incontrovertible. Se deberá analizar el caso concreto.
\end{example}

\subsection{Allanamiento total y parcial: efectos}

% TODO: contenido incompleto o ambiguo en las notas originales (se plantea qué sucede si el allanamiento es parcial, pero la explicación no se completa; el cuadro siguiente distingue únicamente según se cumpla o no la pretensión)
Los efectos del allanamiento dependen de si, además de allanarse, el demandado cumplió o no con la pretensión.

\begin{table}[H]
\centering
\caption{Efectos del allanamiento según el cumplimiento de la pretensión}
\small
\begin{tabularx}{\textwidth}{
    >{\raggedright\arraybackslash}p{0.20\textwidth}
    >{\raggedright\arraybackslash}X
    >{\raggedright\arraybackslash}X
}
\toprule
\textbf{Supuesto} & \textbf{Allanamiento + cumplimiento de la pretensión} & \textbf{Allanamiento sin cumplimiento} \\
\midrule
¿Qué hace el juez? & Dicta una resolución interlocutoria y manda al archivo las actuaciones (art. 307) & Hace lugar a la pretensión y dicta sentencia definitiva para habilitar la ejecución \\
Tipo de resolución & Resolución interlocutoria & Sentencia definitiva \\
Etapa siguiente & Archivo del expediente & Ejecución de sentencia \\
\bottomrule
\end{tabularx}
\end{table}

El artículo 307 se refiere a la resolución interlocutoria en el caso de que se cumpla la pretensión y a la orden de enviar al archivo las actuaciones.

\subsection{La cosa juzgada en el allanamiento}

La resolución que admite el allanamiento \textbf{hace cosa juzgada}. Pero la cosa juzgada opera con respecto a la \textbf{pretensión} y no con respecto a los \textbf{hechos}, porque el allanamiento no importa el reconocimiento de la verdad de los hechos invocados por el actor ni el reconocimiento del derecho invocado.

En consecuencia, no se podrán invocar esos hechos como verdaderos en otro proceso, como si hubiesen sido reconocidos en una sentencia, porque se trata de otra cosa. Esa es una característica propia de este modo anormal de terminar el proceso.

% ----------------------------------------------------------
\section{La transacción (art. 308, CPCCN; arts. 1641 y ss., CCyCN)}

La transacción es un tercer modo anormal de terminar el proceso. En el CPCCN se regula en el artículo 308.

\begin{articulo}{Art. 308 --- CPCCN}
Las partes podrán hacer valer la transacción del derecho en litigio con la presentación del convenio o suscripción de acta ante el juez. Éste se limitará a examinar la concurrencia de los requisitos exigidos por la ley para la validez de la transacción y la homologará o no. En este último caso, continuarán los procedimientos del juicio.
\end{articulo}

\begin{articulo}{Art. 1641 --- Código Civil y Comercial de la Nación (CCyCN)}
La transacción es un contrato por el cual las partes, para evitar el litigio o ponerle fin, haciéndose concesiones recíprocas, extinguen obligaciones dudosas o litigiosas.
\end{articulo}

\subsection{Concepto y naturaleza}

La transacción proviene del ámbito del derecho civil. Es un \textbf{acto jurídico bilateral} por el cual las partes, haciéndose \textbf{concesiones recíprocas}, extinguen obligaciones litigiosas o dudosas. Define la posibilidad de ponerle fin al litigio, que es el marco que el CPCCN incorpora en el artículo 308 como modo anormal de terminar el proceso. Constituye un medio de extinción de las obligaciones entre las partes que se hacen esas concesiones recíprocas.

Es bilateral (requiere el acuerdo de ambas partes) y puede ser judicial o extrajudicial. En el ámbito procesal, la transacción debe presentarse ante el juez para su homologación.

\subsection{Requisitos formales}

La transacción debe hacerse \textbf{por escrito} y presentarse ante el juez en cuyo proceso están en juego los derechos litigiosos, para que éste la homologue o no. También se puede suscribir un acta ante el juez en el cual tramita la causa.

\begin{articulo}{Art. 1643 --- CCyCN (Forma)}
La transacción debe hacerse por escrito. Si recae sobre derechos litigiosos, sólo es eficaz a partir de la presentación del instrumento firmado por los interesados ante el juez de la causa.
\end{articulo}

\subsection{Materias excluidas de la transacción}

\begin{articulo}{Art. 1644 --- CCyCN (Prohibiciones)}
No puede transigirse sobre derechos en los que está comprometido el orden público, ni sobre derechos irrenunciables. Tampoco pueden ser objeto de transacción los derechos sobre las relaciones de familia o el estado de las personas, excepto que se trate de derechos patrimoniales derivados de aquellos o de otros derechos sobre los cuales expresamente este Código admite pactar.
\end{articulo}

Las cuestiones que afectan el \textbf{orden público} no pueden ser objeto de transacción. Respecto de las \textbf{relaciones de familia}, no es posible, por ejemplo, que las partes acuerden considerarse divorciadas y pacten de ese modo su propio divorcio: esto no funciona así.

Lo que se admite dentro de las relaciones de familia como posible objeto de transacción son las cuestiones derivadas que tengan \textbf{contenido patrimonial}, por ejemplo:

\begin{itemize}
    \item las cuestiones sobre alimentos de los hijos menores;
    \item las cuestiones relativas a la división del régimen de comunidad de bienes, en las que las partes pueden acordar cómo dividir o repartir los bienes conyugales.
\end{itemize}

Tampoco pueden transigirse los \textbf{derechos irrenunciables}, como por ejemplo los derechos laborales o de los trabajadores, que no se pueden pactar en esos términos ni ser objeto de transacción.

% ----------------------------------------------------------
\section{La conciliación (art. 309, CPCCN)}

La conciliación se estudió previamente como sistema de solución alternativa de disputas y también constituye un modo anormal de terminar el proceso. Se encuentra prevista en el artículo 309.

\begin{articulo}{Art. 309 --- CPCCN}
Los acuerdos conciliatorios celebrados por las partes ante el juez y homologados por éste tendrán autoridad de cosa juzgada.
\end{articulo}

\begin{articulo}{Art. 36, inc. 2° --- CPCCN}
En cualquier momento del proceso el juez puede convocar a las partes para alcanzar una conciliación.
\end{articulo}

\begin{articulo}{Art. 360, inc. 1° --- CPCCN}
En la audiencia preliminar, el juez intentará conciliar a las partes y acercarlas a los fines de alcanzar un acuerdo.
\end{articulo}

\subsection{Concepto}

Conciliar significa componer, ponerse de acuerdo. % TODO: contenido incompleto o ambiguo en las notas originales (las notas distinguen la conciliación <<extraprocesal o judicial>> y la <<intraprocesal>> sin precisar cada categoría)
Se distingue la conciliación en la fase extraprocesal o judicial y la intraprocesal, para poner en el centro de la escena la figura del juez como quien convoca la conciliación dentro del proceso.

En los dos supuestos se trata de un \textbf{avenimiento amigable} de las partes. En la conciliación intraprocesal, las diferencias de las partes se arreglan ante la presencia del magistrado o a instancia suya, porque es el propio magistrado quien las convoca.

\subsection{Efecto de cosa juzgada}

Si se alcanza la conciliación y se arriba a un acuerdo, el proceso termina por un modo anormal: no por vía de sentencia, sino por la realización de ese acuerdo. El acuerdo \textbf{extingue las pretensiones antagónicas de las partes} y produce el efecto de \textbf{cosa juzgada} una vez que está \textbf{homologado por el juez}.

\subsection{Conciliación y transacción: diferencias esenciales}

La conciliación y la transacción tienen muchos puntos en común, porque en ambas se alcanza un acuerdo que el juez podrá homologar, pero presentan distinciones importantes.

\begin{table}[H]
\centering
\caption{Diferencias entre transacción y conciliación}
\small
\begin{tabularx}{\textwidth}{
    >{\raggedright\arraybackslash}p{0.19\textwidth}
    >{\raggedright\arraybackslash}X
    >{\raggedright\arraybackslash}X
}
\toprule
\textbf{Aspecto} & \textbf{Transacción} & \textbf{Conciliación} \\
\midrule
Iniciativa & Surge de las partes (voluntad propia) & A instancia o convocatoria del juez \\
Ámbito & Judicial o extrajudicial; derechos litigiosos o dudosos & Siempre intraprocesal (dentro del proceso) \\
Naturaleza jurídica & Contrato (arts. 1641 y ss. CCyCN) & Acuerdo procesal ante el magistrado \\
Estado del derecho & Puede recaer sobre derechos dudosos o litigiosos & Solo sobre derechos ya litigiosos (proceso iniciado) \\
Homologación & Puede o no ser homologado por el juez & Se celebra ante el juez y produce cosa juzgada al homologarse \\
\bottomrule
\end{tabularx}
\end{table}

La diferencia vinculada con la \textbf{voluntad} es la siguiente: en la transacción, la voluntad de las partes se manifiesta en el sentido de querer llevar a cabo esa avenencia ante el juez, de modo que la necesidad de transigir surge de las partes. En cambio, la conciliación surge a raíz de una \textbf{propuesta del juez}: se realiza luego de una convocatoria efectuada por el juez en la audiencia preliminar, en virtud de las facultades previstas en el artículo 36, a lo largo de todo el proceso.

% ----------------------------------------------------------

\cornelltitle{modos anormales de terminación: desistimiento, allanamiento, transacción y conciliación}
\cornellblock{
\cqrow{¿Qué es un modo anormal de terminación del proceso?}{Todo aquel que no es la sentencia definitiva. El CPCCN regula cinco supuestos (arts. 304--318): desistimiento, allanamiento, transacción, conciliación y caducidad de la instancia. Algunos autores los denominan también <<modos extraordinarios>>.
}
\cqrow{¿Quién puede desistir y de qué?}{Solo quien articula una pretensión: el actor o el reconviniente. Puede desistir del proceso (renuncia a continuar esa instancia; no extingue el derecho) o del derecho/acción (renuncia al derecho sustancial; genera cosa juzgada y no puede replantearse). El desistimiento del proceso requiere la conformidad del demandado si ya fue notificado el traslado de la demanda.
}
\cqrow{¿Cuál es la diferencia entre desistir del proceso y desistir del derecho?}{El desistimiento del derecho es siempre unilateral y opera con efecto de cosa juzgada material: el derecho queda extinguido y no puede reclamarse en otro proceso. El desistimiento del proceso no genera cosa juzgada material sobre el derecho: el actor puede iniciar otro proceso, siempre que el derecho no esté prescripto.
}
\cqrow{¿Hasta cuándo puede desistirse?}{En cualquier instancia del proceso antes del dictado de la sentencia. Dictada la sentencia, el proceso terminó por su modo normal y el desistimiento ya no tiene objeto.
}
\cqrow{¿Qué legitimación se requiere para desistir?}{Doble legitimación: sustancial (capacidad sobre el derecho en cuestión) y procesal (el representante debe tener poder especial y expreso para desistir; el poder general no alcanza).
}
\cqrow{¿Qué es el allanamiento y a quién corresponde?}{Es el acto jurídico procesal del demandado (o reconvenido) por el cual se somete total o parcialmente a las pretensiones del actor. Es unilateral (no requiere la conformidad del actor), puede ser expreso o tácito, y no implica reconocimiento de la verdad de los hechos ni del derecho.
}
\cqrow{¿Cómo opera el allanamiento con cumplimiento y sin cumplimiento de la pretensión?}{Con cumplimiento: el juez dicta resolución interlocutoria y manda al archivo las actuaciones. Sin cumplimiento: hace lugar a la pretensión y dicta sentencia definitiva para habilitar la ejecución de sentencia.
}
}
\medskip
\cornellblocksum{
\cqrow{¿Qué cosa juzgada produce el allanamiento?}{La resolución que lo admite hace cosa juzgada respecto de la pretensión, pero no respecto de los hechos (que no fueron reconocidos como verdaderos). No pueden invocarse en otro proceso como hechos reconocidos en una sentencia.
}
\cqrow{¿Qué es la transacción y qué la distingue de la conciliación?}{La transacción es un contrato bilateral en que las partes se hacen concesiones recíprocas para extinguir obligaciones litigiosas o dudosas (art. 1641 CCyCN; art. 308 CPCCN). Surge de la voluntad de las partes y puede ser judicial o extrajudicial. La conciliación, en cambio, surge de la convocatoria del juez y es siempre intraprocesal (art. 309 CPCCN).
}
\cqrow{¿Qué materias no pueden ser objeto de transacción?}{Derechos en que esté comprometido el orden público, derechos irrenunciables (como los laborales) y derechos sobre relaciones de familia o estado de las personas. Excepción: derechos patrimoniales derivados (alimentos, división de bienes conyugales).
}
}{Los modos anormales de terminación del proceso, regulados en los arts. 304 a 318 del CPCCN, se oponen al modo normal, que es la sentencia definitiva. El desistimiento es un acto unilateral del actor o reconviniente con efectos distintos: el desistimiento del derecho genera cosa juzgada material, mientras que el del proceso solo extingue la instancia sin afectar el derecho. El allanamiento es el acto unilateral del demandado por el que se somete total o parcialmente a las pretensiones del actor; no importa reconocimiento de hechos ni de derecho y produce cosa juzgada limitada a la pretensión. La transacción (contrato bilateral con concesiones recíprocas, arts. 1641 y ss. CCyCN y art. 308 CPCCN) y la conciliación (acuerdo intraprocesal ante el juez, art. 309 CPCCN) son institutos afines que difieren en la iniciativa (partes frente a juez) y en el ámbito (judicial o extrajudicial frente a siempre intraprocesal).
}

\chapter{La caducidad de la instancia}\label{cap:caducidad}

\section{Marco normativo y concepto}

La caducidad de la instancia está prevista en el CPCCN en los artículos 310 a 318.

\begin{important}
\textbf{Advertencia sobre derecho comparado provincial.} La regulación de la caducidad de la instancia varía entre provincias: no todos los códigos la tratan de igual modo, y el modo en que opera en el proceso judicial depende de cómo esté prevista en el Código Procesal de que se trate. El análisis de este capítulo se circunscribe al sistema del CPCCN, aplicable en el fuero federal y en la Ciudad Autónoma de Buenos Aires.
\end{important}

La caducidad de la instancia es el modo anormal de terminación del proceso que ocupa la mayor cantidad de artículos del capítulo y, en la práctica profesional, en la casuística y en la interpretación de los jueces en los casos concretos, es el instituto que ha tenido más diversas interpretaciones en la jurisprudencia.

\begin{definition}{Caducidad de la instancia}
Extinción del proceso por la inactividad de la parte que tiene la carga de impulsarlo, operada en los plazos previstos en el artículo 310 del CPCCN.
\end{definition}

\subsection{Nota esencial: la inactividad procesal}

La caducidad de la instancia \textbf{no cancela la pretensión}. Entre los modos anormales de terminación del proceso, esto sucede solamente con el desistimiento del proceso y con la caducidad; % TODO: contenido incompleto o ambiguo en las notas originales (las notas afirman que esto <<pasa solamente>> con el desistimiento del proceso y que en los demás modos <<no sucede>>, sin precisar la comparación con la caducidad)
en los demás no ocurre. Por lo tanto, la misma pretensión puede ser promovida en otro proceso judicial, siempre y cuando el derecho \textbf{no esté alcanzado por la prescripción de la acción}. Son términos que muchas veces se confunden.

\subsection{Ámbito de aplicación: procesos dispositivos}

La caducidad de la instancia es aplicable a los \textbf{procesos dispositivos}, porque en ellos el impulso del proceso está a cargo de las partes. La inactividad de la parte que tiene la carga de impulsarlo genera que, una vez vencido el plazo establecido por el Código, el tribunal pueda declarar, de oficio o a pedido de la parte contraria, el cese del curso de esa instancia y, por ende, la extinción del proceso.

Para que se configure la caducidad deben verificarse, entonces:

\begin{itemize}
    \item un proceso dispositivo, en el que las partes tienen la carga de impulsar el proceso y no lo hacen;
    \item el vencimiento de los plazos previstos en el Código;
    \item inactividad procesal durante todos esos plazos, es decir, que la parte no haya impulsado el proceso.
\end{itemize}

\subsection{Fundamento del instituto}

El Código prevé la caducidad de la instancia como modo anormal de terminar el proceso por dos razones:

\begin{enumerate}
    \item El \textbf{interés público}: que los procesos judiciales no permanezcan paralizados indefinidamente en el tiempo, ocupando casilleros en los tribunales. Si la parte no activa el avance del proceso, se puede presumir su \textbf{abandono}.
    \item La relación procesal no solamente vincula a las partes ---actor y demandado--- sino que también incluye al juez, y el órgano jurisdiccional no puede quedar al arbitrio de los tiempos que las partes quieran manejar en sus procesos judiciales. Por eso es necesario poner un límite y un marco a esta circunstancia.
\end{enumerate}

\subsection{Distinción con la prescripción de la acción}

La prescripción de la acción se funda en una presunción de abandono del derecho: durante el tiempo en que la acción podría haberse hecho efectiva a través de una demanda y un reclamo judicial no se lo hizo, lo que impide hacerlo con posterioridad al vencimiento del plazo de prescripción. De modo análogo, la inactividad dentro del proceso judicial de la parte que tiene la carga de impulsarlo importa la presunción de abandono de esa instancia judicial.

\begin{table}[H]
\centering
\caption{Prescripción de la acción y caducidad de la instancia}
\small
\begin{tabularx}{\textwidth}{
    >{\raggedright\arraybackslash}p{0.21\textwidth}
    >{\raggedright\arraybackslash}X
    >{\raggedright\arraybackslash}X
}
\toprule
\textbf{Aspecto} & \textbf{Prescripción de la acción} & \textbf{Caducidad de la instancia} \\
\midrule
Naturaleza & Instituto de derecho sustancial & Instituto de derecho procesal \\
¿Qué se extingue? & El derecho / la acción & La instancia (el proceso en curso) \\
¿Desde cuándo corre? & Desde que nació el derecho exigible & Desde el último acto impulsor del proceso \\
¿Cancela la pretensión? & Sí (si prescribió, no se puede reclamar) & No (puede iniciarse nuevo proceso si el derecho no prescribió) \\
¿Con qué tipo de días? & Los que disponga cada norma sustancial & Días corridos (incluye inhábiles, salvo ferias judiciales) \\
\bottomrule
\end{tabularx}
\end{table}

Son dos institutos diversos: la prescripción de la acción y la caducidad de la instancia. La caducidad de la instancia \textbf{se abre con la promoción de la demanda}, sin importar que se haya notificado o no el traslado de la demanda.

% ----------------------------------------------------------
\section{Plazos de caducidad (art. 310, CPCCN)}

\begin{articulo}{Art. 310 --- CPCCN (Plazos)}
Se producirá la caducidad de la instancia cuando no se instare su curso dentro de los siguientes plazos:
\begin{enumerate}[label=\arabic*°]
    \item Seis meses, en primera instancia.
    \item Tres meses, en segunda o ulterior instancia y en cualquiera de las instancias en el juicio sumarísimo, en el juicio ejecutivo, en las ejecuciones especiales y en los incidentes.
    \item En el que se opere la prescripción de la acción, si fuese menor a los indicados precedentemente.
    \item Un mes, en el incidente de caducidad de instancia.
\end{enumerate}
\end{articulo}

\begin{table}[H]
\centering
\caption{Plazos de caducidad según el art. 310 del CPCCN}
\small
\begin{tabularx}{\textwidth}{
    >{\raggedright\arraybackslash}p{0.12\textwidth}
    >{\raggedright\arraybackslash}X
}
\toprule
\textbf{Plazo} & \textbf{Supuesto} \\
\midrule
6 meses & Primera instancia \\
3 meses & Segunda o ulterior instancia; juicio sumarísimo; juicio ejecutivo; ejecuciones especiales; incidentes \\
Prescripción & El plazo en que se opere la prescripción de la acción, si fuese menor a los indicados \\
1 mes & Incidente de caducidad de instancia \\
\bottomrule
\end{tabularx}
\end{table}

\begin{example}{Plazo de un mes del incidente de caducidad (inc. 4°)}
Cuando la parte contraria acusa la caducidad de la instancia, fundamenta en el caso concreto por qué se encuentra configurada: indica cuál fue el último acto impulsor y que desde allí se encuentra vencido el plazo de seis meses en primera instancia sin que exista ningún acto procesal que tienda a impulsar el proceso. De esa presentación el juez da traslado a la otra parte, que la contesta, y recién después se resuelve, haciendo lugar o rechazando la caducidad.

El inciso 4° dispone que ese incidente ---el traslado, la contestación y la resolución--- también tiene un plazo de caducidad: \textbf{un mes}. Por ejemplo: una vez acusada la caducidad, el juez ordenó el traslado a la otra parte. Si quien acusó la caducidad no le notifica a la otra parte ---no le envía la cédula electrónica a los fines de notificar ese traslado---, entonces, si pasa un mes, quien debía contestar el traslado pero jamás fue notificado podría acusar la caducidad de ese incidente de caducidad de instancia.
\end{example}

% ----------------------------------------------------------
\section{Cómputo del plazo (art. 311, CPCCN)}

\begin{articulo}{Art. 311 --- CPCCN (Cómputo)}
Los plazos señalados en el artículo precedente se computarán desde la fecha de la última petición de las partes o resolución o actuación del juez, secretario u oficial primero, que tenga por efecto impulsar el procedimiento; correrán durante los días inhábiles salvo los que correspondan a las ferias judiciales.
\end{articulo}

El \textbf{último acto impulsor} es el último acto tomado como válido que significa un impulso procesal. No cualquier acto lo es: por ejemplo, la presentación para autorizar a una persona que trabaja en un estudio jurídico a visualizar el expediente no es un acto impulsor que pueda tomarse como válido en términos de impulso procesal. Debe tratarse de un acto que tenga concretamente por finalidad impulsar el proceso hacia las etapas subsiguientes, es decir, hacerlo avanzar y seguir procesando.

Se toma ese último acto válido que signifique impulso procesal y, a partir del \textbf{día siguiente}, se contabiliza el plazo. Los plazos del artículo 310 corren durante los días inhábiles, salvo los que correspondan a las ferias judiciales, según dispone el artículo 311.

\begin{important}
\textbf{Días inhábiles y ferias.} Los plazos de caducidad de instancia corren en \textbf{días corridos} (incluyendo feriados y días inhábiles), a diferencia de los plazos procesales ordinarios, que se cuentan en días hábiles. La única excepción: no se computan los días correspondientes a la feria de verano (enero) ni a la feria de invierno (julio, 15 días).
\end{important}

% ----------------------------------------------------------
\section{Litisconsorcio (art. 312, CPCCN)}

\begin{articulo}{Art. 312 --- CPCCN (Litisconsorcio)}
El impulso del procedimiento por uno de los litisconsortes beneficiará a los demás.
\end{articulo}

El impulso del procedimiento por uno de los litisconsortes se considera que beneficia a todos los demás, porque la \textbf{instancia es indivisible}: no puede haber caducidad de instancia para un litisconsorte y no para el otro. Por lo tanto, si uno impulsa, en términos del sistema dispositivo, el impulso de la instancia beneficia a todos los demás.

% ----------------------------------------------------------
\section{Supuestos de improcedencia (art. 313, CPCCN)}

\begin{articulo}{Art. 313 --- CPCCN (Improcedencia)}
No se producirá la caducidad:
\begin{enumerate}[label=\arabic*°]
    \item En los procedimientos de ejecución de sentencia, salvo si se tratare de incidentes que no guardaren relación estricta con la ejecución procesal forzada propiamente dicha.
    \item En los procesos sucesorios y, en general, en los voluntarios, también salvo los incidentes que en ellos se susciten.
    \item Cuando los procesos estuvieren pendientes de alguna resolución y la demora en dictarla fuere imputable al tribunal.
    \item Si se hubiere llamado autos para sentencia, salvo si se dispusiere prueba de oficio; cuando su producción dependiere de la actividad de las partes, la carga de impulsar el procedimiento existirá desde el momento en que éstas toman conocimiento de las medidas ordenadas.
\end{enumerate}
\end{articulo}

\begin{itemize}
    \item \textbf{Ejecución de sentencia:} no corresponde acusar la caducidad de instancia, salvo si se tratara de incidentes que no guarden relación estricta con la ejecución procesal forzada propiamente dicha.
    \item \textbf{Procesos sucesorios y voluntarios:} tampoco procede la caducidad en procesos sucesorios ni, en general, en cualquier proceso voluntario.
    \item \textbf{Demora imputable al tribunal:} cuando los procesos estuvieran pendientes de alguna resolución y la demora en dictarla fuera imputable al tribunal, la inactividad debe poder imputarse a la parte y no a una demora que corresponda al tribunal.
\end{itemize}

\begin{example}{Demora imputable al tribunal}
Ya se encuentra dictado el auto para sentencia y éste fue firmado, pero la sentencia no se dictó. No es algo que pueda imputarse a la parte, sino que corresponde al propio tribunal.
\end{example}

% ----------------------------------------------------------
\section{Legitimación y modo de operarse (arts. 314 a 316, CPCCN)}

% TODO: contenido incompleto o ambiguo en las notas originales (el encabezado de las notas menciona los arts. 314 a 316, pero solo se transcriben los arts. 315 y 316)
\begin{articulo}{Art. 315 --- CPCCN (Legitimación para pedir la declaración de caducidad)}
Sin perjuicio de lo dispuesto en el artículo siguiente, la declaración de caducidad podrá ser pedida en primera instancia por el demandado; en el incidente, por el contrario a quien lo hubiese promovido; en el recurso, por la parte recurrida. La petición deberá formularse antes de consentir el solicitante cualquier actuación del tribunal o de la parte posterior al vencimiento del plazo legal, y se sustanciará únicamente con un traslado de la parte contraria.
\end{articulo}

\begin{articulo}{Art. 316 --- CPCCN (Modo de operarse)}
La caducidad será declarada de oficio, sin otro trámite que la comprobación de los plazos señalados en el artículo 310, pero antes de que cualquiera de las partes impulse el procedimiento.
\end{articulo}

La caducidad puede pedirse en primera instancia por el \textbf{demandado}; en el incidente, por el \textbf{contrario a quien lo hubiese promovido}; y en el recurso, por la \textbf{parte recurrida}.

La petición de caducidad debe formularse \textbf{antes de consentir cualquier acto} del tribunal o de la otra parte que pueda considerarse impulsor. Si el plazo de seis meses venció pero, luego de ese vencimiento en primera instancia, se produjo un acto que fue consentido, ya no puede acusarse la caducidad por los seis meses transcurridos anteriormente. Es decir, la caducidad puede acusarse siempre que no se consienta ningún acto impulsor posterior al vencimiento, aunque ya se haya operado el vencimiento del plazo.

Asimismo, el pedido de caducidad en la segunda instancia importa el \textbf{desistimiento del recurso interpuesto por el peticionario}, en el caso de que aquel prospere.

\begin{example}{Caducidad de la segunda instancia con recursos de ambas partes}
Las dos partes interpusieron recursos de apelación contra la sentencia de primera instancia, y la parte demandada acusa la caducidad de la segunda instancia respecto del recurso que interpuso la parte actora. Si esa caducidad se considera configurada y se declara la caducidad de la segunda instancia, se tiene por desistido el recurso de apelación interpuesto por el demandado. Es decir, la Cámara no va a tratar ninguno de los dos recursos.
\end{example}

% ----------------------------------------------------------
\section{Efectos de la caducidad (art. 318, CPCCN)}

\begin{articulo}{Art. 318 --- CPCCN (Efectos)}
La caducidad operada en primera o única instancia no extingue la acción, la que podrá ejercitarse en un nuevo juicio, ni perjudica las pruebas producidas, las que podrán hacerse valer en aquel. La caducidad operada en instancias ulteriores acuerda fuerza de cosa juzgada a la resolución recurrida. La caducidad de la instancia principal comprende la reconvención y los incidentes; pero la de éstos no afecta la instancia principal.
\end{articulo}

\subsection{Caducidad en primera o única instancia}

La caducidad operada en primera o única instancia \textbf{no extingue la acción}: ésta puede ejercitarse en un nuevo juicio, y las pruebas producidas en el proceso que caducó no se perjudican, de modo que pueden hacerse valer en el juicio que se inicie después. Esto no es posible cuando el derecho ya está alcanzado por la \textbf{prescripción}, es decir, cuando la acción se encuentra prescripta.

\subsection{Caducidad en instancias ulteriores}

La caducidad operada en instancias ulteriores \textbf{acuerda fuerza de cosa juzgada a la resolución recurrida}. Si lo que caducó es la segunda instancia, la sentencia definitiva dictada por el juez de primera instancia queda con fuerza de sentencia definitiva pasada en autoridad de cosa juzgada: aunque se hayan interpuesto los recursos de apelación correspondientes, al declararse operada la caducidad en la segunda instancia queda firme la sentencia definitiva de primera instancia.

\subsection{Instancia principal, reconvención e incidentes}

La caducidad de la instancia principal \textbf{comprende la reconvención y los incidentes}: si caduca la instancia principal, se afecta todo lo tramitado (reconvenciones, incidentes, etcétera). En cambio, la caducidad de un incidente generado con motivo del proceso principal \textbf{no} genera que se configure la caducidad para el proceso principal.

% ----------------------------------------------------------

\cornelltitle{caducidad de la instancia}
\cornellblock{
\cqrow{¿Qué es la caducidad de la instancia y cuál es su fundamento?}{Es la extinción de la instancia por inactividad procesal de la parte que tiene la carga de impulsar el proceso, transcurridos los plazos del art. 310 CPCCN. Fundamentos: el interés público en no paralizar indefinidamente los procesos y la tutela del órgano jurisdiccional. Se presume el abandono de la instancia.
}
\cqrow{¿En qué procesos opera y qué la diferencia de la prescripción de la acción?}{Opera en los procesos dispositivos, en los que las partes tienen la carga de impulsar el proceso. La prescripción es un instituto de derecho sustancial que extingue el derecho o la acción; la caducidad es un instituto de derecho procesal que extingue la instancia y no cancela la pretensión.
}
\cqrow{¿Cuáles son los plazos de caducidad (art. 310 CPCCN)?}{6 meses en primera instancia. 3 meses en segunda o ulterior instancia, juicio sumarísimo, ejecutivo, ejecuciones especiales e incidentes. El plazo en que se opere la prescripción, si es menor. 1 mes para el incidente de caducidad de instancia.
}
\cqrow{¿Cómo se computan los plazos de caducidad?}{Desde el último acto impulsor del procedimiento (no cualquier acto, sino uno que tenga por finalidad hacer avanzar el proceso), a partir del día siguiente. Corren en días corridos (incluyendo inhábiles y feriados). Solo se excluyen los días de ferias judiciales (enero y julio). A diferencia de otros plazos procesales, que se computan en días hábiles.
}
\cqrow{¿Qué efecto tiene el impulso de un litisconsorte (art. 312)?}{El impulso del procedimiento por uno de los litisconsortes beneficia a los demás, porque la instancia es indivisible.
}
\cqrow{¿Cuándo no procede la caducidad (art. 313)?}{En ejecución de sentencia (salvo incidentes ajenos a ella), en procesos sucesorios y voluntarios (salvo incidentes), cuando la demora sea imputable al tribunal, y cuando se haya llamado autos para sentencia (salvo que se ordene prueba de oficio a cargo de las partes).
}
}
\medskip
\cornellblocksum{
\cqrow{¿Cuándo debe pedirse la caducidad y qué ocurre si se consiente un acto posterior (arts. 315 y 316)?}{Debe pedirse antes de consentir cualquier actuación del tribunal o de la otra parte posterior al vencimiento del plazo. Puede declararse de oficio, pero antes de que cualquiera de las partes impulse el procedimiento. Si se consiente un acto impulsor posterior al vencimiento, ya no puede acusarse la caducidad por el plazo anterior.
}
\cqrow{¿La caducidad extingue la acción?}{No. La caducidad en primera instancia no extingue la acción: puede iniciarse un nuevo proceso y se pueden hacer valer las pruebas producidas. Excepción: si el derecho ya está prescripto. La caducidad en instancias ulteriores sí acuerda fuerza de cosa juzgada a la resolución recurrida (la sentencia de primera instancia queda firme).
}
\cqrow{¿Qué relación hay entre la caducidad de la instancia principal y la de los incidentes?}{La caducidad de la instancia principal comprende la reconvención y los incidentes. Pero la caducidad de un incidente no afecta a la instancia principal.
}
}{La caducidad de la instancia (arts. 310--318 CPCCN) opera por inactividad de la parte que tiene la carga de impulsar el proceso dispositivo, en los plazos de 6 meses (primera instancia) o 3 meses (segunda instancia, sumarísimo, ejecutivo e incidentes), computados en días corridos con exclusión de las ferias judiciales. No extingue la acción en primera instancia (puede reiniciarse el proceso si el derecho no está prescripto), pero en instancias ulteriores acuerda fuerza de cosa juzgada a la resolución recurrida. Es el instituto con mayor casuística jurisprudencial, por la diversidad de situaciones en torno a qué constituye o no un acto impulsor del procedimiento.
}

\chapter{Resoluciones judiciales}\label{cap:resoluciones}

\section{Concepto y clasificación de las resoluciones judiciales}

El tema de las resoluciones judiciales es de suma importancia, aunque muchas veces no es atendido debidamente en la práctica por los tribunales, dado que no siempre se sigue a rajatabla lo que dice el código.

\begin{definition}{Resoluciones judiciales}
Cuando el Código Procesal habla de resoluciones judiciales hace referencia a las diversas formas en las que se pronuncia el tribunal. El tribunal puede realizar desde un acto sumamente complejo, como la \textbf{sentencia definitiva}, hasta un acto procesal sumamente sencillo, como las \textbf{providencias simples}.
\end{definition}

El código agrupa las resoluciones judiciales de acuerdo con su complejidad y con el objetivo que tiene cada una. La clasificación es la siguiente:

\begin{table}[H]
\centering
\small
\renewcommand{\arraystretch}{1.25}
\begin{tabularx}{\textwidth}{
    >{\raggedright\arraybackslash}p{2.6cm}
    >{\raggedright\arraybackslash}p{1.8cm}
    >{\raggedright\arraybackslash}p{2.3cm}
    >{\raggedright\arraybackslash}p{3.0cm}
    >{\raggedright\arraybackslash}X
}
\toprule
\textbf{Tipo de resolución} &
\textbf{Complejidad} &
\textbf{Plazo} &
\textbf{Firma} &
\textbf{Característica central} \\
\midrule
Providencias simples & Mínima & 3 días & Juez, secretario, pro-secretario o jefe de despacho & Mero trámite del proceso \\
Resoluciones interlocutorias & Media-alta & 10 días (1\textordfeminine{} inst.); 15 días (2\textordfeminine{} inst.) & Juez (exclusivamente) & Sustanciación / motivación de la decisión \\
Sentencias homologatorias & Media & 10 días & Juez (exclusivamente) & Modos anormales de terminación del proceso \\
Sentencia definitiva de 1\textordfeminine{} instancia & Máxima & Plazos del art.~34 CPCC & Juez (exclusivamente) & Pone fin normal al proceso \\
Sentencia definitiva de 2\textordfeminine{} instancia & Máxima & 15 días & Tres vocales & Acuerdo de cámara \\
\bottomrule
\end{tabularx}
\caption{Clasificación de las resoluciones judiciales.}
\end{table}
% TODO: contenido incompleto o ambiguo en las notas originales (el plazo de la sentencia definitiva de primera instancia se remite al art. 34 CPCC, que en otro pasaje de los apuntes se vincula con las providencias simples)

\begin{important}
Saber ante qué tipo de resolución se está es fundamental para el litigante, porque cada resolución judicial tiene determinados remedios o recursos que pueden interponerse contra ella en un plazo específico.
\end{important}

\begin{example}{La analogía del cañón y el mosquito}
No se puede matar un mosquito con un cañón, ni con una pistola ni con una ametralladora, porque lo más probable es que el mosquito se vaya volando. Con el recurso ocurre lo mismo: no se pueden atacar ciertos actos del proceso con un cañón cuando muchas veces lo que se necesita es una simple aclaratoria o una revocatoria. Y la revocatoria o reposición no puede interponerse contra cualquier tipo de resolución judicial. Para saber qué recursos están disponibles para atacar el acto, primero hay que saber qué tipo de acto se tiene delante.
\end{example}

Todas las resoluciones judiciales son actos procesales y, como tales, tienen características definidas: deben ser redactadas en idioma español, en tinta negra, fechadas y firmadas por la autoridad competente del tribunal.

\section{Providencias simples}

El artículo 34 del CPCC establece los plazos para el dictado de las providencias simples. El plazo del tribunal es de \textbf{tres días} para expedirse en aquellas cuestiones que hacen al mero trámite de la causa y que no requieren ningún tipo de complejidad. A partir del artículo 161 del CPCC se establecen las distintas resoluciones judiciales que existen.

Las \textbf{providencias simples} son aquellas que hacen al normal desarrollo y avance del proceso, y que no agregan nada relevante: solo hacen que el proceso siga avanzando con su trámite. Son ejemplos:

\begin{itemize}
    \item Agregar algún documento o prueba que ha sido ordenada.
    \item Tener por contestada la demanda.
    \item Correr una vista o un traslado.
    \item Ordenar una notificación.
    \item Autorizar a sacar fotocopias o a consultar la causa.
\end{itemize}

El Código autoriza al secretario del tribunal, al pro-secretario e incluso al jefe de despacho a firmar determinadas providencias simples, para disminuir la carga de trabajo del juez. En general, pueden firmar las vinculadas a la agregación o devolución de escritos, la devolución de escritos presentados fuera de término, la agregación de pruebas y documentos, y las vistas que deben darse al ministerio público.

Sin embargo, determinadas providencias simples deben ser suscriptas necesariamente por el juez: por ejemplo, el reconocimiento de la condición de parte, tener por contestada la demanda y aquellas que tienen una relevancia mayor en el proceso.

\begin{example}{El daño de una providencia simple mal resuelta}
Si, en lugar de proveerse <<téngase por contestada la demanda>>, el juez resuelve que la contestación fue presentada fuera de término (por considerarla extemporánea), se provoca un daño irreparable al litigante. Por eso es tan importante saber que se está ante una providencia simple: el remedio disponible es la \textbf{revocatoria}, que a veces es el único que puede solucionar ese problema.
\end{example}
% TODO: contenido incompleto o ambiguo en las notas originales (redacción confusa del ejemplo sobre rechazo / extemporaneidad de la contestación de demanda)

\section{Resoluciones interlocutorias}

Quizás lo más difícil para los estudiantes, e incluso para muchos abogados, es diferenciar una providencia simple de una resolución interlocutoria. La sentencia definitiva es clara, porque es la que pone fin de modo normal al proceso. Lo que a veces no resulta tan claro es cómo distinguir una providencia simple de una interlocutoria cuando esta última se ha simplificado demasiado por usos y costumbres.

El elemento diferenciador clave es la \textbf{sustanciación}.

\begin{important}
Cuando el Código Procesal indica que lo que caracteriza a las resoluciones interlocutorias es la <<sustanciación>>, no se refiere al traslado (uso forense coloquial), sino a la \textbf{sustancia}, al contenido de la resolución: la motivación y las razones que el juez tuvo que desarrollar para tomar su decisión.
\end{important}

Las resoluciones interlocutorias, por su propio nombre, \textbf{deciden algo}: pueden hacer lugar o rechazar una excepción, convocar o rechazar la convocatoria de un tercero, u ordenar o rechazar una medida cautelar. Para llegar a esa decisión, el juez debe motivar y dar razones.

La palabra <<interlocutorio>> proviene del latín \textit{inter}, que significa <<en medio del debate>>: son decisiones que se toman en medio de la discusión principal.

Los plazos para dictar estas resoluciones son:

\begin{itemize}
    \item Primera instancia: 10 días.
    \item Segunda instancia: 15 días.
\end{itemize}

\begin{example}{La medida cautelar confirma que <<sustanciación>> no es <<traslado>>}
Cuando una parte pide una medida cautelar no hay traslado previo. Sin embargo, la decisión que toma el juez no es una providencia simple: es una sentencia interlocutoria. Esto demuestra que el código no se refiere a la sustanciación como traslado, sino como sustancia, materia o motivación del acto. De otra manera, habría que concluir que una medida cautelar es una providencia simple, lo cual es incorrecto.
\end{example}

Hay resoluciones interlocutorias que pueden llegar a poner fin al proceso (como la sentencia que hace lugar a una excepción de prescripción o a una caducidad de la instancia), aunque lo hacen de un modo anormal, no del modo normal, que solo se alcanza mediante la sentencia definitiva.

\section{Sentencias homologatorias}

Las \textbf{sentencias homologatorias} son resoluciones vinculadas a determinados mecanismos que constituyen \textbf{modos anormales de terminación del proceso}: un allanamiento, un acuerdo conciliatorio o un desistimiento. El código prevé que se dicte una resolución que no llega a ser una sentencia definitiva, pero que tampoco es una providencia simple.

Tienen menos complejidad en su sustancia, pero no dejan de tenerla, porque el código le exige al juez que evalúe si el acto de la parte es correcto y si reúne las condiciones necesarias para lograr el desistimiento, el allanamiento o el acuerdo conciliatorio.

\begin{important}
Los acuerdos transaccionales homologados tienen efectos de cosa juzgada.
\end{important}

\section{Sentencia definitiva de primera instancia}

La sentencia definitiva es el acto más complejo que el código prevé dentro de las resoluciones judiciales. El artículo 163 del CPCC contiene nueve incisos que describen las distintas condiciones que debe reunir.

\subsection{Estructura formal}

Tanto la resolución interlocutoria como la sentencia definitiva suelen comenzar con la fecha y la palabra <<Autos>> o <<Autos y Vistos>>, seguida del resumen de los planteos de las partes, el <<Y considerando>> (donde el tribunal evalúa los planteos) y, finalmente, la parte resolutiva (el fallo propiamente dicho).

\subsection{Contenido del artículo 163 del CPCC}

\begin{formula}{Art. 163 CPCC --- Contenido de la sentencia definitiva de primera instancia}
El artículo 163 establece nueve incisos que describen las condiciones de la sentencia definitiva. Entre los más relevantes se destacan:
\begin{itemize}
    \item (Inc. 5\textordmasculine) Mencionar quiénes son las partes, sus nombres y apellidos, y la relación sucinta de las cuestiones planteadas.
    \item Los fundamentos y considerandos (motivación de la decisión).
    \item (Inc. 5\textordmasculine, 2.\textordmasculine{} párrafo) Posibilidad de fundar la decisión en \textbf{presunciones} no establecidas por la ley, que constituirán prueba cuando se funden en hechos reales y probados y cuando, por su número, precisión, gravedad y concordancia, produzcan convicción sobre la naturaleza del juicio, de conformidad con la regla de la sana crítica.
    \item (Inc. 6\textordmasculine) La decisión expresa, positiva y precisa, de conformidad con las pretensiones deducidas en el juicio (\textbf{principio de congruencia}).
    \item (Inc. 6\textordmasculine, \textit{in fine}) Mérito de los hechos constitutivos, modificativos o extintivos producidos durante la sustanciación del juicio y debidamente probados (\textbf{hechos sobrevinientes}).
    \item (Inc. 7\textordmasculine) El plazo que tiene la parte para cumplir la condena.
    \item (Inc. 8\textordmasculine) Establecer las costas y honorarios.
\end{itemize}
\end{formula}
% TODO: contenido incompleto o ambiguo en las notas originales (se mencionan nueve incisos pero solo se detallan algunos; el inc. 5.º aparece dos veces)

\subsection{El deber de motivación}

Es de gran importancia el deber de motivar las decisiones judiciales, tanto las resoluciones interlocutorias como las sentencias definitivas. No se trata solo de explicar por qué se llega a una resolución determinada.

\begin{important}
La motivación de la decisión sirve para que las partes y, de modo muy indirecto, el resto de la población controlen los actos de gobierno. Las sentencias y las decisiones judiciales son actos de gobierno: el Poder Judicial es uno de los poderes del Estado.
\end{important}

Además, la motivación es la base del \textbf{control recursivo}: para interponer una apelación contra una sentencia, el código exige una crítica concreta y razonada de la decisión que se quiere modificar. Esa crítica solo puede construirse sobre los considerandos, que son la parte donde el juez brinda las explicaciones de por qué llegó al resultado.

\subsection{Presunciones e indiciaria}

\begin{definition}{Presunción (art. 163 CPCC)}
La presunción no es un medio de prueba, sino una forma de dirimir la cuestión cuando el hecho principal no está probado pero se cuenta con una serie de hechos que, por su importancia, gravedad y concordancia, hacen que deba presumirse que aquel hecho no probado de forma directa ha sucedido. Es un mecanismo de razonamiento que la ley brinda al juez para poder dictar sentencia.
\end{definition}

El juez civil no puede mantener silencio y no dictar sentencia porque un hecho no esté probado: puede recurrir a los razonamientos de la prueba presuncional o indiciaria.

Sobre la conducta desplegada por las partes durante el proceso como indicio, véase el capítulo~\ref{cap:indiciaria}.

\subsection{La congruencia y los hechos sobrevinientes}

El inciso 6\textordmasculine{} del artículo 163 establece la decisión expresa, positiva y precisa de conformidad con las pretensiones deducidas en el juicio. Esto no es más ni menos que el \textbf{principio dispositivo} y la \textbf{congruencia}.

\begin{important}
El juez, al dictar la sentencia, debe hacerse cargo de las cuestiones que han sido introducidas. Pero esto también le marca un límite: no puede dar más de lo que se le pidió, ni algo distinto de lo que se le pidió. Hacerlo implicaría violar la congruencia y el derecho de defensa de la parte demandada.
\end{important}

\begin{example}{La incapacidad no reclamada}
En un proceso donde no se reclama el resarcimiento de determinada incapacidad, el juez no podría concederla aunque le parezca razonable. Si lo hiciera, violaría la congruencia, al otorgar más allá de lo solicitado por los litigantes, y además violaría el derecho de defensa de la parte demandada, que no pudo negar que tal incapacidad existió.
\end{example}

Respecto de los \textbf{hechos sobrevinientes}: en juicios que duran dos, tres o cuatro años pueden suceder hechos que modifiquen la relación jurídica objeto del proceso, la extingan, o que hagan que la realidad exterior exija que la decisión del juez se adecue a ese cambio.

\begin{example}{El cuadro de Picasso destruido}
Si en un proceso se reclama la restitución de un cuadro pintado por Picasso y durante el juicio ese cuadro deja de existir (porque se destruyó o desapareció), el juez no deberá condenar a su restitución en la sentencia, orden que no podría cumplirse. Deberá reemplazarla por una suma de dinero que compense el valor del bien que debía serle restituido a esa parte.
\end{example}

\subsection{Las costas}

El inciso 8\textordmasculine{} del artículo 163 establece que el juez debe, tanto en la sentencia definitiva como en las interlocutorias, establecer las costas y los honorarios.

\begin{definition}{Costas}
Las costas comprenden los gastos en que ha incurrido la parte para llevar adelante su reclamo: no solo los honorarios de su abogado, sino todo otro gasto en el que haya debido incurrir (sellados, obtención de documentos, etc.).
\end{definition}

La regla que establece el código para la condena en costas es la del \textbf{vencido}: la parte vencida en el proceso es la que debe hacerse cargo del pago de las costas. Si la demanda prospera, en general el vencido es el demandado; si la demanda es rechazada, el vencido es el actor.

Como toda regla, tiene excepciones. Una posibilidad es la \textbf{distribución en el orden causado}: cada parte paga sus propios gastos y los comunes, por mitades.

\begin{important}
Hay que ser muy cuidadoso con los argumentos para eximirse de costas. Si se dijera <<se creyó con derecho a litigar>> como fundamento, ese argumento es endeble: si la parte no se hubiera creído con derecho a litigar, ya estaría incurriendo en la sanción por temeridad y malicia. Esta posibilidad es bastante excepcional, especialmente en los casos en que se discuten derechos patrimoniales. En cambio, en cuestiones de familia, como los alimentos, es muy usual la distribución en el orden causado, para no afectar el monto del alimento que se reclama.
\end{important}

\cornelltitle{resoluciones judiciales}
\cornellblock{
\cqrow{¿Qué son las resoluciones judiciales y cuál es su clasificación?}{Son las diversas formas en que se pronuncia el tribunal. Se clasifican en: providencias simples (mero trámite, 3 días; pueden firmarlas el secretario, el pro-secretario o el jefe de despacho); resoluciones interlocutorias (deciden cuestiones con motivación, 10 días en primera instancia, solo el juez); sentencias homologatorias (modos anormales de terminación: allanamiento, desistimiento, conciliación); y sentencias definitivas de primera y segunda instancia (el acto más complejo; art.~163 CPCC).
}
\cqrow{¿Cómo se diferencia una providencia simple de una resolución interlocutoria?}{Por la sustanciación, entendida como <<sustancia>> (contenido y motivación) y no como <<traslado>>. La providencia simple solo impulsa el trámite; la interlocutoria requiere que el juez funde su decisión porque está resolviendo algo (admisión o rechazo de una excepción, medida cautelar, etc.). La medida cautelar, que no tiene traslado previo y es interlocutoria, confirma que la sustanciación no equivale a traslado.
}
\cqrow{¿Por qué es importante saber ante qué tipo de resolución se está?}{Porque cada resolución judicial tiene determinados remedios o recursos disponibles, en plazos específicos. No se puede <<matar un mosquito con un cañón>>: a veces se necesita una simple aclaratoria o una revocatoria, y la revocatoria no procede contra cualquier resolución.
}
\cqrow{¿Qué establece el art.~163 CPCC sobre la sentencia definitiva?}{Establece nueve incisos. Entre ellos: mención de las partes y relación sucinta de las cuestiones planteadas; fundamentos; posibilidad de usar presunciones (prueba indiciaria) basadas en hechos reales y probados que, por su número, precisión, gravedad y concordancia, produzcan convicción según la sana crítica; congruencia (decisión expresa, positiva y precisa conforme a las pretensiones deducidas); hechos sobrevinientes debidamente probados; plazo para cumplir la condena; y costas y honorarios.
}
}
\medskip
\cornellblocksum{
\cqrow{¿Para qué sirve el deber de motivación?}{Para que las partes y, de modo muy indirecto, el resto de la población controlen los actos de gobierno, ya que el Poder Judicial es uno de los poderes del Estado. Además, es la base del control recursivo: la crítica concreta y razonada que exige la apelación solo puede construirse sobre los considerandos.
}
\cqrow{¿Qué es una presunción según el art.~163 CPCC?}{No es un medio de prueba, sino una forma de dirimir la cuestión cuando el hecho principal no está probado pero hay hechos que, por su importancia, gravedad y concordancia, hacen presumir que ocurrió. Es un mecanismo de razonamiento que permite al juez dictar sentencia; el juez civil no puede dejar de sentenciar porque un hecho no esté probado.
}
\cqrow{¿Qué es el principio de congruencia?}{El juez debe hacerse cargo de las cuestiones introducidas en el proceso, pero no puede dar más de lo pedido ni algo distinto. Superar ese límite viola la congruencia y el derecho de defensa del demandado, que no pudo contradecir lo que no se le reclamaba.
}
\cqrow{¿Qué son los hechos sobrevinientes?}{Son los hechos que suceden durante el transcurso del proceso y que modifican o extinguen la relación jurídica objeto del litigio, o que exigen que la decisión se adecue a la realidad exterior. El juez debe ajustar la sentencia a esa nueva realidad (por ejemplo, si el bien reclamado desaparece, se reemplaza su restitución por una suma de dinero equivalente).
}
\cqrow{¿Cuál es la regla en materia de costas y cuáles son sus excepciones?}{Las costas comprenden todos los gastos en que incurrió la parte para sostener su reclamo. La regla es la del vencido. Una excepción es la distribución en el orden causado (cada parte paga sus gastos y los comunes por mitades), muy usual en cuestiones de familia como los alimentos; eximirse de costas alegando <<se creyó con derecho a litigar>> es un argumento endeble.
}
}{Las resoluciones judiciales (providencias simples, interlocutorias, homologatorias y sentencias definitivas) se distinguen por su complejidad, plazo, firmante y sustancia, y de esa clasificación depende el recurso procedente. Las interlocutorias y las sentencias exigen motivación, que permite el control de las partes y de la población. La sentencia definitiva debe ajustarse al art.~163 CPCC: puede apoyarse en presunciones, debe ser congruente, atender a los hechos sobrevinientes y resolver sobre costas y honorarios, con la regla del vencido como pauta general.
}

\chapter[Segunda instancia y actuación posterior]{Segunda instancia y actuación posterior a la sentencia}\label{cap:segunda}

%--------------------------------------------------------------

\section{Sentencia definitiva de segunda instancia}

En la cámara (segunda instancia) hay diferencias en el modo en que deben elaborarse las resoluciones. Los recursos que den lugar a resoluciones interlocutorias, firmadas por los tres vocales de cada sala, deben estar redactados de modo impersonal (en persona gramatical impersonal), aunque con estructura similar a la de las resoluciones interlocutorias de primera instancia.

La \textbf{sentencia definitiva de segunda instancia} es un acto complejo, elaborado por un \textbf{acuerdo} de tres jueces: hay un \textbf{vocal preopinante} y dos vocales que pueden adherir al voto del primero o marcar sus diferencias (disidencias). El fallo puede salir dos a uno, y no siempre el vocal preopinante tendrá el voto que determine el sentido de la decisión.

\begin{table}[H]
\centering
\small
\renewcommand{\arraystretch}{1.25}
\begin{tabularx}{\textwidth}{
    >{\raggedright\arraybackslash}p{2.4cm}
    >{\raggedright\arraybackslash}X
    >{\raggedright\arraybackslash}X
}
\toprule
\textbf{Aspecto} &
\textbf{Resoluciones interlocutorias en 2\textordfeminine{} instancia} &
\textbf{Sentencia definitiva en 2\textordfeminine{} instancia} \\
\midrule
Redacción & Impersonal & Primera persona del singular (cada vocal emite su voto) \\
Firmantes & Tres vocales de la sala & Tres vocales (acuerdo) \\
Estructura & Similar a las interlocutorias de 1\textordfeminine{} instancia & Similar a la sentencia definitiva de 1\textordfeminine{} instancia, pero en votos individuales \\
El <<acuerdo>> & No aplica como tal & Reunión teórica de los tres jueces para discutir y decidir \\
\bottomrule
\end{tabularx}
\caption{Resoluciones de segunda instancia.}
\end{table}

El acuerdo es, en teoría, una reunión de los tres jueces para discutir las decisiones e intercambiar opiniones. En la práctica, sin embargo, no se lleva a cabo en la mayoría de los casos: lo que se hace en las cámaras es trabajar con la circulación escrita de los expedientes.

\begin{important}
No puede dictarse una sentencia de segunda instancia si uno de los jueces no está presente. Pueden verse dos firmas siempre que el tercer juez esté en uso de licencia concedida o esté vacante la vocalía, pero las cámaras no tienen la misma facultad que la Corte federal de dictar decisiones cuando se alcanza la mayoría aunque los otros jueces no las suscriban.
\end{important}

Una vez dictada la sentencia, sea de primera o de segunda instancia, el código establece que esto concluye la competencia del juez o del tribunal para expedirse sobre el caso. El juez ya no puede modificar el fallo, estén o no notificadas las partes.

\section{Actuación del juez con posterioridad a la sentencia}

El artículo 166 del CPCC establece cuál puede ser la actuación posterior del juez una vez dictada la sentencia. Se prevé la posibilidad de enmendar errores puramente numéricos o errores evidentes (\textbf{errores materiales}), pero en ningún caso el tribunal o el juez puede modificar lo sustancial de las decisiones ya tomadas.

\begin{example}{Fallo de la Corte federal sobre el silencio en costas}
Durante mucho tiempo se sostuvo que el silencio sobre costas en una sentencia (es decir, que el tribunal omitiera pronunciarse sobre ellas) implicaba que habían sido fijadas en el orden causado. Esto fue aclarado en un fallo de la Corte federal, según el cual al silencio no puede otorgársele ningún sentido. Lo que debe hacer el tribunal es expedirse sobre un tema sobre el que aún no se había expedido: la condena en costas.
\end{example}

\section{El recurso de aclaratoria}

El segundo párrafo del artículo 166 habilita, a pedido de parte, el denominado \textbf{recurso de aclaratoria}: dentro de los tres días y sin sustanciación pueden corregirse errores materiales y aclararse conceptos oscuros sin alterar lo sustancial de la decisión. El plazo para solicitarla es de tres días en primera instancia y de cinco días en segunda instancia. Su régimen completo se desarrolla en el capítulo~\ref{cap:remedios}.

\cornelltitle{segunda instancia y actuación posterior}
\cornellblocksum{
\cqrow{¿Cómo funciona la sentencia definitiva de segunda instancia?}{Es elaborada por un acuerdo de tres jueces: un vocal preopinante y dos vocales que adhieren o disienten. Puede salir dos a uno. La redacción es en primera persona del singular, pues cada vocal emite su voto individual. Las resoluciones interlocutorias de cámara, en cambio, se redactan de modo impersonal.
}
\cqrow{¿Puede dictarse sentencia de segunda instancia sin la presencia de los tres jueces?}{No, salvo que el tercer juez esté en uso de licencia concedida o la vocalía esté vacante. Las cámaras no tienen la facultad de la Corte federal de dictar decisiones al alcanzarse la mayoría aunque los otros jueces no las suscriban.
}
\cqrow{¿Qué efecto tiene el dictado de la sentencia sobre la competencia del juez?}{Concluye la competencia del juez o del tribunal para expedirse sobre el caso; ya no puede modificar el fallo, estén o no notificadas las partes.
}
\cqrow{¿Qué puede hacer el juez una vez dictada la sentencia (art.~166 CPCC)?}{Solo puede enmendar errores puramente numéricos o errores materiales evidentes, y aclarar conceptos oscuros, sin alterar lo sustancial de lo decidido.
}
\cqrow{¿Qué es el recurso de aclaratoria y en qué plazos procede?}{Es el pedido de la parte para que se corrijan errores materiales y se aclaren conceptos oscuros, dentro de los tres días y sin sustanciación según el art.~166. El plazo para solicitarla es de tres días en primera instancia y cinco días en segunda instancia. No puede modificar lo sustancial de la decisión.
}
\cqrow{¿Qué sentido tiene el silencio sobre costas, según el fallo de la Corte federal?}{Al silencio no puede otorgársele ningún sentido (no implica orden causado). El tribunal debe expedirse sobre un tema sobre el que aún no se había expedido: la condena en costas.
}
}{La sentencia definitiva de segunda instancia es el resultado de un acuerdo de tres jueces (vocal preopinante, adhesiones y disidencias), redactada en primera persona del singular. Una vez dictada la sentencia, concluye la competencia del juez o tribunal: el art.~166 CPCC solo permite corregir errores materiales y aclarar conceptos oscuros, sin alterar lo sustancial, mediante la aclaratoria (3 días en primera instancia, 5 en segunda).
}

\chapter{Cosa juzgada}\label{cap:cosajuzgada}

%--------------------------------------------------------------

\section{Concepto e inmutabilidad}

El concepto de cosa juzgada ya fue visto al estudiar las excepciones procesales en el artículo 347; aquí se profundizan algunos aspectos.

\begin{definition}{Cosa juzgada}
Implica que, contra una decisión de un tribunal de primera o de segunda instancia, no hay más recursos disponibles para interponer, o bien el plazo para apelarla ha transcurrido sin que las partes hayan deducido ninguno. En esas condiciones, la decisión adquiere la \textbf{autoridad de cosa juzgada}.
\end{definition}

\begin{important}
La cosa juzgada significa que dentro del patrimonio de la persona ingresa un activo: el derecho a hacer cumplir lo que emana de esa decisión judicial. La realidad previa deja de existir, y lo único que queda es lo que emana de la sentencia.
\end{important}

El abogado debe controlar muy celosamente lo que surge de la decisión, especialmente lo que el tribunal escribe o asienta después de la palabra <<fallo>>, porque eso es lo que en el futuro podrá exigirse que se cumpla con autoridad de cosa juzgada.

Cuando una sentencia de segunda instancia modifica parcialmente la decisión de primera instancia, la cosa juzgada emana de la \textbf{sumatoria de los dos fallos}: las cuestiones que quedaron firmes por no haber sido apeladas en la sentencia de primera instancia se suman a las que emanan de la sentencia firme de segunda instancia.

\section{Cosa juzgada formal y cosa juzgada material}

No todos los procesos tienen la misma característica en cuanto a la calidad de la cosa juzgada. La doctrina diferencia la \textbf{cosa juzgada formal} de la \textbf{cosa juzgada material}.

\begin{table}[H]
\centering
\small
\renewcommand{\arraystretch}{1.25}
\begin{tabularx}{\textwidth}{
    >{\raggedright\arraybackslash}p{2.6cm}
    >{\raggedright\arraybackslash}p{3.0cm}
    >{\raggedright\arraybackslash}X
    >{\raggedright\arraybackslash}X
}
\toprule
\textbf{Tipo} &
\textbf{Proceso típico} &
\textbf{Qué cuestiones abarca} &
\textbf{Posibilidad de juicio posterior} \\
\midrule
Cosa juzgada material & Proceso ordinario (conocimiento pleno) & Todas las cuestiones debatidas extensamente & No (el fondo fue juzgado definitivamente) \\
Cosa juzgada formal & Juicio ejecutivo / caducidad de instancia & Solo las cuestiones discutibles en ese proceso & Sí, en juicio ordinario posterior (sobre lo no juzgado) \\
\bottomrule
\end{tabularx}
\caption{Cosa juzgada material y formal.}
\end{table}

\begin{example}{El juicio ejecutivo}
En un juicio ejecutivo hay cosa juzgada material si se discutió el pago, porque el pago es una excepción que podía oponerse en ese proceso. En cambio, no habrá cosa juzgada material en relación, por ejemplo, con algún vicio de la voluntad que se hubiera tenido al suscribir el título ejecutivo, porque esa cuestión no podía discutirse en el juicio ejecutivo. Respecto de ese aspecto hay cosa juzgada formal, y las cuestiones cuya discusión no puede darse en el juicio ejecutivo podrán discutirse en el juicio ordinario posterior.
\end{example}

\section{Límites subjetivos de la cosa juzgada}

El inciso 6\textordmasculine{} del artículo 347 del CPCC establece la \textbf{excepción de cosa juzgada}. Lo que ataca la excepción es el derecho de la persona que reclama: ese derecho ya fue discutido y decidido en una sentencia previa que está firme, contra la cual ya no hay más recursos para interponer.
% TODO: contenido incompleto o ambiguo en las notas originales (el apunte dice «dividido en una sentencia previa»; se transcribe como «decidido» por el sentido del pasaje)

Los \textbf{límites subjetivos} indican que la cosa juzgada abarca a las personas que pudieron discutir en el marco del proceso previo: actor, demandado, terceros intervinientes y terceros legitimados sobre los cuales se ha discutido una relación sustancial. No puede afectar a sujetos que no han participado de ese proceso, porque no han tenido la posibilidad de defenderse, ofrecer prueba y contradecir.

Esto se vincula con el \textbf{litisconsorcio}:

\begin{important}
En tanto se trate de un litisconsorcio facultativo, nada impide que el juicio llevado adelante produzca una cosa juzgada material que afecte solo a las partes que han participado, sin relevancia respecto de otros sujetos que podrían haber sido litisconsortes facultativos y que no fueron citados.
\end{important}

\section{Límites objetivos de la cosa juzgada}

Los límites objetivos son más difusos. Nadie duda de que la cuestión debatida en el primer proceso es la que ha sido juzgada, pero la lectura del inciso 6\textordmasculine{} del artículo 347 muestra que el alcance es mucho más expansivo de lo que una primera aproximación indicaría.

\begin{formula}{Art. 347, inc. 6\textordmasculine, CPCC --- Límites objetivos}
Para que sea procedente la excepción de cosa juzgada, el examen integral de las dos contiendas debe demostrar que se trata del \textbf{mismo asunto} sometido a decisión judicial. Pero, además, hay cosa juzgada cuando, por existir \textbf{continencia, conexidad, accesoriedad o subsidiariedad}, la sentencia firme ya haya resuelto lo que constituye la materia o la pretensión deducida en el nuevo juicio que se promueve.
\end{formula}

Esto significa que los límites objetivos de la cosa juzgada no solo comprenden lo debatido en el primer proceso, sino también las cuestiones de continencia, conexidad, accesoriedad y subsidiariedad que podrían haberse planteado en él y no se plantearon.

\begin{example}{El accidente de autos y el lucro cesante}
Si en un juicio por un accidente de autos se reclama solo la reparación del automóvil y se resulta ganador, cabe preguntarse si podría iniciarse un segundo proceso reclamando el lucro cesante (porque el auto era utilizado como taxi o remís). A partir del artículo 347, inc.~6\textordmasculine, se podría concluir que existió lo que se denomina \textbf{cosa juzgada implícita}: el lucro cesante era un rubro que podía pedirse en el primer proceso y no se hizo, por lo que precluyó también la posibilidad de hacerlo con posterioridad, al ser una cuestión accesoria o conexa con la discutida en el primer litigio.
\end{example}

Es importante conocer esto porque la cosa juzgada \textbf{puede ser declarada de oficio por el juez}. Puede suceder que un juez no la advierta y que el demandado tampoco oponga la excepción. En el proceso le va mejor al que más sabe y tiene más herramientas para plantear la discusión de estos supuestos.

\cornelltitle{cosa juzgada}
\cornellblock{
\cqrow{¿Qué es la cosa juzgada?}{Es la autoridad que adquiere una decisión judicial cuando contra ella no hay más recursos disponibles o el plazo para interponerlos ha vencido sin que las partes los dedujeran. Produce dos efectos: (1) ingresa al patrimonio del vencedor el derecho a exigir el cumplimiento; (2) la realidad previa deja de existir y solo queda lo que emana de la sentencia (novación procesal). Los acuerdos transaccionales homologados también tienen efectos de cosa juzgada.
}
\cqrow{¿De dónde emana la cosa juzgada cuando la sentencia de segunda instancia modifica parcialmente la de primera?}{De la sumatoria de ambos fallos: las cuestiones que quedaron firmes por no haber sido apeladas en primera instancia se suman a las que emanan de la sentencia firme de segunda instancia. Por eso el abogado debe controlar lo asentado después de la palabra <<fallo>>.
}
\cqrow{¿Cuál es la diferencia entre cosa juzgada formal y material?}{La cosa juzgada material corresponde al proceso ordinario (conocimiento pleno): abarca todas las cuestiones debatidas y no permite un juicio posterior sobre el fondo. La cosa juzgada formal corresponde, por ejemplo, al juicio ejecutivo o a la caducidad de instancia: abarca solo las cuestiones discutibles en ese proceso y permite un juicio ordinario posterior sobre lo no juzgado (por ejemplo, vicios de la voluntad en el juicio ejecutivo).
}
\cqrow{¿Cuáles son los límites subjetivos de la cosa juzgada?}{Solo alcanza a quienes fueron partes del proceso (actor, demandado, terceros intervinientes, terceros legitimados). No puede afectar a quienes no tuvieron la posibilidad de defenderse, ofrecer prueba y contradecir. Se vincula con el litisconsorcio facultativo: la cosa juzgada material afecta solo a las partes que participaron.
}
\cqrow{¿Cuáles son los límites objetivos de la cosa juzgada (art.~347, inc.~6\textordmasculine)?}{No solo las cuestiones debatidas, sino también las conexas, accesorias, subsidiarias o por continencia que pudieron plantearse y no se plantearon. Puede haber cosa juzgada <<implícita>> respecto de rubros accesorios no reclamados (por ejemplo, el lucro cesante tras un juicio por la reparación del automóvil).
}
}
\medskip
\cornellblocksum{
\cqrow{¿Por qué importa que la cosa juzgada pueda declararse de oficio?}{Porque puede ocurrir que el juez no la advierta y que el demandado no oponga la excepción; en el proceso le va mejor a quien más sabe y tiene más herramientas para plantear estos supuestos.
}
}{Una decisión adquiere autoridad de cosa juzgada cuando ya no hay recursos o ha vencido el plazo para interponerlos; el derecho a su cumplimiento ingresa al patrimonio del vencedor. Puede ser material (proceso ordinario, no permite un juicio posterior sobre el fondo) o formal (limitada a lo discutible en ese proceso, como en el juicio ejecutivo). Sus límites subjetivos alcanzan solo a quienes fueron partes; los objetivos, definidos en el art.~347, inc.~6\textordmasculine, se extienden a cuestiones conexas, accesorias, subsidiarias o por continencia que pudieron plantearse.
}

\chapter[Sentencia civil y sentencia penal]{Relación entre sentencia civil y sentencia penal}\label{cap:civilpenal}

%--------------------------------------------------------------

\section{Prelación de la acción penal sobre la civil}

Vinculado con la cosa juzgada está lo que usualmente se denomina la incidencia de la acción penal sobre la acción civil, regulada en el Código Civil y Comercial (CCyC) a partir del artículo 1774.

\begin{formula}{Art. 1774 y ss. CCyC --- Prelación de la acción penal}
Si la acción penal precede a la acción civil o es intentada durante su curso, el dictado de la sentencia definitiva debe suspenderse en el proceso civil hasta que haya conclusión del proceso penal.

Excepciones a la suspensión:
\begin{itemize}
    \item Si median causas de extinción de la acción penal.
    \item Si la dilación del procedimiento penal provoca en los hechos una frustración efectiva del derecho a ser indemnizado en el juicio civil.
    \item Si la acción civil por reparación del daño está fundada en un factor objetivo de responsabilidad (gran cambio del CCyC: en estos casos nada impide que la sentencia civil se dicte aun cuando no se haya dictado sentencia penal).
\end{itemize}
\end{formula}

La razón de ser de esta coordinación es evitar lo que en general se encuentra en los libros como el <<escándalo jurídico>>, que en definitiva es un escándalo lógico.

\begin{important}
El <<escándalo jurídico>> (o lógico) se produce cuando una sentencia en un fuero dice que el autor del hecho fue una persona y otra sentencia dice que esa persona no fue la autora. La coordinación de los dos ordenamientos busca evitar precisamente esta colisión de decisiones.
\end{important}

\section{Influencia de la sentencia penal sobre la civil}

El Código Civil y Comercial dispone que la inexistencia del hecho, la autoría y la ausencia de dolo penal de responsabilidad determinadas en la sentencia penal no pueden desconocerse en la sentencia civil.

\begin{important}
Si la sentencia penal dice que el sindicado como responsable penal no participó, que no es autor del hecho o que el hecho no ocurrió, esto no podrá ser desconocido en la sentencia civil, que deberá fundarse en consecuencia.
\end{important}

Es posible que una absolución penal por ausencia de delito (porque el accionar no está tipificado en el Código Penal) dé, sin embargo, lugar a la reparación civil. En esos casos no habrá ninguna influencia de las consecuencias penales sobre las civiles.

\section{Sentencia penal posterior a la civil}

Puede ocurrir que, en la acción civil, se demande con base en un factor de atribución objetivo, se permita que la acción civil continúe y se dicte una sentencia civil a la que con posterioridad se agregue una sentencia en sede penal.

\begin{formula}{Art. 1780 CCyC --- Sentencia penal posterior a la sentencia civil}
La sentencia penal posterior a la sentencia civil no produce ningún efecto sobre ella, excepto en el caso de revisión. La revisión procede cuando así se solicite y cuando la sentencia civil asigna alcances de cosa juzgada a cuestiones resueltas por la sentencia penal y esta es revisada respecto de esas cuestiones.
\end{formula}

\begin{example}{Absolución penal posterior por inexistencia del hecho}
Si se dictó una sentencia civil condenando al demandado sobre la base de un factor de atribución objetivo (art.~1774, inc.~c, CCyC) y con posterioridad, en la acción penal, el demandado es absuelto por inexistencia del hecho en que se funda esa condena civil, o por no ser su autor, procede la revisión de la condena civil. Aunque la sentencia penal llegó con posterioridad, afirma que el demandado condenado en sede civil no fue autor del hecho, de modo que de ninguna manera podría quedar condenado.
\end{example}

\cornelltitle{sentencia civil y sentencia penal}
\cornellblocksum{
\cqrow{¿En qué consiste la prelación de la acción penal sobre la civil?}{Según el art.~1774 CCyC, si la acción penal precede a la civil o se intenta durante su curso, la sentencia civil debe suspenderse hasta la conclusión del proceso penal, salvo: (a) extinción de la acción penal; (b) frustración efectiva del derecho a ser indemnizado por la dilación del proceso penal; (c) factor objetivo de responsabilidad (gran innovación del CCyC: en ese caso la sentencia civil puede dictarse sin esperar la penal).
}
\cqrow{¿Qué es el <<escándalo jurídico>> y cómo se evita?}{Es un escándalo lógico: ocurre cuando una sentencia dice que el autor del hecho fue una persona y otra sentencia dice que esa persona no fue la autora. La coordinación de los ordenamientos civil y penal busca evitar esa colisión de decisiones.
}
\cqrow{¿Qué efectos produce la sentencia penal anterior sobre la civil?}{La sentencia penal vincula a la civil: si determina la inexistencia del hecho, la falta de autoría o la ausencia de dolo, esas conclusiones no pueden desconocerse en la sentencia civil. Una absolución penal por ausencia de delito (accionar no tipificado) puede, en cambio, dar lugar a la reparación civil sin influencia de las consecuencias penales.
}
\cqrow{¿Qué efectos produce la sentencia penal posterior a la civil (art.~1780 CCyC)?}{En principio, ninguno sobre la sentencia civil ya dictada, salvo en caso de revisión: procede cuando se solicita y la sentencia civil asignó alcances de cosa juzgada a cuestiones resueltas por la sentencia penal que luego son revisadas (por ejemplo, absolución penal por inexistencia del hecho en que se fundó la condena civil).
}
}{La acción penal tiene prelación sobre la civil: la sentencia civil se suspende hasta la conclusión del proceso penal, salvo extinción de la acción penal, frustración del derecho a ser indemnizado por dilación, o factor objetivo de responsabilidad. La sentencia penal anterior no puede desconocerse en lo relativo a la inexistencia del hecho, la autoría y la ausencia de dolo; la posterior no afecta a la civil, salvo revisión. Todo ello busca evitar el escándalo lógico de fallos contradictorios.
}

\section{Síntesis de la sentencia, la cosa juzgada y su coordinación con el proceso penal}
Las resoluciones judiciales constituyen el lenguaje formal a través del cual el tribunal se comunica con las partes y hace avanzar el proceso. Su clasificación (providencias simples, resoluciones interlocutorias, sentencias homologatorias y sentencias definitivas) no es meramente académica: determina qué recursos pueden interponerse, en qué plazo y ante quién. El deber de motivación no solo garantiza el control de las partes sino también el control democrático de los actos del Poder Judicial.

La sentencia definitiva, regulada en el art.~163 CPCC, debe ser expresa, positiva, precisa y congruente con las pretensiones deducidas, pudiendo el juez valerse de presunciones indiciarias y de la conducta procesal de las partes para forjar su convicción.

Una vez firme, la sentencia adquiere autoridad de cosa juzgada, que puede ser material (definitiva) o formal (limitada a lo discutible en ese proceso), y cuyos límites, objetivos y subjetivos, son más amplios de lo que una primera lectura sugiere, pues alcanzan cuestiones conexas, accesorias y subsidiarias (art.~347, inc.~6\textordmasculine).

La coexistencia del proceso civil y el penal exige una coordinación normativa que, con el Código Civil y Comercial, admite la continuación de la acción civil frente a factores objetivos de responsabilidad sin necesidad de esperar la conclusión del proceso penal, buscando en todos los casos evitar el <<escándalo lógico>> de fallos contradictorios.

\partintro{La etapa de impugnación sigue cronológicamente a la etapa de conclusión del proceso, a las resoluciones procesales y a las sentencias. Esta parte desarrolla la \textbf{teoría general de la impugnación}, es decir, el sistema mediante el cual el ordenamiento procesal civil argentino permite atacar las decisiones judiciales erróneas o los actos procesales viciados. Comienza con los fundamentos y las clasificaciones de los actos de impugnación y de los recursos, y con los requisitos que son comunes a todos los recursos. Continúa con la garantía de la doble instancia y los límites del tribunal de alzada, con los remedios que se resuelven ante el mismo juez (aclaratoria y reposición o revocatoria) y con el incidente de nulidad. Finalmente se estudia en detalle el recurso de apelación: su concepto, su concesión, su trámite, sus efectos y las reglas que lo rigen.}
\part{Impugnación y recursos}

\chapter[Teoría general de la impugnación]{Teoría general de la impugnación y clasificación de los recursos}\label{cap:impugnacion}

\section{Ubicación del tema y relación entre impugnación y recurso}

La etapa de impugnación sigue cronológicamente a la etapa de conclusión del proceso, las resoluciones procesales y las sentencias.

Los actos de impugnación y los recursos guardan una relación de \textbf{género y especie}: los \textbf{recursos} son un tipo de acto de impugnación. Existen otros tipos de actos de impugnación, que se desarrollan al analizar las clasificaciones posibles.

La voz \textbf{impugnaci\'on} es un t\'ermino mucho m\'as amplio que el de \textbf{recurso}. Seg\'un Hitter, la impugnaci\'on corresponde m\'as a la Teor\'ia General del Derecho, mientras que el t\'ermino recurso es propio del derecho procesal. La relaci\'on entre ambos es de \textbf{g\'enero a especie}: la impugnaci\'on es el g\'enero y el recurso es una especie de ese g\'enero.

Impugnar quiere decir \textbf{atacar}. Pueden existir actos de impugnaci\'on fuera del proceso judicial y, dentro del proceso, existen distintos tipos de impugnaciones que no siempre son recursos.

\begin{example}{Impugnaci\'on de dictamen pericial}
En la etapa probatoria, si se presenta un dictamen pericial m\'edico y se corre traslado a las partes, \'estas pueden atacarlo mediante un acto procesal escrito de impugnaci\'on, dando los fundamentos de por qu\'e se lo ataca. Ese acto es claramente una impugnaci\'on, pero no es un recurso.
\end{example}

\begin{example}{Impugnaci\'on de liquidaci\'on en ejecuci\'on}
En la etapa de ejecuci\'on del proceso, si se impugna una liquidaci\'on del capital de condena m\'as intereses por entender que los intereses est\'an mal calculados, ese acto tambi\'en es una impugnaci\'on, pero no un recurso.
\end{example}

\begin{important}
S\'olo se est\'a ante un \textbf{recurso} cuando lo que se impugna es una \textbf{resoluci\'on judicial}. Los recursos son mecanismos de impugnaci\'on \textbf{exclusivamente de resoluciones judiciales}: providencias simples, resoluciones interlocutorias y sentencias definitivas (reguladas a partir del art.~160 del C\'odigo Procesal).
\end{important}

\section{Etimología de la palabra «impugnación»}

La palabra \textbf{impugnación} proviene del latín \textit{pugnus}, que significa «puño». De \textit{pugnus} derivaron varias palabras que aluden a la lucha de puños: el pugilato (enfrentamiento a puñetazos), el púgil (el boxeador), la pelea. La noción que subyace a la impugnación es, por lo tanto, la de \textbf{combate, lucha o ataque}.

Aplicado al proceso judicial, cuando se habla de una teoría de la impugnación o de actos de impugnación se hace referencia a todo lo que tiene que ver con atacar o combatir un procedimiento viciado o una decisión equivocada desde la perspectiva de quien impugna. La impugnación es, entonces, todo mecanismo procesal que ataca o combate cualquier acto procesal o procedimiento que lesione el interés de cualquiera de las partes o de un tercero.

\section{Fundamento de la impugnación: el error humano}

El ordenamiento procesal contempla mecanismos de impugnación porque \textbf{existe el error}. Los jueces pueden equivocarse y cometer errores. Del mismo modo, las partes, en los actos procesales que impulsan, pueden cometer irregularidades y realizarlos sin cumplir adecuadamente con las formas. Estos errores de procedimiento o de interpretación de la ley son inevitables, porque están dentro de la naturaleza humana. Los actos de impugnación constituyen la alternativa que busca mitigar los defectos y deficiencias que pueda haber en un proceso o en una resolución judicial.

\textbf{Calamandrei} sosten\'ia que el fundamento de los recursos es, ni m\'as ni menos, el \textbf{error humano}: los jueces son humanos y pueden equivocarse; si se equivocan, debe existir un sistema para subsanar ese error y evitar que se consolide una injusticia.

\begin{note}
\textbf{San Agustín y Cicerón.} La frase «errar es humano» se atribuye habitualmente a San Agustín, aunque tiene antecedentes en Cicerón. Lo grave (lo «diabólico», según San Agustín) es persistir en el error. Precisamente porque el hombre reconoce que se equivoca y puede dar marcha atrás tienen sentido los mecanismos recursivos y toda la teoría de la impugnación.
% TODO: contenido incompleto o ambiguo en las notas originales (la transcripción original menciona «teoría de la evolución», corregida en el propio apunte entre corchetes como «impugnación»)
\end{note}

\begin{note}
\textbf{Francesco Carnelutti.} El procesalista italiano Francesco Carnelutti decía que la amenaza del error pende como \textbf{espada de Damocles} sobre el proceso. Según la leyenda, Damocles tenía sobre su cabeza una espada afilada que pendía de una crin de caballo y que en cualquier momento podía ceder. El error es así el fundamento clásico de toda la teoría de la impugnación.
\end{note}

\textbf{Couture} sosten\'ia que el t\'ermino recurso se vincula con \textit{``volver al camino de partida''}: al recurrir, el tribunal superior vuelve a recorrer un camino ya transitado en la instancia anterior, en la que fue dictada la resoluci\'on impugnada, para revisarlo y, en funci\'on de esa revisi\'on, confirmar la resoluci\'on o modificarla total o parcialmente.

\section{Definición de actos procesales de impugnación}

\begin{definition}{Actos procesales de impugnación (Palacio)}
Los actos procesales de impugnación son aquellos que están dirigidos \textbf{directa e inmediatamente} a provocar la \textbf{modificación o sustitución, total o parcial, de una resolución judicial}, en el mismo proceso en el que ella fue dictada.
\end{definition}

\begin{example}{El sistema de correcciones como un partido mal arbitrado}
Imaginemos un partido de fútbol mal arbitrado. El árbitro comete un error (el juez se equivoca). Los jugadores pueden pedir explicaciones (aclaratoria), pedir que el árbitro reconsidere su decisión en el momento (revocatoria), apelar ante el tribunal deportivo superior (apelación) o, si hubo una grave irregularidad de procedimiento, impugnar la validez del partido entero (incidente de nulidad). Todo ese sistema de correcciones forma la teoría de la impugnación.
\end{example}

\cornelltitle{concepto y fundamento de la impugnación}
\cornellblocksum{
\cqrow{¿Qué relación existe entre los actos de impugnación y los recursos?}{Los actos de impugnación son el género y los recursos son una especie dentro de ese género. Existen otros actos de impugnación que no son recursos.
}
\cqrow{¿Qué significa la palabra «impugnación» según su etimología?}{Proviene del latín \textit{pugnus} (puño) y alude a la lucha, al combate y al ataque. En el proceso, se trata de todo mecanismo que ataca o combate un acto procesal o procedimiento que lesione el interés de las partes o de un tercero.
}
\cqrow{¿Por qué existen los mecanismos de impugnación?}{Porque existe el error: los jueces y las partes pueden equivocarse, y esos errores de procedimiento o de interpretación de la ley son inevitables. Carnelutti decía que la amenaza del error pende como espada de Damocles en el proceso. La impugnación busca mitigar esas deficiencias.
}
\cqrow{¿Cómo se definen los actos procesales de impugnación?}{Son los dirigidos directa e inmediatamente a provocar la modificación o sustitución, total o parcial, de una resolución judicial, en el mismo proceso en el que ella fue dictada (Palacio).
}
}{La impugnación es el sistema que permite atacar actos o resoluciones viciados o equivocados. Su fundamento es la falibilidad humana. El género es el acto de impugnación y los recursos son una de sus especies.
}

\section{Clasificaciones de los actos de impugnación y de los recursos}

Con carácter metodológico general, cabe advertir que se trata de \textbf{clasificaciones posibles}: no hay clasificaciones verdaderas o falsas, buenas o malas. Toda clasificación es una manera de comprender, abarcar y categorizar distintos conceptos. En tanto ayude a comprenderlos mejor, la clasificación es buena; cuando dificulta la comprensión de lo que se quiere categorizar, deja de serlo.


\subsection{Primera clasificación: incidente de nulidad, recursos y otros medios}

La primera clasificación distingue dos grandes grupos, a los que se agregan otros medios menos frecuentes.

\begin{table}[H]
\centering
\small
\begin{tabularx}{\textwidth}{L{0.18\textwidth} >{\raggedright\arraybackslash}X >{\raggedright\arraybackslash}X}
\toprule
 & \textbf{Incidente de nulidad} & \textbf{Recursos} \\
\midrule
\textbf{Normativa y finalidad} &
Arts. 169--174 CPCCN. Corrige errores o vicios en la tramitación del proceso (vicios \textit{in procedendo}). &
Apuntan a la corrección de las decisiones impugnadas. \\
\midrule
\textbf{Ejemplo / tipos} &
Ejemplo típico: notificación defectuosa del traslado de la demanda en un domicilio incorrecto. &
\begin{itemize}[leftmargin=*,nosep]
\item Aclaratoria
\item Revocatoria o reposición
\item Apelación
\item Queja por apelación denegada
\item Recurso extraordinario federal
\item Apelación ordinaria ante la Corte Suprema
\item Inaplicabilidad de ley
\end{itemize} \\
\bottomrule
\end{tabularx}
\end{table}

En el proceso ordinario habitual, los recursos más frecuentes son la \textbf{aclaratoria}, la \textbf{revocatoria} y la \textbf{apelación}. La \textbf{queja} surge como consecuencia de la denegación de la apelación.

\subsection{Otros actos de impugnación (menos frecuentes)}

Si bien no son tan frecuentes, los siguientes mecanismos también integran el sistema de impugnación.

\subsubsection{Replanteo de prueba ante la alzada}

Cuando en el juicio ordinario el juez de primera instancia deniega ciertos medios de prueba, el litigante puede impugnar esa denegatoria en la etapa de revisión de la sentencia. Si la impugnación es adecuada, puede replantear la prueba y producirla ante el tribunal de alzada.

\subsubsection{Acción autónoma de revisión de la cosa juzgada (cosa juzgada írrita)}

Una vez que la sentencia pasa en autoridad de cosa juzgada, la jurisprudencia ha abierto, de modo \textbf{excepcionalísimo} y muy restrictivo, la posibilidad de revisar lo decidido por sentencia firme cuando existen circunstancias excepcionales.

\begin{example}{Juicios de filiación y prueba de ADN}
En otro tiempo no existía la prueba de ADN, por lo que las sentencias dictadas en juicios de filiación no podían valerse de ese medio de prueba. Con el advenimiento de la prueba biológica de altísima precisión que es el ADN, que permite despejar toda duda sobre la cuestión, se abrió la posibilidad de revisar esas sentencias firmes, aun cuando hubieran pasado en autoridad de cosa juzgada. También se aplica cuando, por maniobras fraudulentas, se ocultaron pruebas esenciales durante el proceso.
\end{example}

\subsubsection{Nulidad del laudo arbitral}

Los arts. 767--771 CPCCN prevén las hipótesis de nulidad del laudo arbitral o laudo de amigables componedores como mecanismo de impugnación.
% TODO: contenido incompleto o ambiguo en las notas originales (el docente menciona el art. 771 y, al mismo tiempo, los arts. 767-771)

\subsubsection{Juicio ordinario posterior al juicio ejecutivo}

Permite revisar, en un proceso de conocimiento pleno, una sentencia dictada en juicio ejecutivo, en el que el acotado marco procedimental no permitió ejercer en plenitud las defensas disponibles.

\subsection{Segunda clasificación: remedios y recursos}

La doctrina distingue también entre \textbf{remedios} y \textbf{recursos} en sentido estricto.

\begin{table}[H]
\centering
\small
\begin{tabularx}{\textwidth}{L{0.18\textwidth} >{\raggedright\arraybackslash}X >{\raggedright\arraybackslash}X}
\toprule
 & \textbf{Remedios} & \textbf{Recursos (\textit{stricto sensu})} \\
\midrule
\textbf{Orientación} &
Reparar errores procesales (vicios \textit{in procedendo}). &
Apuntan a que un tribunal jerárquicamente superior reexamine el acto impugnado. \\
\midrule
\textbf{Quién resuelve} &
\textbf{El mismo juez que dictó el acto.} &
Un tribunal jerárquicamente superior. \\
\midrule
\textbf{Vicios que se atacan} &
Vicios \textit{in procedendo}. &
Vicios \textit{in iudicando}: errores en la decisión judicial (interpretación de la ley o de los hechos). \\
\midrule
\textbf{Incluyen} &
Incidente de nulidad, acción autónoma de revisión de cosa juzgada, juicio ordinario posterior, aclaratoria* y revocatoria*. &
Apelación, queja, recurso extraordinario federal, apelación ordinaria ante la Corte, inaplicabilidad de ley. \\
\bottomrule
\end{tabularx}
\end{table}

(*) La aclaratoria y la revocatoria presentan aspectos que podrían involucrarlas dentro de la categoría de remedios, aunque el Código las trate y denomine como «recursos». Técnicamente, en doctrina, sería más adecuado clasificarlas dentro de los remedios.

\subsection{Clasificaci\'on general de los recursos}

Es importante conocer el cuadro completo de los recursos.

\begin{table}[H]
\centering
\begin{tabularx}{\textwidth}{
    >{\raggedright\arraybackslash}p{0.27\textwidth}
    >{\raggedright\arraybackslash}p{0.22\textwidth}
    >{\raggedright\arraybackslash}X
}
\toprule
\textbf{Tipo} & \textbf{Categor\'ia} & \textbf{Ejemplos} \\
\midrule
\textbf{Horizontales} & Remedios & Aclaratoria, revocatoria \\
\addlinespace
\textbf{Verticales ordinarios} & Apelaci\'on & Recurso de apelaci\'on (proceso ordinario) \\
\addlinespace
\textbf{Verticales extraordinarios (comunes)} & Regulados por ley & Inaplicabilidad de ley (fallos plenarios), Recurso Extraordinario Federal (Ley 48) \\
\addlinespace
\textbf{Verticales extraordinarios (excepcionales)} & Creaci\'on pretoriana & Sentencia arbitraria, gravedad institucional, \textit{per saltum} (ya regulado) \\
\addlinespace
\textbf{Especiales} & Queja & Cuando se deniega el recurso de apelaci\'on o el extraordinario \\
\addlinespace
\textbf{Parajudiciales} & Fuera del proceso & Recursos administrativos, recursos arbitrales \\
\bottomrule
\end{tabularx}
\caption{Clasificaci\'on general de los recursos.}
\end{table}

\subsubsection{Recursos horizontales: aclaratoria y revocatoria}

Son horizontales porque se interponen, tramitan y resuelven \textbf{ante el mismo juez} que dict\'o la resoluci\'on que se impugna: se interpone, se funda y es ese mismo juez quien la resuelve. Se dice que, m\'as que recursos, son \textbf{remedios}, porque no existe un nuevo tribunal que vuelva a recorrer el camino de una instancia anterior.

Una caracter\'istica fundamental que los diferencia de la apelaci\'on es que, en los horizontales, el recurso se \textbf{interpone y se funda en el mismo acto}. Esto no sucede en el recurso de apelaci\'on, donde existen dos momentos diferenciados.

\subsubsection{Recursos verticales ordinarios: la apelaci\'on}

La apelaci\'on es un recurso vertical porque se interpone ante el juez de primera instancia (el que dict\'o la resoluci\'on) para que sea la \textbf{c\'amara}, que se encuentra por encima, la que resuelva. Tiene \textbf{dos momentos diferenciados}: la interposici\'on y la fundamentaci\'on, que se desarrollan m\'as adelante.

\cornelltitle{clasificación de los actos de impugnación}
\cornellblocksum{
\cqrow{¿Qué grupos de actos de impugnación distingue la primera clasificación?}{El incidente de nulidad (arts. 169--174 CPCCN, que corrige vicios en la tramitación del proceso) y los recursos (aclaratoria, revocatoria o reposición, apelación, queja por apelación denegada, recurso extraordinario federal, apelación ordinaria ante la Corte Suprema e inaplicabilidad de ley). Se agregan otros medios menos frecuentes.
}
\cqrow{¿Cuáles son los recursos más frecuentes en el proceso ordinario?}{La aclaratoria, la revocatoria y la apelación. La queja surge como consecuencia de la denegación de la apelación.
}
\cqrow{¿Qué otros actos de impugnación menos frecuentes existen?}{El replanteo de prueba ante la alzada, la acción autónoma de revisión de la cosa juzgada, la nulidad del laudo arbitral (arts. 767--771 CPCCN) y el juicio ordinario posterior al juicio ejecutivo.
}
\cqrow{¿Cuándo procede, excepcionalmente, la revisión de la cosa juzgada?}{De modo excepcionalísimo y muy restrictivo, ante circunstancias excepcionales: por ejemplo, la posibilidad de contar con la prueba de ADN en juicios de filiación o el ocultamiento de pruebas esenciales mediante maniobras fraudulentas.
}
\cqrow{¿Cuál es la diferencia entre remedios y recursos \textit{stricto sensu}?}{Los remedios reparan errores procesales (vicios \textit{in procedendo}) y los resuelve el mismo juez que dictó el acto. Los recursos en sentido estricto apuntan a que un tribunal superior reexamine el acto impugnado y atacan los vicios \textit{in iudicando} (errores en la interpretación de la ley o de los hechos).
}
\cqrow{¿Dónde ubica la doctrina a la aclaratoria y a la revocatoria?}{Aunque el Código las denomine «recursos», técnicamente sería más adecuado clasificarlas dentro de los remedios.
}
}{Los actos de impugnación pueden clasificarse en incidente de nulidad, recursos y otros medios menos frecuentes. Otra clasificación distingue remedios (resueltos por el mismo juez, orientados a vicios \textit{in procedendo}) y recursos en sentido estricto (sometidos a un tribunal superior, dirigidos a vicios \textit{in iudicando}).
}

\cornelltitle{fundamentos y clasificación de los recursos}
\cornellblocksum{
\cqrow{¿Qu\'e diferencia hay entre impugnaci\'on y recurso?}{La impugnaci\'on es el g\'enero (atacar un acto, dentro o fuera del proceso); el recurso es una especie de impugnaci\'on y s\'olo procede contra resoluciones judiciales (providencias simples, interlocutorias y sentencias definitivas). Una impugnaci\'on de dictamen pericial o de liquidaci\'on es una impugnaci\'on pero no un recurso.
}
\cqrow{¿Cu\'al es el fundamento de los recursos?}{El error humano (Calamandrei): los jueces pueden equivocarse y debe existir un sistema que subsane el error y evite consolidar una injusticia. Couture vincula el t\'ermino recurso con ``volver al camino de partida'', es decir, que el tribunal superior revise el camino ya transitado para confirmar o modificar la resoluci\'on.
}
\cqrow{¿Por qu\'e la aclaratoria y la revocatoria son recursos horizontales?}{Porque se interponen, tramitan y resuelven ante el mismo juez que dict\'o la resoluci\'on. Por eso se las considera m\'as bien remedios. Se interponen y fundan en el mismo acto.
}
\cqrow{¿Por qu\'e la apelaci\'on es un recurso vertical ordinario?}{Porque se interpone ante el juez de primera instancia, pero lo resuelve la c\'amara, que est\'a por encima. Tiene dos momentos diferenciados: interposici\'on y fundamentaci\'on.
}
}{la impugnaci\'on es el g\'enero y el recurso la especie, limitada a resoluciones judiciales. El fundamento de los recursos es el error humano. Los recursos horizontales se resuelven ante el mismo juez y se interponen y fundan en el mismo acto; la apelaci\'on es un recurso vertical ordinario, que resuelve la c\'amara y tiene dos momentos.
}

\chapter[Requisitos comunes y recurso indiferente]{Requisitos comunes de los recursos y recurso indiferente}\label{cap:requisitos}

\section{Legitimación}

Las \textbf{partes} son las primeras legitimadas para hacer un planteo recursivo: el actor y el demandado. También lo están:
\begin{itemize}
\item los \textbf{terceros}, ya sea por intervención voluntaria u obligada;
\item los \textbf{ministerios públicos} que intervengan en el caso: el fiscal, en cuestiones en las que está comprometido el orden público, y el defensor de menores, en los casos en que hay niños, niñas o adolescentes en el proceso;
\item las \textbf{personas ajenas al proceso} que se ven afectadas en un interés legítimo ante una decisión o actuación procesal; entre ellas, los auxiliares de la justicia que intervienen en el proceso, por ejemplo los peritos cuando se regulan sus honorarios.
\end{itemize}

\begin{example}{Terceros con legitimación recursiva}
Reparticiones públicas o empresas privadas que, en el proceso, incumplen mandatos judiciales (como trabar embargo sobre haberes o no contestar un pedido de informes) y a las que se les impone una multa. Esos terceros, ajenos al proceso, pueden contar con legitimación para plantear un recurso, justamente porque lo decidido en el proceso los afecta de modo directo.
\end{example}

\section{El agravio, gravamen o perjuicio}

Este requisito es fundamental. Los términos \textbf{agravio}, \textbf{gravamen} y \textbf{perjuicio} son sinónimos utilizados indistintamente en doctrina y jurisprudencia.

\begin{definition}{Agravio (gravamen o perjuicio)}
Es la \textbf{insatisfacción total o parcial} de cualquiera de las pretensiones planteadas en virtud de lo decidido; o bien, para la parte demandada, el \textbf{rechazo de las defensas} opuestas. Se trata de la derrota que la decisión judicial provoca frente a los postulados sostenidos en las presentaciones de quien recurre, sea que su postulado apunte a satisfacer una pretensión o a oponer una defensa.
\end{definition}

Sin agravio no hay legitimación para recurrir. El rechazo de la pretensión para el actor, o el rechazo de las defensas para el demandado, constituyen el agravio que habilita el recurso.

\section{El plazo perentorio}

En todos los recursos existe un plazo para plantearlo, que siempre es \textbf{perentorio}, es decir, \textbf{fatal}. El derecho a recurrir debe ejercerse dentro del plazo establecido por el código. Si se deja pasar ese plazo sin plantear el recurso, se pierde la posibilidad futura de formular ese planteo.

\section{Admisibilidad y fundabilidad}

Todos los recursos presentan dos facetas.

\subsection{Admisibilidad}

\begin{definition}{Admisibilidad de un recurso}
Alude a toda la cuestión \textbf{formal} requerida para el planteo impugnativo: el plazo en que debe presentarse; los recaudos formales (si debe ser por escrito o puede ser verbal); qué debe contener el escrito y si debe ser fundado o si basta con el planteo recursivo; y todo lo que tiene que ver con el aspecto formal.
\end{definition}

La admisibilidad también indica si un recurso procede contra determinados actos y contra cuáles no. Es lo que debe analizarse \textbf{antes} de pasar al examen de la sustancia o de lo medular del planteo recursivo.

\subsection{Fundabilidad}

\begin{definition}{Fundabilidad de un recurso}
Apunta al \textbf{fondo} de la cuestión, al aspecto sustancial de la impugnación: la motivación que lleva al sujeto a formular su planteo, la crítica concreta que hace a la decisión del juez o el vicio que señala en alguna actuación procesal. Contempla los agravios, gravámenes o perjuicios que provoca el acto impugnado, en qué afecta a quien recurre y de qué manera se siente perjudicado por lo decidido.
\end{definition}

\begin{important}
\textbf{Orden de análisis.} El análisis de fundabilidad viene \textbf{después} del de admisibilidad. Primero se analiza lo formal: si el recurso procede formalmente contra determinado acto, si se planteó a tiempo y si se planteó de la forma adecuada según el código. Luego se analiza si quien recurre tiene razón en la cuestión que critica y en los motivos que le causan el perjuicio o agravio por el cual se alza contra la decisión o actuación judicial.
\end{important}

\cornelltitle{requisitos comunes de los recursos}
\cornellblocksum{
\cqrow{¿Quiénes tienen legitimación para recurrir?}{Las partes (actor y demandado), los terceros intervinientes (voluntarios u obligados), los ministerios públicos (fiscal, defensor de menores) y las personas ajenas al proceso cuyo interés legítimo resulta afectado (por ejemplo, peritos cuyos honorarios son regulados, o empresas sancionadas por incumplir mandatos judiciales).
}
\cqrow{¿Qué es el agravio o gravamen?}{Es la insatisfacción total o parcial del postulado esgrimido en el proceso (pretensión o defensa). También se denomina gravamen o perjuicio. El rechazo de la pretensión para el actor, o de las defensas para el demandado, constituye el agravio que habilita el recurso. Sin agravio no hay legitimación para recurrir.
}
\cqrow{¿Qué características tiene el plazo para recurrir?}{Es perentorio o fatal: si se deja pasar sin plantear el recurso, se pierde la posibilidad futura de formularlo.
}
\cqrow{¿Cuál es la diferencia entre admisibilidad y fundabilidad de un recurso?}{La admisibilidad se refiere al aspecto formal: plazo, forma de presentación, contenido del escrito, tipo de resolución recurrible. Se analiza primero. La fundabilidad se refiere al fondo: si los agravios son válidos y si la crítica a la decisión judicial es correcta. Se analiza después de superado el juicio de admisibilidad.
}
}{Todo recurso requiere tres presupuestos comunes: legitimación, agravio y plazo perentorio. Presenta dos dimensiones: la admisibilidad formal, que se analiza primero, y la fundabilidad sustancial, que se analiza después.
}

\section{El recurso indiferente}

El recurso indiferente \textbf{no se encuentra tipificado} en el Código Procesal ni en el ordenamiento argentino; no está legislado ni tiene recepción normativa. Algunos autores se refieren a este instituto tomándolo del derecho alemán, en el que la estructura recursiva es bien distinta de la argentina.

Su característica central es que el tribunal \textbf{deja de lado los recaudos de admisibilidad} que tienen todos los planteos impugnativos (plazo, forma, características, tipo de resolución atacada, etc.) en pos de garantizar ciertos derechos fundamentales que están siendo conculcados por alguna decisión judicial. El tribunal de segunda instancia deja de lado los recaudos formales y pasa a conocer y resolver el fondo del asunto, lo sustancial, sin cortapisas ni trabazones de índole formal.

\begin{definition}{Recurso indiferente}
Es aquel que, sin ser el que la ley prescribe expresamente para el caso, o que siéndolo se han omitido elementos formales, produce no obstante los mismos efectos respecto de la posibilidad (procedibilidad) de la vía recursiva que el recurso correctamente articulado.
\end{definition}

\subsection{Casos YPF y Tejida (CSJN)}

\begin{example}{Caso YPF (CSJN), derecho previsional}
El caso que dio andamiaje a este instituto en el ordenamiento argentino fue fallado por la Corte Suprema y es conocido como caso YPF. En el ámbito provincial se había desestimado el planteo de una persona en verdadero estado de vulnerabilidad, que tenía el beneficio de una pensión derivada de su situación de discapacidad. Por defectos en la manera en que realizó los planteos y dedujo los recursos en sede provincial, se había rechazado la posibilidad de revisar y acoger la pretensión. La Corte, dejando de lado esas cuestiones formales, admitió la viabilidad del planteo.
\end{example}

\begin{example}{Caso Tejida (CSJN, marzo de 2018)}
La Corte dijo que, dada la índole peculiar de ciertas pretensiones, compete a los jueces la búsqueda de soluciones que se avengan a éstas, para lo cual deben encauzar los trámites por vías expeditivas, a fin de evitar que la incorrecta utilización de las formas pueda conducir a la frustración de derechos tutelados constitucionalmente, cuya suspensión a las resultas de nuevos trámites es inadmisible.
\end{example}

\cornelltitle{recurso indiferente}
\cornellblocksum{
\cqrow{¿Qué es el recurso indiferente y está regulado en el CPCCN?}{Es una categoría doctrinal que no está tipificada en el Código ni legislada en el ordenamiento argentino. Algunos autores la toman del derecho alemán. Habilita a los tribunales a dejar de lado los requisitos formales de admisibilidad para garantizar derechos fundamentales gravemente afectados, y a resolver el fondo del asunto.
}
\cqrow{¿Qué efectos produce el recurso indiferente?}{Aunque no sea el recurso que la ley prescribe expresamente (o se hayan omitido elementos formales), produce los mismos efectos respecto de la procedibilidad de la vía recursiva que el recurso correctamente articulado.
}
\cqrow{¿Qué casos le sirven de sustento?}{El caso YPF (CSJN), en el que la Corte, dejando de lado cuestiones formales, admitió el planteo de una persona vulnerable con una pensión derivada de su discapacidad; y el caso Tejida (CSJN, marzo de 2018), sobre la búsqueda de soluciones por vías expeditivas para evitar la frustración de derechos constitucionalmente tutelados.
}
}{El recurso indiferente es una construcción jurisprudencial excepcional, no tipificada, para los casos en que las formas frustrarían derechos fundamentales: el tribunal deja de lado los recaudos de admisibilidad y resuelve el fondo.
}

\chapter[Doble instancia y límites del tribunal de alzada]{La garantía de la doble instancia y los límites del tribunal de alzada}\label{cap:doble}

\section{La garantía de la doble instancia}

Cabe preguntarse si existe en el ordenamiento argentino una garantía procesal, amparada por la Constitución Nacional, de que en todo proceso existirá la doble instancia; es decir, la posibilidad de que, por las vías recursivas, toda decisión pueda ser revisada por un órgano diferente de aquel que dictó la resolución.

\subsection{Perspectiva histórica}

La cuestión no es nueva en la historia del proceso. El tema de la doble instancia, incluso la posibilidad de revisiones ulteriores más allá de dos instancias, se remonta a muy vieja data. En el sistema nacional, desde tiempos de la colonia, era reconocida la posibilidad de recurrir a la audiencia, que era la instancia superior a las decisiones de los jueces de primera instancia. Mucho antes, desde la Edad Media y aun desde la época del Imperio Romano, se observa la posibilidad de recurrir las decisiones del pretor ante órganos creados especialmente por el emperador o, en épocas monárquicas, por los reyes, que delegaban en funcionarios de jerarquía la posibilidad de revisar las decisiones de los jueces inferiores.

\subsection{La Constitución Nacional de 1853: ausencia de garantía expresa}

En 1853, la Constitución Nacional no introdujo la doble instancia como garantía del debido proceso. Otros aspectos muy claros están presentes en el art. 18 desde aquella época, pero la doble instancia no está introducida como garantía constitucional: ni siquiera en el aspecto más garantista del ordenamiento, que es el de la persecución de los delitos, ni mucho menos en cuestiones no penales (civiles).

\begin{note}
\textbf{Art. 18 de la Constitución Nacional.} Establece que ningún habitante de la Nación puede ser penado sin juicio previo fundado en ley anterior al hecho del proceso, y que es inviolable la defensa en juicio de la persona y de los derechos. Este artículo, que establece las garantías del debido proceso, no menciona expresamente la doble instancia.
% TODO: contenido incompleto o ambiguo en las notas originales (la cita del art. 18 figura abreviada y con la expresión «hecho del proceso»)
\end{note}

\subsection{El Pacto de San José de Costa Rica y la reforma de 1994}

Recién con la \textbf{Convención Americana de Derechos Humanos} (Pacto de San José de Costa Rica) aparece en el ordenamiento una norma que prevé específicamente el derecho de recurrir del fallo ante juez o tribunal superior. Esto ocurre a raíz de la incorporación del tratado a la Constitución Nacional en 1994, en virtud del \textbf{art. 75, inciso 22}.

\begin{formula}{Art. 8º, inc. 2º, ap. h) CADH}
Establece el derecho de recurrir del fallo ante juez o tribunal superior. El inciso segundo se refiere a las personas inculpadas de un delito, por lo que la garantía está claramente enmarcada en el ámbito penal.
\end{formula}

El art. 75 inc. 22 de la Constitución Nacional otorga jerarquía constitucional a los tratados internacionales de derechos humanos, entre ellos la Convención Americana de Derechos Humanos.

\subsection{El caso Giroldi (CSJN)}

La Corte Suprema ratificó la garantía en materia penal en el precedente \textbf{Giroldi}. Consideró que la incorporación de la Convención Americana de Derechos Humanos a la Constitución Nacional conformaba una garantía que integraba el ordenamiento constitucional y que, por ende, debía respetarse con rango constitucional la doble instancia en materia penal. En materia penal se habla, más que de doble instancia, de \textbf{doble conforme}, para que la decisión condenatoria adquiera validez desde el plano de esta garantía.

\subsection{Materia no penal: ausencia de garantía constitucional}

\begin{important}
La Corte Suprema ha sostenido sistemáticamente que la garantía de la doble instancia \textbf{no tiene jerarquía constitucional en materia no penal}. En materia no penal no hay garantía constitucional de doble instancia, salvo cuando las leyes específicamente lo establezcan.
\end{important}

El \textbf{art. 242 CPCCN} establece, como recaudo de admisibilidad del recurso de apelación, que el monto discutido en el proceso supere cierta suma de dinero (aproximadamente trescientos mil pesos, según las fuentes originales). Toda discusión por una cuestión de menor cuantía no admite apelación, salvo que se trate de cuestiones arancelarias u honorarios.
% TODO: verificar el monto vigente; las fuentes originales indican una cifra aproximada.

Frente al interrogante de si no se priva de la doble instancia a quienes discuten montos menores a 300.000 pesos, la Corte ha dicho que las limitaciones vinculadas con la cuantía establecidas por el legislador son absolutamente razonables y no violan ninguna pauta propia del debido proceso, precisamente porque en materia no penal no existe garantía constitucional de doble instancia.

\section{Límites del conocimiento del tribunal de alzada}

Los tribunales de segunda instancia (cámaras de apelación, tribunales de alzada) sólo son competentes para intervenir por vía de revisión de las decisiones de un juez de primera instancia; no tienen competencia originaria.

\subsection{Primer límite: los escritos constitutivos del proceso}

Al igual que el juez de primera instancia, el tribunal de alzada está limitado por el \textbf{principio de congruencia}, que lo obliga a mantenerse dentro del marco de lo que las partes han postulado en los escritos constitutivos del proceso (pretensión y defensa).

\begin{formula}{Art. 277 CPCCN}
Ante los jueces de segunda instancia no se pueden introducir nuevas pretensiones o defensas no sometidas al conocimiento del juez de primera instancia. Las pretensiones, postulaciones y defensas deducidas ante el juez de primera instancia son las que pueden reeditarse ante el tribunal, con ampliación de fundamentos, pero no pueden plantearse nuevos tópicos no propuestos en primera instancia.
\end{formula}

\subsubsection{Hechos sobrevinientes (\textit{iura novit curia} / hechos nuevos)}
% TODO: contenido incompleto o ambiguo en las notas originales (el apunte rotula este tema como «iura novit curia / hechos nuevos»; anuncia dos tipos de hechos pero sólo se desarrolla uno)

Existen excepciones a este límite. Se anuncian dos tipos de hechos nuevos, de los cuales las fuentes desarrollan el siguiente: los \textbf{hechos posteriores a la sentencia de primera instancia}. Por haber ocurrido con posterioridad, no pudieron ser planteados ante el juez de primera instancia, y el art. 277 los admite.

\begin{formula}{Art. 163, inc. 6º, segunda parte, CPCCN}
Siempre el juez, sea de primera o de segunda instancia, podrá hacer mérito de los hechos constitutivos, modificativos o extintivos de las pretensiones, producidos durante la sustanciación del juicio y debidamente probados.
\end{formula}

\subsection{Segundo límite: los agravios expresados y sus contestaciones}

Además del límite de los escritos constitutivos, el tribunal de alzada está acotado por los propios recursos planteados por las partes: los recursos establecen el límite dentro del cual puede actuar, y los agravios planteados por las partes marcan ese límite de actuación.

\begin{table}[H]
\centering
\small
\begin{tabularx}{\textwidth}{>{\raggedright\arraybackslash}X >{\raggedright\arraybackslash}X}
\toprule
\textbf{Primera limitación} & \textbf{Segunda limitación} \\
\midrule
Los escritos constitutivos del proceso (pretensiones y defensas planteadas en primera instancia). &
Las expresiones de agravios y sus contestaciones (los recursos planteados por las partes). \\
\bottomrule
\end{tabularx}
\end{table}

\section{El principio de \textit{reformatio in peius}}

Este principio tiene íntima relación con los límites del tribunal de alzada.

\begin{definition}{Reformatio in peius}
La sentencia de segunda instancia \textbf{jamás podrá empeorar la situación del único apelante}.
\end{definition}

\begin{example}{Accidente de tránsito y único apelante}
En un juicio de accidente de tránsito, el juez de primera instancia fija una indemnización de 500.000 pesos en favor del actor. La parte demandada no apela, porque considera la suma razonable (o incluso baja). El actor sí apela, porque considera baja la indemnización. El recurso llega a la cámara.

El tribunal de segunda instancia advierte que, a su criterio, la indemnización fue excesiva y que el actor merecía mucho menos. Sin embargo, como el único apelante es el actor (la parte demandada consintió la sentencia), el tribunal de alzada no puede bajar la indemnización: no existe agravio que permita reducirla. A lo sumo podrá confirmar la sentencia.

\textbf{Conclusión:} cuando se apela y se es el único apelante, nunca se puede perder: o se gana más, o se sigue con lo mismo que dio el juez de primera instancia.
\end{example}

El límite por monto del artículo 242 CPCCN como recaudo de admisibilidad de la apelación se analiza con mayor detalle en el capítulo~\ref{cap:apelacion}.

\cornelltitle{doble instancia y límites del tribunal de alzada}
\cornellblocksum{
\cqrow{¿Existe en Argentina una garantía constitucional de doble instancia en materia civil?}{No. En materia no penal, la Corte Suprema sostiene sistemáticamente que no existe garantía constitucional de doble instancia, salvo que las leyes específicamente la establezcan.
}
\cqrow{¿Cómo se incorporó la garantía de doble instancia en materia penal?}{La Constitución de 1853 no la preveía. Surgió con el Pacto de San José de Costa Rica (art. 8º inc. 2º ap. h), incorporado a la Constitución por el art. 75 inc. 22 en 1994, y fue ratificada por la Corte en el caso Giroldi. En materia penal se habla de doble conforme.
}
\cqrow{¿Qué establece el art. 242 CPCCN y por qué no viola el debido proceso?}{Exige que el monto discutido supere cierta suma para admitir la apelación, salvo cuestiones arancelarias u honorarios. La Corte ha dicho que estas limitaciones de cuantía son razonables, porque en materia no penal no existe garantía constitucional de doble instancia.
}
\cqrow{¿Cuáles son los límites del conocimiento del tribunal de alzada?}{Una doble limitación: los escritos constitutivos del proceso (art. 277 CPCCN: no se pueden introducir nuevas pretensiones o defensas) y las expresiones de agravios y sus contestaciones. Se admiten hechos sobrevinientes (art. 163 inc. 6º, segunda parte).
}
\cqrow{¿Qué es el principio de \textit{reformatio in peius}?}{Prohíbe que el tribunal de segunda instancia empeore la situación del único apelante. Si sólo apeló el actor por considerar baja la indemnización y la demandada consintió la sentencia, el tribunal de alzada no puede bajar la indemnización: a lo sumo la confirma o la aumenta.
}
}{La garantía de la doble instancia existe sólo en materia penal con rango constitucional (CADH, art. 75 inc. 22 CN); en materia civil el legislador puede restringirla. El tribunal de alzada tiene una doble limitación (escritos constitutivos y agravios expresados) y no puede perjudicar al único apelante.
}

\chapter[Aclaratoria y reposición]{Remedios ante el mismo juez: aclaratoria y reposición o revocatoria}\label{cap:remedios}

\section{El recurso de aclaratoria}

\begin{definition}{Aclaratoria (Palacio)}
Es un \textbf{remedio} que se concede a las partes para obtener que el mismo juez o tribunal que dictó una resolución subsane las deficiencias materiales o conceptuales que contenga, o la integre de conformidad con las peticiones oportunamente formuladas.
\end{definition}

El \textbf{art. 166, incisos 1º y 2º, CPCCN} regula las actividades del juez después del dictado de la sentencia. El inciso 2º establece que la parte puede plantear la aclaratoria dentro del tercer día de notificada la resolución. El inciso 1º faculta al juez a corregir de oficio cualquier defecto o deficiencia material con anterioridad a la notificación a las partes. En el mismo sentido, el \textbf{art. 36, inciso 6º, CPCCN} establece, entre los deberes y facultades del juez, la facultad de corregir de oficio cualquier defecto o deficiencia material de una resolución, con anterioridad a la notificación a las partes.

\subsection{Supuestos que habilitan la aclaratoria}

\begin{table}[H]
\centering
\small
\begin{tabularx}{\textwidth}{L{0.28\textwidth} >{\raggedright\arraybackslash}X}
\toprule
\textbf{Supuesto} & \textbf{Descripción / ejemplo} \\
\midrule
\textbf{Errores materiales} &
Confundir el nombre de las partes (por ejemplo, escribir «parte actora» cuando se quería decir «parte demandada»); discordancia entre cifras en guarismos y en letras; errores aritméticos en sumas o cálculos. \\
\midrule
\textbf{Errores de copia o transcripción} &
Párrafos «pegados» de otra resolución, especialmente al armar resoluciones con procesadores de texto. Hoy son menos frecuentes que en la época de las máquinas de escribir. \\
\midrule
\textbf{Contradicción entre considerandos y parte resolutiva} &
Los considerandos fundamentan la concesión de un rubro pero la parte resolutiva lo omite, o viceversa. \\
\midrule
\textbf{Conceptos oscuros} &
Discordancia entre la idea que se quiso expresar y las palabras utilizadas, que torna necesaria una aclaración, siempre que ello no implique una modificación sustancial de lo decidido. \\
\midrule
\textbf{Omisión de tratar pretensiones deducidas} &
El juez no trató en la resolución alguna de las pretensiones debatidas en el proceso. La aclaratoria es la oportunidad para ampliar y completar la resolución. \\
\bottomrule
\end{tabularx}
\end{table}

\subsection{Características procesales de la aclaratoria}

\subsubsection{Plazo}

La parte tiene \textbf{tres (3) días} para plantear la aclaratoria, contados desde la notificación de la sentencia o resolución que se pretende aclarar. En el caso de aclaratoria contra sentencia de segunda instancia, el plazo es de \textbf{cinco (5) días}.

\subsubsection{Forma}

Se plantea por escrito, y la fundamentación se hace en el mismo escrito de interposición. Si se trata de una providencia no notificada por cédula, los tres días comienzan a correr desde el martes o viernes posterior al dictado (notificación por ministerio de la ley).

\subsubsection{Sustanciación}

Según el Código, la aclaratoria es resuelta por el juez \textbf{sin traslado previo} a la otra parte (sin sustanciación).

\subsubsection{Relación con el recurso de apelación}

\begin{important}
La aclaratoria \textbf{no interrumpe el plazo para apelar}. La apelación debe interponerse autónomamente dentro del quinto día, sin perjuicio de la aclaratoria planteada. Lo conveniente es plantear la aclaratoria y, antes del quinto día, plantear la apelación; si la aclaratoria suple los defectos y torna innecesaria la apelación, se desistirá eventualmente de ésta. No puede perderse la apelación por haber planteado una aclaratoria.

No cabe la apelación subsidiaria de ese planteo: es aclaratoria o apelación, pero no las dos cosas juntas. \textbf{El único recurso que admite la apelación en subsidio es la revocatoria.}
\end{important}

\cornelltitle{recurso de aclaratoria}
\cornellblocksum{
\cqrow{¿Para qué sirve el recurso de aclaratoria y cuándo no procede?}{Sirve para que el mismo juez corrija errores materiales (nombres, sumas, cálculos), aclare conceptos oscuros o integre la resolución cuando omitió tratar alguna pretensión. No procede para modificar sustancialmente lo decidido: para ello existen la revocatoria o la apelación.
}
\cqrow{¿Qué supuestos habilitan la aclaratoria?}{Errores materiales, errores de copia o transcripción, contradicción entre considerandos y parte resolutiva, conceptos oscuros y omisión de tratar pretensiones deducidas.
}
\cqrow{¿Cuál es el plazo para plantearla?}{Tres días desde la notificación de la resolución; cinco días si se plantea contra una sentencia de segunda instancia.
}
\cqrow{¿Cómo se plantea y cómo se resuelve?}{Por escrito, con la fundamentación en el mismo escrito. El juez la resuelve sin traslado previo a la otra parte.
}
\cqrow{¿Qué relación tiene con el recurso de apelación?}{No interrumpe el plazo para apelar: la apelación debe interponerse autónomamente dentro del quinto día. No cabe la apelación subsidiaria de la aclaratoria; el único recurso que admite apelación en subsidio es la revocatoria.
}
}{La aclaratoria permite al mismo juez corregir errores formales o integrar omisiones sin alterar la esencia de lo decidido. Se plantea por escrito dentro de los tres días (cinco contra sentencia de segunda instancia), no se sustancia y no interrumpe el plazo para apelar.
}

\section{El recurso de reposición o revocatoria}

\begin{definition}{Revocatoria o reposición}
Es el \textbf{remedio} procesal tendiente a que el mismo juez o tribunal que dictó una resolución subsane los agravios que ella ocasiona, \textbf{por contrario imperio}. Se le pide al juez o tribunal que dictó la resolución que la modifique, la cambie o la revise, porque esa resolución causa perjuicios que justifican que corrija la decisión y la modifique en favor de quien hace el planteo.
\end{definition}

Los \textbf{arts. 238 a 241 CPCCN} regulan el recurso: procedencia (art. 238), forma y trámite (arts. 239 y 240) y efectos (art. 241). Se plantea ante el mismo juez que dictó la decisión y es resuelto por ese mismo juez.

\subsection{Recaudos de admisibilidad}

\subsubsection{Tipo de resolución atacable}

Procede únicamente contra \textbf{providencias simples}. Es inadmisible contra sentencias interlocutorias o sentencias definitivas.

\subsubsection{Jueces ante quienes procede}

Procede contra decisiones de jueces de primera instancia y también contra decisiones del presidente de una cámara de apelación (providencias simples de mero trámite adoptadas durante la sustanciación del juicio en segunda instancia, con anterioridad al dictado de la sentencia de alzada).

Cuando la revocatoria se plantea contra una providencia del presidente de la cámara, resuelve el tribunal en pleno (los tres jueces). Esas decisiones del tribunal en pleno hacen ejecutoria: no pueden ser recurridas.

\subsubsection{Providencias que causen o no gravamen irreparable}

\begin{formula}{Art. 238 CPCCN}
El recurso de reposición o revocatoria procede contra providencias simples, causen o no gravamen irreparable.
\end{formula}

El \textbf{gravamen irreparable} es aquel que no puede ser suplido por lo que se decida en la sentencia: la providencia no puede ser corregida durante el curso del proceso ni con la sentencia, sino que, una vez adoptada la decisión (aunque sea en una providencia simple), marca un cambio de rumbo tal que no hay vuelta atrás.

\begin{example}{Contestación de demanda tenida por no presentada}
La providencia simple por la cual se tiene por no contestado el traslado de la demanda en término, por considerarse extemporánea la contestación, deja al demandado sin posibilidad de defensa si no se corrige por vía de revocatoria. Eso es irreparable desde el plano de la defensa. Por ello, causa gravamen irreparable y es susceptible tanto de revocatoria como de apelación.

En cambio, una providencia que ordena agregar copia de un escrito antes de correr traslado puede no causar gravamen irreparable (la cuestión se soluciona presentando la copia), pero aun así es admisible la revocatoria.
\end{example}

\begin{note}
Mientras la revocatoria procede contra providencias simples, causen o no gravamen irreparable, la apelación exige que la providencia simple cause gravamen irreparable para ser admisible.
\end{note}

\subsection{Plazo}

El plazo para plantear la revocatoria es de \textbf{tres (3) días} desde la notificación de la providencia que se pretende impugnar, ya sea por ministerio de la ley o de manera electrónica (art. 135 CPCCN).

\textbf{Excepción:} si la decisión fue adoptada en una audiencia, la revocatoria debe plantearse \textbf{en el mismo acto de la audiencia, verbalmente}. No puede presentarse días después.

\begin{example}{Audiencia preliminar}
En una audiencia preliminar el juez adopta una serie de decisiones. La parte que no esté de acuerdo con alguna de ellas tiene el derecho de plantear la revocatoria en ese mismo momento. No puede hacerlo después de la audiencia, a los tres días.
\end{example}

\subsection{Forma de interposición y fundamentación}

La revocatoria se interpone por escrito (en la mayoría de los casos) o verbalmente (si la providencia fue adoptada en audiencia). Se funda en el mismo escrito o en el mismo acto de su interposición. Se exponen los fundamentos, los agravios y los motivos que llevan a criticar el acto cuya modificación se pretende. El \textbf{art. 239 CPCCN} regula la forma de interposición, incluida la posibilidad de interposición verbal cuando la decisión se adopta en el marco de una audiencia.

\subsection{Trámite y resolución}

Si la decisión atacada fue el resultado de una \textbf{petición formulada por la otra parte}, cuando se plantea la revocatoria el juez debe oír previamente a esa otra parte: debe correrse traslado del planteo de revocatoria a la parte que había pedido la decisión que ahora se impugna. Si la decisión no se originó en la petición de una de las partes, sino que fue una providencia dictada de oficio por el juez, \textbf{no es necesario} correr traslado a la otra parte.

El \textbf{art. 240 CPCCN} regula el trámite, incluidos los supuestos en que procede el traslado a la otra parte y aquellos en que no. El traslado, cuando corresponde, cuenta con un plazo de tres días.

\subsection{Efectos: la ejecutoria y sus excepciones}

Lo que el juez decida respecto de la revocatoria \textbf{hace ejecutoria}: queda firme lo decidido y resulta ejecutable. No es impugnable por vía de otro recurso.

Excepciones a la ejecutoria:
\begin{enumerate}
\item Si la revocatoria fue planteada con \textbf{apelación subsidiaria}: si el juez rechaza la revocatoria, no hace ejecutoria, porque queda pendiente el recurso de apelación interpuesto en subsidio.
\item Si el juez \textbf{admite} la revocatoria: esa decisión puede ser apelada por la contraparte.
\end{enumerate}

\cornelltitle{recurso de reposición o revocatoria}
\cornellblocksum{
\cqrow{¿Qué es la revocatoria o reposición?}{Es el remedio por el cual se pide al mismo juez o tribunal que dictó una resolución que la modifique por contrario imperio, porque le causa un perjuicio al recurrente. Está regulada en los arts. 238 a 241 CPCCN.
}
\cqrow{¿Contra qué tipo de resoluciones procede?}{Sólo contra providencias simples, causen o no gravamen irreparable (art. 238 CPCCN). No procede contra sentencias interlocutorias ni definitivas. Procede contra decisiones de jueces de primera instancia y del presidente de una cámara de apelación. La apelación, en cambio, exige que la providencia simple cause gravamen irreparable.
}
\cqrow{¿Cuál es el plazo y qué excepción tiene?}{Tres días desde la notificación. Si la decisión se adoptó en una audiencia, debe plantearse verbalmente en el mismo acto.
}
\cqrow{¿Cómo se tramita?}{Se interpone y funda por escrito (o verbalmente en audiencia). Si la decisión atacada se originó en una petición de la otra parte, se corre traslado a ésta; si fue dictada de oficio, no es necesario el traslado.
}
\cqrow{¿Qué efectos tiene lo resuelto y qué excepciones existen?}{Hace ejecutoria: queda firme y no es impugnable por otro recurso. No hace ejecutoria si se planteó con apelación subsidiaria y se rechazó la revocatoria, ni si se admite la revocatoria, pues la contraparte puede apelar.
}
}{La revocatoria permite que el mismo juez revea providencias simples. Se plantea dentro de los tres días (o en la audiencia, si allí se adoptó la decisión), admite apelación en subsidio y, como regla, lo resuelto hace ejecutoria.
}

\chapter{El incidente de nulidad}\label{cap:nulidad}

El incidente de nulidad no es un recurso, pero resulta fundamental en el proceso.

\begin{definition}{Incidente de nulidad}
Es el \textbf{remedio} procesal que apunta a la privación de efectos imputada a los actos del proceso que adolecen de algún vicio en sus elementos esenciales y que, por ende, carecen de aptitud para cumplir el fin a que están destinados.
\end{definition}

\begin{example}{Nulidad del traslado de la demanda}
La notificación del traslado de la demanda está dotada, en virtud del art. 339 CPCCN, de una serie de recaudos que hacen a la garantía del derecho de defensa. Es un acto trascendental. Si este acto, por algún vicio o defecto, no cumple su propósito (si el demandado no se entera de que fue emplazado a comparecer a un proceso), el proceso sigue adelante sin que él pueda ejercer su defensa. El incidente de nulidad es el mecanismo que permite atacar ese acto irregular.
\end{example}

\section{Principio general: conservación de los actos procesales}

El proceso tiende a que los actos procesales cumplidos sean válidos. El principio general es el de \textbf{conservación de los actos procesales}: se reputan válidos, salvo que se reúnan los criterios que habilitan la declaración de nulidad de un acto.

\section{Los cinco principios que rigen las nulidades procesales}

Los cinco principios están regulados en los arts. 169 a 174 CPCCN.

\subsection{Principio de especificidad (art. 169, primer párrafo)}

Para que exista nulidad tiene que haber una \textbf{irregularidad formal}, una contradicción a una norma que, de no haber sido observada, apareje la nulidad. Tiene que haber una pauta formal vulnerada. Si en el proceso no se violó ninguna norma formal, no se daría este principio, necesario para que haya nulidad.

\begin{formula}{Art. 169, primer párrafo, CPCCN}
Ningún acto procesal será declarado nulo si la ley no prevé expresamente esa sanción. Sin embargo, la nulidad procederá cuando el acto carezca de los requisitos indispensables para la obtención de su finalidad.
\end{formula}

\subsection{Principio de finalidad o conservación (art. 169, segundo párrafo)}

El acto viciado debe carecer de los requisitos indispensables para obtener su finalidad. Pueden existir hipótesis en las que, pese a la irregularidad del acto cumplido y atacado, éste haya cumplido la finalidad a la que estaba destinado. En ese caso prevalece la \textbf{validez} del acto.

\begin{formula}{Art. 169, segundo párrafo, CPCCN}
Si el acto, no obstante su irregularidad, ha cumplido la finalidad a la que estaba destinado, no procede la declaración de nulidad.
\end{formula}

\subsection{Principio de subsanación (arts. 170 y 172)}
% TODO: contenido incompleto o ambiguo en las notas originales (el apunte rotula este principio con los arts. 170 y 172, pero sólo transcribe el art. 170; el art. 172 se transcribe como principio de trascendencia)

La nulidad es improcedente si el acto viciado fue \textbf{convalidado} o purgado, o fue conocido por quien la invoca y éste no hizo ningún planteo. Se presume que el acto ha sido convalidado si no se promueve el incidente de nulidad dentro del \textbf{quinto día} de conocido.

\begin{formula}{Art. 170 CPCCN}
La nulidad no podrá ser declarada cuando el acto haya sido consentido, aunque fuere tácitamente, por la parte interesada en la declaración. Se considerará que hay consentimiento tácito cuando no se promoviere incidente de nulidad dentro del quinto día de conocido el acto nulo.
\end{formula}

\begin{example}{Emplazado anoticiado de manera irregular pero con conocimiento real}
Una persona se muda de domicilio pero mantiene contacto con quien vive allí. La cédula de traslado de la demanda llega al domicilio anterior. La persona que recibe la cédula le avisa al demandado, quien toma conocimiento pleno. Si el demandado no hace nada y deja transcurrir el plazo de cinco días, el acto queda saneado y ya no procede el incidente de nulidad. El acto irregular cumplió su finalidad (el emplazado supo que era demandado) y fue convalidado.
\end{example}

\subsection{Principio de trascendencia (art. 172)}

Quien plantea el incidente debe \textbf{invocar y demostrar un perjuicio cierto} del que derive su interés en formular el planteo, dimensionando las defensas que no ha podido oponer como consecuencia del acto viciado. La nulidad debe dejar desnuda ante el juez una situación de indefensión de quien la invoca.

\begin{formula}{Art. 172 CPCCN}
El incidente de nulidad es inadmisible cuando quien lo deduce no exprese el perjuicio sufrido del que derive el interés en obtener la declaración y las defensas que no ha podido oponer.
\end{formula}

\begin{important}
\textbf{«Si hay indefensión, hay nulidad. Si no hay indefensión, no hay nulidad.»} El mecanismo de las nulidades apunta a garantizar la defensa, no la regularidad de los actos o el cumplimiento de un formalismo. No existe la nulidad por la mera forma.
\end{important}

\subsection{Principio de protección (art. 171)}

Nadie puede invocar la nulidad cuando fue partícipe o gestor de la nulidad que invoca. Por ejemplo, el actor no puede plantear la nulidad del traslado de la demanda en el domicilio que él mismo denunció y en el que postuló que se hiciera la notificación.

\begin{formula}{Art. 171 CPCCN}
La parte que ha dado lugar a la nulidad no podrá pedir la invalidez del acto realizado.
\end{formula}

\cornelltitle{incidente de nulidad}
\cornellblocksum{
\cqrow{¿Qué es el incidente de nulidad?}{Es el remedio procesal que apunta a la privación de efectos de los actos del proceso que adolecen de algún vicio en sus elementos esenciales y que, por ello, carecen de aptitud para cumplir su fin. Ataca vicios procedimentales (\textit{in procedendo}), como la notificación irregular del traslado de la demanda.
}
\cqrow{¿Qué principio general rige los actos procesales?}{La conservación de los actos procesales: se reputan válidos salvo que se reúnan los criterios que habilitan la declaración de nulidad.
}
\cqrow{¿Qué establece el principio de especificidad?}{Debe haber una norma formal vulnerada (art. 169, primer párrafo). Ningún acto será declarado nulo si la ley no prevé expresamente esa sanción, salvo que el acto carezca de los requisitos indispensables para su finalidad.
}
\cqrow{¿Qué establece el principio de finalidad?}{Si el acto, pese a su irregularidad, cumplió la finalidad a la que estaba destinado, no hay nulidad (art. 169, segundo párrafo).
}
\cqrow{¿Qué establece el principio de subsanación?}{Si el acto fue convalidado, no procede la nulidad (art. 170). Hay consentimiento tácito si no se promueve el incidente dentro del quinto día de conocido el acto.
}
\cqrow{¿Qué establece el principio de trascendencia?}{Quien plantea el incidente debe invocar y demostrar un perjuicio cierto y las defensas que no pudo oponer (art. 172). Debe existir indefensión real: sin indefensión, no hay nulidad.
}
\cqrow{¿Qué establece el principio de protección?}{Quien generó el vicio no puede invocarlo (art. 171).
}
}{El incidente de nulidad ataca vicios procedimentales y exige los cinco principios: especificidad, finalidad, subsanación, trascendencia y protección. La clave es la trascendencia: sin indefensión real, no hay nulidad.
}

\chapter[Recurso de apelación: concepto, procedencia y forma]{El recurso de apelación: concepto, procedencia y forma}\label{cap:apelacion}

% --------------------------------------------------

\section{Definici\'on del recurso de apelaci\'on}

\begin{definition}{Recurso de apelaci\'on}
El recurso de apelaci\'on es un \textbf{acto procesal de impugnaci\'on} que interponen las partes o terceros ante el mismo juez que dict\'o una resoluci\'on, que debe ser una providencia simple, interlocutoria o definitiva, en tanto y en cuanto cause \textbf{gravamen irreparable}, con el objetivo de que el tribunal que est\'a por encima de ese juzgado la revoque total o parcialmente.
\end{definition}

\section{Legitimados para interponer el recurso}

No son s\'olo las partes en sentido estricto (actor o demandado): cualquier sujeto procesal con \textbf{inter\'es acreditado} puede interponer el recurso. Ejemplos:

\begin{itemize}
    \item El Ministerio P\'ublico, si est\'a representando los intereses de un menor o de un incapaz.
    \item Un perito que considera exiguos sus honorarios regulados.
    \item El abogado interviniente, respecto de la regulaci\'on de sus honorarios.
\end{itemize}

El requisito es que exista un \textbf{gravamen irreparable}: concreto, actual, que no se pueda reparar luego con el dictado de la sentencia definitiva.

\section{Resoluciones apelables: la regla de la apelabilidad}

En el proceso ordinario la regla es la \textbf{apelabilidad}: todas las resoluciones son apelables, salvo que el C\'odigo establezca lo contrario. Esto es diferente a lo que ocurre en el proceso ejecutivo o en el sumar\'isimo, donde la apelabilidad est\'a muy restringida.

\begin{important}
\textbf{Inapelable} no es lo mismo que \textbf{irrecurrible}. Si el C\'odigo dice ``irrecurrible'', la resoluci\'on no resiste ning\'un tipo de recurso (ni siquiera revocatoria). Cuando dice ``inapelable'', se refiere \'unicamente a que no resiste el recurso de apelaci\'on; otros recursos podr\'ian interponerse.
\end{important}

\subsection{Art\'iculo 379 CPCCN: inapelabilidad de resoluciones sobre prueba}

% TODO: contenido incompleto o ambiguo en las notas originales (la redacci\'on del segundo p\'arrafo del art. 379 figura as\'i en el apunte)
\begin{formula}{Art. 379 CPCCN}
\textit{Ser\'an inapelables las resoluciones del juez sobre producci\'on, denegaci\'on y sustanciaci\'on de las pruebas. Si se hubiere negado alguna medida, la parte interesada podr\'a solicitar a la c\'amara que la diligencia cuando el expediente le fuere remitido para que conozca del recurso contra la sentencia definitiva.}
\end{formula}

\begin{example}{Denegaci\'on de prueba testimonial en audiencia preliminar}
En la audiencia preliminar, el juez decide denegar la producci\'on de la prueba testimonial por considerarla innecesaria en funci\'on de los hechos controvertidos. La parte actora o demandada no puede apelar esa resoluci\'on, porque el art.~379 la declara inapelable. Sin embargo, podr\'a replantear la cuesti\'on ante la c\'amara cuando el expediente sea elevado para resolver el recurso contra la sentencia definitiva.
\end{example}

\subsection{Art\'iculo 242 CPCCN: l\'imite por monto}

\begin{formula}{Art. 242 CPCCN (actualizado por Acordada 41/2019, vigente desde el 1.º de enero de 2020)}
\textit{Ser\'an inapelables las sentencias definitivas y las dem\'as resoluciones, cualquiera fuere su naturaleza, que se dicten en los procesos en los que el valor cuestionado no supere los \$\,300.000 (pesos trescientos mil). Esta inapelabilidad no rige en materia de alimentos ni respecto de los honorarios regulados en ese proceso.}
\end{formula}

Muchos alumnos pueden tener c\'odigos procesales desactualizados: el monto se va actualizando mediante acordadas de la C\'amara, siendo la \'ultima la Acordada 41/2019.

Se ha cuestionado la constitucionalidad de este art\'iculo por limitar el acceso a la segunda instancia. La Corte Suprema ha interpretado que la \textbf{doble instancia} como principio constitucional-convencional (Pacto de San Jos\'e de Costa Rica) rige s\'olo en materia penal. En materia civil, el legislador puede establecer este l\'imite sin que resulte inconstitucional ni anticonvencional. La finalidad es que los juicios de menor trascendencia econ\'omica tramiten en una sola instancia para ser resueltos m\'as r\'apidamente.

\section{Revocatoria con apelaci\'on en subsidio}

La \'unica combinaci\'on posible de recursos es la \textbf{revocatoria con apelaci\'on en subsidio}, y s\'olo procede ante una \textbf{providencia simple que cause gravamen irreparable}. En ese supuesto, el recurrente interpone la revocatoria y, para el caso de que el juez de primera instancia la rechace, subsidiariamente est\'a la apelaci\'on para ir a la c\'amara.

\begin{example}{Providencia simple que causa gravamen irreparable}
Si se contest\'o la demanda en t\'ermino pero el juzgado la tuvo por extempor\'anea por haber computado mal el plazo de 15 d\'ias (sin tener en cuenta un feriado), eso causa un gravamen irreparable. En ese supuesto: (a) puede interponerse una revocatoria con apelaci\'on en subsidio, o (b) directamente un recurso de apelaci\'on.
\end{example}

\section{Errores que comprende el recurso de apelaci\'on}

\begin{formula}{Art. 253 CPCCN}
El recurso de apelaci\'on comprende tanto los errores \textit{in iudicando} (error en el juzgamiento: fijaci\'on de los hechos o aplicaci\'on del derecho) como los errores \textit{in procedendo} (error en el procedimiento de construcci\'on de la sentencia). Por ello, en el proceso ordinario no existe un recurso de nulidad separado: el recurso de nulidad est\'a comprendido dentro del recurso de apelaci\'on.
\end{formula}

\section{Plazo y forma de interposici\'on}

\begin{table}[H]
\centering
\begin{tabularx}{\textwidth}{
    >{\raggedright\arraybackslash}p{0.25\textwidth}
    >{\raggedright\arraybackslash}p{0.18\textwidth}
    >{\raggedright\arraybackslash}X
}
\toprule
\textbf{Proceso} & \textbf{Plazo} & \textbf{Observaci\'on} \\
\midrule
Proceso ordinario & \textbf{5 d\'ias} & Desde la notificaci\'on (c\'edula o ministerio de ley) \\
\addlinespace
Proceso sumar\'isimo & \textbf{3 d\'ias} & Plazo especial previsto por el c\'odigo \\
\addlinespace
Verbal (en audiencia) & \textbf{Inmediato} & Se deja constancia en el acta \\
\bottomrule
\end{tabularx}
\caption{Plazo de interposici\'on del recurso de apelaci\'on seg\'un el tipo de proceso.}
\end{table}

\section{Los dos momentos del recurso de apelaci\'on}

A diferencia de los recursos horizontales, el recurso de apelaci\'on tiene \textbf{dos momentos diferenciados}.

\begin{table}[H]
\centering
\begin{tabularx}{\textwidth}{
    >{\raggedright\arraybackslash}X
    >{\raggedright\arraybackslash}X
}
\toprule
\textbf{1.º momento: interposici\'on} & \textbf{2.º momento: fundamentaci\'on} \\
\midrule
Dentro de los 5 d\'ias de notificada la resoluci\'on. S\'olo se declara que se interpone el recurso de apelaci\'on contra la resoluci\'on de fecha tal, por causar gravamen irreparable. \textbf{Nada m\'as.} &
Una vez concedido el recurso, en funci\'on de c\'omo fue concedido. Se hace una cr\'itica concreta y razonada de la resoluci\'on, explicando a la c\'amara qu\'e se modifica total o parcialmente y d\'onde est\'a el error del juez de primera instancia. \newline
\textbf{Sentencia definitiva:} expresi\'on de agravios. \newline
\textbf{Otras resoluciones:} memorial. \\
\bottomrule
\end{tabularx}
\caption{Momentos del recurso de apelaci\'on.}
\end{table}

\begin{important}
Si no se fundamenta en tiempo y forma, se tiene por \textbf{desierto} el recurso: es como si no se hubiera interpuesto y la resoluci\'on queda firme.
\end{important}

\cornelltitle{el recurso de apelación}
\cornellblock{
\cqrow{¿Qu\'e es el recurso de apelaci\'on?}{Es un acto procesal de impugnaci\'on que interponen las partes o terceros ante el mismo juez que dict\'o una resoluci\'on (providencia simple, interlocutoria o definitiva), en tanto cause gravamen irreparable, para que el tribunal superior la revoque total o parcialmente.
}
\cqrow{¿Qui\'enes est\'an legitimados para apelar y qu\'e requisito se exige?}{Cualquier sujeto procesal con inter\'es acreditado: por ejemplo, el Ministerio P\'ublico que representa a un menor o incapaz, un perito respecto de sus honorarios o el abogado respecto de la regulaci\'on de los suyos. Se exige un gravamen irreparable, concreto y actual, que no pueda repararse con la sentencia definitiva.
}
\cqrow{¿Cu\'al es la regla sobre apelabilidad y qu\'e diferencia hay entre inapelable e irrecurrible?}{En el proceso ordinario la regla es la apelabilidad. ``Inapelable'' significa que no procede la apelaci\'on, aunque podr\'ian interponerse otros recursos; ``irrecurrible'' significa que la resoluci\'on no resiste ning\'un recurso.
}
\cqrow{¿Qu\'e establece el art.~379 CPCCN?}{Son inapelables las resoluciones del juez sobre producci\'on, denegaci\'on y sustanciaci\'on de las pruebas. La parte interesada puede replantear la cuesti\'on ante la c\'amara cuando el expediente sea remitido para conocer del recurso contra la sentencia definitiva.
}
\cqrow{¿Qu\'e establece el art.~242 CPCCN y por qu\'e no es inconstitucional?}{Son inapelables las resoluciones en procesos cuyo valor cuestionado no supere \$\,300.000 (Acordada 41/2019), salvo en materia de alimentos y honorarios regulados. La Corte Suprema interpret\'o que la doble instancia como principio constitucional-convencional rige s\'olo en materia penal, por lo que en materia civil el l\'imite es v\'alido.
}
\cqrow{¿Qu\'e errores comprende la apelaci\'on?}{Los errores \textit{in iudicando} (fijaci\'on de hechos o aplicaci\'on del derecho) y los errores \textit{in procedendo} (procedimiento de construcci\'on de la sentencia). El recurso de nulidad est\'a comprendido en la apelaci\'on (art.~253).
}
\cqrow{¿Cu\'al es el plazo para interponer la apelaci\'on?}{5 d\'ias en el proceso ordinario, desde la notificaci\'on; 3 d\'ias en el proceso sumar\'isimo; de forma inmediata en el proceso verbal en audiencia, con constancia en el acta.
}
}
\medskip
\cornellblocksum{
\cqrow{¿Cu\'ales son los dos momentos de la apelaci\'on?}{Interposici\'on (dentro de los 5 d\'ias, s\'olo se declara que se apela) y fundamentaci\'on (cr\'itica concreta y razonada, que es expresi\'on de agravios si se apela la sentencia definitiva y memorial en los dem\'as casos). Si no se funda en tiempo y forma, el recurso se tiene por desierto.
}
}{la apelaci\'on procede ante resoluciones que causan gravamen irreparable y su regla es la apelabilidad en el proceso ordinario, con l\'imites como los arts.~379 y 242. Comprende errores de juzgamiento y de procedimiento, se interpone en 5 d\'ias y se funda en un segundo momento; sin fundamentaci\'on el recurso queda desierto. La \'unica combinaci\'on posible de recursos es revocatoria con apelaci\'on en subsidio ante providencias simples.
}

\chapter[Concesión, trámite y efectos de la apelación]{Concesión, trámite y efectos del recurso de apelación}\label{cap:concesion}

Interpuesto el recurso de apelación, corresponde analizar cómo se concede, qué trámite sigue y qué efectos produce. Este capítulo reúne la forma de concesión del recurso (libre o en relación), su trámite (inmediato o diferido) y sus efectos (suspensivo o no suspensivo).

\section{Las tres consecuencias de la concesi\'on}

Al interponerse el recurso de apelaci\'on y ser concedido, se generan tres consecuencias: una forma, un tr\'amite y un efecto.

\begin{table}[H]
\centering
\begin{tabularx}{\textwidth}{
    >{\raggedright\arraybackslash}X
    >{\raggedright\arraybackslash}X
    >{\raggedright\arraybackslash}X
}
\toprule
\textbf{Forma} & \textbf{Tr\'amite} & \textbf{Efecto} \\
\midrule
\textbf{Libre} (amplia) vs. \textbf{en relaci\'on} (restringida) &
\textbf{Inmediato} (regla general) vs. \textbf{diferido} (excepci\'on legal) &
\textbf{Suspensivo} (regla general) vs. \textbf{no suspensivo} o devolutivo (excepci\'on) \\
\bottomrule
\end{tabularx}
\caption{Consecuencias de la concesi\'on del recurso.}
\end{table}

\section{Forma de concesi\'on: libre vs. en relaci\'on}

\begin{table}[H]
\centering
\begin{tabularx}{\textwidth}{
    >{\raggedright\arraybackslash}p{0.22\textwidth}
    >{\raggedright\arraybackslash}X
    >{\raggedright\arraybackslash}X
}
\toprule
 & \textbf{Concesi\'on libre} & \textbf{Concesi\'on en relaci\'on} \\
\midrule
\textbf{Supuesto} & S\'olo cuando se apela la sentencia definitiva de un proceso ordinario & Todos los dem\'as casos \\
\addlinespace
\textbf{Plazo para fundar} & 10 d\'ias (expresi\'on de agravios) & 5 d\'ias (memorial) \\
\addlinespace
\textbf{D\'onde se funda} & Ante la c\'amara & Ante el juzgado de primera instancia \\
\addlinespace
\textbf{Actividad adicional} & S\'i: art.~260 (5 d\'ias paralelos) & No \\
\addlinespace
\textbf{Traslado de fundamentos} & 10 d\'ias (notificaci\'on por ministerio de ley) & 5 d\'ias (notificaci\'on por ministerio de ley) \\
\addlinespace
\textbf{Efecto} & Siempre suspensivo & Depende de cada caso \\
\addlinespace
\textbf{Tr\'amite} & Siempre inmediato & Depende de cada caso \\
\bottomrule
\end{tabularx}
\caption{Concesi\'on libre y concesi\'on en relaci\'on.}
\end{table}

\subsection{Procedimiento cuando se concede en relaci\'on}

\begin{enumerate}
    \item El recurrente queda notificado de la resoluci\'on (por c\'edula o por ministerio de ley).
    \item Dentro de los 5 d\'ias interpone el recurso de apelaci\'on.
    \item El juez dicta una resoluci\'on concediendo en relaci\'on el recurso (se notifica por ministerio de ley).
    \item Desde que queda notificada esa concesi\'on, el recurrente tiene 5 d\'ias para presentar el memorial (fundamentaci\'on).
    \item Del memorial se corre traslado a la otra parte por 5 d\'ias (notificaci\'on por ministerio de ley).
    \item Vencido ese plazo o contestado el traslado, el expediente se eleva a la c\'amara.
    \item La c\'amara, a trav\'es de la sala correspondiente, resuelve directamente el recurso.
\end{enumerate}

\begin{note}
\textbf{Principio de prevenci\'on:} si a lo largo del proceso el expediente ya fue elevado a la c\'amara para resolver otro recurso, se remite a la misma sala que ya intervino. Si es la primera vez, se sortea la sala.
\end{note}

\subsection{Procedimiento cuando se concede libremente}

Este procedimiento es totalmente distinto al anterior.

\begin{enumerate}
    \item Se dicta la sentencia definitiva en un proceso ordinario.
    \item Se notifica por c\'edula. Dentro de los 5 d\'ias se interpone el recurso.
    \item El juez concede libremente todos los recursos interpuestos y el expediente es elevado directamente a la c\'amara.
    \item La sala dicta una resoluci\'on que dice, literalmente: \textit{``P\'onganse los autos a los fines de los art\'iculos 259 y 260 del C\'odigo Procesal. Notif\'iquese.''}
    \item Las partes son notificadas por c\'edula de esa resoluci\'on. Desde esa notificaci\'on corren en paralelo dos plazos:
    \begin{itemize}
        \item 10 d\'ias (art.~259): para expresar agravios.
        \item 5 d\'ias (art.~260): para realizar la actividad adicional en c\'amara.
    \end{itemize}
    \item De la expresi\'on de agravios se corre traslado a la otra parte por 10 d\'ias (notificaci\'on por ministerio de ley).
    \item El tribunal dicta resoluci\'on de autos para sentencia y vota para resolver el recurso.
\end{enumerate}

\subsection{Art\'iculo 260: actividad adicional en c\'amara}

% TODO: contenido incompleto o ambiguo en las notas originales (el apunte transcribe "elevaci\'on" con errores de tipeo en el art. 260)
\begin{formula}{Art. 260 CPCCN, facultades de las partes en segunda instancia (concesi\'on libre)}
\textit{Dentro de los cinco d\'ias de la notificaci\'on de la providencia que dispone la elevaci\'on, los litigantes podr\'an:}
\begin{enumerate}[label=\arabic*.º]
    \item \textit{Fundar los recursos concedidos con tr\'amite diferido.}
    \item \textit{Replantear la prueba denegada o declarada negligente.}
    \item \textit{Acompa\~nar documentos de fecha posterior o desconocidos al inicio del proceso.}
    \item \textit{Invocar hechos nuevos.}
    \item \textit{Pedir la absoluci\'on de posiciones sobre hechos no debatidos en la confesional de primera instancia.}
    \item \textit{Pedir apertura a prueba para la producci\'on del replanteo o de los hechos nuevos.}
\end{enumerate}
\end{formula}

\begin{example}{Replanteo de prueba (art.~260 en relaci\'on con art.~379)}
En la audiencia preliminar el juez deneg\'o la prueba testimonial por considerarla innecesaria. Luego dict\'o sentencia rechazando la demanda por no tener por acreditada la materialidad del accidente. Si bien en su momento esa resoluci\'on era inapelable (art.~379), ahora, al estar el expediente en c\'amara, el recurrente puede replantear la prueba denegada, fundando c\'omo la declaraci\'on de esos testigos presenciales podr\'ia acreditar la existencia del accidente y cambiar el resultado de la sentencia.
\end{example}

\section{L\'imites de la alzada al resolver}

La c\'amara tiene tres l\'imites para resolver:

\begin{enumerate}
    \item \textbf{Expresi\'on de agravios o memorial:} la c\'amara s\'olo entiende en aquello que haya sido materia de cr\'itica concreta y razonada. Lo que no fue objeto de agravios queda firme.
    \item \textbf{Principio de congruencia:} el tribunal no puede considerar hechos que no hayan sido invocados en la demanda, la contestaci\'on, la reconvenci\'on o la contestaci\'on de la reconvenci\'on.
    \item \textbf{\textit{Reformatio in peius}:} la c\'amara no puede empeorar la situaci\'on del recurrente. S\'olo puede modificar en lo que fue impugnado.
\end{enumerate}

\cornelltitle{concesión del recurso}
\cornellblocksum{
\cqrow{¿Qu\'e tres consecuencias genera la concesi\'on del recurso?}{Una forma (libre o en relaci\'on), un tr\'amite (inmediato o diferido) y un efecto (suspensivo o no suspensivo).
}
\cqrow{¿Cu\'ando se concede libremente?}{S\'olo cuando se apela la sentencia definitiva de un proceso ordinario. En todos los dem\'as casos se concede en relaci\'on.
}
\cqrow{¿C\'omo es el procedimiento cuando se concede en relaci\'on?}{Se interpone en 5 d\'ias; el juez concede en relaci\'on; desde la notificaci\'on de la concesi\'on el recurrente tiene 5 d\'ias para presentar el memorial ante el juzgado de primera instancia; se corre traslado por 5 d\'ias; luego el expediente se eleva a la c\'amara, que resuelve directamente. Rige el principio de prevenci\'on en la elecci\'on de la sala.
}
\cqrow{¿C\'omo es el procedimiento cuando se concede libremente?}{Concedido el recurso, el expediente se eleva a la c\'amara y la sala ordena poner los autos a los fines de los arts.~259 y 260. Desde la notificaci\'on por c\'edula corren en paralelo 10 d\'ias para expresar agravios y 5 d\'ias para la actividad adicional. De la expresi\'on de agravios se corre traslado por 10 d\'ias y el tribunal dicta autos para sentencia y vota.
}
\cqrow{¿Qu\'e puede hacerse en c\'amara seg\'un el art.~260?}{Dentro de los 5 d\'ias: fundar recursos concedidos con tr\'amite diferido, replantear prueba denegada o declarada negligente, acompa\~nar documentos posteriores o desconocidos, invocar hechos nuevos, pedir absoluci\'on de posiciones sobre hechos no debatidos en primera instancia y pedir apertura a prueba para el replanteo o los hechos nuevos.
}
\cqrow{¿Cu\'ales son los l\'imites de la alzada?}{La expresi\'on de agravios o memorial (lo que no fue criticado queda firme), el principio de congruencia (no se consideran hechos no invocados en los escritos postulatorios) y la prohibici\'on de \textit{reformatio in peius} (no se puede empeorar la situaci\'on del recurrente).
}
}{la concesi\'on determina forma, tr\'amite y efecto. La forma libre se reserva a la sentencia definitiva del proceso ordinario (fundamentaci\'on ante la c\'amara en 10 d\'ias, con actividad adicional del art.~260); en los dem\'as casos se concede en relaci\'on (memorial ante primera instancia en 5 d\'ias). La c\'amara resuelve dentro de los l\'imites de los agravios, la congruencia y la prohibici\'on de \textit{reformatio in peius}.
}

\section{Tr\'amite: inmediato vs. diferido}

La regla general es que el recurso se concede con \textbf{tr\'amite inmediato}: una vez interpuesto y concedido, de modo continuo se computa el plazo para fundar, se corre traslado y la c\'amara resuelve.

La excepci\'on es el \textbf{tr\'amite diferido}, cuando la ley lo dispone: el momento de fundamentaci\'on, sustanciaci\'on y resoluci\'on se difiere para cuando el expediente es elevado a la c\'amara por el recurso contra la sentencia definitiva.

\begin{important}
Advertencia terminol\'ogica: el legislador usa la expresi\'on ``efecto diferido'' en el c\'odigo, pero en realidad se trata de un tema de tr\'amite, no de efecto. Cuando el c\'odigo dice ``efecto diferido'' debe leerse ``tr\'amite diferido''.
\end{important}

\subsection{Supuestos de tr\'amite diferido}

\begin{itemize}
    \item \textbf{Art\'iculo 366, hechos nuevos:} la resoluci\'on que rechaza el planteo de hechos nuevos puede apelarse con tr\'amite diferido. La fundamentaci\'on se hace cuando el expediente llega a la c\'amara para resolver el recurso contra la sentencia definitiva. Finalidad: evitar que el expediente sea elevado varias veces antes de la sentencia.
    \item \textbf{Art\'iculo 69, costas y honorarios incidentales:} las apelaciones sobre costas y honorarios regulados a lo largo del proceso se conceden con tr\'amite diferido para el mismo momento. Finalidad: celeridad y econom\'ia procesal.
\end{itemize}

\section{Efectos: suspensivo vs. no suspensivo (devolutivo)}

El efecto hace referencia a qu\'e sucede con la resoluci\'on impugnada: si se suspenden o no sus efectos. La regla general es el \textbf{efecto suspensivo}: al interponerse el recurso, se suspende la ejecutoriedad de la resoluci\'on hasta que la c\'amara resuelva. La raz\'on es que la c\'amara puede modificar o revocar la resoluci\'on, por lo que tiene sentido no ejecutarla mientras est\'a siendo revisada.

\subsection{Origen hist\'orico del t\'ermino ``devolutivo''}

El t\'ermino ``devolutivo'' que usa el c\'odigo tiene origen en el derecho romano: el \'unico con jurisdicci\'on era el emperador, quien la delegaba en los jueces. Cuando una parte impugnaba una decisi\'on, se ``devolv\'ia'' la jurisdicci\'on al emperador. En el sistema actual, todos los jueces (desde el de primera instancia hasta la Corte Suprema) tienen la misma funci\'on jurisdiccional; lo que var\'ia es el \'ambito de su competencia. Por eso el t\'ermino ``devolutivo'' puede llevar a confusi\'on.

\begin{important}
Cuando el c\'odigo dice ``devolutivo'', debe leerse \textbf{``no suspensivo''}.
\end{important}

\subsection{Supuestos de efecto no suspensivo}

\begin{itemize}
    \item \textbf{Alimentos (sentencia que admite):} la sentencia que fija alimentos a cargo del demandado puede ejecutarse aunque est\'e pendiente el recurso de apelaci\'on del demandado ante la c\'amara.
    \item \textbf{Medidas cautelares (art.~198):} la resoluci\'on que concede una medida cautelar (por ejemplo, un embargo) puede apelarse por el afectado, pero el recurso tiene efecto no suspensivo: el embargo se traba y contin\'ua trab\'andose sin importar que el recurso est\'e pendiente ante la c\'amara.
    \item \textbf{Beneficio de litigar sin gastos (art.~81):} la resoluci\'on que concede el beneficio tiene efecto no suspensivo: se aplica de inmediato, aunque haya recurso en tr\'amite.
\end{itemize}

\section{Art\'iculo 250: procedimiento seg\'un el tipo de resoluci\'on}

Cuando el recurso se concede con efecto no suspensivo, el procedimiento var\'ia seg\'un el tipo de resoluci\'on impugnada.

\begin{table}[H]
\centering
\begin{tabularx}{\textwidth}{
    >{\raggedright\arraybackslash}X
    >{\raggedright\arraybackslash}X
}
\toprule
\textbf{La resoluci\'on es una sentencia} & \textbf{La resoluci\'on es una interlocutoria} \\
\midrule
El expediente se eleva a la c\'amara para que pueda revisar todo el proceso y resolver el recurso. Para ejecutar lo resuelto en la sentencia, se forma un \textbf{incidente de ejecuci\'on de sentencia} que queda en primera instancia (con copia certificada de la sentencia). \newline
\textit{Expediente principal $\rightarrow$ c\'amara. Incidente de ejecuci\'on $\rightarrow$ primera instancia.} &
El expediente queda en primera instancia para continuar el proceso (la sentencia definitiva a\'un no fue dictada). Se forma un \textbf{incidente de apelaci\'on} con las copias indicadas por el juez, que se eleva a la c\'amara. El apelante tiene 5 d\'ias para acompa\~nar esas copias. \newline
\textit{Expediente principal $\rightarrow$ primera instancia. Incidente de apelaci\'on $\rightarrow$ c\'amara.} \\
\bottomrule
\end{tabularx}
\caption{Procedimiento del art\'iculo 250 seg\'un el tipo de resoluci\'on (efecto no suspensivo).}
\end{table}

Si el apelante no acompa\~na las copias para formar el incidente dentro de los 5 d\'ias, se declara la \textbf{deserci\'on} del recurso: es como si no se hubiera interpuesto.

\cornelltitle{trámite y efectos de la apelación}
\cornellblocksum{
\cqrow{¿Qu\'e es el tr\'amite diferido?}{Es la excepci\'on legal a la regla del tr\'amite inmediato: la fundamentaci\'on, sustanciaci\'on y resoluci\'on del recurso se difieren para cuando el expediente sea elevado a la c\'amara por el recurso contra la sentencia definitiva.
}
\cqrow{¿Cu\'ales son los supuestos de tr\'amite diferido y su finalidad?}{El art.~366 (rechazo del planteo de hechos nuevos) y el art.~69 (costas y honorarios incidentales). Su finalidad es la celeridad y la econom\'ia procesal, evitando que el expediente se eleve varias veces antes de la sentencia.
}
\cqrow{¿Qu\'e significa ``efecto diferido'' en el c\'odigo?}{Se trata de un tema de tr\'amite y no de efecto: donde el c\'odigo dice ``efecto diferido'' debe leerse ``tr\'amite diferido''.
}
\cqrow{¿Qu\'e es el efecto suspensivo?}{Al interponerse el recurso se suspende la ejecutoriedad de la resoluci\'on impugnada hasta que la c\'amara resuelva. Es la regla general, porque la c\'amara puede modificar o revocar la resoluci\'on.
}
\cqrow{¿Qu\'e es el efecto no suspensivo y c\'omo debe leerse ``devolutivo''?}{Es la excepci\'on: la resoluci\'on puede ejecutarse aunque el recurso est\'e pendiente. El t\'ermino ``devolutivo'' proviene del derecho romano (devoluci\'on de la jurisdicci\'on al emperador) y debe leerse como ``no suspensivo''. Ejemplos: alimentos, medidas cautelares (art.~198) y beneficio de litigar sin gastos (art.~81).
}
\cqrow{¿Qu\'e ocurre, seg\'un el art.~250, con el efecto no suspensivo?}{Si se apela una sentencia, el expediente va a la c\'amara y se forma un incidente de ejecuci\'on que queda en primera instancia. Si se apela una interlocutoria, el expediente queda en primera instancia y se forma un incidente de apelaci\'on que se eleva a la c\'amara, con copias que el apelante debe acompa\~nar en 5 d\'ias; de lo contrario, el recurso se declara desierto por deserci\'on.
}
}{el tr\'amite es inmediato como regla y diferido por excepci\'on legal (arts.~366 y 69); el efecto es suspensivo como regla y no suspensivo (devolutivo) por excepci\'on (alimentos, cautelares, beneficio de litigar sin gastos). Con efecto no suspensivo, el art.~250 organiza el expediente seg\'un se apele una sentencia o una interlocutoria, y la falta de copias en 5 d\'ias produce la deserci\'on.
}

\chapter[Reglas y mecanismos especiales de la apelación]{Reglas, ejercicios y mecanismos especiales de la apelación}\label{cap:reglas}

\section{Las seis reglas del recurso de apelación}

Las reglas estudiadas permiten sintetizar el r\'egimen del recurso de apelaci\'on.

\begin{longtable}{
    >{\raggedright\arraybackslash}p{0.13\textwidth}
    >{\raggedright\arraybackslash}p{0.22\textwidth}
    >{\raggedright\arraybackslash}p{0.50\textwidth}
}
\toprule
\textbf{N.º} & \textbf{Regla} & \textbf{Contenido} \\
\midrule
\endfirsthead
\toprule
\textbf{N.º} & \textbf{Regla} & \textbf{Contenido} \\
\midrule
\endhead
\bottomrule
\endlastfoot

\textbf{Regla 1} & \textbf{Alcance del recurso} & El recurso de apelaci\'on comprende tanto los errores \textit{in iudicando} (fijaci\'on de hechos o aplicaci\'on del derecho) como los errores \textit{in procedendo} (procedimiento de construcci\'on de la sentencia). El recurso de nulidad est\'a incluido dentro del de apelaci\'on (art.~253 CPCCN). \\
\addlinespace
\textbf{Regla 2} & \textbf{Plazo de interposici\'on} & El plazo para interponer el recurso es de 5 d\'ias, salvo que la ley disponga lo contrario (por ejemplo, proceso sumar\'isimo: 3 d\'ias). \\
\addlinespace
\textbf{Regla 3} & \textbf{Forma de concesi\'on} & El recurso se concede siempre en relaci\'on, salvo un \'unico supuesto en que se concede libremente: cuando lo que se apela es la sentencia definitiva dictada en un proceso ordinario. \\
\addlinespace
\textbf{Regla 4} & \textbf{Libre: siempre suspensivo e inmediato} & Cuando el recurso se concede libremente (sentencia definitiva de proceso ordinario), siempre es con efecto suspensivo y tr\'amite inmediato. \\
\addlinespace
\textbf{Regla 5} & \textbf{Tr\'amite} & El recurso se concede siempre con tr\'amite inmediato, salvo que la ley disponga expresamente el tr\'amite diferido (por ejemplo, arts.~366 y 69 CPCCN). \\
\addlinespace
\textbf{Regla 6} & \textbf{Efecto} & El recurso se concede siempre con efecto suspensivo, salvo que la ley disponga lo contrario (efecto no suspensivo o devolutivo; por ejemplo, alimentos, medidas cautelares y beneficio de litigar sin gastos). \\
\end{longtable}

\cornelltitle{las seis reglas de la apelación}
\cornellblocksum{
\cqrow{¿Cu\'ales son las reglas sobre alcance y plazo?}{La apelaci\'on comprende errores \textit{in iudicando} e \textit{in procedendo}, con el recurso de nulidad incluido (art.~253). El plazo es de 5 d\'ias, salvo disposici\'on legal en contrario, como en el proceso sumar\'isimo (3 d\'ias).
}
\cqrow{¿Cu\'ales son las reglas sobre concesi\'on libre?}{El recurso se concede siempre en relaci\'on, salvo cuando se apela la sentencia definitiva de un proceso ordinario, en cuyo caso se concede libremente, siempre con efecto suspensivo y tr\'amite inmediato.
}
\cqrow{¿Cu\'ales son las reglas sobre tr\'amite y efecto?}{El tr\'amite es siempre inmediato salvo disposici\'on expresa de tr\'amite diferido (arts.~366 y 69). El efecto es siempre suspensivo salvo disposici\'on legal en contrario (alimentos, medidas cautelares, beneficio de litigar sin gastos).
}
}{cada regla del recurso de apelaci\'on admite excepciones legales: plazo de 5 d\'ias (salvo 3 en el sumar\'isimo), concesi\'on en relaci\'on (salvo la libre para la sentencia definitiva del proceso ordinario), tr\'amite inmediato (salvo diferido) y efecto suspensivo (salvo no suspensivo).
}

\section{Admisibilidad vs. fundabilidad: \textquestiondown qui\'en analiza qu\'e?}

\begin{exercise}{Admisibilidad y fundabilidad}
\textbf{Pregunta:} \textit{¿Cu\'al va a ser el juez o tribunal que analiza, en el caso del recurso de apelaci\'on, los requisitos de admisibilidad del recurso y los de fundabilidad?}

\medskip
\textbf{Respuesta:}
\begin{itemize}
    \item \textbf{Admisibilidad $\rightarrow$ juez de primera instancia:} analiza si el recurso fue interpuesto en el plazo legal, si la resoluci\'on es susceptible de apelaci\'on y si quien lo interpone es el legitimado. Lo hace al momento de conceder o denegar el recurso.
    \item \textbf{Fundabilidad $\rightarrow$ c\'amara:} analiza si la expresi\'on de agravios o el memorial constituyen una cr\'itica concreta y razonada, y decide si hace lugar o no al recurso (confirma, modifica o revoca la resoluci\'on).
\end{itemize}
\end{exercise}

\section{Caso pr\'actico: contestaci\'on de demanda declarada extempor\'anea}

\begin{exercise}{Contestaci\'on de demanda extempor\'anea}
\textbf{Supuesto:} se act\'ua como abogados del demandado en un proceso ordinario de da\~nos y perjuicios por mala praxis m\'edica. Se tienen 15 d\'ias para contestar demanda y se la contesta el d\'ia 16, dentro del plazo de gracia (las dos primeras horas del d\'ia, antes de las 9:30 a.\,m.). El d\'ia siguiente, el juzgado emite esta resoluci\'on:

\begin{quote}
\textit{``Buenos Aires, 22 de junio de 2020. DESGLOSESE por extempor\'aneo.''} Firma: Mar\'ia Eugenia G\'omez, Secretaria.
\end{quote}

El juzgado no tuvo en cuenta un feriado intermedio y por eso comput\'o mal el plazo. \textbf{¿Qu\'e se hace?}
\end{exercise}

Pueden plantearse las siguientes respuestas:

\begin{itemize}
    \item Interponer una revocatoria (art.~238 y ss.): respuesta incorrecta.
    \item Interponer revocatoria con apelaci\'on en subsidio: respuesta incorrecta.
    \item Interponer directamente un recurso de apelaci\'on: respuesta incorrecta.
\end{itemize}

\begin{important}
Todas son incorrectas porque la resoluci\'on \textbf{no est\'a firmada por el juez, sino por la secretaria del juzgado}. El secretario y el prosecretario act\'uan por delegaci\'on. S\'olo el juez tiene jurisdicci\'on y s\'olo el juez puede revocar por contrario imperio. La resoluci\'on no es susceptible ni de revocatoria ni de apelaci\'on. \textbf{La v\'ia correcta es la revocaci\'on nominada del art\'iculo 38 ter del CPCCN.}
\end{important}

\section{Revocaci\'on nominada: art\'iculo 38 ter CPCCN}

\begin{formula}{Art. 38 ter CPCCN}
\textit{Dentro del plazo de tres d\'ias desde que se notifican de la resoluci\'on dictada por el secretario o prosecretario administrativo, las partes pueden requerir al juez que deje sin efecto lo dispuesto por el secretario o prosecretario administrativo.}
\end{formula}

Diferencias con la revocatoria del art.~238:

\begin{table}[H]
\centering
\begin{tabularx}{\textwidth}{
    >{\raggedright\arraybackslash}p{0.25\textwidth}
    >{\raggedright\arraybackslash}X
    >{\raggedright\arraybackslash}X
}
\toprule
 & \textbf{Revocatoria (art.~238)} & \textbf{Revocaci\'on nominada (art.~38 ter)} \\
\midrule
\textbf{Plazo} & 5 d\'ias & \textbf{3 d\'ias} \\
\addlinespace
\textbf{Fundamentaci\'on} & Simult\'anea a la interposici\'on & Simult\'anea a la interposici\'on \\
\addlinespace
\textbf{Resoluci\'on que ataca} & Providencias simples dictadas por el juez & Resoluciones dictadas por secretario o prosecretario \\
\addlinespace
\textbf{Qui\'en resuelve} & El mismo juez & El juez (revisa lo dispuesto por el auxiliar) \\
\bottomrule
\end{tabularx}
\caption{Revocatoria y revocaci\'on nominada.}
\end{table}

\section{Reposici\'on \textit{in extremis}: creaci\'on pretoriana}

La reposici\'on \textit{in extremis} es un mecanismo impugnatorio \textbf{no regulado en el c\'odigo}, creado por la jurisprudencia. No debe confundirse ni con la revocatoria del art.~238 ni con la revocaci\'on del art.~38 ter.

\begin{important}
Requisitos de la reposici\'on \textit{in extremis}:
\begin{itemize}
    \item Se trata de una resoluci\'on que \textbf{no} es una providencia simple.
    \item No existe otro remedio ni recurso posible.
    \item El error es de \'indole \textbf{sustancial} (de fondo), no meramente formal.
    \item El error es \textbf{notorio}: de mantenerse, se consolidar\'ia una injusticia muy grande para alguno de los litigantes.
    \item Se interpone y se funda en el mismo acto. Plazo: 3 d\'ias.
\end{itemize}
\end{important}

\begin{example}{Sentencia inapelable por monto con error sustancial}
Una sentencia definitiva dictada en un proceso donde la demanda es de \$\,200.000 (menor al m\'inimo del art.~242) es inapelable. Si esa sentencia tiene un error sustancial en la fijaci\'on de los hechos que constituye una injusticia notoria, la parte afectada puede interponer la reposici\'on \textit{in extremis} dentro de los 3 d\'ias, fund\'andola en ese mismo acto.
\end{example}

\subsection{Jurisprudencia sobre la reposici\'on \textit{in extremis}}

\begin{note}
\textbf{C\'amara Civil, Sala E (sumario citado en las notas).} Las resoluciones del tribunal de alzada no son, en principio, susceptibles de recurso de reposici\'on, por no revestir el car\'acter de providencias simples. Este principio reconoce excepciones en circunstancias excepcionales: cuando media un error esencial en la apreciaci\'on de los antecedentes del caso y el mantenimiento de la situaci\'on conducir\'ia a un resultado injusto que es deber de los jueces modificar, para preservar a los litigantes.
\end{note}

La Corte Suprema de Justicia tambi\'en ha dicho que sus decisiones no son susceptibles de recurso de reposici\'on, aunque ese principio puede reconocer excepciones cuando se trata de situaciones serias e inequ\'ivocas que demuestran con nitidez manifiesta el error que se pretende subsanar.

\cornelltitle{ejercicios y mecanismos especiales de la apelación}
\cornellblocksum{
\cqrow{¿Qu\'e analiza el juez de primera instancia y qu\'e la c\'amara?}{El juez de primera instancia analiza la admisibilidad (plazo legal, resoluci\'on susceptible de apelaci\'on y legitimado) al conceder o denegar el recurso. La c\'amara analiza la fundabilidad: si la expresi\'on de agravios o el memorial son una cr\'itica concreta y razonada, y decide si confirma, modifica o revoca.
}
\cqrow{¿Por qu\'e no procede revocatoria ni apelaci\'on contra una resoluci\'on firmada por el secretario?}{Porque el secretario y el prosecretario act\'uan por delegaci\'on: s\'olo el juez tiene jurisdicci\'on y s\'olo \'el puede revocar por contrario imperio.
}
\cqrow{¿Qu\'e es la revocaci\'on nominada del art.~38 ter?}{Es el mecanismo por el cual, dentro de los 3 d\'ias de notificada la resoluci\'on dictada por el secretario o prosecretario, las partes piden al juez que la deje sin efecto. Se interpone y funda en el mismo acto.
}
\cqrow{¿Qu\'e diferencia hay entre el art.~38 ter y la revocatoria del art.~238?}{La revocatoria (art.~238) tiene plazo de 5 d\'ias y ataca providencias simples dictadas por el juez; la revocaci\'on nominada (art.~38 ter) tiene plazo de 3 d\'ias y ataca resoluciones del secretario o prosecretario. En ambas la fundamentaci\'on es simult\'anea a la interposici\'on.
}
\cqrow{¿Cu\'ando procede la reposici\'on \textit{in extremis}?}{Cuando la resoluci\'on no es una providencia simple, no existe otro remedio o recurso posible, el error es sustancial y notorio, y su mantenimiento consagrar\'ia una injusticia grande. Se interpone y funda en el mismo acto, en 3 d\'ias. Es de creaci\'on pretoriana y no est\'a regulada en el c\'odigo.
}
\cqrow{¿Qu\'e han dicho la C\'amara Civil (Sala E) y la Corte Suprema sobre la reposici\'on?}{Que sus resoluciones no son, en principio, susceptibles de reposici\'on, pero ese principio admite excepciones ante un error esencial o ante situaciones serias e inequ\'ivocas que demuestran con nitidez el error que se pretende subsanar.
}
}{la admisibilidad se analiza en primera instancia y la fundabilidad en la c\'amara. Contra resoluciones del secretario o prosecretario procede la revocaci\'on nominada del art.~38 ter (3 d\'ias, fundamentaci\'on simult\'anea). La reposici\'on \textit{in extremis}, de creaci\'on pretoriana, opera cuando no existe otro remedio, el error es sustancial y notorio y su mantenimiento implicar\'ia una injusticia grave. Identificar correctamente el tipo de resoluci\'on y qui\'en la firm\'o es la clave para elegir la v\'ia impugnativa correcta.
}

\section{Síntesis general de la impugnación y los recursos}

\begin{important}
Esta parte desarrolla la \textbf{teoría general de la impugnación}, que constituye el sistema mediante el cual el ordenamiento procesal civil argentino permite atacar las decisiones judiciales erróneas o los actos procesales viciados. El fundamento último es la falibilidad humana: los jueces y las partes pueden equivocarse, y ese error requiere mecanismos de corrección.
\end{important}

\cornelltitle{general de la impugnación y los recursos}
\cornellblock{
\cqrow{¿Qué es un acto de impugnación y cuál es la diferencia con un recurso?}{Los actos de impugnación son el género: todo mecanismo procesal dirigido a atacar un acto o procedimiento que lesione el interés de las partes o de terceros. Los recursos son una especie dentro de ese género. Otros tipos son el incidente de nulidad, la acción autónoma de revisión de cosa juzgada, el juicio ordinario posterior, etc.
}
\cqrow{¿Por qué existen los mecanismos de impugnación?}{Porque el error es inherente a la condición humana. Los jueces y las partes pueden equivocarse. Carnelutti decía que la amenaza del error pende como espada de Damocles en el proceso. La impugnación existe para mitigar esas deficiencias.
}
\cqrow{¿Cuál es la diferencia entre vicio \textit{in procedendo} y vicio \textit{in iudicando}?}{El vicio \textit{in procedendo} es un defecto en el procedimiento (por ejemplo, una notificación irregular) y se ataca por incidentes de nulidad o remedios. El vicio \textit{in iudicando} es un error en la decisión judicial (interpretación de la ley o de los hechos) y se ataca por recursos en sentido estricto, ante un tribunal superior.
}
\cqrow{¿Quiénes tienen legitimación para recurrir y qué es el agravio?}{Las partes, los terceros intervinientes, los ministerios públicos y las personas ajenas al proceso con interés legítimo afectado. El agravio es la insatisfacción total o parcial del postulado esgrimido; sin agravio no hay legitimación para recurrir.
}
\cqrow{¿Existe garantía constitucional de doble instancia en materia civil?}{No. En materia no penal no existe garantía constitucional de doble instancia. En materia penal sí, desde la incorporación del Pacto de San José de Costa Rica (art. 8º inc. 2º ap. h) a la Constitución vía art. 75 inc. 22 CN, ratificada por el caso Giroldi.
}
\cqrow{¿Qué es la \textit{reformatio in peius}?}{Prohíbe que el tribunal de segunda instancia empeore la situación del único apelante.
}
\cqrow{¿Qué es el recurso indiferente?}{Una categoría doctrinal no tipificada en el CPCCN que habilita a los tribunales a dejar de lado los requisitos formales de admisibilidad para garantizar derechos fundamentales gravemente afectados. Se aplica excepcionalmente y tiene sustento en los casos «YPF» y «Tejida» de la CSJN.
}
}
\medskip
\cornellblocksum{
\cqrow{¿Para qué sirve la aclaratoria y contra qué procede la revocatoria?}{La aclaratoria sirve para que el mismo juez corrija errores materiales, aclare conceptos oscuros o integre la resolución; no procede para modificar sustancialmente lo decidido. La revocatoria procede sólo contra providencias simples, causen o no gravamen irreparable (art. 238 CPCCN), y no contra sentencias interlocutorias ni definitivas.
}
\cqrow{¿Cuáles son los cinco principios de las nulidades procesales?}{1) Especificidad (art. 169, primer párrafo): debe haber una norma formal vulnerada. 2) Finalidad (art. 169, segundo párrafo): si el acto cumplió su fin, no hay nulidad. 3) Subsanación (art. 170): si el acto fue convalidado (no se plantea en cinco días), no procede. 4) Trascendencia (art. 172): debe existir indefensión real; sin indefensión, no hay nulidad. 5) Protección (art. 171): quien generó el vicio no puede invocarlo.
}
}{Esta parte desarrolla la teoría general de la impugnación. El género es el acto de impugnación (que abarca el incidente de nulidad, los recursos y otros medios atípicos) y la especie más frecuente son los recursos. La doctrina diferencia remedios (resueltos por el mismo juez, orientados a vicios procedimentales) y recursos en sentido estricto (sometidos a un tribunal superior, dirigidos a vicios en la decisión). Todo recurso requiere legitimación, agravio y plazo perentorio, y presenta dos dimensiones: admisibilidad formal y fundabilidad sustancial. La garantía de la doble instancia existe con rango constitucional sólo en materia penal. El tribunal de alzada tiene una doble limitación (los escritos constitutivos y los agravios expresados) y no puede perjudicar al único apelante. El recurso indiferente es una construcción jurisprudencial excepcional. La aclaratoria corrige errores formales o integra omisiones sin alterar la esencia de lo decidido; la revocatoria permite que el mismo juez revea providencias simples. El incidente de nulidad ataca vicios procedimentales y exige los cinco principios; sin indefensión real, no hay nulidad.
}

\partintro{Esta parte aborda la ejecución de sentencia en el proceso civil y comercial argentino (Libro Tercero del Código Procesal Civil y Comercial de la Nación). Se estudia su naturaleza jurídica como etapa del proceso de conocimiento, sus diferencias con el juicio ejecutivo, los requisitos para ejecutar una sentencia (que debe estar consentida o ejecutoriada), el procedimiento según el tipo de obligación (dar dinero líquido, ilíquido, hacer, no hacer o dar cosas), las excepciones admisibles conforme al artículo 506 del CPCCN y los títulos que habilitan este proceso especial de ejecución.

\textbf{Hilo conductor.} Así como en un partido de fútbol primero se define el reglamento, luego se juega y finalmente se ejecuta el resultado (el equipo ganador cobra el trofeo), el proceso judicial tiene una etapa de conocimiento donde se discute y un fallo que solo cobra vida útil cuando se ejecuta. El recorrido de esta parte corresponde a ese último tramo: cómo hacer efectiva la sentencia.

Luego del estudio de la ejecución de la sentencia, la parte desarrolla el juicio ejecutivo y el cumplimiento de la sentencia de remate.

\textbf{El juicio ejecutivo.} Se desarrolla el juicio ejecutivo como proceso de conocimiento sumario: concepto, diferencia con el proceso ordinario, títulos ejecutivos (completos e incompletos), preparación de la vía ejecutiva, embargo ejecutivo, excepciones admisibles, sentencia de trance y remate, apelación, costas y ejecuciones especiales.

\textbf{Cumplimiento de la sentencia de remate.} Se examina la etapa coercitiva del proceso: embargo ejecutorio, subasta de bienes muebles e inmuebles, designación del martillero, publicación de edictos, informes registrales, base de la subasta, acto de remate, aprobación, depósito del saldo, levantamiento de gravámenes, distribución del producido, sobreseimiento, nulidades y adquisición del título.

\begin{example}{del cheque no pagado}
Un vecino firmó un cheque para pagar una refacción pero luego no lo pagó. El acreedor no necesita iniciar un largo juicio ordinario: lleva directamente ese cheque al tribunal (juicio ejecutivo), el juez embarga de inmediato y, si el deudor no paga ni opone excepciones válidas, se dicta sentencia de trance y remate. Una vez firme esa sentencia comienza el cumplimiento de la sentencia de remate: el acreedor ejecuta los bienes del deudor ---por ejemplo, su departamento--- a través de una subasta judicial dirigida por un martillero, y el producido de esa subasta se aplica al crédito.
\end{example}

Ambos son los dos momentos de un mismo proceso: primero se discute, con limitaciones, si se debe, y luego se cobra por la fuerza.}
\part{Ejecución de la sentencia y juicio ejecutivo}

\chapter[Naturaleza y encuadre de la ejecución]{Naturaleza jurídica y encuadre de la ejecución de sentencia}\label{cap:ejec-naturaleza}

\section{Ubicación normativa en el Código Procesal}

La ejecución de sentencia se encuentra regulada en el Libro Tercero del CPCCN, que trata de los procesos de ejecución. Dentro de este libro se distinguen tres títulos:

\begin{itemize}
    \item \textbf{Título I:} ejecución de sentencias.
    \item \textbf{Título II:} juicio ejecutivo.
    \item \textbf{Título III:} ejecuciones especiales (hipotecaria, prendaria, fiscal y comercial).
\end{itemize}

\section{La ejecución de sentencia no es un proceso independiente}

Si bien el Código incluye a la ejecución de sentencia como un proceso de ejecución más, y por ello cierta doctrina la considera un proceso independiente del principal en el cual fue dictada la sentencia, la posición que aquí se sostiene es que la ejecución de sentencia es tan solo \textbf{una etapa más dentro del proceso de conocimiento} en el que fue dictada.

El proceso de conocimiento comprende las siguientes etapas:

\begin{enumerate}
    \item \textbf{Etapa introductoria:} demanda, contestación, reconvención y excepciones.
    \item \textbf{Etapa probatoria:} medios de prueba para formar convicción en el juez.
    \item \textbf{Etapa decisoria:} valoración de la prueba y dictado de la sentencia.
    \item \textbf{Etapa recursiva:} mecanismos impugnativos.
\end{enumerate}

Una vez agotada la etapa recursiva se accede a la ejecución. El hecho de que el Código la trate como un proceso de ejecución es lo que permite a la doctrina contraria sostener su independencia, pero para la posición aquí adoptada se trata de una etapa del proceso de conocimiento.

\section{Diferencias entre ejecución de sentencia y juicio ejecutivo}

La \textbf{intimación de pago} no es un requisito esencial en la ejecución de sentencia, mientras que sí lo es en el juicio ejecutivo, que es un proceso de ejecución obviamente independiente. La ejecución de sentencia no requiere intimación de pago porque ello se cubre con la notificación de la sentencia al demandado. En cambio, el juicio ejecutivo exige la intimación de pago porque es la que permite al ejecutado oponer excepciones para defenderse: al ser intimado de pago puede deducir las excepciones del artículo 544, mediante las cuales puede cuestionar el título base de la ejecución.

\begin{formula}{Art. 544 CPCCN}
Excepciones admisibles en el juicio ejecutivo. El ejecutado puede oponer, al ser intimado de pago, estas excepciones para cuestionar el título base de la ejecución: falsedad e inhabilidad de título, entre otras.
\end{formula}

Ambos procesos tienen en común que en ellos se ejecutan títulos, pero con notables diferencias. En el juicio ejecutivo se ejecutan \textbf{títulos ejecutivos}, que si bien documentan derechos, presumiblemente no tienen la certeza que posee el título ejecutorio, que es la sentencia.

\begin{table}[H]
\centering
\small
\begin{tabularx}{\textwidth}{
    >{\raggedright\arraybackslash}p{0.19\textwidth}
    >{\raggedright\arraybackslash}X
    >{\raggedright\arraybackslash}X
}
\toprule
\textbf{Aspecto} & \textbf{Juicio ejecutivo} & \textbf{Ejecución de sentencia} \\
\midrule
Tipo de título &
Título ejecutivo (derecho presunto). &
Título ejecutorio: sentencia (derecho cierto). \\[0.3em]

Intimación de pago &
Esencial. Habilita las excepciones del art.~544. &
No es requisito. La sentencia ya fue notificada. \\[0.3em]

Obligación que instrumenta &
Solo obligación de dar cantidad de dinero líquida o fácilmente liquidable. &
Obligaciones de dar, hacer o no hacer. \\[0.3em]

Embargo &
Optativo (puede trabarse o no). &
Absolutamente esencial para la ejecución. \\[0.3em]

Excepciones admisibles &
Art.~544 (más amplias, pueden abrirse a prueba). &
Art.~506 (sumamente limitadas, solo hechos posteriores a la sentencia). \\[0.3em]

Apelaciones durante el proceso &
Concedidas con trámite diferido hasta apelar la sentencia de remate. &
Concedidas con trámite diferido hasta apelar la sentencia de venta. \\
\bottomrule
\end{tabularx}
\caption{Diferencias entre el juicio ejecutivo y la ejecución de sentencia.}
\end{table}

\section{La sentencia como título ejecutorio: el derecho cierto}

La sentencia también constituye un título, pero un título que contiene un \textbf{derecho cierto}: ese derecho ha sido discutido con la suficiente amplitud en el juicio de conocimiento en el que se dictó la sentencia, y la \textbf{cosa juzgada} impide la discusión de hechos anteriores a ella.

Quien tiene a su favor una sentencia tiene, por ello, un título que contiene un derecho cierto. Por eso se lo llama \textbf{título ejecutorio}: trae aparejada una coacción inmediata para agredir el patrimonio del condenado que no cumple la sentencia. El condenado dispone de algunas defensas, pero sumamente limitadas, que se canalizan mediante la oposición de las excepciones que consagra el artículo 506 del CPCCN.

\begin{formula}{Art. 506 CPCCN}
Excepciones admisibles en la ejecución de sentencia. Solo se consideran legítimas: 1) falsedad de la ejecutoria; 2) prescripción de la ejecutoria; 3) pago; 4) quita, espera o remisión. En todos los casos, los hechos deben ser posteriores al dictado de la sentencia y probarse con constancias del expediente o documentos emanados del ejecutante.
\end{formula}

\cornelltitle{naturaleza y encuadre de la ejecución de sentencia}
\cornellblocksum{
\cqrow{¿Dónde se regula la ejecución de sentencia?}{En el Libro Tercero del CPCCN (procesos de ejecución), Título I. El Título II regula el juicio ejecutivo y el Título III las ejecuciones especiales (hipotecaria, prendaria, fiscal y comercial).
}
\cqrow{¿Por qué la ejecución de sentencia no es un proceso independiente?}{Porque es la culminación lógica del proceso de conocimiento: una etapa más dentro de él. Está íntimamente ligada a la sentencia que se ejecuta, es competente el mismo juez que conoció en la causa y la intimación de pago no es un requisito esencial, ya que se cubre con la notificación de la sentencia. El Código la trata como proceso de ejecución, y ello permite a la doctrina contraria sostener su independencia.
}
\cqrow{¿Qué diferencia al título ejecutorio del título ejecutivo?}{El título ejecutivo (juicio ejecutivo) contiene un derecho \textbf{presunto}, que puede discutirse con mayor amplitud (art.~544). El título ejecutorio (sentencia) contiene un derecho \textbf{cierto}, porque ya fue debatido con amplitud en el proceso de conocimiento. La cosa juzgada impide discutir hechos anteriores al dictado de la sentencia.
}
\cqrow{¿Qué papel cumple el embargo en cada proceso?}{En el juicio ejecutivo el embargo es optativo; en la ejecución de sentencia es absolutamente esencial.
}
\cqrow{¿Qué excepciones se admiten en cada proceso?}{En el juicio ejecutivo, las del art.~544 (más amplias, pueden abrirse a prueba). En la ejecución de sentencia, las del art.~506 (sumamente limitadas, solo hechos posteriores a la sentencia).
}
}{La ejecución de sentencia es una etapa del proceso de conocimiento y no un proceso autónomo. Ejecuta un título ejecutorio (la sentencia) que contiene un derecho cierto amparado por la cosa juzgada, a diferencia del juicio ejecutivo, que ejecuta títulos con derecho presunto, exige intimación de pago y admite las excepciones más amplias del art.~544. En la ejecución de sentencia el embargo es esencial y las excepciones se limitan a las del art.~506.
}

\chapter{Requisitos para ejecutar la sentencia}\label{cap:ejec-requisitos}

\section{La sentencia consentida}

La principal condición que debe cumplir la sentencia para ser ejecutada, conforme al artículo 499, es estar \textbf{consentida o ejecutoriada}.

\begin{definition}{Sentencia consentida}
Una sentencia está consentida cuando ha sido debidamente notificada a las partes y no ha sido impugnada. El consentimiento puede ser \textbf{expreso} (por ejemplo, mediante un escrito en el que se manifiesta que se consiente la sentencia) o \textbf{tácito} (cuando ha vencido el plazo para impugnarla sin que ninguna de las partes lo haya hecho).
\end{definition}

\begin{formula}{Art. 499 CPCCN}
Requisitos para ejecutar una sentencia. La sentencia debe estar consentida o ejecutoriada. Vencido el plazo de cumplimiento fijado en la sentencia ---o al día siguiente de quedar firme si el juez no fijó plazo especial--- y a pedido de parte, el juez procede a la ejecución.
\end{formula}

\section{La sentencia ejecutoriada}

\begin{definition}{Sentencia ejecutoriada}
Una sentencia está ejecutoriada cuando ha sido impugnada mediante los mecanismos impugnativos que admitía y ya no admite nuevos recursos.
\end{definition}
% TODO: contenido incompleto o ambiguo en las notas originales (la frase original sobre los recursos que ya no admite resulta ilegible; se adoptó la formulación del cuadro Cornell del apunte)

Tanto si ha sido consentida como si ha quedado ejecutoriada, la sentencia se encuentra \textbf{firme}.

\section{Plazo vencido, incumplimiento e instancia de parte}

Una vez firme la sentencia, y vencido el plazo para su cumplimiento, el juez procede a la ejecución a pedido de parte. El plazo puede ser fijado expresamente por el juez en la sentencia; si nada se dice al respecto, el cumplimiento es exigible al día siguiente de aquel en que la sentencia quedó firme.

Se requiere \textbf{inescindiblemente} el pedido de la parte vencedora para proceder a ejecutar la sentencia, ya que rige el \textbf{principio dispositivo}.

Además, para poder iniciar la etapa de ejecución la sentencia debe hallarse \textbf{incumplida} por el condenado. Si vencido el plazo la sentencia está prácticamente cumplida y se promueve la ejecución, esta será rechazada porque ha habido cumplimiento. Se entiende que hay cumplimiento aun cuando este se produzca por la aplicación de \textbf{astreintes} (multa progresiva a favor del vencedor que la parte puede reclamar para que el ejecutado cumpla su obligación); en ese caso no se trata de ejecución de sentencia sino de cumplimiento de la sentencia.

\begin{important}
Requisitos para ejecutar una sentencia: (1) sentencia consentida o ejecutoriada; (2) vencimiento del plazo de cumplimiento; (3) incumplimiento del condenado; (4) instancia de parte.
\end{important}

\section{Excepción: la ejecución provisoria}

Como excepción, el Código permite en algunas situaciones ejecutar la sentencia aunque no esté firme. Se trata de una \textbf{ejecución provisoria}, sujeta a lo que en definitiva suceda, generalmente por vías recursivas.

El principio es que todos los recursos de apelación se conceden con \textbf{efecto suspensivo} (apelada la sentencia, se suspende su cumplimiento), salvo que la ley establezca lo contrario. Existe la posibilidad de que la apelación se conceda con \textbf{efecto devolutivo}, y en esos casos la sentencia puede ejecutarse aunque esté apelada.

\subsection{Alimentos}

\begin{formula}{Art. 647 CPCCN}
Efecto del recurso en materia de alimentos. El recurso de apelación contra la sentencia que condena al pago de alimentos se concede con efecto devolutivo. Esto permite al alimentado ejecutar la sentencia de primera instancia aunque esté apelada, formándose incidente de apelación con el testimonio de la sentencia de primera instancia.
\end{formula}

El alimentado beneficiado con una sentencia que condena al alimentante a pagar alimentos puede ejecutarla aunque el alimentante la apele, dado que el recurso se concede con efecto devolutivo. Se protege así un derecho esencial: la necesidad del alimentado de contar con los alimentos para su subsistencia.

\subsection{Proceso sumarísimo}

\begin{formula}{Art. 498 CPCCN}
Proceso sumarísimo. El recurso de apelación contra la sentencia definitiva se concede con efecto devolutivo, salvo que su cumplimiento pudiere ocasionar al condenado un perjuicio irreparable, en cuyo caso el juez podrá concederlo con efecto suspensivo.
\end{formula}

Por lo tanto, la sentencia del juicio sumarísimo se cumple aunque esté apelada, salvo perjuicio irreparable para el condenado.

\subsection{Recurso extraordinario}

Una situación similar se presenta cuando se interpone recurso extraordinario contra una sentencia de cámara que confirma la sentencia de primera instancia. Si las dos sentencias son coincidentes, aunque la de segunda instancia haya sido objeto de recurso extraordinario, se permite la ejecución de la sentencia dando \textbf{fianza} de responder por si la Corte revocara luego la sentencia a través del recurso extraordinario.

\section{Ejecución parcial de la sentencia}

La segunda parte del artículo 499 permite ejecutar parcialmente una sentencia cuando una de sus partes ha sido consentida, aunque se haya interpuesto recurso respecto de otros rubros. Se puede ejecutar parcialmente la sentencia respecto de la parte que quedó firme.

No se trata de una excepción a la ejecutoriedad de la sentencia no consentida: lo que se permite es que, habiéndose cuestionado una parte del fallo, la otra parte, independiente y que puede cuantificarse separadamente, haya quedado firme y pueda llevarse adelante su ejecución.

\begin{example}{Accidente de tránsito}
El juez hizo lugar a los rubros daño emergente, lucro cesante y daño moral, y la parte demandada apela solamente los rubros relativos al daño moral y al lucro cesante. Firme la suma fijada por daño emergente, se puede proceder a ejecutar parcialmente la sentencia por los rubros que quedaron firmes.
\end{example}

Para ello se libra un \textbf{testimonio} en el que consta que determinado rubro, por determinado importe, ha quedado firme, de modo que pueda ser ejecutado en primera instancia, mientras el expediente sube a la cámara para conocer del recurso por los demás rubros cuestionados.

\cornelltitle{requisitos para ejecutar la sentencia}
\cornellblocksum{
\cqrow{¿Cuándo se considera consentida una sentencia?}{Cuando fue debidamente notificada y no fue impugnada en el plazo, o cuando fue consentida expresamente por escrito. El consentimiento puede ser expreso o tácito. % TODO: contenido incompleto o ambiguo en las notas originales (el cuadro Cornell del apunte agrega un tercer supuesto: apelación cuyo recurso se abandonó, por ejemplo por no fundarlo en tiempo; no figura en la transcripción original)
Se agrega, como tercer supuesto, que la sentencia fue apelada pero el recurso se abandonó (por ejemplo, por no fundarlo en tiempo).
}
\cqrow{¿Qué significa que la sentencia esté ejecutoriada?}{Que fue impugnada por todos los recursos que admitía y ya no admite nuevos embates: se agotaron los mecanismos impugnativos disponibles.
}
\cqrow{¿Cuáles son los requisitos para ejecutar una sentencia?}{1) Sentencia consentida o ejecutoriada. 2) Vencido el plazo de cumplimiento (fijado por el juez o el día siguiente a quedar firme). 3) Incumplimiento del condenado. 4) Instancia de parte. Rige el principio dispositivo: sin pedido de parte no hay ejecución.
}
\cqrow{¿Cuándo se puede ejecutar una sentencia que no está firme?}{Excepcionalmente, mediante ejecución provisoria, cuando la apelación se concede con efecto devolutivo: a) alimentos (art.~647 CPCCN); b) proceso sumarísimo (art.~498, salvo perjuicio irreparable); c) recurso extraordinario contra sentencia de cámara confirmatoria de la de primera instancia, dando fianza de responder si la Corte revoca la sentencia.
}
\cqrow{¿Cuál es la diferencia entre ejecución total y ejecución parcial?}{La ejecución parcial procede cuando solo una parte del fallo quedó firme y el resto fue apelado. Se libra testimonio acreditando que ese rubro quedó consentido, se ejecuta en primera instancia y el expediente sube a cámara por los rubros cuestionados. Por ejemplo: en un accidente de tránsito, el daño emergente queda firme mientras se apelan el lucro cesante y el daño moral.
}
}{Para ejecutar una sentencia debe estar consentida o ejecutoriada, haber vencido el plazo de cumplimiento, haber incumplido el condenado y existir pedido de parte. Excepcionalmente puede ejecutarse una sentencia no firme (alimentos, proceso sumarísimo, recurso extraordinario con doble sentencia coincidente y fianza), y puede ejecutarse parcialmente el rubro que quedó firme.
}

\chapter[Competencia y procedimiento de la ejecución]{Competencia y procedimiento general de la ejecución}\label{cap:ejec-procedimiento}

\section{Competencia: el juez que dictó la sentencia}

\begin{formula}{Art. 501 CPCCN}
Competencia para la ejecución. Inc.~1: el juez que dictó la sentencia es competente para entender en su ejecución. Inc.~2: si el objeto de la ejecución lo impone, puede encomendarse a un juez de otra competencia territorial la realización de determinadas diligencias, lo que no importa una excepción a las reglas de competencia. Inc.~3: es competente el juez del proceso principal cuando existan conexiones directas entre causas sucesivas.
\end{formula}

El artículo 501 establece que resulta competente el juez que dictó la sentencia. Ello se vincula con la idea de que el proceso no debe entenderse aislado de la ejecución de sentencia: el proceso principal tiene una etapa más, la de ejecución, y por lógica consecuencia es competente para entender en ella quien entendió en todo el proceso hasta el dictado de la sentencia.

\subsection{Delegación a un juez de otra competencia territorial}

El inciso 2 permite que, si así lo impone el objeto de la ejecución, se encomiende a un juez de otra competencia territorial una determinada diligencia para ejecutar la sentencia. Esto no es una excepción a las reglas de competencia, sino una encomienda que un juez de una jurisdicción hace a otro: una delegación de la competencia para una diligencia determinada.

\begin{example}{Inmueble en extraña jurisdicción}
Debe otorgarse la posesión de un inmueble ubicado en extraña jurisdicción, o llevarse adelante un desalojo en extraña jurisdicción. El juez que dictó la sentencia de desalojo y tiene asiento en Buenos Aires no puede trasladarse a ejecutarla a la ciudad de Corral; por delegación interviene el juez con competencia territorial en relación con el lugar del inmueble respecto del cual debe darse la tenencia a la parte actora.
\end{example}
% TODO: contenido incompleto o ambiguo en las notas originales (el nombre de la ciudad «Corral» figura así en el apunte)

\subsection{Conexión directa entre causas sucesivas}

El inciso 3 establece que es competente el juez del principal si median conexiones directas entre causas sucesivas.

\begin{example}{Desalojo y daños y perjuicios}
Se inició como proceso principal un desalojo y, como consecuencia de él, se inició luego un proceso de daños y perjuicios derivados de esa ocupación. El juez del desalojo puede entender en ambos.
\end{example}

\section{Condena a pagar cantidad líquida}

Debe diferenciarse si la sentencia condena al pago de una cantidad líquida o de una cantidad ilíquida.

\begin{definition}{Condena a cantidad líquida}
La sentencia condena al pago de una cantidad líquida cuando esta está expresada numéricamente o es fácilmente determinable mediante una simple operación aritmética.
\end{definition}

\begin{example}{Capital más interés}
Si la sentencia condena a pagar un capital de 500.000 pesos con más el interés del 40\,\% anual, la sentencia se entiende líquida aunque no esté expresado su monto final, dado que puede calcularse con una simple operación aritmética que incluso puede realizar el juez para determinar el monto.
\end{example}

\section{La traba del embargo}

El embargo es absolutamente esencial: la ejecución de sentencia pretende, ante el incumplimiento del condenado, individualizar bienes del deudor para proceder a su venta y realización efectiva, eventualmente a través de la subasta. Los tipos de embargo (preventivo, ejecutivo y ejecutorio) se distinguen en el capítulo~\ref{cap:embargo}.

El juez tiene todo el poder del Estado, \textbf{imperium}, para hacer cumplir coercitivamente sus decisiones. El obligado que no cumple voluntariamente es embargado; una vez afectado un bien de su patrimonio al cumplimiento de la sentencia mediante un embargo ejecutorio, y previos los trámites que correspondan para ejercer algún tipo de defensa, ese bien es subastado para pagar al acreedor vencedor.

\subsection{Bienes inmuebles}

Para pedir la traba de un embargo sobre un inmueble debe acompañarse un \textbf{informe de dominio} donde conste la titularidad registral en cabeza del deudor. Acreditada la titularidad dominial, el juzgado ordena trabar embargo por la cantidad líquida que consta en la sentencia, más la suma que presupueste el juzgado para responder a intereses y costas.

Si el inmueble está en la misma jurisdicción, el embargo se diligencia por \textbf{oficio} firmado por el secretario al Registro de la Propiedad Inmueble. Si está en extraña jurisdicción (por ejemplo, un inmueble en la ciudad de La Plata con embargo ordenado por un juez de la ciudad de Buenos Aires), se traba mediante el llamado \textbf{testimonio ley 22.172}.

\begin{formula}{Ley 22.172}
Comunicaciones entre jueces de distintas jurisdicciones. Establece las formas en las que se comunican los jueces de distintas jurisdicciones para llevar a cabo determinadas diligencias, en este caso la inscripción del embargo ante el Registro de la Propiedad Inmueble de otra provincia. El testimonio debe ser firmado por el juez y por el secretario, tener sello de agua y otras características establecidas en la ley.
\end{formula}

\subsection{Bienes muebles registrables}

Si se trata de un mueble registrable (por ejemplo, un automóvil, una embarcación o una aeronave), se embarga por \textbf{oficio} firmado por el secretario. Se acompaña un informe de dominio expedido por los registros, como el del automotor, que acredite la titularidad en cabeza del demandado; el juez ordena el embargo y el oficio para que se inscriba ante el registro correspondiente.

Si es en extraña jurisdicción, también se tramita por testimonio de la ley 22.172, firmado por el juez y el secretario, con constancia en el texto del testimonio de los recaudos que exige la ley.

\subsection{Bienes muebles no registrables}

Son muebles no registrables, por ejemplo, las máquinas y determinadas herramientas que no sean de uso indispensable según la profesión del deudor (las que sí lo sean serían inembargables). Se embargan mediante \textbf{mandamiento de embargo} firmado por el secretario. Debe combinarse la fecha con el oficial de justicia, encargado de llevarlo a cabo, quien se constituye en el domicilio del deudor o donde estén los bienes y traba el embargo mediante un acta que labra en el momento.

\section{Efectos del embargo}

El embargo concede \textbf{prioridad como primer embargante} para valerse del producido del bien embargado cuando se efectúe la subasta y se convierta en dinero. Importa una \textbf{indisponibilidad}, una afectación del bien al proceso.

\begin{important}
Que el bien esté embargado no significa que no pueda ser vendido. Si el deudor vende el bien indicando que está embargado, un tercero puede adquirirlo, pero deberá soportar las consecuencias del embargo, porque este persigue a la cosa aunque la transfiera el deudor. La indisponibilidad consiste en que no puede disponerse del bien como si estuviera libre de gravámenes.
\end{important}

Si, por el tipo de bien que deba subastarse, se quiere desapoderar directamente al deudor, debe pedirse otra medida además del embargo: el \textbf{secuestro} del bien, para trasladarlo y quitarlo de la esfera de dominio de su dueño.

\section{Citación de venta}

Una vez embargado un bien y afectado al proceso, corresponde continuar la ejecución citando de venta al deudor. Aunque en la etapa de ejecución el conocimiento del juez es mínimo, no se excluye cierto grado de conocimiento a través de las excepciones que el deudor puede oponer al progreso de la ejecución; de ahí la necesidad de la citación de venta.

La secuencia que establece el Código es: \textbf{primero embargar} (sin bienes embargados no hay sobre qué realizar la sentencia) y \textbf{después citar de venta} por cédula al deudor, para que comparezca a oponer excepciones dentro de \textbf{cinco días} desde que fue notificado. En la práctica a veces ambos actos se hacen conjuntamente: por ejemplo, mediante un mandamiento que traba embargo sobre bienes muebles y deja constancia de que se cita de venta al deudor para que oponga las excepciones del artículo 506.

Si el deudor no opone excepciones, el juez dicta directamente la sentencia de venta o resolución de venta mandando llevar adelante la ejecución, y se pasa a las reglas de cumplimiento de la sentencia de remate del juicio ejecutivo (artículos 559 y siguientes), que se tratan en el capítulo~\ref{cap:remate}.

\section{Las excepciones del artículo 506}

\subsection{Falsedad de la ejecutoria}

La falsedad de la ejecutoria supone que se está ejecutando una sentencia adulterada: porque la firma no pertenece al juez, o porque se enmendó (sobreraspado, adulterado) el monto que originalmente constaba en la sentencia.

No es común, pero podría darse cuando se pide un testimonio, que sale de la esfera del expediente y puede ser objeto de adulteración. No puede imponerse al ejecutado el cumplimiento de una obligación que no se corresponde con la que concretamente se dictaminó en la sentencia: debe haber una correspondencia perfecta entre lo que se ejecuta y lo decidido por sentencia firme en la causa.

\subsection{Prescripción de la ejecutoria}

Se aplica el plazo genérico del artículo 2560 del Código Civil y Comercial de la Nación (CCCN), de \textbf{cinco años}, por no existir uno específico. Antes era de diez años por el artículo 4023 del Código de Vélez; ahora se redujo a cinco.

\begin{formula}{Art. 2560 del CCCN}
Plazo genérico de prescripción: cinco años, excepto que esté previsto uno diferente en la legislación local. Es el plazo aplicable para la prescripción de la ejecutoria cuando no existe plazo específico.
\end{formula}

El vencedor no puede dejar pasar el plazo que quiera para ejecutar la sentencia: a partir del día en que resulta exigible el cumplimiento, tiene cinco años para ejecutarla, ya que si no realiza actos que impulsen el proceso de ejecución opera la prescripción de la ejecutoria.

\begin{important}
Aunque haya transcurrido el plazo de prescripción, el vencedor puede deducir la pretensión de ejecución. Si la otra parte no opone la excepción de prescripción de ejecutoria en el plazo que corresponde (cinco días), aunque estuviera vencido el plazo de prescripción liberatoria, la sentencia podrá ejecutarse eficazmente. Rige el principio dispositivo: la prescripción debe ser opuesta por la parte que se beneficiaría con ella y no puede ser decidida de oficio por el juez.
\end{important}

\subsection{Pago}

Si el deudor canceló la obligación que surge de la sentencia, puede defenderse acreditando el pago mediante documentos que emanen del vencedor.

\begin{important}
Todas las excepciones admisibles en la ejecución de sentencia deben fundarse en \textbf{hechos posteriores a la sentencia}, dado que los anteriores están cubiertos por la eficacia de la cosa juzgada y la preclusión impide volver sobre ellos. Además, hay una limitación en los medios de prueba: solo se permiten \textbf{constancias del expediente} o \textbf{documentos emanados del vencedor}, con exclusión de cualquier otro medio probatorio.
\end{important}

\subsection{Quita, espera y remisión}

El inciso 4 del artículo 506 se refiere a las excepciones de quita, espera o remisión.

\begin{itemize}
    \item \textbf{Quita:} después de dictada la sentencia, el vencedor realiza una quita del monto que se le debe. El ejecutado se excepciona acompañando el documento que da cuenta de la quita, y se reduce la ejecución y, eventualmente, el embargo.
    \item \textbf{Espera:} el vencedor otorga un plazo de cumplimiento al ejecutado y, no obstante, ejecuta la sentencia antes de su vencimiento. El ejecutado opone la excepción de espera; si se hace lugar, el ejecutante debe aguardar a que pase el plazo para poder ejecutar la sentencia.
    \item \textbf{Remisión:} consiste, por decirlo así, en el perdón de la deuda. Técnicamente, remisión es devolver al deudor el título donde consta la obligación, pero esto no es posible aquí porque existe un expediente. Si se perdona la deuda, se acompaña el instrumento correspondiente y el juez hace lugar a la excepción, ordena levantar el embargo.
\end{itemize}
% TODO: contenido incompleto o ambiguo en las notas originales (la explicación original de por qué no es posible devolver el título en la remisión resulta ambigua)

\begin{example}{Quita}
La sentencia es de 600.000 pesos por daños y perjuicios y existe un documento emanado del vencedor, posterior a la sentencia, que le hace una quita de 200.000 pesos. El ejecutado puede excepcionarse acompañando ese documento, y la ejecución y, eventualmente, el embargo se reducen a 400.000 pesos.
\end{example}

\section{Excepciones no enumeradas admitidas por la jurisprudencia}

El artículo 506 dice que solo se consideran legítimas las excepciones que enumera, pero la jurisprudencia entiende que no pueden dejarse de lado algunas otras, como la \textbf{inhabilidad de título}. Si el ejecutante pretende ejecutar a quien no es el demandado condenado, el pretendido ejecutado puede oponer la excepción de falta de legitimación pasiva, que se incluye dentro de la inhabilidad de título.

Lo mismo se ha permitido respecto de otras excepciones vinculadas con presupuestos procesales, como la \textbf{incompetencia}: si se va a ejecutar la sentencia ante un juez incompetente, el deudor puede excepcionarse aunque la excepción no esté señalada expresamente entre las del artículo 506.

\section{Trámite de las excepciones y sentencia de venta}

Opuesta alguna excepción, se da traslado al ejecutante por cinco días para que la conteste. Una vez contestada, el juez debe resolver. % TODO: contenido incompleto o ambiguo en las notas originales (el apunte atribuye «resolver» al ejecutado; por el contexto corresponde al juez)

\begin{itemize}
    \item \textbf{Hace lugar a las excepciones:} ordena levantar el embargo. Si es la excepción de pago o de prescripción de ejecutoria, levanta el embargo; si es la de quita, adecua el embargo en la forma que corresponda; si es la de espera, hace lugar a la excepción pero deja trabado el embargo para que el ejecutante pueda continuar cuando se verifique el plazo correspondiente.
    \item \textbf{Rechaza las excepciones:} dicta la sentencia de venta mandando llevar adelante la ejecución.
\end{itemize}

La sentencia de venta puede ser apelada por el ejecutado. El recurso se concede con \textbf{efecto suspensivo y trámite inmediato}, pero si el ejecutante otorga \textbf{fianza} para responder por daños y perjuicios, puede transformar el efecto de la apelación en devolutivo y ejecutar la sentencia no obstante estar apelada por el ejecutado.

Si el juez ordena levantar el embargo, puede apelar el ejecutante. Esa apelación va en relación, con efecto suspensivo y trámite inmediato. Una vez firme la resolución, se continúa con el trámite de cumplimiento de la sentencia de remate del juicio ejecutivo, que es una etapa común al juicio ejecutivo y a la ejecución de sentencia.

\cornelltitle{competencia y procedimiento general de la ejecución}
\cornellblock{
\cqrow{¿Qué juez es competente para ejecutar la sentencia?}{El juez que dictó la sentencia (art.~501 CPCCN), porque la ejecución es una etapa más del proceso principal. Puede encomendar diligencias determinadas a un juez de otra competencia territorial, lo que no es excepción a las reglas de competencia sino una delegación. Si hay conexión directa entre causas sucesivas, es competente el juez del proceso principal.
}
\cqrow{¿Cuándo una condena es a cantidad líquida?}{Cuando la cantidad está expresada numéricamente o es fácilmente determinable mediante una simple operación aritmética, aunque no esté expresado el monto final.
}
\cqrow{¿Por qué el embargo es esencial en la ejecución de sentencia?}{Porque el procedimiento tiene por finalidad afectar bienes del deudor para su posterior subasta. Sin bienes afectados al proceso no puede continuarse la ejecución. En el juicio ejecutivo el embargo es optativo; en la ejecución de sentencia es absolutamente esencial.
}
\cqrow{¿Cómo se traba el embargo según el tipo de bien?}{Inmueble (misma jurisdicción): oficio del secretario al Registro de la Propiedad, con informe de dominio. Inmueble (otra jurisdicción): testimonio ley 22.172, firmado por juez y secretario, con sello de agua. Mueble registrable (automóvil, embarcación, aeronave): oficio del secretario con informe de dominio del registro correspondiente. Mueble no registrable: mandamiento de embargo firmado por el secretario, diligenciado por oficial de justicia mediante acta.
}
\cqrow{¿Qué efectos tiene el embargo?}{Prioridad como primer embargante sobre el producido de la subasta e indisponibilidad (afectación del bien al proceso). El bien puede ser vendido, pero el adquirente soporta el embargo, que persigue a la cosa. Para desapoderar al deudor se requiere además el secuestro.
}
\cqrow{¿Qué es la citación de venta y cuándo se produce?}{Es el acto que habilita al ejecutado a oponer excepciones. Se practica después del embargo (aunque en la práctica puede hacerse en el mismo mandamiento) y se notifica por cédula. El ejecutado tiene cinco días para oponer las excepciones del art.~506 CPCCN.
}
}
\medskip
\cornellblocksum{
\cqrow{¿Cuáles son las excepciones del art.~506 CPCCN?}{1) Falsedad de la ejecutoria (adulteración de la sentencia). 2) Prescripción de la ejecutoria (cinco años, art.~2560 CCCN; debe ser opuesta por la parte, no de oficio). 3) Pago (documentado con constancias emanadas del vencedor). 4) Quita, espera o remisión (documentos posteriores a la sentencia emanados del ejecutante). Condición común: hechos posteriores a la sentencia; prueba solo por constancias del expediente o documentos del ejecutante.
}
\cqrow{¿Qué excepciones admite la jurisprudencia además de las del art.~506?}{Inhabilidad de título (falta de legitimación pasiva: ejecutar a quien no fue condenado) e incompetencia (ejecutar ante juez incompetente).
}
\cqrow{¿Qué pasa si el ejecutado no opone excepciones?}{El juez dicta directamente la sentencia de venta mandando llevar adelante la ejecución, y se aplican las reglas de cumplimiento de la sentencia de remate del juicio ejecutivo (arts.~559 y ss. CPCCN; capítulo~\ref{cap:remate}).
}
\cqrow{¿Cómo se apela la sentencia de venta?}{Con efecto suspensivo y trámite inmediato; si el ejecutante otorga fianza para responder por daños y perjuicios, puede transformarse en devolutivo y ejecutarse la sentencia pese a la apelación.
}
}{Es competente el juez que dictó la sentencia, con posibilidad de delegar diligencias en jueces de otra jurisdicción. Ante una condena líquida se traba embargo (según el tipo de bien) y luego se cita de venta al deudor, que tiene cinco días para oponer las excepciones del art.~506 (falsedad, prescripción, pago, quita, espera o remisión), fundadas en hechos posteriores a la sentencia y probadas solo con constancias del expediente o documentos del ejecutante; la jurisprudencia admite además inhabilidad de título e incompetencia. Sin excepciones, o rechazadas estas, se dicta la sentencia de venta, apelable con efecto suspensivo salvo fianza.
}

\chapter[Condena ilíquida y liquidación]{Condena a cantidad ilíquida y práctica de la liquidación}\label{cap:ejec-iliquida}

\section{Cuándo corresponde practicar liquidación}

Cuando la sentencia no establece un monto determinado que debe abonar el condenado, se plantea cómo proceder. Ocurre, por ejemplo, con los intereses (la sentencia fija la tasa pero no el monto), con los frutos que deben restituirse o con los daños y perjuicios cuyo derecho se reconoce al actor pero que, por distintas cuestiones procesales, el juez no puede fijar en el acto del fallo.

\begin{formula}{Art. 165 CPCCN}
Condena al pago de frutos, intereses o daños. El juez, al pronunciar sentencia, debe fijar el importe de los frutos, intereses, daños y perjuicios con arreglo a lo dispuesto por el artículo 165. Cuando ello no sea posible, el juez debe fijar las bases para que se practique la liquidación conforme a las pautas que establezca.
\end{formula}
% TODO: contenido incompleto o ambiguo en las notas originales (el texto del recuadro del apunte para el art. 165 se autorreferencia)

El artículo 165 exige en primer término que el juez determine la cantidad en forma líquida por estos rubros, pero reconoce que pueden existir impedimentos; en ese caso exige que al menos fije las pautas para que se establezca posteriormente en una liquidación.

\begin{important}
La liquidación debe estar en consonancia y ser correlativa con lo establecido en la sentencia. No pueden incluirse rubros ni irse más allá o por fuera de lo solicitado en las pretensiones articuladas al iniciar el proceso, pues sería nugatoria la ejecución que se llevara a cabo con base en una liquidación que incluyera conceptos indebidos.
\end{important}

\section{Quién practica la liquidación y en qué plazo}

\begin{formula}{Art. 503 CPCCN (liquidación)}
Liquidación. El acreedor podrá practicar la liquidación dentro del plazo de diez días contados desde que la sentencia sea exigible. Dará traslado a la otra parte por cinco días. Si hay conformidad, o si vencido el plazo el deudor no la impugna, el juez aprobará la liquidación presentada, disponiendo que se la tenga como base para el embargo y la subasta.
\end{formula}
% TODO: contenido incompleto o ambiguo en las notas originales (el apunte cita el art. 503 tanto para la liquidación como para la condena a hacer, sección 5.1)

Es una \textbf{facultad del vencedor} practicar la liquidación, para lo cual dispone de diez días contados desde que sea exigible el cumplimiento, es decir, desde que la sentencia quedó firme. Si el juez fijó un plazo especial, el plazo del vencedor para practicar la liquidación comienza al día siguiente del vencimiento del plazo para cumplir.

De la liquidación se da traslado a la otra parte, por cédula y por cinco días. Si la contraparte consiente expresamente o guarda silencio (consentimiento tácito, por no impugnar), el juez aprueba la liquidación y se continúa como en el proceso con cantidad líquida: se traba embargo sobre el monto de la liquidación y lo que eventualmente se presupueste para responder a intereses y costas.

\section{Impugnación de la liquidación}

\begin{formula}{Art. 504 CPCCN}
Impugnación de la liquidación. Si hubiere impugnación a la liquidación, se formará incidente siguiendo las normas de los artículos 178 y siguientes. Las partes deberán ofrecer la prueba de la que intenten valerse en sus respectivos escritos de impugnación y contestación.
\end{formula}

El vencido puede impugnar la liquidación porque entiende que no corresponden los rubros o porque, aun correspondiendo, están mal o defectuosamente calculados. Se forma entonces el incidente que exige el artículo 504, siguiendo las pautas del artículo 178 y siguientes. Las partes deben ofrecer la prueba de la que intenten valerse para fundar la impugnación o la contestación de la impugnación que haga el ejecutante. Vencidos los plazos y contestada la impugnación, el juez resuelve y, en su caso, aprueba la liquidación que corresponda.

\section{Debate doctrinal sobre la apelación de la liquidación aprobada}

Existen dos posiciones doctrinarias sobre el efecto con que se concede la apelación contra la resolución que aprueba la liquidación:

\begin{table}[H]
\centering
\small
\begin{tabularx}{\textwidth}{
    >{\raggedright\arraybackslash}p{0.17\textwidth}
    >{\raggedright\arraybackslash}X
    >{\raggedright\arraybackslash}X
}
\toprule
 & \textbf{Alsina} & \textbf{Palacio} \\
\midrule
Fundamento &
Por tratarse de un incidente, el recurso se concede en relación. &
Se aplica el art.~509, que dispone que todas las apelaciones en ejecución de sentencia se conceden con trámite diferido. \\[0.3em]

Efecto &
Suspensivo y trámite inmediato. &
Trámite diferido. \\[0.3em]

Consecuencias &
Se suspende el cumplimiento de la resolución y el expediente va a la cámara para su resolución; una vez confirmada, vuelve a primera instancia y permite trabar embargo y continuar la ejecución. &
La liquidación aprobada, aunque apelada por el ejecutado, sirve de base para trabar el embargo y continuar la ejecución; la apelación se funda recién junto con una eventual apelación contra la sentencia de venta. \\
\bottomrule
\end{tabularx}
\caption{Posiciones doctrinarias sobre el recurso contra la liquidación aprobada.}
\end{table}

Esta discusión varía fundamentalmente la solución del caso.

\cornelltitle{condena ilíquida y liquidación}
\cornellblocksum{
\cqrow{¿Qué sucede cuando la sentencia condena a una cantidad ilíquida?}{Corresponde practicar liquidación (art.~503 CPCCN). El vencedor tiene diez días desde que la sentencia es exigible. Se da traslado a la otra parte por cinco días. Si hay conformidad o silencio, el juez aprueba la liquidación y se sigue con el trámite de embargo. Si hay impugnación, se forma incidente (art.~504, reglas del art.~178 y ss.).
}
\cqrow{¿Qué debe hacer el juez según el art.~165 cuando no puede fijar el monto?}{Debe fijar las bases o pautas para que se practique la liquidación posteriormente. La liquidación debe ser correlativa con la sentencia y con las pretensiones articuladas al iniciar el proceso.
}
\cqrow{¿Cuál es el debate sobre la apelación de la liquidación aprobada?}{Alsina: al ser un incidente, se concede en relación, con efecto suspensivo y trámite inmediato. Palacio: rige el art.~509 y se concede con trámite diferido, de modo que la liquidación sirve de base para el embargo y la apelación se funda junto con la apelación de la sentencia de venta.
}
}{Si la condena es ilíquida, el vencedor practica la liquidación en diez días desde que la sentencia es exigible; se corre traslado por cinco días y, si hay conformidad o silencio, se aprueba y se sigue con el embargo; si hay impugnación se tramita un incidente (arts.~503 y 504). Sobre el efecto de la apelación de la liquidación aprobada, Alsina y Palacio sostienen soluciones opuestas.
}

\chapter{Obligaciones de hacer, no hacer y dar cosas}\label{cap:ejec-hacer}

\section{Condena a hacer alguna cosa}

\begin{formula}{Art. 503 CPCCN (condena a hacer)}
Si la condena es a hacer alguna cosa y el obligado no la realiza dentro del plazo correspondiente, el ejecutante puede optar por: a) que la cosa sea realizada por un tercero a costa del ejecutado; o b) que se le paguen los daños y perjuicios provenientes de la inejecución.
\end{formula}
% TODO: contenido incompleto o ambiguo en las notas originales (el apunte cita el art. 503 también para la liquidación, sección 4.2)

Si el deudor cumple haciendo aquello a lo que fue obligado dentro del plazo correspondiente, no hay ejecución. Si no lo hace, el ejecutante puede promover la ejecución de sentencia, optando conforme al artículo 503. Debe advertirse que el deudor no puede ser obligado por la fuerza física a hacer algo contra su voluntad. El vencedor, ejecutante, puede elegir entre solicitar que la prestación sea realizada por un \textbf{tercero} o que se obligue al deudor a resarcir los \textbf{daños y perjuicios} provenientes de la inejecución.

Si opta por que lo realice un tercero, hay que estimar el costo de esa actividad (no una liquidación sino una \textbf{estimación}) y dar traslado al vencido para que se expida sobre si ese valor se corresponde con el objeto de la obligación que ahora realizará un tercero a su costa. La opción es del acreedor, pero la actividad es a costa del ejecutado.

\section{Condena a otorgar escritura pública}

\begin{formula}{Art. 502 CPCCN}
Condena a otorgar escritura pública. Si la sentencia condena al otorgamiento de escritura pública y este no es cumplido por el ejecutado, el juez la suscribirá a su costa, sin perjuicio de que si llegara a devenir de imposible cumplimiento por alguna circunstancia, esa escrituración se transforme en la obligación de resarcir los daños y perjuicios provenientes de esa imposibilidad.
\end{formula}

Si la sentencia condena al otorgamiento de escritura pública y el ejecutado no la cumple, el procedimiento continúa con quien fue designado para suscribirla a su costa. Si el cumplimiento deviniera imposible, la condena escrituraria se transforma en la obligación de resarcir los daños y perjuicios derivados de esa imposibilidad, que se fijan por el mismo procedimiento que establece el artículo 513.

\section{Condena a no hacer}

Cuando se trata de una condena a no hacer y el ejecutado la incumple, es decir, hace algo contrario a lo que estaba obligado, el acreedor puede ejercer la opción entre solicitar que se \textbf{reponga el estado de cosas anterior} o que se traduzca en la \textbf{indemnización de daños y perjuicios}, conforme al artículo 513.

\begin{formula}{Art. 513 CPCCN}
Daños y perjuicios por inejecución. Los daños y perjuicios derivados del incumplimiento se determinarán conforme a las pautas de los artículos 503 y 504 (liquidación con vía incidental) o bien a través de juicio ordinario o sumarísimo, según lo determine el juez de manera irrecurrible.
\end{formula}

Los daños y perjuicios se determinan, entonces, mediante el procedimiento de liquidación, con la derivación que hace el Código a las normas de los incidentes, o por juicio ordinario o sumarísimo, según lo determine el juez, y en forma irrecurrible.

\section{Condena a dar cosas: mandamiento de desapoderamiento}

Si la sentencia no se cumple voluntariamente, pueden aplicarse \textbf{astreintes}, es decir, multas progresivas conminatorias a favor de la otra parte para obligar a cumplir al ejecutado.

Si la obligación de dar alguna cosa no es cumplida voluntariamente, el vencedor puede ejecutar la sentencia mediante un \textbf{mandamiento} que ordena librar el juez para \textbf{desapoderar}, aun contra la voluntad del deudor, al vencido de aquello que fue ordenado por sentencia que se restituya al vencedor. En el mismo momento en que se diligencia el mandamiento para desapoderar de la cosa al ejecutado, se lo cita para que oponga las excepciones del artículo 506.

Si no puede cumplirse la entrega de la cosa por alguna circunstancia (porque ha desaparecido o es imposible su cumplimiento), el vencedor puede solicitar que se compense con una suma equivalente al valor de la cosa. La determinación de cuánto vale la cosa y cuánto debe compensarse al ejecutante se realiza también por las pautas de los artículos 503 y 504, por la vía incidental o por juicio ordinario o sumarísimo, según establezca el juez y de manera irrecurrible.

\section{Liquidaciones especiales}

\begin{formula}{Art. 516 CPCCN}
Liquidaciones especiales. Cuando las liquidaciones impliquen cuentas complicadas de lenta o difícil justificación o que requieran conocimientos especiales, podrán someterse a la decisión de peritos árbitros o amigables componedores. El laudo que emitan será vinculante y obligatorio para las partes. Lo mismo se aplica para liquidaciones de sociedades conyugales o determinación del carácter propio o ganancial de bienes.
\end{formula}

Cuando las liquidaciones implican cuentas complicadas de lenta o difícil justificación, o que requieren conocimientos especiales, pueden dejarse a la decisión de peritos árbitros o, en su caso, de amigables componedores. Tanto el perito árbitro como el amigable componedor emiten un laudo que es vinculante y obligatorio para las partes, y que debe respetar incluso el juez en estas cuestiones técnicas fijadas por los peritos en el marco de la competencia que se les atribuye.

\section{Otros títulos ejecutables por este procedimiento}

\begin{formula}{Art. 500 CPCCN}
Títulos ejecutables por el procedimiento de ejecución de sentencia: 1) las sentencias de los tribunales judiciales y los laudos de tribunales arbitrales (aunque estos carecen de imperium); 2) las transacciones o acuerdos homologados; 3) (inciso incorporado por ley 26.589) los acuerdos celebrados en mediación con la firma certificada del mediador, salvo que incluyan menores o incapaces, en cuyo caso requieren homologación judicial; 4) las multas procesales; 5) los honorarios regulados judicialmente.
\end{formula}
% TODO: contenido incompleto o ambiguo en las notas originales (el apunte menciona en la explicación oral el «inciso cuarto» para los acuerdos de mediación y el «inciso segundo» para las multas procesales, lo que no coincide con la numeración del recuadro)

\begin{itemize}
    \item \textbf{Sentencias judiciales y laudos arbitrales:} los tribunales arbitrales tienen jurisdicción pero carecen de imperium para hacer cumplir coercitivamente sus decisiones.
    \item \textbf{Transacciones o acuerdos homologados:} las transacciones son las que realizan las partes extinguiendo obligaciones litigiosas o dudosas; los acuerdos homologados pueden provenir, por ejemplo, de una conciliación llevada a cabo por el juez en la audiencia preliminar.
    \item \textbf{Acuerdos de mediación con firma certificada del mediador:} incorporados por la reforma de la ley 26.589. Su incumplimiento permite saltar todas las etapas del proceso de conocimiento y acceder directamente a la última, la ejecución de sentencia, para exigir el cumplimiento según las pautas ya vistas. Si el acuerdo incluye menores o incapaces, su representante debe solicitar al juez la homologación, con intervención del Ministerio Pupilar, para proteger esos intereses.
    \item \textbf{Multas procesales:} por ejemplo, las impuestas por temeridad y malicia a favor de una de las partes, que esta puede ejecutar por este proceso.
    \item \textbf{Honorarios regulados judicialmente:} su cobro también se ejecuta por este procedimiento.
\end{itemize}

\begin{note}
Paradójicamente, el acuerdo conciliatorio celebrado ante el mismo juez en la audiencia preliminar requiere homologación judicial, mientras que el acuerdo celebrado ante el mediador, con firma certificada, puede ejecutarse sin homologación (salvo que incluya menores o incapaces).
\end{note}

\section{Reflexión final}

El Código establece un procedimiento de ejecución de sentencia \textbf{abreviado}, limitado en cuanto a la apelación y a las defensas, dada la trascendencia del título ejecutorio, que instrumenta un derecho cierto. Permite transitar una etapa con un mínimo de conocimiento y con una mínima posibilidad de prueba por parte del ejecutado, y hacer efectiva la sentencia de modo de satisfacer el derecho que pretendía el actor al iniciar la demanda, cuya pretensión no se agota en una mera sentencia declarativa de la certeza de su derecho, sino que requiere efectivo cumplimiento.

Se recuerda, además, la afirmación de un autor según la cual, para tener derecho, hay que contar con la norma que lo acuerde, saber exponerlo, un juez que quiera reconocerlo y un deudor que pueda pagarlo.

\cornelltitle{hacer, no hacer, dar cosas y otros títulos}
\cornellblock{
\cqrow{¿Qué ocurre con la condena a hacer algo?}{El ejecutante puede optar (art.~503): a) que un tercero realice la cosa a costa del ejecutado (se estima el costo y se da traslado al vencido); b) el pago de daños y perjuicios por inejecución (se cuantifican por liquidación/incidente o juicio ordinario/sumarísimo, según determine el juez, de manera irrecurrible). El deudor no puede ser obligado por la fuerza física a hacer.
}
\cqrow{¿Qué ocurre con la condena a otorgar escritura pública?}{Si el ejecutado no la cumple, el juez la suscribe a su costa (art.~502). Si deviniera de imposible cumplimiento, se transforma en la obligación de resarcir los daños y perjuicios, fijados por el procedimiento del art.~513.
}
\cqrow{¿Qué pasa si la condena es a no hacer algo y se incumple?}{El acreedor puede optar (art.~513): a) la reposición al estado anterior; b) la indemnización por daños y perjuicios, cuantificados por liquidación/incidente o juicio ordinario/sumarísimo según determine el juez de manera irrecurrible.
}
\cqrow{¿Cómo se ejecuta la condena a dar cosas?}{Pueden aplicarse astreintes. Si no se cumple voluntariamente, mediante mandamiento de desapoderamiento librado por el juez, al tiempo que se cita al ejecutado para oponer las excepciones del art.~506. Si la entrega es imposible, se compensa con una suma equivalente al valor de la cosa, determinada por los arts.~503 y 504 o por juicio ordinario o sumarísimo.
}
\cqrow{¿Cómo se resuelven las liquidaciones especiales (art.~516)?}{Cuando implican cuentas complicadas o conocimientos especiales, pueden someterse a peritos árbitros o amigables componedores, cuyo laudo es vinculante y obligatorio para las partes e incluso para el juez.
}
\cqrow{¿Qué títulos habilitan la ejecución de sentencia además de la sentencia judicial?}{Art.~500 CPCCN: 1) laudos arbitrales (el árbitro carece de imperium); 2) transacciones y acuerdos homologados; 3) acuerdos de mediación con firma certificada del mediador (ley 26.589), salvo menores o incapaces, que requieren homologación; 4) multas procesales; 5) honorarios regulados judicialmente.
}
}
\medskip
\cornellblocksum{
\cqrow{Paradoja del art.~500: acuerdo judicial vs.\ acuerdo de mediación}{El acuerdo conciliatorio celebrado ante el mismo juez en la audiencia preliminar requiere homologación judicial para ser ejecutable. En cambio, el acuerdo de mediación con firma certificada del mediador no la requiere (salvo menores o incapaces).
}
}{Según el tipo de obligación, la ejecución adopta distintos carriles: en la condena a hacer, el ejecutante opta entre un tercero a costa del deudor o los daños y perjuicios; en la escritura pública, el juez suscribe; en la condena a no hacer, entre la reposición al estado anterior o la indemnización; en la condena a dar cosas, el mandamiento de desapoderamiento, con compensación en dinero si la entrega es imposible. Además de la sentencia, se ejecutan por este procedimiento laudos, transacciones o acuerdos homologados, acuerdos de mediación certificados, multas procesales y honorarios regulados.
}

\section{Síntesis de la ejecución de sentencia}
La ejecución de sentencia es la etapa final del proceso de conocimiento ---no un proceso autónomo---, mediante la cual el Estado pone su imperium al servicio del vencedor para hacer efectivo un derecho declarado cierto por sentencia firme.

A diferencia del juicio ejecutivo, que ejecuta títulos con derechos presuntos y admite defensas más amplias (art.~544 CPCCN), la ejecución de sentencia opera sobre un título ejecutorio que ya superó el debate exhaustivo del proceso de conocimiento: la cosa juzgada material impide discutir hechos anteriores. Por eso las excepciones admisibles (art.~506 CPCCN) se reducen a la falsedad de la ejecutoria, la prescripción, el pago y la quita, espera o remisión, todas fundadas en hechos \textbf{posteriores} a la sentencia y probables solo con documentos del expediente o del ejecutante.

El procedimiento requiere que la sentencia esté consentida o ejecutoriada, que haya vencido el plazo de cumplimiento y que exista incumplimiento del condenado, todo a instancia de parte. El embargo es esencial: traba bienes para su posterior subasta. Según el tipo de obligación condenada (dar dinero, hacer, no hacer, dar cosas), el proceso adopta distintos carriles: liquidación (si la condena es ilíquida), tercero a costa del deudor o indemnización (si la condena es a hacer o no hacer) y mandamiento de desapoderamiento (si es a dar cosas).

Otros títulos equiparados a la sentencia ---laudos arbitrales, transacciones homologadas, acuerdos de mediación certificados, multas procesales y honorarios--- también se ejecutan por este procedimiento, lo que subraya su carácter de herramienta central en la tutela jurisdiccional efectiva.

\chapter{El juicio ejecutivo}\label{cap:juicioejecutivo}

\chaptermark{El juicio ejecutivo}

%----------------------------------------------------------
\section{Ubicación en el código procesal}

El juicio ejecutivo se localiza en la parte del código referida al proceso de ejecución. Queda integrado, dentro de esa parte, a continuación del tema de la ejecución de sentencias.

El código agrupa tanto la ejecución de sentencias ---el cumplimiento de la decisión de un juicio ordinario de conocimiento pleno--- como el juicio ejecutivo, a partir del artículo 520. Ambos mecanismos comparten luego la regulación referida a la realización de los bienes del deudor.

\begin{quote}
\textbf{Art. 520 CPCCN.} A partir de este artículo el código regula el juicio ejecutivo.
\end{quote}

\begin{important}
El juicio ejecutivo no debe confundirse con el proceso de ejecución de sentencias, ya que responden a una lógica totalmente distinta.
\end{important}

%----------------------------------------------------------
\section{Concepto y naturaleza jurídica}

\subsection{El proceso de conocimiento sumario}

La ejecución de sentencia tiene como antecedente un conocimiento pleno, es decir, exhaustivo en relación con la causa de la obligación. En el juicio ejecutivo, en cambio, se está ante un \textbf{juicio de conocimiento sumario}, de conocimiento acotado o, si se quiere, superficial.

% TODO: contenido incompleto o ambiguo en las notas originales (el apunte menciona "el de la provincia de González"; se interpreta como referencia al código procesal provincial)
No debe confundirse esta noción con el proceso sumario previsto, por ejemplo, en la provincia de Buenos Aires, o el que antes existía en el código procesal civil y comercial. Cuando el código nacional o el código provincial hablan de juicio sumario, se refieren a un juicio de conocimiento pleno. En cambio, cuando se afirma que el juicio ejecutivo es un proceso de conocimiento sumario, se indica que el conocimiento del juez sobre el conflicto está acotado, disminuido o seccionado: no va más allá del título cartular que se lleva a ejecutar.

\subsection{La opción por la rapidez}

El juicio ejecutivo es un proceso en el cual se opta por la rapidez en detrimento de la exhaustividad en el conocimiento. Su razón de ser es el tiempo: el legislador quiso que, en ciertos casos en que existen determinadas prestaciones, el acreedor tenga a su alcance ---ante el incumplimiento del deudor--- una herramienta rápida para recobrar el crédito o las sumas de dinero adeudadas, sin tener que someterse, como regla, a un juicio de conocimiento pleno.

En este proceso no hay traslado de demanda, excepciones como las del juicio ordinario, apertura a prueba, alegatos ni sentencia en los términos del ordinario. Tampoco rige la regla de que todas las decisiones son apelables: sucede lo contrario. Es un proceso al cual sólo se puede acceder ante la presencia de un título ejecutivo.

%----------------------------------------------------------
\section{Diferencia con la ejecución de sentencias}

\begin{table}[H]
\centering
\renewcommand{\arraystretch}{1.25}
\begin{tabularx}{\textwidth}{Y Y}
\toprule
\textbf{Ejecución de sentencias} & \textbf{Juicio ejecutivo} \\
\midrule
Requiere un juicio ordinario previo con conocimiento pleno. & Sólo requiere un título ejecutivo. \\
La sentencia es el antecedente del trámite ejecutorio. & El título cartular es el fundamento directo. \\
El deudor pudo discutir exhaustivamente la causa. & La causa de la obligación no puede discutirse en el proceso. \\
La cosa juzgada es material para todas las cuestiones. & La cosa juzgada es formal respecto de las cuestiones causales. \\
El embargo resultante es «ejecutorio». & El embargo resultante es «ejecutivo». \\
\bottomrule
\end{tabularx}
\end{table}

\subsection{El desdoblamiento del conocimiento del conflicto}

El juicio ejecutivo desdobla el conocimiento del conflicto y opta por un mecanismo procesal ágil y rápido mediante el cual ---con menos cantidad de trámites--- el acreedor, a partir de la existencia de un título que traiga aparejada fuerza ejecutiva, puede hacerse más rápidamente del crédito, en detrimento de las defensas que puede oponer el deudor. Esas defensas sólo pueden referirse a cuestiones que emerjan del título cartular que se ejecuta.

Para resguardar el derecho de defensa del deudor, se le concede la posibilidad de discutir todas las cuestiones que quedan fuera del juicio ejecutivo en un \textbf{juicio ordinario posterior}, en el cual la discusión ya no está acotada a las cuestiones que surjan del título.

\subsection{La cosa juzgada en el juicio ejecutivo}

En alguna medida, se dice que la cosa juzgada que emana del juicio ejecutivo es una \textbf{cosa juzgada formal}: lo es porque existen cuestiones causales que no pueden debatirse en él. Sin embargo, respecto de las cuestiones que sí pueden debatirse en el juicio ejecutivo ---por ejemplo, el pago--- no se trata de cosa juzgada formal, porque el pago es una excepción que sí puede interponerse en el marco del ejecutivo. Para el pago, entonces, el juicio ejecutivo arroja una \textbf{cosa juzgada material}: no podrá discutirse en el ordinario posterior una cuestión que sí podía plantearse en el ejecutivo.

%----------------------------------------------------------
\section{Requisitos para acceder a la vía ejecutiva}

A partir del artículo 523, el código enumera los títulos que tienen fuerza ejecutiva por ley. Para acceder a la vía ejecutiva se requieren dos condiciones esenciales:

\begin{enumerate}
    \item Contar con un título que traiga aparejada la ejecución (fuerza ejecutiva).
    \item Que exista una obligación exigible de dar cantidades líquidas de dinero o fácilmente liquidables.
\end{enumerate}

\subsection{La exigibilidad}

La obligación debe tener un plazo vencido; es decir, debe ser exigible, de modo que el acreedor pueda pedir el cobro. Debe tratarse de una obligación de dar sumas de dinero y no de una obligación de hacer o de no hacer. El juicio ejecutivo queda, por el momento, limitado a obligaciones de dar sumas de dinero.

\subsection{La liquidez o fácil liquidabilidad}

Una suma es \textbf{líquida} cuando la cantidad de dinero adeudada está expresada en números. El código también admite sumas \textbf{fácilmente liquidables}: por ejemplo, si se trata de cánones mensuales de alquileres, la suma líquida puede obtenerse mediante una liquidación sencilla en el marco del proceso ejecutivo.

\subsection{El caso de la moneda extranjera}

Una situación bastante común se presenta en las hipotecas pactadas en moneda extranjera. En Argentina, los dólares tuvieron distintos tratamientos jurídicos a lo largo del tiempo: en un momento se los consideró una obligación de dar cosas y, en otro, una obligación equivalente a dar moneda corriente.

% TODO: contenido incompleto o ambiguo en las notas originales (el apunte indica que actualmente se la trata "como una obligación de dar" sin precisar de qué clase)
Actualmente se la trata como una obligación de dar, por lo que, conforme el artículo 520, corresponde convertir la moneda extranjera a su equivalente nacional para obtener la suma líquida por la cual se puede ejecutar.

%----------------------------------------------------------
\section{Títulos ejecutivos --- Artículo 523}

El artículo 523 realiza una larga enumeración de supuestos en los que se puede acceder a la vía ejecutiva. La vía ejecutiva es siempre una opción para el acreedor: podría optar directamente por el juicio ordinario (art.~521).

\begin{quote}
\textbf{Art. 523 CPCCN --- Títulos con fuerza ejecutiva:}
\begin{enumerate}[label=\arabic*.º),leftmargin=2.2em]
    \item El instrumento público presentado en forma.
    \item El instrumento privado suscripto por el obligado, reconocido judicialmente, o cuya firma estuviese certificada por escribano con intervención del obligado y registrada la certificación en el protocolo.
    \item La confesión de deuda líquida y exigible prestada ante el juez competente para conocer en la ejecución.
    \item La cuenta aprobada o reconocida como consecuencia del procedimiento establecido en el artículo 525.
    \item La letra de cambio, factura de crédito, cobranza bancaria de factura de crédito, vale o pagaré, el cheque y la constancia de saldo deudor en cuenta corriente bancaria, de acuerdo con las condiciones establecidas por las leyes en la materia.
    \item El crédito por alquileres o arrendamiento de muebles.
    \item Los demás títulos que tuvieran fuerza ejecutiva por ley y no estén sujetos a un procedimiento especial.
\end{enumerate}
\end{quote}

\subsection{El inciso 7.º: apertura a títulos especiales}

El código abre la posibilidad de acceder a la vía ejecutiva a cualquier otro título al que una ley especial le dé fuerza ejecutiva. Es el caso, por ejemplo, del cobro de algunas matrículas en ciertos colegios profesionales, las deudas fiscales y otros créditos que gozan de este beneficio: la entidad emite un certificado de deuda ---por ejemplo, las matrículas de los abogados en los colegios profesionales--- y ese instrumento constituye un título ejecutivo que emana directamente del acreedor.

También tiene fuerza ejecutiva el crédito por expensas ---las llamadas \textbf{ejecuciones de expensas}---, que se conforma, por ejemplo, con el certificado de deuda suscripto por el administrador del consorcio.

%----------------------------------------------------------
\section{Títulos incompletos y preparación de la vía ejecutiva}

Dentro de la enumeración de los títulos ejecutivos se encuentran los denominados \textbf{títulos incompletos}: instrumentos que no tienen por sí mismos fuerza ejecutiva, sino que requieren algún trámite previo para lograr la \textbf{preparación de la vía ejecutiva}.

\begin{example}{el contrato de locación}
Cuando el locador quiere ejecutar los alquileres pactados en un contrato de locación, del título no surge cuál es la suma líquida que pretende ejecutar, porque al firmarse el contrato no se sabía en qué mes el deudor dejaría de pagar el canon locativo. Por eso, de modo previo, se cita al inquilino al tribunal para que acompañe el último recibo de pago del alquiler: si no lo trae, se tienen por adeudados todos los posteriores y, con una cuenta sencilla, se obtiene la suma fácilmente liquidable.
\end{example}

Si el contrato de locación tampoco tiene la firma certificada por escribano, el inquilino deberá reconocer o desconocer la firma que se le atribuye. La sanción por no concurrir es que se lo tendrá por reconocido. Si concurre y niega la firma, y luego se prueba mediante pericial que la firma le pertenece, le corresponde una multa del 30\,\% más de la deuda que se reclama.

%----------------------------------------------------------
\section{El embargo ejecutivo}

Cuando se tiene un título completo con fuerza ejecutiva ---o un título incompleto que luego se preparó y adquirió fuerza ejecutiva--- el paso siguiente es directamente la intimación al deudor para que oponga excepciones y el embargo. El deudor es embargado desde el comienzo del proceso ejecutivo.

El \textbf{embargo ejecutivo} es el que se traba a consecuencia de la fuerza ejecutiva de un título cartular.

Los tres tipos de embargo ---preventivo, ejecutivo y ejecutorio--- se distinguen según el momento procesal en que se traban y se desarrollan en el capítulo~\ref{cap:embargo}.

%----------------------------------------------------------
\section{La intimación de pago --- Artículo 531}

La intimación de pago equivale, en el juicio ejecutivo, al traslado de demanda del juicio ordinario, pero se realiza mediante un \textbf{mandamiento}. El mandamiento, que dice «intimación de pago y citación de remate», se libra al domicilio del deudor.

Se otorgan cinco días hábiles judiciales para interponer las excepciones que el deudor se crea con derecho a oponer. El mandamiento puede estar acompañado del embargo domiciliario ---en cuyo caso el oficial de justicia va al domicilio, intima de pago y embarga los muebles embargables--- o ser simplemente un mandamiento de intimación de pago, y el embargo se tramita por otro carril: si se trata de inmuebles, se libra oficio directamente al Registro de la Propiedad Inmueble; si se trata de cuentas bancarias, a la entidad bancaria donde esté registrada la cuenta.

\begin{quote}
\textbf{Art. 531 CPCCN.} Regula la intimación de pago y la forma del mandamiento.

Son irrenunciables para el deudor (no pueden ser objeto de acuerdos procesales): 1) la intimación de pago; 2) la citación para oponer excepciones; 3) la sentencia.
\end{quote}

%----------------------------------------------------------
\section{Ampliación de la ejecución --- Artículos 540 y 541}

En el curso del juicio ejecutivo ---especialmente en las ejecuciones de expensas--- pueden vencer nuevos períodos de la misma obligación mientras el juicio tramita. El código permite ampliar la ejecución a esos nuevos períodos.

\begin{quote}
\textbf{Art. 540 CPCCN.} Ampliación por vencimientos anteriores a la sentencia: se reitera la intimación y los nuevos períodos pasan a integrar la suma líquida reclamada.

\textbf{Art. 541 CPCCN.} Vencimientos posteriores a la sentencia: sólo se permite que el deudor traiga el recibo de pago; no se abre nueva discusión.
\end{quote}

%----------------------------------------------------------
\section{Excepciones admisibles --- Artículo 544}

Una vez notificado, el ejecutado tiene cinco días para interponer las excepciones previstas y limitadas por el artículo 544, que son enunciadas de modo taxativo. Si no opone excepciones, el juez dicta la sentencia denominada de \textbf{trance y remate}.

\begin{quote}
\textbf{Art. 544 CPCCN --- Excepciones admisibles en el juicio ejecutivo:}
\begin{enumerate}[label=\arabic*.º),leftmargin=2.2em]
    \item Incompetencia.
    \item Falta de personería en el ejecutante o en el ejecutado, por carecer de capacidad civil para estar en juicio.
    \item Litispendencia.
    \item Falsedad o inhabilidad del título con el que se pide la ejecución.
    \item Prescripción.
    \item Pago total o parcial documentado.
    \item Compensación de crédito líquido que resulte de documento que traiga aparejada la ejecución.
    \item Quita, espera, remisión, novación, transacción, conciliación o cosa juzgada.
\end{enumerate}
\end{quote}

\begin{note}
La falta de legitimación ---omitida por el código--- se ha incorporado jurisprudencial y doctrinariamente a través de la excepción de inhabilidad de título.
\end{note}

%----------------------------------------------------------
\section{Falsedad e inhabilidad del título (art.~544, inc.~4.º)}

\subsection{Falsedad material y falsedad ideológica}

Es fundamental distinguir entre la falsedad material y la falsedad ideológica.

\begin{table}[H]
\centering
\renewcommand{\arraystretch}{1.25}
\begin{tabularx}{\textwidth}{Y Y}
\toprule
\textbf{Falsedad material (admisible)} & \textbf{Falsedad ideológica (inadmisible en el ejecutivo)} \\
\midrule
Adulteración física del documento. & Cuestiona la causa de la obligación. \\
Ejemplo: tachaduras, sobreescrituras, fecha alterada, nombre agregado. & Ejemplo: alegar que no se recibió el dinero del préstamo. \\
Se puede probar con pericial caligráfica. & Debe discutirse en el juicio ordinario posterior. \\
Cosa juzgada material para ese punto. & Produce sólo cosa juzgada formal en el ejecutivo. \\
\bottomrule
\end{tabularx}
\end{table}

La falsedad a la que se refiere la excepción es la \textbf{adulteración física del documento}: que haya sido tachado, enmendado o sobrescrito, que se le hayan agregado condiciones después de firmado, que se haya agregado un nombre o que se haya cambiado o sobrescrito la fecha. Toda adulteración o falsedad material del documento es una excepción admisible en el marco del juicio ejecutivo. Lo que no es admisible es la llamada \textbf{falsedad ideológica}.

\subsection{El caso extremo: ordinarización del proceso}

\begin{example}{hipoteca firmada bajo amenazas}
Una persona firma una hipoteca bajo coacción con amenazas ---con un revólver--- y no recibe nada a cambio. Luego los acreedores ---en este caso, delincuentes--- ejecutan la hipoteca, subastan la casa y se quedan con ese dinero.
\end{example}

En casos extremos como este, la jurisprudencia ha admitido la \textbf{ordinarización del proceso ejecutivo}: se permite que el juicio ordinario corra de modo paralelo al ejecutivo. En ese ordinario, el deudor podría pedir como medida cautelar la suspensión del juicio ejecutivo o, al menos, que ambas sentencias se dicten de modo conjunto.

\subsection{La inhabilidad del título}

La excepción de inhabilidad de título permite cuestionar que el instrumento presentado no trae aparejada la ejecución, lo cual obliga al juez a examinar nuevamente el título en la sentencia. También se utiliza para plantear la falta de legitimación activa o pasiva: que el acreedor no es la persona que surge del título, o que el ejecutado no es quien debe.

\begin{important}
Requisito formal (art.~544, inc.~4.º): tanto la excepción de falsedad como la de inhabilidad del título sólo pueden interponerse siempre que el ejecutado niegue la existencia de la deuda.
\end{important}

%----------------------------------------------------------
\section{Apertura a prueba y sentencia de trance y remate}

El juez analiza las excepciones, rechaza conforme el artículo 547 las que no fueron autorizadas por ley y evalúa las restantes para verificar si cumplen los requisitos. Si son de puro derecho, las resolverá en la sentencia de trance y remate. Si corresponde, puede abrirse prueba por un plazo breve.

La sentencia del juicio ejecutivo se denomina \textbf{sentencia de trance y remate} y debe dictarse en un plazo de diez días, porque se equipara a una resolución interlocutoria.

\begin{quote}
\textbf{Art. 547 CPCCN.} El juez rechaza \emph{in limine} las excepciones no previstas legalmente.
\end{quote}

%----------------------------------------------------------
\section{Apelación e inapelabilidad en el juicio ejecutivo}

En el juicio ejecutivo la regla es contraria a la del juicio ordinario: la \textbf{inapelabilidad} es la norma y sólo son apelables las decisiones que el código autoriza expresamente. Cualquier apelación que se deduzca en el trámite ---salvo las autorizadas--- será denegada. La regla de la inapelabilidad es mucho más estricta para el ejecutado en el marco del juicio ejecutivo.

\begin{quote}
\textbf{Art. 554 CPCCN --- Casos en que procede la apelación de la sentencia de trance y remate:}
\begin{enumerate}[label=\alph*),leftmargin=2.2em]
    \item Cuando se desestiman sin más las excepciones por no cumplirse los requisitos del art.~544.
    \item Cuando las excepciones hubiesen tramitado como de puro derecho.
    \item Cuando se hubiese producido prueba respecto de las opuestas.
    % TODO: contenido incompleto o ambiguo en las notas originales (la redacción del inciso d) del apunte es confusa: "cuyo ordinario posterior")
    \item Cuando la sentencia versare sobre puntos ajenos al ámbito natural del proceso o que cause gravamen irreparable cuyo ordinario posterior.
\end{enumerate}
\end{quote}

El recurso aplicable es la \textbf{apelación en relación}, más acotada que la apelación libre.

%----------------------------------------------------------
\section{Costas y juicio ordinario posterior}

Las costas siguen la regla general: las soporta la parte vencida. La salvedad es que, si el ejecutado logra que se rechace una parte de la ejecución ---por ejemplo, acreditando el pago parcial documentado---, no se imponen las costas de modo proporcional al porcentaje rechazado: la regla es que prospera la ejecución y no hay eximición parcial de costas.

\begin{quote}
\textbf{Art. 553 CPCCN --- Condiciones para el juicio ordinario posterior:} hasta que no esté concluido el juicio ejecutivo no se podría promover el ordinario posterior (salvo los casos de ordinarización paralela ya vistos).
\end{quote}

%----------------------------------------------------------
\section{Ejecuciones especiales --- Artículo 595 y siguientes}

A partir del artículo 595 el código regula las ejecuciones especiales: distintos supuestos de subastas de bienes muebles, inmuebles y otras cuestiones. Se prevén reglas específicas para algunos tipos de ejecución que se basan en el trámite del juicio ejecutivo pero tienen condiciones especiales.

\subsection{La ejecución hipotecaria}

Mediante la Ley 24.441 se creó la denominada hipoteca privada: una ejecución que se apoya en la jurisdicción para realizar algunos trámites pero que se lleva adelante fundamentalmente de modo privado.

Esta regulación recibe objeciones, porque omite el dictado de una sentencia: si el deudor no opone una excepción, el juez no tiene que dictar sentencia. Ello viola directamente la Constitución Nacional, que establece que nadie puede ser privado de su propiedad sino en virtud de una sentencia previa. Si el deudor no pone excepciones, aun cuando debería haber sentencia, se ejecutan los bienes del deudor.

% TODO: contenido incompleto o ambiguo en las notas originales (frase ininteligible en el apunte: "Esto es algo que el adorno vio")
Esta deficiencia no fue subsanada hasta la fecha.

\cornelltitle{juicio ejecutivo}
\cornellblock{
\cqrow{¿Qué es el juicio ejecutivo y cómo se ubica en el código?}{Es un proceso de conocimiento sumario (acotado) regulado a partir del art.~520 CPCCN, dentro del proceso de ejecución. Prioriza la rapidez sobre la exhaustividad: el acreedor accede directamente al embargo y al cobro sin necesidad de un juicio ordinario previo.
}
\cqrow{¿Cuáles son los requisitos para acceder a la vía ejecutiva?}{1) Tener un título ejecutivo (completo o preparado). 2) Que la obligación sea de dar sumas de dinero. 3) Que sea líquida o fácilmente liquidable. 4) Que sea exigible (plazo vencido).
}
\cqrow{¿Qué son los títulos completos e incompletos?}{Completos: traen aparejada la ejecución por sí solos (art.~523: pagaré, cheque, escritura pública, etc.). Incompletos: requieren un trámite previo de preparación de la vía ejecutiva (por ejemplo, el contrato de locación, que requiere la citación del inquilino para obtener la suma líquida).
}
\cqrow{¿Cuáles son los tres tipos de embargo?}{Preventivo: como medida cautelar (requiere verosimilitud y peligro en la demora). Ejecutivo: en el juicio ejecutivo, por la fuerza del título cartular, sin sentencia previa. Ejecutorio: resultado de una sentencia firme en ejecución.
}
\cqrow{¿Qué excepciones admite el juicio ejecutivo?}{Las del art.~544, taxativas: incompetencia, falta de personería, litispendencia, falsedad o inhabilidad del título, prescripción, pago documentado, compensación con título ejecutivo, y quita, espera, remisión, novación, transacción, conciliación o cosa juzgada. La falta de legitimación se incorpora por vía de la inhabilidad de título.
}
\cqrow{¿Qué diferencia hay entre falsedad material e ideológica?}{Material (admisible en el ejecutivo): adulteración física del documento (tachaduras, sobreescrituras, fecha alterada). Ideológica (inadmisible): cuestiona la causa de la obligación (por ejemplo, «firmé pero no recibí el dinero»); debe discutirse en el juicio ordinario posterior.
}
\cqrow{¿Qué es la cosa juzgada formal en el juicio ejecutivo?}{Las cuestiones causales no debatibles en el ejecutivo no quedan cubiertas por cosa juzgada material: pueden discutirse en el ordinario posterior. En cambio, las cuestiones que sí se debatieron (por ejemplo, el pago) generan cosa juzgada material y no se pueden replantear.
}
}
\medskip
\cornellblocksum{
\cqrow{¿Cómo es la apelación en el juicio ejecutivo?}{La regla es la inapelabilidad (inversa a la del juicio ordinario). Sólo son apelables las resoluciones que el código autoriza expresamente (art.~554). El recurso aplicable es la apelación en relación, más acotada que la libre.
}
}{En el juicio ejecutivo el legislador sacrifica la exhaustividad del conocimiento en aras de la celeridad: el acreedor que cuenta con un título ejecutivo ---completo o debidamente preparado--- puede agredir el patrimonio del deudor de manera inmediata, sin aguardar un proceso ordinario de pleno conocimiento. Las excepciones que el deudor puede oponer son taxativas (art.~544 CPCCN) y excluyen la discusión de la causa de la obligación, reservada para el juicio ordinario posterior. La cosa juzgada resultante es, en consecuencia, mixta: formal respecto de las cuestiones causales no debatidas y material respecto de aquellas que sí pudieron ventilarse (como el pago).
}

\chapter{Cumplimiento de la sentencia de remate}\label{cap:remate}

\chaptermark{Cumplimiento de la sentencia de remate}

%----------------------------------------------------------
\section{La etapa coercitiva del proceso}

Hasta la sentencia definitiva, toda la discusión ---la demanda, la contestación, la prueba, la etapa recursiva--- tuvo como meta obtener un reconocimiento de derechos en sede jurisdiccional. Corresponde ahora llevar a la realidad lo dispuesto en el proceso judicial: los derechos reconocidos deben tener su correlato en prestaciones o conductas de las personas en la vida de relación, como pagar la deuda reconocida judicialmente o hacer o dejar de hacer alguna actividad.

Los elementos de la jurisdicción pueden enunciarse en tres: la \textbf{notio} (conocer el conflicto), la \textbf{documentación} y la \textbf{coercio}, es decir, la coerción: el Estado impone su decisión por la fuerza cuando las personas no han cumplido voluntariamente la decisión jurisdiccional.

\subsection{Títulos asimilables a la sentencia}

Tanto el cumplimiento de la sentencia de trance y remate como el procedimiento de ejecución de sentencia ---que se utiliza cuando el juicio ordinario está concluido--- sirven también para los títulos asimilables a la sentencia:

\begin{itemize}
    \item Acuerdos de mediación.
    \item Honorarios regulados judicialmente.
    \item Sentencias extranjeras reconocidas por el mecanismo de exequátur.
    \item Otros que las leyes hayan asimilado.
\end{itemize}

\subsection{Inapelabilidad del ejecutado --- Artículo 560}

El ejecutado que llegó a esta etapa tuvo amplia posibilidad de defensa durante todo el proceso previo. Por eso, en el cumplimiento de la sentencia de remate la regla es la \textbf{inapelabilidad para el ejecutado}.

\begin{quote}
\textbf{Art. 560 CPCCN.} Son inapelables para el ejecutado todas las resoluciones que se dicten durante el trámite de cumplimiento de la sentencia de remate.

Excepciones: las resoluciones que no puedan constituir objeto del juicio ordinario posterior; las que, debiendo ser objeto del ordinario posterior, ya fueron debatidas; las que se relacionen con el carácter de parte del ejecutado; y los casos del art.~550 y 4, inc.~4 (puntos ajenos al ámbito natural del proceso que causan gravamen irreparable).
\end{quote}
% TODO: contenido incompleto o ambiguo en las notas originales (la remisión "art. 550 y 4, inc. 4" figura así en el apunte)

%----------------------------------------------------------
\section{El embargo ejecutorio: presupuesto del remate}

No es posible avanzar hacia el remate de una cosa si no ha sido embargada previamente. Se trata del \textbf{embargo ejecutorio}, que es el resultado de la traba que se realiza una vez que existe una condena firme en sede jurisdiccional.

Si lo que se va a percibir son sumas de dinero ya embargadas, no hay que hacer ningún remate: con la firmeza de la sentencia de venta se practica una liquidación y se ordena el giro para que el acreedor se haga de ese dinero. El remate cobra relevancia cuando hay un bien ---mueble o inmueble--- que debe subastarse para transformarlo en dinero.

%----------------------------------------------------------
\section{Bienes subastables}

Se puede subastar todo aquello que esté en el comercio, todo aquello que pueda ser vendido.

También es posible subastar derechos hereditarios. Cuanto mayor sea la determinación del acervo hereditario, con mayor precisión se podrá establecer el valor de ese derecho y, de tal manera, fijar la base de la venta.

%----------------------------------------------------------
\section{El martillero: designación y funciones}

El encargado de llevar adelante la subasta a las órdenes del juez es el \textbf{martillero} o corredor inmobiliario, que se anota en las listas de las cámaras para ser designado. Actualmente, en el fuero civil de la Nación, la designación se realiza directamente por sorteo por computadora.

\subsection{Obligaciones del martillero}

\begin{enumerate}
    \item Cumplir las funciones personalmente: no puede delegarlas.
    \item Responder sólo a las órdenes del juez, no de las partes.
    \item Aceptar el cargo obligatoriamente.
    \item Informar el valor fiscal y de mercado del bien a subastar.
    \item Exhibir el inmueble a los eventuales compradores.
    \item Publicar los edictos.
    \item Depositar en el proceso las sumas percibidas en el acto de subasta.
    \item Puede pedir un anticipo de fondos al ejecutante para cubrir los gastos del trámite.
\end{enumerate}

%----------------------------------------------------------
\section{La liga: distorsión del sistema de subastas}

\begin{example}{el fenómeno de la «liga»}
Una serie de personas se dedica a distorsionar lo que se obtiene en las subastas judiciales.

\textbf{Mecanismo:} impiden que concurran eventuales compradores interesados en un bien particular, de modo de intentar adquirirlo por el menor valor posible ---si es posible, por la base---. Luego realizan la «verdadera subasta» fuera del tribunal: el verdadero comprador paga a la liga la diferencia entre lo pagado en sede judicial y el valor real.

\textbf{Impacto:} esa diferencia ---que puede oscilar entre un 20 y un 30\,\% del valor de un inmueble--- es la que el deudor no obtiene. El acreedor cobra menos y el deudor ---si había un remanente--- también pierde.

\textbf{Ejemplo numérico:} si el inmueble vale \$100.000 y se subasta a \$60.000 cuando podría haberse subastado a \$90.000, esos \$30.000 quedan en poder de la liga.
\end{example}

%----------------------------------------------------------
\section{Subastas electrónicas}

Se señala como una necesidad impostergable la modernización del sistema: resulta fuera de contexto que en la práctica se siga sacando a subasta un inmueble en una oficina, con todos los interesados presentes durante unas dos horas para hacer la puja en forma presencial. Los anuncios podrían hacerse por internet en lugar de por edictos, que son costosos, y las subastas podrían realizarse de modo digital o electrónico.

La Provincia de Buenos Aires tiene una ley de subasta electrónica que se está implementando gradualmente. La subasta no dura dos horas sino quince días: distintos interesados toman conocimiento del inmueble por internet, se registran y hacen ofertas, que se actualizan a medida que aparece un mejor postor.

Las ventajas son múltiples: elimina a la liga, obtiene valores mucho más altos y reales y beneficia tanto al acreedor ---que cobra más--- como al deudor ---que puede obtener un remanente luego de pagada la deuda---.

%----------------------------------------------------------
\section{Trámites previos a la subasta de inmuebles}

\subsection{Informes registrales y de deudas}

Cuando el juez ordena la subasta y designa al martillero, ordena también que las partes obtengan informes de todos los entes ---fundamentalmente las municipalidades y los organismos de recaudación provincial, como ARBA en Buenos Aires--- sobre las deudas que tienen los inmuebles. También se solicitan informes de deudas por servicios: agua (AySA), gas, luz, etc.

\begin{itemize}
    \item \textbf{Informes registrales:} se requieren al Registro de la Propiedad Inmueble sobre embargos, inhibiciones y gravámenes (hipotecas) que pesen sobre el inmueble y sobre el ejecutado.
    \item \textbf{Validez de los informes:} 60 días. No es necesario actualizarlos antes de la subasta, salvo que haya transcurrido mucho tiempo por alguna suspensión del proceso.
    \item \textbf{Inmuebles en propiedad horizontal:} debe requerirse también la deuda por expensas.
\end{itemize}

\subsection{Citación de los acreedores con gravámenes}

Este trámite es muy importante, porque no es lo mismo subastar un inmueble que no tiene gravamen que subastar uno que tiene una hipoteca previa. Cuando el inmueble tiene hipotecas o embargos, el acreedor, antes de cobrar, deberá respetar las preferencias y privilegios de esos gravámenes.

Una vez que el inmueble se subasta y el embargo pasa al producido de la subasta ---por subrogación real---, el juez debe repartir ese dinero de acuerdo con los privilegios y preferencias de los gravámenes y embargos. Por eso se notifica a todos los acreedores antes de la subasta, para que puedan presentarse, defender el precio de venta y controlar la legalidad del proceso.

%----------------------------------------------------------
\section{Los edictos y la publicidad --- Artículo 566}

\begin{quote}
\textbf{Art. 566 CPCCN.} El remate se anunciará por dos días en el Boletín Oficial y en otro diario de amplia circulación del lugar donde se ubica el inmueble.
\end{quote}

El edicto debe contener obligatoriamente:

\begin{itemize}
    \item Expediente, carátula y juzgado donde tramita.
    \item Fecha, hora y lugar de la subasta.
    \item Días y horas en que el martillero exhibirá el inmueble a los interesados.
    \item Deudas del inmueble (servicios, tasas, contribuciones municipales, expensas).
    \item Gravámenes que pesan sobre el bien.
    \item Obligación del adquirente de pagar en el acto la seña y la comisión del martillero.
    \item Modalidades especiales fijadas por el juez para la venta.
\end{itemize}

Se observa que la referencia a los diarios como medio de comunicación más relevante resulta anacrónica, dado que poca gente sigue comprando diarios, si se compara con el tiempo en que la norma fue redactada. Sin embargo, la norma continúa vigente y debe cumplirse.

El martillero puede también solicitar al tribunal autorización para realizar publicidad adicional ---por ejemplo, por internet--- para atraer más compradores.

%----------------------------------------------------------
\section{La base de la subasta --- Artículo 578}

El martillero debe informar al juez el valor fiscal del inmueble y su valor de mercado. Muchas veces los valores fiscales están muy desconectados del verdadero valor de mercado de una propiedad, es decir, del valor que podría llegar a obtenerse en una venta.

\begin{formula}{base de la subasta (art.~578 CPCCN)}
Si no hay acuerdo de partes, la base de venta se fijará en \textbf{dos tercios de la valuación fiscal actualizada}.
\end{formula}

\begin{note}
En general, el juez se aparta de ese valor fiscal para impedir que los bienes sean malvendidos, fijando una base menor al valor de mercado para atraer distintos compradores. Las partes ---y los acreedores con gravámenes--- pueden impugnar la tasación informada por el martillero.
\end{note}

%----------------------------------------------------------
\section{El acto de remate}

Hay distintos resultados posibles: que haya un comprador, que el remate fracase por falta de postores o que el propio acreedor ejecutante quiera adquirir el bien.

\subsection{Cuando hay comprador}

Si hay comprador, el martillero le hace firmar en el acto de la subasta el \textbf{boleto de venta judicial}, donde consta el precio de compra, lo que paga de seña, lo que paga de comisión del martillero y algún otro impuesto que corresponda a la venta. El martillero retiene su comisión y deposita el dinero obtenido en el proceso, informando el nombre e identidad del comprador.

El comprador debe constituir domicilio electrónico (art.~41 CPCCN) para ser notificado de las decisiones judiciales.

Secuencia posterior al acto de remate:

\begin{enumerate}
    \item El martillero informa el resultado al tribunal.
    \item Se da conocimiento del resultado a acreedor y deudor, que tienen cinco días para formular objeciones.
    \item Sin objeciones (o desestimadas), el juez aprueba el remate.
    \item El comprador tiene cinco días para depositar el saldo de precio: el 70\,\% restante, dado que el 30\,\% se paga como seña en el acto.
\end{enumerate}

\subsection{Crítica a la exigencia del dinero en efectivo}

Se señala que exigir el 30\,\% en efectivo el día de la subasta está fuera de contexto: obliga a las personas a concurrir con dinero que quizá no usarían, con todos los riesgos que ello trae. Para los fines de la efectividad de la subasta, resulta bastante complejo.

%----------------------------------------------------------
\section{Aprobación y nulidades}

El plazo para plantear nulidades es de cinco días contados desde que se conoce el resultado de la subasta. Para quien estuvo presente, el cómputo se inicia desde el propio acto, independientemente de cuándo el martillero informe al tribunal.

Las nulidades que pueden plantearse en esta etapa son las \textbf{nulidades propias del acto}: las ocurridas ese día, no las nulidades procesales anteriores. Estas últimas ---las que se fueron produciendo desde que se ordenó la subasta, como la base o el edicto--- debían cuestionarse dentro de los cinco días desde que se conocieron en el expediente. En general, la única nulidad que podrá plantearse es algún incumplimiento en el acto propio de la subasta.

\begin{important}
\textbf{Principio de trascendencia:} quien articula la nulidad debe acreditar el perjuicio. Sin perjuicio concreto, no procede la nulidad: la falta menor no hace caer el acto. La trascendencia de la nulidad debe tener entidad suficiente para hacer caer la venta.
\end{important}

%----------------------------------------------------------
\section{El sobreseimiento del juicio --- Artículo 583}

Aun hecha la subasta, el deudor tiene una posibilidad más de recuperar el bien: el ejecutado puede, antes de que el comprador en subasta deposite el saldo de precio, pedir el \textbf{sobreseimiento}. Una vez que el comprador depositó el saldo de precio, se considera perfeccionada la venta.

Para que proceda el sobreseimiento, el deudor debe depositar:

\begin{enumerate}
    \item El capital adeudado.
    \item Los intereses presupuestados.
    \item Las costas del proceso.
    \item Todos los gastos.
    \item Un importe a favor del comprador integrado por la comisión que pagó al martillero, los sellados del boleto y el equivalente a la seña depositada.
\end{enumerate}

%----------------------------------------------------------
\section{Título del adquirente e inscripción registral}

Así como en una compraventa de inmueble fuera de los tribunales hay una escritura, dentro del proceso se extiende un \textbf{testimonio} que contiene toda la referencia y los antecedentes de la subasta judicial y las distintas resoluciones del juez, para que el adquirente en subasta lo inscriba en el Registro de la Propiedad Inmueble. Ese testimonio será el nuevo título del bien. Puede obtenerse sin necesidad de escribano: es un trámite más del proceso, que extiende el secretario del tribunal.

\subsection{Levantamiento de gravámenes}

El adquirente pide el levantamiento de todas las medidas dispuestas contra el inmueble ---hipotecas, embargos--- al solo efecto de escritura, para que el título ingrese y el inmueble quede liberado. Los acreedores que tenían esos gravámenes deben presentarse en el proceso a hacer valer sus preferencias y privilegios en la distribución del producido.

%----------------------------------------------------------
\section{Distribución del producido: la liquidación}

Una vez depositado el saldo de precio, se practica la liquidación: capital, intereses, costas y gastos. Se pone en conocimiento del deudor ---que tiene un plazo para impugnarla--- y, una vez aprobada, se entrega al ejecutante, cancelando el crédito.

El deudor puede tener un \textbf{remanente}: si el producido de la subasta supera la deuda total (capital, intereses y costas), el excedente le corresponde al deudor.

%----------------------------------------------------------
\section{Subasta fracasada y acreedor comprador}

\subsection{Fracaso por falta de postores}

Si no hay postores, se realizará una nueva subasta, en general intentando bajar la base para atraer nuevos interesados (arts.~592 y 593 CPCCN).

\subsection{El acreedor como comprador}

El acreedor puede concurrir a comprar él mismo el inmueble que se subasta y \textbf{compensar su crédito con el precio}. En general, los acreedores ejecutantes piden al juez que, en caso de resultar adquirentes, los exima de depositar la seña. Es un gran beneficio para ese comprador: no tiene que concurrir con el dinero en la mano, expuesto a la inseguridad.

Esta posibilidad también podría ser solicitada por algún acreedor embargante de otro juicio que deba concurrir a la subasta del inmueble en el proceso en que se va a realizar, para ser eximido del depósito de la seña.

%----------------------------------------------------------
\section{Los plenarios: diferencia entre fueros}

\begin{important}
\textbf{Plenario de la Cámara Comercial:} el adquirente debe pagar las deudas que están expresadas en el edicto.

\textbf{Plenario de la Cámara Civil:} el adquirente no tiene que afrontar todas esas deudas, en tanto alcance el dinero que se obtiene en la subasta.

\textbf{Consecuencia práctica:} según se subaste en un juzgado civil o comercial, la ecuación económica para el comprador cambia radicalmente. Un abogado debe verificar los plenarios vigentes antes de aconsejar a su cliente.
\end{important}

%----------------------------------------------------------
\section{Desocupación del inmueble --- Artículo 589}

\begin{quote}
\textbf{Art. 589 CPCCN.} Como regla, no hay desalojo de los ocupantes del inmueble hasta que no se hubiese pagado el saldo de precio y hecho la tradición de la cosa al adquirente. Se exceptúan los casos en que se haya pactado otra cosa, como en las hipotecas, donde el hipotecado generalmente se obliga a desocupar dos días después del auto de subasta.
\end{quote}

\cornelltitle{cumplimiento de la sentencia de remate}
\cornellblocksum{
\cqrow{¿Qué es el cumplimiento de la sentencia de remate?}{Es la etapa coercitiva del proceso: una vez firme la sentencia de trance y remate (o de venta en el ordinario), el Estado impone su decisión por la fuerza ejecutando los bienes del deudor. Presupuesto indispensable: el embargo ejecutorio.
}
\cqrow{¿Cuáles son las funciones del martillero?}{Auxiliar del tribunal: evalúa el bien, informa valor fiscal y de mercado, lo exhibe a los compradores, publica edictos, conduce el acto de remate, suscribe el boleto de venta judicial, retiene su comisión y deposita el producido. No puede recibir órdenes de las partes.
}
\cqrow{¿Qué es «la liga» y cómo opera?}{Un grupo organizado que impide la concurrencia de compradores interesados al acto de remate para adquirir el bien por el mínimo (la base) y luego realiza una subasta privada fuera del tribunal. La diferencia de precio (20--30\,\% del valor) queda en sus manos, en perjuicio del acreedor y del deudor.
}
\cqrow{¿Cuál es la base de la subasta según el art.~578?}{Si no hay acuerdo de partes: dos tercios de la valuación fiscal actualizada. En la práctica, el juez suele apartarse de ese valor fiscal (muy desconectado del mercado) y fija una base menor al valor de mercado para atraer compradores sin que el bien sea malvendido.
}
\cqrow{¿Cuáles son los pasos tras el acto de remate?}{1) El martillero informa al tribunal. 2) Cinco días para objeciones. 3) Aprobación del remate. 4) Cinco días para depositar el saldo de precio. 5) Posibilidad de sobreseimiento por el deudor antes del depósito. 6) Liquidación, aprobación y pago al acreedor. 7) Testimonio e inscripción del título del adquirente.
}
\cqrow{¿Qué es el sobreseimiento del juicio (art.~583)?}{El ejecutado puede recuperar el bien después del remate pero antes de que el comprador deposite el saldo de precio. Para ello debe pagar: capital, intereses, costas, gastos, comisión del martillero, sellados del boleto y el equivalente a la seña.
}
\cqrow{¿Por qué importan los plenarios sobre deudas en las subastas?}{La Cámara Comercial establece que el adquirente paga sólo las deudas expresadas en el edicto. La Cámara Civil establece que todas las deudas se abonan con el producido de la subasta (el adquirente no las asume si alcanza el dinero). La diferencia cambia radicalmente la ecuación económica para el comprador.
}
}{Una vez firme la sentencia de trance y remate, el Estado despliega su poder coercitivo pleno: se traba el embargo ejecutorio, se designa un martillero, se practican los informes registrales, se publican edictos, se notifica a los acreedores concurrentes y se realiza el acto de subasta. El producido se distribuye según los privilegios y preferencias de los gravámenes. El adquirente en subasta obtiene su título a través de un testimonio extendido por el secretario del tribunal, sin necesidad de escritura pública. El sistema presenta dos grandes deficiencias estructurales: la persistencia del mecanismo presencial, que alimenta el fenómeno de la «liga», y la divergencia entre los plenarios de la Cámara Civil y la Cámara Comercial sobre las deudas asumidas por el adquirente, que el operador jurídico debe consultar en cada caso para brindar un asesoramiento económicamente preciso.
}

\partintro{Los sistemas cautelares son herramientas procesales que permiten al juez \textbf{conservar} el estado de las cosas mientras dura el proceso, evitando que la sentencia llegue tarde o resulte inútil. Esta parte comienza con los fundamentos, la clasificación, los presupuestos y los caracteres de las medidas cautelares, continúa con el sistema asegurativo genérico y con el ámbito excepcional de la tutela anticipada, y concluye con el estudio de las medidas cautelares en particular: el embargo, el secuestro, la intervención judicial, la anotación de litis, la prohibición de contratar y las medidas de protección de personas y de familia. Las medidas en particular se estudian siempre en conexión con la parte general (verosimilitud del derecho, peligro en la demora y contracautela).

El proceso civil cuenta con diversas herramientas para garantizar la eficacia de la sentencia, del mismo modo en que un médico dispone de distintos instrumentos para tratar diferentes dolencias: desde el embargo, la más usada y análoga a un seguro sobre un bien, hasta las medidas de familia, orientadas a proteger derechos fundamentales de los niños y los vínculos familiares.}
\part{Sistemas cautelares}

\chapter[Fundamentos de los sistemas cautelares]{Fundamentos de los sistemas cautelares}\label{cap:fundamentos-cautelares}

% TODO: verificar consistencia entre fuentes originales (uno de los archivos fuente numera de manera inconsistente la unidad de sistemas cautelares).

\section{Importancia de los sistemas cautelares}

Los sistemas cautelares constituyen un tema de gran trascendencia para el ejercicio profesional, porque son herramientas que posibilitan a los abogados, a quienes trabajan dentro del Poder Judicial y a los funcionarios cumplir con la \textbf{tutela judicial efectiva}.

Sin los sistemas cautelares sería necesario esperar a la sentencia en todos los procesos judiciales, lo que haría que muchos derechos quedaran vulnerados. En ausencia de cautelares, muchas sentencias solamente servirían para «colgarse en un cuadro».

Se trata de un tema de gran importancia para el justiciable y para todo operador jurídico que deba emprender la gestión de un conflicto, ya que estas herramientas permiten, frente a un peligro inminente o una situación afligente, dar una respuesta desde el sistema de justicia en un tiempo distinto del que insumiría un proceso hasta llegar a su fin, es decir, hasta la sentencia.

Los sistemas cautelares son herramientas procesales que permiten al juez \textbf{conservar} el estado de las cosas mientras dura el proceso, evitando que la sentencia llegue tarde o resulte inútil. Puede pensarse el proceso judicial como una carrera: los sistemas cautelares funcionan como cercos que impiden que alguien modifique el camino antes de que se declare al ganador.

\section{La actuación de la ley en el proceso}

Al estudiar la jurisdicción y sus distintas etapas (momentos jurisdiccionales) se distinguen tres modalidades de actuación de la ley en el proceso.

\begin{table}[H]
\centering
\small
\begin{tabularx}{\textwidth}{>{\raggedright\arraybackslash}p{0.2\textwidth} Y Y}
\toprule
\textbf{Modo de actuación} & \textbf{Cuándo ocurre} & \textbf{En qué etapa} \\
\midrule
Para \textbf{conocer} & Cuando la jurisdicción conoce los hechos & Etapa introductoria (escritos postulatorios), probatoria (producción de prueba) y alegatos \\
Para \textbf{ejecutar} & Cuando el juez ejecuta la sentencia & Luego de la sentencia firme y consentida (ejecutoriada) \\
Para \textbf{conservar / asegurar} & Antes o durante el proceso, ante urgencia & En cualquier momento en que se requiera respuesta inmediata \\
\bottomrule
\end{tabularx}
\end{table}

Esta distinción corresponde a las tres finalidades que, según Calamandrei, tiene la actuación de la ley en el proceso: que el juez conozca, que ejecute y que conserve.

La actuación de la ley con \textbf{carácter conservatorio} es la que interesa para el estudio de los sistemas cautelares. Esta actuación abarca tanto la función conservativa como la innovativa:

\begin{itemize}
    \item \textbf{Conservar para mantener un determinado statu quo:} cuando el riesgo proviene de que se altere la situación existente.
    \item \textbf{Conservar alterando:} cuando el riesgo proviene de que se mantenga una determinada situación o statu quo y esa conservación es la que genera el daño.
\end{itemize}

\section{El poder cautelar}

El \textit{iter} procesal no puede soslayarse: un proceso judicial insume necesariamente tiempo. Si se transcurre todo el \textit{iter} procesal hasta llegar a la sentencia, es probable que situaciones vigentes, o que requieran una respuesta rápida, inmediata y efectiva, no tengan solución. Los sistemas cautelares vienen a dar respuesta a esas situaciones.

\begin{definition}{Poder cautelar}
Respuesta efectiva e inmediata que se da en el marco de un proceso judicial prudente, con el fin de respetar la tutela judicial efectiva con relación a los derechos discutidos en ese proceso.
\end{definition}

Sus objetivos son:

\begin{itemize}
    \item Contrarrestar los efectos negativos del transcurso del tiempo en el proceso judicial.
    \item Respetar la tutela judicial efectiva de los derechos discutidos.
    \item Mantener la igualdad de las partes en el proceso.
\end{itemize}

\begin{important}
El legislador crea los sistemas cautelares para contrarrestar el efecto que el tiempo genera en el proceso ---lo que Calamandrei denominaba el ``ordinario iter procesal''--- y para mantener la \textbf{igualdad de las partes ante la jurisdicción}. Se busca dar respuesta a situaciones urgentes, garantizar una tutela judicial efectiva y evitar que la sentencia resulte estéril por tardía.
\end{important}

\begin{example}{Camacho Acosta (CSJN)}
Fallo de la Corte Suprema de Justicia de la Nación que se trabajará más adelante al estudiar la prohibición de innovar. Un trabajador pierde el antebrazo y necesita una prótesis. La prótesis exige una respuesta inmediata de la jurisdicción: no se puede esperar a la sentencia para decidir si la prótesis le correspondía o no. La jurisdicción debe anticipar esa decisión a fin de resguardar la salud y el derecho de esa persona a contar con la prótesis. Es el ejemplo paradigmático de la necesidad de las medidas cautelares.
\end{example}

Otros supuestos frecuentes en los que se requieren sistemas cautelares son:

\begin{itemize}
    \item La protección de una persona.
    \item La conservación de una prueba para el proceso.
    \item Mantener el giro comercial de una empresa.
    \item La aplicación de una norma inconstitucional a un caso concreto (solicitud de medida cautelar para suspender sus efectos mientras tramita el proceso).
\end{itemize}

\section{Cautela y condena}

\textbf{Cautela no es condena}: son dos términos distintos. El término \textit{cautelar} deriva de «cauto», quien observa con cuidado, y también de «caución», vinculado con «precaver» y con lo «precavido». Todos estos términos remiten a la prevención y a la conservación, es decir, a poder cautelar una determinada situación de hecho o de derecho que se dé en un caso concreto.

\begin{important}
La cautela se mueve en un ámbito \textbf{provisorio}, que nada tiene que ver con la condena. La condena se da únicamente en la sentencia, luego de pasar por todas las etapas del proceso.
\end{important}

\begin{table}[H]
\centering
\begin{tabularx}{\textwidth}{>{\raggedright\arraybackslash}p{0.47\textwidth} Y}
\toprule
\textbf{Cautela} & \textbf{Condena} \\
\midrule
Provisoria & Definitiva \\
Anterior a la sentencia & Es la sentencia definitiva \\
Basada en la verosimilitud del derecho & Basada en certeza jurídica plena \\
Puede modificarse o levantarse & No puede modificarse (cosa juzgada) \\
Tramita \textit{inaudita parte} & Con plena bilateralidad \\
Requiere contracautela & No requiere contracautela \\
Sujeta a caducidad & No caduca \\
Instrumental, accesoria & Principal, autónoma \\
\bottomrule
\end{tabularx}
\end{table}

La sentencia de mérito implica el máximo momento de la jurisdicción: fijar los hechos del proceso, elevarlos a las normas abstractas y crear la norma individual para el caso concreto, luego del conocimiento del juez, de los hechos, de la producción de la prueba y de los alegatos. Ese proceso previo a la sentencia es ajeno a la cautela, que tiene otro nivel de conocimiento, inferior, fundado en la verosimilitud del derecho.

\cornelltitle{fundamentos de los sistemas cautelares}
\cornellblocksum{
\cqrow{¿Por qué son importantes los sistemas cautelares?}{Porque permiten cumplir con la tutela judicial efectiva: sin ellos habría que esperar a la sentencia en todos los procesos y muchos derechos quedarían vulnerados; muchas sentencias solo servirían para «colgarse en un cuadro». Permiten dar una respuesta al justiciable, frente a un peligro inminente o una situación afligente, en un tiempo distinto del que insume el proceso hasta la sentencia.
}
\cqrow{¿Cuáles son los tres modos en que la ley actúa en el proceso?}{1) Para \textbf{conocer} (etapas introductoria, probatoria y de alegatos). 2) Para \textbf{ejecutar} (luego de la sentencia firme y consentida). 3) Para \textbf{conservar o asegurar} (ante urgencia, en cualquier momento en que se requiera respuesta inmediata); este último es el que interesa a los sistemas cautelares.
}
\cqrow{¿Qué significa conservar manteniendo el statu quo y conservar alterando?}{Conservar para mantener un statu quo se da cuando el riesgo proviene de que se altere la situación existente. Conservar alterando se da cuando el riesgo proviene de que se mantenga la situación y esa conservación es la que genera el daño.
}
\cqrow{¿Qué es el poder cautelar y cuál es su finalidad?}{Es la respuesta efectiva e inmediata que se da en el marco de un proceso judicial prudente para respetar la tutela judicial efectiva de los derechos discutidos. Busca contrarrestar los efectos negativos del transcurso del tiempo, respetar esa tutela y mantener la igualdad de las partes.
}
\cqrow{¿Qué ilustra el caso Camacho Acosta (CSJN)?}{La necesidad de las medidas cautelares: un trabajador que perdió el antebrazo necesita una prótesis y la jurisdicción debe anticipar la decisión, sin esperar a la sentencia, para resguardar su salud y su derecho.
}
\cqrow{¿Qué significa «cautelar» y en qué se diferencia de la condena?}{Deriva de «cauto» y de «caución» (precaver), lo que remite a prevención y conservación. La cautela es provisoria, anterior a la sentencia, basada en la verosimilitud, modificable, \textit{inaudita parte}, con contracautela, sujeta a caducidad e instrumental. La condena es definitiva, es la sentencia, se basa en certeza plena, no puede modificarse (cosa juzgada), tiene plena bilateralidad, no requiere contracautela, no caduca y es principal y autónoma.
}
}{la ley actúa en el proceso para conocer, ejecutar y conservar o asegurar. El poder cautelar es la respuesta efectiva e inmediata que contrarresta el efecto del tiempo procesal, respeta la tutela judicial efectiva y mantiene la igualdad de las partes. La cautela, provisoria y fundada en la verosimilitud, se distingue de la condena, definitiva y fundada en la certeza plena que solo se alcanza con la sentencia.
}

\cornelltitle{marco general de los sistemas cautelares}
\cornellblocksum{
\cqrow{¿Cuáles son las tres finalidades de la ley en el proceso según Calamandrei?}{La ley actúa para que el juez conozca (etapa introductoria y probatoria), para que ejecute (etapa de ejecución, para que sus decisiones sean cumplidas) y para que conserve (asegurar, proteger y preservar, que es el campo de los sistemas cautelares).
}
\cqrow{¿Qué significa que el juez ``conserve''?}{Conservar puede ser mantener o alterar una situación de hecho o de derecho, según de dónde provenga el perjuicio: se ordena mantener cuando el perjuicio proviene de alterar la situación, y a la inversa.
}
\cqrow{¿Por qué existen los sistemas cautelares?}{Para contrarrestar el efecto del tiempo en el proceso (el ``ordinario iter procesal'' de Calamandrei), mantener la igualdad de las partes ante la jurisdicción, responder a situaciones urgentes con tutela judicial efectiva y evitar una sentencia estéril por tardía.
}
\cqrow{¿Cómo se clasifican los sistemas cautelares?}{En ámbito normal (asegurativo probatorio, asegurativo garantizador, precautorio tradicional y asegurativo genérico), donde la pretensión cautelar no coincide con la sustancial; y ámbito excepcional (tutela anticipada), donde se superpone total o parcialmente con ella.
}
\cqrow{¿Qué presupuestos deben acreditarse para peticionar una medida cautelar?}{Verosimilitud en el derecho (probabilidad, no certeza), peligro en la demora (perjuicio inminente o irreparable), contracautela (garantía por los daños que pudiera causar la traba incorrecta) y los presupuestos procesales del Art.~195 CPCC.
}
\cqrow{¿Qué diferencia hay entre lo cautelar y la sentencia en cuanto a su carácter?}{Lo cautelar es provisorio y puede modificarse, reducirse, ampliarse o dejarse sin efecto según varíen las circunstancias; la sentencia es definitiva.
}
}{Los sistemas cautelares permiten al juez conservar la situación de hecho o de derecho durante el proceso, contrarrestando el efecto del tiempo y preservando la igualdad de las partes. Se dividen en un ámbito normal y un ámbito excepcional (tutela anticipada). Para obtener cualquier medida se acreditan verosimilitud en el derecho, peligro en la demora, contracautela y los presupuestos procesales del Art.~195 CPCC, y lo resuelto es siempre provisorio, a diferencia de la sentencia.
}

\chapter[Clasificación de los sistemas cautelares]{Clasificación de los sistemas cautelares}\label{cap:clasificacion-cautelar}

\section{El tratamiento como sistema}

Los institutos cautelares se estudian como un \textbf{sistema} porque, al recorrer el Código Procesal Civil y Comercial, se encuentra que la ley actúa con carácter conservativo o asegurativo en distintos lugares y con distintas modalidades. Agrupar esas manifestaciones en subsistemas permite comprenderlas en su conjunto y distinguirlas, a los fines de entender el sistema del código procesal en todos aquellos aspectos en los que la ley actúa de este modo.

\section{Ámbito normal y ámbito excepcional}

Los sistemas cautelares se clasifican en dos grandes ámbitos.

\begin{table}[H]
\centering
\begin{tabularx}{\textwidth}{>{\raggedright\arraybackslash}p{0.25\textwidth} Y}
\toprule
\textbf{Ámbito} & \textbf{Tipo de sistema} \\
\midrule
Normal & Sistemas asegurativos probatorios \\
Normal & Sistemas precautorios tradicionales \\
Normal & Sistemas asegurativos garantizadores \\
Normal & Sistemas asegurativos genéricos \\
Excepcional & Tutela anticipada \\
Excepcional & Sistemas protectores \\
\bottomrule
\end{tabularx}
\end{table}

Según otra formulación presente en las fuentes, el criterio que distingue ambos ámbitos es la relación entre la pretensión cautelar y la pretensión sustancial:

\noindent\begin{tabularx}{\textwidth}{@{}>{\raggedright\arraybackslash}p{2.6cm} Y Y@{}}
\toprule
\textbf{Ámbito} & \textbf{Sistemas} & \textbf{Característica principal} \\
\midrule
Normal (típico) &
Asegurativo probatorio; asegurativo garantizador; precautorio tradicional; asegurativo genérico. &
Están regulados en la ley; la pretensión cautelar \textbf{no} coincide con la pretensión sustancial. \\
\addlinespace
Excepcional &
Tutela anticipada. &
La pretensión cautelar se superpone total o parcialmente con la pretensión sustancial. \\
\bottomrule
\end{tabularx}

% TODO: verificar consistencia entre fuentes originales (una fuente incluye los sistemas protectores en el ámbito excepcional; otra menciona únicamente la tutela anticipada).

\section{Ámbito normal}

\subsection{Sistemas asegurativos probatorios}

Estos sistemas propenden a la \textbf{mayor eficacia del proceso}. Su finalidad es resguardar fuentes de prueba que pueden perderse con el transcurso del tiempo: si no se permite a la ley actuar en el proceso de este modo, hay prueba que se perdería.

\begin{normativa}{Artículo 326 --- CPCCN: prueba anticipada}
Regula la \textbf{prueba anticipada}. Es allí donde se encuentra regulado el sistema asegurativo probatorio. La ley actúa en el proceso con carácter asegurativo o conservatorio, permitiendo adelantar el momento en el cual se produce una prueba que, de no producirse en ese momento, se perdería al aguardar la etapa probatoria ordinaria del proceso.
\end{normativa}

Si en la etapa probatoria ordinaria se pierde una prueba porque no se permitió su adelantamiento, la parte que la tenía pierde la posibilidad de acreditar el fundamento de su pretensión, lo que afecta la igualdad de las partes en el proceso.

\subsection{Sistemas precautorios tradicionales}

Agrupan las medidas cautelares clásicas: el embargo, la anotación de litis, el secuestro, entre otras. Apuntan fundamentalmente a resguardar bienes.

\begin{important}
Su finalidad es resguardar el efectivo cumplimiento de la sentencia. Si no se traba un embargo para hacer efectiva la sentencia, esa sentencia solamente servirá para «colgarse en un cuadro», porque no podrá ejecutarse.
\end{important}

\begin{table}[H]
\centering
\begin{tabularx}{\textwidth}{>{\raggedright\arraybackslash}p{0.3\textwidth} Y}
\toprule
\textbf{Sistema} & \textbf{Finalidad} \\
\midrule
Asegurativo probatorio & Propender a la mayor eficacia del proceso (resguardar prueba) \\
Precautorio tradicional & Resguardar el efectivo cumplimiento de la sentencia (resguardar bienes) \\
\bottomrule
\end{tabularx}
\end{table}

\subsection{Sistemas asegurativos garantizadores}

Estos sistemas garantizan a la parte contraria, es decir, a quien puede ser afectado por la medida cautelar, frente a la posibilidad de que la medida sea revocada o no reconocida en la sentencia. El instrumento central es la \textbf{contracautela}.

Si un juez dicta una medida y luego se acredita que no fue correctamente dictada, el sistema asegurativo garantizador, es decir, la contracautela, resguarda los daños y perjuicios que la medida pueda haber ocasionado al afectado.

\begin{normativa}{Artículo 509 --- CPCCN: fianza para ejecutar}
Exige una fianza para poder seguir avanzando con la ejecución, a fin de resguardar los derechos de la parte contraria en caso de que la sentencia sea revocada.
\end{normativa}

\begin{normativa}{Artículo 555 --- CPCCN: juicio ejecutivo}
También exige fianza para ejecutar, por las mismas razones que el artículo 509: resguardar a quien soportó la resolución contraria en caso de reversión del proceso.
\end{normativa}

\subsection{Sistemas asegurativos genéricos}

Son el cuarto tipo dentro del ámbito normal.

\begin{normativa}{Artículo 228 --- CPCCN: inhibición general de bienes}
Es una medida residual que se otorga cuando no existen bienes para embargar. Pertenece al sistema asegurativo genérico dentro del ámbito normal de los sistemas cautelares.
\end{normativa}

\begin{normativa}{Artículo 230 --- CPCCN: prohibición de innovar}
Es la medida elegida como eje explicativo del sistema cautelar, porque permite el acceso a la tutela anticipada (ámbito excepcional) y porque presenta dos facetas: la medida de no innovar y la medida innovativa.
\end{normativa}

\begin{table}[H]
\centering
\small
\begin{tabularx}{\textwidth}{>{\raggedright\arraybackslash}p{0.2\textwidth} Y Y}
\toprule
\textbf{Faceta} & \textbf{Cuándo se aplica} & \textbf{Qué se pide} \\
\midrule
Medida de \textbf{no innovar} & Cuando el perjuicio proviene de la modificación del statu quo & Que no se innove; que se conserve la situación existente \\
Medida \textbf{innovativa} & Cuando el perjuicio proviene de mantener el statu quo & Que se innove; que se modifique la situación para evitar el daño irreparable \\
\bottomrule
\end{tabularx}
\end{table}

En ambos casos se conserva: la medida innovativa se da cuando el perjuicio va a provenir de mantener el statu quo; entonces se pide que se innove, porque si se mantiene la situación como está se genera el daño irreparable que necesita una respuesta inmediata.

\begin{normativa}{Artículo 232 --- CPCCN: medida cautelar genérica o innominada}
Contempla la posibilidad de peticionar medidas cautelares no previstas específicamente en la ley, cuando las circunstancias del caso lo requieren.
\end{normativa}

También se menciona la cautelar autónoma, proveniente del ámbito del contencioso-administrativo. Estos institutos se desarrollan en el capítulo~\ref{cap:generico}.

\section{Ámbito excepcional}

En el ámbito excepcional se encuentran la tutela anticipada y los sistemas protectores.

\begin{definition}{Tutela anticipada}
Pronunciamiento jurisdiccional de carácter provisorio a través del cual se adelanta, en todo o en parte, aquello que debió ser objeto de la sentencia de mérito, a fin de evitar un daño actual o inminente, o una situación afligente. A diferencia de las cautelares ordinarias, puede recaer sobre el mismo objeto de la sentencia definitiva.
\end{definition}

La tutela anticipada se desarrolla en el capítulo~\ref{cap:tutela}.
% TODO: verificar el desarrollo de los sistemas protectores, que las fuentes mencionan sin desarrollarlos en este punto.

\section{Panorama de los sistemas cautelares}

Los sistemas cautelares se desarrollan del siguiente modo:
\begin{itemize}
    \item \textbf{Parte general:} qué son, su tratamiento como sistema, los distintos tipos creados por el legislador, los presupuestos sustanciales y procesales para peticionar y obtener una medida cautelar, su trámite, las facultades del juez, los derechos del demandado y las características propias de las resoluciones dictadas en materia cautelar (capítulos~\ref{cap:fundamentos-cautelares} a~\ref{cap:caracteres}).
    \item \textbf{Sistemas precautorios tradicionales:} embargo, inventario, anotación de litis, secuestro y demás medidas del sistema precautorio tradicional reguladas en el Código Procesal (capítulos~\ref{cap:embargo} a~\ref{cap:proteccion}).
    \item \textbf{Sistema asegurativo probatorio:} es el que el legislador creó a través de la producción de prueba anticipada, regulada en el art.~326 CPCC.
    \item \textbf{Sistema asegurativo garantizador:} vinculado con la contracautela y la fianza (art.~258 CPCC).
    \item \textbf{Sistema asegurativo genérico, ámbito excepcional y tutela anticipada} (capítulos~\ref{cap:generico} y~\ref{cap:tutela}).
\end{itemize}
% TODO: contenido incompleto o ambiguo en las notas originales (el sistema asegurativo garantizador y el Art. 258 CPCC sólo figuran en el índice de una de las fuentes; no se desarrollan).

\cornelltitle{clasificación de los sistemas cautelares}
\cornellblock{
\cqrow{¿Por qué se estudian los sistemas cautelares como un «sistema»?}{Porque en distintos lugares del Código Procesal se verifican manifestaciones en las que la ley actúa con carácter conservatorio o asegurativo. Agruparlas en subsistemas permite comprenderlas en su conjunto, distinguirlas y entender el sistema del código procesal en esos aspectos.
}
\cqrow{¿Cómo se clasifican los sistemas cautelares según su ámbito?}{Ámbito normal: asegurativos probatorios, precautorios tradicionales, asegurativos garantizadores y asegurativos genéricos. Ámbito excepcional: tutela anticipada y sistemas protectores.
}
\cqrow{¿Qué resguardan los sistemas asegurativos probatorios?}{Propenden a la mayor eficacia del proceso, resguardando fuentes de prueba que pueden perderse con el tiempo. Se encuentran regulados en el art.~326 CPCCN (prueba anticipada), que permite adelantar la producción de una prueba que se perdería al aguardar la etapa probatoria ordinaria.
}
\cqrow{¿Cuál es la finalidad de los sistemas precautorios tradicionales?}{Resguardar el efectivo cumplimiento de la sentencia, fundamentalmente mediante el resguardo de bienes (embargo, anotación de litis, secuestro, entre otras). Sin embargo, la sentencia solo serviría para «colgarse en un cuadro».
}
\cqrow{¿Qué garantizan los sistemas asegurativos garantizadores?}{Garantizan a la parte afectada por la medida frente a la posibilidad de que sea revocada o no reconocida en la sentencia. El instrumento central es la contracautela; los arts.~509 y 555 CPCCN exigen fianza para ejecutar.
}
\cqrow{¿Qué comprenden los sistemas asegurativos genéricos?}{La inhibición general de bienes (art.~228, medida residual cuando no hay bienes para embargar), la prohibición de innovar (art.~230) y la medida cautelar genérica o innominada (art.~232). También se menciona la cautelar autónoma del ámbito contencioso-administrativo.
}
\cqrow{¿Cuáles son las dos facetas de la prohibición de innovar?}{Medida de no innovar: cuando el perjuicio proviene de modificar el statu quo (se pide conservar la situación). Medida innovativa: cuando el perjuicio proviene de mantener el statu quo (se pide modificarlo para evitar el daño irreparable).
}
}
\medskip
\cornellblocksum{
\cqrow{¿Qué es la tutela anticipada y por qué pertenece al ámbito excepcional?}{Es un pronunciamiento jurisdiccional provisorio que adelanta, en todo o en parte, lo que debió ser objeto de la sentencia de mérito, para evitar un daño actual o inminente o una situación afligente. A diferencia de las cautelares ordinarias, puede recaer sobre el mismo objeto de la sentencia definitiva.
}
}{los sistemas cautelares se agrupan en un ámbito normal (probatorios, precautorios tradicionales, garantizadores y genéricos) y un ámbito excepcional (tutela anticipada y sistemas protectores). Cada subsistema responde a una finalidad distinta: resguardar prueba, asegurar el cumplimiento de la sentencia, proteger al afectado por la medida o admitir medidas genéricas como la prohibición de innovar, que presenta una faceta de no innovar y otra innovativa.
}

\chapter[Medidas cautelares y presupuestos]{Medidas cautelares tradicionales y sus presupuestos}\label{cap:presupuestos}

\section{Concepto de medidas cautelares}

\begin{definition}{Medidas cautelares}
Actos procesales del órgano jurisdiccional, adoptados en el curso de un proceso o previos a él, para:
\begin{itemize}
    \item Asegurar bienes o pruebas.
    \item Mantener una determinada situación de hecho.
    \item Brindar seguridad o protección a una persona.
\end{itemize}
En todos los casos implican un \textbf{anticipo provisorio}, que no es definitivo, y que constituye uno de los caracteres esenciales de las medidas cautelares.
\end{definition}

\section{Presupuestos de las medidas cautelares}

Las medidas cautelares tienen presupuestos propios que determinan su procedencia. Estos presupuestos se exigen tanto para las medidas del sistema precautorio tradicional como para las del sistema asegurativo genérico.

\begin{table}[H]
\centering
\begin{tabularx}{\textwidth}{>{\raggedright\arraybackslash}p{0.25\textwidth} Y}
\toprule
\textbf{Tipo de presupuesto} & \textbf{Contenido} \\
\midrule
\textbf{Sustanciales} & Verosimilitud en el derecho y peligro en la demora; como consecuencia, la contracautela \\
\textbf{Procesales} & Enumerados en el art.~195 CPCCN: derecho a asegurar, medida solicitada, disposición legal aplicable y requisitos específicos de la medida \\
\bottomrule
\end{tabularx}
\end{table}

\section{Presupuestos sustanciales}

\subsection{Verosimilitud en el derecho (\textit{fumus boni iuris})}

El viejo adagio latino \textit{fumus boni iuris} se traduce como el «humo de buen derecho». Habla de la \textbf{apariencia}: es un juicio de valor que hace la jurisdicción, según el cual lo que aparenta es que la pretensión va a tener un buen resultado en la sentencia. Frente a esa apariencia, y frente al juicio de valor sobre el derecho que invoca quien peticiona la medida, el juez considera que está dado el presupuesto sustancial para admitir la cautelar. No tiene relación con la certeza a la que debe acceder el juez cuando dicta la sentencia definitiva.

\begin{important}
La verosimilitud del derecho \textbf{no es certeza}. Es un conocimiento inferior y provisional, que se basa en la apariencia de que la pretensión del peticionante está jurídicamente fundada. La certeza plena se exige solo para la sentencia definitiva.
\end{important}

\subsection{Peligro en la demora}

Es el temor fundado de que, con el transcurso del tiempo que insume un proceso judicial, el derecho que se invoca se torne estéril, es decir, que no se pueda dictar una sentencia eficaz. Se configura cuando existe un riesgo fundado de que, de no decretarse la cautelar, la sentencia que en su momento se dicte resulte ineficaz o de cumplimiento imposible.

\subsection{La contracautela}

\begin{definition}{Contracautela}
Garantía que debe otorgar quien pide la medida, para responder por los daños y perjuicios y las costas que pueda generar si la medida resulta luego rechazada en la acción de fondo o revocada por un tribunal superior.
\end{definition}

La contracautela guarda relación con los otros presupuestos y funciona como \textbf{vasos comunicantes}:

\begin{formula}{Contracautela como vasos comunicantes}
A mayor verosimilitud en el derecho y mayor peligro en la demora, menor será la contracautela. A menor verosimilitud en el derecho y menor peligro en la demora, mayor será la contracautela.
\end{formula}

\begin{table}[H]
\centering
\small
\begin{tabularx}{\textwidth}{>{\raggedright\arraybackslash}p{0.17\textwidth} Y Y}
\toprule
\textbf{Tipo} & \textbf{En qué consiste} & \textbf{Ejemplo} \\
\midrule
\textbf{Real} & Garantía sobre bienes concretos & Depósito en dinero, hipoteca sobre un bien, embargo sobre un inmueble \\
\textbf{Personal} & Fianza otorgada por una persona con reconocida capacidad económica & Fianza de un tercero solvente \\
\textbf{Juratoria} & Juramento del peticionante de responder por los daños & Supuestos de máxima verosimilitud en el derecho; es la caución menos exigente \\
\bottomrule
\end{tabularx}
\end{table}

\section{Presupuestos procesales: artículo 195 CPCCN}

\begin{normativa}{Artículo 195 --- CPCCN: oportunidad y presupuesto}
Las providencias cautelares podrán ser solicitadas antes o después de deducida la demanda, a menos que de la ley resultare que ésta debe entablarse previamente. El escrito deberá expresar: el derecho que se pretende asegurar, la medida que se pide, la disposición de la ley en que se funde, y el cumplimiento de los requisitos que correspondan en particular a la medida requerida.
\end{normativa}

Los cuatro presupuestos procesales que enumera el artículo 195 son:

\begin{enumerate}
    \item El derecho que se pretende asegurar.
    \item La medida que se pide (embargo, inhibición, secuestro, etc.).
    \item La disposición legal en que se funda.
    \item El cumplimiento de los requisitos que correspondan en particular a la medida requerida.
\end{enumerate}

\cornelltitle{presupuestos de las medidas cautelares}
\cornellblocksum{
\cqrow{¿Qué son las medidas cautelares?}{Actos procesales del órgano jurisdiccional, adoptados en el curso de un proceso o previos a él, para asegurar bienes o pruebas, mantener una determinada situación de hecho o brindar seguridad o protección a una persona. En todos los casos implican un anticipo provisorio, que no es definitivo.
}
\cqrow{¿Cuáles son los presupuestos de las medidas cautelares?}{Sustanciales: verosimilitud en el derecho y peligro en la demora, con la contracautela como consecuencia. Procesales: los enumerados en el art.~195 CPCCN.
}
\cqrow{¿Qué es el \textit{fumus boni iuris}?}{Es el «humo de buen derecho»: la apariencia de que la pretensión del peticionante tiene fundamento jurídico suficiente. Es un juicio de valor de la jurisdicción que no exige certeza plena, sino verosimilitud, y que habilita el dictado de la cautelar.
}
\cqrow{¿Qué es el peligro en la demora?}{Es el temor fundado de que, con el transcurso del tiempo que insume el proceso, el derecho invocado se torne estéril y la sentencia resulte ineficaz o de cumplimiento imposible.
}
\cqrow{¿Qué es la contracautela y de qué tipos puede ser?}{Es la garantía que debe otorgar quien pide la medida para responder por daños, perjuicios y costas si la medida es rechazada en la acción de fondo o revocada por un tribunal superior. Puede ser real (garantía sobre bienes concretos), personal (fianza de una persona con reconocida capacidad económica) o juratoria (juramento de responder por los daños; la menos exigente).
}
\cqrow{¿Cómo funciona la contracautela como vasos comunicantes?}{A mayor verosimilitud y mayor peligro en la demora, menor contracautela; a menor verosimilitud y menor peligro en la demora, mayor contracautela.
}
\cqrow{¿Qué debe expresar el escrito según el art.~195 CPCCN?}{1) El derecho que se pretende asegurar. 2) La medida que se pide. 3) La disposición legal en que se funda. 4) El cumplimiento de los requisitos particulares de la medida requerida.
}
}{las medidas cautelares son actos procesales que implican un anticipo provisorio. Proceden cuando concurren los presupuestos sustanciales (verosimilitud en el derecho y peligro en la demora), de los que se sigue la contracautela, que opera en forma inversa a la intensidad de aquellos, y los presupuestos procesales del art.~195 CPCCN.
}

\chapter[Caracteres y clasificación de las medidas cautelares]{Caracteres particulares y clasificación de las medidas cautelares}\label{cap:caracteres}

\section{Provisoriedad: artículo 202 CPCCN}

\begin{normativa}{Artículo 202 --- CPCCN: provisoriedad}
Las medidas cautelares son provisorias. Subsistirán mientras duren las circunstancias que las determinaron. En cualquier momento en que éstas cesaren, se podrá requerir su levantamiento.
\end{normativa}

Esta característica diferencia esencialmente a la cautela de la sentencia: cuando cambian las circunstancias que dieron origen al dictado de la medida cautelar, ésta puede modificarse o levantarse. La sentencia, en cambio, pone fin al proceso y no puede alterarse.

Lo que se decide en materia cautelar es \textbf{provisorio}: la medida puede dejarse sin efecto, modificarse, reducirse o ampliarse a lo largo del proceso, según se modifiquen o no las circunstancias de hecho y de derecho que existían al momento en que el juez la ordenó.

\begin{important}
No deben confundirse ambos planos: lo que se resuelve en la \textbf{sentencia} es \textbf{definitivo}; lo que se resuelve en materia \textbf{cautelar} es \textbf{provisorio}.
\end{important}

\section{Modificabilidad: artículo 203 CPCCN}

\begin{normativa}{Artículo 203 --- CPCCN: modificación}
El afectado por una medida cautelar puede pedir su sustitución por otra que le resulte menos gravosa, ofreciendo bienes que resulten igualmente suficientes. También puede peticionarse la ampliación o mejora de la medida.
\end{normativa}

\begin{example}{Sustitución de un embargo}
Si se embargó un inmueble y este embargo genera un daño innecesario al afectado, éste puede ofrecer otro bien a embargo que resulte igualmente suficiente al que fue embargado; la medida puede sustituirse para evitar el daño que genera la medida dictada.
\end{example}

\section{Facultades del juez: artículo 204 CPCCN}

\begin{normativa}{Artículo 204 --- CPCCN: facultades del juez}
El juez, para evitar perjuicios o gravámenes innecesarios al titular de los bienes, podrá disponer una medida precautoria distinta de la solicitada, o limitarla, teniendo en cuenta la importancia del derecho que se intentare proteger.
\end{normativa}

El juez tiene plenas facultades para modificar o limitar la medida cautelar solicitada, siempre resguardando el derecho que se intenta proteger. Esto también distingue a la cautelar de la sentencia definitiva.

\section{Caducidad: artículo 207 CPCCN}

\begin{normativa}{Artículo 207 --- CPCCN: caducidad de las medidas cautelares}
Cuando las medidas cautelares se hayan peticionado antes del inicio del proceso principal, a los 10 días siguientes de su traba (o de su inscripción registral o notificación, según corresponda), debe promoverse la demanda. La inhibición general de bienes y el embargo se extinguen a los 5 años desde la fecha de inscripción en el registro respectivo. Para evitar la extinción, debe pedirse la reinscripción de la medida antes del vencimiento de ese plazo, lo que genera un nuevo período de 5 años.
\end{normativa}

Mecanismo de reinscripción: antes de que venzan los 5 años de la inscripción original, se solicita la reinscripción; el oficio correspondiente se libra y se inscribe en el registro, generando un nuevo plazo de 5 años desde esa nueva inscripción.

\section{Tramitación \textit{inaudita parte} y recursos: artículo 198 CPCCN}

\begin{normativa}{Artículo 198 --- CPCCN: cumplimiento de la medida y recursos}
Las medidas precautorias se decretarán y cumplirán sin audiencia de la otra parte (\textit{inaudita parte}). Ningún incidente planteado por el destinatario de la medida podrá detener su cumplimiento. Si el afectado no hubiese tomado conocimiento de las medidas con motivo de su ejecución, se les notificarán personalmente o por cédula dentro de los 3 días. Quien hubiese obtenido la medida será responsable de los perjuicios que generare la demora en notificarla. La providencia que admitiere o denegare una medida cautelar será recurrible por vía de reposición; también será admisible la apelación subsidiaria o directa. El recurso de apelación, en caso de admitirse la medida, será concedido con efecto devolutivo (no suspensivo).
\end{normativa}

La medida cautelar tramita \textit{inaudita parte}, lo que también significa que no se está ante una hipótesis de condena, sino de cautela. Es obligación de quien obtuvo la cautelar notificarla dentro del plazo de 3 días de trabada, bajo responsabilidad por los perjuicios que genere la demora.

\begin{important}
El recurso de apelación en materia cautelar se concede con efecto \textbf{devolutivo} (no suspensivo): la vigencia de la medida cautelar continúa mientras se sustancia el recurso. Aquí radica la importancia de la contracautela: si el tribunal superior revoca la resolución, la contracautela permite responder por los daños durante el tiempo en que la medida estuvo vigente (art.~199 CPCCN).
\end{important}

\section{Clasificación sobre bienes, entes ideales y personas}

Para hacer más didáctico el estudio, las medidas cautelares se clasifican según el objeto sobre el que recaen. Esta clasificación es meramente didáctica y el detalle de cada medida se desarrolla en los capítulos sobre medidas cautelares en particular.

\begin{table}[H]
\centering
\small
\begin{tabularx}{\textwidth}{>{\raggedright\arraybackslash}p{0.2\textwidth} Y Y}
\toprule
\textbf{Objeto} & \textbf{Grado / tipo} & \textbf{Medida específica} \\
\midrule
\textbf{Bienes} & Individualización del bien & Inventario / anotación de litis \\
\textbf{Bienes} & Indisposición del bien & Embargo / prohibición de contratar \\
\textbf{Bienes} & Desapoderamiento del bien & Secuestro / interventor recaudador \\
\textbf{Entes ideales} (empresas / sociedades) & Interferencia simple & Veedores, interventores informativos \\
\textbf{Entes ideales} & Intervención compleja & Administración judicial de la sociedad \\
\textbf{Personas} & Sistema genérico / residual & Guarda de personas (art.~234 CPCCN) \\
\textbf{Personas} & Protección específica & Protección contra violencia familiar (Ley 24.417) \\
\bottomrule
\end{tabularx}
\end{table}

\begin{normativa}{Artículo 234 --- CPCCN: guarda de personas}
Regula la guarda de personas como medida cautelar residual sobre personas. Se aplica cuando las circunstancias del caso lo requieren.
\end{normativa}

\begin{normativa}{Ley 24.417 --- Protección contra la Violencia Familiar}
Esta ley, de aplicación en el ámbito de Capital Federal, prevé medidas cautelares específicas de protección para personas en situación de violencia doméstica.
\end{normativa}

\cornelltitle{caracteres y clasificación de las medidas cautelares}
\cornellblock{
\cqrow{¿Qué establece el art.~202 CPCCN sobre la provisoriedad?}{Las medidas cautelares son provisorias: subsisten mientras duren las circunstancias que las determinaron y, cuando estas cesan, se puede requerir su levantamiento. Esto las diferencia de la sentencia, que pone fin al proceso y no puede alterarse.
}
\cqrow{¿Qué permite el art.~203 CPCCN (modificabilidad)?}{Que el afectado pida la sustitución de la medida por otra menos gravosa, ofreciendo bienes igualmente suficientes, y que se peticione la ampliación o mejora de la medida.
}
\cqrow{¿Qué facultades otorga el art.~204 CPCCN al juez?}{Para evitar perjuicios o gravámenes innecesarios al titular de los bienes, puede disponer una medida distinta de la solicitada o limitarla, teniendo en cuenta la importancia del derecho que se intenta proteger.
}
\cqrow{¿Qué establece el art.~207 CPCCN sobre la caducidad?}{Si la cautelar se peticionó antes del inicio del proceso principal, la demanda debe promoverse a los 10 días siguientes de su traba (o de su inscripción registral o notificación, según corresponda). La inhibición general de bienes y el embargo se extinguen a los 5 años desde su inscripción registral, salvo que se pida la reinscripción antes del vencimiento, lo que genera un nuevo período de 5 años.
}
\cqrow{¿Qué significa que la cautelar tramita \textit{inaudita parte} (art.~198 CPCCN)?}{Que se decreta y cumple sin audiencia de la otra parte, y ningún incidente del destinatario puede detener su cumplimiento. Si el afectado no tomó conocimiento por la ejecución, se le notifica personalmente o por cédula dentro de los 3 días, y quien obtuvo la medida responde por los perjuicios que cause la demora en notificarla.
}
\cqrow{¿Cómo se recurre la resolución cautelar y qué efecto tiene la apelación?}{La providencia que admite o deniega la medida es recurrible por reposición, con apelación subsidiaria o directa. Si se admite la medida, la apelación se concede con efecto devolutivo (no suspensivo): la medida sigue vigente mientras se sustancia el recurso, lo que explica la importancia de la contracautela.
}
}
\medskip
\cornellblocksum{
\cqrow{¿Cómo se clasifican las medidas cautelares según su objeto?}{Sobre bienes (individualización: inventario / anotación de litis; indisposición: embargo / prohibición de contratar; desapoderamiento: secuestro / interventor recaudador). Sobre entes ideales (interferencia simple: veedores e interventores informativos; intervención compleja: administración judicial de la sociedad). Sobre personas (guarda de personas, art.~234 CPCCN; protección contra violencia familiar, Ley 24.417).
}
}{las medidas cautelares son provisorias, modificables y sujetas a caducidad; tramitan \textit{inaudita parte} y su apelación se concede con efecto devolutivo, de modo que la contracautela cobra importancia. Estos caracteres las distinguen de la sentencia definitiva. Además, se clasifican didácticamente según recaigan sobre bienes, entes ideales o personas, y cada medida se estudia en detalle en los capítulos sobre medidas cautelares en particular.
}

\chapter[Sistema asegurativo genérico]{Sistema asegurativo genérico (ámbito normal)}\label{cap:generico}

\section{Caracterización general}

El sistema asegurativo genérico comprende cuatro medidas:

\begin{enumerate}
    \item \textbf{Inhibición general de bienes} (Art.~228 CPCC).
    \item \textbf{Medida cautelar genérica} (Art.~232 CPCC).
    \item \textbf{Prohibición de innovar}, que incluye una faceta innovativa y otra no innovativa (Art.~230 CPCC).
    \item \textbf{Cautela autónoma}, creación pretoriana luego regulada por la Ley 26.854.
\end{enumerate}

\begin{important}
\textbf{Carácter residual o subsidiario.} Las medidas de este sistema sólo prosperan cuando las medidas del sistema precautorio tradicional (embargo, secuestro, etc.) no resultan aptas para la situación de hecho y de derecho que se pretende proteger en el caso concreto.
\end{important}

\section{Inhibición general de bienes (Art.~228 CPCC)}

\begin{quote}
\small
\textbf{Artículo 228 --- Código Procesal Civil y Comercial de la Nación.} \textit{Inhibición general de bienes. En todos los casos en que habiendo lugar a embargo éste no pudiere hacerse efectivo por no conocerse bienes del deudor, o por no cubrir éstos el importe del crédito reclamado, podrá solicitarse contra aquél la inhibición general de enajenar y gravar bienes, la que se deberá dejar sin efecto siempre que presentare a embargo bienes suficientes o diere caución bastante. El que solicitare la inhibición deberá expresar el nombre, apellido y domicilio del deudor, así como todo otro dato que pueda individualizar al inhibido, sin perjuicio de los demás requisitos que impongan las leyes. La inhibición sólo producirá efecto desde la fecha de su anotación salvo para los casos en que el dominio se hubiere transmitido con anterioridad, de acuerdo con lo dispuesto en la legislación general. No concederá preferencia sobre las anotadas con posterioridad.}
\end{quote}

\subsection{Definición y finalidad}

\begin{definition}{Inhibición general de bienes}
Medida cautelar genérica, regulada en un solo artículo (el 228 CPCC), mediante la cual se pretende que el demandado o afectado por la medida \textbf{no venda ni grave} bienes que formen parte de su patrimonio.
\end{definition}

\subsection{Carácter residual: cuándo procede}

La inhibición general de bienes \textbf{no es una medida autónoma o principal} como el embargo, sino residual. Procede en los mismos casos en que procede el embargo, pero únicamente cuando se da alguno de estos supuestos:

\begin{enumerate}
    \item No existen bienes para embargar (el deudor no tiene bienes identificados).
    \item Existen bienes para embargar, pero no resultan suficientes para responder o garantizar el crédito que se pretende asegurar.
\end{enumerate}

\begin{example}{Embargo insuficiente}
Si se debe asegurar una suma equivalente a un millón de pesos y sólo pudieron embargarse quinientos mil pesos en una cuenta bancaria, sin otros bienes para embargar, queda por garantizar el cumplimiento de los quinientos mil pesos restantes. La inhibición general de bienes puede ir junto con el embargo de manera residual, porque los bienes embargados no resultan suficientes.
\end{example}

\subsection{Diferencias con el embargo}

\medskip
\renewcommand{\arraystretch}{1.3}
\noindent\begin{tabularx}{\textwidth}{@{}>{\raggedright\arraybackslash}p{2.8cm} Y Y@{}}
\toprule
\textbf{Aspecto} & \textbf{Embargo} & \textbf{Inhibición general de bienes} \\
\midrule
Determinación &
\textbf{Determinado:} recae sobre un bien identificado y particular. &
\textbf{Indeterminada:} afecta a todos los bienes presentes y futuros del patrimonio del deudor. \\
\addlinespace
Tipo de bienes &
Todo tipo de bienes: muebles, inmuebles, registrables y no registrables. &
Sólo bienes \textbf{registrables} (porque debe inscribirse para producir efectos). \\
\addlinespace
Preferencia temporal &
El primer embargante tiene derecho de preferencia sobre los posteriores. &
No existe derecho de preferencia; no rige el principio de prioridad temporal. \\
\addlinespace
Habilita ejecución forzada &
Sí: permite avanzar hacia el remate y cobro del crédito. &
No: no habilita la ejecución forzada porque no hay bienes identificados de los que cobrar. \\
\addlinespace
Extinción &
5 años desde la inscripción. &
5 años desde la inscripción (igual que el embargo). \\
\addlinespace
Caducidad si es previa al proceso &
10 días para iniciar el proceso principal. &
10 días para iniciar el proceso principal (igual que el embargo). \\
\bottomrule
\end{tabularx}
\renewcommand{\arraystretch}{1}

\subsection{Alcance temporal y efectos}

La inhibición produce efectos desde su inscripción en adelante. Es posible que se desconozcan los bienes que tiene el deudor o demandado, o que éste no tenga bienes al momento de la medida y luego ingresen bienes a su patrimonio, que quedarán alcanzados por la inhibición.

Tanto el embargo como la inhibición se extinguen a los \textbf{cinco años} de su inscripción. Para mantener la validez de la medida, antes de su vencimiento debe peticionarse al juez la \textbf{reinscripción}, a fin de evitar su extinción.

\subsection{Procedimiento}

\begin{enumerate}
    \item Se peticiona que se dicte la inhibición general de bienes en determinados registros, acreditando los presupuestos cautelares.
    \item El juez ordena y decreta la inhibición, identificando al demandado por nombre completo y número de DNI.
    \item El letrado confecciona un oficio dirigido al registro correspondiente (por ejemplo, el Registro de la Propiedad Inmueble), que se deja en el juzgado para su revisión y firma.
    \item El letrado diligencia el oficio ingresándolo al registro para que se inscriba la inhibición.
\end{enumerate}

Suele pedirse informes en los registros pertinentes para conocer qué bienes tiene el demandado a su nombre y con su DNI: el Registro de la Propiedad Inmueble de la Ciudad de Buenos Aires y de la Provincia de Buenos Aires, y el Registro de la Propiedad Automotor, de alcance nacional.

\subsection{Cese de la inhibición}

La inhibición puede levantarse antes de los cinco años en los siguientes casos:

\begin{enumerate}
    \item Si el deudor ofrece un bien a embargo que sea suficiente para garantizar el crédito.
    \item Si el deudor ofrece una garantía suficiente (en el ámbito de una medida cautelar).
\end{enumerate}

El pedido de levantamiento por sustitución, formulado por el deudor, se traslada al actor o acreedor, y luego el juez resuelve; fundamentalmente porque el bien ofrecido a embargo debe ser suficiente para garantizar el crédito en cuestión.

\section{Medida cautelar genérica o innominada (Art.~232 CPCC)}

\begin{quote}
\small
\textbf{Artículo 232 --- Código Procesal Civil y Comercial de la Nación.} \textit{Fuera de los casos previstos en los artículos precedentes, quien tuviere motivos para temer que durante el tiempo anterior al reconocimiento judicial de su derecho, éste pudiera sufrir un perjuicio inminente o irreparable, podrá solicitar las medidas urgentes que, según las circunstancias, fuesen más aptas para asegurar provisionalmente el cumplimiento de la sentencia.}
\end{quote}

La medida cautelar genérica establece que, fuera de los casos previstos en los artículos precedentes (todas las medidas cautelares ya estudiadas), quien tenga motivos para temer que, antes de la sentencia, su derecho pueda sufrir un perjuicio inminente o irreparable, puede solicitar las medidas urgentes que según las circunstancias resulten más aptas para asegurar provisionalmente el cumplimiento de la sentencia.

Se la llama \textbf{genérica} o \textbf{innominada} porque habilita a la jurisdicción a dictar cualquier medida, ya sea sola o aislada, o en complemento de las tradicionales, que mejor responda a las situaciones de hecho y de derecho que se pretende asegurar y proteger antes de la sentencia.

\begin{important}
La medida cautelar genérica rige para asegurar cualquier tipo de bienes, personas o derechos. Es la que más amplía las posibilidades de peticionar una medida cautelar cuando las reguladas anteriormente no resultan adecuadas y suficientes.
\end{important}

\section{Prohibición de innovar (Art.~230 CPCC)}

\subsection{La prohibición de innovar como medida cautelar matriz}

Se entiende que la prohibición de innovar es un sistema asegurativo \textbf{matriz} de los restantes, la ``medida cautelar madre'' de todas las demás. Muchos autores atribuyen ese carácter al embargo, pero se considera más adecuado afirmar que el sistema matriz es la prohibición de innovar, porque permite cautelar no sólo bienes, sino también personas y derechos, según los hechos del caso concreto.

\begin{quote}
\small
\textbf{Artículo 230 --- Código Procesal Civil y Comercial de la Nación.} \textit{Podrá decretarse la prohibición de innovar en toda clase de juicios siempre que: 1. El derecho resulte verosímil. 2. Existiere peligro de que si se mantuviera o alterara, en su caso, la situación de hecho o de derecho, la modificación pudiera influir en la sentencia o convirtiera su ejecución en ineficaz. 3. La cautela no pudiere obtenerse por medio de otra medida precautoria.}
\end{quote}

\subsection{Las dos facetas de la prohibición de innovar}

El inciso 2 del Art.~230 revela las dos facetas de esta medida. Algunos autores las consideran dos vías distintas (medida innovativa y de no innovar); sin embargo, según cómo las reguló el legislador, ambas son facetas de la prohibición de innovar.

\medskip
\renewcommand{\arraystretch}{1.3}
\noindent\begin{tabularx}{\textwidth}{@{}>{\raggedright\arraybackslash}p{2.6cm} Y Y@{}}
\toprule
\textbf{Faceta} & \textbf{Cuándo procede} & \textbf{Qué ordena el juez} \\
\midrule
Innovativa (medida innovativa) &
El perjuicio proviene de \textbf{mantener} el statu quo existente. &
Ordena \textbf{alterar} la situación de hecho o de derecho. \\
\addlinespace
No innovativa (medida de no innovar) &
El perjuicio proviene de \textbf{alterar} el statu quo existente. &
Ordena \textbf{mantener} la situación de hecho o de derecho. \\
\bottomrule
\end{tabularx}
\renewcommand{\arraystretch}{1}

\subsection{Los dos ámbitos en que opera}

Además de sus dos facetas, la prohibición de innovar puede operar en dos ámbitos distintos:

\begin{enumerate}
    \item \textbf{Ámbito normal:} cuando la pretensión cautelar \textbf{no} se superpone con la pretensión sustancial del proceso principal.
    \item \textbf{Ámbito excepcional:} cuando la pretensión cautelar \textbf{se superpone} total o parcialmente con la pretensión sustancial; esto configura la \textbf{tutela anticipada}.
\end{enumerate}

\section{Cautela autónoma}

\subsection{Ubicación en el sistema y origen}

La cautela autónoma es la cuarta medida del sistema asegurativo genérico. Es propia del derecho administrativo y de la justicia en lo contencioso-administrativo; no está regulada en el Código Procesal Civil y Comercial de la Nación. En un primer momento no tuvo regulación y fue una \textbf{creación pretoriana} jurisprudencial; luego recibió una regulación legal bastante cuestionable.

\subsection{En qué consiste}

\begin{definition}{Cautela autónoma}
Petición que el justiciable hace al juez para que \textbf{suspenda los efectos de un acto administrativo} hasta que se resuelva el recurso administrativo interpuesto contra ese acto y que, de algún modo, habilitó la instancia judicial, es decir, hasta que se agote la vía administrativa.
\end{definition}

\subsection{La presunción de legitimidad (Art.~12 Ley 19.549)}

\begin{quote}
\small
\textbf{Artículo 12 --- Ley 19.549 (Procedimiento Administrativo).} \textit{El acto administrativo goza de presunción de legitimidad; su fuerza ejecutoria faculta a la Administración a ponerlo en práctica por sus propios medios, salvo que la ley o la naturaleza del acto exigieren la intervención judicial, e impide que los recursos que interpongan los administrados suspendan su ejecución y efectos, salvo que una norma expresa establezca lo contrario. Sin embargo, la Administración podrá, de oficio o a pedido de parte y mediante resolución fundada, suspender la ejecución por razones de interés público, o para evitar perjuicios graves al interesado.}
\end{quote}

El acto administrativo goza de presunción de legitimidad. Por ello, aunque el administrado interponga un recurso administrativo para que el acto sea revisado, ese recurso no suspende los efectos del acto, que puede ser ejecutado. Esto se vincula con los efectos suspensivos y no suspensivos de los recursos: el recurso administrativo no tiene efecto suspensivo.

Esta situación puede generar un daño o perjuicio irreparable en el administrado; por ejemplo, si el acto ordenó abonar una suma o multa de importancia trascendente. En ese caso, estando pendiente de resolución el recurso administrativo, el administrado puede acudir a la justicia y peticionar la cautela autónoma, es decir, que el juez disponga una medida que suspenda los efectos del acto administrativo hasta que se resuelva el recurso.

\subsection{Por qué se llama ``autónoma''}

Se llama autónoma porque se trata de un procedimiento administrativo y la medida es autónoma respecto de él: no hay un proceso judicial en curso. Es instrumental y accesoria, porque se vincula con la suspensión de los efectos del acto administrativo.

\subsection{Creación pretoriana y trámite original}

Originalmente se promovían amparos para obtener este tipo de cautelas, hasta que la Corte resolvió que el amparo no era la vía idónea, y entonces se creó la cautela autónoma. El trámite original se encuadró en el Art.~230 CPCC (prohibición de innovar) y los presupuestos requeridos eran:

\begin{enumerate}
    \item Verosimilitud del derecho.
    \item Peligro en la demora.
\end{enumerate}

El carácter de la medida era \textbf{más excepcional} que el de las medidas ordinarias, dado que el acto administrativo goza de presunción de legitimidad.

\subsection{La Ley 26.854: medidas cautelares contra el Estado}

Lo estudiado hasta aquí sobre medidas cautelares no rige cuando el Estado está del otro lado. En ese caso debe tenerse en cuenta la \textbf{Ley 26.854}, que regula cuestiones particulares, distintas de las ya vistas, cuando el justiciable pretende una medida cautelar contra el Estado.

\begin{important}
\textbf{Artículo 13 --- Ley 26.854 (suspensión de efectos de un acto estatal).} Regula la cautela autónoma en el fuero contencioso-administrativo. Para peticionar la suspensión de los efectos de un acto estatal ya no bastan la verosimilitud del derecho y el peligro en la demora. Se exige:
\begin{enumerate}
    \item Acreditar sumariamente perjuicios graves que impliquen imposible reparación posterior.
    \item Que tramite \textbf{bilateralizada}: el juez debe convocar previamente a la entidad estatal a presentar un informe.
    \item Que el administrado haya solicitado previamente en sede administrativa la suspensión del acto; sólo ante la negativa o el silencio de cinco días queda habilitada la vía judicial.
\end{enumerate}
\end{important}

En cuanto al trámite, según el Art.~4 de la ley, cualquier medida cautelar peticionada contra el Estado deja de tramitar \textbf{inaudita parte}: siempre, antes de resolver, el juez debe convocar a la entidad estatal, que presenta un informe pormenorizado de la cuestión peticionada, y luego resuelve.

Esta regulación podría ser atacada por ciertas inconstitucionalidades y, al igual que la regulación de la tutela anticipada de urgencia que se analiza más adelante, muchas veces se legisla en contra del justiciable.

\cornelltitle{sistema asegurativo genérico}
\cornellblock{
\cqrow{¿Qué es el carácter residual o subsidiario del sistema asegurativo genérico?}{Las medidas del sistema (inhibición general de bienes, medida cautelar genérica, prohibición de innovar) sólo proceden cuando las medidas del sistema precautorio tradicional (embargo, secuestro, etc.) no resultan aptas para proteger la situación de hecho y de derecho del caso concreto.
}
\cqrow{¿Cuándo procede la inhibición general de bienes?}{En los mismos casos en que procede el embargo, pero sólo cuando no existen bienes para embargar o cuando los bienes embargables no alcanzan para garantizar el crédito. Busca que el demandado no venda ni grave bienes de su patrimonio.
}
\cqrow{¿En qué se diferencia la inhibición del embargo?}{La inhibición es indeterminada (afecta bienes presentes y futuros), alcanza sólo bienes registrables, no otorga preferencia temporal y no habilita la ejecución forzada. El embargo es determinado, recae sobre cualquier tipo de bien, otorga preferencia al primer embargante y habilita el remate. Ambos se extinguen a los 5 años de su inscripción y, si son previos al proceso, caducan a los 10 días si no se inicia el proceso principal.
}
\cqrow{¿Cómo cesa o se mantiene la inhibición?}{Se levanta antes de los 5 años si el deudor ofrece un bien suficiente a embargo o una garantía suficiente; se corre traslado al acreedor y el juez resuelve. Para mantenerla debe reinscribirse antes de su vencimiento.
}
\cqrow{¿Qué es la medida cautelar genérica o innominada?}{Es la del Art.~232 CPCC: permite solicitar las medidas urgentes que según las circunstancias sean más aptas para asegurar provisionalmente el cumplimiento de la sentencia, fuera de los casos previstos en los artículos anteriores, ya sea de forma aislada o complementando las tradicionales. Sirve para cautelar bienes, personas o derechos.
}
\cqrow{¿Cuáles son las dos facetas de la prohibición de innovar?}{Innovativa: el perjuicio proviene de mantener el statu quo, por lo que el juez ordena alterarlo. No innovativa: el perjuicio proviene de alterarlo, por lo que el juez ordena mantenerlo. Ambas surgen del inciso 2 del Art.~230 CPCC.
}
\cqrow{¿Por qué se considera a la prohibición de innovar la medida cautelar matriz?}{Porque permite cautelar no sólo bienes sino también personas y derechos, según los hechos del caso concreto.
}
}
\medskip
\cornellblocksum{
\cqrow{¿Qué es la cautela autónoma y por qué surge?}{Es la petición al juez de suspender los efectos de un acto administrativo hasta que se resuelva el recurso administrativo interpuesto. Surge porque, según el Art.~12 de la Ley 19.549, el acto goza de presunción de legitimidad y los recursos no suspenden su ejecución, lo que puede causar un perjuicio irreparable. Fue una creación pretoriana (antes se usaba el amparo, que la Corte consideró no idóneo).
}
\cqrow{¿Qué cambió con la Ley 26.854?}{Para las medidas contra el Estado, el trámite deja de ser inaudita parte (el juez convoca a la entidad, que presenta un informe) y, para suspender un acto estatal (Art.~13), se exige acreditar perjuicios graves de imposible reparación posterior y haber solicitado antes la suspensión en sede administrativa, con negativa o silencio de 5 días.
}
}{El sistema asegurativo genérico es residual y subsidiario: comprende la inhibición general de bienes (Art.~228, indeterminada, sobre bienes registrables, sin preferencia ni ejecución forzada, por 5 años), la medida cautelar genérica (Art.~232, la medida más apta según las circunstancias), la prohibición de innovar (Art.~230, con facetas innovativa y no innovativa, medida matriz) y la cautela autónoma (suspensión de efectos de actos administrativos, hoy condicionada por la Ley 26.854 con trámite bilateralizado y requisitos agravados).
}

\chapter{Ámbito excepcional: tutela anticipada}\label{cap:tutela}

\section{Qué determina que un ámbito sea excepcional}

El ámbito excepcional se da cuando, a través de la petición y concesión de la cautela, se adelanta en todo o en parte aquello que luego será decidido en la sentencia de mérito. La pretensión cautelar se superpone total o parcialmente con la pretensión sustancial (la invocada en la demanda, que es el objeto del proceso). Lo que se peticiona o se concede cautelarmente es total o parcialmente igual a lo que se pretende que se reconozca al dictar sentencia.

\begin{definition}{Tutela anticipada}
Sistema cautelar en virtud del cual la jurisdicción dicta una medida para proteger y asegurar anticipadamente (de modo \textbf{provisorio}), total o parcialmente, aquello que después se decidirá de modo \textbf{definitivo} al dictar sentencia.
\end{definition}

\section{¿El juez no estaría prejuzgando?}

Cabe preguntarse si es válido este sistema o si el juez estaría prejuzgando al otorgarlo antes de tiempo. La clave está en no confundir el carácter \textbf{provisorio} de lo que se decide cautelarmente con el carácter \textbf{definitivo} de lo que se decide en la sentencia.

La superposición es posible porque la tutela es provisoria: para otorgarla el juez sólo exige que se acredite verosimilitud en el derecho y, en virtud del peligro en la demora y de la urgencia, protege y asegura. Luego, en el debate que se desarrolla a lo largo del proceso (defensa del afectado por la medida, pruebas que se produzcan), al dictar sentencia el juez puede resolver en el mismo sentido, pero ya con carácter definitivo, o en sentido contrario, y entonces queda sin efecto lo que se concedió provisionalmente.

\medskip
\renewcommand{\arraystretch}{1.3}
\noindent\begin{tabularx}{\textwidth}{@{}>{\raggedright\arraybackslash}p{2.6cm} Y Y Y@{}}
\toprule
\textbf{Momento} & \textbf{Decisión} & \textbf{Estándar de prueba} & \textbf{Carácter} \\
\midrule
Al otorgar la tutela anticipada &
Protege provisionalmente lo que será luego objeto de sentencia. &
Verosimilitud en el derecho y peligro en la demora. &
\textbf{Provisorio:} puede modificarse. \\
\addlinespace
Al dictar sentencia definitiva &
Decide con pleno conocimiento del caso. &
Certeza, basada en la prueba producida. &
\textbf{Definitivo:} cosa juzgada. \\
\bottomrule
\end{tabularx}
\renewcommand{\arraystretch}{1}

\medskip
Cuando el juez otorga la tutela anticipada de modo provisorio hay \textbf{cautela}; cuando decide en la sentencia definitiva sobre esa misma pretensión hay \textbf{condena}.

\section{Vías de materialización de la tutela anticipada}

La tutela anticipada se presenta a través de dos vías principales:

\begin{enumerate}
    \item \textbf{Sistemas cautelares (medidas cautelares):} mediante la prohibición de innovar (Art.~230 CPCC), la medida cautelar genérica (Art.~232 CPCC) o cualquier medida cautelar. El juez actúa dentro del ámbito excepcional por la superposición de pretensiones.
    \item \textbf{Subsistemas expresamente creados por el legislador:} normas sustanciales o procesales que regulan explícitamente la anticipación de la tutela para situaciones específicas.
\end{enumerate}

\section{Vías legales específicas de tutela anticipada}

\subsection{Alimentos provisorios (Art.~544 CCyC)}

\begin{quote}
\small
\textbf{Artículo 544 --- Código Civil y Comercial de la Nación.} \textit{Desde el principio de la causa o en el transcurso de ella el juez puede decretar la prestación de alimentos provisionales.}
\end{quote}

\begin{example}{Juicio de alimentos}
Una madre, en representación de sus hijos menores, demanda a su cónyuge, padre de los menores, para que se fije una cuota alimentaria a favor de los hijos, porque el progenitor demandado no abona ninguna suma por ese concepto. Dado el tiempo que demoran estos procesos y el perjuicio que supondría que durante años los menores no contaran con cuota alguna, el legislador creó el sistema de tutela anticipada del Art.~544 CCyC, denominado \textbf{alimentos provisorios}.

El juez puede otorgar de modo provisional un alimento a lo largo del proceso. Esa pretensión es igual, de algún modo, a la pretensión sustancial que el juez debe decidir en la sentencia; sin embargo, es plenamente viable y válida porque los alimentos provisionales son provisorios. Al dictar sentencia, luego de conocer las pruebas de las partes, el juez podrá fijar un alimento similar al provisional, uno considerablemente mayor o menor, o rechazar la demanda.
\end{example}

\subsection{Entrega anticipada del inmueble en el desalojo (Art.~680 bis CPCC)}

\begin{quote}
\small
\textbf{Artículo 680 bis --- Código Procesal Civil y Comercial de la Nación.} \textit{En los casos en que la acción de desalojo se dirija contra el intruso, en cualquier estado del juicio, después de trabada la litis, el juez puede imponer la inmediata entrega del inmueble, si el derecho invocado fuere verosímil, previa caución brindada por el actor por los eventuales daños y perjuicios que se pudieran irrogar.}
\end{quote}

\begin{example}{Proceso de desalojo contra el intruso}
La pretensión sustancial principal en el proceso de desalojo es la desocupación y entrega del inmueble. El Art.~680 bis permite que se adelante cautelarmente aquello que también es objeto de la sentencia de mérito: la desocupación y entrega anticipada del inmueble. Es provisorio, porque luego el juez tomará una decisión definitiva al dictar sentencia.

Por eso el Art.~680 bis exige verosimilitud en el derecho y que se conceda una caución (contracautela) para responder por los eventuales daños y perjuicios que pudiera generar al demandado una sentencia posterior en contra.
\end{example}

\section{Jurisprudencia de la Corte Suprema}

\subsection{Camacho Acosta (CSJN 320:1633): faceta innovativa}

\textbf{Hechos del caso.} El señor Camacho Acosta era un trabajador que prestaba servicios en una empresa. A raíz de un accidente de trabajo le amputaron el brazo izquierdo. En ese momento no estaba vigente ninguna póliza de aseguradora de riesgos de trabajo, por lo que promovió contra la empresa un juicio por daños y perjuicios reclamando todos los daños provenientes del accidente.

\textbf{Pretensión cautelar.} Camacho Acosta necesitaba colocar en su brazo izquierdo la prótesis con cierta cercanía en el tiempo a la amputación, porque si esperaba a la sentencia correría el riesgo de que la prótesis ya no pudiera ser colocada (por rechazo anatómico) y, además, no contar con ella le dificultaba reinsertarse en el mercado laboral. Peticionó una medida \textbf{innovativa} dentro del Art.~230 CPCC: que se le otorgara el dinero para acceder a la compra de la prótesis.

\textbf{Superposición de pretensiones.} El otorgamiento del dinero para la prótesis se superponía \emph{parcialmente} con la pretensión sustancial (la indemnización integral), lo que ubica el caso en el ámbito excepcional: tutela anticipada.

\textbf{Resolución.} El juez de primera instancia y la Sala J de la Cámara Civil rechazaron el pedido. Al llegar el expediente a la Corte por recurso extraordinario, la CSJN hizo lugar a la medida innovativa, reconociendo la tutela anticipada, al reunirse en el caso verosimilitud en el derecho, peligro en la demora y contracautela, y para evitar el daño irreparable que podía producirse si recién al dictarse la sentencia pudiera acceder a la prótesis.

\begin{important}
Lo único que agrega la Corte respecto de la tutela anticipada es que en estos casos se requiere \textbf{mayor prudencia} de los jueces al valorar los presupuestos y los hechos del caso concreto. Queda librado al juez, según su interpretación y valoración de los hechos, otorgar o no esta tutela judicial efectiva.
\end{important}

\subsubsection{Inversión del contradictorio}

Lo típico de estos sistemas es que se \textbf{invierte el contradictorio}: la jurisdicción primero protege y luego conoce. Se tutela primero para evitar el daño irreparable, y luego viene la etapa del contradictorio y la decisión definitiva, con certeza, al dictar sentencia. Si primero no se protegiera y se procediera a conocer, la sentencia podría ser estéril y el reconocimiento de derechos llegaría tarde.

\subsection{Clarín (CSJN 333:1885): faceta no innovativa}

\textbf{Contexto.} Hace aproximadamente diez años se dictó la Ley 26.522, conocida como ``Ley de Medios'', cuya finalidad era evitar la monopolización de los distintos medios de comunicación y telecomunicaciones. Las empresas que contaban con una gran cantidad de licencias debían desinvertir y vender parte de ellas. Esto afectaba de modo muy importante al Grupo Clarín, que debía iniciar, dentro del plazo legal, un proceso de desinversión.

\textbf{Proceso principal.} El Grupo Clarín planteó un proceso en el que discutió la constitucionalidad de los artículos 41 y 161 de la ley, que eran los que lo obligaban a desinvertir.

\textbf{Pretensión cautelar.} Solicitó una tutela anticipada dentro de la prohibición de innovar (Art.~230 CPCC), en su faceta \textbf{no innovativa}: que no se le aplicara la ley y no se alterara la situación existente hasta que se resolviera la pretensión sustancial sobre la constitucionalidad de los artículos. Si no se concedía y la sentencia declaraba la inconstitucionalidad, llegaría tarde: el Grupo ya habría tenido que desinvertir.

\textbf{Resolución y enseñanza.} La CSJN hizo lugar a la prohibición de innovar en su faceta no innovativa, y la ley no se aplicó mientras duró el proceso. Sin embargo, en la sentencia definitiva la Corte resolvió que los artículos 41 y 161 \textbf{eran constitucionales}, decisión contraria a lo protegido cautelarmente.
% TODO: contenido incompleto o ambiguo en las notas originales (la cita del docente sobre este fallo contiene una frase confusa: "el juez consideró autor necesario otorgar esta tutela").

El caso permite ver, por un lado, cómo opera la faceta de no innovar de la prohibición y, por otro, con claridad, el \textbf{carácter provisorio} de este tipo de medidas: se otorgó la tutela para evitar un daño muy difícil de reparar y, una vez producido el conocimiento a lo largo del proceso, se decidió que las normas eran constitucionales.

\section{Debate doctrinario sobre la regulación de la tutela anticipada}

\subsection{La posición crítica: regulación independiente en las provincias}

Una parte de la doctrina entiende que la tutela anticipada \textbf{no es un sistema cautelar} y que no resulta adecuado que cobre vida a través de la prohibición de innovar del Art.~230. Sostiene que debe regularse expresamente de manera distinta de las medidas cautelares, porque adelantar total o parcialmente, de modo provisorio y cautelar, lo que será objeto de la sentencia de mérito implicaría un instituto distinto.

Esta posición llevó a varias provincias ---como San Luis, San Juan y Chaco--- a regular la tutela anticipada de modo independiente en sus códigos procesales, con requisitos más exigentes. La crítica expresada es que ya no bastaría acreditar los presupuestos que exige la Corte (verosimilitud en el derecho, peligro en la demora, contracautela y mayor prudencia del juez): se establecen numerosos y rígidos requisitos adicionales que entorpecen el acceso del justiciable a esta tutela.

\subsection{Artículo 403 del Proyecto de reforma del CPCC: tutela anticipada de urgencia}

El proyecto de reforma del Código Procesal Civil y Comercial de la Nación regula en su artículo 403 la llamada ``tutela anticipada de urgencia'', con requisitos notablemente más exigentes:

\medskip
\renewcommand{\arraystretch}{1.3}
\noindent\begin{tabularx}{\textwidth}{@{}>{\raggedright\arraybackslash}p{2.8cm} Y Y@{}}
\toprule
\textbf{Requisito} & \textbf{Régimen actual (CSJN, Camacho Acosta)} & \textbf{Proyecto, Art.~403} \\
\midrule
Presupuesto temporal &
Peligro en la demora. &
Urgencia (estándar más elevado). \\
\addlinespace
Tipo de perjuicio &
Perjuicio inminente o irreparable. &
Frustración del derecho y daño irreparable. \\
\addlinespace
Prueba del derecho &
Verosimilitud (probabilidad). &
Debe ofrecer pruebas para acreditar la probabilidad del derecho. \\
\addlinespace
Contracautela &
Sí. &
Sí. \\
\addlinespace
Materias alcanzadas &
Cualquier tipo de derecho. &
Sólo prestaciones dinerarias. \\
\addlinespace
Trámite &
Inaudita parte: el afectado se defiende después. &
Con audiencia previa: primero se cita al demandado y luego se resuelve. \\
\bottomrule
\end{tabularx}
\renewcommand{\arraystretch}{1}

\medskip
Se critica la limitación a prestaciones dinerarias como incomprensible: se pierde una tutela judicial efectiva para situaciones vinculadas con otros derechos fundamentales y no sólo con cuestiones netamente dinerarias.

También se critica el trámite con audiencia previa por la demora que puede implicar. Si el demandado aún no se presentó en el proceso, debe notificárselo en su domicilio real (por ejemplo, Bahía Blanca, cuando el proceso tramita en la Ciudad de Buenos Aires); a veces es necesario diligenciar varias cédulas a través de la oficina de notificaciones, lo que puede demorar meses, para otorgar una tutela que en general requiere cierta urgencia.

\subsection{Principios vulnerados según la crítica}

Si la decisión es provisoria, si el afectado puede luego defenderse, si existe una contracautela como garantía por los daños y perjuicios, y si el instituto funciona del modo visto como medida cautelar dentro de la doctrina de la Corte, se cuestiona por qué dictar regulaciones que atentan contra principios reconocidos en tratados de derechos humanos con jerarquía constitucional:

\begin{itemize}
    \item \textbf{Tutela judicial efectiva:} consagrada en los tratados internacionales de derechos humanos con jerarquía constitucional.
    \item \textbf{Principio pro homine:} debe estarse siempre a la interpretación más favorable a la persona.
    \item \textbf{Principio de progresividad:} una vez que los justiciables adquieren cierto nivel de protección, el legislador no puede regular de modo regresivo en su contra.
\end{itemize}

Según esta posición, los sistemas que regulan la tutela anticipada de urgencia de este modo podrían ser tachados de \textbf{inconstitucionales y anticonvencionales}, porque el legislador crea un sistema apartándose de los principios de tutela judicial efectiva y de progresividad.

\cornelltitle{ámbito excepcional y tutela anticipada}
\cornellblock{
\cqrow{¿Qué define al ámbito excepcional de los sistemas cautelares?}{La superposición total o parcial de la pretensión cautelar con la pretensión sustancial (el objeto del proceso principal), a diferencia del ámbito normal, donde ambas son distintas.
}
\cqrow{¿Qué es la tutela anticipada?}{Es el sistema cautelar por el cual la jurisdicción protege y asegura anticipadamente y de modo provisorio, en todo o en parte, lo que luego se decidirá definitivamente en la sentencia de mérito.
}
\cqrow{¿Por qué no hay prejuzgamiento si se adelanta lo que decidirá la sentencia?}{Porque lo decidido cautelarmente es provisorio: se otorga con verosimilitud en el derecho y peligro en la demora, y luego, tras el debate y la prueba, la sentencia puede resolver en el mismo sentido con carácter definitivo o en sentido contrario, dejando sin efecto lo concedido.
}
\cqrow{¿Qué diferencia hay entre cautela y condena?}{Si el juez otorga la tutela anticipada de modo provisorio hay cautela; si decide en la sentencia definitiva sobre esa misma pretensión hay condena, que es definitiva y con autoridad de cosa juzgada.
}
\cqrow{¿Cuáles son las vías de materialización de la tutela anticipada?}{(1) Los sistemas cautelares: prohibición de innovar (Art.~230), medida cautelar genérica (Art.~232) u otra medida cautelar. (2) Subsistemas creados por el legislador: alimentos provisorios (Art.~544 CCyC) y entrega anticipada del inmueble en el desalojo contra el intruso (Art.~680 bis CPCC).
}
\cqrow{¿Qué resolvió la Corte en Camacho Acosta (CSJN 320:1633)?}{Hizo lugar a una medida innovativa (Art.~230 CPCC) como tutela anticipada: se anticipó el dinero para la prótesis del trabajador al que se le amputó el brazo, por verosimilitud, peligro en la demora y contracautela. Agregó que estos casos exigen mayor prudencia del juez al valorar presupuestos y hechos.
}
\cqrow{¿Qué resolvió la Corte en Clarín (CSJN 333:1885) y qué demuestra?}{Hizo lugar a la prohibición de no innovar: no se aplicó la Ley 26.522 al Grupo Clarín mientras duró el proceso. En la sentencia definitiva declaró constitucionales los artículos 41 y 161, lo que muestra el carácter provisorio de la tutela anticipada y la faceta no innovativa de la prohibición de innovar.
}
}
\medskip
\cornellblocksum{
\cqrow{¿Qué significa la inversión del contradictorio?}{En la tutela anticipada la jurisdicción primero protege y luego conoce, para evitar un daño irreparable o una sentencia estéril por tardía; el contradictorio y la decisión definitiva vienen después.
}
\cqrow{¿Qué críticas se formulan a la regulación independiente de la tutela anticipada?}{Que las provincias que la regulan y el Art.~403 del proyecto agregan requisitos numerosos y rígidos (urgencia, pruebas, sólo prestaciones dinerarias, audiencia previa) que entorpecen el acceso a la tutela y pueden demorar meses. Se afirma que vulneran la tutela judicial efectiva, el principio pro homine y la progresividad, y podrían ser inconstitucionales y anticonvencionales.
}
}{El ámbito excepcional se configura cuando la pretensión cautelar se superpone total o parcialmente con la sustancial; esto es la tutela anticipada, válida porque es provisoria y se otorga con verosimilitud y peligro en la demora, con contracautela y mayor prudencia judicial. Se materializa por vías cautelares (prohibición de innovar, medida genérica) y por subsistemas legales (alimentos provisorios, desalojo anticipado). La Corte aplicó la faceta innovativa en Camacho Acosta y la no innovativa en Clarín, donde la sentencia final fue contraria a lo cautelarmente protegido. La regulación independiente con requisitos agravados es criticada por contrariar la tutela judicial efectiva, el principio pro homine y la progresividad.
}

\chapter{El embargo}\label{cap:embargo}

\section{El embargo como medida cautelar por antonomasia}

\begin{definition}{Embargo}
La doctrina ha dicho generalmente que el embargo es la \textbf{medida cautelar por antonomasia}: es la más usada y la más habitual, y consiste en la \textbf{afectación e individualización de un bien} ---mueble o inmueble, registrable o no--- al resultado de un proceso.
\end{definition}

El bien así individualizado queda vinculado directamente con el proceso, queda a disposición del juez y opera como \textbf{prenda concreta} para garantizar la efectividad del resultado del juicio.

Ese bien se encuentra afectado ante la eventualidad de que se concrete la realización de bienes, es decir, la necesidad de sacar a la venta bienes para pagar la deuda reclamada en el proceso. Será el bien fundamental a la hora de procederse a la subasta pública como modo de garantizar el pago de la deuda.

\section{Tipos de embargo}

El ordenamiento reconoce tres tipos de embargo, que se distinguen según el momento procesal en que se traban:

\begin{table}[H]
\centering
\small
\begin{tabularx}{\textwidth}{>{\raggedright\arraybackslash}p{2.6cm} >{\raggedright\arraybackslash}p{3.6cm} >{\raggedright\arraybackslash}X}
\toprule
\textbf{Tipo} & \textbf{Momento} & \textbf{Normativa / características} \\
\midrule
\textbf{Embargo preventivo} & Antes del proceso o durante su tramitación & Arts.\ 209 a 212 CPCCN. Eminentemente cautelar. Exige los tres recaudos: verosimilitud del derecho, peligro en la demora y contracautela suficiente. \\
\addlinespace
\textbf{Embargo ejecutivo} & Una vez admitido el título en el juicio ejecutivo & Se traba cuando el juez reconoce el título como admisible y libra la orden de intimación de pago. Individualiza el bien que será ejecutado si no se paga. \\
\addlinespace
\textbf{Embargo ejecutorio} & En el proceso de ejecución de sentencia & Se dicta cuando la sentencia está confirmada y se avanza en la concreción del pago. Es el paso previo a la subasta del bien. Art.\ 551 CPCCN (sentencia de trance y remate). \\
\bottomrule
\end{tabularx}
\caption{Tipos de embargo}
\end{table}

\subsection{El embargo preventivo}

El embargo preventivo es el eminentemente cautelar. Puede trabarse aun antes del proceso o durante la secuela del proceso, siempre que se reúnan los tres recaudos desarrollados en el capítulo~\ref{cap:presupuestos}: \textbf{verosimilitud del derecho}, \textbf{peligro en la demora} y \textbf{contracautela} adecuada o suficiente.

Los artículos de referencia son los siguientes:
\begin{itemize}
    \item \textbf{Art.\ 209 CPCCN:} establece los supuestos de procedencia del embargo preventivo.
    \item \textbf{Art.\ 210 CPCCN:} amplía los supuestos y complementa el artículo anterior.
    \item \textbf{Art.\ 211 CPCCN:} procedimiento para la traba del embargo.
    \item \textbf{Art.\ 212 CPCCN:} hipótesis específicas derivadas de la secuela del proceso.
\end{itemize}

\subsection{El embargo ejecutivo}

En el trámite del proceso de ejecución existe una primera etapa, que es el análisis del título. No siempre el título trae aparejada, por sí mismo, la ejecución: a veces es necesario preparar la vía ejecutiva. En esos casos procede el embargo preventivo, porque todavía no hay un título ejecutivo declarado admisible.

Una vez que el juez advierte que se trata de un título \textit{prima facie} hábil ---ya sea porque se preparó la vía ejecutiva o porque objetivamente reúne todos los recaudos propios del título---, el título se considera admisible y se libra la orden de intimación de pago. A partir de ese momento procede el \textbf{embargo ejecutivo}.

Por lo tanto, en un juicio ejecutivo puede haber un embargo preventivo trabado antes del inicio del proceso o con el inicio del proceso, si este no comienza directamente con la intimación de pago porque es necesario preparar la vía ejecutiva. Una vez preparada la vía ejecutiva y admitido el título, el embargo pasa a ser ejecutivo, ya sea porque el preventivo se transforma en ejecutivo, ya sea porque se traba directamente como embargo ejecutivo.

\subsection{El embargo ejecutorio}

El embargo ejecutorio es el que se traba en el marco del proceso de ejecución de sentencia. Se dicta cuando una sentencia ya ha sido confirmada por el superior y se empieza a avanzar en la concreción del pago derivado de una sentencia condenatoria. Con sentencia firme y ejecutoriada se admite la traba de un embargo que ya no es preventivo ---porque el derecho está consolidado y reconocido--- ni ejecutivo ---porque no se traba en el trámite de ejecución de un título ejecutivo---, sino \textbf{ejecutorio}, propio del ámbito de la ejecución de sentencia.

\begin{important}
\textbf{Art.\ 551 CPCCN (sentencia de trance y remate).} A partir de su dictado, los embargos anteriores se transforman en ejecutorios por efecto de la ley. La transformación opera automáticamente, muchas veces sin disposición concreta del tribunal, por la etapa del proceso en que se encuentra la causa. El embargo ejecutorio es el paso previo a la subasta del bien.
\end{important}

\section{Efectos de la traba del embargo}

El embargo es una medida procesal que se dicta en el proceso, pero que tiene efectos regulados en los códigos sustanciales. Trae consecuencias civiles y penales, que se ponen en juego en caso de incumplimiento de las obligaciones del depositario de la cosa embargada.

\begin{important}
\textbf{Consecuencias extralegales del embargo.} El depositario del bien embargado tiene la obligación de conservar la cosa, no disminuir su valor y ponerla a disposición del tribunal cuando sea requerido. El Código Penal tipifica delitos vinculados al incumplimiento de las obligaciones del depositario.
\end{important}

Los efectos fundamentales son los siguientes:
\begin{itemize}
    \item El embargo individualiza el bien, que queda afectado a la garantía y vinculado al proceso.
    \item El embargo \textbf{no produce la desapropiación} del bien por parte del afectado. El afectado puede usar y gozar de la cosa, en tanto no la haga perder valor.
    \item Tratándose de bienes registrables, el embargo \textbf{no produce la imposibilidad de enajenación} del bien: el bien puede ser vendido.
\end{itemize}

\subsection{El principio \textit{nemo plus iuris} y el embargo}

Por el principio \textit{nemo plus iuris} ---nadie puede transmitir sobre una cosa un derecho mayor al que posee sobre ella---, quien posee un bien embargado lo posee en tanto se encuentra embargado. Si lo transfiere, lo transmite embargado, y quien lo adquiere adquiere sobre la cosa el mismo derecho que tenía su propietario.

\subsection{El plenario «Sertop» --- Doctrina judicial obligatoria}

\begin{important}
\textbf{Plenario «Sertop»} (Cámara Nacional de Apelaciones en lo Civil). Constituye doctrina judicial obligatoria en la Ciudad Autónoma de Buenos Aires.

\textbf{Regla:} el adquirente de un bien embargado asume la totalidad de la deuda ---no solo el monto por el cual se trabó la medida---, incluyendo ampliaciones, intereses y costas del proceso.

En el fuero \textbf{comercial} este criterio no rige: la transmisión se limita al monto de la cautela.
\end{important}

Ante la duda que genera el criterio que se empleará en este tipo de transmisiones, difícilmente alguien esté dispuesto a adquirir una cosa embargada. Por ello, el embargo suele operar como una suerte de traba a la libre circulación de los bienes. Generalmente quienes adquieren bienes embargados especulan con ello, pagando un precio significativamente menor al de mercado, o tienen algún tipo de vinculación con el sujeto embargado.

\section{Supuestos de procedencia del embargo preventivo (arts.\ 209 y 210 CPCCN)}

La enumeración legal es meramente enunciativa y orientativa, no taxativa: las situaciones que habilitan el embargo pueden ser muchas más, y no es necesario que deban encuadrarse en alguno de los presupuestos que menciona el código.

Los supuestos típicos del art.\ 209 CPCCN son:
\begin{enumerate}[label=\textbf{\arabic*.}]
    \item \textbf{Deudor sin domicilio en la República} (inciso 1°). Se presume el peligro en la demora, porque puede ser dificultoso encontrarle bienes para responder.
    \item \textbf{Reconocimiento de la deuda a través de instrumento público.} Se presume reunida la verosimilitud y el peligro en la demora.
    \item \textbf{Instrumento privado cuyas firmas son abonadas por dos testigos} (información sumaria). El aval de dos testigos que declaran haber presenciado la firma o conocer la rúbrica del deudor permite dictar la medida sin mayores miramientos.
    \item \textbf{Contratos bilaterales} instrumentados públicamente o en instrumento privado con dos testigos (información sumaria), que avalan que la prestación a cargo del acreedor fue cumplida y resta cumplir la contraprestación.
\end{enumerate}

\subsection{Hipótesis procesales del artículo 212 CPCCN}

El art.\ 212 describe tres hipótesis en las que la secuela del proceso determina la procedencia del embargo preventivo:
\begin{itemize}
    \item \textbf{Inciso 1° (remisión al art.\ 63):} declaración de rebeldía del demandado. El hecho de que la parte sea declarada rebelde da motivo formal para la traba del embargo preventivo.
    \item \textbf{Inciso 2°:} confesión de la existencia de la deuda en audiencia confesional, o declaración de confeso ficto (por no comparecer a la audiencia habiendo sido citado). En ambos casos procede el embargo sin mayores recaudos.
    \item \textbf{Inciso 3°:} existencia de sentencia definitiva condenatoria, aunque esté apelada. Si el juez de primera instancia condenó al pago, el embargo preventivo está plenamente configurado; la eventual revisión por la Cámara no impide la medida.
\end{itemize}

\begin{note}
Cuando el demandado no contesta la demanda, por efecto del art.\ 356 se presume la veracidad de los hechos descriptos, lo que también habilita la traba. % TODO: contenido incompleto o ambiguo en las notas originales (las notas aluden a «hechos ilícitos descriptos»)
\end{note}

La existencia de una sentencia judicial condenatoria, aunque se encuentre apelada, otorga una mayor verosimilitud del derecho. La Cámara puede revisar y revocar la sentencia; sin embargo, ya existe un juez de primera instancia que consideró que el deudor debe pagar la suma reclamada, y es lógico que en ese contexto se dicte el embargo preventivo como modo de garantizar el cumplimiento de la sentencia ya dictada.

\section{Bienes inembargables --- Artículo 219 CPCCN}

\begin{enumerate}[label=\textbf{Inciso \arabic*°:}, leftmargin=2.4cm]
    \item El lecho cotidiano del deudor, su esposa e hijos; las ropas y bienes muebles de indispensable uso. Según el criterio judicial, si el deudor tiene cinco televisores puede reconocerse la inembargabilidad de uno, y los demás quedan embargables. También son inembargables las herramientas u objetos que el deudor utiliza para el ejercicio de su profesión u oficio.
    \item Los sepulcros, salvo que la deuda sea por el saldo del precio de compra del sepulcro o por bienes adquiridos para su construcción.
    \item Los bienes declarados inembargables por leyes especiales (jubilaciones, pensiones, salarios, etc.).
\end{enumerate}

\subsection{Inembargabilidad del salario}

El trabajador tiene en su salario un porcentaje de inembargabilidad establecido por la Ley de Contrato de Trabajo, que establece que solo puede ser embargado hasta el \textbf{20\,\%} de sus ingresos.

\begin{example}{Límite del 20\,\% y deuda alimentaria}
Si el salario del deudor es de \$30.000, el 20\,\% equivale a \$6.000. Sin embargo, si el deudor tiene tres hijos menores y la madre reclama alimentos, \$6.000 puede resultar insuficiente. En ese supuesto, el juez puede excepcionalmente elevar el límite de embargabilidad al 40\,\% o 50\,\% del ingreso, dejando de lado el límite de la Ley de Contrato de Trabajo, siempre como modo de garantizar la subsistencia del trabajador.
\end{example}

Un criterio similar rige para jubilaciones y pensiones, que en principio son inembargables, salvo en caso de deudas alimentarias, donde el juez puede establecer el porcentaje embargable garantizando la subsistencia del jubilado o pensionado.

\section{Caducidad del embargo --- Artículo 207 CPCCN}

En materia de embargo rige, como en general en todas las medidas cautelares, lo establecido por el art.\ 207 del código procesal sobre la caducidad de la traba.

\begin{formula}{Caducidad (art.\ 207 CPCCN)}
\begin{itemize}
    \item \textbf{Plazo:} 10 días desde la traba de la medida para iniciar el proceso principal.
    \item En procesos sujetos a mediación previa, el inicio de la mediación interrumpe el plazo de caducidad (equivale al inicio del proceso).
    \item \textbf{Embargos registrales} (Registro del Automotor): caducan a los \textbf{cinco años}. Debe pedirse la reinscripción dentro de ese plazo para mantener subsistente la medida.
\end{itemize}
\end{formula}

\cornelltitle{el embargo}
\cornellblock{
\cqrow{¿Qué es el embargo y por qué se lo llama «medida cautelar por antonomasia»?}{Es la medida cautelar más usada y habitual. Consiste en la afectación e individualización de un bien ---mueble o inmueble, registrable o no--- al resultado de un proceso. El bien queda vinculado al proceso, a disposición del juez, y opera como prenda concreta para garantizar el resultado del juicio y, llegado el caso, la subasta pública.
}
\cqrow{¿Cuáles son los tres tipos de embargo?}{
}
\cqrow{¿Qué efectos produce la traba del embargo?}{Individualiza el bien y lo afecta a la garantía. No desapropia al afectado, que puede usar y gozar de la cosa mientras no la haga perder valor. En bienes registrables no impide la enajenación. El depositario debe conservar la cosa, no disminuir su valor y ponerla a disposición del tribunal; su incumplimiento tiene consecuencias civiles y penales.
}
\cqrow{¿Puede venderse un bien embargado? ¿Cuáles son las consecuencias?}{Sí, pero por \textit{nemo plus iuris} se transmite embargado y el adquirente obtiene el mismo derecho que tenía el propietario. Según el plenario «Sertop» (doctrina obligatoria en la Ciudad Autónoma de Buenos Aires), el adquirente asume la totalidad de la deuda, incluyendo ampliaciones, intereses y costas; en el fuero comercial la transmisión se limita al monto de la cautela. Por ello difícilmente alguien adquiere bienes embargados y el embargo opera como traba a la libre circulación.
}
\cqrow{¿Cuáles son los supuestos de procedencia del embargo preventivo?}{La enumeración de los arts.\ 209 y 210 es enunciativa, no taxativa. Art.\ 209: deudor sin domicilio en la República; reconocimiento de deuda por instrumento público; instrumento privado con firmas abonadas por dos testigos; contratos bilaterales con prestación del acreedor cumplida. Art.\ 212: rebeldía, confesión o confeso ficto, y sentencia condenatoria aunque esté apelada.
}
\cqrow{¿Cuáles son los principales bienes inembargables?}{Art.\ 219 CPCCN: el lecho cotidiano y las ropas y bienes muebles de indispensable uso; los sepulcros (con la salvedad indicada); y los bienes declarados inembargables por leyes especiales (jubilaciones, pensiones, salarios). El salario es embargable hasta el 20\,\%, límite que el juez puede elevar excepcionalmente en deudas alimentarias.
}
}
\medskip
\cornellblocksum{
\cqrow{¿Cuándo caduca el embargo?}{Art.\ 207 CPCCN: a los 10 días desde la traba si no se inicia el proceso principal (la mediación previa interrumpe el plazo). Los embargos registrales sobre automotores caducan a los cinco años, salvo reinscripción dentro de ese plazo. \\




%==============================================================
}
}{El embargo individualiza un bien y lo afecta al proceso sin desapropiar al deudor. Se clasifica según el estadio procesal en preventivo, ejecutivo y ejecutorio. El adquirente de un bien embargado lo recibe con la afectación (\textit{nemo plus iuris}; plenario «Sertop»). Rigen límites de embargabilidad (art.\ 219, límite del 20\,\% del salario) y la caducidad del art.\ 207.
}

\chapter[Secuestro e intervención judicial]{Secuestro e intervención judicial}\label{cap:secuestro}

\section{El secuestro}
%==============================================================

\subsection{Concepto y ámbito de aplicación}

El secuestro está orientado, como medida cautelar, a desapoderar al deudor de la disposición material y jurídica sobre una cosa de su propiedad.

\begin{important}
El secuestro solo opera respecto de bienes \textbf{muebles}. No existe secuestro de bienes inmuebles. Puede recaer sobre bienes muebles registrables (autos, motos) y no registrables. Los animales (ganado, caballos de carrera) también pueden ser objeto de secuestro.
\end{important}

\subsection{Tipos de secuestro}

\subsubsection{Secuestro como medida cautelar propiamente dicha (art.\ 221 CPCCN)}

Es excepcional y no tan habitual como el embargo. De acuerdo con las características del bien, las del deudor o la situación en que se encuentra el bien, en algunas hipótesis procede el secuestro como único modo de garantizar la intangibilidad de la cosa.

Procede cuando la cosa puede ser deteriorada por el uso continuo del deudor, o cuando el deudor ha dado muestras de falta de cuidado y ha permitido que el bien pierda su valor.

\begin{example}{Procedencia del secuestro como cautelar}
\textbf{Ejemplo 1.} La cosa es una obra de arte almacenada en un depósito con humedad, que no recibe los cuidados indispensables para su preservación. Es preferible el secuestro al embargo para evitar el deterioro del bien.

\textbf{Ejemplo 2.} Se trata de aparatología para realizar estudios técnicos o científicos que requiere un mantenimiento que el deudor no realiza, generando un deterioro progresivo.
\end{example}

\subsubsection{Secuestro como sanción}

La doctrina reconoce que el secuestro a veces es una sanción derivada de la conducta del deudor o de quien incumple una obligación procesal. El código lo prevé en dos hipótesis:

\textbf{Art.\ 128 CPCCN --- Secuestro por retención de expediente.} Cuando un letrado o un auxiliar del tribunal (escribano, perito) retira un expediente en préstamo y no lo devuelve en el plazo fijado por el juez:
\begin{enumerate}
    \item Se intima a devolverlo.
    \item Si no lo devuelve en el plazo de la intimación, se dispone el secuestro del expediente por oficial de justicia, con allanamiento de morada si fuera necesario.
\end{enumerate}
Es una hipótesis de secuestro como sanción por incumplimiento de la obligación de restituir.

\textbf{Art.\ 323 CPCCN, incisos 1° a 5° --- Diligencias preliminares.} En las diligencias preliminares puede pedirse a la parte demandada que exhiba un objeto, un inmueble, un testamento o los títulos de una cosa vendida. Si se incumple la intimación de exhibición, una de las sanciones previstas es el secuestro de la cosa, título o documento de que se trate.

\subsubsection{Secuestro como paso previo a la subasta}

Una vez que el bien es embargado en el marco de un juicio de ejecución y es necesario proceder a su venta en pública subasta, porque ya se agotaron los pasos procesales previos, el secuestro actúa como presupuesto de la subasta pública de los bienes embargados.

Si el bien es un inmueble, no corresponde el secuestro, porque este no procede respecto de todo lo que sea inmueble. Los bienes muebles, sean registrables o no, tienen como paso previo a la subasta el secuestro.

\cornelltitle{el secuestro}
\cornellblocksum{
\cqrow{¿Qué es el secuestro y sobre qué bienes recae?}{Es una medida cautelar orientada a desapoderar al deudor de la disposición material y jurídica sobre una cosa de su propiedad. Solo recae sobre bienes muebles, registrables o no, incluidos los animales. No existe secuestro de inmuebles.
}
\cqrow{¿Cuándo procede el secuestro como medida cautelar (art.\ 221 CPCCN)?}{Excepcionalmente, cuando por las características del bien, del deudor o de la situación del bien es el único modo de garantizar la intangibilidad de la cosa: si puede deteriorarse por el uso continuo del deudor o si el deudor ha mostrado falta de cuidado y permitió que pierda valor (por ejemplo, una obra de arte en un depósito con humedad).
}
\cqrow{¿En qué hipótesis el secuestro opera como sanción?}{En dos: art.\ 128 CPCCN, cuando un letrado o auxiliar del tribunal no devuelve un expediente tras la intimación (secuestro por oficial de justicia, con allanamiento si fuera necesario); y art.\ 323 CPCCN, cuando se incumple la intimación a exhibir en diligencias preliminares.
}
\cqrow{¿Qué función cumple el secuestro respecto de la subasta?}{Es presupuesto de la subasta pública de los bienes muebles embargados en un juicio de ejecución. Los inmuebles no se secuestran. \\




%==============================================================
}
}{El secuestro desapodera al deudor de la disposición material y jurídica de bienes muebles. Es excepcional como cautelar (art.\ 221), opera como sanción (arts.\ 128 y 323) y constituye paso previo a la subasta de los muebles embargados. No procede sobre inmuebles.
}

\section{La intervención judicial}
%==============================================================

\subsection{Concepto y fundamento (arts.\ 222 a 227 CPCCN)}

\begin{definition}{Intervención judicial}
Consiste en la designación, por parte del juez, de un \textbf{auxiliar externo al tribunal} para que se involucre en la administración o gestión de un ente colectivo, un patrimonio o una empresa. El interventor \textbf{no tiene facultades de gobierno ni de disposición de bienes}. Su función es garantizar una saludable administración o brindar información al juez.
\end{definition}

\subsection{Tipos de intervención judicial}

\subsubsection{Interventor recaudador (art.\ 223 CPCCN)}

El interventor recaudador es un auxiliar del tribunal designado por el juez, cuya misión es recaudar fondos del giro habitual de la administración de una empresa, un comercio o un patrimonio.

Suele darse cuando la parte deudora no tiene un patrimonio significativo pero tiene un comercio o consultorio. El interventor concurre cotidianamente al establecimiento y, del dinero que ingresa, retiene una suma para afectarla al pago de la deuda reclamada en el proceso.

\begin{example}{Empresa de transporte de ómnibus}
Una empresa de ómnibus de larga distancia que cobra sus ingresos por ventanilla (venta de pasajes) tiene una deuda de \$20.000. Embargar un ómnibus y sacarlo a subasta resulta un despropósito, porque el gasto de la subasta supera la deuda.

\textbf{Solución:} se designa un interventor a ventanilla que retiene del ingreso cotidiano la suma correspondiente, logrando así el cumplimiento de la deuda de forma proporcional y eficiente.

En el caso de empresas de colectivos con tarjeta SUBE, se puede directamente comunicar por oficio a la administración de la gestión de las tarjetas para que retengan lo cobrado.
\end{example}

\subsubsection{Interventor informante (art.\ 224 CPCCN)}

Está orientado a que el interventor que concurre al comercio de la empresa de la parte afectada dé noticia al juez acerca de la manera en que se lleva a cabo la administración, el patrimonio que posee la empresa, el estado de deudas y todo otro dato que el juez necesite. Aporta así información valiosa al proceso sobre la gestión de un comercio, empresa o patrimonio, que puede ser necesaria para decisiones que deban tomarse en el proceso.

\begin{example}{Sucesión con empresa}
En un juicio sucesorio, el patrimonio del causante incluye una empresa o comercio del que era titular. Para distribuir la herencia entre los herederos es necesario contar con información precisa sobre ingresos, egresos, bienes y deudas de ese patrimonio. En esos casos la intervención del informante es necesaria para «poner blanco sobre negro» acerca del alcance de los bienes a distribuir y los ingresos que genera el patrimonio.
\end{example}

\subsubsection{Veedor (art.\ 115 Ley 19.550 --- Ley de Sociedades Comerciales)}

En doctrina se discute qué es el veedor a diferencia del informante. Prácticamente es lo mismo: la Ley de Sociedades denomina veedor al informante, y la diferencia tiene que ver más con una cuestión terminológica que con el alcance de la gestión.

En ambos casos ---interventor informante del art.\ 224 CPCCN o veedor del art.\ 115 de la Ley 19.550--- el alcance de su gestión y la información que debe brindar está establecido con mayor precisión por lo que dispone el juez en la resolución que ordena la intervención que por lo preestablecido en las normas del código.

\subsubsection{Interventor administrador o co-administrador}

El interventor gestiona la administración ---en casos puntuales que el juez comercial establece--- o la cogestiona con administradores propios de la empresa. Se da un paso más: ya no es un informante, sino que realiza cuestiones vinculadas con la administración cotidiana.

Suele suceder cuando la empresa está en estado prefalencial y es necesario evitar que, por maniobras de administración, se propicie la insolvencia a través del retiro de dinero por parte de los socios u otras maniobras perjudiciales.

\begin{important}
\textbf{Límites del interventor.} La intervención no permite la disposición de bienes ni el gobierno de la empresa. Solo en el caso del administrador o co-administrador está habilitada la administración. Los actos de disposición o gobierno deben ser siempre adoptados por el juez, con intervención del síndico, cumpliendo los pasos procesales del proceso falencial. El administrador nunca tiene facultades para actos de disposición sin disposición expresa del juez.
\end{important}

\subsubsection{Quiénes pueden ser designados interventores}

Se aconseja que sean personas idóneas o técnicas en el negocio de que se trate. Generalmente se los busca entre administradores de empresa, contadores o incluso abogados para ciertas gestiones, como la de recaudación en un comercio.

Los abogados, contadores y licenciados en administración de empresas que pretenden ofrecer sus servicios para este tipo de intervenciones se anotan en los listados que elabora anualmente la Cámara Civil y la Cámara Comercial, presentando su currículum y poniéndose a disposición de los tribunales.

\cornelltitle{la intervención judicial}
\cornellblocksum{
\cqrow{¿Qué es la intervención judicial?}{Es la designación por el juez de un auxiliar externo al tribunal para que se involucre en la administración o gestión de un ente colectivo, un patrimonio o una empresa, a fin de garantizar una saludable administración o brindar información al juez.
}
\cqrow{¿Qué facultades tiene el interventor judicial?}{No tiene facultades de gobierno ni de disposición de bienes. Solo el administrador o co-administrador está habilitado para la administración. Los actos de disposición o gobierno los adopta siempre el juez, con intervención del síndico en el proceso falencial.
}
\cqrow{¿Qué hace el interventor recaudador (art.\ 223)?}{Recauda fondos del giro habitual de una empresa, comercio o patrimonio, reteniendo del dinero que ingresa una suma para afectarla al pago de la deuda. Es útil cuando el deudor no tiene patrimonio significativo pero sí un comercio o consultorio (por ejemplo, el interventor a ventanilla en una empresa de ómnibus).
}
\cqrow{¿Qué diferencia hay entre el interventor informante y el veedor?}{Prácticamente ninguna: el informante del art.\ 224 CPCCN y el veedor del art.\ 115 de la Ley 19.550 difieren más por la terminología que por el alcance de su gestión, que está determinado sobre todo por la resolución judicial que ordena la intervención.
}
\cqrow{¿Cuándo actúa un interventor administrador o co-administrador?}{En casos puntuales que establece el juez comercial, cuando la empresa está en estado prefalencial y es necesario evitar maniobras de administración que propicien la insolvencia. Realiza tareas de administración cotidiana, sola o junto con los administradores de la empresa.
}
\cqrow{¿Quiénes pueden ser interventores?}{Personas idóneas o técnicas en el negocio (administradores de empresa, contadores, abogados), inscriptas en los listados que elaboran anualmente la Cámara Civil y la Cámara Comercial. \\




%==============================================================
}
}{La intervención judicial coloca a un auxiliar del tribunal en la gestión de una empresa o patrimonio, sin facultades de gobierno ni de disposición. Sus variantes son el recaudador (art.\ 223), el informante o veedor (art.\ 224 CPCCN; art.\ 115 Ley 19.550) y el administrador o co-administrador; los actos de disposición corresponden siempre al juez.
}

\chapter[Anotación de litis y prohibición de contratar]{Anotación de litis y prohibición de contratar}\label{cap:anotacion}

\section{La anotación de litis}
%==============================================================

\subsection{Concepto y fundamento}

Este tema se conecta con el principio del tercero adquirente de buena fe a título oneroso, estudiado en Derecho Civil I. Entre los efectos de la nulidad de actos jurídicos se encuentra volver las cosas hacia atrás: ante la nulidad, todo vuelve al \textit{statu quo ante}. Si el acto jurídico era la compraventa o la donación de un inmueble y es declarado nulo, el inmueble vuelve al patrimonio de quien lo enajenó nulamente.

\begin{important}
\textbf{Art.\ 392 CCyCN --- Tercero adquirente de buena fe a título oneroso.} Cuando existe un tercero adquirente de buena fe a título oneroso, la ley privilegia su situación: este tercero, que no fue parte en el acto declarado nulo y no tenía nada que ver con él, no tiene por qué cargar con las consecuencias de las operaciones viciadas realizadas anteriormente respecto del bien. Por lo tanto, los efectos de la nulidad no se extienden a su respecto.
\end{important}

\begin{definition}{Anotación de litis}
Es una medida cautelar que, adoptada siempre sobre bienes \textbf{registrables}, publicita \textit{erga omnes} (con alcance a todos) la existencia de un proceso en trámite en el que se discute la validez o el alcance de una registración existente respecto de ese bien.

\textbf{Consecuencia:} al ser pública la existencia del juicio, nadie puede invocar luego la condición de adquirente de buena fe, porque nadie puede alegar buena fe cuando sabe que hay un litigio en trámite respecto del bien.
\end{definition}

Quien adquiere el bien en esos términos sabe que adquiere un bien en riesgo de un litigio que puede resolverse en favor de uno u otro y que, según cómo salga el juicio, tal vez todo vuelva atrás y el bien deba devolverse a un anterior propietario que lo enajenó en virtud de un acto viciado.

\subsection{Requisitos de procedencia}

\begin{itemize}
    \item Solo procede contra bienes \textbf{registrables}.
    \item Debe haber una \textbf{demanda iniciada}. A diferencia del embargo, que puede trabarse antes del proceso, la anotación de litis exige en principio que exista juicio ya iniciado (se admite que el inicio de la mediación equivale al inicio del proceso).
    \item El juicio debe tratarse de una pretensión que, de prosperar, implicaría la modificación de un asiento registral respecto del bien (ya sea de titularidad u otra cuestión vinculada a lo que se anota registralmente).
\end{itemize}

\subsection{Supuestos típicos de procedencia}

\begin{table}[H]
\centering
\small
\begin{tabularx}{\textwidth}{>{\raggedright\arraybackslash}p{3.2cm} >{\raggedright\arraybackslash}p{3.0cm} >{\raggedright\arraybackslash}X}
\toprule
\textbf{Tipo de juicio} & \textbf{Normativa} & \textbf{Particularidad} \\
\midrule
\textbf{Usucapión / prescripción adquisitiva} & Art.\ 1905 CCyCN & El juez ordena la anotación de oficio (sin pedido de parte) al disponer el traslado de la demanda. \\
\addlinespace
\textbf{Expropiación} & Arts.\ 24 y 14 Ley 21.499 & Una de las primeras medidas al iniciarse el proceso es la anotación de litis, para impedir que el propietario venda el bien y luego un tercero invoque buena fe frente a la expropiación. \\
\addlinespace
\textbf{Escrituración} & CPCCN & Busca modificar el asiento registral para que el adquirente por boleto pase a ser titular registral. \\
\addlinespace
\textbf{Nulidad de escritura} & Art.\ 392 CCyCN & Supuesto clásico de vicios del consentimiento o actos irregulares que intentan cancelar un asiento registral. \\
\addlinespace
\textbf{Simulación} & CCyCN & Caso típico de procedencia de la anotación de litis. \\
\addlinespace
\textbf{Interdicto de adquirir} & Art.\ 609 CPCCN & Hipótesis propia de la anotación de litis. \\
\bottomrule
\end{tabularx}
\caption{Supuestos típicos de anotación de litis}
\end{table}

\subsection{Extinción de la anotación de litis}

La anotación de litis \textbf{no se rige por el art.\ 207 CPCCN}. Se extingue por:
\begin{enumerate}
    \item \textbf{Cumplimiento de la sentencia:} si la demanda prospera, el cumplimiento de la sentencia determina la modificación registral buscada y concluye la anotación.
    \item \textbf{Desestimación del reclamo:} cuando el proceso concluye por sentencia desestimatoria, se agota la anotación de litis. Debe comunicarse al registro.
    \item \textbf{Modos anormales de terminación:} caducidad de instancia, desistimiento o transacción. Cualquiera de ellos trae aparejada la extinción de la anotación.
\end{enumerate}

\cornelltitle{la anotación de litis}
\cornellblocksum{
\cqrow{¿Qué es la anotación de litis y qué protege?}{Es una medida cautelar sobre bienes registrables que publicita erga omnes la existencia de un proceso en el que se discute la validez o el alcance de una registración. Al ser pública, impide que alguien invoque luego buena fe como adquirente, ya que nadie puede alegar buena fe sabiendo que hay un litigio en trámite sobre el bien.
}
\cqrow{¿Qué relación tiene con el art.\ 392 CCyCN?}{El art.\ 392 privilegia al tercero adquirente de buena fe a título oneroso, a quien no se extienden los efectos de la nulidad. La anotación de litis neutraliza esa buena fe al hacer pública la existencia del juicio.
}
\cqrow{¿Cuáles son los requisitos de procedencia?}{Que se trate de un bien registrable; que haya demanda iniciada (el inicio de la mediación equivale al inicio del proceso); y que la pretensión, de prosperar, implique la modificación de un asiento registral.
}
\cqrow{¿En qué supuestos procede la anotación de litis de oficio?}{En la usucapión (art.\ 1905 CCyCN), donde el juez la ordena sin pedido de parte al disponer el traslado de la demanda. En la expropiación (arts.\ 24 y 14 Ley 21.499) es una de las primeras medidas al iniciarse el proceso.
}
\cqrow{¿Qué otros juicios son supuestos típicos?}{Escrituración, nulidad de escritura (art.\ 392 CCyCN), simulación e interdicto de adquirir (art.\ 609 CPCCN).
}
\cqrow{¿Cómo se extingue la anotación de litis?}{No se rige por el art.\ 207 CPCCN. Se extingue por cumplimiento de la sentencia, por desestimación del reclamo (con comunicación al registro) o por modos anormales de terminación del proceso (caducidad de instancia, desistimiento o transacción). \\




%==============================================================
}
}{La anotación de litis opera sobre bienes registrables como mecanismo de publicidad que neutraliza la buena fe del tercero adquirente. Exige juicio iniciado cuya pretensión implique modificar un asiento registral, es ordenada de oficio en usucapión, es una de las primeras medidas en expropiación y se extingue por cumplimiento de la sentencia, desestimación o modos anormales de terminación.
}

\section{La prohibición de contratar}
%==============================================================

\begin{note}
La prohibición de innovar (art.\ 230 CPCCN) se analiza en el capítulo~\ref{cap:generico}, donde es objeto de una explicación más detenida. Aquí se desarrolla directamente la prohibición de contratar (art.\ 231 CPCCN).
\end{note}

\subsection{Concepto y procedencia (art.\ 231 CPCCN)}

Es una cautela que en determinados casos resulta procedente, orientada a evitar que, con relación a un bien que es objeto del proceso o que se encuentra embargado en el proceso, se realicen actos jurídicos que puedan disminuir o condicionar la cautela o el trámite del cumplimiento de la sentencia.

\begin{example}{Ejecución hipotecaria}
En el marco de un juicio de ejecución hipotecaria, además del embargo del bien, suele adoptarse la prohibición de contratar para evitar que el deudor hipotecario introduzca personas al inmueble mediante un contrato de locación o comodato, dificultando el avance de las medidas de cumplimiento de la sentencia en la etapa de subasta. La prohibición de contratar puede notificarse también al locatario existente, para hacerle saber que no podrá prorrogarse ni ampliarse el contrato próximo a vencer.
\end{example}

\subsection{Notificación y publicidad}

La prohibición de contratar debe:
\begin{itemize}
    \item Si se trata de un bien registrable, asentarse en el registro para su publicidad \textit{erga omnes}.
    \item Notificarse al titular de la cosa o propietario.
    \item Poder notificarse también a cualquier tercero que se sospeche que pueda ser la persona que contrate con el propietario.
\end{itemize}

\subsection{Caducidad especial (art.\ 231 CPCCN)}

\begin{important}
La prohibición de contratar tiene una caducidad especial, más breve que la del art.\ 207: \textbf{cinco (5) días}. Este plazo corre desde el \textbf{dictado} de la medida (no desde su traba). Solo se aplica si la prohibición se adopta \textbf{antes del inicio del proceso}. Si se adopta durante el curso del proceso ---incluso en la etapa conclusional--- no rige este plazo especial.
\end{important}

\cornelltitle{la prohibición de contratar}
\cornellblocksum{
\cqrow{¿Para qué sirve la prohibición de contratar?}{Para evitar que, respecto de un bien objeto del proceso o embargado en él, se realicen actos jurídicos que puedan disminuir o condicionar la cautela o el cumplimiento de la sentencia. Por ejemplo, en una ejecución hipotecaria evita que el deudor introduzca personas al inmueble mediante locación o comodato.
}
\cqrow{¿Cómo se notifica y se publicita?}{Se asienta en el registro si el bien es registrable (publicidad erga omnes), se notifica al titular o propietario y puede notificarse a cualquier tercero que se sospeche pueda contratar con él, como el locatario existente, a quien se le hace saber que no podrá prorrogarse ni ampliarse el contrato.
}
\cqrow{¿Qué caducidad especial tiene y cuándo se aplica?}{Cinco días, contados desde el dictado de la medida (no desde su traba), y solo si se adopta antes del inicio del proceso. Si se adopta durante el proceso, incluso en la etapa conclusional, no rige ese plazo especial. \\




%==============================================================
}
}{La prohibición de contratar protege la integridad del bien durante el proceso evitando actos jurídicos que condicionen la cautela o la ejecución de la sentencia. Debe publicitarse y notificarse, y tiene una caducidad especial de cinco días desde su dictado cuando es previa al inicio del proceso.
}

\chapter[Protección de personas y familia]{Protección de personas y medidas cautelares en materia de familia}\label{cap:proteccion}

%==============================================================

\section{La protección de personas y el artículo 234 CPCCN}

El análisis de esta medida en el ámbito procesal obedece a una tradición antigua, que incluía en el antiguo art.\ 234 del código procesal una serie de potestades cautelares del juez referidas a la protección de personas. Uno de sus incisos daba amplias facultades al juez para tomar medidas de cambio de guarda y de disposición respecto de niños, en hipótesis en que se observara o se sospechara que estaban sufriendo algún tipo de maltrato, violencia o falta de cuidado.

\begin{important}
\textbf{Ley 26.061 --- Sistema Nacional de Protección Integral de Niñas, Niños y Adolescentes.} Esta ley \textbf{derogó} el inciso del art.\ 234 CPCCN que otorgaba facultades cautelares a los jueces respecto de niños. Como consecuencia, sacó de la órbita del Poder Judicial toda la gestión de cuidado y protección de situaciones de maltrato, abuso o descuido de niñas, niños y adolescentes, y las colocó en el ámbito de la \textbf{administración pública}, a través de las Defensorías de Niños, Niñas y Adolescentes.

Los llamados a intervenir ahora son la escuela, el hospital, los centros de salud, la policía y cualquier organismo que detecte una situación de riesgo.
\end{important}

Lo que sí procede y debe ser judicializado, de acuerdo con los parámetros de la ley, es el \textbf{control de legalidad} de las medidas que la administración ---las defensorías--- adopta respecto del cuidado o el abrigo de situaciones de riesgo de los niños. En tanto el órgano de la administración adopte medidas excepcionales de separación del niño de su ámbito familiar primario, estas deben ser supervisadas y controladas en cuanto a su legalidad, para que no se prolonguen en el tiempo más allá de lo necesario.

\textbf{Estado actual del art.\ 234 CPCCN.} El grueso de la actividad protectora que poseía con anterioridad ha quedado transferida a la administración pública a través del Sistema Nacional de Protección Integral de Niños, Niñas y Adolescentes (Ley 26.061). El art.\ 234 solo conserva una función tuitiva de los jueces en materia de \textbf{mayores de edad con restricciones a la capacidad por problemáticas de salud mental}.

\section{Potestades cautelares en materia de violencia}

\subsection{Ley 24.417 --- Ley de Violencia Familiar}

Establece una serie de medidas cautelares que los jueces pueden adoptar para erradicar cualquier tipo de situación de agresión padecida por personas mayores de edad (hombres, mujeres y niños involucrados con personas mayores de edad). % TODO: contenido incompleto o ambiguo en las notas originales (redacción del ámbito de aplicación)
Incluye medidas de distanciamiento, prohibición de acercamiento, fijación provisional de alimentos, regímenes de comunicación y cuidado personal de los niños.

\subsection{Ley 26.485 --- Ley de Protección de Violencia contra la Mujer}

Refuerza y amplía las medidas de la Ley 24.417 en materia de perspectiva de género. Otorga al sistema judicial potestades cautelares de amplísima gama, que permiten tomar medidas de distanciamiento, prohibición de acercamiento, prestación alimentaria, fijación de regímenes de comunicación entre padres e hijos, establecimiento del cuidado personal de los niños y medidas patrimoniales vinculadas a la distribución de bienes.

\section{Medidas cautelares en el proceso de familia}

En materia de derecho de familia, la exigencia de los presupuestos de verosimilitud del derecho, peligro en la demora y contracautela está superada por el vínculo familiar que existe entre las personas involucradas.

\begin{table}[H]
\centering
\small
\begin{tabularx}{\textwidth}{>{\raggedright\arraybackslash}p{3.2cm} >{\raggedright\arraybackslash}X}
\toprule
\textbf{Presupuesto} & \textbf{Aplicación en procesos de familia} \\
\midrule
\textbf{Verosimilitud del derecho} & Surge del propio vínculo familiar (conyugal, de pareja, filial). No es necesario evaluar más que la existencia del vínculo como presupuesto que satisface este recaudo. \\
\addlinespace
\textbf{Peligro en la demora} & Está dado \textit{in re ipsa} (en la propia naturaleza de las cosas). La urgencia de cautelar el cuidado de los hijos, los alimentos provisorios, la vivienda y el régimen de comunicación es evidente por los derechos fundamentales comprometidos. \\
\addlinespace
\textbf{Contracautela} & Prácticamente en la generalidad de los casos no existe contracautela exigible. La caución juratoria está implícita y ni siquiera se pide, ya que se torna innecesaria. \\
\bottomrule
\end{tabularx}
\caption{Presupuestos cautelares en los procesos de familia}
\end{table}

\subsection{Las cautelares de familia en el Código Civil y Comercial de la Nación}

Las cautelares de familia se encuentran en su gran mayoría insertas en el código de fondo, el Código Civil y Comercial de la Nación. Esto no es algo nuevo ni nació con la reforma del año 2015, sino que ya venía con el código anterior.

El legislador siempre ha tenido en cuenta que ciertos aspectos del conflicto familiar compete al legislador nacional garantizarlos, y no ha querido dejarlos librados a lo que pueda establecer el legislador local provincial. Por eso plasmó en el código de fondo normas eminentemente operativas vinculadas con medidas cautelares propias del derecho de familia, que abarcan:
\begin{itemize}
    \item La atribución del hogar familiar.
    \item La garantía de subsistencia alimentaria provisoria mientras se define en un proceso de conocimiento las prestaciones alimentarias para los hijos menores y, eventualmente, entre cónyuges.
    \item La definición de la manera en que los padres ejercerán el cuidado personal de los hijos.
    \item El régimen de comunicación de los hijos con sus padres.
    \item Las medidas de carácter patrimonial vinculadas con la distribución de bienes como consecuencia de la extinción del régimen matrimonial.
\end{itemize}

\begin{note}
\textbf{Lectura recomendada.} Se recomienda como lectura obligatoria el artículo publicado en el sitio web \url{www.jorgearojas.com.ar}: «Las medidas cautelares en el Código Civil y Comercial de la Nación», del Dr.\ Jorge A.\ Rojas. Analiza las medidas cautelares nominadas e innominadas del CCyCN, las heredadas del código anterior, las absolutamente novedosas y aspectos de tutelas anticipatorias contempladas en el Código Civil y Comercial.
\end{note}

\cornelltitle{protección de personas y familia}
\cornellblocksum{
\cqrow{¿Qué cambió respecto de la protección de niños con la Ley 26.061?}{La ley derogó el inciso del art.\ 234 CPCCN que daba facultades cautelares a los jueces sobre niños. La gestión de protección ante maltrato, abuso o descuido pasó a la administración pública (Defensorías de Niños, Niñas y Adolescentes), con intervención de escuelas, hospitales, centros de salud, policía y otros organismos. Al Poder Judicial le queda el control de legalidad de las medidas excepcionales de separación del niño de su ámbito familiar, para que no se prolonguen más de lo necesario.
}
\cqrow{¿Qué función conserva hoy el art.\ 234 CPCCN?}{Solo una función tuitiva de los jueces respecto de mayores de edad con restricciones a la capacidad por problemáticas de salud mental.
}
\cqrow{¿Qué medidas permiten las leyes 24.417 y 26.485?}{Otorgan al Poder Judicial potestades cautelares amplias en materia de violencia familiar y de género: distanciamiento, prohibición de acercamiento, alimentos provisorios, regímenes de comunicación y cuidado personal de los niños y, en la Ley 26.485, también medidas patrimoniales vinculadas a la distribución de bienes. La Ley 26.485 refuerza y amplía la Ley 24.417 con perspectiva de género.
}
\cqrow{¿Qué particularidades tienen los presupuestos de las cautelares en familia?}{La verosimilitud surge del propio vínculo familiar; el peligro en la demora está dado \textit{in re ipsa} por los derechos fundamentales comprometidos; y en general no hay contracautela exigible, pues la caución juratoria está implícita.
}
\cqrow{¿Dónde se regulan las cautelares de familia y qué abarcan?}{En su gran mayoría en el Código Civil y Comercial de la Nación (ya ocurría con el código anterior), con normas operativas sobre atribución del hogar familiar, alimentos provisorios, cuidado personal de los hijos, régimen de comunicación y medidas patrimoniales tras la extinción del régimen matrimonial. \\




%--------------------------------------------------------------
% BIBLIOGRAFÍA (solo la referencia que figura en el apunte)
%--------------------------------------------------------------
}
}{Con la Ley 26.061 la adopción de medidas de protección de niños se trasladó a la administración pública, reservando a los jueces el control de legalidad. Las leyes 24.417 y 26.485 otorgan al Poder Judicial amplias potestades cautelares en violencia familiar y de género. En los procesos de familia los presupuestos cautelares se satisfacen por el propio vínculo, por el peligro en la demora \textit{in re ipsa} y por la ausencia de contracautela exigible.
}

\chapter*{Bibliografía}
\addcontentsline{toc}{chapter}{Bibliografía}
\markboth{Bibliografía}{Bibliografía}
\begin{itemize}
    \item Rojas y Falcón. \textit{Cómo se hace un alegato}. Ed. Abeledo Perrot.
    \item Rojas, Jorge A. «Las medidas cautelares en el Código Civil y Comercial de la Nación». Disponible en: \url{www.jorgearojas.com.ar}.
\end{itemize}

\end{document}
